% Generated by scripts/build_readers.py; change the reviewed sources and rebuild.
\def\EditionMode{1}
\def\EditionLabel{Corrected reader with apparatus}
\def\EditionStatus{WORKING EDITION: 132 of 706 source pages reviewed and included}
\documentclass[10pt,oneside]{book}
\usepackage{fontspec}
\setmainfont{Latin Modern Roman}
\setsansfont{Latin Modern Sans}
\usepackage{amsmath,amssymb,mathtools}
\usepackage{unicode-math}
\setmathfont{Latin Modern Math}
\usepackage{xcolor,tikz-cd,lua-ul,marginnote,pdflscape}
\usetikzlibrary{arrows.meta,decorations.pathmorphing,calc,shapes.geometric}
\tikzcdset{ampersand replacement=\&}
\tikzcdset{every label/.append style={text=black}}
\usepackage[a4paper,left=18mm,right=45mm,top=18mm,bottom=22mm,marginparwidth=28mm,marginparsep=7mm]{geometry}
\usepackage[unicode,colorlinks=true,urlcolor=black,linkcolor=black]{hyperref}
\usepackage{bookmark}
\setlength{\parindent}{0pt}
\setlength{\parskip}{4pt plus 1pt minus 1pt}
\setlength{\emergencystretch}{1.5em}
\raggedbottom
\definecolor{PSHand}{RGB}{25,58,132}
\definecolor{PSDelete}{RGB}{146,37,37}
\definecolor{PSTypedOld}{RGB}{95,95,95}
\ifnum\EditionMode=2\colorlet{blue}{black}\else\colorlet{blue}{PSHand}\fi
\RenewDocumentCommand{\textcircled}{m}{\tikz[baseline=(PSCircle.base)]\node[draw,circle,inner sep=.5pt,line width=.4pt](PSCircle){#1};}
\newif\ifPSInApparatus
\newcounter{psnote}
\newcounter{psvariant}
\newcounter{pslinkmark}
\newif\ifPSInLink
\newsavebox{\PSNoteBox}
\newdimen\PSNoteLimit
{\footnotesize\global\PSNoteLimit=3\baselineskip}
\newcommand{\CurrentScanURL}{}
\newcommand{\CurrentKey}{}
\newcommand{\CurrentScan}{}
\newcommand{\pssource}[3]{\psflushmargins\clearpage\gdef\CurrentKey{#1}\gdef\CurrentScan{#2}\gdef\CurrentScanURL{#3}\setcounter{psvariant}{0}\typeout{PS-SOURCE: #1}}
\newcommand{\psdivision}[2]{\pdfbookmark[0]{#1}{#2}\par\noindent{\large\bfseries #1\par}\medskip}
\newcommand{\pslinkmarker}[2]{\stepcounter{pslinkmark}\latelua{texio.write_nl('log','PS-LINK-SHIP|#1|\CurrentKey|\arabic{pslinkmark}|#2')}}
\newcommand{\pslink}[2]{\ifPSInLink\pslinkmarker{#1}{nested}#2\else\begingroup\PSInLinktrue\href{\CurrentScanURL}{\pslinkmarker{#1}{outer}#2}\endgroup\fi}
\newcommand{\scanpage}[1]{\ifnum\EditionMode=2\leavevmode\marginpar{\scriptsize\sffamily\pslink{scanpage}{scan #1}}\else\par\noindent{\small\sffamily\pslink{scanpage}{[scan #1]}}\par\medskip\fi}
\newcommand{\mspage}[1]{\ifnum\EditionMode=2\leavevmode\marginpar{\scriptsize\sffamily p.\,#1}\else{\small\sffamily[p.\,#1]}\fi}
\newcommand{\hand}[1]{{\ifnum\EditionMode=2\color{black}\else\color{PSHand}\fi #1}}
\let\over\relax
\NewDocumentCommand{\over}{O{} m m}{{\textcolor{PSTypedOld}{#2}\ensuremath{^{\!\nearrow}}\if\relax\detokenize{#1}\relax\hand{#3}\else #3\fi}}
\NewDocumentCommand{\inserted}{O{} m}{{\if\relax\detokenize{#1}\relax\hand{\ensuremath{\langle}#2\ensuremath{\rangle}}\else\ensuremath{\langle}#2\ensuremath{\rangle}\fi}}
\newcommand{\deleted}[1]{{\color{PSDelete}\ifmmode\text{[del.: }#1\text{]}\else[\textsf{del.:} #1]\fi}}
\newcommand{\xout}[1]{{\color{PSTypedOld}\ifmmode\text{[x: }#1\text{]}\else[\textsf{x:} #1]\fi}}
\let\PSMathUnderline\underline
\RenewDocumentCommand{\underline}{+m}{\ifmmode\PSMathUnderline{#1}\else\underLine{#1}\fi}
\NewDocumentCommand{\uline}{+m}{\underline{#1}}
\NewDocumentCommand{\uuline}{+m}{\ifmmode\PSMathUnderline{\PSMathUnderline{#1}}\else\underLine{\underLine{#1}}\fi}
\newunderlinetype\PSWaveLine{\leaders\hbox{\lower.7ex\hbox{\tikz[x=1pt,y=1pt,baseline=0pt]{\draw[line width=.3pt] (0,0) .. controls (1,1) and (2,1) .. (3,0) .. controls (4,-1) and (5,-1) .. (6,0);}}}}
\NewDocumentCommand{\uwave}{+m}{{\PSWaveLine #1}}
\newcommand{\unsure}[1]{\pslink{unsure}{[#1\ifmmode\,?\else\,?\fi]}}
\newcommand{\illegible}[1]{\ifmmode\text{\pslink{illegible}{[illeg.\ #1]}}\else\pslink{illegible}{[illeg.\ #1]}\fi}
\NewDocumentCommand{\diagram}{+m}{\ifmmode\vcenter{\hbox{#1}}\else\par\medskip{\centering #1\par}\medskip\psflushmargins\fi}
\newenvironment{diagramblock}{\par\begin{center}}{\end{center}\par}
\newcommand{\PSMathNotes}{}
\makeatletter
\NewDocumentCommand{\psplacednote}{m +m}{%
      \sbox{\PSNoteBox}{\parbox{\marginparwidth}{\footnotesize\raggedright\hand{\textsf{m#1.} #2}}}%
      \ifdim\dimexpr\ht\PSNoteBox+\dp\PSNoteBox\relax<\PSNoteLimit
        \leavevmode\marginpar{\usebox{\PSNoteBox}}%
      \else
        \par\nobreak\smallskip\hrule\nobreak\smallskip
        {\footnotesize\textsf{Marginal note m#1.}\quad\hand{#2}\par}
        \nobreak\smallskip\hrule\smallskip
      \fi}
\newcommand{\psflushmargins}{\ifmmode\else\ifinner\else\ifx\PSMathNotes\@empty\else\let\PSTempNotes\PSMathNotes\global\let\PSMathNotes\@empty\PSTempNotes\fi\fi\fi}
\NewDocumentCommand{\margin}{+m}{%
  \ifPSInApparatus\ifmmode\hand{\text{[marg.: }#1\text{]}}\else\hand{\textsf{[marg.: }#1\textsf{]}}\fi
  \else\stepcounter{psnote}%
    \ifmmode
      \text{\textsuperscript{m\arabic{psnote}}}%
      \begingroup\edef\PSQueue{\endgroup\noexpand\g@addto@macro\noexpand\PSMathNotes{\noexpand\psplacednote{\arabic{psnote}}{\unexpanded{$#1$}}}}\PSQueue
    \else
      \textsuperscript{m\arabic{psnote}}\psplacednote{\arabic{psnote}}{#1}%
    \fi
  \fi}
\AtBeginDocument{\everymath\expandafter{\the\everymath\aftergroup\psflushmargins}\everydisplay\expandafter{\the\everydisplay\aftergroup\psflushmargins}}
\makeatother
\newcommand{\psnumber}[2]{\expandafter\gdef\csname psnum@#1\endcsname{#2}}
\NewDocumentCommand{\var}{+m +m m m}{\ifnum\EditionMode=0 #1\else\ifnum\EditionMode=1\ifmmode\text{\textsuperscript{\tiny\csname psnum@#4\endcsname}}\else\textsuperscript{\tiny\csname psnum@#4\endcsname}\fi\fi #2\fi}
\newcommand{\psapparatusstart}{\par\bigskip\hrule\smallskip\PSInApparatustrue\begingroup\footnotesize\raggedright\textsf{Apparatus for scan \CurrentScan}\par}
\newcommand{\psapparatusend}{\par\endgroup\PSInApparatusfalse}
\NewDocumentCommand{\psentry}{m m m +m +m m}{\par\noindent\textbf{#2. #3.}\quad\textit{Typescript:} #4\quad\textit{Corrected:} \if\relax\detokenize{#5}\relax\textnormal{[omitted]}\else #5\fi\quad #6\par}
\newcommand{\psfrontmatter}{%
\hypersetup{pageanchor=false}\begin{titlepage}\raggedright
{\Large Alexander Grothendieck\par}\vspace{12mm}
{\Huge Pursuing Stacks\par}\medskip{\LARGE 1983\par}\vspace{12mm}
{\Large\EditionLabel\par}\vspace{8mm}
{\large\EditionStatus\par}\vfill
Université de Montpellier, Grothendieck archive, cotes 134-2--134-8.\par
Including the 26 source pages of R. Brown's translation at cote 134-2, sites 115--143.\par
\end{titlepage}\hypersetup{pageanchor=true}
\chapter*{Editorial statement}
AI-transcribed by Claude Opus 5.5 (Anthropic, effort high); checked against the scans, and this edition built, by GPT-6 Astra (OpenAI, ChatGPT Pro mode); prompts and records by Claude (Anthropic).

The witness consulted and not reproduced is the LaTeX edition of Mateo Carmona with the collaboration of Ulrik Buchholtz (arXiv 2111.01000 v2, 2021), extended from the 2015 edition by ``the scrivener'', CC0 1.0.

The translation in the appendix is R. Brown's, prepared with assistance from L. Breen and J-L. Loday, as his note says; the rights in it belong to its translator or his heirs.

The rights in the original work belong to its author or the author's representatives or heirs. Where we hold rights, in the AI-produced transcription, translation, corrections and typesetting, they are dedicated to the public domain under CC0 1.0. This edition is published for educational, non-profit purposes, and we claim no rights we do not hold. If you hold rights in this work and object to its publication here, please open a GitHub issue and we will take it down.

\section*{How to read this form}
\ifnum\EditionMode=0
The diplomatic form preserves typed and handwritten wording, cancellations, insertions, overwritten readings, underlining, doubts, and original page numbers. Blue marks handwriting; red labels handwritten cancellations; gray labels typewriter cancellations and underlying overwritten text. Angle brackets mark insertions; a superscript arrow separates an overwritten reading from its replacement. Typed insertions and typed overwriting remain plain. The disposition of long marginal notes is identified explicitly.
\else
The corrected text applies the author's insertions, replacements, cancellations, and recorded transpositions. It retains his English and notation, including consistent forms such as ``standart'', ``developped'', ``it's'' for ``its'', and ``A la''. An isolated slip is corrected only when specifically recorded and justified. Mathematical anomalies are preserved and discussed in the guide.
\ifnum\EditionMode=1
Blue continues to identify handwriting. A numbered call identifies each maximal editorial operation, with its full diplomatic and corrected readings and reason in the apparatus following the source page. Internal layers of a wholly replaced or deleted passage are preserved within that entry rather than duplicated as separate entries. These entries are generated from the same variants that govern the running text.
\else
This reading form omits the apparatus, cancelled material, superseded readings, insertion brackets, and color distinctions. It retains the resulting text, marginal notes, underlining, and unresolved doubts. Scan and source-page references appear in the margin.
\fi\fi

Every scan mark and remaining doubt links to its source page in the archive viewer. Doubts are editorial reading marks, not emendations. The source-page breaks identify the archival sequence; a source page may occupy more than one printed page. A partial working edition is explicitly labelled and contains only reviewed source pages.
\clearpage}

\begin{document}
\psfrontmatter
\pssource{ps_134-5_001}{134-5/1}{https://grothendieck.umontpellier.fr/web/viewer.html?file=../134-5.pdf\#page=2}
\psdivision{Cote 134-5}{division-134-5}
\psnumber{ps_134-5_001-v0001}{1}
\psnumber{ps_134-5_001-v0002}{2}
\psnumber{ps_134-5_001-v0003}{3}
\psnumber{ps_134-5_001-v0004}{4}
\psnumber{ps_134-5_001-v0005}{5}
\psnumber{ps_134-5_001-v0006}{6}
\psnumber{ps_134-5_001-v0007}{7}
\psnumber{ps_134-5_001-v0008}{8}
\psnumber{ps_134-5_001-v0009}{9}
\psnumber{ps_134-5_001-v0010}{10}
\psnumber{ps_134-5_001-v0011}{11}
\psnumber{ps_134-5_001-v0012}{12}
\psnumber{ps_134-5_001-v0013}{13}
\psnumber{ps_134-5_001-v0014}{14}
\psnumber{ps_134-5_001-v0015}{15}
%% BEGIN SOURCE ps_134-5_001
% Pursuing Stacks (cotes 134-2 to 134-8): archive cote 134-5, site page 1 (PDF page 2), typescript page 188.
% typescript. Draft compared: ps4.tex:1-63.
% AI-transcribed by Claude Opus 5.5 (claude-opus-5-5, adaptive thinking, effort high), Message Batches API, batch ps2 (msgbatch_01NVwr1Yuo8qJf5rykeGBXPK), 2026-11-12; reviewed against both supplied scan halves; see CHECK_DECISIONS.tsv.
\scanpage{134-5/1}
\mspage{188}
11.4.

6\var{\over{4}{7}}{\hand{7}}{writing over}{ps_134-5_001-v0001}.    A little more pondering and scribbling finally seems to show
that the \var{\deleted{real}}{}{deletion}{ps_134-5_001-v0002} key for an understanding of modelizers is'nt really the
notion of contractibility,  but rather the notion of aspheric objects
(besides,  of course,  the notion of weak equivalence).  At the same time
it appears that the notion of an \uline{aspheric} \uline{map} in (Cat),  more specifical-%
ly of a $\underline{W}$-aspheric``map''\var{\xout{xxxx}}{}{typewriter deletion}{ps_134-5_001-v0003} (i.e. functor between small categories) is
a lot more important than \var{\inserted{being just}}{\hand{being just}}{insertion}{ps_134-5_001-v0004} a highly expedient technical convenience,  as
it has been so far  -  it is indeed one of the basic notions of the theo-%
ry of modelizers we got into.  As a matter of fact,  I should have known
this for a number of weaks already, \var{\inserted{ever}}{\hand{ever}}{insertion}{ps_134-5_001-v0005} since I did some scribbling about
the plausible notion of ``morphism'' between test-categories (as well as
their weak and strong variants),  and readily convinced myself that the
natural ``morphisms'' here were nothing else but the aspheric functors bet-%
ween those categories.  I kind of forgot about this,  as it did'nt seem
too urgent to start moving \var{\inserted{around}}{\hand{around}}{insertion}{ps_134-5_001-v0006} the test category I was working with.  If I \var{\xout{xx}}{}{typewriter deletion}{ps_134-5_001-v0007}
had been a little more systematic in grinding through the usual functori-%
alities,  as soon as a significant notion (such as the various test noti-%
ons)  appears,  I presumably would have hit upon the crucial point about
modelizers and so-called ``asphericity structures'' a lot sooner,  without
going through the long-winded detour of homotopy structures\hand{,} and \var{\inserted{the still}}{\hand{the still}}{insertion}{ps_134-5_001-v0008} extreme-%
ly special types of test functors suggested by contractibility assumpti-%
ons.  However \var{\xout{it}}{}{typewriter deletion}{ps_134-5_001-v0009} I believe that most of the work I went through,  \var{\xout{xxx}}{}{typewriter deletion}{ps_134-5_001-v0010}
although irrelevant for the ``asphericity story'' itself,  will still be
useful,  especially when it comes to pinpointing the so-called ``canonic\unsure{\hand{al}}
modelizers,  whose modelizing structure is intrinsicly determined by the
category structure.

First thing now which we've to do is to have a closer look at the
meaning of asphericity for \var{\inserted{a}}{\hand{a}}{insertion}{ps_134-5_001-v0011} functor\var{\xout{s}}{}{typewriter deletion}{ps_134-5_001-v0012} between  small categories.  There is
no reason whatever to \var{\xout{restrict to}}{}{typewriter deletion}{ps_134-5_001-v0013} make any restrictions on these catego-%
ries besides smallness  (namely the cardinals of the sets of objects and
arrows being in the ``universe'' we are implicitly working in,  or even
more stringently still,  A and B to be in (Cat)  namely to be objects
of that universe).  Thus we will not assume A and B to be test categories
or the like.  \var{\xout{However,}}{}{typewriter deletion}{ps_134-5_001-v0014} We will be led to consider,  for \uline{any} small cate-%
gory A ,  the localization of A\^{} with respect to $\underline{W}$-equivalences,  which
I'll denote by $(\mathrm{Hot}^{\hand{\underline{W}}}_A)$ \hand{or simply $(\mathrm{Hot}_A)$:}
\[ \hand{(1)} \qquad \hand{(\mathrm{Hot}^{\underline{W}}_A) or} \ (\mathrm{Hot}^{\var{\deleted{\illegible{1}}}{}{deletion}{ps_134-5_001-v0015}}_A) \ = \ \underline{W}_A^{-1}A\hat{} \quad . \]

\margin{1}

%% END SOURCE ps_134-5_001
\ifnum\EditionMode=1
\psapparatusstart
\psentry{ps_134-5_001-v0001}{1}{writing over}{\over{4}{7}}{\hand{7}}{The superposed handwritten reading replaces the underlying reading.}
\psentry{ps_134-5_001-v0002}{2}{deletion}{\deleted{real}}{}{The author cancels this material by hand; it leaves the corrected running text.}
\psentry{ps_134-5_001-v0003}{3}{typewriter deletion}{\xout{xxxx}}{}{The typewriter strikes out this material; it leaves the corrected running text.}
\psentry{ps_134-5_001-v0004}{4}{insertion}{\inserted{being just}}{\hand{being just}}{The handwritten insertion is incorporated at its indicated position; insertion brackets leave the running text.}
\psentry{ps_134-5_001-v0005}{5}{insertion}{\inserted{ever}}{\hand{ever}}{The handwritten insertion is incorporated at its indicated position; insertion brackets leave the running text.}
\psentry{ps_134-5_001-v0006}{6}{insertion}{\inserted{around}}{\hand{around}}{The handwritten insertion is incorporated at its indicated position; insertion brackets leave the running text.}
\psentry{ps_134-5_001-v0007}{7}{typewriter deletion}{\xout{xx}}{}{The typewriter strikes out this material; it leaves the corrected running text.}
\psentry{ps_134-5_001-v0008}{8}{insertion}{\inserted{the still}}{\hand{the still}}{The handwritten insertion is incorporated at its indicated position; insertion brackets leave the running text.}
\psentry{ps_134-5_001-v0009}{9}{typewriter deletion}{\xout{it}}{}{The typewriter strikes out this material; it leaves the corrected running text.}
\psentry{ps_134-5_001-v0010}{10}{typewriter deletion}{\xout{xxx}}{}{The typewriter strikes out this material; it leaves the corrected running text.}
\psentry{ps_134-5_001-v0011}{11}{insertion}{\inserted{a}}{\hand{a}}{The handwritten insertion is incorporated at its indicated position; insertion brackets leave the running text.}
\psentry{ps_134-5_001-v0012}{12}{typewriter deletion}{\xout{s}}{}{The typewriter strikes out this material; it leaves the corrected running text.}
\psentry{ps_134-5_001-v0013}{13}{typewriter deletion}{\xout{restrict to}}{}{The typewriter strikes out this material; it leaves the corrected running text.}
\psentry{ps_134-5_001-v0014}{14}{typewriter deletion}{\xout{However,}}{}{The typewriter strikes out this material; it leaves the corrected running text.}
\psentry{ps_134-5_001-v0015}{15}{deletion}{$\deleted{\illegible{1}}$}{}{The author cancels this material by hand; it leaves the corrected running text.}
\psapparatusend
\fi
\pssource{ps_134-5_002}{134-5/2}{https://grothendieck.umontpellier.fr/web/viewer.html?file=../134-5.pdf\#page=3}
\psnumber{ps_134-5_002-v0001}{1}
\psnumber{ps_134-5_002-v0002}{2}
\psnumber{ps_134-5_002-v0003}{3}
\psnumber{ps_134-5_002-v0004}{4}
\psnumber{ps_134-5_002-v0005}{5}
\psnumber{ps_134-5_002-v0006}{6}
\psnumber{ps_134-5_002-v0007}{7}
\psnumber{ps_134-5_002-v0008}{8}
\psnumber{ps_134-5_002-v0009}{9}
\psnumber{ps_134-5_002-v0010}{10}
\psnumber{ps_134-5_002-v0011}{11}
\psnumber{ps_134-5_002-v0012}{12}
\psnumber{ps_134-5_002-v0013}{13}
\psnumber{ps_134-5_002-v0014}{14}
\psnumber{ps_134-5_002-v0015}{15}
\psnumber{ps_134-5_002-v0016}{16}
\psnumber{ps_134-5_002-v0017}{17}
\psnumber{ps_134-5_002-v0018}{18}
\psnumber{ps_134-5_002-v0019}{19}
%% BEGIN SOURCE ps_134-5_002
% Pursuing Stacks (cotes 134-2 to 134-8): archive cote 134-5, site page 2 (PDF page 3), typescript page 189.
% typescript. Draft compared: ps4.tex:47-121.
% AI-transcribed by Claude Opus 5.5 (claude-opus-5-5, adaptive thinking, effort high), Message Batches API, batch ps2 (msgbatch_01NVwr1Yuo8qJf5rykeGBXPK), 2026-11-12; reviewed against both supplied scan halves; see CHECK_DECISIONS.tsv.
\scanpage{134-5/2}
\mspage{189}

These categories, I suspect, are quite interesting in themselves, and
they merit to be understood. Thus One of Quillen's results asserts that
(at least for $\underline{W} = W_{(\mathrm{Cat})}$ = ordinary weak aquivalences, but presumably
his arguments \var{\inserted[typed]{will}}{will}{typed insertion}{ps_134-5_002-v0001} carry over to an arbitrary $\underline{W}$ ) in case A is a product
category \hand{$\Delta$}xA$_{o}$ , where A$_{o}$ is any small category and \hand{$\Delta$} the category of
standart, \var{\inserted[typed]{ordered}}{ordered}{typed insertion}{ps_134-5_002-v0002} simplices, then A\^{} is a closed model category admitting as
weak equivalences the set $\underline{W}_{A}$\hand{$^{(*)}$}; and hence \hand{(}Hot$_{A}$\hand{)} , the corresponding
\var{\margin{(*) 3.5. This, I quickly became aware, is a misrepresentation of Quillen's result - the ``weak equivalences'' he introduced are a lot stronger than $\underline{W}_{A}$. I'll have to come back upon this soon enough! \deleted{However, if A is \unsure{aspheric}, it is true that its product by any test category is again a test category, see below} X}}{\margin{(*) 3.5. This, I quickly became aware, is a misrepresentation of Quillen's result - the ``weak equivalences'' he introduced are a lot stronger than $\underline{W}_{A}$. I'll have to come back upon this soon enough!  X}}{marginal note revised}{ps_134-5_002-v0003}
``homotopy category'', admits familiar homotopy constructions, including
the two types of Dold-Puppe exact sequences, tied up with loop- and sus-%
pension functors. It is very hard to believe that this should be a spe-%
cial feature of the category \hand{$\Delta$} as \var{\inserted[typed]{the}}{the}{typed insertion}{ps_134-5_002-v0004} multiplying factor - surely any
test category or strict test category instead should do as well. As we'll
check below, the product of a \var{\inserted[typed]{local}}{local}{typed insertion}{ps_134-5_002-v0005} test category with any \var{\xout{($\underline{W}$-)aspheric}}{}{typewriter deletion}{ps_134-5_002-v0006} \var{\inserted[typed]{\deleted{non-empty}}}{}{typed insertion}{ps_134-5_002-v0007}
category A$_{o}$ is again a local test category, hence a test-category if
both factors are ($\underline{W}$-) aspheric. This suggests that maybe \var{\inserted[typed]{for}}{for}{typed insertion}{ps_134-5_002-v0008} any local test
category A, the corresponding A\^{} is a closed model category - but it
is'nt even clear yet if the same does'nt hold for \uline{any} small category A
whatever ! It's surely something worth looking at.

As we'll see presently, it is a tautology more or less that a
functor
\[ \text{(1)} \qquad i\colon A \hand{\longrightarrow} B \quad , \]
giving rise to a functor
\[ \text{(2)} \qquad i^{*}\colon B\hat{} \hand{\longrightarrow} A\hat{} \]
\hand{(}commuting to all types of direct and inverse limits,\hand{)} induces a functor
on the localizations
\[ \hand{\overline{\textcolor{black}{i^{*}}}} \ : \ (\mathrm{Hot}_{B}) \hand{\longrightarrow} (\mathrm{Hot})_{A}\text{\var{\xout{*}}{}{typewriter deletion}{ps_134-5_002-v0009}} \quad , \]
\var{\inserted{provided}}{\hand{provided}}{insertion}{ps_134-5_002-v0010} \var{\deleted{iff}}{}{deletion}{ps_134-5_002-v0011} i is \var{\inserted[typed]{$\underline{W}$-}}{$\underline{W}$-}{typed insertion}{ps_134-5_002-v0012} aspheric. Thus
\[ A \hand{\longmapsto} (\mathrm{Hot}_{A}) \]
can be viewed as a functor with respect to A, provided we take as ``mor-%
phisms'' between\allowbreak{}``objects''\allowbreak{}A the \uline{aspheric} functors \var{\inserted{only}}{\hand{only}}{insertion}{ps_134-5_002-v0013} - i.e. it is a functor
on the \var{\inserted[typed]{sub}}{sub}{typed insertion}{ps_134-5_002-v0014}category \var{\xout{(Cat)}}{}{typewriter deletion}{ps_134-5_002-v0015} \var{\xout{xxx}}{}{typewriter deletion}{ps_134-5_002-v0016} \var{\xout{asph}}{}{typewriter deletion}{ps_134-5_002-v0017} $(\mathrm{Cat})_{\underline{W}\text{-asph}}$ \var{\inserted[typed]{of (Cat),}}{of (Cat),}{typed insertion}{ps_134-5_002-v0018} having same objects as (Cat),
but with maps restricted to \var{\inserted{be}}{\hand{be}}{insertion}{ps_134-5_002-v0019} $\underline{W}$-aspheric ones.

We'll denote by
\[ \mathrm{Hot}(\underline{W}) = \underline{W}^{-1}(\mathrm{Cat}) \]
the homotopy category defined in terms of the basic localizer $\underline{W}$ . For
any small category A , we get a commutative diagram

\margin{2}

%% END SOURCE ps_134-5_002
\ifnum\EditionMode=1
\psapparatusstart
\psentry{ps_134-5_002-v0001}{1}{typed insertion}{\inserted[typed]{will}}{will}{The typed insertion is incorporated at its indicated position; insertion brackets leave the running text.}
\psentry{ps_134-5_002-v0002}{2}{typed insertion}{\inserted[typed]{ordered}}{ordered}{The typed insertion is incorporated at its indicated position; insertion brackets leave the running text.}
\psentry{ps_134-5_002-v0003}{3}{marginal note revised}{\margin{(*) 3.5. This, I quickly became aware, is a misrepresentation of Quillen's result - the ``weak equivalences'' he introduced are a lot stronger than $\underline{W}_{A}$. I'll have to come back upon this soon enough! \deleted{However, if A is \unsure{aspheric}, it is true that its product by any test category is again a test category, see below} X}}{\margin{(*) 3.5. This, I quickly became aware, is a misrepresentation of Quillen's result - the ``weak equivalences'' he introduced are a lot stronger than $\underline{W}_{A}$. I'll have to come back upon this soon enough!  X}}{The author's internal cancellations and replacements in this marginal note are applied at its source-line call.}
\psentry{ps_134-5_002-v0004}{4}{typed insertion}{\inserted[typed]{the}}{the}{The typed insertion is incorporated at its indicated position; insertion brackets leave the running text.}
\psentry{ps_134-5_002-v0005}{5}{typed insertion}{\inserted[typed]{local}}{local}{The typed insertion is incorporated at its indicated position; insertion brackets leave the running text.}
\psentry{ps_134-5_002-v0006}{6}{typewriter deletion}{\xout{($\underline{W}$-)aspheric}}{}{The typewriter strikes out this material; it leaves the corrected running text.}
\psentry{ps_134-5_002-v0007}{7}{typed insertion}{\inserted[typed]{\deleted{non-empty}}}{}{The typed insertion is incorporated at its indicated position; insertion brackets leave the running text.}
\psentry{ps_134-5_002-v0008}{8}{typed insertion}{\inserted[typed]{for}}{for}{The typed insertion is incorporated at its indicated position; insertion brackets leave the running text.}
\psentry{ps_134-5_002-v0009}{9}{typewriter deletion}{\xout{*}}{}{The typewriter strikes out this material; it leaves the corrected running text.}
\psentry{ps_134-5_002-v0010}{10}{insertion}{\inserted{provided}}{\hand{provided}}{The handwritten insertion is incorporated at its indicated position; insertion brackets leave the running text.}
\psentry{ps_134-5_002-v0011}{11}{deletion}{\deleted{iff}}{}{The author cancels this material by hand; it leaves the corrected running text.}
\psentry{ps_134-5_002-v0012}{12}{typed insertion}{\inserted[typed]{$\underline{W}$-}}{$\underline{W}$-}{The typed insertion is incorporated at its indicated position; insertion brackets leave the running text.}
\psentry{ps_134-5_002-v0013}{13}{insertion}{\inserted{only}}{\hand{only}}{The handwritten insertion is incorporated at its indicated position; insertion brackets leave the running text.}
\psentry{ps_134-5_002-v0014}{14}{typed insertion}{\inserted[typed]{sub}}{sub}{The typed insertion is incorporated at its indicated position; insertion brackets leave the running text.}
\psentry{ps_134-5_002-v0015}{15}{typewriter deletion}{\xout{(Cat)}}{}{The typewriter strikes out this material; it leaves the corrected running text.}
\psentry{ps_134-5_002-v0016}{16}{typewriter deletion}{\xout{xxx}}{}{The typewriter strikes out this material; it leaves the corrected running text.}
\psentry{ps_134-5_002-v0017}{17}{typewriter deletion}{\xout{asph}}{}{The typewriter strikes out this material; it leaves the corrected running text.}
\psentry{ps_134-5_002-v0018}{18}{typed insertion}{\inserted[typed]{of (Cat),}}{of (Cat),}{The typed insertion is incorporated at its indicated position; insertion brackets leave the running text.}
\psentry{ps_134-5_002-v0019}{19}{insertion}{\inserted{be}}{\hand{be}}{The handwritten insertion is incorporated at its indicated position; insertion brackets leave the running text.}
\psapparatusend
\fi
\pssource{ps_134-5_003}{134-5/3}{https://grothendieck.umontpellier.fr/web/viewer.html?file=../134-5.pdf\#page=4}
\psnumber{ps_134-5_003-v0001}{1}
\psnumber{ps_134-5_003-v0002}{2}
\psnumber{ps_134-5_003-v0003}{3}
\psnumber{ps_134-5_003-v0004}{4}
\psnumber{ps_134-5_003-v0005}{5}
\psnumber{ps_134-5_003-v0006}{6}
\psnumber{ps_134-5_003-v0007}{7}
\psnumber{ps_134-5_003-v0008}{8}
\psnumber{ps_134-5_003-v0009}{9}
\psnumber{ps_134-5_003-v0010}{10}
\psnumber{ps_134-5_003-v0011}{11}
\psnumber{ps_134-5_003-v0012}{12}
\psnumber{ps_134-5_003-v0013}{13}
\psnumber{ps_134-5_003-v0014}{14}
%% BEGIN SOURCE ps_134-5_003
% Pursuing Stacks (cotes 134-2 to 134-8): archive cote 134-5, site page 3 (PDF page 4), typescript page 190.
% typescript. Draft compared: ps4.tex:98-173.
% AI-transcribed by Claude Opus 5.5 (claude-opus-5-5, adaptive thinking, effort high), Message Batches API, batch ps2 (msgbatch_01NVwr1Yuo8qJf5rykeGBXPK), 2026-11-12; reviewed against both supplied scan halves; see CHECK_DECISIONS.tsv.
\scanpage{134-5/3}
\mspage{190}

% Hand-drawn arrows are marked blue; typed objects remain plain.
\diagram{\begin{tikzcd}
A\hat{} \arrow[color=blue, r, "i_A"] \arrow[color=blue, d, swap, "\hand{\gamma_A}"] \& (Cat) \arrow[color=blue, d, "\hand{\gamma_{\underline{W}} \text{ or } \gamma}"] \\
(Hot_A) \arrow[color=blue, r, "\overline{i}_A"] \& (H\unsure{\hand{\hat{o}}}t(\underline{W}))
\end{tikzcd}}
(3) \hfill ,

we denote by
\[ \hand{\varphi}_A= \hand{\gamma i_A = \overline{i}_A\gamma_A} : A\hat{} \hand{\longrightarrow} (Hot(\underline{W})) \]
the corresponding composition.

Coming back to the case of a functor i and corresponding $i^*$
((1) and (2)), the functors $i^*$ , $i_A$ , $i_B$ do \uline{not} give rise to a commu-%
tative triangle, but to a triangle \hand{\uline{\textcolor{black}{with commutation morphism}}} $\hand{\lambda}_i$ :
\diagram{\begin{tikzcd}[column sep=small]
B\hat{} \arrow[color=blue, rr, "i^*"] \arrow[color=blue, dr, swap, "i_B", ""{name=iB, right}] \& \& A\hat{} \arrow[color=blue, dl, "i_A", ""{name=iA, left}] \\
\& (Cat) \arrow[color=blue, bend right=10, from=iA, to=iB, swap, "\hand{\lambda_i}"] \&
\end{tikzcd}}
\hfill ,

i.e. for any F in $B\hat{}$ , we get a map \var{\xout{between}}{}{typewriter deletion}{ps_134-5_003-v0001}
\[ (4) \qquad \hand{\lambda}_i(F) : i_Ai^*(F) = A_{/i^{*}(F)} \hand{\longrightarrow} i_B(F) = B_{/F} \quad , \]
functorial in F for variable F. \unsure{T\hand{h}he|The} first hand side of (4), also
written simply $A_{/F}$ when there is no ambiguity for i , \var{\xout{ixxth}}{}{typewriter deletion}{ps_134-5_003-v0002} can be
interpreted as the category of pairs
\[ (a,p) \ , \quad a \in \mathrm{Ob}\, A \ , \quad p: i(a) \hand{\to} F \quad , \]
where p is a map in $B\hat{}$, B identified as usual to a full subcategory of
$B\hat{}$ (hence i(a) identified with an \unsure{object|obbject} of $B\hat{}$ ). The \var{\xout{functor}}{}{typewriter deletion}{ps_134-5_003-v0003} \var{\inserted[typed]{map}}{map}{typed insertion}{ps_134-5_003-v0004} $\hand{\lambda}_i(F)$
for fixed F is of course the \var{\xout{functxxxxxx}}{}{typewriter deletion}{ps_134-5_003-v0005} functor
\[ (a,p) \hand{\longmapsto} (i(a),p) \quad : \quad \hand{A_{/F} \longrightarrow B_{/F} .} \]

The topological significance of course is clear: interpreting i as
defining a"map" or morphism of the corresponding topoi $A\hat{}$ and $B\hat{}$, having
$i^*$ as inverse image functor, an object F of $B\hat{}$ gives rise to an \uline{induced}
\uline{topos} $B\hat{}_{/F} \hand{\overset{\approx}{\longrightarrow}} (B_{/F})\hat{}$ , and the restriction of the ``topos above'' $A\hat{}$
to the induced $B\hat{}_{/F}$ , or equivalently the result of base change
$B\hat{}_{/F} \hand{\to} B\hat{}$ , gives rise to the \var{\inserted[typed]{induced}}{induced}{typed insertion}{ps_134-5_003-v0006} morphism of topoi $A\hat{}_{/i^*(F)} \hand{\longrightarrow} B\hat{}_{/F}$ ,
represented precisely by the map $\hand{\lambda}_i(F)$ in (Cat).

The condition \var{\xout{taht}}{}{typewriter deletion}{ps_134-5_003-v0007} of $\underline{W}$-asphericity on i may be expressed in many-%
fold ways, as properties o\var{\over{f}{f}}{\hand{f}}{writing over}{ps_134-5_003-v0008} either one of the three aspects
\[ \hand{\lambda}_i \ , \quad i^* \ , \quad i_Ai^* \]
of the situation created by i , \var{\deleted{and}}{}{deletion}{ps_134-5_003-v0009} \var{\inserted{with respect to}}{\hand{with respect to}}{insertion}{ps_134-5_003-v0010} the localizing sets $\underline{W}$ , $\underline{W}_A$ , $\underline{W}_B$ \hand{,} \var{\inserted{or to the notion of aspheric objects.}}{\hand{or to the notion of aspheric objects.}}{insertion}{ps_134-5_003-v0011} \margin{$\checkmark$}
As $\underline{W}_A$ is defined in terms of $\underline{W}$ as just the inverse image of the latter
by $i_{\var{\over{A}{A}}{\hand{A}}{writing over}{ps_134-5_003-v0012}}$ , \var{\inserted{and the same for aspheric objects,}}{\hand{and the same for aspheric objects,}}{insertion}{ps_134-5_003-v0013} \hand{(}it turns out that each of the conditions we \var{\xout{ax}}{}{typewriter deletion}{ps_134-5_003-v0014}
\margin{3}

%% END SOURCE ps_134-5_003
\ifnum\EditionMode=1
\psapparatusstart
\psentry{ps_134-5_003-v0001}{1}{typewriter deletion}{\xout{between}}{}{The typewriter strikes out this material; it leaves the corrected running text.}
\psentry{ps_134-5_003-v0002}{2}{typewriter deletion}{\xout{ixxth}}{}{The typewriter strikes out this material; it leaves the corrected running text.}
\psentry{ps_134-5_003-v0003}{3}{typewriter deletion}{\xout{functor}}{}{The typewriter strikes out this material; it leaves the corrected running text.}
\psentry{ps_134-5_003-v0004}{4}{typed insertion}{\inserted[typed]{map}}{map}{The typed insertion is incorporated at its indicated position; insertion brackets leave the running text.}
\psentry{ps_134-5_003-v0005}{5}{typewriter deletion}{\xout{functxxxxxx}}{}{The typewriter strikes out this material; it leaves the corrected running text.}
\psentry{ps_134-5_003-v0006}{6}{typed insertion}{\inserted[typed]{induced}}{induced}{The typed insertion is incorporated at its indicated position; insertion brackets leave the running text.}
\psentry{ps_134-5_003-v0007}{7}{typewriter deletion}{\xout{taht}}{}{The typewriter strikes out this material; it leaves the corrected running text.}
\psentry{ps_134-5_003-v0008}{8}{writing over}{\over{f}{f}}{\hand{f}}{The superposed handwritten reading replaces the underlying reading.}
\psentry{ps_134-5_003-v0009}{9}{deletion}{\deleted{and}}{}{The author cancels this material by hand; it leaves the corrected running text.}
\psentry{ps_134-5_003-v0010}{10}{insertion}{\inserted{with respect to}}{\hand{with respect to}}{The handwritten insertion is incorporated at its indicated position; insertion brackets leave the running text.}
\psentry{ps_134-5_003-v0011}{11}{insertion}{\inserted{or to the notion of aspheric objects.}}{\hand{or to the notion of aspheric objects.}}{The handwritten insertion is incorporated at its indicated position; insertion brackets leave the running text.}
\psentry{ps_134-5_003-v0012}{12}{writing over}{$\over{A}{A}$}{$\hand{A}$}{The superposed handwritten reading replaces the underlying reading.}
\psentry{ps_134-5_003-v0013}{13}{insertion}{\inserted{and the same for aspheric objects,}}{\hand{and the same for aspheric objects,}}{The handwritten insertion is incorporated at its indicated position; insertion brackets leave the running text.}
\psentry{ps_134-5_003-v0014}{14}{typewriter deletion}{\xout{ax}}{}{The typewriter strikes out this material; it leaves the corrected running text.}
\psapparatusend
\fi
\pssource{ps_134-5_004}{134-5/4}{https://grothendieck.umontpellier.fr/web/viewer.html?file=../134-5.pdf\#page=5}
\psnumber{ps_134-5_004-v0001}{1}
\psnumber{ps_134-5_004-v0002}{2}
\psnumber{ps_134-5_004-v0003}{3}
\psnumber{ps_134-5_004-v0004}{4}
\psnumber{ps_134-5_004-v0005}{5}
\psnumber{ps_134-5_004-v0006}{6}
\psnumber{ps_134-5_004-v0007}{7}
\psnumber{ps_134-5_004-v0008}{8}
\psnumber{ps_134-5_004-v0009}{9}
\psnumber{ps_134-5_004-v0010}{10}
\psnumber{ps_134-5_004-v0011}{11}
\psnumber{ps_134-5_004-v0012}{12}
\psnumber{ps_134-5_004-v0013}{13}
\psnumber{ps_134-5_004-v0014}{14}
%% BEGIN SOURCE ps_134-5_004
% Pursuing Stacks (cotes 134-2 to 134-8): archive cote 134-5, site page 4 (PDF page 5), typescript page 191.
% typescript. Draft compared: ps4.tex:157-216.
% AI-transcribed by Claude Opus 5.5 (claude-opus-5-5, adaptive thinking, effort high), Message Batches API, batch ps2 (msgbatch_01NVwr1Yuo8qJf5rykeGBXPK), 2026-11-12; reviewed against both supplied scan halves; see CHECK_DECISIONS.tsv.
\scanpage{134-5/4}
\mspage{191}

are led to express on $i^*$ , can be formulated equivalently in terms of
the composition $i_Ai^*$ . I'll restrict to formulate these in terms of $i^*$
only, which will be the form most adapted to the use we are going to
make later of the notion of a $\underline{W}$-aspheric map, when introducing the so-%
called ``asphericity structures''\hand{/}and corresponding ``testing functors''.


27.5.

68. It has been over six weeks now that I did'nt write down any notes.
The reason for this is that I felt the story of asphericity s\var{\over{\unsure{r}}{t}}{\hand{t}}{writing over}{ps_134-5_004-v0001}ruct\var{\over{\unsure{i}}{u}}{\hand{u}}{writing over}{ps_134-5_004-v0002}res
and canonical modelizers was going to come now without any problem, al-%
most as a matter of routine \unsure{tow} write it down with some care - therefore
I started doing some scratchwork on a few questions which had been around
but kept in the background since the beginning, \var{\inserted{and}}{\hand{and}}{insertion}{ps_134-5_004-v0003} which were a lot less
clear in my mind. Some reflection \var{\deleted{on these}}{}{deletion}{ps_134-5_004-v0004} was needed anyhow, before it
would make much sense to start writing down anything on these. Finally it
took longer than expected, as usual\var{\unsure{\deleted{,}}}{}{deletion}{ps_134-5_004-v0005}\hand{-} partly because (as usual\var{\deleted{,}}{}{deletion}{ps_134-5_004-v0006} too\hand{!}) a
few surprises would turn up on my way. Also, I finally allowed myself to
become distracted by some reflection on the``Lego-Teichmüller construction
game'', \var{\deleted{especially}}{}{deletion}{ps_134-5_004-v0007} \var{\inserted{and pretty much so during}}{\hand{and pretty much so during}}{insertion}{ps_134-5_004-v0008} last week. The occasion was a series of \var{\xout{talks}}{}{typewriter deletion}{ps_134-5_004-v0009} unformal
talks Y. has been giving in Molino's seminar, on Thurston\unsure{'}s \var{\xout{work on the}}{}{typewriter deletion}{ps_134-5_004-v0010}
hyperbolic geometry game and his compactification of Teichmüller space.
Y. was getting interested again in mathematics after a five year's inter-%
ruption. He must have heard about my seminar last year on ``anabelian
algebraic geometry'' and the ``Teichmüller tower'', and suggested I might
drop in to get an idea about Thurston's work. This \var{\inserted{work}}{\hand{work}}{insertion}{ps_134-5_004-v0011} indeed appears as \var{\xout{very}}{}{typewriter deletion}{ps_134-5_004-v0012}
\margin{$\ast$} closely related in various respects to my sporadic reflections\hand{$\ast$} of the last
two years, just with a diametrically opposed emphasis - mine being on
the algebro-geometric and arithmetic aspects of ``moduli'' of algebraic
curves, his on hyperbolic riemanian geometry and the simply connected
\var{\inserted[typed]{transcendental}}{transcendental}{typed insertion}{ps_134-5_004-v0013} Teichmüller space\hand{s}\var{\unsure{\deleted{,}}}{}{deletion}{ps_134-5_004-v0014} \hand{(}rather than the algebraic modular varieties\hand{)}. The main
intersection appears to be interest in surface surgery and the relation of
this to the Teichmüller modular group. I took the occasion to try and
recollect about the Lego-Teichmüller game, which I had thought of last
year as a plausible,very concrete way for modelizing and visualizing the
whole tower of Teichmüller groupoids $T_{g,\hand{\nu}}$ and the main operations among
these, especially the ``cutting'' and``gluing'' operations. The very unfor-%
mal talk I gave was mainly intended for Y. as a matter of fact, and it

\margin{4}

%% END SOURCE ps_134-5_004
\ifnum\EditionMode=1
\psapparatusstart
\psentry{ps_134-5_004-v0001}{1}{writing over}{\over{\unsure{r}}{t}}{\hand{t}}{The superposed handwritten reading replaces the underlying reading.}
\psentry{ps_134-5_004-v0002}{2}{writing over}{\over{\unsure{i}}{u}}{\hand{u}}{The superposed handwritten reading replaces the underlying reading.}
\psentry{ps_134-5_004-v0003}{3}{insertion}{\inserted{and}}{\hand{and}}{The handwritten insertion is incorporated at its indicated position; insertion brackets leave the running text.}
\psentry{ps_134-5_004-v0004}{4}{deletion}{\deleted{on these}}{}{The author cancels this material by hand; it leaves the corrected running text.}
\psentry{ps_134-5_004-v0005}{5}{deletion}{\unsure{\deleted{,}}}{}{All material inside this wrapper is cancelled. Its original doubt and presentation remain in the diplomatic text and apparatus; no empty doubt mark or page label is inserted into the corrected running text.}
\psentry{ps_134-5_004-v0006}{6}{deletion}{\deleted{,}}{}{The author cancels this material by hand; it leaves the corrected running text.}
\psentry{ps_134-5_004-v0007}{7}{deletion}{\deleted{especially}}{}{The author cancels this material by hand; it leaves the corrected running text.}
\psentry{ps_134-5_004-v0008}{8}{insertion}{\inserted{and pretty much so during}}{\hand{and pretty much so during}}{The handwritten insertion is incorporated at its indicated position; insertion brackets leave the running text.}
\psentry{ps_134-5_004-v0009}{9}{typewriter deletion}{\xout{talks}}{}{The typewriter strikes out this material; it leaves the corrected running text.}
\psentry{ps_134-5_004-v0010}{10}{typewriter deletion}{\xout{work on the}}{}{The typewriter strikes out this material; it leaves the corrected running text.}
\psentry{ps_134-5_004-v0011}{11}{insertion}{\inserted{work}}{\hand{work}}{The handwritten insertion is incorporated at its indicated position; insertion brackets leave the running text.}
\psentry{ps_134-5_004-v0012}{12}{typewriter deletion}{\xout{very}}{}{The typewriter strikes out this material; it leaves the corrected running text.}
\psentry{ps_134-5_004-v0013}{13}{typed insertion}{\inserted[typed]{transcendental}}{transcendental}{The typed insertion is incorporated at its indicated position; insertion brackets leave the running text.}
\psentry{ps_134-5_004-v0014}{14}{deletion}{\unsure{\deleted{,}}}{}{All material inside this wrapper is cancelled. Its original doubt and presentation remain in the diplomatic text and apparatus; no empty doubt mark or page label is inserted into the corrected running text.}
\psapparatusend
\fi
\pssource{ps_134-5_005}{134-5/5}{https://grothendieck.umontpellier.fr/web/viewer.html?file=../134-5.pdf\#page=6}
\psnumber{ps_134-5_005-v0001}{1}
\psnumber{ps_134-5_005-v0002}{2}
\psnumber{ps_134-5_005-v0003}{3}
\psnumber{ps_134-5_005-v0004}{4}
\psnumber{ps_134-5_005-v0005}{5}
\psnumber{ps_134-5_005-v0006}{6}
\psnumber{ps_134-5_005-v0007}{7}
\psnumber{ps_134-5_005-v0008}{8}
\psnumber{ps_134-5_005-v0009}{9}
\psnumber{ps_134-5_005-v0010}{10}
\psnumber{ps_134-5_005-v0011}{11}
\psnumber{ps_134-5_005-v0012}{12}
\psnumber{ps_134-5_005-v0013}{13}
\psnumber{ps_134-5_005-v0014}{14}
\psnumber{ps_134-5_005-v0015}{15}
\psnumber{ps_134-5_005-v0016}{16}
\psnumber{ps_134-5_005-v0017}{17}
\psnumber{ps_134-5_005-v0018}{18}
\psnumber{ps_134-5_005-v0019}{19}
\psnumber{ps_134-5_005-v0020}{20}
\psnumber{ps_134-5_005-v0021}{21}
\psnumber{ps_134-5_005-v0022}{22}
\psnumber{ps_134-5_005-v0023}{23}
\psnumber{ps_134-5_005-v0024}{24}
\psnumber{ps_134-5_005-v0025}{25}
\psnumber{ps_134-5_005-v0026}{26}
\psnumber{ps_134-5_005-v0027}{27}
\psnumber{ps_134-5_005-v0028}{28}
\psnumber{ps_134-5_005-v0029}{29}
\psnumber{ps_134-5_005-v0030}{30}
\psnumber{ps_134-5_005-v0031}{31}
\psnumber{ps_134-5_005-v0032}{32}
\psnumber{ps_134-5_005-v0033}{33}
\psnumber{ps_134-5_005-v0034}{34}
\psnumber{ps_134-5_005-v0035}{35}
\psnumber{ps_134-5_005-v0036}{36}
\psnumber{ps_134-5_005-v0037}{37}
\psnumber{ps_134-5_005-v0038}{38}
\psnumber{ps_134-5_005-v0039}{39}
\psnumber{ps_134-5_005-v0040}{40}
\psnumber{ps_134-5_005-v0041}{41}
\psnumber{ps_134-5_005-v0042}{42}
\psnumber{ps_134-5_005-v0043}{43}
\psnumber{ps_134-5_005-v0044}{44}
\psnumber{ps_134-5_005-v0045}{45}
\psnumber{ps_134-5_005-v0046}{46}
%% BEGIN SOURCE ps_134-5_005
% Pursuing Stacks (cotes 134-2 to 134-8): archive cote 134-5, site page 5 (PDF page 6), typescript page 192.
% typescript. Draft compared: ps4.tex:203-262.
% AI-transcribed by Claude Opus 5.5 (claude-opus-5-5, adaptive thinking, effort high), Message Batches API, batch ps2 (msgbatch_01NVwr1Yuo8qJf5rykeGBXPK), 2026-11-12; reviewed against both supplied scan halves; see CHECK_DECISIONS.tsv.
\scanpage{134-5/5}
\mspage{\hand{192}}
was an agreeable surprise to notice that the message \var{\inserted{this time}}{\hand{this time}}{insertion}{ps_134-5_005-v0001} was getting through.
\var{\xout{He has been the first person,x\allowbreak{}x\allowbreak{}x\allowbreak{}x\allowbreak{}x\allowbreak{}x\allowbreak{}}}{}{typewriter deletion}{ps_134-5_005-v0002} For the five or six years since my
attention became attracted by the fascinating melting-pot of key structu-%
res \var{\xout{and key situations intersecting}}{}{typewriter deletion}{ps_134-5_005-v0003} in geometry, topology, arithmetic\var{\xout{s}}{}{typewriter deletion}{ps_134-5_005-v0004},
discrete and algebraic groups, intertwining tightly in a kind of
very basic Galois-Teichmüller theory\var{\xout{?}}{}{typewriter deletion}{ps_134-5_005-v0005}, Y. \var{\xout{xxxxxxxxxx}}{}{typewriter deletion}{ps_134-5_005-v0006}\var{\inserted{has been}}{\hand{has been}}{insertion}{ps_134-5_005-v0007} the very first
\var{\xout{the has}}{}{typewriter deletion}{ps_134-5_005-v0008} person I met \var{\inserted[typed]{so far}}{so far}{typed insertion}{ps_134-5_005-v0009} to have a feeling\var{\inserted[typed]{for}}{for}{typed insertion}{ps_134-5_005-v0010}\var{\xout{x}}{}{typewriter deletion}{ps_134-5_005-v0011}(a not yet dulled instinct I might
say\var{\xout{,}}{}{typewriter deletion}{ps_134-5_005-v0012}\var{\inserted{for sensing)}}{\hand{for sensing)}}{insertion}{ps_134-5_005-v0013} \var{\deleted{allowing him to scent,)}}{}{deletion}{ps_134-5_005-v0014} the extraordinary riches\hand{,} opening up \var{\inserted[typed]{here}}{here}{typed insertion}{ps_134-5_005-v0015}\var{\deleted{for}}{}{deletion}{ps_134-5_005-v0016} inves-%
tigation. T\hand{h}e series of talks I had given  in a tentative seminar last year
\var{\inserted[typed]{had}}{had}{typed insertion}{ps_134-5_005-v0017} turned short, by lack of any active interest and participation of \var{\inserted[typed]{anyone among}}{anyone among}{typed insertion}{ps_134-5_005-v0018} the hand-%
full of \var{\inserted{mere}}{\hand{mere}}{insertion}{ps_134-5_005-v0019} listeners. And the two or three occasions I had \var{\xout{in}}{}{typewriter deletion}{ps_134-5_005-v0020} the years be\var{\xout{f}}{}{typewriter deletion}{ps_134-5_005-v0021}\hand{-}
fore to tell about the matter of \var{\xout{anabelian algebraic geometry}}{}{typewriter deletion}{ps_134-5_005-v0022} two-dimensi
onal maps\var{\inserted[typed]{(``cartes'')}}{(``cartes'')}{typed insertion}{ps_134-5_005-v0023} and their amazing \unsure{algebro-arithmetic|algebro-aithmetic} implications,  to a few
highbrow colleagues with uncomparably wider background and knowhow than
anyone around here, I met \var{\xout{only}}{}{typewriter deletion}{ps_134-5_005-v0024} with polite interest,  or polite indiffe-%
rence which is the same.  As there was nobody around \var{\over{\unsure{anywhere|anyway}}{anyhow}}{\hand{anyhow}}{writing over}{ps_134-5_005-v0025} to take any
interest in th\var{\over{e}{ese}}{\hand{ese}}{writing over}{ps_134-5_005-v0026} juicy\var{\deleted{, basic}}{}{deletion}{ps_134-5_005-v0027} greenlands nobody would care to see,  becau-%
se there was no text-book nor any official seminar notes to prove they
existed,  after a few years I \var{\inserted[typed]{finally}}{finally}{typed insertion}{ps_134-5_005-v0028} set off myself for a preliminary voyage.

\var{\xout{\unsure{ww}}}{}{typewriter deletion}{ps_134-5_005-v0029} I thought it was going to take me a week or two to tour it and \var{\xout{xxx}}{}{typewriter deletion}{ps_134-5_005-v0030}
kind of recense resources. It took me five months instead \var{\inserted{of intensive work,}}{\hand{of intensive work,}}{insertion}{ps_134-5_005-v0031} and \var{\inserted[typed]{two}}{two}{typed insertion}{ps_134-5_005-v0032} impressive
heaps of notes (baptized  ``La Longue Marche à travers la théorie de
Galois'' ),  to get a \var{\xout{preliminary}}{}{typewriter deletion}{ps_134-5_005-v0033} first,  approximative grasp of \var{\inserted[typed]{some}}{some}{typed insertion}{ps_134-5_005-v0034} the
main structures and relationships involved, \var{\xout{xxxxxx}}{}{typewriter deletion}{ps_134-5_005-v0035}\var{\inserted[typed]{The}}{The}{typed insertion}{ps_134-5_005-v0036} main emphasis \var{\inserted[typed]{was (still is)}}{was (still is)}{typed insertion}{ps_134-5_005-v0037} upon \margin{/}
an understanding of the\var{\xout{x}}{}{typewriter deletion}{ps_134-5_005-v0038} action of profinite Galois-groups (foremost
among which $\mathrm{Gal}_{\overline{Q}/Q}$ and the subgroups of finite index) on non-commutative
profinite fundamental groups,  and primarily on fundamental groups of
algebraic curves,\var{\deleted{and}}{}{deletion}{ps_134-5_005-v0039}\var{\inserted[typed]{-increasingly too on those of}}{-increasingly too on those of}{typed insertion}{ps_134-5_005-v0040} modular varieties (more accurately, \var{\inserted[typed]{modular}}{modular}{typed insertion}{ps_134-5_005-v0041} \uline{multiplicities}\hand{)}
for such curves -  \var{\deleted{namely}}{}{deletion}{ps_134-5_005-v0042}\var{\inserted[typed]{on the}}{on the}{typed insertion}{ps_134-5_005-v0043} profinite completions of the Teichmüller group\hand{\unsure{s.}}
The voyage was the most rewarding and exciting I had in mathematics so
far  -  and still it became very clear that it was just like a first
glimpse upon a wholly new landscape - one landscape surely among count-%
less others of a\var{\xout{n}}{}{typewriter deletion}{ps_134-5_005-v0044} continent unknown,  eager to be discovered.

This was in the first half of the year 1981  -  just two years ago,
\var{\inserted[typed]{it turns out,}}{it turns out,}{typed insertion}{ps_134-5_005-v0045} but it looks almost infinitely remote,  because such a lot of things to\hand{ok}
place since.  Looking back,  it turns out there have been since roughly
four main \var{\inserted{alternating}}{\hand{alternating}}{insertion}{ps_134-5_005-v0046} periods of reflection,  one period of reflection on personal
matters alternating with one on mathematics.  The next mathematical
reflexion started with a long digression on tame topology and the ``dépl\unsure{o}\hand{y-}

\margin{5}

%% END SOURCE ps_134-5_005
\ifnum\EditionMode=1
\psapparatusstart
\psentry{ps_134-5_005-v0001}{1}{insertion}{\inserted{this time}}{\hand{this time}}{The handwritten insertion is incorporated at its indicated position; insertion brackets leave the running text.}
\psentry{ps_134-5_005-v0002}{2}{typewriter deletion}{\xout{He has been the first person,x\allowbreak{}x\allowbreak{}x\allowbreak{}x\allowbreak{}x\allowbreak{}x\allowbreak{}}}{}{The typewriter strikes out this material; it leaves the corrected running text.}
\psentry{ps_134-5_005-v0003}{3}{typewriter deletion}{\xout{and key situations intersecting}}{}{The typewriter strikes out this material; it leaves the corrected running text.}
\psentry{ps_134-5_005-v0004}{4}{typewriter deletion}{\xout{s}}{}{The typewriter strikes out this material; it leaves the corrected running text.}
\psentry{ps_134-5_005-v0005}{5}{typewriter deletion}{\xout{?}}{}{The typewriter strikes out this material; it leaves the corrected running text.}
\psentry{ps_134-5_005-v0006}{6}{typewriter deletion}{\xout{xxxxxxxxxx}}{}{The typewriter strikes out this material; it leaves the corrected running text.}
\psentry{ps_134-5_005-v0007}{7}{insertion}{\inserted{has been}}{\hand{has been}}{The handwritten insertion is incorporated at its indicated position; insertion brackets leave the running text.}
\psentry{ps_134-5_005-v0008}{8}{typewriter deletion}{\xout{the has}}{}{The typewriter strikes out this material; it leaves the corrected running text.}
\psentry{ps_134-5_005-v0009}{9}{typed insertion}{\inserted[typed]{so far}}{so far}{The typed insertion is incorporated at its indicated position; insertion brackets leave the running text.}
\psentry{ps_134-5_005-v0010}{10}{typed insertion}{\inserted[typed]{for}}{for}{The typed insertion is incorporated at its indicated position; insertion brackets leave the running text.}
\psentry{ps_134-5_005-v0011}{11}{typewriter deletion}{\xout{x}}{}{The typewriter strikes out this material; it leaves the corrected running text.}
\psentry{ps_134-5_005-v0012}{12}{typewriter deletion}{\xout{,}}{}{The typewriter strikes out this material; it leaves the corrected running text.}
\psentry{ps_134-5_005-v0013}{13}{insertion}{\inserted{for sensing)}}{\hand{for sensing)}}{The handwritten insertion is incorporated at its indicated position; insertion brackets leave the running text.}
\psentry{ps_134-5_005-v0014}{14}{deletion}{\deleted{allowing him to scent,)}}{}{The author cancels this material by hand; it leaves the corrected running text.}
\psentry{ps_134-5_005-v0015}{15}{typed insertion}{\inserted[typed]{here}}{here}{The typed insertion is incorporated at its indicated position; insertion brackets leave the running text.}
\psentry{ps_134-5_005-v0016}{16}{deletion}{\deleted{for}}{}{The author cancels this material by hand; it leaves the corrected running text.}
\psentry{ps_134-5_005-v0017}{17}{typed insertion}{\inserted[typed]{had}}{had}{The typed insertion is incorporated at its indicated position; insertion brackets leave the running text.}
\psentry{ps_134-5_005-v0018}{18}{typed insertion}{\inserted[typed]{anyone among}}{anyone among}{The typed insertion is incorporated at its indicated position; insertion brackets leave the running text.}
\psentry{ps_134-5_005-v0019}{19}{insertion}{\inserted{mere}}{\hand{mere}}{The handwritten insertion is incorporated at its indicated position; insertion brackets leave the running text.}
\psentry{ps_134-5_005-v0020}{20}{typewriter deletion}{\xout{in}}{}{The typewriter strikes out this material; it leaves the corrected running text.}
\psentry{ps_134-5_005-v0021}{21}{typewriter deletion}{\xout{f}}{}{The typewriter strikes out this material; it leaves the corrected running text.}
\psentry{ps_134-5_005-v0022}{22}{typewriter deletion}{\xout{anabelian algebraic geometry}}{}{The typewriter strikes out this material; it leaves the corrected running text.}
\psentry{ps_134-5_005-v0023}{23}{typed insertion}{\inserted[typed]{(``cartes'')}}{(``cartes'')}{The typed insertion is incorporated at its indicated position; insertion brackets leave the running text.}
\psentry{ps_134-5_005-v0024}{24}{typewriter deletion}{\xout{only}}{}{The typewriter strikes out this material; it leaves the corrected running text.}
\psentry{ps_134-5_005-v0025}{25}{writing over}{\over{\unsure{anywhere|anyway}}{anyhow}}{\hand{anyhow}}{The superposed handwritten reading replaces the underlying reading.}
\psentry{ps_134-5_005-v0026}{26}{writing over}{\over{e}{ese}}{\hand{ese}}{The superposed handwritten reading replaces the underlying reading.}
\psentry{ps_134-5_005-v0027}{27}{deletion}{\deleted{, basic}}{}{The author cancels this material by hand; it leaves the corrected running text.}
\psentry{ps_134-5_005-v0028}{28}{typed insertion}{\inserted[typed]{finally}}{finally}{The typed insertion is incorporated at its indicated position; insertion brackets leave the running text.}
\psentry{ps_134-5_005-v0029}{29}{typewriter deletion}{\xout{\unsure{ww}}}{}{The typewriter strikes out this material; it leaves the corrected running text.}
\psentry{ps_134-5_005-v0030}{30}{typewriter deletion}{\xout{xxx}}{}{The typewriter strikes out this material; it leaves the corrected running text.}
\psentry{ps_134-5_005-v0031}{31}{insertion}{\inserted{of intensive work,}}{\hand{of intensive work,}}{The handwritten insertion is incorporated at its indicated position; insertion brackets leave the running text.}
\psentry{ps_134-5_005-v0032}{32}{typed insertion}{\inserted[typed]{two}}{two}{The typed insertion is incorporated at its indicated position; insertion brackets leave the running text.}
\psentry{ps_134-5_005-v0033}{33}{typewriter deletion}{\xout{preliminary}}{}{The typewriter strikes out this material; it leaves the corrected running text.}
\psentry{ps_134-5_005-v0034}{34}{typed insertion}{\inserted[typed]{some}}{some}{The typed insertion is incorporated at its indicated position; insertion brackets leave the running text.}
\psentry{ps_134-5_005-v0035}{35}{typewriter deletion}{\xout{xxxxxx}}{}{The typewriter strikes out this material; it leaves the corrected running text.}
\psentry{ps_134-5_005-v0036}{36}{typed insertion}{\inserted[typed]{The}}{The}{The typed insertion is incorporated at its indicated position; insertion brackets leave the running text.}
\psentry{ps_134-5_005-v0037}{37}{typed insertion}{\inserted[typed]{was (still is)}}{was (still is)}{The typed insertion is incorporated at its indicated position; insertion brackets leave the running text.}
\psentry{ps_134-5_005-v0038}{38}{typewriter deletion}{\xout{x}}{}{The typewriter strikes out this material; it leaves the corrected running text.}
\psentry{ps_134-5_005-v0039}{39}{deletion}{\deleted{and}}{}{The author cancels this material by hand; it leaves the corrected running text.}
\psentry{ps_134-5_005-v0040}{40}{typed insertion}{\inserted[typed]{-increasingly too on those of}}{-increasingly too on those of}{The typed insertion is incorporated at its indicated position; insertion brackets leave the running text.}
\psentry{ps_134-5_005-v0041}{41}{typed insertion}{\inserted[typed]{modular}}{modular}{The typed insertion is incorporated at its indicated position; insertion brackets leave the running text.}
\psentry{ps_134-5_005-v0042}{42}{deletion}{\deleted{namely}}{}{The author cancels this material by hand; it leaves the corrected running text.}
\psentry{ps_134-5_005-v0043}{43}{typed insertion}{\inserted[typed]{on the}}{on the}{The typed insertion is incorporated at its indicated position; insertion brackets leave the running text.}
\psentry{ps_134-5_005-v0044}{44}{typewriter deletion}{\xout{n}}{}{The typewriter strikes out this material; it leaves the corrected running text.}
\psentry{ps_134-5_005-v0045}{45}{typed insertion}{\inserted[typed]{it turns out,}}{it turns out,}{The typed insertion is incorporated at its indicated position; insertion brackets leave the running text.}
\psentry{ps_134-5_005-v0046}{46}{insertion}{\inserted{alternating}}{\hand{alternating}}{The handwritten insertion is incorporated at its indicated position; insertion brackets leave the running text.}
\psapparatusend
\fi
\pssource{ps_134-5_006}{134-5/6}{https://grothendieck.umontpellier.fr/web/viewer.html?file=../134-5.pdf\#page=7}
\psnumber{ps_134-5_006-v0001}{1}
\psnumber{ps_134-5_006-v0002}{2}
\psnumber{ps_134-5_006-v0003}{3}
\psnumber{ps_134-5_006-v0004}{4}
\psnumber{ps_134-5_006-v0005}{5}
\psnumber{ps_134-5_006-v0006}{6}
\psnumber{ps_134-5_006-v0007}{7}
\psnumber{ps_134-5_006-v0008}{8}
\psnumber{ps_134-5_006-v0009}{9}
\psnumber{ps_134-5_006-v0010}{10}
\psnumber{ps_134-5_006-v0011}{11}
\psnumber{ps_134-5_006-v0012}{12}
\psnumber{ps_134-5_006-v0013}{13}
\psnumber{ps_134-5_006-v0014}{14}
\psnumber{ps_134-5_006-v0015}{15}
\psnumber{ps_134-5_006-v0016}{16}
\psnumber{ps_134-5_006-v0017}{17}
\psnumber{ps_134-5_006-v0018}{18}
\psnumber{ps_134-5_006-v0019}{19}
\psnumber{ps_134-5_006-v0020}{20}
\psnumber{ps_134-5_006-v0021}{21}
\psnumber{ps_134-5_006-v0022}{22}
\psnumber{ps_134-5_006-v0023}{23}
\psnumber{ps_134-5_006-v0024}{24}
%% BEGIN SOURCE ps_134-5_006
% Pursuing Stacks (cotes 134-2 to 134-8): archive cote 134-5, site page 6 (PDF page 7), typescript page 193.
% typescript. Draft compared: ps4.tex:249-307.
% AI-transcribed by Claude Opus 5.5 (claude-opus-5-5, adaptive thinking, effort high), Message Batches API, batch ps2 (msgbatch_01NVwr1Yuo8qJf5rykeGBXPK), 2026-11-12; reviewed against both supplied scan halves; see CHECK_DECISIONS.tsv.
\scanpage{134-5/6}
\mspage{193}

ment'' (``unfolding'') of stratified structures,  as a leading thread toward\unsure{s}
% NOTE: the right edge of the page is cut off; the last letters of several lines are hidden
a heuristic understanding of the natural stratification of \var{\inserted[typed]{the}}{the}{typed insertion}{ps_134-5_006-v0001} Mumford-Delign\unsure{e}
compac\unsure{ti}fications of modular multiplicities $M_{g,\hand{\nu}}$ (for curves of genus $g$
\unsure{andowed|endowed} with \hand{$\nu$} points).  This then led to the ``anabelian seminar'' which
\var{\xout{has}}{}{typewriter deletion}{ps_134-5_006-v0002} turned short,  last spring.  Then a month or two sickness,  intensive
meditation \var{\xout{and}}{}{typewriter deletion}{ps_134-5_006-v0003} for three or four months,  a few more months for settling
some important personal matter;  and now,since february,another mathemati%
-cal reflection started. \hand{$\rfloor$}I \var{\xout{have}}{}{typewriter deletion}{ps_134-5_006-v0004} am unable to tell the meaning of this
alternation of periods of meditation on personal matters and periods of
mathematical reflection,  which has been going through my life for the
last seven years,  more and more,\var{\xout{\unsure{like the}}}{}{typewriter deletion}{ps_134-5_006-v0005}- very much like the unceasing
up and down of waves,  or like a steady breathing going through my life,
without any attempt any longer of controlling it one way or other.  One
common moving force surely is the inborn curiosity  -  a thirst \var{\over{of}{for}}{\hand{for}}{writing over}{ps_134-5_006-v0006} gettin\unsure{g}
acquainted with the juicy things of th\var{\over{is}{e}}{\hand{e}}{writing over}{ps_134-5_006-v0007} \var{\inserted[typed]{unexhaustible}}{unexhaustible}{typed insertion}{ps_134-5_006-v0008} world,  whether they be the brea-%
thing body of the Beloved,  or the evasive substance of one's own life,
or the much less evasive substance of mathematical structures and their
delicate interplay.  This thirst in itself is of a nature quite different
of the ego's  -  it is the thirst of life to know about itself,  a primal
creative force which,one suspects,has been around forever,  long before
a human ego  -  a bundle \var{\inserted{mainly}}{\hand{mainly}}{insertion}{ps_134-5_006-v0009} of fears,  of inhibitions and self-deceptions -
came into being.  Still I am aware that the ego is strongly involved in
the particular way in which the creative force expresses itself,  in my
own life or anyone's else's (when this force is allowed to come into pla\unsure{y}
at all ...).  The motivations behind any strong energy investment,  and
more particularly so when in it \var{\deleted{in}}{}{deletion}{ps_134-5_006-v0010}\var{\inserted[typed]{is}}{is}{typed insertion}{ps_134-5_006-v0011} an activity attached with any kind of
social status or prestige, are a lot more complex and ego-\var{\deleted{\unsure{loaded}}}{}{deletion}{ps_134-5_006-v0012}\var{\inserted{driven}}{\hand{driven}}{insertion}{ps_134-5_006-v0013} than
one generally cares to admit.  True,  ambition by itself is \var{\deleted{utterly}}{}{deletion}{ps_134-5_006-v0014} power-%
less for discovering or understanding \var{\inserted[typed]{or perceiving}}{or perceiving}{typed insertion}{ps_134-5_006-v0015} \var{\over{an}{an}}{\hand{an}}{writing over}{ps_134-5_006-v0016}ything substantial whatever,
neither a mathematical relationship nor the perfume of a flower.  In \var{\deleted{the}}{}{deletion}{ps_134-5_006-v0017}
moments of work and \var{\inserted[typed]{of}}{of}{typed insertion}{ps_134-5_006-v0018} discovery,  in any creative moments in life \var{\inserted[typed]{as a matter of fact}}{as a matter of fact}{typed insertion}{ps_134-5_006-v0019}, ambitio\unsure{n}
is absent; the Artisan is a keen interest,  which is just one of the mani-%
fold aspects of love.  What \var{\inserted[typed]{however}}{however}{typed insertion}{ps_134-5_006-v0020} pushes us so relentlessly to work again,and
so often causes our life \var{\inserted[typed]{and passion}}{and passion}{typed insertion}{ps_134-5_006-v0021} gradually to dry off and become unsensitive,
even to the kindred passion of a fellow-being\hand{,} or to the un\unsure{sus}pected,
bea\var{\inserted[typed]{u}}{u}{typed insertion}{ps_134-5_006-v0022}ties and mysteries of the very field we are supposed to be ploughing
-  this force is neither love nor keen interest for the things \var{\deleted{an}}{}{deletion}{ps_134-5_006-v0023}\var{\inserted[typed]{and}}{and}{typed insertion}{ps_134-5_006-v0024} beings
in this world.  It is interesting enough though,  and surely deserves
a close look !

\margin{6}

%% END SOURCE ps_134-5_006
\ifnum\EditionMode=1
\psapparatusstart
\psentry{ps_134-5_006-v0001}{1}{typed insertion}{\inserted[typed]{the}}{the}{The typed insertion is incorporated at its indicated position; insertion brackets leave the running text.}
\psentry{ps_134-5_006-v0002}{2}{typewriter deletion}{\xout{has}}{}{The typewriter strikes out this material; it leaves the corrected running text.}
\psentry{ps_134-5_006-v0003}{3}{typewriter deletion}{\xout{and}}{}{The typewriter strikes out this material; it leaves the corrected running text.}
\psentry{ps_134-5_006-v0004}{4}{typewriter deletion}{\xout{have}}{}{The typewriter strikes out this material; it leaves the corrected running text.}
\psentry{ps_134-5_006-v0005}{5}{typewriter deletion}{\xout{\unsure{like the}}}{}{The typewriter strikes out this material; it leaves the corrected running text.}
\psentry{ps_134-5_006-v0006}{6}{writing over}{\over{of}{for}}{\hand{for}}{The superposed handwritten reading replaces the underlying reading.}
\psentry{ps_134-5_006-v0007}{7}{writing over}{\over{is}{e}}{\hand{e}}{The superposed handwritten reading replaces the underlying reading.}
\psentry{ps_134-5_006-v0008}{8}{typed insertion}{\inserted[typed]{unexhaustible}}{unexhaustible}{The typed insertion is incorporated at its indicated position; insertion brackets leave the running text.}
\psentry{ps_134-5_006-v0009}{9}{insertion}{\inserted{mainly}}{\hand{mainly}}{The handwritten insertion is incorporated at its indicated position; insertion brackets leave the running text.}
\psentry{ps_134-5_006-v0010}{10}{deletion}{\deleted{in}}{}{The author cancels this material by hand; it leaves the corrected running text.}
\psentry{ps_134-5_006-v0011}{11}{typed insertion}{\inserted[typed]{is}}{is}{The typed insertion is incorporated at its indicated position; insertion brackets leave the running text.}
\psentry{ps_134-5_006-v0012}{12}{deletion}{\deleted{\unsure{loaded}}}{}{The author cancels this material by hand; it leaves the corrected running text.}
\psentry{ps_134-5_006-v0013}{13}{insertion}{\inserted{driven}}{\hand{driven}}{The handwritten insertion is incorporated at its indicated position; insertion brackets leave the running text.}
\psentry{ps_134-5_006-v0014}{14}{deletion}{\deleted{utterly}}{}{The author cancels this material by hand; it leaves the corrected running text.}
\psentry{ps_134-5_006-v0015}{15}{typed insertion}{\inserted[typed]{or perceiving}}{or perceiving}{The typed insertion is incorporated at its indicated position; insertion brackets leave the running text.}
\psentry{ps_134-5_006-v0016}{16}{writing over}{\over{an}{an}}{\hand{an}}{The superposed handwritten reading replaces the underlying reading.}
\psentry{ps_134-5_006-v0017}{17}{deletion}{\deleted{the}}{}{The author cancels this material by hand; it leaves the corrected running text.}
\psentry{ps_134-5_006-v0018}{18}{typed insertion}{\inserted[typed]{of}}{of}{The typed insertion is incorporated at its indicated position; insertion brackets leave the running text.}
\psentry{ps_134-5_006-v0019}{19}{typed insertion}{\inserted[typed]{as a matter of fact}}{as a matter of fact}{The typed insertion is incorporated at its indicated position; insertion brackets leave the running text.}
\psentry{ps_134-5_006-v0020}{20}{typed insertion}{\inserted[typed]{however}}{however}{The typed insertion is incorporated at its indicated position; insertion brackets leave the running text.}
\psentry{ps_134-5_006-v0021}{21}{typed insertion}{\inserted[typed]{and passion}}{and passion}{The typed insertion is incorporated at its indicated position; insertion brackets leave the running text.}
\psentry{ps_134-5_006-v0022}{22}{typed insertion}{\inserted[typed]{u}}{u}{The typed insertion is incorporated at its indicated position; insertion brackets leave the running text.}
\psentry{ps_134-5_006-v0023}{23}{deletion}{\deleted{an}}{}{The author cancels this material by hand; it leaves the corrected running text.}
\psentry{ps_134-5_006-v0024}{24}{typed insertion}{\inserted[typed]{and}}{and}{The typed insertion is incorporated at its indicated position; insertion brackets leave the running text.}
\psapparatusend
\fi
\pssource{ps_134-5_007}{134-5/7}{https://grothendieck.umontpellier.fr/web/viewer.html?file=../134-5.pdf\#page=8}
\psnumber{ps_134-5_007-v0001}{1}
\psnumber{ps_134-5_007-v0002}{2}
\psnumber{ps_134-5_007-v0003}{3}
\psnumber{ps_134-5_007-v0004}{4}
\psnumber{ps_134-5_007-v0005}{5}
\psnumber{ps_134-5_007-v0006}{6}
\psnumber{ps_134-5_007-v0007}{7}
\psnumber{ps_134-5_007-v0008}{8}
\psnumber{ps_134-5_007-v0009}{9}
\psnumber{ps_134-5_007-v0010}{10}
\psnumber{ps_134-5_007-v0011}{11}
\psnumber{ps_134-5_007-v0012}{12}
\psnumber{ps_134-5_007-v0013}{13}
\psnumber{ps_134-5_007-v0014}{14}
\psnumber{ps_134-5_007-v0015}{15}
\psnumber{ps_134-5_007-v0016}{16}
\psnumber{ps_134-5_007-v0017}{17}
\psnumber{ps_134-5_007-v0018}{18}
\psnumber{ps_134-5_007-v0019}{19}
\psnumber{ps_134-5_007-v0020}{20}
\psnumber{ps_134-5_007-v0021}{21}
\psnumber{ps_134-5_007-v0022}{22}
%% BEGIN SOURCE ps_134-5_007
% Pursuing Stacks (cotes 134-2 to 134-8): archive cote 134-5, site page 7 (PDF page 8), typescript page 194.
% typescript. Draft compared: ps4.tex:293-349.
% AI-transcribed by Claude Opus 5.5 (claude-opus-5-5, adaptive thinking, effort high), Message Batches API, batch ps2 (msgbatch_01NVwr1Yuo8qJf5rykeGBXPK), 2026-11-12; reviewed against both supplied scan halves; see CHECK_DECISIONS.tsv.
\scanpage{134-5/7}
\mspage{194}

I thought I was starting a retrospective of six weeks of scratchwork
on homotopical algebra  -  and it turns out to be a \var{\inserted[typed]{(very)}}{(very)}{typed insertion}{ps_134-5_007-v0001} short retrospective
rather of the up-and-down movement of \var{\inserted[typed]{my}}{my}{typed insertion}{ps_134-5_007-v0002} mathematical interests and invest-%
ments\hand{|}during the last years.  Doubtlessly the very strongest\unsure{\hand{,}} \var{\inserted[typed]{attraction\unsure{,}}}{attraction\unsure{,}}{typed insertion}{ps_134-5_007-v0003} the greatest
fascination goes with this ``new world'' of anabelian algebraic geometry.
It may seem strange that instead,  I am indulging in this lengthy digressi-%
-on on homotopical algebra,  which is almost wholly irrelevant I feel for
the Galois-Teichm\var{\xout{\unsure{u}}}{}{typewriter deletion}{ps_134-5_007-v0004}üller story.  The reason\hand{|}is surely an inner reluctance,
an unreadiness to embark upon a long-term voyage,  well knowing that it is
\var{\inserted[typed]{so}}{so}{typed insertion}{ps_134-5_007-v0005} enticing \var{\xout{xxxxxxx,}}{}{typewriter deletion}{ps_134-5_007-v0006}that I may well be caught in this game for a number of
years - not doing anything else \var{\xout{from}}{}{typewriter deletion}{ps_134-5_007-v0007} day and night than making love with
mathematics,  and maybe sleeping and eating now and then.  I have gone
through this a number of times, \var{\inserted[typed]{and}}{and}{typed insertion}{ps_134-5_007-v0008} at times I thought I was through.  Finally
I came to admit and to accept,  two years ago,  I was not through yet -
this was during the months of meditation after that ``long march through
Galois theory'' -  which had been\hand{,} too\hand{,} a wholly unexpected fit of mathemati-%
cal passion, not to say frenzy. \var{\deleted{Finally}}{}{deletion}{ps_134-5_007-v0009} \var{\inserted{And}}{\hand{And}}{insertion}{ps_134-5_007-v0010} \var{\over{D}{d}}{\hand{d}}{writing over}{ps_134-5_007-v0011}uring the last weeks,  just
reflecting a little here and there upon the Teichmüller-Légo game and
it's arithmetical implications,  I let myself be caught again by this
fascination \var{\xout{ofxxx}}{}{typewriter deletion}{ps_134-5_007-v0012} - it is becoming kind of clear \var{\inserted[typed]{now}}{now}{typed insertion}{ps_134-5_007-v0013} that I am going to finish
writing up those notes on algebra,  almost like some homework that has got
to be done (anyhow I like to finish when I started \var{\xout{xx}}{}{typewriter deletion}{ps_134-5_007-v0014} something)  -  and
as soon as I'm through with the notes,  back to geom\unsure{\var{\over{w}{e}}{\hand{e}}{writing over}{ps_134-5_007-v0015}}try in the long last !
Also,  the idea is in the air for the last few months  -  since I decided
to publish these unformal notes on stacks or whatever it'll turn out to be -
that I may well \var{\inserted[typed]{go on}}{go on}{typed insertion}{ps_134-5_007-v0016} the same way,  writing up and publishing unformal notes
on other topics,  including tame topology and anabelian algebraic geometry.
In contrast to the present notes,  I got heaps of scratchwork done \var{\inserted[typed]{on these}}{on these}{typed insertion}{ps_134-5_007-v0017} in the
years before  -  in this respect time is even riper for me to ramble \var{\inserted{``publicly'',}}{\hand{``publicly'',}}{insertion}{ps_134-5_007-v0018} \var{\deleted{on}}{}{deletion}{ps_134-5_007-v0019}
\var{\deleted{\xout{these};}}{}{deletion}{ps_134-5_007-v0020} than on stacks and homotopy theory\unsure{\var{\over{.}{!}}{\hand{!}}{writing over}{ps_134-5_007-v0021}}

From Y.  who looked through a lot of litterature on the subject,  it
strikes me (agreeably of course) that nobody yet hit upon ``the'' natural
presentation of the Teichmüller groupoids,  which kind of imposes itself
quite \var{\over{naturall}{forcibl}}{\hand{forcibl}}{writing over}{ps_134-5_007-v0022}y in the set-up I let myself be guided by.  Technically spea-%
king (and this will rejoice Ronnie Brown I'm sure!),  I suspect one main rea-%
son why this is so,  is that people are accustomed to work with fundamental
groups and generators and relations for these and stick to it,  even in
contexts when this is wholly unadequate,  namely when you get a clear

\margin{7}

%% END SOURCE ps_134-5_007
\ifnum\EditionMode=1
\psapparatusstart
\psentry{ps_134-5_007-v0001}{1}{typed insertion}{\inserted[typed]{(very)}}{(very)}{The typed insertion is incorporated at its indicated position; insertion brackets leave the running text.}
\psentry{ps_134-5_007-v0002}{2}{typed insertion}{\inserted[typed]{my}}{my}{The typed insertion is incorporated at its indicated position; insertion brackets leave the running text.}
\psentry{ps_134-5_007-v0003}{3}{typed insertion}{\inserted[typed]{attraction\unsure{,}}}{attraction\unsure{,}}{The typed insertion is incorporated at its indicated position; insertion brackets leave the running text.}
\psentry{ps_134-5_007-v0004}{4}{typewriter deletion}{\xout{\unsure{u}}}{}{The typewriter strikes out this material; it leaves the corrected running text.}
\psentry{ps_134-5_007-v0005}{5}{typed insertion}{\inserted[typed]{so}}{so}{The typed insertion is incorporated at its indicated position; insertion brackets leave the running text.}
\psentry{ps_134-5_007-v0006}{6}{typewriter deletion}{\xout{xxxxxxx,}}{}{The typewriter strikes out this material; it leaves the corrected running text.}
\psentry{ps_134-5_007-v0007}{7}{typewriter deletion}{\xout{from}}{}{The typewriter strikes out this material; it leaves the corrected running text.}
\psentry{ps_134-5_007-v0008}{8}{typed insertion}{\inserted[typed]{and}}{and}{The typed insertion is incorporated at its indicated position; insertion brackets leave the running text.}
\psentry{ps_134-5_007-v0009}{9}{deletion}{\deleted{Finally}}{}{The author cancels this material by hand; it leaves the corrected running text.}
\psentry{ps_134-5_007-v0010}{10}{insertion}{\inserted{And}}{\hand{And}}{The handwritten insertion is incorporated at its indicated position; insertion brackets leave the running text.}
\psentry{ps_134-5_007-v0011}{11}{writing over}{\over{D}{d}}{\hand{d}}{The superposed handwritten reading replaces the underlying reading.}
\psentry{ps_134-5_007-v0012}{12}{typewriter deletion}{\xout{ofxxx}}{}{The typewriter strikes out this material; it leaves the corrected running text.}
\psentry{ps_134-5_007-v0013}{13}{typed insertion}{\inserted[typed]{now}}{now}{The typed insertion is incorporated at its indicated position; insertion brackets leave the running text.}
\psentry{ps_134-5_007-v0014}{14}{typewriter deletion}{\xout{xx}}{}{The typewriter strikes out this material; it leaves the corrected running text.}
\psentry{ps_134-5_007-v0015}{15}{writing over}{\over{w}{e}}{\hand{e}}{The superposed handwritten reading replaces the underlying reading.}
\psentry{ps_134-5_007-v0016}{16}{typed insertion}{\inserted[typed]{go on}}{go on}{The typed insertion is incorporated at its indicated position; insertion brackets leave the running text.}
\psentry{ps_134-5_007-v0017}{17}{typed insertion}{\inserted[typed]{on these}}{on these}{The typed insertion is incorporated at its indicated position; insertion brackets leave the running text.}
\psentry{ps_134-5_007-v0018}{18}{insertion}{\inserted{``publicly'',}}{\hand{``publicly'',}}{The handwritten insertion is incorporated at its indicated position; insertion brackets leave the running text.}
\psentry{ps_134-5_007-v0019}{19}{deletion}{\deleted{on}}{}{The author cancels this material by hand; it leaves the corrected running text.}
\psentry{ps_134-5_007-v0020}{20}{deletion}{\deleted{\xout{these};}}{}{The author cancels this material by hand; it leaves the corrected running text.}
\psentry{ps_134-5_007-v0021}{21}{writing over}{\over{.}{!}}{\hand{!}}{The superposed handwritten reading replaces the underlying reading.}
\psentry{ps_134-5_007-v0022}{22}{writing over}{\over{naturall}{forcibl}}{\hand{forcibl}}{The superposed handwritten reading replaces the underlying reading.}
\psapparatusend
\fi
\pssource{ps_134-5_008}{134-5/8}{https://grothendieck.umontpellier.fr/web/viewer.html?file=../134-5.pdf\#page=9}
\psnumber{ps_134-5_008-v0001}{1}
\psnumber{ps_134-5_008-v0002}{2}
\psnumber{ps_134-5_008-v0003}{3}
\psnumber{ps_134-5_008-v0004}{4}
\psnumber{ps_134-5_008-v0005}{5}
\psnumber{ps_134-5_008-v0006}{6}
\psnumber{ps_134-5_008-v0007}{7}
\psnumber{ps_134-5_008-v0008}{8}
\psnumber{ps_134-5_008-v0009}{9}
\psnumber{ps_134-5_008-v0010}{10}
\psnumber{ps_134-5_008-v0011}{11}
\psnumber{ps_134-5_008-v0012}{12}
\psnumber{ps_134-5_008-v0013}{13}
\psnumber{ps_134-5_008-v0014}{14}
\psnumber{ps_134-5_008-v0015}{15}
\psnumber{ps_134-5_008-v0016}{16}
\psnumber{ps_134-5_008-v0017}{17}
%% BEGIN SOURCE ps_134-5_008
% Pursuing Stacks (cotes 134-2 to 134-8): archive cote 134-5, site page 8 (PDF page 9), typescript page 195.
% typescript. Draft compared: ps4.tex:337-392.
% AI-transcribed by Claude Opus 5.5 (claude-opus-5-5, adaptive thinking, effort high), Message Batches API, batch ps2 (msgbatch_01NVwr1Yuo8qJf5rykeGBXPK), 2026-11-12; reviewed against both supplied scan halves; see CHECK_DECISIONS.tsv.
\scanpage{134-5/8}
\mspage{195}
description by generators and relations only when working simultaneously
with a whole bunch of base-points chosen with care -- or equivalently,
working in the algebraic context of \uline{groupoids} , rather than groups. Cho-%
sing paths for connecting the basepoints natural to the situation to \var{\inserted{just}}{\hand{just}}{insertion}{ps_134-5_008-v0001} one
among them, and reducing the groupoid to a single group, will then hope-%
lessly destroy the structure and inner symmetries of the situation, and
result in a mess of generators and relations no one dares to write down,
because everyone feels they won't be of any use whatever, and just confu-%
se the picture rather than clarify it. I have known such perplexity my-%
self a long time ago, \var{\inserted{namely in}}{\hand{namely in}}{insertion}{ps_134-5_008-v0002} \var{\deleted{in \xout{situations}}}{}{deletion}{ps_134-5_008-v0003} van Kampen type situations, whose
only understandable formulation is in terms of (amalgamated sums of)
groupoids. Still, standing habits of thought are very strong, and
during the long march through Galois theory, two years ago, it took me
weeks and months trying to formulate everything in terms of groups or
``exterior groups'' (i.e. groups ``up to inner automorphism'')\hand{,} and finally
\var{\deleted{\hand{\unsure{progressively}}}}{}{deletion}{ps_134-5_008-v0004}
\var{\xout{convincing myself only groupoids would fit nicely.}}{}{typewriter deletion}{ps_134-5_008-v0005} learning the lesson
and letting myself be convinced progressively, not to say reluctantly,
that \var{\xout{only}}{}{typewriter deletion}{ps_134-5_008-v0006} groupoids only would fit nicely. Another ``technical point'' of
course is the basic fact(and the wealth of intuitions accompanying it)
that the Teichmüller groups are \var{\xout{indeed}}{}{typewriter deletion}{ps_134-5_008-v0007} fundamental groups indeed -- a
fact ignored it seems by most geometers, because the \var{\inserted[typed]{natural}}{natural}{typed insertion}{ps_134-5_008-v0008} ``spaces'' they are
fundamental groups of are not topological\var{\xout{ly}}{}{typewriter deletion}{ps_134-5_008-v0009} spaces, but \var{\xout{multiplicities}}{}{typewriter deletion}{ps_134-5_008-v0010}
the modular ``multiplicities'' $M_{g,\hand{\nu}}$ - namely \hand{\uline{\textcolor{black}{topoi}}}\hand{!} The ``points'' of these
``spaces'' are just the structures being investigated (namely algebraic
curves of type $(g,\hand{\nu})$), and the (finite) automorphism groups of these
``points''enter into the picture in a very crucial way. They can be \var{\inserted{adequately}}{\hand{adequately}}{insertion}{ps_134-5_008-v0011} chosen
as part of the system of basic generators for the Teichmüller groupoid
$T_{g,\hand{\nu}}$ , \var{\deleted{which}}{}{deletion}{ps_134-5_008-v0012} \var{\inserted{The latter of course}}{\hand{The latter of course}}{insertion}{ps_134-5_008-v0013} is essentially (up to suitable restriction of base-points)
\var{\inserted{just}}{\hand{just}}{insertion}{ps_134-5_008-v0014} the fundamental groupoid of $M_{g,\hand{\nu}}$. It is through this interpretation of
the Teichmüller groups or groupoïds that it becomes clear that the profi-%
nite Galois group $\mathrm{Gal}_{\overline{\underline{Q}}/\underline{Q}}$ operates on the profinite completion of these
and \var{\inserted{of}}{\hand{of}}{insertion}{ps_134-5_008-v0015} their various variants, and \var{\inserted{this}}{\hand{this}}{insertion}{ps_134-5_008-v0016} (it turns out) in a way respecting the
manifold structures and relationships tying them \var{\inserted[typed]{tightly}}{tightly}{typed insertion}{ps_134-5_008-v0017} together.

\margin{8}

%% END SOURCE ps_134-5_008
\ifnum\EditionMode=1
\psapparatusstart
\psentry{ps_134-5_008-v0001}{1}{insertion}{\inserted{just}}{\hand{just}}{The handwritten insertion is incorporated at its indicated position; insertion brackets leave the running text.}
\psentry{ps_134-5_008-v0002}{2}{insertion}{\inserted{namely in}}{\hand{namely in}}{The handwritten insertion is incorporated at its indicated position; insertion brackets leave the running text.}
\psentry{ps_134-5_008-v0003}{3}{deletion}{\deleted{in \xout{situations}}}{}{The author cancels this material by hand; it leaves the corrected running text.}
\psentry{ps_134-5_008-v0004}{4}{deletion}{\deleted{\hand{\unsure{progressively}}}}{}{The author cancels this material by hand; it leaves the corrected running text.}
\psentry{ps_134-5_008-v0005}{5}{typewriter deletion}{\xout{convincing myself only groupoids would fit nicely.}}{}{The typewriter strikes out this material; it leaves the corrected running text.}
\psentry{ps_134-5_008-v0006}{6}{typewriter deletion}{\xout{only}}{}{The typewriter strikes out this material; it leaves the corrected running text.}
\psentry{ps_134-5_008-v0007}{7}{typewriter deletion}{\xout{indeed}}{}{The typewriter strikes out this material; it leaves the corrected running text.}
\psentry{ps_134-5_008-v0008}{8}{typed insertion}{\inserted[typed]{natural}}{natural}{The typed insertion is incorporated at its indicated position; insertion brackets leave the running text.}
\psentry{ps_134-5_008-v0009}{9}{typewriter deletion}{\xout{ly}}{}{The typewriter strikes out this material; it leaves the corrected running text.}
\psentry{ps_134-5_008-v0010}{10}{typewriter deletion}{\xout{multiplicities}}{}{The typewriter strikes out this material; it leaves the corrected running text.}
\psentry{ps_134-5_008-v0011}{11}{insertion}{\inserted{adequately}}{\hand{adequately}}{The handwritten insertion is incorporated at its indicated position; insertion brackets leave the running text.}
\psentry{ps_134-5_008-v0012}{12}{deletion}{\deleted{which}}{}{The author cancels this material by hand; it leaves the corrected running text.}
\psentry{ps_134-5_008-v0013}{13}{insertion}{\inserted{The latter of course}}{\hand{The latter of course}}{The handwritten insertion is incorporated at its indicated position; insertion brackets leave the running text.}
\psentry{ps_134-5_008-v0014}{14}{insertion}{\inserted{just}}{\hand{just}}{The handwritten insertion is incorporated at its indicated position; insertion brackets leave the running text.}
\psentry{ps_134-5_008-v0015}{15}{insertion}{\inserted{of}}{\hand{of}}{The handwritten insertion is incorporated at its indicated position; insertion brackets leave the running text.}
\psentry{ps_134-5_008-v0016}{16}{insertion}{\inserted{this}}{\hand{this}}{The handwritten insertion is incorporated at its indicated position; insertion brackets leave the running text.}
\psentry{ps_134-5_008-v0017}{17}{typed insertion}{\inserted[typed]{tightly}}{tightly}{The typed insertion is incorporated at its indicated position; insertion brackets leave the running text.}
\psapparatusend
\fi
\pssource{ps_134-5_009}{134-5/9}{https://grothendieck.umontpellier.fr/web/viewer.html?file=../134-5.pdf\#page=10}
\psnumber{ps_134-5_009-v0001}{1}
\psnumber{ps_134-5_009-v0002}{2}
\psnumber{ps_134-5_009-v0003}{3}
\psnumber{ps_134-5_009-v0004}{4}
\psnumber{ps_134-5_009-v0005}{5}
\psnumber{ps_134-5_009-v0006}{6}
\psnumber{ps_134-5_009-v0007}{7}
\psnumber{ps_134-5_009-v0008}{8}
\psnumber{ps_134-5_009-v0009}{9}
\psnumber{ps_134-5_009-v0010}{10}
\psnumber{ps_134-5_009-v0011}{11}
\psnumber{ps_134-5_009-v0012}{12}
\psnumber{ps_134-5_009-v0013}{13}
\psnumber{ps_134-5_009-v0014}{14}
\psnumber{ps_134-5_009-v0015}{15}
\psnumber{ps_134-5_009-v0016}{16}
\psnumber{ps_134-5_009-v0017}{17}
\psnumber{ps_134-5_009-v0018}{18}
%% BEGIN SOURCE ps_134-5_009
% Pursuing Stacks (cotes 134-2 to 134-8): archive cote 134-5, site page 9 (PDF page 10), typescript page 196.
% typescript. Draft compared: ps4.tex:373-441.
% AI-transcribed by Claude Opus 5.5 (claude-opus-5-5, adaptive thinking, effort high), Message Batches API, batch ps2 (msgbatch_01NVwr1Yuo8qJf5rykeGBXPK), 2026-11-12; reviewed against both supplied scan halves; see CHECK_DECISIONS.tsv.
\scanpage{134-5/9}
\mspage{196}

29.5.

69. Before resuming \var{\inserted{more}}{\hand{more}}{insertion}{ps_134-5_009-v0001} \var{\deleted{with}}{}{deletion}{ps_134-5_009-v0002} technical work again, I would like to have a
short retrospective of the last six weeks' scratchwork, now lying on my
desk as a thickly bunch of scratchnotes\hand{,} nobody but I could possibly make
any sense of. \hand{//}The first thing I had on my mind has been there now for
nearly twenty years - ever since it had become clear, in the SGA 5
seminar on L-functions and a propos the formalism of traces in terms of
derived categories, that Verdier's set-up of derived categories was un-%
sufficient for formulating adeq\var{\over{\unsure{ua}}{ua}}{\hand{ua}}{writing over}{ps_134-5_009-v0003}tely \var{\deleted{even}}{}{deletion}{ps_134-5_009-v0004} some rather evident situati-%
ons\var{\xout{,}}{}{typewriter deletion}{ps_134-5_009-v0005} and relationships, such as the addition formu\hand{l}a for traces, or the
multiplicative formula for determinants. It then became appearant that
the derived category of an abelian category (say) was to\hand{o} coarse an object
in various respects, that it had to be complemented by similar``triangu-%
lated categories''(such as the \var{\xout{\illegible{1}}}{}{typewriter deletion}{ps_134-5_009-v0006} derived category of a suitable category
of ``triangles'' of complexes, or t\var{\over{\unsure{h}}{h}}{\hand{h}}{writing over}{ps_134-5_009-v0007}e whole bunch of derived categories
of categories of filtered complexes of order n with variable n), closely
connected to it. Deligne and Illusie had both set out, independently,
to work out \var{\deleted{a}}{}{deletion}{ps_134-5_009-v0008}\var{\inserted{some}}{\hand{some}}{insertion}{ps_134-5_009-v0009} set-up meeting the most urgent requirements\hand{.} (Illusie's
treatment in terms of filtered complexes was written down and published
in his thesis six years later (Springer Lecture Notes N\textsuperscript{\underline{o}} 239)). While
adequate for the main tasks then at hand, neither treatment was really
wholly satisfactory to my taste. One main feature I believe making me
feel unconfortable, was that the extra categories which had to be intro-%
duced,to \var{\xout{compl\unsure{ement}}}{}{typewriter deletion}{ps_134-5_009-v0010} \var{\inserted[typed]{round up some}}{round up some}{typed insertion}{ps_134-5_009-v0011}\var{\inserted{what}}{\hand{what}}{insertion}{ps_134-5_009-v0012} a stripped-and-naked triangulated categor\var{\over{\unsure{ies}}{y}}{\hand{y}}{writing over}{ps_134-5_009-v0013}\hand{,}were
triangulated categories in their own right, in Verdier's sense, but
remaining nearly as stripped \var{\xout{as the}}{}{typewriter deletion}{ps_134-5_009-v0014} by themselves as the initial trian-%
gulated category they were intended to provide clothing for. In other
words, \unsure{theye|there} was a lack of inner stability in the formalism, making it
appear as very much provisional still. Also, while interested in asso-%
ciating to an abelian category a handy sequence of ``filtered derived
categories'', Illusie made no attempt to pin down what exactly the
inner structure of the object he had arrived at was - unlike Verdier,
who \var{\inserted[typed]{had}}{had}{typed insertion}{ps_134-5_009-v0015} introduced, alongside with the notion of a derived category of an
abelian category, a general notion of triangulated categories, into
which these derived categories would fit. The obvious idea which was in
my head by then for avoiding such shortcomings, was \var{\xout{to associate to}}{}{typewriter deletion}{ps_134-5_009-v0016}\var{\inserted[typed]{that}}{that}{typed insertion}{ps_134-5_009-v0017} an
abelian category $\underline{A}$ gave rise, not only to the \var{\xout{usual}}{}{typewriter deletion}{ps_134-5_009-v0018} single usual
derived category D($\underline{A}$) of Verdier, but also, for every type of diagrams,

\margin{9}

%% END SOURCE ps_134-5_009
\ifnum\EditionMode=1
\psapparatusstart
\psentry{ps_134-5_009-v0001}{1}{insertion}{\inserted{more}}{\hand{more}}{The handwritten insertion is incorporated at its indicated position; insertion brackets leave the running text.}
\psentry{ps_134-5_009-v0002}{2}{deletion}{\deleted{with}}{}{The author cancels this material by hand; it leaves the corrected running text.}
\psentry{ps_134-5_009-v0003}{3}{writing over}{\over{\unsure{ua}}{ua}}{\hand{ua}}{The superposed handwritten reading replaces the underlying reading.}
\psentry{ps_134-5_009-v0004}{4}{deletion}{\deleted{even}}{}{The author cancels this material by hand; it leaves the corrected running text.}
\psentry{ps_134-5_009-v0005}{5}{typewriter deletion}{\xout{,}}{}{The typewriter strikes out this material; it leaves the corrected running text.}
\psentry{ps_134-5_009-v0006}{6}{typewriter deletion}{\xout{\illegible{1}}}{}{The typewriter strikes out this material; it leaves the corrected running text.}
\psentry{ps_134-5_009-v0007}{7}{writing over}{\over{\unsure{h}}{h}}{\hand{h}}{The superposed handwritten reading replaces the underlying reading.}
\psentry{ps_134-5_009-v0008}{8}{deletion}{\deleted{a}}{}{The author cancels this material by hand; it leaves the corrected running text.}
\psentry{ps_134-5_009-v0009}{9}{insertion}{\inserted{some}}{\hand{some}}{The handwritten insertion is incorporated at its indicated position; insertion brackets leave the running text.}
\psentry{ps_134-5_009-v0010}{10}{typewriter deletion}{\xout{compl\unsure{ement}}}{}{The typewriter strikes out this material; it leaves the corrected running text.}
\psentry{ps_134-5_009-v0011}{11}{typed insertion}{\inserted[typed]{round up some}}{round up some}{The typed insertion is incorporated at its indicated position; insertion brackets leave the running text.}
\psentry{ps_134-5_009-v0012}{12}{insertion}{\inserted{what}}{\hand{what}}{The handwritten insertion is incorporated at its indicated position; insertion brackets leave the running text.}
\psentry{ps_134-5_009-v0013}{13}{writing over}{\over{\unsure{ies}}{y}}{\hand{y}}{The superposed handwritten reading replaces the underlying reading.}
\psentry{ps_134-5_009-v0014}{14}{typewriter deletion}{\xout{as the}}{}{The typewriter strikes out this material; it leaves the corrected running text.}
\psentry{ps_134-5_009-v0015}{15}{typed insertion}{\inserted[typed]{had}}{had}{The typed insertion is incorporated at its indicated position; insertion brackets leave the running text.}
\psentry{ps_134-5_009-v0016}{16}{typewriter deletion}{\xout{to associate to}}{}{The typewriter strikes out this material; it leaves the corrected running text.}
\psentry{ps_134-5_009-v0017}{17}{typed insertion}{\inserted[typed]{that}}{that}{The typed insertion is incorporated at its indicated position; insertion brackets leave the running text.}
\psentry{ps_134-5_009-v0018}{18}{typewriter deletion}{\xout{usual}}{}{The typewriter strikes out this material; it leaves the corrected running text.}
\psapparatusend
\fi
\pssource{ps_134-5_010}{134-5/10}{https://grothendieck.umontpellier.fr/web/viewer.html?file=../134-5.pdf\#page=11}
\psnumber{ps_134-5_010-v0001}{1}
\psnumber{ps_134-5_010-v0002}{2}
\psnumber{ps_134-5_010-v0003}{3}
\psnumber{ps_134-5_010-v0004}{4}
\psnumber{ps_134-5_010-v0005}{5}
\psnumber{ps_134-5_010-v0006}{6}
\psnumber{ps_134-5_010-v0007}{7}
\psnumber{ps_134-5_010-v0008}{8}
\psnumber{ps_134-5_010-v0009}{9}
\psnumber{ps_134-5_010-v0010}{10}
\psnumber{ps_134-5_010-v0011}{11}
\psnumber{ps_134-5_010-v0012}{12}
\psnumber{ps_134-5_010-v0013}{13}
\psnumber{ps_134-5_010-v0014}{14}
\psnumber{ps_134-5_010-v0015}{15}
\psnumber{ps_134-5_010-v0016}{16}
\psnumber{ps_134-5_010-v0017}{17}
%% BEGIN SOURCE ps_134-5_010
% Pursuing Stacks (cotes 134-2 to 134-8): archive cote 134-5, site page 10 (PDF page 11), typescript page 197.
% typescript. Draft compared: ps4.tex:425-486.
% AI-transcribed by Claude Opus 5.5 (claude-opus-5-5, adaptive thinking, effort high), Message Batches API, batch ps2 (msgbatch_01NVwr1Yuo8qJf5rykeGBXPK), 2026-11-12; reviewed against both supplied scan halves; see CHECK_DECISIONS.tsv.
\scanpage{134-5/10}
\mspage{197}
\var{\inserted{to}}{\hand{to}}{insertion}{ps_134-5_010-v0001} the derived category of the abelian category of all \underline{A}-valued diagrams
of this type. In precise terms, for any small category I, we get
the categ\var{\inserted[typed]{o}}{o}{typed insertion}{ps_134-5_010-v0002}ry D(\underline{Hom}(I,\underline{A})), depending functorially in a contravariant
way on I. Rewriting this category $D_{\underline{A}}(I)$ say, the idea was to consider
(*) \[ I \hand{\mapsto} D_{\underline{A}}(I) \quad , \]
possibly \var{\inserted[typed]{with I suitably}}{with I suitably}{typed insertion}{ps_134-5_010-v0003} restricted \var{\xout{xx}}{}{typewriter deletion}{ps_134-5_010-v0004} (for instance to finite categories, or to finite
ordered sets, corresponding to finite \uline{commutative} diagrams), as embo-%
dying the "full" triangulated structure defined by \underline{A} . This of course
at once raises a number of questions, such as recovering the usual
triangulated structure of $D(\underline{A}) = D_{\underline{A}}(e)$ ( e the final object of (Cat))
in terms of (*), and pinning down too the relevant formal properties
(and possibly even extra structure) one had to assume on (*). I had
never so far taken the time to sit down and play around some and see
how this goes through, expecting that surely someone else would do it
some day and I would be informed - but apparently in the last eighteen
years nobody ever was interested. Also, it had been rather clear from
the start that Verdier's constructions could be adapted and \var{\inserted{did}}{\hand{did}}{insertion}{ps_134-5_010-v0005} make sense
for non-commutative homotopy set-ups, which was also appearant in betw-%
ween the lines in Gabriel-Zisman's \var{\xout{xxxxxxxxxf}}{}{typewriter deletion}{ps_134-5_010-v0006} book on the foundations
of homotopy theory, and a lot more explicitly in Quillen's axiomatizati-%
on of homotopical algebra. This \var{\inserted{axiomatization}}{\hand{axiomatization}}{insertion}{ps_134-5_010-v0007} I found very appealing indeed - and
right now \var{\inserted{still}}{\hand{still}}{insertion}{ps_134-5_010-v0008} his little book is my most congenial and main source of infor-%
mation on foundational matters of homotopical algebra. I reme\var{\inserted[typed]{m}}{m}{typed insertion}{ps_134-5_010-v0009}ber though
\var{\inserted[typed]{my}}{my}{typed insertion}{ps_134-5_010-v0010} being a little disappointed at Quillen's not caring either to pursue the
matter of what exactly a "non commutative triangulated category structu\var{\xout{\unsure{re}}}{}{typewriter deletion}{ps_134-5_010-v0011}-%
re" (of the type he was getting from his model categories) was, just
contenting himself tomumble a few words about existence of "higher
structure" (t\var{\over{hen}{han}}{\hand{han}}{writing over}{ps_134-5_010-v0012} just the Dold-Puppe sequences), which \var{\inserted{(he implied)}}{\hand{(he implied)}}{insertion}{ps_134-5_010-v0013} needed to be under-%
stood. I felt of course that presumably the variance formalism (*) should
furnish any kind of "higher" structure one was looking for, but it was'
nt really my business to check.

It still is'nt, however I did some homework on (*) - it was the
first thing indeed I looked at in these six weeks, and some main featu-%
res came out very readily indeed. It turns out that the main formal
variance property to demand on (*), presumably even the only one , is
that for a given map \var{\inserted{$f$:}}{\hand{$f$:}}{insertion}{ps_134-5_010-v0014} I $\hand{\to}$ J on the indexing categories of diagram-types
I and J , the corresponding functor
\[ \hand{f^*:} D_{\var{\deleted{\underline{A}}}{}{deletion}{ps_134-5_010-v0015}}(J) \hand{\longrightarrow} D_{\var{\deleted{\underline{A}}}{}{deletion}{ps_134-5_010-v0016}}(I) \]
\hand{\uline{\textcolor{black}{should have both a left and a right adjoint}}}, \var{\inserted{say $f_!$ and $f_*$.}}{\hand{say $f_!$ and $f_*$.}}{insertion}{ps_134-5_010-v0017} In case J= e , the two
\margin{10}

%% END SOURCE ps_134-5_010
\ifnum\EditionMode=1
\psapparatusstart
\psentry{ps_134-5_010-v0001}{1}{insertion}{\inserted{to}}{\hand{to}}{The handwritten insertion is incorporated at its indicated position; insertion brackets leave the running text.}
\psentry{ps_134-5_010-v0002}{2}{typed insertion}{\inserted[typed]{o}}{o}{The typed insertion is incorporated at its indicated position; insertion brackets leave the running text.}
\psentry{ps_134-5_010-v0003}{3}{typed insertion}{\inserted[typed]{with I suitably}}{with I suitably}{The typed insertion is incorporated at its indicated position; insertion brackets leave the running text.}
\psentry{ps_134-5_010-v0004}{4}{typewriter deletion}{\xout{xx}}{}{The typewriter strikes out this material; it leaves the corrected running text.}
\psentry{ps_134-5_010-v0005}{5}{insertion}{\inserted{did}}{\hand{did}}{The handwritten insertion is incorporated at its indicated position; insertion brackets leave the running text.}
\psentry{ps_134-5_010-v0006}{6}{typewriter deletion}{\xout{xxxxxxxxxf}}{}{The typewriter strikes out this material; it leaves the corrected running text.}
\psentry{ps_134-5_010-v0007}{7}{insertion}{\inserted{axiomatization}}{\hand{axiomatization}}{The handwritten insertion is incorporated at its indicated position; insertion brackets leave the running text.}
\psentry{ps_134-5_010-v0008}{8}{insertion}{\inserted{still}}{\hand{still}}{The handwritten insertion is incorporated at its indicated position; insertion brackets leave the running text.}
\psentry{ps_134-5_010-v0009}{9}{typed insertion}{\inserted[typed]{m}}{m}{The typed insertion is incorporated at its indicated position; insertion brackets leave the running text.}
\psentry{ps_134-5_010-v0010}{10}{typed insertion}{\inserted[typed]{my}}{my}{The typed insertion is incorporated at its indicated position; insertion brackets leave the running text.}
\psentry{ps_134-5_010-v0011}{11}{typewriter deletion}{\xout{\unsure{re}}}{}{The typewriter strikes out this material; it leaves the corrected running text.}
\psentry{ps_134-5_010-v0012}{12}{writing over}{\over{hen}{han}}{\hand{han}}{The superposed handwritten reading replaces the underlying reading.}
\psentry{ps_134-5_010-v0013}{13}{insertion}{\inserted{(he implied)}}{\hand{(he implied)}}{The handwritten insertion is incorporated at its indicated position; insertion brackets leave the running text.}
\psentry{ps_134-5_010-v0014}{14}{insertion}{\inserted{$f$:}}{\hand{$f$:}}{The handwritten insertion is incorporated at its indicated position; insertion brackets leave the running text.}
\psentry{ps_134-5_010-v0015}{15}{deletion}{$\deleted{\underline{A}}$}{}{The author cancels this material by hand; it leaves the corrected running text.}
\psentry{ps_134-5_010-v0016}{16}{deletion}{$\deleted{\underline{A}}$}{}{The author cancels this material by hand; it leaves the corrected running text.}
\psentry{ps_134-5_010-v0017}{17}{insertion}{\inserted{say $f_!$ and $f_*$.}}{\hand{say $f_!$ and $f_*$.}}{The handwritten insertion is incorporated at its indicated position; insertion brackets leave the running text.}
\psapparatusend
\fi
\pssource{ps_134-5_011}{134-5/11}{https://grothendieck.umontpellier.fr/web/viewer.html?file=../134-5.pdf\#page=12}
\psnumber{ps_134-5_011-v0001}{1}
\psnumber{ps_134-5_011-v0002}{2}
\psnumber{ps_134-5_011-v0003}{3}
\psnumber{ps_134-5_011-v0004}{4}
\psnumber{ps_134-5_011-v0005}{5}
\psnumber{ps_134-5_011-v0006}{6}
\psnumber{ps_134-5_011-v0007}{7}
\psnumber{ps_134-5_011-v0008}{8}
\psnumber{ps_134-5_011-v0009}{9}
\psnumber{ps_134-5_011-v0010}{10}
\psnumber{ps_134-5_011-v0011}{11}
\psnumber{ps_134-5_011-v0012}{12}
\psnumber{ps_134-5_011-v0013}{13}
\psnumber{ps_134-5_011-v0014}{14}
\psnumber{ps_134-5_011-v0015}{15}
\psnumber{ps_134-5_011-v0016}{16}
\psnumber{ps_134-5_011-v0017}{17}
\psnumber{ps_134-5_011-v0018}{18}
%% BEGIN SOURCE ps_134-5_011
% Pursuing Stacks (cotes 134-2 to 134-8): archive cote 134-5, site page 11 (PDF page 12), typescript page 198.
% typescript. Draft compared: ps4.tex:473-539.
% AI-transcribed by Claude Opus 5.5 (claude-opus-5-5, adaptive thinking, effort high), Message Batches API, batch ps2 (msgbatch_01NVwr1Yuo8qJf5rykeGBXPK), 2026-11-12; reviewed against both supplied scan halves; see CHECK_DECISIONS.tsv.
\scanpage{134-5/11}
\mspage{198}

functors we get from D(I)  to  D(e)\var{\inserted{$=\mathbb{D}$}}{\hand{$=\mathbb{D}$}}{insertion}{ps_134-5_011-v0001}( the ``stripped'' triangulated catego-%
ry) can be viewed as \hand{\uline{\textcolor{black}{a substitute for taking}}},  \hand{\uline{\textcolor{black}{respectively}}},  \hand{\uline{\textcolor{black}{direct and}}}
\hand{\uline{\textcolor{black}{inverse limits in T}}} (for a system of objects indexed by I),  \hand{\uline{\textcolor{black}{which in}}}
\hand{\uline{\textcolor{black}{the usual sense dont generally exist in}}} \hand{$\mathbb{D}_{o}$} (except just finite sums and
products).  These operations admit as important special cases (when I
is either\var{\inserted{one}}{\hand{one}}{insertion}{ps_134-5_011-v0002} of the two\var{\inserted{mutually dual}}{\hand{mutually dual}}{insertion}{ps_134-5_011-v0003} categories

\margin{(**)}
% Hand-drawn arrows are marked blue; typed objects remain plain.
\diagram{
\begin{tikzcd}[column sep=small,row sep=tiny]
 \& b \\ a \ar[color=blue, ur] \ar[color=blue, dr] \& \\ \& c
\end{tikzcd}
\quad\text{or}\quad
\begin{tikzcd}[column sep=small,row sep=tiny]
 \& b \ar[color=blue, dl] \\ a \& \\ \& c \ar[color=blue, ul]
\end{tikzcd}
\quad ,
}

the operation of \hand{(}bi\unsure{n}ary\hand{)} amalgamated sums or fibered products,  and
hence also of taking ``cofibers'' and \var{\xout{f}}{}{typewriter deletion}{ps_134-5_011-v0004} ``fibers'' of maps,  in the \var{\xout{Sxxxxxx}}{}{typewriter deletion}{ps_134-5_011-v0005}
\var{\xout{Sxxxx}}{}{typewriter deletion}{ps_134-5_011-v0006} sense introduced by Cartan-Serre in homotopy theory about thirty
years ago.  I also checked that the two mutually dual Dold-Puppe sequen-%
ces follow quite formally from the set-up.  One just has to fit in a
suitable extra axiom to insure the usual exactness properties for these
sequences.

Except in the commutative case when starting with an abelian cate-%
gory\var{\inserted[typed]{as above}}{as above}{typed insertion}{ps_134-5_011-v0007}, I did not check however that \var{\xout{xxx}}{}{typewriter deletion}{ps_134-5_011-v0008} there is indeed such ``higher
variance structure'' in the usual cases\hand{,} when a typical ``triangulated
category'' in some sense or other turns up,  for instance from a model
category in Quillen's sense.  What I did check \var{\deleted{in}}{}{deletion}{ps_134-5_011-v0009}\var{\inserted[typed]{though}}{though}{typed insertion}{ps_134-5_011-v0010} this last case,  under
a mild additional assumption which seems verified in all practical cases,
is the existence of the operation  $f_! = \hand{\int}_I$  (``integration'' ) and
$f_* = \hand{\prod}_I$(cointegration) for the special case $f\colon I \hand{\longrightarrow} e$,  when I is either
one of the two categories\var{\inserted{(**)}}{\hand{(**)}}{insertion}{ps_134-5_011-v0011} above.  I expect that working some more,  one
should get under the same assumptions \var{\deleted{the \xout{xxxx}}}{}{deletion}{ps_134-5_011-v0012} at least the existence of
$f_!$ and $f_{\hand{*}}$  for any map $f\colon I \hand{\longrightarrow} J$ between\var{\inserted{finite}}{\hand{finite}}{insertion}{ps_134-5_011-v0013} ordered sets.

My main interest of course at present is in the category (Hot)
itself,  more generally in Hot($\underline{W}$) = $\underline{W}^{-1}$(Cat),  where $\underline{W}$ is a ``basic loca-%
lizer''.  More generally,  if (M,W)  is any modelizer say,  the natural
thing to do,  paraphrasing (*),  is for any indexing category I to
endow  $\underline{\mathrm{Hom}}$(I,M)  with the set of arrows $W_I$  defined by componentwise
belonging to W ,  and to define
\[ D_{(M,W)}(I) = D(I) \overset{\mathrm{def}}{=} W_I^{-1}\underline{\mathrm{Hom}}(I,M) \quad , \]
with the obvious contravariant dependance on I,  denoted by $f^*$  for
$f\colon I \hand{\rightarrow} J$.  The question  then arises as to the existence of left and
right adjoints,  $f_!$  and $f_*$ .  In case we take M = (Cat),  the existence
of \var{\xout{xx}}{}{typewriter deletion}{ps_134-5_011-v0014} $f_!$  goes through with am\var{\inserted[typed]{a}}{a}{typed insertion}{ps_134-5_011-v0015}\unsure{z}ing smoothness :  interpreting a\var{\deleted{n}}{}{deletion}{ps_134-5_011-v0016}\var{\inserted[typed]{``model''}}{``model''}{typed insertion}{ps_134-5_011-v0017} object
of  $\underline{\mathrm{Hom}}$(I,(Cat))  \var{\xout{in terms of a}}{}{typewriter deletion}{ps_134-5_011-v0018} namely a functor
\[ I \hand{\longrightarrow} (Cat) \]

\margin{11}

%% END SOURCE ps_134-5_011
\ifnum\EditionMode=1
\psapparatusstart
\psentry{ps_134-5_011-v0001}{1}{insertion}{\inserted{$=\mathbb{D}$}}{\hand{$=\mathbb{D}$}}{The handwritten insertion is incorporated at its indicated position; insertion brackets leave the running text.}
\psentry{ps_134-5_011-v0002}{2}{insertion}{\inserted{one}}{\hand{one}}{The handwritten insertion is incorporated at its indicated position; insertion brackets leave the running text.}
\psentry{ps_134-5_011-v0003}{3}{insertion}{\inserted{mutually dual}}{\hand{mutually dual}}{The handwritten insertion is incorporated at its indicated position; insertion brackets leave the running text.}
\psentry{ps_134-5_011-v0004}{4}{typewriter deletion}{\xout{f}}{}{The typewriter strikes out this material; it leaves the corrected running text.}
\psentry{ps_134-5_011-v0005}{5}{typewriter deletion}{\xout{Sxxxxxx}}{}{The typewriter strikes out this material; it leaves the corrected running text.}
\psentry{ps_134-5_011-v0006}{6}{typewriter deletion}{\xout{Sxxxx}}{}{The typewriter strikes out this material; it leaves the corrected running text.}
\psentry{ps_134-5_011-v0007}{7}{typed insertion}{\inserted[typed]{as above}}{as above}{The typed insertion is incorporated at its indicated position; insertion brackets leave the running text.}
\psentry{ps_134-5_011-v0008}{8}{typewriter deletion}{\xout{xxx}}{}{The typewriter strikes out this material; it leaves the corrected running text.}
\psentry{ps_134-5_011-v0009}{9}{deletion}{\deleted{in}}{}{The author cancels this material by hand; it leaves the corrected running text.}
\psentry{ps_134-5_011-v0010}{10}{typed insertion}{\inserted[typed]{though}}{though}{The typed insertion is incorporated at its indicated position; insertion brackets leave the running text.}
\psentry{ps_134-5_011-v0011}{11}{insertion}{\inserted{(**)}}{\hand{(**)}}{The handwritten insertion is incorporated at its indicated position; insertion brackets leave the running text.}
\psentry{ps_134-5_011-v0012}{12}{deletion}{\deleted{the \xout{xxxx}}}{}{The author cancels this material by hand; it leaves the corrected running text.}
\psentry{ps_134-5_011-v0013}{13}{insertion}{\inserted{finite}}{\hand{finite}}{The handwritten insertion is incorporated at its indicated position; insertion brackets leave the running text.}
\psentry{ps_134-5_011-v0014}{14}{typewriter deletion}{\xout{xx}}{}{The typewriter strikes out this material; it leaves the corrected running text.}
\psentry{ps_134-5_011-v0015}{15}{typed insertion}{\inserted[typed]{a}}{a}{The typed insertion is incorporated at its indicated position; insertion brackets leave the running text.}
\psentry{ps_134-5_011-v0016}{16}{deletion}{\deleted{n}}{}{The author cancels this material by hand; it leaves the corrected running text.}
\psentry{ps_134-5_011-v0017}{17}{typed insertion}{\inserted[typed]{``model''}}{``model''}{The typed insertion is incorporated at its indicated position; insertion brackets leave the running text.}
\psentry{ps_134-5_011-v0018}{18}{typewriter deletion}{\xout{in terms of a}}{}{The typewriter strikes out this material; it leaves the corrected running text.}
\psapparatusend
\fi
\pssource{ps_134-5_012}{134-5/12}{https://grothendieck.umontpellier.fr/web/viewer.html?file=../134-5.pdf\#page=13}
\psnumber{ps_134-5_012-v0001}{1}
\psnumber{ps_134-5_012-v0002}{2}
\psnumber{ps_134-5_012-v0003}{3}
\psnumber{ps_134-5_012-v0004}{4}
\psnumber{ps_134-5_012-v0005}{5}
\psnumber{ps_134-5_012-v0006}{6}
\psnumber{ps_134-5_012-v0007}{7}
\psnumber{ps_134-5_012-v0008}{8}
\psnumber{ps_134-5_012-v0009}{9}
\psnumber{ps_134-5_012-v0010}{10}
\psnumber{ps_134-5_012-v0011}{11}
\psnumber{ps_134-5_012-v0012}{12}
\psnumber{ps_134-5_012-v0013}{13}
\psnumber{ps_134-5_012-v0014}{14}
\psnumber{ps_134-5_012-v0015}{15}
\psnumber{ps_134-5_012-v0016}{16}
%% BEGIN SOURCE ps_134-5_012
% Pursuing Stacks (cotes 134-2 to 134-8): archive cote 134-5, site page 12 (PDF page 13), typescript page 199.
% typescript. Draft compared: ps4.tex:525-585.
% AI-transcribed by Claude Opus 5.5 (claude-opus-5-5, adaptive thinking, effort high), Message Batches API, batch ps2 (msgbatch_01NVwr1Yuo8qJf5rykeGBXPK), 2026-11-12; reviewed against both supplied scan halves; see CHECK_DECISIONS.tsv.
\scanpage{134-5/12}
\mspage{199}

in terms of a cofibered category X over I
\[ p:\ X \hand{\longrightarrow} I \quad , \]
and assuming for simplicity f cofiber\var{\over{\unsure{ed}}{ing}}{\hand{ing}}{writing over}{ps_134-5_012-v0001} too, $f_{\hand{!}}(X)$ is just X itself,
the total category of the cofibering, viewed as a (cofibe\hand{red}) \hand{category}
over J by using the fu\hand{n}ctor $q = f\hand{\circ}\hand{p}$ ! This appl\hand{ies for instan}ce when J
is the final category, and yields the operation of``integration of
homotopy types'' $\hand{\int}_I$ , in terms of the total category of a \var{\xout{fib}}{}{typewriter deletion}{ps_134-5_012-v0002} cofibered
category over I. If we want to rid ourselves from any extra assumption
on f, we can describe D(I) (up to equivalence) in terms of the category
$(\mathrm{Cat})_{/I}$ of categories X over I (not necessarily cofibered over I), $W_I$
being replaced by the corresponding notion of ``$\underline{W}$-\uline{equiva\hand{l}ence\unsure{s}} \var{\xout{\unsure{over}}}{}{typewriter deletion}{ps_134-5_012-v0003} \uline{relati\hand{ve}
to I} '' for maps $u : X \hand{\to} Y$ of objects of (Cat) over I , by which we mean
a map u such that the localized maps
\[ u_{/i} \ :\ X_{/i} \hand{\longrightarrow} Y_{/i} \]
are in $\underline{W}$ , for any i in I. Regarding now any category X over I as a
category over J by means of $f\hand{\circ}p$ , this is \var{\inserted{clearly}}{\hand{clearly}}{insertion}{ps_134-5_012-v0004} compatible with the relative
weak equivalences $\underline{W}_I$ and $\underline{W}_J$ , and yields by localization the looked-for
functor $f_{\hand{!}}$ .

This amazingly simple construction and \unsure{interpretation|interpratation} of the basic
$f_{\hand{!}}$ and $\hand{\int}_I$ operations is one main reward, it appears, for working with
the ``basic modelizer'' (Cat), which in this occurence, as in the whole
test- and asphericity story, quite evidently deserves his name. It has
turned out \var{\inserted{since}}{\hand{since}}{insertion}{ps_134-5_012-v0005} that in some other respects - for instance, paradoxically,
when it comes to the question of relationship between this lofty integra-%
ion operation, and true honest amalgamated sums - the modelizers $A^{\wedge}$
associated to test categories \var{\inserted[typed]{$A$}}{$A$}{typed insertion}{ps_134-5_012-v0006} (\var{\inserted[typed]{namely}}{namely}{typed insertion}{ps_134-5_012-v0007} the so-called ``elementary modelizers'') are
more convenient tools then (Cat). Thus, it appears very doubtful still
that (Cat) is a ``model category'' in Quillen's sense, in any reasonable
way (with $\underline{W}$ of course as the set of ``weak equivalences''). I finally got
the feeling that \var{\xout{for}}{}{typewriter deletion}{ps_134-5_012-v0008} a good mastery of the basic aspects of homotopy types
and \var{\inserted[typed]{of}}{of}{typed insertion}{ps_134-5_012-v0009} \var{\xout{\illegible{2}}}{}{typewriter deletion}{ps_134-5_012-v0010} basic relationships among these, will require mainly
great ``aisance'' in playing \var{\xout{with}}{}{typewriter deletion}{ps_134-5_012-v0011} around with a number of available descrip-%
tions of homotopy types by models, no one among which (not even \hand{by} mode\hand{ls}
in (Cat), and surely \var{\xout{not}}{}{typewriter deletion}{ps_134-5_012-v0012} still less \var{\inserted{by}}{\hand{by}}{insertion}{ps_134-5_012-v0013} semisimplicial structures) being ade-%
quate for replacing all others.

As for the $f_{\hand{*}}$ and cointegration $\hand{\prod}_I$ operations among the catego-%
ries D(I), except in the very special case noted above (corresponding to
fibered products), I did not hit upon any ready-to-use candidate for it,
and I doubt there is any. I do believe the\var{\over{y}{se}}{\hand{se}}{writing over}{ps_134-5_012-v0014} \var{\inserted{operations}}{\hand{operations}}{insertion}{ps_134-5_012-v0015} exist indeed, and \var{\inserted{I}}{\hand{I}}{insertion}{ps_134-5_012-v0016} even have

\margin{12}

%% END SOURCE ps_134-5_012
\ifnum\EditionMode=1
\psapparatusstart
\psentry{ps_134-5_012-v0001}{1}{writing over}{\over{\unsure{ed}}{ing}}{\hand{ing}}{The superposed handwritten reading replaces the underlying reading.}
\psentry{ps_134-5_012-v0002}{2}{typewriter deletion}{\xout{fib}}{}{The typewriter strikes out this material; it leaves the corrected running text.}
\psentry{ps_134-5_012-v0003}{3}{typewriter deletion}{\xout{\unsure{over}}}{}{The typewriter strikes out this material; it leaves the corrected running text.}
\psentry{ps_134-5_012-v0004}{4}{insertion}{\inserted{clearly}}{\hand{clearly}}{The handwritten insertion is incorporated at its indicated position; insertion brackets leave the running text.}
\psentry{ps_134-5_012-v0005}{5}{insertion}{\inserted{since}}{\hand{since}}{The handwritten insertion is incorporated at its indicated position; insertion brackets leave the running text.}
\psentry{ps_134-5_012-v0006}{6}{typed insertion}{\inserted[typed]{$A$}}{$A$}{The typed insertion is incorporated at its indicated position; insertion brackets leave the running text.}
\psentry{ps_134-5_012-v0007}{7}{typed insertion}{\inserted[typed]{namely}}{namely}{The typed insertion is incorporated at its indicated position; insertion brackets leave the running text.}
\psentry{ps_134-5_012-v0008}{8}{typewriter deletion}{\xout{for}}{}{The typewriter strikes out this material; it leaves the corrected running text.}
\psentry{ps_134-5_012-v0009}{9}{typed insertion}{\inserted[typed]{of}}{of}{The typed insertion is incorporated at its indicated position; insertion brackets leave the running text.}
\psentry{ps_134-5_012-v0010}{10}{typewriter deletion}{\xout{\illegible{2}}}{}{The typewriter strikes out this material; it leaves the corrected running text.}
\psentry{ps_134-5_012-v0011}{11}{typewriter deletion}{\xout{with}}{}{The typewriter strikes out this material; it leaves the corrected running text.}
\psentry{ps_134-5_012-v0012}{12}{typewriter deletion}{\xout{not}}{}{The typewriter strikes out this material; it leaves the corrected running text.}
\psentry{ps_134-5_012-v0013}{13}{insertion}{\inserted{by}}{\hand{by}}{The handwritten insertion is incorporated at its indicated position; insertion brackets leave the running text.}
\psentry{ps_134-5_012-v0014}{14}{writing over}{\over{y}{se}}{\hand{se}}{The superposed handwritten reading replaces the underlying reading.}
\psentry{ps_134-5_012-v0015}{15}{insertion}{\inserted{operations}}{\hand{operations}}{The handwritten insertion is incorporated at its indicated position; insertion brackets leave the running text.}
\psentry{ps_134-5_012-v0016}{16}{insertion}{\inserted{I}}{\hand{I}}{The handwritten insertion is incorporated at its indicated position; insertion brackets leave the running text.}
\psapparatusend
\fi
\pssource{ps_134-5_013}{134-5/13}{https://grothendieck.umontpellier.fr/web/viewer.html?file=../134-5.pdf\#page=14}
\psnumber{ps_134-5_013-v0001}{1}
\psnumber{ps_134-5_013-v0002}{2}
\psnumber{ps_134-5_013-v0003}{3}
\psnumber{ps_134-5_013-v0004}{4}
\psnumber{ps_134-5_013-v0005}{5}
\psnumber{ps_134-5_013-v0006}{6}
\psnumber{ps_134-5_013-v0007}{7}
\psnumber{ps_134-5_013-v0008}{8}
\psnumber{ps_134-5_013-v0009}{9}
\psnumber{ps_134-5_013-v0010}{10}
\psnumber{ps_134-5_013-v0011}{11}
\psnumber{ps_134-5_013-v0012}{12}
%% BEGIN SOURCE ps_134-5_013
% Pursuing Stacks (cotes 134-2 to 134-8): archive cote 134-5, site page 13 (PDF page 14), typescript page 200.
% typescript. Draft compared: ps4.tex:570-626.
% AI-transcribed by Claude Opus 5.5 (claude-opus-5-5, adaptive thinking, effort high), Message Batches API, batch ps2 (msgbatch_01NVwr1Yuo8qJf5rykeGBXPK), 2026-11-12; reviewed against both supplied scan halves; see CHECK_DECISIONS.tsv.
\scanpage{134-5/13}
\mspage{200}
\var{\xout{axratherxgeneralx}}{}{typewriter deletion}{ps_134-5_013-v0001}
\var{\xout{setx\allowbreak{}upx\allowbreak{}inx\allowbreak{}mindx\allowbreak{}(includingx\allowbreak{}Quillen'sx\allowbreak{})x\allowbreak{}x\allowbreak{}withx\allowbreak{}thex\allowbreak{}"mildx\allowbreak{}restriction"x\allowbreak{})x\allowbreak{}}}{}{typewriter deletion}{ps_134-5_013-v0002} in
\var{\xout{which}}{}{typewriter deletion}{ps_134-5_013-v0003} mind a rather general condition on a pair (M,W) with W\hand{$\subset$}\,Fl(M), for
both basic operations $f_!$ and $f_*$ to exist between the \unsure{carresponding|corresponding} cate\var{\deleted{t}}{}{deletion}{ps_134-5_013-v0004}-%
gories $D_{(M,W)}(I)$ -- but to establish this \var{\deleted{fact}}{}{deletion}{ps_134-5_013-v0005} \var{\inserted{expectation,}}{\hand{expectation,}}{insertion}{ps_134-5_013-v0006} may require a good amount
of work. I'll come back upon these matters in due course.

There arises of course the question of giving a suitable name to the
structure $I \hand{\mapsto} D(I)$ I arrived at, which seems to embody at least some
main features of a satisfactory notion of a ``triangulated category'' (not
necessarily commutative), gradually emerging from darkness. I have
thought of calling such a structure a ``\hand{\uline{\textcolor{black}{derivator}}}'', with the implication
that it's main function is to furnish us with a somehow\hand{,} ``complete'' bunch
(in terms of a rounded-up self-contained formalism) of categories D(I),
which are being looked at as ``\hand{\uline{\textcolor{black}{derived}}} \hand{\uline{\textcolor{black}{categories}}}'' in some sense or other.
The only way I know of for constructing such derivator, is as above in
terms of a pair (M,W), submitted to suitable conditions for ensuring
existence of $f_!$ and $f_*$ , at least \hand{when} f \var{\inserted{is}}{\hand{is}}{insertion}{ps_134-5_013-v0007} any map between finite ordered
sets. We may look upon D(I) as a \var{\xout{sub}}{}{typewriter deletion}{ps_134-5_013-v0008} refinement and substitue for the
notion of family of objects of $D(e) = D_o$ indexed by I, and the integra-%
tion and cointegration operations from D(I) to $D_o$ as substitutes (interms
of these finer objects) of direct and inverse limits in $D_o$ . When tempted
to think of these latter operations (with values in $D_o$ ) as the basic
structure involved, one cannot help \var{\inserted{though,}}{\hand{though,}}{insertion}{ps_134-5_013-v0009} looking for the same kind of structu-%
re on any one of these subsidiary categories D(I), as these are being
thought of as derived categories in their own right. It then appears at
once that the \var{\inserted{``more}}{\hand{``more}}{insertion}{ps_134-5_013-v0010} refined substitues''\hand{/}for J-indexed systems of objects of D(I)
are just the objects of D(IxJ), and the corresponding integration and
cointegration operations
\[ D(IxJ) \hand{\longrightarrow} D(I) \]
are nothing but $p_!$ and $p_*$ , where p: IxJ \hand{$\longrightarrow$} I is the projection. Thus
one is inevitably conducted to look at operations $f_!$ and $f_*$ instead of
merely integration and cointegration -- thus providing for the ``inner
stability'' of the structure described, as I had been l\hand{ooking for from}
the very start.

The notion of integration of homotopy types appears here \hand{as a} natur\hand{\unsure{al|el}}
by-product of an attempt to grasp the ``full structure'' of a triangulated
cate\hand{go}ry. However, I had been feeling the need for such a notion of
integration of homotopy types for about one or two years already,
\hand{(}without any clear \var{\deleted{notion}}{}{deletion}{ps_134-5_013-v0011} \var{\inserted{idea yet}}{\hand{idea yet}}{insertion}{ps_134-5_013-v0012} that this operation should be one out of two

\margin{13}

%% END SOURCE ps_134-5_013
\ifnum\EditionMode=1
\psapparatusstart
\psentry{ps_134-5_013-v0001}{1}{typewriter deletion}{\xout{axratherxgeneralx}}{}{The typewriter strikes out this material; it leaves the corrected running text.}
\psentry{ps_134-5_013-v0002}{2}{typewriter deletion}{\xout{setx\allowbreak{}upx\allowbreak{}inx\allowbreak{}mindx\allowbreak{}(includingx\allowbreak{}Quillen'sx\allowbreak{})x\allowbreak{}x\allowbreak{}withx\allowbreak{}thex\allowbreak{}"mildx\allowbreak{}restriction"x\allowbreak{})x\allowbreak{}}}{}{The typewriter strikes out this material; it leaves the corrected running text.}
\psentry{ps_134-5_013-v0003}{3}{typewriter deletion}{\xout{which}}{}{The typewriter strikes out this material; it leaves the corrected running text.}
\psentry{ps_134-5_013-v0004}{4}{deletion}{\deleted{t}}{}{The author cancels this material by hand; it leaves the corrected running text.}
\psentry{ps_134-5_013-v0005}{5}{deletion}{\deleted{fact}}{}{The author cancels this material by hand; it leaves the corrected running text.}
\psentry{ps_134-5_013-v0006}{6}{insertion}{\inserted{expectation,}}{\hand{expectation,}}{The handwritten insertion is incorporated at its indicated position; insertion brackets leave the running text.}
\psentry{ps_134-5_013-v0007}{7}{insertion}{\inserted{is}}{\hand{is}}{The handwritten insertion is incorporated at its indicated position; insertion brackets leave the running text.}
\psentry{ps_134-5_013-v0008}{8}{typewriter deletion}{\xout{sub}}{}{The typewriter strikes out this material; it leaves the corrected running text.}
\psentry{ps_134-5_013-v0009}{9}{insertion}{\inserted{though,}}{\hand{though,}}{The handwritten insertion is incorporated at its indicated position; insertion brackets leave the running text.}
\psentry{ps_134-5_013-v0010}{10}{insertion}{\inserted{``more}}{\hand{``more}}{The handwritten insertion is incorporated at its indicated position; insertion brackets leave the running text.}
\psentry{ps_134-5_013-v0011}{11}{deletion}{\deleted{notion}}{}{The author cancels this material by hand; it leaves the corrected running text.}
\psentry{ps_134-5_013-v0012}{12}{insertion}{\inserted{idea yet}}{\hand{idea yet}}{The handwritten insertion is incorporated at its indicated position; insertion brackets leave the running text.}
\psapparatusend
\fi
\pssource{ps_134-5_014}{134-5/14}{https://grothendieck.umontpellier.fr/web/viewer.html?file=../134-5.pdf\#page=15}
\psnumber{ps_134-5_014-v0001}{1}
\psnumber{ps_134-5_014-v0002}{2}
\psnumber{ps_134-5_014-v0003}{3}
\psnumber{ps_134-5_014-v0004}{4}
\psnumber{ps_134-5_014-v0005}{5}
\psnumber{ps_134-5_014-v0006}{6}
\psnumber{ps_134-5_014-v0007}{7}
\psnumber{ps_134-5_014-v0008}{8}
\psnumber{ps_134-5_014-v0009}{9}
\psnumber{ps_134-5_014-v0010}{10}
\psnumber{ps_134-5_014-v0011}{11}
\psnumber{ps_134-5_014-v0012}{12}
\psnumber{ps_134-5_014-v0013}{13}
\psnumber{ps_134-5_014-v0014}{14}
\psnumber{ps_134-5_014-v0015}{15}
\psnumber{ps_134-5_014-v0016}{16}
\psnumber{ps_134-5_014-v0017}{17}
\psnumber{ps_134-5_014-v0018}{18}
%% BEGIN SOURCE ps_134-5_014
% Pursuing Stacks (cotes 134-2 to 134-8): archive cote 134-5, site page 14 (PDF page 15), typescript page 201.
% typescript. Draft compared: ps4.tex:612-670.
% AI-transcribed by Claude Opus 5.5 (claude-opus-5-5, adaptive thinking, effort high), Message Batches API, batch ps2 (msgbatch_01NVwr1Yuo8qJf5rykeGBXPK), 2026-11-12; reviewed against both supplied scan halves; see CHECK_DECISIONS.tsv.
\scanpage{134-5/14}
\mspage{201}
main ingredients of a (\var{\inserted{by then}}{\hand{by then}}{insertion}{ps_134-5_014-v0001} still very misty) notion of a triangulated category
of sorts\hand{)}. This feeling arose from my  ponderings on stratified structures,
and the ``screwing together'' of such structures in terms of simple building
blocks (essentially,  various types of ``tubes'' associated to such structu\var{\xout{r}}{}{typewriter deletion}{ps_134-5_014-v0002}%
res,  related to each other by various \var{\inserted[typed]{proper}}{proper}{typed insertion}{ps_134-5_014-v0003} maps which are either inclusions,
or  -  in the equisingular case \var{\inserted[typed]{at any rate}}{at any rate}{typed insertion}{ps_134-5_014-v0004} -  fiber maps).  This\hand{``}screwing together
\var{\inserted{operation''}}{\hand{operation''}}{insertion}{ps_134-5_014-v0005} could be expressed as being a direct limit of a certain finite system of
spaces.   In the cases I was most interested in (namely the Mumford-Deligne
compactifications $M^{\wedge}_{g,\hand{\nu}}$  of the modular multiplicities  $M_{g,\hand{\nu}}$),  these
spaces or\hand{``}tubes\hand{''} have exceedingly simple homotopy types  -  they \var{\xout{wwww}}{}{typewriter deletion}{ps_134-5_014-v0006} \var{\inserted[typed]{are}}{are}{typed insertion}{ps_134-5_014-v0007} just
$K(\hand{\pi},1)$-spaces,  where \var{\xout{ex}}{}{typewriter deletion}{ps_134-5_014-v0008} \var{\inserted[typed]{each}}{each}{typed insertion}{ps_134-5_014-v0009} \hand{$\pi$} is a Teichmüller-type discrete group (practical-%
ly,  a product of usual Teichmüller groups).  It then occured to me that
the \var{\inserted[typed]{whole}}{whole}{typed insertion}{ps_134-5_014-v0010} homotopy type of $M^{\wedge}_{g,\hand{\nu}}$ ,  or of any locally closed union of strata\hand{,}
or (more generally still) of ``the'' tubular neighbourhood of \var{\xout{the}}{}{typewriter deletion}{ps_134-5_014-v0011} such a
union in any larger one\hand{,} \var{\inserted[typed]{etc}}{etc}{typed insertion}{ps_134-5_014-v0012} -  that all these homotopy types should be ex-%
pressible in terms of the given \var{\xout{direct}}{}{typewriter deletion}{ps_134-5_014-v0013} system of spaces,  and more accura-%
tely still,  just in terms of the corresponding system of fundamental
groupoids (embodying their homotopy types).  In this situation,  what I
was mainly out for,  was precisely an accurate and workeable description
of this direct system of groupoids  (  which could be viewed as just one
section of the whole ``Teichmüller tower'' of Teichmüller groupoïds ...).
Thus,  it was a rewarding extra feature of the situation (by then just an
expectation\hand{,} as a matter of fact),  that such a description should at the
same time yie\var{\over{\unsure{l}}{l}}{\hand{l}}{writing over}{ps_134-5_014-v0014}d a ``purely algebraic'' description of the homotopy types
of all the spaces (rather,  multiplicities,to be wholly accurate) which
I could think of in terms of the natural stratification of $M^{\wedge}_{g,\hand{\nu}}$ .  The\hand{re}
was an awareness that this operation on homotopy types could not be
described simply in terms of a functor  I \hand{$\longrightarrow$} (Hot) ,  where I is the in-%
dexing category,  that a functor \var{\inserted{$i\mapsto X_i$ :}}{\hand{$i\mapsto X_i$ :}}{insertion}{ps_134-5_014-v0015} I \hand{$\longrightarrow$} M  (where M is some model category
such as (Spaces)  or (Cat) ) should be available in order to define an
``integra\var{\over{\unsure{l}}{t}}{\hand{t}}{writing over}{ps_134-5_014-v0016}ed homotopy type''  \hand{$\int$}$X_i$ .  This justified feeling got somewhat
blurred lately,  for a little while,  by the definitely unreasonable ex-%
pectation that finite limits should exist in (Hot) after all,  why not !
It's enough to have a look though (which probably I did years ago and then
forgot in the meanwhile) to ma\var{\over{\unsure{k}}{k}}{\hand{k}}{writing over}{ps_134-5_014-v0017}e sure they dont ...

\quad Whether or not this notion of ``integration of homotopy types'' is
more or less well known already under some name or other,  is'nt quite
clear to me.  It is'nt familiar to Ronnie Brown visibly,  but it seems he
heard about such a kind of thing,  without \var{\inserted{his}}{\hand{his}}{insertion}{ps_134-5_014-v0018} being specific about it.  It
\margin{14}

%% END SOURCE ps_134-5_014
\ifnum\EditionMode=1
\psapparatusstart
\psentry{ps_134-5_014-v0001}{1}{insertion}{\inserted{by then}}{\hand{by then}}{The handwritten insertion is incorporated at its indicated position; insertion brackets leave the running text.}
\psentry{ps_134-5_014-v0002}{2}{typewriter deletion}{\xout{r}}{}{The typewriter strikes out this material; it leaves the corrected running text.}
\psentry{ps_134-5_014-v0003}{3}{typed insertion}{\inserted[typed]{proper}}{proper}{The typed insertion is incorporated at its indicated position; insertion brackets leave the running text.}
\psentry{ps_134-5_014-v0004}{4}{typed insertion}{\inserted[typed]{at any rate}}{at any rate}{The typed insertion is incorporated at its indicated position; insertion brackets leave the running text.}
\psentry{ps_134-5_014-v0005}{5}{insertion}{\inserted{operation''}}{\hand{operation''}}{The handwritten insertion is incorporated at its indicated position; insertion brackets leave the running text.}
\psentry{ps_134-5_014-v0006}{6}{typewriter deletion}{\xout{wwww}}{}{The typewriter strikes out this material; it leaves the corrected running text.}
\psentry{ps_134-5_014-v0007}{7}{typed insertion}{\inserted[typed]{are}}{are}{The typed insertion is incorporated at its indicated position; insertion brackets leave the running text.}
\psentry{ps_134-5_014-v0008}{8}{typewriter deletion}{\xout{ex}}{}{The typewriter strikes out this material; it leaves the corrected running text.}
\psentry{ps_134-5_014-v0009}{9}{typed insertion}{\inserted[typed]{each}}{each}{The typed insertion is incorporated at its indicated position; insertion brackets leave the running text.}
\psentry{ps_134-5_014-v0010}{10}{typed insertion}{\inserted[typed]{whole}}{whole}{The typed insertion is incorporated at its indicated position; insertion brackets leave the running text.}
\psentry{ps_134-5_014-v0011}{11}{typewriter deletion}{\xout{the}}{}{The typewriter strikes out this material; it leaves the corrected running text.}
\psentry{ps_134-5_014-v0012}{12}{typed insertion}{\inserted[typed]{etc}}{etc}{The typed insertion is incorporated at its indicated position; insertion brackets leave the running text.}
\psentry{ps_134-5_014-v0013}{13}{typewriter deletion}{\xout{direct}}{}{The typewriter strikes out this material; it leaves the corrected running text.}
\psentry{ps_134-5_014-v0014}{14}{writing over}{\over{\unsure{l}}{l}}{\hand{l}}{The superposed handwritten reading replaces the underlying reading.}
\psentry{ps_134-5_014-v0015}{15}{insertion}{\inserted{$i\mapsto X_i$ :}}{\hand{$i\mapsto X_i$ :}}{The handwritten insertion is incorporated at its indicated position; insertion brackets leave the running text.}
\psentry{ps_134-5_014-v0016}{16}{writing over}{\over{\unsure{l}}{t}}{\hand{t}}{The superposed handwritten reading replaces the underlying reading.}
\psentry{ps_134-5_014-v0017}{17}{writing over}{\over{\unsure{k}}{k}}{\hand{k}}{The superposed handwritten reading replaces the underlying reading.}
\psentry{ps_134-5_014-v0018}{18}{insertion}{\inserted{his}}{\hand{his}}{The handwritten insertion is incorporated at its indicated position; insertion brackets leave the running text.}
\psapparatusend
\fi
\pssource{ps_134-5_015}{134-5/15}{https://grothendieck.umontpellier.fr/web/viewer.html?file=../134-5.pdf\#page=16}
\psnumber{ps_134-5_015-v0001}{1}
\psnumber{ps_134-5_015-v0002}{2}
\psnumber{ps_134-5_015-v0003}{3}
\psnumber{ps_134-5_015-v0004}{4}
\psnumber{ps_134-5_015-v0005}{5}
\psnumber{ps_134-5_015-v0006}{6}
\psnumber{ps_134-5_015-v0007}{7}
\psnumber{ps_134-5_015-v0008}{8}
\psnumber{ps_134-5_015-v0009}{9}
\psnumber{ps_134-5_015-v0010}{10}
\psnumber{ps_134-5_015-v0011}{11}
\psnumber{ps_134-5_015-v0012}{12}
\psnumber{ps_134-5_015-v0013}{13}
\psnumber{ps_134-5_015-v0014}{14}
\psnumber{ps_134-5_015-v0015}{15}
\psnumber{ps_134-5_015-v0016}{16}
\psnumber{ps_134-5_015-v0017}{17}
\psnumber{ps_134-5_015-v0018}{18}
\psnumber{ps_134-5_015-v0019}{19}
%% BEGIN SOURCE ps_134-5_015
% Pursuing Stacks (cotes 134-2 to 134-8): archive cote 134-5, site page 15 (PDF page 16), typescript page 202.
% typescript. Draft compared: ps4.tex:657-722.
% AI-transcribed by Claude Opus 5.5 (claude-opus-5-5, adaptive thinking, effort high), Message Batches API, batch ps2 (msgbatch_01NVwr1Yuo8qJf5rykeGBXPK), 2026-11-12; reviewed against both supplied scan halves; see CHECK_DECISIONS.tsv.
\scanpage{134-5/15}
\mspage{202}
was the episodic correspondance with him which finally pushed me \var{\inserted[typed]{last}}{last}{typed insertion}{ps_134-5_015-v0001} \var{\xout{in}}{}{typewriter deletion}{ps_134-5_015-v0002} janua-%
ry to sit down for an afternoon and try to figure out what there actually
was, in a lengthy and somewhat rambling letter to Illusie (who does'nt
seem to have heard at all about such operatio\var{\over{n}{ns}}{\hand{ns}}{writing over}{ps_134-5_015-v0003}). This preliminary reflec-%
tion proved quite useful lately, I'll have to come back anyhow upon some
of the specific features of integration of homotopy types later, and there
is not much point dwelling on it any longer \var{\inserted[typed]{at present.}}{at present.}{typed insertion}{ps_134-5_015-v0004} \var{\deleted{\xout{in this retrospective}}}{}{deletion}{ps_134-5_015-v0005}

\margin{70.} This \hand{was} a \hand{(}\var{\inserted{finally}}{\hand{finally}}{insertion}{ps_134-5_015-v0006} somewhat lengthy\hand{!)} review of ponderings which did'nt take me
more than just a few days, because it was about things some of which were
on my mind for a long time indeed. It took \var{\xout{me}}{}{typewriter deletion}{ps_134-5_015-v0007} a lot more work to try to
carry through the standard homotopy constructions, giving rise to the
Dold-Puppe sequences, within the basic modelizer (Cat). Most of the work
arose, it now seems to me, out of a block I got ( I could'nt tell why)
against the K\hand{an}\var{\inserted[typed]{type}}{type}{typed insertion}{ps_134-5_015-v0008} condition on complexes, so I tried hard to get along
without anything of the sort. I kind of fooled myself into believing that
I was forced to do so, because I was working in an axiomatic set-up depen-%
dant upon the``basic localizer'' $\underline{W}$ , so the Kan condition would'nt be rele-%
vant anyhow. The main point was to get, for any map f: $X \hand{\to} Y$ in (Cat),
a factorization
\[ X \hand{\xrightarrow{\textcolor{black}{i}}} Z \hand{\xrightarrow{\textcolor{black}{p}}} Y \quad , \]
where p has the property that base change by p transforms weak equivalences
into weak equivalences\hand{,} \var{\deleted{and}}{}{deletion}{ps_134-5_015-v0009} (visibly a Serre-fibration type condition),
and i satisfies the dual condition with respect to co\hand{-}base change\hand{;} and
moreover where either p , or i can be assumed to be in $\underline{W}$ i.e. \var{\inserted[typed]{to}}{to}{typed insertion}{ps_134-5_015-v0010} be a weak
equivalence. (This\hand{,} by the way\hand{,} is the extra condition \hand{on a pair} (M,W) I
have been referring to \var{\deleted{on}}{}{deletion}{ps_134-5_015-v0011} \var{\inserted[typed]{above}}{above}{typed insertion}{ps_134-5_015-v0012} (page 200)\var{\deleted{\unsure{)}}}{}{deletion}{ps_134-5_015-v0013}\hand{,} for \var{\inserted{(hopefully)}}{\hand{(hopefully)}}{insertion}{ps_134-5_015-v0014} getting $f_{\hand{!}}$ and $f_{*}$ \hand{operations.)}
At present, I do not yet know whether such factorization \var{\xout{is}}{}{typewriter deletion}{ps_134-5_015-v0015} always
exists for a map in (Cat), without even \var{\xout{assumin}}{}{typewriter deletion}{ps_134-5_015-v0016} demanding that either i
or p should be in $\underline{W}$ . \hand{$\Vert$}

I first devoted a lot of attention to Serre-type
conditions on maps in (Cat), which turned out quite rewarding - with
the impression \var{\over{\unsure{o}}{of}}{\hand{of}}{writing over}{ps_134-5_015-v0017} arrivi\hand{ng} at a coherent and nicely auto-dual picture of
cohomology properties of functors i.e. maps in (Cat), as far as these
were concerned with \hand{\uline{\textcolor{black}{base change \var{\over{b}{b}}{\hand{b}}{writing over}{ps_134-5_015-v0018}ehaviour}}} (and \hand{\uline{\textcolor{black}{not}}} co-base change). Here
I was guided by work done long ago to get the étale cohomology theory off
the ground, and where the two main theorems achieving this aim were pre-%
cisely two theorems of commutation of ``higher direct images'' $R^{i}g_{\hand{*}}$ with
respect to base-change by a map h - namely, \var{\inserted{it is OK}}{\hand{it is OK}}{insertion}{ps_134-5_015-v0019} when either \hand{g} is \hand{\uline{\textcolor{black}{proper}}},
or h is \hand{\uline{\textcolor{black}{smooth}}}. It was rather natural then to introduce the notion of

\margin{15}

%% END SOURCE ps_134-5_015
\ifnum\EditionMode=1
\psapparatusstart
\psentry{ps_134-5_015-v0001}{1}{typed insertion}{\inserted[typed]{last}}{last}{The typed insertion is incorporated at its indicated position; insertion brackets leave the running text.}
\psentry{ps_134-5_015-v0002}{2}{typewriter deletion}{\xout{in}}{}{The typewriter strikes out this material; it leaves the corrected running text.}
\psentry{ps_134-5_015-v0003}{3}{writing over}{\over{n}{ns}}{\hand{ns}}{The superposed handwritten reading replaces the underlying reading.}
\psentry{ps_134-5_015-v0004}{4}{typed insertion}{\inserted[typed]{at present.}}{at present.}{The typed insertion is incorporated at its indicated position; insertion brackets leave the running text.}
\psentry{ps_134-5_015-v0005}{5}{deletion}{\deleted{\xout{in this retrospective}}}{}{The author cancels this material by hand; it leaves the corrected running text.}
\psentry{ps_134-5_015-v0006}{6}{insertion}{\inserted{finally}}{\hand{finally}}{The handwritten insertion is incorporated at its indicated position; insertion brackets leave the running text.}
\psentry{ps_134-5_015-v0007}{7}{typewriter deletion}{\xout{me}}{}{The typewriter strikes out this material; it leaves the corrected running text.}
\psentry{ps_134-5_015-v0008}{8}{typed insertion}{\inserted[typed]{type}}{type}{The typed insertion is incorporated at its indicated position; insertion brackets leave the running text.}
\psentry{ps_134-5_015-v0009}{9}{deletion}{\deleted{and}}{}{The author cancels this material by hand; it leaves the corrected running text.}
\psentry{ps_134-5_015-v0010}{10}{typed insertion}{\inserted[typed]{to}}{to}{The typed insertion is incorporated at its indicated position; insertion brackets leave the running text.}
\psentry{ps_134-5_015-v0011}{11}{deletion}{\deleted{on}}{}{The author cancels this material by hand; it leaves the corrected running text.}
\psentry{ps_134-5_015-v0012}{12}{typed insertion}{\inserted[typed]{above}}{above}{The typed insertion is incorporated at its indicated position; insertion brackets leave the running text.}
\psentry{ps_134-5_015-v0013}{13}{deletion}{\deleted{\unsure{)}}}{}{The author cancels this material by hand; it leaves the corrected running text.}
\psentry{ps_134-5_015-v0014}{14}{insertion}{\inserted{(hopefully)}}{\hand{(hopefully)}}{The handwritten insertion is incorporated at its indicated position; insertion brackets leave the running text.}
\psentry{ps_134-5_015-v0015}{15}{typewriter deletion}{\xout{is}}{}{The typewriter strikes out this material; it leaves the corrected running text.}
\psentry{ps_134-5_015-v0016}{16}{typewriter deletion}{\xout{assumin}}{}{The typewriter strikes out this material; it leaves the corrected running text.}
\psentry{ps_134-5_015-v0017}{17}{writing over}{\over{\unsure{o}}{of}}{\hand{of}}{The superposed handwritten reading replaces the underlying reading.}
\psentry{ps_134-5_015-v0018}{18}{writing over}{\over{b}{b}}{\hand{b}}{The superposed handwritten reading replaces the underlying reading.}
\psentry{ps_134-5_015-v0019}{19}{insertion}{\inserted{it is OK}}{\hand{it is OK}}{The handwritten insertion is incorporated at its indicated position; insertion brackets leave the running text.}
\psapparatusend
\fi
\pssource{ps_134-5_016}{134-5/16}{https://grothendieck.umontpellier.fr/web/viewer.html?file=../134-5.pdf\#page=17}
\psnumber{ps_134-5_016-v0001}{1}
\psnumber{ps_134-5_016-v0002}{2}
\psnumber{ps_134-5_016-v0003}{3}
\psnumber{ps_134-5_016-v0004}{4}
\psnumber{ps_134-5_016-v0005}{5}
\psnumber{ps_134-5_016-v0006}{6}
\psnumber{ps_134-5_016-v0007}{7}
\psnumber{ps_134-5_016-v0008}{8}
\psnumber{ps_134-5_016-v0009}{9}
\psnumber{ps_134-5_016-v0010}{10}
\psnumber{ps_134-5_016-v0011}{11}
\psnumber{ps_134-5_016-v0012}{12}
\psnumber{ps_134-5_016-v0013}{13}
\psnumber{ps_134-5_016-v0014}{14}
\psnumber{ps_134-5_016-v0015}{15}
\psnumber{ps_134-5_016-v0016}{16}
\psnumber{ps_134-5_016-v0017}{17}
\psnumber{ps_134-5_016-v0018}{18}
\psnumber{ps_134-5_016-v0019}{19}
\psnumber{ps_134-5_016-v0020}{20}
\psnumber{ps_134-5_016-v0021}{21}
\psnumber{ps_134-5_016-v0022}{22}
\psnumber{ps_134-5_016-v0023}{23}
\psnumber{ps_134-5_016-v0024}{24}
\psnumber{ps_134-5_016-v0025}{25}
\psnumber{ps_134-5_016-v0026}{26}
\psnumber{ps_134-5_016-v0027}{27}
\psnumber{ps_134-5_016-v0028}{28}
\psnumber{ps_134-5_016-v0029}{29}
\psnumber{ps_134-5_016-v0030}{30}
%% BEGIN SOURCE ps_134-5_016
% Pursuing Stacks (cotes 134-2 to 134-8): archive cote 134-5, site page 16 (PDF page 17), typescript page 203.
% typescript. Draft compared: ps4.tex:710-769.
% AI-transcribed by Claude Opus 5.5 (claude-opus-5-5, adaptive thinking, effort high), Message Batches API, batch ps2 (msgbatch_01NVwr1Yuo8qJf5rykeGBXPK), 2026-11-12; reviewed against both supplied scan halves; see CHECK_DECISIONS.tsv.
\scanpage{134-5/16}
\mspage{203}

(cohomological) \hand{\uline{\textcolor{black}{smoothness}}} and (coh.) \hand{\uline{\textcolor{black}{properness}}} of a map in (Cat), by
the obvious base-change propert\var{\over{y}{ies}}{\hand{ies}}{writing over}{ps_134-5_016-v0001}. It turned out that these could be readi-%
ly \var{\over{ch}{cha}}{\hand{cha}}{writing over}{ps_134-5_016-v0002}racterized by suitable asphericity conditions, which are formally
quite similar to the well known valuative criteri\var{\over{um}{on}}{\hand{on}}{writing over}{ps_134-5_016-v0003} for a map of finite
type between noetherian schemes to be \var{\deleted{either}}{}{deletion}{ps_134-5_016-v0004} ``universally open'' (which can
be viewed as a ``purely topological'' variant of the notion of smoothness),
\var{\deleted{or}}{}{deletion}{ps_134-5_016-v0005}\var{\inserted{resp.}}{\hand{resp.}}{insertion}{ps_134-5_016-v0006} ``universally closed'' \var{\deleted{(for instance,}}{}{deletion}{ps_134-5_016-v0007} \var{\inserted{or rather, more stringently,}}{\hand{or rather, more stringently,}}{insertion}{ps_134-5_016-v0008} proper\var{\deleted{)}}{}{deletion}{ps_134-5_016-v0009}. These conditions, more-%
over, are trivially satisfied when f: X$\hand{\longrightarrow}$Y turns X into a category over Y
which is \var{\deleted{either \xout{co}}}{}{deletion}{ps_134-5_016-v0010}fibered over Y (i.e. definable in terms of a \var{\inserted{contravariant}}{\hand{contravariant}}{insertion}{ps_134-5_016-v0011} pseudo-func-%
tor $Y^{\circ}\hand{\longrightarrow}$(Cat)) , \var{\deleted{or}}{}{deletion}{ps_134-5_016-v0012} \var{\inserted{resp.}}{\hand{resp.}}{insertion}{ps_134-5_016-v0013} \var{\inserted{co}}{\hand{co}}{insertion}{ps_134-5_016-v0014}fibered over Y (namely, defin\var{\xout{\unsure{a}}}{}{typewriter deletion}{ps_134-5_016-v0015}able in terms of a
\var{\xout{contravariant}}{}{typewriter deletion}{ps_134-5_016-v0016} pseudo-functor $Y$\var{\deleted{$^{\circ}$}}{}{deletion}{ps_134-5_016-v0017}$\hand{\longrightarrow}$(Cat)). If we call such maps in (Cat)
``\var{\xout{co}}{}{typewriter deletion}{ps_134-5_016-v0018}fibrations'' and ``\var{\inserted[typed]{co}}{co}{typed insertion}{ps_134-5_016-v0019}fibrations'' (very much in conflict, alas, with \var{\xout{\unsure{the}}}{}{typewriter deletion}{ps_134-5_016-v0020}
Quillen's neat set-up of \var{\xout{co}}{}{typewriter deletion}{ps_134-5_016-v0021}fibrations-\var{\inserted[typed]{co}}{co}{typed insertion}{ps_134-5_016-v0022}fibrations !) , it turns out that
\hand{\uline{\textcolor{black}{fibrations are smooth}}}, \hand{\uline{\textcolor{black}{cofibrations are proper}}}. This is all the more re-%
markab\unsure{k|l}e, as the notions of fibered and cofibered categories were intro-%
duced with a view upon ``large categories'', in order to pin down some
standard properties met in all situations of ``base change'' \hand{(}and \var{\inserted[typed]{later}}{later}{typed insertion}{ps_134-5_016-v0023} the dual
situation of co-base change \hand{)} - the main motivation for this being the
need to fomulate with a minimum of precision a set-up for \hand{``}descent techni-%
ques\hand{''} in algebraic geometry. (These techniques, as well as the cohomologi-%
cal base change theorems, make visibly a sense too in the context of
analytic spaces say, or of topological spaces, but they don't seem to
have been assimilated yet by geometers outside of algebraic geometry\hand{)}
That these typically ``general non-sense'' notions should have such precise
topological implications came as a complete surprise ! As a consequence,
a bifibration (namely a map which is both a fibration and a cofibration)
is smooth and proper and hence a \hand{(}cohomological) ``Serre fibration'', for
instance the sheaves $R^{i}f_{*}(F)$, when F is a constant \var{\inserted[typed]{abelian}}{abelian}{typed insertion}{ps_134-5_016-v0024} sheaf on X \hand{(}i.e. a\var{\xout{n}}{}{typewriter deletion}{ps_134-5_016-v0025}
\var{\inserted[typed]{constant}}{constant}{typed insertion}{ps_134-5_016-v0026} abelian group object in $X^{\wedge}$\hand{)}, namely $y\hand{\mapsto}H^{i}(X_{/y},F)$ , are \hand{\uline{\textcolor{black}{local}}} \hand{\uline{\textcolor{black}{systems}}} \var{\xout{\illegible{1}}}{}{typewriter deletion}{ps_134-5_016-v0027}
on Y i.e. factor through the fundamental groupoid of Y. \unsure{*}

A \var{\xout{still}}{}{typewriter deletion}{ps_134-5_016-v0028} greater surprise \var{\over{st}{st}}{\hand{st}}{writing over}{ps_134-5_016-v0029}ill was the\var{\xout{\illegible{1}}}{}{typewriter deletion}{ps_134-5_016-v0030} \hand{\uline{\textcolor{black}{duality between the notions of
properness and smoothness}}}: just as a map f: X$\hand{\longrightarrow}$Y in (Cat) is a cofibra-%
tion iff the ``dual'' map $f^{\circ}: X^{\circ}\hand{\longrightarrow}Y^{\circ}$ is a fibration, it turns out that
f \hand{\uline{\textcolor{black}{is proper iff}}} $f^{\circ}$ \hand{\uline{\textcolor{black}{is smooth}}}. This was a really startling fact, and it
caused me to wonder, in the context of more general topoi than just those
of the type $X^{\wedge}$, whether there was'nt \hand{\uline{\textcolor{black}{a notion of duality generalizing}}}
\hand{\uline{\textcolor{black}{the relationship between two topoi $X^{\wedge}$ and}}} \hand{$X^{\vee}$} $\overset{def}{=}$ $(X^{\circ})^{\wedge}$ . Indeed, the\unsure{se}
two categories of sheaves can be described intrinsically one in terms of
the other, up to equivalence, by a natural pairing

\margin{16}

%% END SOURCE ps_134-5_016
\ifnum\EditionMode=1
\psapparatusstart
\psentry{ps_134-5_016-v0001}{1}{writing over}{\over{y}{ies}}{\hand{ies}}{The superposed handwritten reading replaces the underlying reading.}
\psentry{ps_134-5_016-v0002}{2}{writing over}{\over{ch}{cha}}{\hand{cha}}{The superposed handwritten reading replaces the underlying reading.}
\psentry{ps_134-5_016-v0003}{3}{writing over}{\over{um}{on}}{\hand{on}}{The superposed handwritten reading replaces the underlying reading.}
\psentry{ps_134-5_016-v0004}{4}{deletion}{\deleted{either}}{}{The author cancels this material by hand; it leaves the corrected running text.}
\psentry{ps_134-5_016-v0005}{5}{deletion}{\deleted{or}}{}{The author cancels this material by hand; it leaves the corrected running text.}
\psentry{ps_134-5_016-v0006}{6}{insertion}{\inserted{resp.}}{\hand{resp.}}{The handwritten insertion is incorporated at its indicated position; insertion brackets leave the running text.}
\psentry{ps_134-5_016-v0007}{7}{deletion}{\deleted{(for instance,}}{}{The author cancels this material by hand; it leaves the corrected running text.}
\psentry{ps_134-5_016-v0008}{8}{insertion}{\inserted{or rather, more stringently,}}{\hand{or rather, more stringently,}}{The handwritten insertion is incorporated at its indicated position; insertion brackets leave the running text.}
\psentry{ps_134-5_016-v0009}{9}{deletion}{\deleted{)}}{}{The author cancels this material by hand; it leaves the corrected running text.}
\psentry{ps_134-5_016-v0010}{10}{deletion}{\deleted{either \xout{co}}}{}{The author cancels this material by hand; it leaves the corrected running text.}
\psentry{ps_134-5_016-v0011}{11}{insertion}{\inserted{contravariant}}{\hand{contravariant}}{The handwritten insertion is incorporated at its indicated position; insertion brackets leave the running text.}
\psentry{ps_134-5_016-v0012}{12}{deletion}{\deleted{or}}{}{The author cancels this material by hand; it leaves the corrected running text.}
\psentry{ps_134-5_016-v0013}{13}{insertion}{\inserted{resp.}}{\hand{resp.}}{The handwritten insertion is incorporated at its indicated position; insertion brackets leave the running text.}
\psentry{ps_134-5_016-v0014}{14}{insertion}{\inserted{co}}{\hand{co}}{The handwritten insertion is incorporated at its indicated position; insertion brackets leave the running text.}
\psentry{ps_134-5_016-v0015}{15}{typewriter deletion}{\xout{\unsure{a}}}{}{The typewriter strikes out this material; it leaves the corrected running text.}
\psentry{ps_134-5_016-v0016}{16}{typewriter deletion}{\xout{contravariant}}{}{The typewriter strikes out this material; it leaves the corrected running text.}
\psentry{ps_134-5_016-v0017}{17}{deletion}{\deleted{$^{\circ}$}}{}{The author cancels this material by hand; it leaves the corrected running text.}
\psentry{ps_134-5_016-v0018}{18}{typewriter deletion}{\xout{co}}{}{The typewriter strikes out this material; it leaves the corrected running text.}
\psentry{ps_134-5_016-v0019}{19}{typed insertion}{\inserted[typed]{co}}{co}{The typed insertion is incorporated at its indicated position; insertion brackets leave the running text.}
\psentry{ps_134-5_016-v0020}{20}{typewriter deletion}{\xout{\unsure{the}}}{}{The typewriter strikes out this material; it leaves the corrected running text.}
\psentry{ps_134-5_016-v0021}{21}{typewriter deletion}{\xout{co}}{}{The typewriter strikes out this material; it leaves the corrected running text.}
\psentry{ps_134-5_016-v0022}{22}{typed insertion}{\inserted[typed]{co}}{co}{The typed insertion is incorporated at its indicated position; insertion brackets leave the running text.}
\psentry{ps_134-5_016-v0023}{23}{typed insertion}{\inserted[typed]{later}}{later}{The typed insertion is incorporated at its indicated position; insertion brackets leave the running text.}
\psentry{ps_134-5_016-v0024}{24}{typed insertion}{\inserted[typed]{abelian}}{abelian}{The typed insertion is incorporated at its indicated position; insertion brackets leave the running text.}
\psentry{ps_134-5_016-v0025}{25}{typewriter deletion}{\xout{n}}{}{The typewriter strikes out this material; it leaves the corrected running text.}
\psentry{ps_134-5_016-v0026}{26}{typed insertion}{\inserted[typed]{constant}}{constant}{The typed insertion is incorporated at its indicated position; insertion brackets leave the running text.}
\psentry{ps_134-5_016-v0027}{27}{typewriter deletion}{\xout{\illegible{1}}}{}{The typewriter strikes out this material; it leaves the corrected running text.}
\psentry{ps_134-5_016-v0028}{28}{typewriter deletion}{\xout{still}}{}{The typewriter strikes out this material; it leaves the corrected running text.}
\psentry{ps_134-5_016-v0029}{29}{writing over}{\over{st}{st}}{\hand{st}}{The superposed handwritten reading replaces the underlying reading.}
\psentry{ps_134-5_016-v0030}{30}{typewriter deletion}{\xout{\illegible{1}}}{}{The typewriter strikes out this material; it leaves the corrected running text.}
\psapparatusend
\fi
\pssource{ps_134-5_017}{134-5/17}{https://grothendieck.umontpellier.fr/web/viewer.html?file=../134-5.pdf\#page=18}
\psnumber{ps_134-5_017-v0001}{1}
\psnumber{ps_134-5_017-v0002}{2}
\psnumber{ps_134-5_017-v0003}{3}
\psnumber{ps_134-5_017-v0004}{4}
\psnumber{ps_134-5_017-v0005}{5}
\psnumber{ps_134-5_017-v0006}{6}
\psnumber{ps_134-5_017-v0007}{7}
\psnumber{ps_134-5_017-v0008}{8}
\psnumber{ps_134-5_017-v0009}{9}
\psnumber{ps_134-5_017-v0010}{10}
\psnumber{ps_134-5_017-v0011}{11}
\psnumber{ps_134-5_017-v0012}{12}
\psnumber{ps_134-5_017-v0013}{13}
\psnumber{ps_134-5_017-v0014}{14}
%% BEGIN SOURCE ps_134-5_017
% Pursuing Stacks (cotes 134-2 to 134-8): archive cote 134-5, site page 17 (PDF page 18), typescript page 204.
% typescript. Draft compared: ps4.tex:756-816.
% AI-transcribed by Claude Opus 5.5 (claude-opus-5-5, adaptive thinking, effort high), Message Batches API, batch ps2 (msgbatch_01NVwr1Yuo8qJf5rykeGBXPK), 2026-11-12; reviewed against both supplied scan halves; see CHECK_DECISIONS.tsv.
\scanpage{134-5/17}
\mspage{204}

\[X^{\wedge}\mathrm{x}X^{\hand{\vee}}\hand{\longrightarrow}\ \text{(Sets)}\]
commuting\var{\inserted{componentwise}}{\hand{componentwise}}{insertion}{ps_134-5_017-v0001} with (small) direct \var{\xout{sums}}{}{typewriter deletion}{ps_134-5_017-v0002}\var{\inserted[typed]{limits}}{limits}{typed insertion}{ps_134-5_017-v0003}, and inducing an equivalence between
either factor with the category of \var{\xout{\illegible{1}}}{}{typewriter deletion}{ps_134-5_017-v0004} "co-sheaves" on the other, namely
covariant functors to (Sets) commuting with direct limits. But it is'nt
\var{\over{a}{a}}{\hand{a}}{writing over}{ps_134-5_017-v0005}t all clear, starting with an arbitrary topos $\underline{A}$ say, whether the catego-%
ry\var{\inserted[typed]{$\underline{A}\hand{'}$}}{$\underline{A}\hand{'}$}{typed insertion}{ps_134-5_017-v0006} of all cosheaves on $\underline{A}$ is again a topos, and still less whether $\underline{A}$ can
be recovered\var{\inserted{(up to equivalence)}}{\hand{(up to equivalence)}}{insertion}{ps_134-5_017-v0007} in terms of $\underline{A}'$ as the category of all cosheaves on $\underline{A}'$ .

To come back though to the factorization problem raised above (p.202),
the main trouble here is that, except the case of an isomorphism $i\colon X\hand{\xrightarrow{\sim}}Z$,
I was unable to pin down a \var{\inserted[typed]{single}}{single}{typed insertion}{ps_134-5_017-v0008} case of a \var{\xout{functorx}}{}{typewriter deletion}{ps_134-5_017-v0009} map $i$ in (Cat) such that
co-base change by $i$ transforms weak equivalences (in the usual sense say)
into weak equivalences. O\var{\over{n}{n}}{\hand{n}}{writing over}{ps_134-5_017-v0010}e candidate I had in mind, the so-called
``\hand{\uline{\textcolor{black}{open immersions}}}'' , \var{\xout{inducingx}}{}{typewriter deletion}{ps_134-5_017-v0011} namely functors $i\colon X\hand{\longrightarrow}Z$ inducing an isomor-%
phism between $X$ and a ``sieve'' (or ``crible'' ) in $Z$ (corresponding to an
open sub-topos of $Z^{\wedge}$) , and dually the ``closed immersions'', finally have
turned out delusive - a disappointment maybe, but \var{\inserted[typed]{still more}}{still more}{typed insertion}{ps_134-5_017-v0012} a big relief to find
out at last how the score was ! Almost immediately in the wake of this
negative result in (Cat), and in close connection with the fairly well
understood \hand{$\int$} substitute for amalgamated sums in (Cat), came the big com-%
pensation valid\var{\xout{x}}{}{typewriter deletion}{ps_134-5_017-v0013} in any category $A^{\wedge}$ , namely the fact that for a cocar-%
tesian square
% Hand-drawn arrows are marked blue; typed objects remain plain.
\diagram{\begin{tikzcd}
Y \ar[color=blue, r,hook,"i"] \ar[color=blue, d,swap,"g"] \& X \ar[color=blue, d,"\hand{f}"] \\
Y' \ar[color=blue, r,"i'"] \& X'
\end{tikzcd}}
in $A^{\wedge}$, where $i$ is a monomorphism, if $i$ (resp. $g$ ) is in $\underline{W}_A$ , so is
$i'$ (resp. $f$). This implies that \hand{\uline{\textcolor{black}{co-base change by a monomorphism in $A^{\wedge}$}}}
\hand{\uline{\textcolor{black}{transforms weak equivalences into weak equivalences}}}. The common main fact
behind these statements is that for a diagram as above (without assuming
$i$ nor $g$ to be in $\underline{W}_A$ )\hand{,} $X'$ can be interpreted up to weak equivalence as
the``integral\var{\xout{xx}}{}{typewriter deletion}{ps_134-5_017-v0014}'' of the diagram
\diagram{\begin{tikzcd}
Y \ar[color=blue, r] \ar[color=blue, d] \& X \\
Y' \&
\end{tikzcd}}
,
more accurately\hand{,} the natural map in (Cat)
\[\hand{\int}\ \diagram{\begin{tikzcd}
A_{/Y} \ar[color=blue, r] \ar[color=blue, d] \& A_{/X} \\
A_{/Y'} \&
\end{tikzcd}}\ \hand{\longrightarrow}\ A_{/X'}\]
is in $\underline{W}$ .

The corresponding statement in (Cat) itself, even if $i$ is supposed

\margin{17}

%% END SOURCE ps_134-5_017
\ifnum\EditionMode=1
\psapparatusstart
\psentry{ps_134-5_017-v0001}{1}{insertion}{\inserted{componentwise}}{\hand{componentwise}}{The handwritten insertion is incorporated at its indicated position; insertion brackets leave the running text.}
\psentry{ps_134-5_017-v0002}{2}{typewriter deletion}{\xout{sums}}{}{The typewriter strikes out this material; it leaves the corrected running text.}
\psentry{ps_134-5_017-v0003}{3}{typed insertion}{\inserted[typed]{limits}}{limits}{The typed insertion is incorporated at its indicated position; insertion brackets leave the running text.}
\psentry{ps_134-5_017-v0004}{4}{typewriter deletion}{\xout{\illegible{1}}}{}{The typewriter strikes out this material; it leaves the corrected running text.}
\psentry{ps_134-5_017-v0005}{5}{writing over}{\over{a}{a}}{\hand{a}}{The superposed handwritten reading replaces the underlying reading.}
\psentry{ps_134-5_017-v0006}{6}{typed insertion}{\inserted[typed]{$\underline{A}\hand{'}$}}{$\underline{A}\hand{'}$}{The typed insertion is incorporated at its indicated position; insertion brackets leave the running text.}
\psentry{ps_134-5_017-v0007}{7}{insertion}{\inserted{(up to equivalence)}}{\hand{(up to equivalence)}}{The handwritten insertion is incorporated at its indicated position; insertion brackets leave the running text.}
\psentry{ps_134-5_017-v0008}{8}{typed insertion}{\inserted[typed]{single}}{single}{The typed insertion is incorporated at its indicated position; insertion brackets leave the running text.}
\psentry{ps_134-5_017-v0009}{9}{typewriter deletion}{\xout{functorx}}{}{The typewriter strikes out this material; it leaves the corrected running text.}
\psentry{ps_134-5_017-v0010}{10}{writing over}{\over{n}{n}}{\hand{n}}{The superposed handwritten reading replaces the underlying reading.}
\psentry{ps_134-5_017-v0011}{11}{typewriter deletion}{\xout{inducingx}}{}{The typewriter strikes out this material; it leaves the corrected running text.}
\psentry{ps_134-5_017-v0012}{12}{typed insertion}{\inserted[typed]{still more}}{still more}{The typed insertion is incorporated at its indicated position; insertion brackets leave the running text.}
\psentry{ps_134-5_017-v0013}{13}{typewriter deletion}{\xout{x}}{}{The typewriter strikes out this material; it leaves the corrected running text.}
\psentry{ps_134-5_017-v0014}{14}{typewriter deletion}{\xout{xx}}{}{The typewriter strikes out this material; it leaves the corrected running text.}
\psapparatusend
\fi
\pssource{ps_134-5_018}{134-5/18}{https://grothendieck.umontpellier.fr/web/viewer.html?file=../134-5.pdf\#page=19}
\psnumber{ps_134-5_018-v0001}{1}
\psnumber{ps_134-5_018-v0002}{2}
\psnumber{ps_134-5_018-v0003}{3}
\psnumber{ps_134-5_018-v0004}{4}
\psnumber{ps_134-5_018-v0005}{5}
\psnumber{ps_134-5_018-v0006}{6}
\psnumber{ps_134-5_018-v0007}{7}
\psnumber{ps_134-5_018-v0008}{8}
\psnumber{ps_134-5_018-v0009}{9}
\psnumber{ps_134-5_018-v0010}{10}
\psnumber{ps_134-5_018-v0011}{11}
\psnumber{ps_134-5_018-v0012}{12}
\psnumber{ps_134-5_018-v0013}{13}
\psnumber{ps_134-5_018-v0014}{14}
\psnumber{ps_134-5_018-v0015}{15}
\psnumber{ps_134-5_018-v0016}{16}
\psnumber{ps_134-5_018-v0017}{17}
\psnumber{ps_134-5_018-v0018}{18}
\psnumber{ps_134-5_018-v0019}{19}
\psnumber{ps_134-5_018-v0020}{20}
\psnumber{ps_134-5_018-v0021}{21}
\psnumber{ps_134-5_018-v0022}{22}
\psnumber{ps_134-5_018-v0023}{23}
\psnumber{ps_134-5_018-v0024}{24}
%% BEGIN SOURCE ps_134-5_018
% Pursuing Stacks (cotes 134-2 to 134-8): archive cote 134-5, site page 18 (PDF page 19), typescript page 205.
% typescript. Draft compared: ps4.tex:797-864.
% AI-transcribed by Claude Opus 5.5 (claude-opus-5-5, adaptive thinking, effort high), Message Batches API, batch ps2 (msgbatch_01NVwr1Yuo8qJf5rykeGBXPK), 2026-11-12; reviewed against both supplied scan halves; see CHECK_DECISIONS.tsv.
\scanpage{134-5/18}
\mspage{205}

to be an open immersion\var{\inserted{in (Cat)}}{\hand{in (Cat)}}{insertion}{ps_134-5_018-v0001} say (or a closed immersion,  which amounts to the
same by duality),  is definitely false,  in other words:  while\var{\inserted[typed]{the modelizer}}{the modelizer}{typed insertion}{ps_134-5_018-v0002} (Cat)
allows for a remarkably simple description of integration of homotopy
types,  as seen in the previous section,  in the\var{\inserted[typed]{basic}}{basic}{typed insertion}{ps_134-5_018-v0003} case\var{\inserted[typed]{however}}{however}{typed insertion}{ps_134-5_018-v0004} of an integration
% Hand-drawn arrows are marked blue; typed objects remain plain.
\hand{$\int$}\diagram{\begin{tikzcd}[sep=small] Y \ar[color=blue, r,"\hand{i}"] \ar[color=blue, d,"\hand{g}"'] \& X \\ Y' \& \end{tikzcd}} \hand{(}corresponding to \hand{``}amalgamated sums\hand{''}\hand{)}, and even \var{\deleted{\unsure{if}}}{}{deletion}{ps_134-5_018-v0005}\var{\inserted{when}}{\hand{when}}{insertion}{ps_134-5_018-v0006} i: Y\hand{$\rightarrow$}X is an
open or closed immersion,  this operation does definitely \hand{\uline{\textcolor{black}{not}}} correspond
to the operation of taking\var{\inserted[typed]{the}}{the}{typed insertion}{ps_134-5_018-v0007} amalgamated sum\var{\deleted{\unsure{s}}}{}{deletion}{ps_134-5_018-v0008}\var{\inserted[typed]{X\hand{'}}}{X\hand{'}}{typed insertion}{ps_134-5_018-v0009} in (Cat).  \var{\deleted{\unsure{I}}}{}{deletion}{ps_134-5_018-v0010}It does though \margin{a vertical line in the left margin, from this line to ``resp. a cofibration''}
when X\hand{$\rightarrow$}X' is smooth resp. proper,  for instance if it is a fibration
resp. a cofibration,  and this in fact implies the positive result in $A^{\wedge}$
noted above.  \var{\xout{It is possible that it works too if}}{}{typewriter deletion}{ps_134-5_018-v0011}  This condition\var{\inserted{moreover}}{\hand{moreover}}{insertion}{ps_134-5_018-v0012} is satis-%
fied if g: Y\hand{$\rightarrow$}Y' is equally an open resp. closed immersion,  in which ca-%
se the situation is just the \var{\xout{situation}}{}{typewriter deletion}{ps_134-5_018-v0013} one of an ambient X'\hand{,} and two open
resp. closed subobjects X and Y'\hand{,}with intersection Y .  This is a useful
result,  but wholly unsufficient for the factorization problem we were
after in (Cat),  with a view of performing the standard homotopy construc-%
tions in (Cat) itself.  It may be true\var{\inserted{still}}{\hand{still}}{insertion}{ps_134-5_018-v0014} that if i: Y\hand{$\rightarrow$} X is not only an
open or closed immersion,  but a weak equivalence as well,  that then
i': Y' \hand{$\rightarrow$}X' is equally a weak equivalence,  or what amounts\var{\inserted[typed]{to}}{to}{typed insertion}{ps_134-5_018-v0015} the same,  that
X' can be identified up to weak equivalence with the homotopy integral
(which indeed,  up to weak equivalence,  just reduces to Y');  but I have
been unable so far to clear up th\var{\over{e}{is}}{\hand{is}}{writing over}{ps_134-5_018-v0016} matter.  If true,  this would be quite
a useful result,  but still unsufficient it seems in order to carry out \hand{the}
standard homotopy constructions in (Cat) itself.

To sum up,  the main drawback of (Cat) as a modelizer \var{\over{\illegible{1}}{is}}{\hand{is}}{writing over}{ps_134-5_018-v0017} that,  \var{\xout{excep}}{}{typewriter deletion}{ps_134-5_018-v0018}
except in very special cases which are just too restricted,  amalgamated
sums in (Cat) don't have a reasonable meaning in terms of homotopy opera-%
tions  -  whereas in a category $A^{\wedge}$,  a topos indeed where therefore amal-%
gamated sums have as good exactness properties as if working in (Sets),
these amalgmated sums (\var{\deleted{when \unsure{it} correspond\over{ing}{ing} to}}{}{deletion}{ps_134-5_018-v0019}\var{\inserted{for}}{\hand{for}}{insertion}{ps_134-5_018-v0020} a two-arrow diagram with one
arrow a monomorphism) \hand{\uline{\textcolor{black}{do}}} have a homotopy-theoretic meaning.  This finally
seems to force us,  in order to develop some of the basic structure in
Hot($\underline{W}$) ,  to leave the haven of the basic modelizer (Cat),  and work in an
elementary modelizer $A^{\wedge}$\hand{/}instead,  where A is some $\underline{W}$-test category.  This
then brought me back finally to the question:wether these modelizers are
\hand{\uline{\textcolor{black}{closed model categories}}} in Quillen's sense,  when we take of course for
``weak equivalences''  $\underline{W}_A$ ,  and\var{\inserted{moreover}}{\hand{moreover}}{insertion}{ps_134-5_018-v0021} as ``cofibrations'' (in the sense of Quillen'
s set-up)  just the monomorphisms.  Relying heavily upon the result o\var{\over{\unsure{f}}{n}}{\hand{n}}{writing over}{ps_134-5_018-v0022}
monomorphisms in $A^{\wedge}$\hand{/}stated above,  it seems to come out that we do get a
closed model category indeed  -  and even a simplicial model category,
if we \var{\over{\illegible{1}}{a}}{\hand{a}}{writing over}{ps_134-5_018-v0023}re out for this.  \var{\deleted{\unsure{T}}}{}{deletion}{ps_134-5_018-v0024}There is still  a cardinality question to be

\margin{18}

%% END SOURCE ps_134-5_018
\ifnum\EditionMode=1
\psapparatusstart
\psentry{ps_134-5_018-v0001}{1}{insertion}{\inserted{in (Cat)}}{\hand{in (Cat)}}{The handwritten insertion is incorporated at its indicated position; insertion brackets leave the running text.}
\psentry{ps_134-5_018-v0002}{2}{typed insertion}{\inserted[typed]{the modelizer}}{the modelizer}{The typed insertion is incorporated at its indicated position; insertion brackets leave the running text.}
\psentry{ps_134-5_018-v0003}{3}{typed insertion}{\inserted[typed]{basic}}{basic}{The typed insertion is incorporated at its indicated position; insertion brackets leave the running text.}
\psentry{ps_134-5_018-v0004}{4}{typed insertion}{\inserted[typed]{however}}{however}{The typed insertion is incorporated at its indicated position; insertion brackets leave the running text.}
\psentry{ps_134-5_018-v0005}{5}{deletion}{\deleted{\unsure{if}}}{}{The author cancels this material by hand; it leaves the corrected running text.}
\psentry{ps_134-5_018-v0006}{6}{insertion}{\inserted{when}}{\hand{when}}{The handwritten insertion is incorporated at its indicated position; insertion brackets leave the running text.}
\psentry{ps_134-5_018-v0007}{7}{typed insertion}{\inserted[typed]{the}}{the}{The typed insertion is incorporated at its indicated position; insertion brackets leave the running text.}
\psentry{ps_134-5_018-v0008}{8}{deletion}{\deleted{\unsure{s}}}{}{The author cancels this material by hand; it leaves the corrected running text.}
\psentry{ps_134-5_018-v0009}{9}{typed insertion}{\inserted[typed]{X\hand{'}}}{X\hand{'}}{The typed insertion is incorporated at its indicated position; insertion brackets leave the running text.}
\psentry{ps_134-5_018-v0010}{10}{deletion}{\deleted{\unsure{I}}}{}{The author cancels this material by hand; it leaves the corrected running text.}
\psentry{ps_134-5_018-v0011}{11}{typewriter deletion}{\xout{It is possible that it works too if}}{}{The typewriter strikes out this material; it leaves the corrected running text.}
\psentry{ps_134-5_018-v0012}{12}{insertion}{\inserted{moreover}}{\hand{moreover}}{The handwritten insertion is incorporated at its indicated position; insertion brackets leave the running text.}
\psentry{ps_134-5_018-v0013}{13}{typewriter deletion}{\xout{situation}}{}{The typewriter strikes out this material; it leaves the corrected running text.}
\psentry{ps_134-5_018-v0014}{14}{insertion}{\inserted{still}}{\hand{still}}{The handwritten insertion is incorporated at its indicated position; insertion brackets leave the running text.}
\psentry{ps_134-5_018-v0015}{15}{typed insertion}{\inserted[typed]{to}}{to}{The typed insertion is incorporated at its indicated position; insertion brackets leave the running text.}
\psentry{ps_134-5_018-v0016}{16}{writing over}{\over{e}{is}}{\hand{is}}{The superposed handwritten reading replaces the underlying reading.}
\psentry{ps_134-5_018-v0017}{17}{writing over}{\over{\illegible{1}}{is}}{\hand{is}}{The superposed handwritten reading replaces the underlying reading.}
\psentry{ps_134-5_018-v0018}{18}{typewriter deletion}{\xout{excep}}{}{The typewriter strikes out this material; it leaves the corrected running text.}
\psentry{ps_134-5_018-v0019}{19}{deletion}{\deleted{when \unsure{it} correspond\over{ing}{ing} to}}{}{The author cancels this material by hand; it leaves the corrected running text.}
\psentry{ps_134-5_018-v0020}{20}{insertion}{\inserted{for}}{\hand{for}}{The handwritten insertion is incorporated at its indicated position; insertion brackets leave the running text.}
\psentry{ps_134-5_018-v0021}{21}{insertion}{\inserted{moreover}}{\hand{moreover}}{The handwritten insertion is incorporated at its indicated position; insertion brackets leave the running text.}
\psentry{ps_134-5_018-v0022}{22}{writing over}{\over{\unsure{f}}{n}}{\hand{n}}{The superposed handwritten reading replaces the underlying reading.}
\psentry{ps_134-5_018-v0023}{23}{writing over}{\over{\illegible{1}}{a}}{\hand{a}}{The superposed handwritten reading replaces the underlying reading.}
\psentry{ps_134-5_018-v0024}{24}{deletion}{\deleted{\unsure{T}}}{}{The author cancels this material by hand; it leaves the corrected running text.}
\psapparatusend
\fi
\pssource{ps_134-5_019}{134-5/19}{https://grothendieck.umontpellier.fr/web/viewer.html?file=../134-5.pdf\#page=20}
\psnumber{ps_134-5_019-v0001}{1}
\psnumber{ps_134-5_019-v0002}{2}
\psnumber{ps_134-5_019-v0003}{3}
\psnumber{ps_134-5_019-v0004}{4}
\psnumber{ps_134-5_019-v0005}{5}
\psnumber{ps_134-5_019-v0006}{6}
\psnumber{ps_134-5_019-v0007}{7}
\psnumber{ps_134-5_019-v0008}{8}
\psnumber{ps_134-5_019-v0009}{9}
\psnumber{ps_134-5_019-v0010}{10}
\psnumber{ps_134-5_019-v0011}{11}
\psnumber{ps_134-5_019-v0012}{12}
\psnumber{ps_134-5_019-v0013}{13}
\psnumber{ps_134-5_019-v0014}{14}
\psnumber{ps_134-5_019-v0015}{15}
\psnumber{ps_134-5_019-v0016}{16}
\psnumber{ps_134-5_019-v0017}{17}
\psnumber{ps_134-5_019-v0018}{18}
\psnumber{ps_134-5_019-v0019}{19}
\psnumber{ps_134-5_019-v0020}{20}
\psnumber{ps_134-5_019-v0021}{21}
\psnumber{ps_134-5_019-v0022}{22}
\psnumber{ps_134-5_019-v0023}{23}
\psnumber{ps_134-5_019-v0024}{24}
\psnumber{ps_134-5_019-v0025}{25}
\psnumber{ps_134-5_019-v0026}{26}
\psnumber{ps_134-5_019-v0027}{27}
\psnumber{ps_134-5_019-v0028}{28}
%% BEGIN SOURCE ps_134-5_019
% Pursuing Stacks (cotes 134-2 to 134-8): archive cote 134-5, site page 19 (PDF page 20), typescript page 206.
% typescript. Draft compared: ps4.tex:852-914.
% AI-transcribed by Claude Opus 5.5 (claude-opus-5-5, adaptive thinking, effort high), Message Batches API, batch ps2 (msgbatch_01NVwr1Yuo8qJf5rykeGBXPK), 2026-11-12; reviewed against both supplied scan halves; see CHECK_DECISIONS.tsv.
\scanpage{134-5/19}
\mspage{206}
settled to get\var{\inserted[typed]{the Quillen factorizations in}}{the Quillen factorizations in}{typed insertion}{ps_134-5_019-v0001} the general case,  but this should not be too serious a
difficulty I feel.\var{\over{x}{\unsure{\&}}}{\hand{\unsure{\&}}}{writing over}{ps_134-5_019-v0002} What\var{\inserted{however}}{\hand{however}}{insertion}{ps_134-5_019-v0003} makes me still feel a little unhappy in all this\hand{,}
is\var{\inserted[typed]{rather}}{rather}{typed insertion}{ps_134-5_019-v0004} that I did not get a direct proof for an elementary modelizer being a
closed model category  -  I finally have to make a reduction to the known
case of semisimplicial complexes,  settled by Quillen in his notes.  This
\var{\inserted{detour}}{\hand{detour}}{insertion}{ps_134-5_019-v0005} looks rather artificial  -  it is the \var{\xout{only}}{}{typewriter deletion}{ps_134-5_019-v0006}\var{\inserted[typed]{first}}{first}{typed insertion}{ps_134-5_019-v0007} instance,  and presumably the
last one,  where the theory I am digging out seems to depend \var{\over{o}{o}}{\hand{o}}{writing over}{ps_134-5_019-v0008}n semi-simpli
-cial techniques,\var{\inserted{techniques}}{\hand{techniques}}{insertion}{ps_134-5_019-v0009} which moreover I dont\var{\inserted{really}}{\hand{really}}{insertion}{ps_134-5_019-v0010} know and am not really eager to
swallow.  It's just a prejudice maybe,  \var{\xout{but I}}{}{typewriter deletion}{ps_134-5_019-v0011} a block maybe against the
semi-simplicial approach which I never really liked nor assimilated  -
but I do have the feeling that the more refined and specific semi-simplici-%
al techniques and notions (such as minimal fibrations,  used in Quillen's
proof though,  alas !)  are \var{\deleted{\xout{\illegible{1}}}}{}{deletion}{ps_134-5_019-v0012} \hand{irrelevan}t for an understanding of
the main structures entering homotopy theory and homo\hand{to}pical algebra.  As
for the notion of a Kan complex or a Kan fibration  -  namely just a
``fibration'' in Quillen's axiomatic set-up,which I was \var{\deleted{\xout{glad}}}{}{deletion}{ps_134-5_019-v0013} finally glad to
find\hand{,}``ready foruse''  -  I \var{\deleted{finally}}{}{deletion}{ps_134-5_019-v0014} came to convince myself\var{\inserted{at last}}{\hand{at last}}{insertion}{ps_134-5_019-v0015} that it was a
basic notion indeed,  and it was no use trying to by\var{\deleted{e}}{}{deletion}{ps_134-5_019-v0016}-pass it at all price.
Thus I took to the opposite,  and tried to \var{\xout{develop}}{}{typewriter deletion}{ps_134-5_019-v0017} pin down a Quillen\var{\inserted{-}}{\hand{-}}{insertion}{ps_134-5_019-v0018}type
factorization theorem,  and \var{\xout{the}}{}{typewriter deletion}{ps_134-5_019-v0019} his\hand{cha}racteristic sea-saw game \var{\xout{of}}{}{typewriter deletion}{ps_134-5_019-v0020}\var{\inserted[typed]{between}}{between}{typed insertion}{ps_134-5_019-v0021} right
and left lifting properties,  in \var{\deleted{the}}{}{deletion}{ps_134-5_019-v0022}\var{\inserted{as}}{\hand{as}}{insertion}{ps_134-5_019-v0023} great\var{\deleted{est}}{}{deletion}{ps_134-5_019-v0024} generality\var{\inserted{as}}{\hand{as}}{insertion}{ps_134-5_019-v0025} I could manage.

\hand{30.5.}

71.  The \var{\deleted{bulk of}}{}{deletion}{ps_134-5_019-v0026} scratchwork done since last month has of course conside-%
rably cleared up the prospects of my present pre-stacks reflection on
homotopy models,  in which I unsuspectingly embarked three months ago.
\var{\over{\illegible{1}}{Now}}{\hand{Now}}{writing over}{ps_134-5_019-v0027} I would like to sketch a provisional working program for the notes
still ahead.

A)  Write down at last the story of aspericity structures and canonical
modelizers,  as I was about to when I interrupted the notes to do my \var{\deleted{scrat}}{}{deletion}{ps_134-5_019-v0028}
scratchwork.

B)  Study the basic modelizer (Cat),  and the common properties of
elementary modelizers $A^{\wedge}$,  with a main emphasis upon base change and
co-base change properties,  and upon Quillen type factorization questions.
Here it will be useful to dwell somewhat on the ``homotopy integral'' variant
of taking amalgamated sums in (Cat),  on the analoguous constructions for
topoi,  and how these compare to the usual amalgamated sums, including the
interesting case of topological spaces.  It turns out that the homotopy
\margin{19}

%% END SOURCE ps_134-5_019
\ifnum\EditionMode=1
\psapparatusstart
\psentry{ps_134-5_019-v0001}{1}{typed insertion}{\inserted[typed]{the Quillen factorizations in}}{the Quillen factorizations in}{The typed insertion is incorporated at its indicated position; insertion brackets leave the running text.}
\psentry{ps_134-5_019-v0002}{2}{writing over}{\over{x}{\unsure{\&}}}{\hand{\unsure{\&}}}{The superposed handwritten reading replaces the underlying reading.}
\psentry{ps_134-5_019-v0003}{3}{insertion}{\inserted{however}}{\hand{however}}{The handwritten insertion is incorporated at its indicated position; insertion brackets leave the running text.}
\psentry{ps_134-5_019-v0004}{4}{typed insertion}{\inserted[typed]{rather}}{rather}{The typed insertion is incorporated at its indicated position; insertion brackets leave the running text.}
\psentry{ps_134-5_019-v0005}{5}{insertion}{\inserted{detour}}{\hand{detour}}{The handwritten insertion is incorporated at its indicated position; insertion brackets leave the running text.}
\psentry{ps_134-5_019-v0006}{6}{typewriter deletion}{\xout{only}}{}{The typewriter strikes out this material; it leaves the corrected running text.}
\psentry{ps_134-5_019-v0007}{7}{typed insertion}{\inserted[typed]{first}}{first}{The typed insertion is incorporated at its indicated position; insertion brackets leave the running text.}
\psentry{ps_134-5_019-v0008}{8}{writing over}{\over{o}{o}}{\hand{o}}{The superposed handwritten reading replaces the underlying reading.}
\psentry{ps_134-5_019-v0009}{9}{insertion}{\inserted{techniques}}{\hand{techniques}}{The handwritten insertion is incorporated at its indicated position; insertion brackets leave the running text.}
\psentry{ps_134-5_019-v0010}{10}{insertion}{\inserted{really}}{\hand{really}}{The handwritten insertion is incorporated at its indicated position; insertion brackets leave the running text.}
\psentry{ps_134-5_019-v0011}{11}{typewriter deletion}{\xout{but I}}{}{The typewriter strikes out this material; it leaves the corrected running text.}
\psentry{ps_134-5_019-v0012}{12}{deletion}{\deleted{\xout{\illegible{1}}}}{}{The author cancels this material by hand; it leaves the corrected running text.}
\psentry{ps_134-5_019-v0013}{13}{deletion}{\deleted{\xout{glad}}}{}{The author cancels this material by hand; it leaves the corrected running text.}
\psentry{ps_134-5_019-v0014}{14}{deletion}{\deleted{finally}}{}{The author cancels this material by hand; it leaves the corrected running text.}
\psentry{ps_134-5_019-v0015}{15}{insertion}{\inserted{at last}}{\hand{at last}}{The handwritten insertion is incorporated at its indicated position; insertion brackets leave the running text.}
\psentry{ps_134-5_019-v0016}{16}{deletion}{\deleted{e}}{}{The author cancels this material by hand; it leaves the corrected running text.}
\psentry{ps_134-5_019-v0017}{17}{typewriter deletion}{\xout{develop}}{}{The typewriter strikes out this material; it leaves the corrected running text.}
\psentry{ps_134-5_019-v0018}{18}{insertion}{\inserted{-}}{\hand{-}}{The handwritten insertion is incorporated at its indicated position; insertion brackets leave the running text.}
\psentry{ps_134-5_019-v0019}{19}{typewriter deletion}{\xout{the}}{}{The typewriter strikes out this material; it leaves the corrected running text.}
\psentry{ps_134-5_019-v0020}{20}{typewriter deletion}{\xout{of}}{}{The typewriter strikes out this material; it leaves the corrected running text.}
\psentry{ps_134-5_019-v0021}{21}{typed insertion}{\inserted[typed]{between}}{between}{The typed insertion is incorporated at its indicated position; insertion brackets leave the running text.}
\psentry{ps_134-5_019-v0022}{22}{deletion}{\deleted{the}}{}{The author cancels this material by hand; it leaves the corrected running text.}
\psentry{ps_134-5_019-v0023}{23}{insertion}{\inserted{as}}{\hand{as}}{The handwritten insertion is incorporated at its indicated position; insertion brackets leave the running text.}
\psentry{ps_134-5_019-v0024}{24}{deletion}{\deleted{est}}{}{The author cancels this material by hand; it leaves the corrected running text.}
\psentry{ps_134-5_019-v0025}{25}{insertion}{\inserted{as}}{\hand{as}}{The handwritten insertion is incorporated at its indicated position; insertion brackets leave the running text.}
\psentry{ps_134-5_019-v0026}{26}{deletion}{\deleted{bulk of}}{}{The author cancels this material by hand; it leaves the corrected running text.}
\psentry{ps_134-5_019-v0027}{27}{writing over}{\over{\illegible{1}}{Now}}{\hand{Now}}{The superposed handwritten reading replaces the underlying reading.}
\psentry{ps_134-5_019-v0028}{28}{deletion}{\deleted{scrat}}{}{The author cancels this material by hand; it leaves the corrected running text.}
\psapparatusend
\fi
\pssource{ps_134-5_020}{134-5/20}{https://grothendieck.umontpellier.fr/web/viewer.html?file=../134-5.pdf\#page=21}
\psnumber{ps_134-5_020-v0001}{1}
\psnumber{ps_134-5_020-v0002}{2}
\psnumber{ps_134-5_020-v0003}{3}
\psnumber{ps_134-5_020-v0004}{4}
\psnumber{ps_134-5_020-v0005}{5}
\psnumber{ps_134-5_020-v0006}{6}
\psnumber{ps_134-5_020-v0007}{7}
\psnumber{ps_134-5_020-v0008}{8}
\psnumber{ps_134-5_020-v0009}{9}
\psnumber{ps_134-5_020-v0010}{10}
\psnumber{ps_134-5_020-v0011}{11}
\psnumber{ps_134-5_020-v0012}{12}
%% BEGIN SOURCE ps_134-5_020
% Pursuing Stacks (cotes 134-2 to 134-8): archive cote 134-5, site page 20 (PDF page 21), typescript page 207.
% typescript. Draft compared: ps4.tex:902-963.
% AI-transcribed by Claude Opus 5.5 (claude-opus-5-5, adaptive thinking, effort high), Message Batches API, batch ps2 (msgbatch_01NVwr1Yuo8qJf5rykeGBXPK), 2026-11-12; reviewed against both supplied scan halves; see CHECK_DECISIONS.tsv.
\scanpage{134-5/20}
\mspage{207}

integral variant for amalgamated sums is essentially \var{\over{c}{ch}}{\hand{ch}}{writing over}{ps_134-5_020-v0001}aracterized by a
Mayer-Vietoris type long exact sequence for cohomology,  and the cases
when the homotopy construction turns out to be equivalent to usual amalga-%
mated sums,  are  just those when such a Mayer-Vietoris sequence exists
for the latter.  An interesting and typical case is for topological spaces,
taking the amalgamated sum for a diagram
% Hand-drawn arrows are marked blue; typed objects remain plain.
\diagram{\begin{tikzcd}[sep=small]
Y \ar[color=blue, r,hook,"i"] \ar[color=blue, d,swap,"g"] \& X \\
Y' \&
\end{tikzcd}\quad ,}
when $i$ is a closed immersion and $g$ is proper  (which is also the basic type
of amalgamations which occur in the ``unfo\var{\over[typed]{r}{l}}{l}{typewriter overwriting}{ps_134-5_020-v0002}ding'' of stratified structures\hand{)}.

In the course of the last weeks' reflections,  there has taken place
also a substantial clarification concerning the relevant properties of a
\var{basic localizer $\underline{W}$ and how most of these\hand{,} follow from just the first three,\hand{$\swarrow$}
% the words "follow from just the first three," are circled by hand; a hand-drawn arrow carries them to after "of W," in the next line
including \inserted[typed]{\hand{\uline{\textcolor{black}{strong}}}} saturation of $\underline{W}$\hand{,} }{basic localizer $\underline{W}$ and how most of these\hand{,}
% the words "follow from just the first three," are circled by hand; a hand-drawn arrow carries them to after "of W," in the next line
including \hand{\uline{\textcolor{black}{strong}}} saturation of $\underline{W}$\hand{,} follow from just the first three, }{transposition}{ps_134-5_020-v0003}(a question which kept turning up like a nuisan-%
ce throughout the notes !).  This should be among the very first things to
write down in this part of the reflections,  as $\underline{W}$ after all is \hand{\uline{\textcolor{black}{the}}} one
axiomatic data upon which the whole set-up depends.

C)   A reflexion on the main common features of the various contexts
met with so far having a ``homotopy theory'' flavor,  with a hope to \var{\xout{be able}}{}{typewriter deletion}{ps_134-5_020-v0004}
work out at least some of the main features of an all-encompassing new
structure,  along the lines of Verdier's (commutative) theory of derived
categories and triangulated categories.  The basic idea here,  for the time
being,  seems to be the notion of a \hand{\uline{\textcolor{black}{derivator}}},  which should account for
all the kind of structure dealt with in Verdier's set-up,  as we\var{\over[typed]{e}{l}}{l}{typewriter overwriting}{ps_134-5_020-v0005}l as in
Deligne's and Illusie\hand{'}s later elaborations.  T\var{\over[typed]{h}{h}}{h}{typewriter overwriting}{ps_134-5_020-v0006}ere seems to be however some
important extra features which are not accounted for by the mere derivator,
such as external Hom's with values in (Hot) or some closely related cate-%
gory,  and the formalism of basic invariants (such as $\hand{\pi}_i$ , $H_i$ or $H^i$ ),
with values in suitable categories (often abelian one's),  which among
others allow to check weak equivalence.  Such features seem to be invariab-%
ly around in all cases I know of,  and they need to be understood I feel.

Coming back to (Hot) itself and to modelizers \var{\inserted[typed]{$(M,W)$}}{$(M,W)$}{typed insertion}{ps_134-5_020-v0007} giving rise to it,
there is the puzzling question about when exactly can we assert that
taking group objects of $M$ and weak equivalences between these (namely
group\var{\inserted[typed]{-object}}{-object}{typed insertion}{ps_134-5_020-v0008} homomorphisms which are also in $W$),  we get by localization a cate-%
gory equivalent to the category \var{\deleted{\unsure{(Hot.)}}}{}{deletion}{ps_134-5_020-v0009} of pointed \var{\inserted[typed]{0-connected}}{0-connected}{typed insertion}{ps_134-5_020-v0010} homotopy types.  This is
a well known basic fact when we take as models semisimplicial complexes or
(I \var{\over[typed]{\unsure{q}}{g}}{g}{typewriter overwriting}{ps_134-5_020-v0011}uess) topological spaces - \var{\inserted{a fact,}}{\hand{a fact,}}{insertion}{ps_134-5_020-v0012} closely connected to the game of associa-%

\margin{20}

%% END SOURCE ps_134-5_020
\ifnum\EditionMode=1
\psapparatusstart
\psentry{ps_134-5_020-v0001}{1}{writing over}{\over{c}{ch}}{\hand{ch}}{The superposed handwritten reading replaces the underlying reading.}
\psentry{ps_134-5_020-v0002}{2}{typewriter overwriting}{\over[typed]{r}{l}}{l}{The superposed typed reading replaces the underlying reading.}
\psentry{ps_134-5_020-v0003}{3}{transposition}{basic localizer $\underline{W}$ and how most of these\hand{,} follow from just the first three,\hand{$\swarrow$}
% the words "follow from just the first three," are circled by hand; a hand-drawn arrow carries them to after "of W," in the next line
including \inserted[typed]{\hand{\uline{\textcolor{black}{strong}}}} saturation of $\underline{W}$\hand{,} }{basic localizer $\underline{W}$ and how most of these\hand{,}
% the words "follow from just the first three," are circled by hand; a hand-drawn arrow carries them to after "of W," in the next line
including \hand{\uline{\textcolor{black}{strong}}} saturation of $\underline{W}$\hand{,} follow from just the first three, }{A handwritten circle and arrow move the clause follow from just the first three to after including strong saturation of W; the typed order is retained in the diplomatic text.}
\psentry{ps_134-5_020-v0004}{4}{typewriter deletion}{\xout{be able}}{}{The typewriter strikes out this material; it leaves the corrected running text.}
\psentry{ps_134-5_020-v0005}{5}{typewriter overwriting}{\over[typed]{e}{l}}{l}{The superposed typed reading replaces the underlying reading.}
\psentry{ps_134-5_020-v0006}{6}{typewriter overwriting}{\over[typed]{h}{h}}{h}{The superposed typed reading replaces the underlying reading.}
\psentry{ps_134-5_020-v0007}{7}{typed insertion}{\inserted[typed]{$(M,W)$}}{$(M,W)$}{The typed insertion is incorporated at its indicated position; insertion brackets leave the running text.}
\psentry{ps_134-5_020-v0008}{8}{typed insertion}{\inserted[typed]{-object}}{-object}{The typed insertion is incorporated at its indicated position; insertion brackets leave the running text.}
\psentry{ps_134-5_020-v0009}{9}{deletion}{\deleted{\unsure{(Hot.)}}}{}{The author cancels this material by hand; it leaves the corrected running text.}
\psentry{ps_134-5_020-v0010}{10}{typed insertion}{\inserted[typed]{0-connected}}{0-connected}{The typed insertion is incorporated at its indicated position; insertion brackets leave the running text.}
\psentry{ps_134-5_020-v0011}{11}{typewriter overwriting}{\over[typed]{\unsure{q}}{g}}{g}{The superposed typed reading replaces the underlying reading.}
\psentry{ps_134-5_020-v0012}{12}{insertion}{\inserted{a fact,}}{\hand{a fact,}}{The handwritten insertion is incorporated at its indicated position; insertion brackets leave the running text.}
\psapparatusend
\fi
\pssource{ps_134-5_021}{134-5/21}{https://grothendieck.umontpellier.fr/web/viewer.html?file=../134-5.pdf\#page=22}
\psnumber{ps_134-5_021-v0001}{1}
\psnumber{ps_134-5_021-v0002}{2}
\psnumber{ps_134-5_021-v0003}{3}
\psnumber{ps_134-5_021-v0004}{4}
\psnumber{ps_134-5_021-v0005}{5}
\psnumber{ps_134-5_021-v0006}{6}
\psnumber{ps_134-5_021-v0007}{7}
\psnumber{ps_134-5_021-v0008}{8}
\psnumber{ps_134-5_021-v0009}{9}
\psnumber{ps_134-5_021-v0010}{10}
\psnumber{ps_134-5_021-v0011}{11}
\psnumber{ps_134-5_021-v0012}{12}
\psnumber{ps_134-5_021-v0013}{13}
%% BEGIN SOURCE ps_134-5_021
% Pursuing Stacks (cotes 134-2 to 134-8): archive cote 134-5, site page 21 (PDF page 22), typescript page 208.
% typescript. Draft compared: ps4.tex:950-1010.
% AI-transcribed by Claude Opus 5.5 (claude-opus-5-5, adaptive thinking, effort high), Message Batches API, batch ps2 (msgbatch_01NVwr1Yuo8qJf5rykeGBXPK), 2026-11-12; reviewed against both supplied scan halves; see CHECK_DECISIONS.tsv.
\scanpage{134-5/21}
\mspage{208}

ting to any topological group it's``\hand{\uline{\textcolor{black}{classifying space}}}'',  defined ``up to homo-%
topy''.  I suspect the same should hold in any elementary modelizer $A^{\wedge}$,
$A$ a test category,  at least in the ``strict'' case\var{\inserted{namely,}}{\hand{namely,}}{insertion}{ps_134-5_021-v0001} when $A^{\wedge}$\hand{/}is totally asphe-%
ric,  i.e.  the canonical functor  $A^{\wedge}\hand{\longrightarrow}$ (Hot)  commutes to finite products.
The corresponding statement for the basic localizer (Cat) itself however is
definitely false\unsure{.|,}  Group objects in (Cat) are indeed very interesting and
well-known beings  (introduced,  I understand from Ronnie Brown, by Henry
Whitehead long time ago,  under the somewhat misleading name of ``crossed
modules''),  yet they embody not arbitrary pointed 0-connected homotopy
types,  but merely those for which \hand{$\pi$}$_i(X) = 0$ for $i \hand{>} 2$ .  Thus we get only
2-truncated homotopy types  -  and presumably,  \var{\xout{using}}{}{typewriter deletion}{ps_134-5_021-v0002} starting with n-Cat
instead of (Cat) as a modelizer,  we then should get \unsure{(n+1)}-truncated homoto-%
py types.  This ties in with the observation that taking group objects\var{\inserted{either}}{\hand{either}}{insertion}{ps_134-5_021-v0003} in
(Cat),  or in the full subcategory  (Groupoids)  of the latter,  amounts to
the same  -  and similarly surely for (n-Cat);  on the other hand it has
been kind of clear from the very beginning of this reflection that \var{\xout{always}}{}{typewriter deletion}{ps_134-5_021-v0004}\var{\inserted[typed]{at any}}{at any}{typed insertion}{ps_134-5_021-v0005}
\hand{\uline{\textcolor{black}{finite}}} level,  groupoids and n-groupoids (n finite) will only yield \hand{\uline{\textcolor{black}{trunca-}}}%
\hand{\uline{\textcolor{black}{ted}}} homotopy types.

A related intriguing question is when exactly does a modelizer (M,W)
give rise to a \hand{\uline{\textcolor{black}{Dold-Puppe theorem}}}  -  namely when do we get an actual \hand{\uline{\textcolor{black}{equi-}}}%
\hand{\uline{\textcolor{black}{valence}}} between the category of \hand{\uline{\textcolor{black}{abelian}}} group objects of M,  and the catego-%
ry of chain complexes of abelian groups ?  \var{\over{t}{T}}{\hand{T}}{writing over}{ps_134-5_021-v0006}he original statement was in
case  $M = \hand{\Delta}^{\wedge} =$  semisimplicial complexes,  and doubtlessly it was one main
impetus for the sudden invasion of homotopy and cohomology by semisimplicial
calculus  -  so much so it seems that for many people,  ``homotopy'' has be-%
come synonymous to ``semisimplicial algebra''.  The impression that semi-sim-%
plicial complexes is the \var{\over[typed]{g}{G}}{G}{typewriter overwriting}{ps_134-5_021-v0007}od-given ground for doing homotopy and even cohomo-%
logy,  comes out rather strong also in Quillen's foundational notes, and
\var{\inserted[typed]{in}}{in}{typed insertion}{ps_134-5_021-v0008} Illusie\hand{'}s thesis.  Still there are too some cubical theory chaps I heard,
who surely must have noticed long ago that the Dold-Puppe theorem is valid
equally \var{\xout{in the}}{}{typewriter deletion}{ps_134-5_021-v0009} for cubical complexes (I could hardly imagine that it\var{\inserted{possibly}}{\hand{possibly}}{insertion}{ps_134-5_021-v0010} \hand{co}uld'
\var{\xout{not be}}{}{typewriter deletion}{ps_134-5_021-v0011} nt).  Now it turns out that semisimplicial and cubical complexes are
part of a trilogy,  together with\var{\inserted{so-called}}{\hand{so-called}}{insertion}{ps_134-5_021-v0012} ``\hand{\uline{\textcolor{black}{hemispherical}}} \hand{\uline{\textcolor{black}{complexes}}}'' ,  which look
a lot simpler still,  with just two boundary operators and one degeneracy
in each dimension.  They can be viewed as  embodying the ``primitive struc-%
ture'' of an \hand{$\infty$}- groupoid,  the boundary operators being the``source'' and
``target'' maps,  and the degeneracy the map associating to any i-object the
corresponding ``identity''.  \hand{(} I hit upon this structure among the first
examples of test categories and elementary modelizers,  and have been told
since by Ronnie Brown that he has\var{\inserted{already}}{\hand{already}}{insertion}{ps_134-5_021-v0013} \hand{known} for a while these models,  under

\margin{21}

%% END SOURCE ps_134-5_021
\ifnum\EditionMode=1
\psapparatusstart
\psentry{ps_134-5_021-v0001}{1}{insertion}{\inserted{namely,}}{\hand{namely,}}{The handwritten insertion is incorporated at its indicated position; insertion brackets leave the running text.}
\psentry{ps_134-5_021-v0002}{2}{typewriter deletion}{\xout{using}}{}{The typewriter strikes out this material; it leaves the corrected running text.}
\psentry{ps_134-5_021-v0003}{3}{insertion}{\inserted{either}}{\hand{either}}{The handwritten insertion is incorporated at its indicated position; insertion brackets leave the running text.}
\psentry{ps_134-5_021-v0004}{4}{typewriter deletion}{\xout{always}}{}{The typewriter strikes out this material; it leaves the corrected running text.}
\psentry{ps_134-5_021-v0005}{5}{typed insertion}{\inserted[typed]{at any}}{at any}{The typed insertion is incorporated at its indicated position; insertion brackets leave the running text.}
\psentry{ps_134-5_021-v0006}{6}{writing over}{\over{t}{T}}{\hand{T}}{The superposed handwritten reading replaces the underlying reading.}
\psentry{ps_134-5_021-v0007}{7}{typewriter overwriting}{\over[typed]{g}{G}}{G}{The superposed typed reading replaces the underlying reading.}
\psentry{ps_134-5_021-v0008}{8}{typed insertion}{\inserted[typed]{in}}{in}{The typed insertion is incorporated at its indicated position; insertion brackets leave the running text.}
\psentry{ps_134-5_021-v0009}{9}{typewriter deletion}{\xout{in the}}{}{The typewriter strikes out this material; it leaves the corrected running text.}
\psentry{ps_134-5_021-v0010}{10}{insertion}{\inserted{possibly}}{\hand{possibly}}{The handwritten insertion is incorporated at its indicated position; insertion brackets leave the running text.}
\psentry{ps_134-5_021-v0011}{11}{typewriter deletion}{\xout{not be}}{}{The typewriter strikes out this material; it leaves the corrected running text.}
\psentry{ps_134-5_021-v0012}{12}{insertion}{\inserted{so-called}}{\hand{so-called}}{The handwritten insertion is incorporated at its indicated position; insertion brackets leave the running text.}
\psentry{ps_134-5_021-v0013}{13}{insertion}{\inserted{already}}{\hand{already}}{The handwritten insertion is incorporated at its indicated position; insertion brackets leave the running text.}
\psapparatusend
\fi
\pssource{ps_134-5_022}{134-5/22}{https://grothendieck.umontpellier.fr/web/viewer.html?file=../134-5.pdf\#page=23}
\psnumber{ps_134-5_022-v0001}{1}
\psnumber{ps_134-5_022-v0002}{2}
\psnumber{ps_134-5_022-v0003}{3}
\psnumber{ps_134-5_022-v0004}{4}
\psnumber{ps_134-5_022-v0005}{5}
\psnumber{ps_134-5_022-v0006}{6}
\psnumber{ps_134-5_022-v0007}{7}
\psnumber{ps_134-5_022-v0008}{8}
\psnumber{ps_134-5_022-v0009}{9}
\psnumber{ps_134-5_022-v0010}{10}
\psnumber{ps_134-5_022-v0011}{11}
\psnumber{ps_134-5_022-v0012}{12}
\psnumber{ps_134-5_022-v0013}{13}
%% BEGIN SOURCE ps_134-5_022
% Pursuing Stacks (cotes 134-2 to 134-8): archive cote 134-5, site page 22 (PDF page 23), typescript page 209.
% typescript. Draft compared: ps4.tex:998-1055.
% AI-transcribed by Claude Opus 5.5 (claude-opus-5-5, adaptive thinking, effort high), Message Batches API, batch ps2 (msgbatch_01NVwr1Yuo8qJf5rykeGBXPK), 2026-11-12; reviewed against both supplied scan halves; see CHECK_DECISIONS.tsv.
\scanpage{134-5/22}
\mspage{209}
the similar name of ``globular complexes''). Roughly speaking, it can be
said that the three types of complexes correspond to three ``series'' of
\hand{\uline{\textcolor{black}{regular}}} \hand{\uline{\textcolor{black}{cellular}}} \hand{\uline{\textcolor{black}{subdivision\var{\over{s}{s}}{\hand{s}}{writing over}{ps_134-5_022-v0001}}}} of all spheres $S_n$ \var{\deleted{\unsure{,n}}}{}{deletion}{ps_134-5_022-v0002} where the two-dimensio-%
nal pieces are respectively (by increasing order of ``intricacy'') bigons,
triangles, and squares. It should'nt be hard to show these are the only
series of regular cellular subdivisions of all spheres (one in each dimensi-%
on), such that for any cell of such subdivision, the induced subdivision
of the bounding sphere should still be in the series (up to isomorphism).
The existence moreover of suitable ``degeneracy'' maps, which merit a care-%
ful general definition in this context of cellular subdivisions of spheres,
is an important \var{\inserted[typed]{common extra}}{common extra}{typed insertion}{ps_134-5_022-v0003} \var{\xout{striking}}{}{typewriter deletion}{ps_134-5_022-v0004} feature of the three basic contexts, whose exact
significance I have not quite understood still. \hand{\illegible{1}} To come back to Dold-%
Puppe, sure enough it is still valid in the hemispherical context. Wri-%
ting down the equivalence of categories in explicit terms comes out with
baffling simplicity. I wrote it down without even looking for it, in a
letter to Ronnie Brown, while explaining in a PS the ``yoga'' of associating
to a chain complex of abelian groups an $\infty$-groupoid with additive struc-%
ture (a so-called ``Picard\var{\xout{xx}}{}{typewriter deletion}{ps_134-5_022-v0005} category'' but within the context of $\infty$-catego-%
ries or $\infty$-groupoids\hand{,} rather than usual categories).

On the other hand, taking multico\var{\inserted[typed]{m}}{m}{typed insertion}{ps_134-5_022-v0006}plexes instead of simple ones, \var{\xout{whi}}{}{typewriter deletion}{ps_134-5_022-v0007}
which still can be interpreted as working in a category $A\hat{}$ for suitable
test category $A$ (namely, a product category of categories of the type
\hand{$\Box$}, \hand{$\triangle$} and \hand{$\Diamond$}), it is clear that Dold-Puppe's theorem \var{\deleted{is no longer true}}{}{deletion}{ps_134-5_022-v0008}
as originally stated \var{\inserted{is no longer true \deleted{\unsure{in}} these;}}{\hand{is no longer true  these;}}{insertion}{ps_134-5_022-v0009} in such case the category of abelian group objects
of $A\hat{}$ is equivalent to a category of \hand{\uline{\textcolor{black}{multiple}}} chain complexes. This shows
that definitely, among all possible elementary modelizers, the three in
our trilogy are distinguished indeed, as giving rise to a Dold-Puppe
theorem. A thorough understanding of this theorem would imply, I feel,
\var{\inserted{an understanding}}{\hand{an understanding}}{insertion}{ps_134-5_022-v0010} of which exactly are the modelizers, or elementary modelizers at any rate,
giving rise to such a theorem. I would'nt be too surprised if it turned
out that the three we got are the only ones, up to equivalence.

Another common feature of these modelizers, is that they allow for
a sweeping computational description of cohomology (or homology) invariants
in terms of the so-called\hand{``}boundary operations\hand{''}. This is visibly connected
to the \hand{(}strongly intuitive\hand{)} t\var{\over{y}{i}}{\hand{i}}{writing over}{ps_134-5_022-v0011}\var{\over{e}{e}}{\hand{e}}{writing over}{ps_134-5_022-v0012} between these kind of models, \var{\over{\illegible{1}}{and}}{\hand{and}}{writing over}{ps_134-5_022-v0013} cellular
subdivisions of spheres. However, at this level (unlike the Dold-Puppe
story) the regularity feature of the cellular subdivisions we got, and
the fact that we allow for just one (up to isomorphism) in each dimension,

\margin{22}

%% END SOURCE ps_134-5_022
\ifnum\EditionMode=1
\psapparatusstart
\psentry{ps_134-5_022-v0001}{1}{writing over}{\over{s}{s}}{\hand{s}}{The superposed handwritten reading replaces the underlying reading.}
\psentry{ps_134-5_022-v0002}{2}{deletion}{\deleted{\unsure{,n}}}{}{The author cancels this material by hand; it leaves the corrected running text.}
\psentry{ps_134-5_022-v0003}{3}{typed insertion}{\inserted[typed]{common extra}}{common extra}{The typed insertion is incorporated at its indicated position; insertion brackets leave the running text.}
\psentry{ps_134-5_022-v0004}{4}{typewriter deletion}{\xout{striking}}{}{The typewriter strikes out this material; it leaves the corrected running text.}
\psentry{ps_134-5_022-v0005}{5}{typewriter deletion}{\xout{xx}}{}{The typewriter strikes out this material; it leaves the corrected running text.}
\psentry{ps_134-5_022-v0006}{6}{typed insertion}{\inserted[typed]{m}}{m}{The typed insertion is incorporated at its indicated position; insertion brackets leave the running text.}
\psentry{ps_134-5_022-v0007}{7}{typewriter deletion}{\xout{whi}}{}{The typewriter strikes out this material; it leaves the corrected running text.}
\psentry{ps_134-5_022-v0008}{8}{deletion}{\deleted{is no longer true}}{}{The author cancels this material by hand; it leaves the corrected running text.}
\psentry{ps_134-5_022-v0009}{9}{insertion}{\inserted{is no longer true \deleted{\unsure{in}} these;}}{\hand{is no longer true  these;}}{The handwritten insertion is incorporated at its indicated position; insertion brackets leave the running text.}
\psentry{ps_134-5_022-v0010}{10}{insertion}{\inserted{an understanding}}{\hand{an understanding}}{The handwritten insertion is incorporated at its indicated position; insertion brackets leave the running text.}
\psentry{ps_134-5_022-v0011}{11}{writing over}{\over{y}{i}}{\hand{i}}{The superposed handwritten reading replaces the underlying reading.}
\psentry{ps_134-5_022-v0012}{12}{writing over}{\over{e}{e}}{\hand{e}}{The superposed handwritten reading replaces the underlying reading.}
\psentry{ps_134-5_022-v0013}{13}{writing over}{\over{\illegible{1}}{and}}{\hand{and}}{The superposed handwritten reading replaces the underlying reading.}
\psapparatusend
\fi
\pssource{ps_134-5_023}{134-5/23}{https://grothendieck.umontpellier.fr/web/viewer.html?file=../134-5.pdf\#page=24}
\psnumber{ps_134-5_023-v0001}{1}
\psnumber{ps_134-5_023-v0002}{2}
\psnumber{ps_134-5_023-v0003}{3}
\psnumber{ps_134-5_023-v0004}{4}
\psnumber{ps_134-5_023-v0005}{5}
\psnumber{ps_134-5_023-v0006}{6}
\psnumber{ps_134-5_023-v0007}{7}
\psnumber{ps_134-5_023-v0008}{8}
\psnumber{ps_134-5_023-v0009}{9}
\psnumber{ps_134-5_023-v0010}{10}
\psnumber{ps_134-5_023-v0011}{11}
\psnumber{ps_134-5_023-v0012}{12}
\psnumber{ps_134-5_023-v0013}{13}
\psnumber{ps_134-5_023-v0014}{14}
\psnumber{ps_134-5_023-v0015}{15}
%% BEGIN SOURCE ps_134-5_023
% Pursuing Stacks (cotes 134-2 to 134-8): archive cote 134-5, site page 23 (PDF page 24), typescript page 210.
% typescript. Draft compared: ps4.tex:1043-1101.
% AI-transcribed by Claude Opus 5.5 (claude-opus-5-5, adaptive thinking, effort high), Message Batches API, batch ps2 (msgbatch_01NVwr1Yuo8qJf5rykeGBXPK), 2026-11-12; reviewed against both supplied scan halves; see CHECK_DECISIONS.tsv.
\scanpage{134-5/23}
\mspage{210}

seems to be irrelevant.  It might be worth while to write down with care
what exactly is needed,  in order to define,  in terms of a bunch of
cellular structures\var{\xout{,}}{}{typewriter deletion}{ps_134-5_023-v0001} of spheres,  a corresponding test-category (hopefully
even a strict one),  and a way of computing in terms of boundary operators
the cohomology of the corresponding models.  The more delicate point here
may be that if we really want to get an actual test category,  not just a
\var{\xout{strict}}{}{typewriter deletion}{ps_134-5_023-v0002} weak one,  we should have \var{\inserted[typed]{``enough''}}{``enough''}{typed insertion}{ps_134-5_023-v0003} degeneracy maps between our cellular struc-%
tures,  which might well prove an extremely stringent requirement.  \var{\over{A}{A}}{\hand{A}}{writing over}{ps_134-5_023-v0004}gain,
it will be interesting if we can meet it otherwise than just sticking to
our trilogy.

D) ``Back to topoi''.  They have been my main intuitive leading thread in
the reflections so far,  but have remained somewhat implicit most times.
When working with categories  $X$ as models for homotopy types,  we have been
thinking in reality of the associated topos $X\hat{}$.  In the same way,  when re-%
lativizing \var{the \deleted{(Hot)} construction \inserted{of (Hot)} as a derived category over an object $I$
of (Cat)\hand{,} \hand{$\swarrow$}%
% The handwritten circle and arrow move the construction-of-Hot phrase to this position.
 }{over an object $I$
of (Cat) the  construction \hand{of (Hot)} as a derived category\hand{,} %
% The handwritten circle and arrow move the construction-of-Hot phrase to this position.
 }{transposition}{ps_134-5_023-v0005}namely working with categories over $I$ as models for a derived
category $D(I)$,  the leading intuition again has been to look at $X$ over $I$
as the topos $X\hat{}$\hand{/}over $I\hat{}$.  This gives strong suggestions as to defining
$D(I)$ not only in terms of a given category $I$ ,  but also in terms of an
arbitrary topos $T$  as well,  standing for $I\hat{}$,  and look \var{\over{to}{to}}{\hand{to}}{writing over}{ps_134-5_023-v0006} what extent the
$f^{*},f_{\hand{!}},f_{*}$  formalism of derivators extends to the case when $f$ is a ``map''
between topoi.  The definition of \var{\inserted[typed]{some}}{some}{typed insertion}{ps_134-5_023-v0007} $D(T)$ for a topos $T$ can be given in a
number of ways it seems.  As far as I know,  the only one which has been
\var{\inserted{written and}}{\hand{written and}}{insertion}{ps_134-5_023-v0008} used so far consists in taking \var{\xout{the}}{}{typewriter deletion}{ps_134-5_023-v0009} a suitable derived category of the cate-%
gory of semisimplicial objects of $T$ .  This is done in Illusie's thesis,
where there is no mention though of $f_{\hand{!}}$ or $f_{*}$ functors  -  which,  maybe,
should rather be written $Lf_{\hand{!}}$  and $Rf_{*}$ ,  suggesting they are respectively
left and right derived functors of the more familiar $f_{\hand{!}}$ and $f_{*}$  functors
for sheaves (\var{\inserted{the}}{\hand{the}}{insertion}{ps_134-5_023-v0010} left and right adjoints of the inverse image functor $f^{*}$  for
sheaves of sets).  One would expect\hand{,} at  best\hand{,} $Lf_{\hand{!}}$  to exist when $f_{\hand{!}}$ itself
does,  namely \var{\inserted{when}}{\hand{when}}{insertion}{ps_134-5_023-v0011} $f^{*}$  commutes not only with finite inverse limits,  but with
infinite products as well.  As for $Rf_{*}$ ,  whereas there is no problem for
the existence of $f_{*}$ itself,  \var{\inserted[typed]{already}}{already}{typed insertion}{ps_134-5_023-v0012} \var{\deleted{\xout{still even}}}{}{deletion}{ps_134-5_023-v0013} in the case of a morphism of topoi
coming from a map in (Cat),  a map say between finite ordered sets,  the
existence of $Rf_{*}$  has still to be established.  Thus presumably I'll con-%
tent myself \var{\over{wi}{wi}}{\hand{wi}}{writing over}{ps_134-5_023-v0014}th writing down and comparing a few tentative definitions of
$D(T)$ and make some reasonable guesses as to it's \var{\xout{functorial}}{}{typewriter deletion}{ps_134-5_023-v0015} variances.
Intuitively,  the objects of $D(T)$  may be viewed as \hand{``}sheaves of homotopy
types over $T$\hand{''} , or \hand{``}relative homotopy types over $T$\hand{''},  or \hand{``}non commutative
chain complexes over $T$ up to quasi-isomorphism\hand{''}.  As in the case when $T$ is

\margin{23}

%% END SOURCE ps_134-5_023
\ifnum\EditionMode=1
\psapparatusstart
\psentry{ps_134-5_023-v0001}{1}{typewriter deletion}{\xout{,}}{}{The typewriter strikes out this material; it leaves the corrected running text.}
\psentry{ps_134-5_023-v0002}{2}{typewriter deletion}{\xout{strict}}{}{The typewriter strikes out this material; it leaves the corrected running text.}
\psentry{ps_134-5_023-v0003}{3}{typed insertion}{\inserted[typed]{``enough''}}{``enough''}{The typed insertion is incorporated at its indicated position; insertion brackets leave the running text.}
\psentry{ps_134-5_023-v0004}{4}{writing over}{\over{A}{A}}{\hand{A}}{The superposed handwritten reading replaces the underlying reading.}
\psentry{ps_134-5_023-v0005}{5}{transposition}{the \deleted{(Hot)} construction \inserted{of (Hot)} as a derived category over an object $I$
of (Cat)\hand{,} \hand{$\swarrow$}%
% The handwritten circle and arrow move the construction-of-Hot phrase to this position.
 }{over an object $I$
of (Cat) the  construction \hand{of (Hot)} as a derived category\hand{,} %
% The handwritten circle and arrow move the construction-of-Hot phrase to this position.
 }{The handwritten circle and arrow move the construction-of-Hot phrase after of (Cat), retaining its internal deletion and insertion.}
\psentry{ps_134-5_023-v0006}{6}{writing over}{\over{to}{to}}{\hand{to}}{The superposed handwritten reading replaces the underlying reading.}
\psentry{ps_134-5_023-v0007}{7}{typed insertion}{\inserted[typed]{some}}{some}{The typed insertion is incorporated at its indicated position; insertion brackets leave the running text.}
\psentry{ps_134-5_023-v0008}{8}{insertion}{\inserted{written and}}{\hand{written and}}{The handwritten insertion is incorporated at its indicated position; insertion brackets leave the running text.}
\psentry{ps_134-5_023-v0009}{9}{typewriter deletion}{\xout{the}}{}{The typewriter strikes out this material; it leaves the corrected running text.}
\psentry{ps_134-5_023-v0010}{10}{insertion}{\inserted{the}}{\hand{the}}{The handwritten insertion is incorporated at its indicated position; insertion brackets leave the running text.}
\psentry{ps_134-5_023-v0011}{11}{insertion}{\inserted{when}}{\hand{when}}{The handwritten insertion is incorporated at its indicated position; insertion brackets leave the running text.}
\psentry{ps_134-5_023-v0012}{12}{typed insertion}{\inserted[typed]{already}}{already}{The typed insertion is incorporated at its indicated position; insertion brackets leave the running text.}
\psentry{ps_134-5_023-v0013}{13}{deletion}{\deleted{\xout{still even}}}{}{The author cancels this material by hand; it leaves the corrected running text.}
\psentry{ps_134-5_023-v0014}{14}{writing over}{\over{wi}{wi}}{\hand{wi}}{The superposed handwritten reading replaces the underlying reading.}
\psentry{ps_134-5_023-v0015}{15}{typewriter deletion}{\xout{functorial}}{}{The typewriter strikes out this material; it leaves the corrected running text.}
\psapparatusend
\fi
\pssource{ps_134-5_024}{134-5/24}{https://grothendieck.umontpellier.fr/web/viewer.html?file=../134-5.pdf\#page=25}
\psnumber{ps_134-5_024-v0001}{1}
\psnumber{ps_134-5_024-v0002}{2}
\psnumber{ps_134-5_024-v0003}{3}
\psnumber{ps_134-5_024-v0004}{4}
\psnumber{ps_134-5_024-v0005}{5}
\psnumber{ps_134-5_024-v0006}{6}
\psnumber{ps_134-5_024-v0007}{7}
\psnumber{ps_134-5_024-v0008}{8}
\psnumber{ps_134-5_024-v0009}{9}
\psnumber{ps_134-5_024-v0010}{10}
\psnumber{ps_134-5_024-v0011}{11}
\psnumber{ps_134-5_024-v0012}{12}
\psnumber{ps_134-5_024-v0013}{13}
\psnumber{ps_134-5_024-v0014}{14}
\psnumber{ps_134-5_024-v0015}{15}
\psnumber{ps_134-5_024-v0016}{16}
\psnumber{ps_134-5_024-v0017}{17}
\psnumber{ps_134-5_024-v0018}{18}
\psnumber{ps_134-5_024-v0019}{19}
\psnumber{ps_134-5_024-v0020}{20}
%% BEGIN SOURCE ps_134-5_024
% Pursuing Stacks (cotes 134-2 to 134-8): archive cote 134-5, site page 24 (PDF page 25), typescript page 211.
% typescript. Draft compared: ps4.tex:1090-1146.
% AI-transcribed by Claude Opus 5.5 (claude-opus-5-5, adaptive thinking, effort high), Message Batches API, batch ps2 (msgbatch_01NVwr1Yuo8qJf5rykeGBXPK), 2026-11-12; reviewed against both supplied scan halves; see CHECK_DECISIONS.tsv.
\scanpage{134-5/24}
\mspage{211}
the final (one point) topos, \var{\over{\unsure{where}}{whence}}{\hand{whence}}{writing over}{ps_134-5_024-v0001} D(T) is just the homotopy category
(Hot), there \var{\xout{is}}{}{typewriter deletion}{ps_134-5_024-v0002} \var{\inserted[typed]{must be}}{must be}{typed insertion}{ps_134-5_024-v0003} of course a vast variety of ways of defining D(T) in term\unsure{s}
of model categories, and I would like to review some which seem significan\unsure{t.}
% NOTE: the right edge of the scan cuts off the ends of several lines
In fact, one cannot help but looking at two mutually dual groups of model
categories, giving rise to \var{\inserted[typed]{(at least)}}{(at least)}{typed insertion}{ps_134-5_024-v0004} two definitely non equivalent derived categori\hand{\unsure{es}}
D(T) and D'(T) say, which,in case T=I\^{}, would correspond to D(I) and
D(I\var{\deleted{\illegible{1}}}{}{deletion}{ps_134-5_024-v0005}\var{\inserted[typed]{${}^{o}$}}{${}^{o}$}{typed insertion}{ps_134-5_024-v0006}). A typical model category for the former \var{\xout{would be}}{}{typewriter deletion}{ps_134-5_024-v0007} is made up with
Illusie's semisimplicial sheaves; a typical model category for \var{\deleted{\xout{the latter}}}{}{deletion}{ps_134-5_024-v0008}
D'(T), on the other hand, should be made up with 1-stacks on T (in Giraud\hand{'}
s sense), \var{\over{f}{f}}{\hand{f}}{writing over}{ps_134-5_024-v0009}or a suitable notion of weak equivalence between these.

Maybe the most natural models of all, in this context, for ``relativ\unsure{e}
homotopy types over the topos T\hand{''}, should be \hand{\uline{\textcolor{black}{topoI}}} \var{\inserted[typed]{X}}{X}{typed insertion}{ps_134-5_024-v0010} over T (generalizing \var{\inserted{the}}{\hand{the}}{insertion}{ps_134-5_024-v0011}
categories over an ``indexing category'' I). The only trouble with this
point of view though is that the best we can hope for, in terms of Illu-%
sie's D(T) say, is that a topos X over T gives rise to a proobject of
D(T), which needs not come from an object of D(T) itself i.e. is not neces-%
sarily \var{\xout{locally}}{}{typewriter deletion}{ps_134-5_024-v0012} ``essentially constant.'' This brings us back to the simpler
\var{\inserted[typed]{and still more basic}}{and still more basic}{typed insertion}{ps_134-5_024-v0013} question of associating a pro-homotopy type, namely a pro-object of (Hot),
to any topos X - namely \var{\inserted[typed]{back}}{back}{typed insertion}{ps_134-5_024-v0014} to the \hand{\v{\textcolor{black}{C}}}ech-Verdier-Artin-Mazur construction.
This has been handled so far using the semisimplicial models for (Hot),
I suspect though that using (Cat) as a modelizer will give a more elegant
treatment, as already suggested earlier. Whichever way we choose to get
the basic functor
\[ \text{Topoi} \hand{\longrightarrow} \text{Pro(Hot)} \quad , \]
this functor will allow us, given a basic localizer $\underline{W}$ , to define
$\underline{W}$-equivalences between topoi\hand{\unsure{,}} as maps \var{\xout{such that}}{}{typewriter deletion}{ps_134-5_024-v0015} which become isomorphisms
under the composition of the basic functor above, and the canonical
functor
\[ \text{Pro(Hot)} \hand{\longrightarrow} \text{Pro(Hot}\var{\over{\illegible{1}}{(}}{\hand{(}}{writing over}{ps_134-5_024-v0016}\underline{W})) \]
deduced from the localization functor \var{\xout{Hot}}{}{typewriter deletion}{ps_134-5_024-v0017}
\[ \text{(Hot)} \hand{\longrightarrow} \text{(Hot(}\underline{W})) \ . \]
(It turns out, using Quillen's theory, that usual weak equivalence is
indeed the finest of all possible basic localizers $\underline{W}$ , hence Hot($\underline{W}$) is
indeed a l\var{\over{o}{o}}{\hand{o}}{writing over}{ps_134-5_024-v0018}\var{\over{\illegible{1}}{ca}}{\hand{ca}}{writing over}{ps_134-5_024-v0019}lization of (Hot).)

Maybe it is not too unreasonable to expect that all, or most homotopy
constructions, involving a topos T, can be expressed \var{\xout{by}}{}{typewriter deletion}{ps_134-5_024-v0020} replacing T by
it's image in Pro(Hot) , or in Pro(Hot($\underline{W}$)) if the constructions are rela-%
tive to a given basic localizer $\underline{W}$ . The very first example one would like

\margin{24}

%% END SOURCE ps_134-5_024
\ifnum\EditionMode=1
\psapparatusstart
\psentry{ps_134-5_024-v0001}{1}{writing over}{\over{\unsure{where}}{whence}}{\hand{whence}}{The superposed handwritten reading replaces the underlying reading.}
\psentry{ps_134-5_024-v0002}{2}{typewriter deletion}{\xout{is}}{}{The typewriter strikes out this material; it leaves the corrected running text.}
\psentry{ps_134-5_024-v0003}{3}{typed insertion}{\inserted[typed]{must be}}{must be}{The typed insertion is incorporated at its indicated position; insertion brackets leave the running text.}
\psentry{ps_134-5_024-v0004}{4}{typed insertion}{\inserted[typed]{(at least)}}{(at least)}{The typed insertion is incorporated at its indicated position; insertion brackets leave the running text.}
\psentry{ps_134-5_024-v0005}{5}{deletion}{\deleted{\illegible{1}}}{}{The author cancels this material by hand; it leaves the corrected running text.}
\psentry{ps_134-5_024-v0006}{6}{typed insertion}{\inserted[typed]{${}^{o}$}}{${}^{o}$}{The typed insertion is incorporated at its indicated position; insertion brackets leave the running text.}
\psentry{ps_134-5_024-v0007}{7}{typewriter deletion}{\xout{would be}}{}{The typewriter strikes out this material; it leaves the corrected running text.}
\psentry{ps_134-5_024-v0008}{8}{deletion}{\deleted{\xout{the latter}}}{}{The author cancels this material by hand; it leaves the corrected running text.}
\psentry{ps_134-5_024-v0009}{9}{writing over}{\over{f}{f}}{\hand{f}}{The superposed handwritten reading replaces the underlying reading.}
\psentry{ps_134-5_024-v0010}{10}{typed insertion}{\inserted[typed]{X}}{X}{The typed insertion is incorporated at its indicated position; insertion brackets leave the running text.}
\psentry{ps_134-5_024-v0011}{11}{insertion}{\inserted{the}}{\hand{the}}{The handwritten insertion is incorporated at its indicated position; insertion brackets leave the running text.}
\psentry{ps_134-5_024-v0012}{12}{typewriter deletion}{\xout{locally}}{}{The typewriter strikes out this material; it leaves the corrected running text.}
\psentry{ps_134-5_024-v0013}{13}{typed insertion}{\inserted[typed]{and still more basic}}{and still more basic}{The typed insertion is incorporated at its indicated position; insertion brackets leave the running text.}
\psentry{ps_134-5_024-v0014}{14}{typed insertion}{\inserted[typed]{back}}{back}{The typed insertion is incorporated at its indicated position; insertion brackets leave the running text.}
\psentry{ps_134-5_024-v0015}{15}{typewriter deletion}{\xout{such that}}{}{The typewriter strikes out this material; it leaves the corrected running text.}
\psentry{ps_134-5_024-v0016}{16}{writing over}{$\over{\illegible{1}}{(}$}{$\hand{(}$}{The superposed handwritten reading replaces the underlying reading.}
\psentry{ps_134-5_024-v0017}{17}{typewriter deletion}{\xout{Hot}}{}{The typewriter strikes out this material; it leaves the corrected running text.}
\psentry{ps_134-5_024-v0018}{18}{writing over}{\over{o}{o}}{\hand{o}}{The superposed handwritten reading replaces the underlying reading.}
\psentry{ps_134-5_024-v0019}{19}{writing over}{\over{\illegible{1}}{ca}}{\hand{ca}}{The superposed handwritten reading replaces the underlying reading.}
\psentry{ps_134-5_024-v0020}{20}{typewriter deletion}{\xout{by}}{}{The typewriter strikes out this material; it leaves the corrected running text.}
\psapparatusend
\fi
\pssource{ps_134-5_025}{134-5/25}{https://grothendieck.umontpellier.fr/web/viewer.html?file=../134-5.pdf\#page=26}
\psnumber{ps_134-5_025-v0001}{1}
\psnumber{ps_134-5_025-v0002}{2}
\psnumber{ps_134-5_025-v0003}{3}
\psnumber{ps_134-5_025-v0004}{4}
\psnumber{ps_134-5_025-v0005}{5}
\psnumber{ps_134-5_025-v0006}{6}
\psnumber{ps_134-5_025-v0007}{7}
\psnumber{ps_134-5_025-v0008}{8}
\psnumber{ps_134-5_025-v0009}{9}
\psnumber{ps_134-5_025-v0010}{10}
\psnumber{ps_134-5_025-v0011}{11}
\psnumber{ps_134-5_025-v0012}{12}
\psnumber{ps_134-5_025-v0013}{13}
\psnumber{ps_134-5_025-v0014}{14}
%% BEGIN SOURCE ps_134-5_025
% Pursuing Stacks (cotes 134-2 to 134-8): archive cote 134-5, site page 25 (PDF page 26), typescript page 212.
% typescript. Draft compared: ps4.tex:1132-1189.
% AI-transcribed by Claude Opus 5.5 (claude-opus-5-5, adaptive thinking, effort high), Message Batches API, batch ps2 (msgbatch_01NVwr1Yuo8qJf5rykeGBXPK), 2026-11-12; reviewed against both supplied scan halves; see CHECK_DECISIONS.tsv.
\scanpage{134-5/25}
\mspage{212}
to look up in this respect,  is surely D(T),  defined say à la Illusie.

In the present context,  the main point of the property of \hand{\uline{\textcolor{black}{local asphe-}}}%
\hand{\uline{\textcolor{black}{ricity}}} for a topos (cf.  section 35)  is that the corresponding prohomotopy
type is essentially constant,  i.e. the topos defines an actual homotopy
type.  Thus locally aspheric topoi  and weak equivalences between these
should \var{\over{b}{b}}{\hand{b}}{writing over}{ps_134-5_025-v0001}e eligible models for homotopy types  (more accurately,  make up
a modelizer)\hand{,} just as the basic modelizer (Cat) contained in it.  The cor-%
responding statements,  when introducing a basic loca\var{\over{l}{l}}{\hand{l}}{writing over}{ps_134-5_025-v0002}izer  $\underline{W}$ ,  should be
equally valid.  One might expect too a relative variant for the notion of
local asphericity or $\underline{W}$-asphericity,  in case of a topos X over a given
base topos T (whom we may have to suppose already locally aspheric),  with
the \var{\deleted{property}}{}{deletion}{ps_134-5_025-v0003}\var{\inserted{implication}}{\hand{implication}}{insertion}{ps_134-5_025-v0004} that the corresponding object of Pro(D(T)) should be again
essentially constant.

A last question I would like to mention here is about the meaning of
the notion of a so-called ``modelizing topos'',  introduced in a somewhat
formal way in section 35,  as a \var{\over{l}{l}}{\hand{l}}{writing over}{ps_134-5_025-v0005}oc ally aspheric \var{\inserted[typed]{and aspheric}}{and aspheric}{typed insertion}{ps_134-5_025-v0006} topos T such that the
Lawvere element $L_T$ of T is aspheric over the final object. (We assumed at
first,  moreover,  that T be even totally\var{\xout{x}}{}{typewriter deletion}{ps_134-5_025-v0007} aspheric,  but soon after the
point of view \var{\inserted[typed]{and terminology}}{and terminology}{typed insertion}{ps_134-5_025-v0008} shifted a little and the totally aspheric case was referred \var{\xout{t}}{}{typewriter deletion}{ps_134-5_025-v0009}
to as a \hand{\uline{\textcolor{black}{strictly}}} modelizing topos\hand{,} \var{\inserted{cf. page 68.}}{\hand{cf. page 68.}}{insertion}{ps_134-5_025-v0010})  As made clear there,  the expectation
suggested by the terminology is of course that such a topos should indeed
be a modelizer,  when endowed with the usual notion of weak equivalence.
This still makes sense and seems plausible enough,  when usual weak equi-%
valence is replaced by the weaker notion defined in terms of an arbitrary
basic localizer $\underline{W}$ .   A related question is to get a feeling for how res-%
trictive the conditions put on a modelizing topos \var{\deleted{\xout{\illegible{1}}}}{}{deletion}{ps_134-5_025-v0011} are.  It is clear
that nearly all topoi met with in practice,  including those associated to
the more common topological spaces (such as locally contractible ones) are
locally aspheric  - \var{\inserted{but}}{\hand{but}}{insertion}{ps_134-5_025-v0012} what about the condition on the Lawvere element ?
For instance, taking a topological space admitting a \var{\inserted[typed]{finite}}{finite}{typed insertion}{ps_134-5_025-v0013} triangulation and
which is aspheric i.e. contractible,  is the corresponding topos \var{\xout{modeli}}{}{typewriter deletion}{ps_134-5_025-v0014}
modelizing ?

\margin{25}

%% END SOURCE ps_134-5_025
\ifnum\EditionMode=1
\psapparatusstart
\psentry{ps_134-5_025-v0001}{1}{writing over}{\over{b}{b}}{\hand{b}}{The superposed handwritten reading replaces the underlying reading.}
\psentry{ps_134-5_025-v0002}{2}{writing over}{\over{l}{l}}{\hand{l}}{The superposed handwritten reading replaces the underlying reading.}
\psentry{ps_134-5_025-v0003}{3}{deletion}{\deleted{property}}{}{The author cancels this material by hand; it leaves the corrected running text.}
\psentry{ps_134-5_025-v0004}{4}{insertion}{\inserted{implication}}{\hand{implication}}{The handwritten insertion is incorporated at its indicated position; insertion brackets leave the running text.}
\psentry{ps_134-5_025-v0005}{5}{writing over}{\over{l}{l}}{\hand{l}}{The superposed handwritten reading replaces the underlying reading.}
\psentry{ps_134-5_025-v0006}{6}{typed insertion}{\inserted[typed]{and aspheric}}{and aspheric}{The typed insertion is incorporated at its indicated position; insertion brackets leave the running text.}
\psentry{ps_134-5_025-v0007}{7}{typewriter deletion}{\xout{x}}{}{The typewriter strikes out this material; it leaves the corrected running text.}
\psentry{ps_134-5_025-v0008}{8}{typed insertion}{\inserted[typed]{and terminology}}{and terminology}{The typed insertion is incorporated at its indicated position; insertion brackets leave the running text.}
\psentry{ps_134-5_025-v0009}{9}{typewriter deletion}{\xout{t}}{}{The typewriter strikes out this material; it leaves the corrected running text.}
\psentry{ps_134-5_025-v0010}{10}{insertion}{\inserted{cf. page 68.}}{\hand{cf. page 68.}}{The handwritten insertion is incorporated at its indicated position; insertion brackets leave the running text.}
\psentry{ps_134-5_025-v0011}{11}{deletion}{\deleted{\xout{\illegible{1}}}}{}{The author cancels this material by hand; it leaves the corrected running text.}
\psentry{ps_134-5_025-v0012}{12}{insertion}{\inserted{but}}{\hand{but}}{The handwritten insertion is incorporated at its indicated position; insertion brackets leave the running text.}
\psentry{ps_134-5_025-v0013}{13}{typed insertion}{\inserted[typed]{finite}}{finite}{The typed insertion is incorporated at its indicated position; insertion brackets leave the running text.}
\psentry{ps_134-5_025-v0014}{14}{typewriter deletion}{\xout{modeli}}{}{The typewriter strikes out this material; it leaves the corrected running text.}
\psapparatusend
\fi
\pssource{ps_134-5_026}{134-5/26}{https://grothendieck.umontpellier.fr/web/viewer.html?file=../134-5.pdf\#page=27}
\psnumber{ps_134-5_026-v0001}{1}
\psnumber{ps_134-5_026-v0002}{2}
\psnumber{ps_134-5_026-v0003}{3}
\psnumber{ps_134-5_026-v0004}{4}
\psnumber{ps_134-5_026-v0005}{5}
\psnumber{ps_134-5_026-v0006}{6}
\psnumber{ps_134-5_026-v0007}{7}
\psnumber{ps_134-5_026-v0008}{8}
\psnumber{ps_134-5_026-v0009}{9}
\psnumber{ps_134-5_026-v0010}{10}
\psnumber{ps_134-5_026-v0011}{11}
\psnumber{ps_134-5_026-v0012}{12}
\psnumber{ps_134-5_026-v0013}{13}
%% BEGIN SOURCE ps_134-5_026
% Pursuing Stacks (cotes 134-2 to 134-8): archive cote 134-5, site page 26 (PDF page 27), typescript page 213.
% typescript. Draft compared: ps4.tex:1171-1239.
% AI-transcribed by Claude Opus 5.5 (claude-opus-5-5, adaptive thinking, effort high), Message Batches API, batch ps2 (msgbatch_01NVwr1Yuo8qJf5rykeGBXPK), 2026-11-12; reviewed against both supplied scan halves; see CHECK_DECISIONS.tsv.
\scanpage{134-5/26}
\mspage{213}

10.6.

72. In the long last, \var{\over[typed]{w}{W}}{W}{typewriter overwriting}{ps_134-5_026-v0001}e'll come back now to the ``asphericity game''\hand{!}
\var{\deleted{again !}}{}{deletion}{ps_134-5_026-v0002} Let's take up the exposition at the point where we stopped
two months ago (section 67, p.188 ). We were then about to reformulate
in various ways the property of asphericity, more specifically
\underline{W}-asphericity, for a given map
\[ i : A \hand{\longrightarrow} B \]
in (Cat). To this end, we introduced the corresponding diagram of maps
in (Cat) with ``commutation morphism'' $\hand{\lambda}_i$ :
% Hand-drawn arrows are marked blue; typed objects remain plain.
\diagram{\begin{tikzcd}[column sep=small]
B^{\wedge} \arrow[color=blue, rr, "i^{*}"] \arrow[color=blue, dr, "i_B"', ""{name=iB, right}] \& \& A^{\wedge} \arrow[color=blue, dl, "i_A", ""{name=iA, left}] \\
\& (\mathrm{Cat}) \arrow[color=blue, bend right=10, from=iA, to=iB, "\text{\hand{$\lambda_i$}}"'] \&
\end{tikzcd}}
\hfill .

As already stated, ,the asphericity condition on $i$ can be expressed \var{\xout{\illegible{1}}}{}{typewriter deletion}{ps_134-5_026-v0003}
in a variety of equivalent ways, as a condition on either of the three
``aspects''
\[ \hand{\lambda}_i \ , \ i^{*} \ , \ i_A i^{*} \]
of the situation created by $i$, with respect to the localizing sets
$\underline{W}$ , $\underline{W}_A$ , $\underline{W}_B$ in the three categories under consideration, or with res-%
pect to the notion of \underline{W}-aspheric objects in these. As $\underline{W}_A$\hand{,} and aspheric
objects in $A^{\wedge}$\hand{,}are defined \var{\inserted{respectively}}{\hand{respectively}}{insertion}{ps_134-5_026-v0004} in terms of \underline{W} and aspheric objects in (Cat)
via the functor \hand{\uline{\textcolor{black}{$i_A$ ,}}} it turns out that the \var{\xout{characterizations}}{}{typewriter deletion}{ps_134-5_026-v0005}\var{\inserted[typed]{formulations}}{formulations}{typed insertion}{ps_134-5_026-v0006} in terms of
properties of $i_A i^{*}$ reduce trivially to the corresponding formulations
in terms of $i^{*}$ , which will be the most directly useful for\var{\xout{m}}{}{typewriter deletion}{ps_134-5_026-v0007} our \var{\xout{pur}}{}{typewriter deletion}{ps_134-5_026-v0008}
\var{\xout{later}}{}{typewriter deletion}{ps_134-5_026-v0009} purpose for later work - hence we'll omit them, and focus atten-%
tion instead on $\hand{\lambda}_i$ and $i^{*}$ . Notations are those of loc. cit. , in par-%
ticular we are working with the categories
\[ \mathrm{Hot}^{\underline{W}}_{A} \quad \text{or} \quad \mathrm{Hot}_A \overset{\mathrm{def}}{=} \underline{W}_A^{-1} A^{\wedge} \]
and
\[ \mathrm{Hot}(\underline{W}) = \underline{W}^{-1} (\mathrm{Cat}) \quad , \]
where \underline{W} is a given ``basic localizer'' , namely a set of arrows in (Cat)
satisfying certain conditions.

It seems worthwhile here to be careful to state which exactly are
the properties of \underline{W} we are going to use in the ``asphericity game'' - there%
fore I'm going to list them again here, \var{\xout{\unsure{Once again I under}}}{}{typewriter deletion}{ps_134-5_026-v0010} using a \var{\deleted{desig-}}{}{deletion}{ps_134-5_026-v0011}
\var{\deleted{nation}}{}{deletion}{ps_134-5_026-v0012}\var{\inserted{labelling}}{\hand{labelling}}{insertion}{ps_134-5_026-v0013} which\hand{,} hopefully\hand{,} will not have to be changed again:

Loc 1 ) ``Mild saturation'' (cf. page \hand{59}) ,

Loc 2 ) ``Homotopy condition'' :$\hand{\Delta}_1 x X \hand{\rightarrow} X$ is in \underline{W} for any $X$ in (Cat) ,

\margin{26}

%% END SOURCE ps_134-5_026
\ifnum\EditionMode=1
\psapparatusstart
\psentry{ps_134-5_026-v0001}{1}{typewriter overwriting}{\over[typed]{w}{W}}{W}{The superposed typed reading replaces the underlying reading.}
\psentry{ps_134-5_026-v0002}{2}{deletion}{\deleted{again !}}{}{The author cancels this material by hand; it leaves the corrected running text.}
\psentry{ps_134-5_026-v0003}{3}{typewriter deletion}{\xout{\illegible{1}}}{}{The typewriter strikes out this material; it leaves the corrected running text.}
\psentry{ps_134-5_026-v0004}{4}{insertion}{\inserted{respectively}}{\hand{respectively}}{The handwritten insertion is incorporated at its indicated position; insertion brackets leave the running text.}
\psentry{ps_134-5_026-v0005}{5}{typewriter deletion}{\xout{characterizations}}{}{The typewriter strikes out this material; it leaves the corrected running text.}
\psentry{ps_134-5_026-v0006}{6}{typed insertion}{\inserted[typed]{formulations}}{formulations}{The typed insertion is incorporated at its indicated position; insertion brackets leave the running text.}
\psentry{ps_134-5_026-v0007}{7}{typewriter deletion}{\xout{m}}{}{The typewriter strikes out this material; it leaves the corrected running text.}
\psentry{ps_134-5_026-v0008}{8}{typewriter deletion}{\xout{pur}}{}{The typewriter strikes out this material; it leaves the corrected running text.}
\psentry{ps_134-5_026-v0009}{9}{typewriter deletion}{\xout{later}}{}{The typewriter strikes out this material; it leaves the corrected running text.}
\psentry{ps_134-5_026-v0010}{10}{typewriter deletion}{\xout{\unsure{Once again I under}}}{}{The typewriter strikes out this material; it leaves the corrected running text.}
\psentry{ps_134-5_026-v0011}{11}{deletion}{\deleted{desig-}}{}{The author cancels this material by hand; it leaves the corrected running text.}
\psentry{ps_134-5_026-v0012}{12}{deletion}{\deleted{nation}}{}{The author cancels this material by hand; it leaves the corrected running text.}
\psentry{ps_134-5_026-v0013}{13}{insertion}{\inserted{labelling}}{\hand{labelling}}{The handwritten insertion is incorporated at its indicated position; insertion brackets leave the running text.}
\psapparatusend
\fi
\pssource{ps_134-5_027}{134-5/27}{https://grothendieck.umontpellier.fr/web/viewer.html?file=../134-5.pdf\#page=28}
\psnumber{ps_134-5_027-v0001}{1}
\psnumber{ps_134-5_027-v0002}{2}
\psnumber{ps_134-5_027-v0003}{3}
\psnumber{ps_134-5_027-v0004}{4}
\psnumber{ps_134-5_027-v0005}{5}
\psnumber{ps_134-5_027-v0006}{6}
\psnumber{ps_134-5_027-v0007}{7}
\psnumber{ps_134-5_027-v0008}{8}
\psnumber{ps_134-5_027-v0009}{9}
\psnumber{ps_134-5_027-v0010}{10}
\psnumber{ps_134-5_027-v0011}{11}
\psnumber{ps_134-5_027-v0012}{12}
\psnumber{ps_134-5_027-v0013}{13}
\psnumber{ps_134-5_027-v0014}{14}
\psnumber{ps_134-5_027-v0015}{15}
\psnumber{ps_134-5_027-v0016}{16}
\psnumber{ps_134-5_027-v0017}{17}
\psnumber{ps_134-5_027-v0018}{18}
\psnumber{ps_134-5_027-v0019}{19}
\psnumber{ps_134-5_027-v0020}{20}
%% BEGIN SOURCE ps_134-5_027
% Pursuing Stacks (cotes 134-2 to 134-8): archive cote 134-5, site page 27 (PDF page 28), typescript page 214.
% typescript. Draft compared: ps4.tex:1221-1317.
% AI-transcribed by Claude Opus 5.5 (claude-opus-5-5, adaptive thinking, effort high), Message Batches API, batch ps2 (msgbatch_01NVwr1Yuo8qJf5rykeGBXPK), 2026-11-12; reviewed against both supplied scan halves; see CHECK_DECISIONS.tsv.
\scanpage{134-5/27}
\mspage{214}

Loc 3 ) ``Localization condition'' : If $X,Y$ are objects in (Cat) over
an object $\var{\over{A}{\mathcal{S}}}{\hand{\mathcal{S}}}{writing over}{ps_134-5_027-v0001}$, and $u\colon X\hand{\to}Y$ an $\var{\over{A}{\mathcal{S}}}{\hand{\mathcal{S}}}{writing over}{ps_134-5_027-v0002}$-morphism such that for any a in
$\var{\over{A}{\mathcal{S}}}{\hand{\mathcal{S}}}{writing over}{ps_134-5_027-v0003}$, the induced $u_{/a} : X_{/a}\hand{\to}Y_{/a}$ is in $\underline{W}$ , then so is u .

I should be noted that the ``mild saturation'' condition is slightly
weaker than the saturation condition introduced later (p. 101 , conditi-%
ons a')b')c')), namely in condition c') (if $f\colon X\hand{\to}Y$ and $g\colon Y\hand{\to}X$ are
\var{\over{\unsure{j}}{s}}{\hand{s}}{writing over}{ps_134-5_027-v0004}uch that $gf,fg \hand{\in} \underline{W}$ , then $f,g \hand{\in} \underline{W}$ ) we restrict to the case when gf is
the identity, namely \var{\deleted{\illegible{1}}}{}{deletion}{ps_134-5_027-v0005} f an inclusion and g a retraction upon the corres%
ponding subob\var{\over{b}{j}}{\hand{j}}{writing over}{ps_134-5_027-v0006}ect. On the other hand, in what follows we \var{\deleted{\illegible{1}}}{}{deletion}{ps_134-5_027-v0007}are going to
use Loc 3 ) only in case $Y=\var{\over{A}{\mathcal{S}}}{\hand{\mathcal{S}}}{writing over}{ps_134-5_027-v0008}$ and \var{\deleted{\illegible{1}}}{}{deletion}{ps_134-5_027-v0009} $Y\hand{\to}\var{\over{A}{\mathcal{S}}}{\hand{\mathcal{S}}}{writing over}{ps_134-5_027-v0010}$ is the identity. I stated
the condition in greater generality than needed for the time being,
in view of later convenience - as later it will have to be used in full
strength. It is not clear whether the weaker form \var{\inserted[typed]{of Loc 3}}{of Loc 3}{typed insertion}{ps_134-5_027-v0011} (plus of course Loc 1
and Loc 2 ) implies already the stronger. We'll see later a number of
nice further properties \var{\inserted[typed]{of $\underline{W}$}}{of $\underline{W}$}{typed insertion}{ps_134-5_027-v0012} implied by these we are going to work with for
the time being - including \hand{\uline{\textcolor{black}{strong}}} saturation of $\underline{W}$ , namely that $\underline{W}$ is
the set of arrows made invertible by the localization functor
\[ (\mathrm{Cat}) \hand{\longrightarrow} \underline{W}^{-1}(\mathrm{Cat}) = \mathrm{Hot}(\underline{W}) \quad . \]

\hand{\uline{\textcolor{black}{Proposition}}}\hand{1}.Let as above $i\colon A\hand{\to}B$ be a map in (Cat). Consider the
following conditions on i:

(i) For any F in $B^{\wedge}$, $\hand{\lambda}_i(F)\colon A_{/F} \hand{\to} B_{/F}$ is in $\underline{W}$ .

(i') For any F as above, $\hand{\lambda}_i(F)$ is $\underline{W}$-aspheric (i.e. satisfies the
assumption on u in Loc 3) above, when $Y\hand{\to}\var{\over{A}{\mathcal{S}}}{\hand{\mathcal{S}}}{writing over}{ps_134-5_027-v0013}$ is an identity).

(i") Same \var{\deleted{as}}{}{deletion}{ps_134-5_027-v0014} (i), but restricting to F = b in B .

(ii) For any F in $B^{\wedge}$, F $\underline{W}$-aspheric $\hand{\Longrightarrow}$ $i^{*}(F)$ $\underline{W}$-aspheric .
(ii') For any F in $B^{\wedge}$, F $\underline{W}$-aspheric $\hand{\Longleftrightarrow}$ $i^{*}(F)$ $\underline{W}$-aspheric .

(iii) For any b in B, $A_{/b}$ ( $\overset{\text{def}}{=}$ $A_{/i^*(b)}$ ) is $\underline{W}$-aspheric, i.e. i \var{\inserted[typed]{is $\underline{W}$-aspheric.}}{is $\underline{W}$-aspheric.}{typed insertion}{ps_134-5_027-v0015}

(iv) For any map f in $B^{\wedge}$ , $f \hand{\in} \underline{W}_B \hand{\Longrightarrow} i^{*}(f) \hand{\in} \underline{W}_A$ .

(iv') For any f as above, $f \hand{\in} \underline{W}_B \hand{\Longleftrightarrow} i^{*}(f) \hand{\in} \underline{W}_A$ .

(v) Condition (iv) holds i.e. $i^{*}$ induces a functor
\[ \hand{\overline{\textcolor{black}{i^{*}}}} : \mathrm{Ho}\var{\over{t}{t}}{\hand{t}}{writing over}{ps_134-5_027-v0016}_B \hand{\longrightarrow} \mathrm{Hot}_A \hand{,} \]
and \var{\inserted{moreover}}{\hand{moreover}}{insertion}{ps_134-5_027-v0017} the latter is an \hand{\uline{\textcolor{black}{equivalence}}}.

The conditions (i) up to (iii) are all equivalent, \var{\inserted[typed]{call this \xout{condi}}}{call this }{typed insertion}{ps_134-5_027-v0018} \var{\xout{\illegible{3}}}{}{typewriter deletion}{ps_134-5_027-v0019}
\var{\xout{\illegible{1}}}{}{typewriter deletion}{ps_134-5_027-v0020} set of conditions \hand{(}As\hand{)} ( $\underline{W}$-asphericity). We moreover have the
following implications between \hand{(}As\hand{)} and the remaining conditions (iv) to
(v) :

\margin{27}

%% END SOURCE ps_134-5_027
\ifnum\EditionMode=1
\psapparatusstart
\psentry{ps_134-5_027-v0001}{1}{writing over}{$\over{A}{\mathcal{S}}$}{$\hand{\mathcal{S}}$}{The superposed handwritten reading replaces the underlying reading.}
\psentry{ps_134-5_027-v0002}{2}{writing over}{$\over{A}{\mathcal{S}}$}{$\hand{\mathcal{S}}$}{The superposed handwritten reading replaces the underlying reading.}
\psentry{ps_134-5_027-v0003}{3}{writing over}{$\over{A}{\mathcal{S}}$}{$\hand{\mathcal{S}}$}{The superposed handwritten reading replaces the underlying reading.}
\psentry{ps_134-5_027-v0004}{4}{writing over}{\over{\unsure{j}}{s}}{\hand{s}}{The superposed handwritten reading replaces the underlying reading.}
\psentry{ps_134-5_027-v0005}{5}{deletion}{\deleted{\illegible{1}}}{}{The author cancels this material by hand; it leaves the corrected running text.}
\psentry{ps_134-5_027-v0006}{6}{writing over}{\over{b}{j}}{\hand{j}}{The superposed handwritten reading replaces the underlying reading.}
\psentry{ps_134-5_027-v0007}{7}{deletion}{\deleted{\illegible{1}}}{}{The author cancels this material by hand; it leaves the corrected running text.}
\psentry{ps_134-5_027-v0008}{8}{writing over}{$\over{A}{\mathcal{S}}$}{$\hand{\mathcal{S}}$}{The superposed handwritten reading replaces the underlying reading.}
\psentry{ps_134-5_027-v0009}{9}{deletion}{\deleted{\illegible{1}}}{}{The author cancels this material by hand; it leaves the corrected running text.}
\psentry{ps_134-5_027-v0010}{10}{writing over}{$\over{A}{\mathcal{S}}$}{$\hand{\mathcal{S}}$}{The superposed handwritten reading replaces the underlying reading.}
\psentry{ps_134-5_027-v0011}{11}{typed insertion}{\inserted[typed]{of Loc 3}}{of Loc 3}{The typed insertion is incorporated at its indicated position; insertion brackets leave the running text.}
\psentry{ps_134-5_027-v0012}{12}{typed insertion}{\inserted[typed]{of $\underline{W}$}}{of $\underline{W}$}{The typed insertion is incorporated at its indicated position; insertion brackets leave the running text.}
\psentry{ps_134-5_027-v0013}{13}{writing over}{$\over{A}{\mathcal{S}}$}{$\hand{\mathcal{S}}$}{The superposed handwritten reading replaces the underlying reading.}
\psentry{ps_134-5_027-v0014}{14}{deletion}{\deleted{as}}{}{The author cancels this material by hand; it leaves the corrected running text.}
\psentry{ps_134-5_027-v0015}{15}{typed insertion}{\inserted[typed]{is $\underline{W}$-aspheric.}}{is $\underline{W}$-aspheric.}{The typed insertion is incorporated at its indicated position; insertion brackets leave the running text.}
\psentry{ps_134-5_027-v0016}{16}{writing over}{$\over{t}{t}$}{$\hand{t}$}{The superposed handwritten reading replaces the underlying reading.}
\psentry{ps_134-5_027-v0017}{17}{insertion}{\inserted{moreover}}{\hand{moreover}}{The handwritten insertion is incorporated at its indicated position; insertion brackets leave the running text.}
\psentry{ps_134-5_027-v0018}{18}{typed insertion}{\inserted[typed]{call this \xout{condi}}}{call this }{The typed insertion is incorporated at its indicated position; insertion brackets leave the running text.}
\psentry{ps_134-5_027-v0019}{19}{typewriter deletion}{\xout{\illegible{3}}}{}{The typewriter strikes out this material; it leaves the corrected running text.}
\psentry{ps_134-5_027-v0020}{20}{typewriter deletion}{\xout{\illegible{1}}}{}{The typewriter strikes out this material; it leaves the corrected running text.}
\psapparatusend
\fi
\pssource{ps_134-5_028}{134-5/28}{https://grothendieck.umontpellier.fr/web/viewer.html?file=../134-5.pdf\#page=29}
\psnumber{ps_134-5_028-v0001}{1}
\psnumber{ps_134-5_028-v0002}{2}
\psnumber{ps_134-5_028-v0003}{3}
\psnumber{ps_134-5_028-v0004}{4}
\psnumber{ps_134-5_028-v0005}{5}
\psnumber{ps_134-5_028-v0006}{6}
\psnumber{ps_134-5_028-v0007}{7}
\psnumber{ps_134-5_028-v0008}{8}
\psnumber{ps_134-5_028-v0009}{9}
\psnumber{ps_134-5_028-v0010}{10}
\psnumber{ps_134-5_028-v0011}{11}
%% BEGIN SOURCE ps_134-5_028
% Pursuing Stacks (cotes 134-2 to 134-8): archive cote 134-5, site page 28 (PDF page 29), typescript page 215.
% typescript. Draft compared: ps4.tex:1301-1375.
% AI-transcribed by Claude Opus 5.5 (claude-opus-5-5, adaptive thinking, effort high), Message Batches API, batch ps2 (msgbatch_01NVwr1Yuo8qJf5rykeGBXPK), 2026-11-12; reviewed against both supplied scan halves; see CHECK_DECISIONS.tsv.
\scanpage{134-5/28}
\mspage{215}

% Hand-drawn arrows are marked blue; typed objects remain plain.
\margin{(*)}
\diagram{
\begin{tikzcd}
\text{(As)} \ar[color=blue, rr, Rightarrow, dashed, "\hand{if A,B ps.test}"] \ar[color=blue, dr, Rightarrow] \& \& \text{(v)} \ar[color=blue, ddl, Rightarrow] \\
\& \text{(iv')} \ar[color=blue, d, Rightarrow] \& \\
\& \text{(iv)} \ar[color=blue, uul, Rightarrow, dashed, bend left, "\hand{if A,B \underline{W}-asph.}"'] \&
\end{tikzcd}
} ,

where the implication (iv)$\Rightarrow$(As)   is subject to  A,B  being $\underline{W}$-aspheric,
and  (As)$\Rightarrow$(v) to  A,B  being ``pseudo-test categories''\var{\inserted[typed]{(for $\underline{W}$)}}{(for $\underline{W}$)}{typed insertion}{ps_134-5_028-v0001}, namely the cano-%
nical functors
\[ \hand{\overline{\textcolor{black}{\mathrm{i}}}}_A : \var{\deleted{\unsure{4}}}{}{deletion}{ps_134-5_028-v0002}\ \mathrm{Hot}_A \hand{\longrightarrow} \mathrm{Hot}(\underline{W}) \quad , \quad \hand{\overline{\textcolor{black}{\mathrm{i}}}}_B \ : \ \mathrm{Hot}_B \hand{\longrightarrow} \mathrm{Hot}(\underline{W}) \]
being equivalences.

\hand{\uline{\textcolor{black}{Corollary}}}:  Assume that A and B are pseudo test categories,  and are
$\underline{W}$-aspheric.  Then  all conditions (i) to (v) of the proposition above are
equivalent.

\hand{\uline{\textcolor{black}{Remark\var{\xout{x}}{}{typewriter deletion}{ps_134-5_028-v0003}1}}}.If we admit strong saturation of $\underline{W}$  (which will be proved later)\hand{,}
it follows at once that \var{\xout{the assumption}}{}{typewriter deletion}{ps_134-5_028-v0004} a pseudo-test category is \var{\xout{already}}{}{typewriter deletion}{ps_134-5_028-v0005}
necessarily $\underline{W}$-aspheric  -  hence the conclusion of the corollary holds
assuming only A ahd B are pseudo test categories.  Of course it holds a
fortiori if A and B are weak test categories,  or even test categories.
\var{\inserted{Also, strong saturation implies that in (*) above, we have even the implication : (v) $\Longrightarrow$ (iv'), stronger than (v) $\Longrightarrow$ (iv)}}{\hand{Also, strong saturation implies that in (*) above, we have even the implication : (v) $\Longrightarrow$ (iv'), stronger than (v) $\Longrightarrow$ (iv)}}{insertion}{ps_134-5_028-v0006}

Proof of proposition.   It is purely formal  -  for the first part,
\var{\inserted[typed]{it follows from}}{it follows from}{typed insertion}{ps_134-5_028-v0007}\var{\xout{we get}}{}{typewriter deletion}{ps_134-5_028-v0008} the diagram of tautological implications
\diagram{
\begin{tikzcd}[column sep=small]
 \& \text{(i')} \ar[color=blue, d, Rightarrow] \& \\
 \& \text{(i)} \ar[color=blue, dl, Rightarrow] \ar[color=blue, dr, Rightarrow] \& \\
\begin{array}{c}\var{\xout{\unsure{ixixix}}}{}{typewriter deletion}{ps_134-5_028-v0009}\\ \var{\xout{(ii}}{}{typewriter deletion}{ps_134-5_028-v0010}\\ \text{(ii')}\end{array} \ar[color=blue, d, Rightarrow] \& \& \text{(i'')} \ar[color=blue, dl, Leftrightarrow] \\
\text{(ii)} \ar[color=blue, r, Rightarrow] \& \text{(iii)} \ar[color=blue, r, Rightarrow] \& \text{(i')}
\end{tikzcd}
} \hand{,}

where the implication  (i')$\Rightarrow$(i)  is contained in the assumption Loc 3
on $\underline{W}$ .  \var{\xout{Onx\allowbreak{}thex\allowbreak{}otherx\allowbreak{}handx\allowbreak{}wex\allowbreak{}immediatelyx\allowbreak{}getx\allowbreak{}x\allowbreak{}(iii)x\allowbreak{}x\allowbreak{}x\allowbreak{}x\allowbreak{}(ii)x\allowbreak{}x\allowbreak{}}}{}{typewriter deletion}{ps_134-5_028-v0011} The impli-%
cations of the diagram  (*)  are about as formal  -  there is no point I
guess wfiting it out here.

\hand{\uline{\textcolor{black}{Remark 2}}}   Using the canonical functors
\[ \hand{\phi}_A \ : \ A^{\wedge} \hand{\to} \mathrm{Hot}(\underline{W}) \quad , \quad \hand{\phi}_B \ : \ B^{\wedge} \hand{\longrightarrow} \mathrm{Hot}(\underline{W}) \quad , \]
there are two other amusing versions still of (i)\hand{,} which  look a lot wea-%
ker still,  and are equivalent however to (i) i.e. to $\underline{W}$-asphericity of i,
namely :

(vi)    For any F in $B^{\wedge}$ ,  $\hand{\phi}_B(F) = \hand{\gamma}(B_{/F})$  and  $\hand{\phi}_A(i^{*}(F)) = \hand{\gamma}(A_{/F})$
are isomorphic objects of  $\mathrm{Hot}(\underline{W})$ .

\margin{28}

%% END SOURCE ps_134-5_028
\ifnum\EditionMode=1
\psapparatusstart
\psentry{ps_134-5_028-v0001}{1}{typed insertion}{\inserted[typed]{(for $\underline{W}$)}}{(for $\underline{W}$)}{The typed insertion is incorporated at its indicated position; insertion brackets leave the running text.}
\psentry{ps_134-5_028-v0002}{2}{deletion}{$\deleted{\unsure{4}}$}{}{The author cancels this material by hand; it leaves the corrected running text.}
\psentry{ps_134-5_028-v0003}{3}{typewriter deletion}{\xout{x}}{}{The typewriter strikes out this material; it leaves the corrected running text.}
\psentry{ps_134-5_028-v0004}{4}{typewriter deletion}{\xout{the assumption}}{}{The typewriter strikes out this material; it leaves the corrected running text.}
\psentry{ps_134-5_028-v0005}{5}{typewriter deletion}{\xout{already}}{}{The typewriter strikes out this material; it leaves the corrected running text.}
\psentry{ps_134-5_028-v0006}{6}{insertion}{\inserted{Also, strong saturation implies that in (*) above, we have even the implication : (v) $\Longrightarrow$ (iv'), stronger than (v) $\Longrightarrow$ (iv)}}{\hand{Also, strong saturation implies that in (*) above, we have even the implication : (v) $\Longrightarrow$ (iv'), stronger than (v) $\Longrightarrow$ (iv)}}{The handwritten insertion is incorporated at its indicated position; insertion brackets leave the running text.}
\psentry{ps_134-5_028-v0007}{7}{typed insertion}{\inserted[typed]{it follows from}}{it follows from}{The typed insertion is incorporated at its indicated position; insertion brackets leave the running text.}
\psentry{ps_134-5_028-v0008}{8}{typewriter deletion}{\xout{we get}}{}{The typewriter strikes out this material; it leaves the corrected running text.}
\psentry{ps_134-5_028-v0009}{9}{typewriter deletion}{$\xout{\unsure{ixixix}}$}{}{The typewriter strikes out this material; it leaves the corrected running text.}
\psentry{ps_134-5_028-v0010}{10}{typewriter deletion}{$\xout{(ii}$}{}{The typewriter strikes out this material; it leaves the corrected running text.}
\psentry{ps_134-5_028-v0011}{11}{typewriter deletion}{\xout{Onx\allowbreak{}thex\allowbreak{}otherx\allowbreak{}handx\allowbreak{}wex\allowbreak{}immediatelyx\allowbreak{}getx\allowbreak{}x\allowbreak{}(iii)x\allowbreak{}x\allowbreak{}x\allowbreak{}x\allowbreak{}(ii)x\allowbreak{}x\allowbreak{}}}{}{The typewriter strikes out this material; it leaves the corrected running text.}
\psapparatusend
\fi
\pssource{ps_134-5_029}{134-5/29}{https://grothendieck.umontpellier.fr/web/viewer.html?file=../134-5.pdf\#page=30}
\psnumber{ps_134-5_029-v0001}{1}
\psnumber{ps_134-5_029-v0002}{2}
\psnumber{ps_134-5_029-v0003}{3}
\psnumber{ps_134-5_029-v0004}{4}
\psnumber{ps_134-5_029-v0005}{5}
\psnumber{ps_134-5_029-v0006}{6}
\psnumber{ps_134-5_029-v0007}{7}
\psnumber{ps_134-5_029-v0008}{8}
\psnumber{ps_134-5_029-v0009}{9}
%% BEGIN SOURCE ps_134-5_029
% Pursuing Stacks (cotes 134-2 to 134-8): archive cote 134-5, site page 29 (PDF page 30), typescript page 216.
% typescript. Draft compared: ps4.tex:1359-1429.
% AI-transcribed by Claude Opus 5.5 (claude-opus-5-5, adaptive thinking, effort high), Message Batches API, batch pilot (msgbatch_0189Nhf1xB615YWnztubtLfc), 2026-11-12; reviewed against both supplied scan halves; see CHECK_DECISIONS.tsv.
\scanpage{134-5/29}
\mspage{216}

(vi') Same as (vi), with F restricted to be an object b in B.

Indeed, we have of course (i)$\hand{\Rightarrow}$(vi)$\hand{\Rightarrow}$(vi'), but also (vi')$\hand{\Rightarrow}$(iii)
if we admit strong saturation of $\underline{W}$, which implies that an object X of
(Cat) is $\underline{W}$-aspheric iff it's image in Hot($\underline{W}$) is a final object.

\hand{\uline{\textcolor{black}{Remark}}} 3. If we don't assume A and B to be $\underline{W}$-aspheric, the implications
\[ \text{(As)} \hand{\Longrightarrow} \text{(iv')} \hand{\Longrightarrow} \text{(iv)} \]
are both strict. As an illustration of this point, take for B a discre-%
te category (defined in terms of a set of indices I = Ob B), thus a
\margin{$(A_b)_{b\in I}$} \var{\xout{functor of a}}{}{typewriter deletion}{ps_134-5_029-v0001} category A over B is essentially the same as a family $(A_{\hand{b}})_{\hand{b\in I}}$
of objects of (Cat) indexed by I. In terms of this\var{\xout{e}}{}{typewriter deletion}{ps_134-5_029-v0002} family, we see at
once that the three conditions above on i: A$\hand{\rightarrow}$B mean respectively
(a) that all categories $A_b$ are $\underline{W}$-aspheric, (b) that all categories
$A_b$ are non-empty, and (c) condition vacuous. This example brings to
mind that \hand{\uline{\textcolor{black}{condition}}} (iv) \hand{\uline{\textcolor{black}{is a Serre-fibration type condition}}} , we'll
come back upon this condition when studying \var{\xout{\illegible{1}}}{}{typewriter deletion}{ps_134-5_029-v0003} homotopy properties
of maps in (Cat)\hand{,} with special emphasis on base change questions. Like-%
wise, condition (iv') appears as a strengthening of such Serre-type
condition, to the effect that the restriction of A over any connected
component of B should be moreover non-empty. As an example, we may take
the projection BxC $\hand{\rightarrow}$ \var{\over{\unsure{C}}{B}}{\hand{B}}{writing over}{ps_134-5_029-v0004} , where C is any non-empty object in (Cat).

I would like now to dwell still a little on the case, of special
interest of course for the modelizing story, when A and B are test cate-%
gories or something of the kind. The formulation (i) of the asphericity
condition \var{\over{of}{for}}{\hand{for}}{writing over}{ps_134-5_029-v0005} the functor i can be expressed by stating that in the
the functor $\hand{\overline{\textcolor{black}{i^{*}}}}$ \var{\xout{on}}{}{typewriter deletion}{ps_134-5_029-v0006} \var{\inserted[typed]{between}}{between}{typed insertion}{ps_134-5_029-v0007} the localizations $\mathrm{Hot}_B$ and $\mathrm{Hot}_A$ exists, and gives
rise (via \hand{$\lambda$}$_i$) to a commutation morphism which is an \hand{\uline{\textcolor{black}{isomorphism}}}
% Hand-drawn arrows are marked blue; typed objects remain plain.
\diagram{
\begin{tikzcd}[column sep=small]
\mathrm{Hot}_B \arrow[color=blue, rr, "\hand{\overline{\textcolor{black}{i}^{*}}}"] \arrow[color=blue, dr, swap, "\hand{\overline{i_B}}", ""{name=iB, right}] \& \& \mathrm{Hot}_A \arrow[color=blue, dl, "\hand{\overline{i_A}}", ""{name=iA, left}] \\
\& \mathrm{Hot}(\underline{W}) \arrow[color=blue, bend right=10, from=iA, to=iB, "\hand{\overline{\lambda_i}}"', "\hand{\sim}"] \&
\end{tikzcd}
}
\margin{$\overset{\overline{\lambda_i}}{\underset{\sim}{\longleftarrow}}$}
.

This shows, when A and B are pseudo-test categories for $\underline{W}$ , i.e. the
functors $\hand{\overline{\textcolor{black}{i_A}}}$ and $\hand{\overline{\textcolor{black}{i_B}}}$ are equivalences, that the functor $\hand{\overline{\textcolor{black}{i^{*}}}}$ deduced
from the \var{\inserted{$\underline{W}$-}}{\hand{$\underline{W}$-}}{insertion}{ps_134-5_029-v0008}aspheric \var{\xout{functor}}{}{typewriter deletion}{ps_134-5_029-v0009} map i: A$\hand{\rightarrow}$B , \hand{\uline{\textcolor{black}{does not depend}}} (up to cano-%
nical isomorphism) on the choice of i , and can be described as the com-%
position of $\hand{\overline{\textcolor{black}{i_B}}}$ followed by a quasi-inverse for $\hand{\overline{\textcolor{black}{i_A}}}$ . This of course
is very much in keeping with the ``inspiring assumption'' (section 28),
which just means that up to unique isomorphism, there is indeed but \hand{\uline{\textcolor{black}{one}}}

\margin{29}

%% END SOURCE ps_134-5_029
\ifnum\EditionMode=1
\psapparatusstart
\psentry{ps_134-5_029-v0001}{1}{typewriter deletion}{\xout{functor of a}}{}{The typewriter strikes out this material; it leaves the corrected running text.}
\psentry{ps_134-5_029-v0002}{2}{typewriter deletion}{\xout{e}}{}{The typewriter strikes out this material; it leaves the corrected running text.}
\psentry{ps_134-5_029-v0003}{3}{typewriter deletion}{\xout{\illegible{1}}}{}{The typewriter strikes out this material; it leaves the corrected running text.}
\psentry{ps_134-5_029-v0004}{4}{writing over}{\over{\unsure{C}}{B}}{\hand{B}}{The superposed handwritten reading replaces the underlying reading.}
\psentry{ps_134-5_029-v0005}{5}{writing over}{\over{of}{for}}{\hand{for}}{The superposed handwritten reading replaces the underlying reading.}
\psentry{ps_134-5_029-v0006}{6}{typewriter deletion}{\xout{on}}{}{The typewriter strikes out this material; it leaves the corrected running text.}
\psentry{ps_134-5_029-v0007}{7}{typed insertion}{\inserted[typed]{between}}{between}{The typed insertion is incorporated at its indicated position; insertion brackets leave the running text.}
\psentry{ps_134-5_029-v0008}{8}{insertion}{\inserted{$\underline{W}$-}}{\hand{$\underline{W}$-}}{The handwritten insertion is incorporated at its indicated position; insertion brackets leave the running text.}
\psentry{ps_134-5_029-v0009}{9}{typewriter deletion}{\xout{functor}}{}{The typewriter strikes out this material; it leaves the corrected running text.}
\psapparatusend
\fi
\pssource{ps_134-5_030}{134-5/30}{https://grothendieck.umontpellier.fr/web/viewer.html?file=../134-5.pdf\#page=31}
\psnumber{ps_134-5_030-v0001}{1}
\psnumber{ps_134-5_030-v0002}{2}
\psnumber{ps_134-5_030-v0003}{3}
\psnumber{ps_134-5_030-v0004}{4}
\psnumber{ps_134-5_030-v0005}{5}
\psnumber{ps_134-5_030-v0006}{6}
\psnumber{ps_134-5_030-v0007}{7}
\psnumber{ps_134-5_030-v0008}{8}
\psnumber{ps_134-5_030-v0009}{9}
\psnumber{ps_134-5_030-v0010}{10}
\psnumber{ps_134-5_030-v0011}{11}
\psnumber{ps_134-5_030-v0012}{12}
\psnumber{ps_134-5_030-v0013}{13}
\psnumber{ps_134-5_030-v0014}{14}
\psnumber{ps_134-5_030-v0015}{15}
\psnumber{ps_134-5_030-v0016}{16}
\psnumber{ps_134-5_030-v0017}{17}
\psnumber{ps_134-5_030-v0018}{18}
%% BEGIN SOURCE ps_134-5_030
% Pursuing Stacks (cotes 134-2 to 134-8): archive cote 134-5, site page 30 (PDF page 31), typescript page 217.
% typescript. Draft compared: ps4.tex:1417-1496.
% AI-transcribed by Claude Opus 5.5 (claude-opus-5-5, adaptive thinking, effort high), Message Batches API, batch ps2 (msgbatch_01NVwr1Yuo8qJf5rykeGBXPK), 2026-11-12; reviewed against both supplied scan halves; see CHECK_DECISIONS.tsv.
\scanpage{134-5/30}
\mspage{217}
equivalence from Hot$_B$ to Hot$_A$  (both categories being equivalent to
Hot($\underline{W}$)).  Here we are thinking of course of the extension of the``assump-%
tion'' contemplated earlier,  when usual weak equivalence is replaced by
a basic localizer $\underline{W}$ as above.  It seems plausible that the assumption
holds true in all cases  -  anyhow we did'nt have at any moment to make
explicit use of it. (besides at a moment dr\unsure{a}wing inspiration from it ...\hand{)}
\hand{\uline{\textcolor{black}{Proposition}}} 2.   Let  i: $A\hand{\to} B$ be a map in (Cat).

a)  If i is $\underline{W}$-aspheric,  and A\^{}is totally $\underline{W}$-aspheric,  then B\^{}is to\unsure{t}-%
tally $\underline{W}$-aspheric too.

b)  Assume \var{\inserted{A $\underline{W}$-aspheric, and that}}{\hand{A $\underline{W}$-aspheric, and that}}{insertion}{ps_134-5_030-v0001} B is \var{\deleted{a}}{}{deletion}{ps_134-5_030-v0002}a $\underline{W}$-test category,  admitting \var{\deleted{a}}{}{deletion}{ps_134-5_030-v0003} \var{\inserted[typed]{the separating}}{the separating}{typed insertion}{ps_134-5_030-v0004} $\underline{W}_B$-homotopy  inter-%
val\var{\inserted[typed]{$\underline{I}=$}}{$\underline{I}=$}{typed insertion}{ps_134-5_030-v0005}$(I,\hand{\delta}_o,\hand{\delta}_1)$, \var{\inserted[typed]{satisfying the homotopy condition \xout{\illegible{1}} \hand{(T H 1)} of page \hand{50.}}}{satisfying the homotopy condition  \hand{(T H 1)} of page \hand{50.}}{typed insertion}{ps_134-5_030-v0006} Consider the following conditions:

\var{\over{\illegible{1}}{1}}{\hand{1}}{writing over}{ps_134-5_030-v0007})  $i^*(I)$ is $\underline{W}$-aspheric over $e_{A^\wedge}$ . \var{\deleted{and A is W-aspheric.}}{}{deletion}{ps_134-5_030-v0008}

2)  i is $\underline{W}$-aspheric ,

3)  $i^*(I)$  is $\underline{W}$-aspheric ,  \var{\deleted{and A is $\underline{W}$-aspheric}}{}{deletion}{ps_134-5_030-v0009}

We have the implic ations
\[ 1)\ \hand{\Longrightarrow}\ 2)\ \hand{\Longrightarrow}\ 3) \quad , \]
hence if A is totally $\underline{W}$-aspheric (hence  3)\hand{$\Longrightarrow$}1) )  all three conditions
are equivalent.

Proof:  \var{\xout{Fxxx}}{}{typewriter deletion}{ps_134-5_030-v0010} a)  We have to prove that if b,b' are in B,  then
their product in B\^{} is $\underline{W}$-aspheric ,  but by assumption on i we know that
the images of b,b' by $i^*$  are $\underline{W}$-aspheric,  hence (as A\^{}is totally $\underline{W}$-as-%
pheric)  their product
\[ i^*(b) \ x\ i^*(b') \ = i^*(bxb') \]
is $\underline{W}$-aspheric too, hence so is bxb'  by criterion (ii') of prop. above.

b)  The homotopy condition \var{\inserted{(TH1)}}{\hand{(TH1)}}{insertion}{ps_134-5_030-v0011} referred to\var{\xout{\illegible{1}}}{}{typewriter deletion}{ps_134-5_030-v0012} means that all objects of
B are  $\underline{I}$-contractib\var{\over{\unsure{le}}{le}}{\hand{le}}{writing over}{ps_134-5_030-v0013}.  As $i^*$ commutes with finite products,  it follows
that the objects $i^*$(b) of A\^{}, for b in B,  are $i^*(\underline{I})$-contractible.  When
$i^*(I)$ is $\underline{W}$-aspheric over $e_{A^\wedge}$,  this implies that \var{\xout{\unsure{e(b)xxx}}}{}{typewriter deletion}{ps_134-5_030-v0014} so are the
objects $i^*$(b),  \var{\xout{hence}}{}{typewriter deletion}{ps_134-5_030-v0015} a fortiori they are  $\underline{W}$-aspheric  (as by assumption
$e_{A^\wedge}$ is $\underline{W}$-aspheric ).  Thus  1)\hand{$\Longrightarrow$}2),  and  2)\hand{$\Longrightarrow$}3)  is trivial.
\hand{\uline{\textcolor{black}{Corollary}}} :  Let  i: $A\hand{\to} B$ be a $\underline{W}$-aspheric map in (Cat),  assume that
A is totally $\underline{W}$-aspheric,  and B is a local $\underline{W}$-test category.  Then both
A and B are strict $\underline{W}$-test categories.

Indeed,  \var{\inserted[typed]{\xout{xx} by a) above we see that B is totally \xout{\unsure{Wx(x}}}}{ by a) above we see that B is totally }{typed insertion}{ps_134-5_030-v0016} \var{\deleted{\xout{\unsure{asx\allowbreak{}ix\allowbreak{}isx\allowbreak{}inx\allowbreak{}$\underline{W}$x\allowbreak{}x\allowbreak{}andx\allowbreak{}Ax\allowbreak{}isx\allowbreak{}$\underline{W}$-aspheric,x\allowbreak{}x\allowbreak{}Bx\allowbreak{}isx\allowbreak{}a,x\allowbreak{}x\allowbreak{}hencex\allowbreak{}Bx\allowbreak{}is}}}}{}{deletion}{ps_134-5_030-v0017}
\var{\deleted{\xout{\unsure{x\allowbreak{}x\allowbreak{}beingx\allowbreak{}ax\allowbreak{}$\underline{W}$-testx\allowbreak{}category,x\allowbreak{}x\allowbreak{}andx\allowbreak{}beingx\allowbreak{}ax\allowbreak{}x\allowbreak{}x\allowbreak{}}}}}{}{deletion}{ps_134-5_030-v0018} $\underline{W}$-aspheric,  hence B is a
strict $\underline{W}$-test category.  In order to prove that so is A,  we only have
to show that  A\^{}admits a separating homotopy interval for $\underline{W}_A$ .  By assump
\margin{30}

%% END SOURCE ps_134-5_030
\ifnum\EditionMode=1
\psapparatusstart
\psentry{ps_134-5_030-v0001}{1}{insertion}{\inserted{A $\underline{W}$-aspheric, and that}}{\hand{A $\underline{W}$-aspheric, and that}}{The handwritten insertion is incorporated at its indicated position; insertion brackets leave the running text.}
\psentry{ps_134-5_030-v0002}{2}{deletion}{\deleted{a}}{}{The author cancels this material by hand; it leaves the corrected running text.}
\psentry{ps_134-5_030-v0003}{3}{deletion}{\deleted{a}}{}{The author cancels this material by hand; it leaves the corrected running text.}
\psentry{ps_134-5_030-v0004}{4}{typed insertion}{\inserted[typed]{the separating}}{the separating}{The typed insertion is incorporated at its indicated position; insertion brackets leave the running text.}
\psentry{ps_134-5_030-v0005}{5}{typed insertion}{\inserted[typed]{$\underline{I}=$}}{$\underline{I}=$}{The typed insertion is incorporated at its indicated position; insertion brackets leave the running text.}
\psentry{ps_134-5_030-v0006}{6}{typed insertion}{\inserted[typed]{satisfying the homotopy condition \xout{\illegible{1}} \hand{(T H 1)} of page \hand{50.}}}{satisfying the homotopy condition  \hand{(T H 1)} of page \hand{50.}}{The typed insertion is incorporated at its indicated position; insertion brackets leave the running text.}
\psentry{ps_134-5_030-v0007}{7}{writing over}{\over{\illegible{1}}{1}}{\hand{1}}{The superposed handwritten reading replaces the underlying reading.}
\psentry{ps_134-5_030-v0008}{8}{deletion}{\deleted{and A is W-aspheric.}}{}{The author cancels this material by hand; it leaves the corrected running text.}
\psentry{ps_134-5_030-v0009}{9}{deletion}{\deleted{and A is $\underline{W}$-aspheric}}{}{The author cancels this material by hand; it leaves the corrected running text.}
\psentry{ps_134-5_030-v0010}{10}{typewriter deletion}{\xout{Fxxx}}{}{The typewriter strikes out this material; it leaves the corrected running text.}
\psentry{ps_134-5_030-v0011}{11}{insertion}{\inserted{(TH1)}}{\hand{(TH1)}}{The handwritten insertion is incorporated at its indicated position; insertion brackets leave the running text.}
\psentry{ps_134-5_030-v0012}{12}{typewriter deletion}{\xout{\illegible{1}}}{}{The typewriter strikes out this material; it leaves the corrected running text.}
\psentry{ps_134-5_030-v0013}{13}{writing over}{\over{\unsure{le}}{le}}{\hand{le}}{The superposed handwritten reading replaces the underlying reading.}
\psentry{ps_134-5_030-v0014}{14}{typewriter deletion}{\xout{\unsure{e(b)xxx}}}{}{The typewriter strikes out this material; it leaves the corrected running text.}
\psentry{ps_134-5_030-v0015}{15}{typewriter deletion}{\xout{hence}}{}{The typewriter strikes out this material; it leaves the corrected running text.}
\psentry{ps_134-5_030-v0016}{16}{typed insertion}{\inserted[typed]{\xout{xx} by a) above we see that B is totally \xout{\unsure{Wx(x}}}}{ by a) above we see that B is totally }{The typed insertion is incorporated at its indicated position; insertion brackets leave the running text.}
\psentry{ps_134-5_030-v0017}{17}{deletion}{\deleted{\xout{\unsure{asx\allowbreak{}ix\allowbreak{}isx\allowbreak{}inx\allowbreak{}$\underline{W}$x\allowbreak{}x\allowbreak{}andx\allowbreak{}Ax\allowbreak{}isx\allowbreak{}$\underline{W}$-aspheric,x\allowbreak{}x\allowbreak{}Bx\allowbreak{}isx\allowbreak{}a,x\allowbreak{}x\allowbreak{}hencex\allowbreak{}Bx\allowbreak{}is}}}}{}{The author cancels this material by hand; it leaves the corrected running text.}
\psentry{ps_134-5_030-v0018}{18}{deletion}{\deleted{\xout{\unsure{x\allowbreak{}x\allowbreak{}beingx\allowbreak{}ax\allowbreak{}$\underline{W}$-testx\allowbreak{}category,x\allowbreak{}x\allowbreak{}andx\allowbreak{}beingx\allowbreak{}ax\allowbreak{}x\allowbreak{}x\allowbreak{}}}}}{}{The author cancels this material by hand; it leaves the corrected running text.}
\psapparatusend
\fi
\pssource{ps_134-5_031}{134-5/31}{https://grothendieck.umontpellier.fr/web/viewer.html?file=../134-5.pdf\#page=32}
\psnumber{ps_134-5_031-v0001}{1}
\psnumber{ps_134-5_031-v0002}{2}
\psnumber{ps_134-5_031-v0003}{3}
\psnumber{ps_134-5_031-v0004}{4}
\psnumber{ps_134-5_031-v0005}{5}
\psnumber{ps_134-5_031-v0006}{6}
\psnumber{ps_134-5_031-v0007}{7}
\psnumber{ps_134-5_031-v0008}{8}
%% BEGIN SOURCE ps_134-5_031
% Pursuing Stacks (cotes 134-2 to 134-8): archive cote 134-5, site page 31 (PDF page 32), typescript page 218.
% typescript. Draft compared: ps4.tex:1480-1551.
% AI-transcribed by Claude Opus 5.5 (claude-opus-5-5, adaptive thinking, effort high), Message Batches API, batch ps2 (msgbatch_01NVwr1Yuo8qJf5rykeGBXPK), 2026-11-12; reviewed against both supplied scan halves; see CHECK_DECISIONS.tsv.
\scanpage{134-5/31}
\mspage{218}

tion on B,  there is \var{\xout{such an interval}}{}{typewriter deletion}{ps_134-5_031-v0001} \var{\inserted{a}}{\hand{a}}{insertion}{ps_134-5_031-v0002} separating $\underline{W}_B$-homotopy interval
$\underline{I} = (I,\hand{\delta}_0,\hand{\delta}_1)$ in $B^\wedge$.  The exactness properties of $i^*$  imply that
$i^*(\underline{I})$  is a separating interval in $A^\wedge$,  \var{\xout{and}}{}{typewriter deletion}{ps_134-5_031-v0003} the asphericity condition on
i implies that moreover $i^*(I)$ is $\underline{W}$-aspheric,  hence $\underline{W}$-aspheric over $e_{A^\wedge}$
% A pen spacing slash separates the following A-hat from is; it is not a mathematical slash.
as $A^\wedge$ is  totally $\underline{W}$-aspheric,  qed.

To finish \var{\xout{these}}{}{typewriter deletion}{ps_134-5_031-v0004} with the more formal properties of the notion of
$\underline{W}$-aspheric maps in (Cat),  let's give a list of the standard stability
conditions for this notion,  with respect notably to composition,  base
change,  and cartesian products:

\hand{\uline{\textcolor{black}{Proposition 3}}}.   a)  Consider two maps  \var{\xout{AxxxxxB}}{}{typewriter deletion}{ps_134-5_031-v0005}
\[ A \overset{i}{\hand{\longrightarrow}} B \overset{j}{\hand{\longrightarrow}} C \]
in (Cat).   Then if  i and j are $\underline{W}$-aspheric,  so is ji.  If ji and i are
$\underline{W}$-aspheric,  so is j .  Any isomorphism in (Cat) is $\underline{W}$-aspheric.

b)   Let
% Hand-drawn arrows are marked blue; typed objects remain plain.
% A long pen curve runs beneath this square and partly over the next line; its function is unresolved.
\diagram{\begin{tikzcd}
A \& A' \arrow[color=blue, l, "f"'] \arrow[color=blue, d, "i'"] \\
B \arrow[color=blue, from=u, "i"'] \& B' \arrow[color=blue, l, "g"]
\end{tikzcd}}
be a cartesian square in (Cat),  assume i is $\underline{W}$-aspheric and g is \var{\xout{xx}}{}{typewriter deletion}{ps_134-5_031-v0006}fibe-%
ring  (for instance an induction functor  $B_{/F} \hand{\to} B$ ,  with  F in $B^\wedge$).
Then i' is $\underline{W}$-aspheric  (and, of course,  f is equally fibering).

c)   Let
\[ i: A \hand{\to} B , \quad i' : A' \hand{\to} B' \]
be two W-aspheric maps in (Cat),  then
\[ i\mathrm{x}i' : A\mathrm{x}A' \hand{\to} B\mathrm{x}B' \]
is $\underline{W}$-aspheric.

\hand{\uline{\textcolor{black}{Proof}}}.  Property a) is formal,  in terms of the criterion (ii') of
prop. 1 .  Property c) follows formally from the criterion (iii),  and
the canonical isomorphism
\[ (A\mathrm{x}A')_{/b\mathrm{x}b'} \hand{\simeq} (A_{/b})\mathrm{x}(A'_{/b'}) \quad , \]
and the fact that a product of two $\underline{W}$-aspheric objects of (Cat) is again
$\underline{W}$-aspheric  (cf.  prop. of page 167,  making use of  \var{\xout{Loc 3)}}{}{typewriter deletion}{ps_134-5_031-v0007} the localiza-%
tion condition  Loc 3)  on $\underline{W}$  in it's full generality).   Another proof
goes via  a) and b),  by reducing first (using a)) to the case when
either i or i' are identities,  and using the fact that a projection map
$C\mathrm{x}B \hand{\to} B$ in (Cat) is a fibration.

We are left with proving property b).  For this,  we note that the
asphericity condition (iii) on a map $i: A \hand{\to} B$  just means that \var{\xout{the}}{}{typewriter deletion}{ps_134-5_031-v0008} for
\margin{31}

%% END SOURCE ps_134-5_031
\ifnum\EditionMode=1
\psapparatusstart
\psentry{ps_134-5_031-v0001}{1}{typewriter deletion}{\xout{such an interval}}{}{The typewriter strikes out this material; it leaves the corrected running text.}
\psentry{ps_134-5_031-v0002}{2}{insertion}{\inserted{a}}{\hand{a}}{The handwritten insertion is incorporated at its indicated position; insertion brackets leave the running text.}
\psentry{ps_134-5_031-v0003}{3}{typewriter deletion}{\xout{and}}{}{The typewriter strikes out this material; it leaves the corrected running text.}
\psentry{ps_134-5_031-v0004}{4}{typewriter deletion}{\xout{these}}{}{The typewriter strikes out this material; it leaves the corrected running text.}
\psentry{ps_134-5_031-v0005}{5}{typewriter deletion}{\xout{AxxxxxB}}{}{The typewriter strikes out this material; it leaves the corrected running text.}
\psentry{ps_134-5_031-v0006}{6}{typewriter deletion}{\xout{xx}}{}{The typewriter strikes out this material; it leaves the corrected running text.}
\psentry{ps_134-5_031-v0007}{7}{typewriter deletion}{\xout{Loc 3)}}{}{The typewriter strikes out this material; it leaves the corrected running text.}
\psentry{ps_134-5_031-v0008}{8}{typewriter deletion}{\xout{the}}{}{The typewriter strikes out this material; it leaves the corrected running text.}
\psapparatusend
\fi
\pssource{ps_134-5_032}{134-5/32}{https://grothendieck.umontpellier.fr/web/viewer.html?file=../134-5.pdf\#page=33}
\psnumber{ps_134-5_032-v0001}{1}
\psnumber{ps_134-5_032-v0002}{2}
\psnumber{ps_134-5_032-v0003}{3}
\psnumber{ps_134-5_032-v0004}{4}
\psnumber{ps_134-5_032-v0005}{5}
\psnumber{ps_134-5_032-v0006}{6}
\psnumber{ps_134-5_032-v0007}{7}
\psnumber{ps_134-5_032-v0008}{8}
\psnumber{ps_134-5_032-v0009}{9}
%% BEGIN SOURCE ps_134-5_032
% Pursuing Stacks (cotes 134-2 to 134-8): archive cote 134-5, site page 32 (PDF page 33), typescript page 219.
% typescript. Draft compared: ps4.tex:1536-1601.
% AI-transcribed by Claude Opus 5.5 (claude-opus-5-5, adaptive thinking, effort high), Message Batches API, batch ps2 (msgbatch_01NVwr1Yuo8qJf5rykeGBXPK), 2026-11-12; reviewed against both supplied scan halves; see CHECK_DECISIONS.tsv.
\scanpage{134-5/32}
\mspage{219}

any base change of the type
\[ B_{/b} \hand{\longrightarrow} B \quad , \]
where b is in B,  the corresponding map
\[ i_{/b} \ : \ A\mathrm{x}_{B}B_{/b} \hand{\simeq} A_{/b} \hand{\longrightarrow} B_{/b} \]
is in $\underline{W}$ .  Applying this to the case of $i' : A' \hand{\longrightarrow} B'$ ,  and an  object
b' in B',  and denoting by b it's image in B,  \var{\deleted{and}}{}{deletion}{ps_134-5_032-v0001} using transitivity of
base change,  we get a cartesian square
% Hand-drawn arrows are marked blue; typed objects remain plain.
\diagram{
\begin{tikzcd}
A_{/b} \arrow[color=blue, d, "\hand{i_{/b}}"] \& A'_{/b'} \arrow[color=blue, l] \arrow[color=blue, d, "\hand{i'_{/b'}}"] \\
B_{/b} \& B'_{/b'} \arrow[color=blue, l]
\end{tikzcd}
\quad ,
}
and we got to prove $i'_{/b'}$ is in $\underline{W}$  i.e. $A'_{/b'}$ is $\underline{W}$-aspheric,  using
the fact that we know the same holds for $i_{/b}$  i.e. $A_{/b}$ is $\underline{W}$-aspheric.
Thus all we have to prove is that the first horizontal arrow
is in $\underline{W}$ .  But we check at once that the condition that $B' \hand{\longrightarrow} B$ is fiber-%
ring  implies that the induced functor  \var{\xout{B}}{}{typewriter deletion}{ps_134-5_032-v0002}
\[ B'_{/b'} \hand{\longrightarrow} B_{/b} \]
is fibering too,  and moreover \var{\over{b}{h}}{\hand{h}}{writing over}{ps_134-5_032-v0003}as fibers which have final objects.  Hen-%
ce by base change,  the same \var{\deleted{conditions}}{}{deletion}{ps_134-5_032-v0004}\var{\inserted{properties}}{\hand{properties}}{insertion}{ps_134-5_032-v0005} hold for  \var{\xout{A' $_{/b'}$}}{}{typewriter deletion}{ps_134-5_032-v0006}
\[ A'_{/b'} \hand{\rightarrow} A_{/b} \quad . \]
This reminds us of the ``fibration condition''  L 5  (page 164),  which
should ensure that a fibration with $\underline{W}$-aspheric fibers is in $\underline{W}$ .  We did
not include this axiom among the assumptions\hand{(}recalled above\hand{)}  we want to
make on $\underline{W}$ .  However it turns out that the assumptions we do make here
imply already the fibration condition,  as well as the dual condition
on cofibrations.  We'll give a proof later  -  in order not to diverge
at present from our main purpose.  There will not be any vicious circle,
as all we're going to use of prop.3  for the formalism of asphericity
structures  and canonical modelizers is the first part a).  I included
b) and c) for the sake of completeness,  and because \var{\over{\illegible{1}}{\unsure{c}}}{\hand{\unsure{c}}}{writing over}{ps_134-5_032-v0007}) is useful for
dealing with products of two categories,  notably of two test categories
-  a theme which has been long pending,  and on which I would like to
digress next,  before getting invol\var{\deleted{\unsure{l}}}{}{deletion}{ps_134-5_032-v0008}ved with asphericity structures.

\hand{\uline{\textcolor{black}{Remark 4}}} .  In part a) of proposition,  if j and ji are $\underline{W}$-aspheric,  we
cannot conclude that i is . \hand{\illegible{1}} \var{\over{i}{I}}{\hand{I}}{writing over}{ps_134-5_032-v0009}f C is the final object of (Cat),  this
means that a map between $\underline{W}$-aspheric objects in (Cat) need not be
$\underline{W}$-aspheric .

\margin{32}

%% END SOURCE ps_134-5_032
\ifnum\EditionMode=1
\psapparatusstart
\psentry{ps_134-5_032-v0001}{1}{deletion}{\deleted{and}}{}{The author cancels this material by hand; it leaves the corrected running text.}
\psentry{ps_134-5_032-v0002}{2}{typewriter deletion}{\xout{B}}{}{The typewriter strikes out this material; it leaves the corrected running text.}
\psentry{ps_134-5_032-v0003}{3}{writing over}{\over{b}{h}}{\hand{h}}{The superposed handwritten reading replaces the underlying reading.}
\psentry{ps_134-5_032-v0004}{4}{deletion}{\deleted{conditions}}{}{The author cancels this material by hand; it leaves the corrected running text.}
\psentry{ps_134-5_032-v0005}{5}{insertion}{\inserted{properties}}{\hand{properties}}{The handwritten insertion is incorporated at its indicated position; insertion brackets leave the running text.}
\psentry{ps_134-5_032-v0006}{6}{typewriter deletion}{\xout{A' $_{/b'}$}}{}{The typewriter strikes out this material; it leaves the corrected running text.}
\psentry{ps_134-5_032-v0007}{7}{writing over}{\over{\illegible{1}}{\unsure{c}}}{\hand{\unsure{c}}}{The superposed handwritten reading replaces the underlying reading.}
\psentry{ps_134-5_032-v0008}{8}{deletion}{\deleted{\unsure{l}}}{}{The author cancels this material by hand; it leaves the corrected running text.}
\psentry{ps_134-5_032-v0009}{9}{writing over}{\over{i}{I}}{\hand{I}}{The superposed handwritten reading replaces the underlying reading.}
\psapparatusend
\fi
\pssource{ps_134-5_033}{134-5/33}{https://grothendieck.umontpellier.fr/web/viewer.html?file=../134-5.pdf\#page=34}
\psnumber{ps_134-5_033-v0001}{1}
\psnumber{ps_134-5_033-v0002}{2}
\psnumber{ps_134-5_033-v0003}{3}
\psnumber{ps_134-5_033-v0004}{4}
\psnumber{ps_134-5_033-v0005}{5}
\psnumber{ps_134-5_033-v0006}{6}
\psnumber{ps_134-5_033-v0007}{7}
\psnumber{ps_134-5_033-v0008}{8}
\psnumber{ps_134-5_033-v0009}{9}
\psnumber{ps_134-5_033-v0010}{10}
\psnumber{ps_134-5_033-v0011}{11}
\psnumber{ps_134-5_033-v0012}{12}
\psnumber{ps_134-5_033-v0013}{13}
\psnumber{ps_134-5_033-v0014}{14}
\psnumber{ps_134-5_033-v0015}{15}
\psnumber{ps_134-5_033-v0016}{16}
%% BEGIN SOURCE ps_134-5_033
% Pursuing Stacks (cotes 134-2 to 134-8): archive cote 134-5, site page 33 (PDF page 34), typescript page 220.
% typescript. Draft compared: ps4.tex:1578-1653.
% AI-transcribed by Claude Opus 5.5 (claude-opus-5-5, adaptive thinking, effort high), Message Batches API, batch ps2 (msgbatch_01NVwr1Yuo8qJf5rykeGBXPK), 2026-11-12; reviewed against both supplied scan halves; see CHECK_DECISIONS.tsv.
\scanpage{134-5/33}
\mspage{220}

11.6.

73. I am not quite through yet with generalities on asphericity crite-%
ria for a map in (Cat), it turns out - it was just getting prohibitive-%
ly late last night to go on ! \margin{\P} From now on, I'll drop the qualifying
$\underline{W}$ when speaking of asphericity, test categories, modelizers and the
like, as by now it is well understood I guess there is a given $\underline{W}$\var{\xout{\illegible{1}}}{}{typewriter deletion}{ps_134-5_033-v0001} around
in all we are doing. It'll be enough to be specific \var{\over{\unsure{an}}{in}}{\hand{in}}{writing over}{ps_134-5_033-v0002} those (presumably
rare) instances when working with more than one basic localizer.

Coming back to the last remark in yesterday\hand{'}s notes, a good illustra-%
tion is the case of a functor $i : A \hand{\rightarrow} B$ of aspheric objects of (Cat),
when A is \var{\xout{the}}{}{typewriter deletion}{ps_134-5_033-v0003}\var{\inserted[typed]{a}}{a}{typed insertion}{ps_134-5_033-v0004} final \var{\xout{category}}{}{typewriter deletion}{ps_134-5_033-v0005} object in (Cat) i.e. a one-point discrete
category. Then i \hand{\uline{\textcolor{black}{is aspheric iff}}} i(a) \hand{\uline{\textcolor{black}{is an initial object of}}} B - an
extremely stringent \var{\inserted{extra}}{\hand{extra}}{insertion}{ps_134-5_033-v0006} condition indeed !

In part a) of proposition 3 above, stating that isomorphisms in
(Cat) are aspheric, it would have been timely to be more generous -
indeed, \hand{\uline{\textcolor{black}{any equivalence of categories in}}} (Cat) \hand{\uline{\textcolor{black}{is aspheric}}}. This asso-%
ciates immediately with a map of topoi which is an equivalence being
(trivially so) aspheric (\var{\xout{for}}{}{typewriter deletion}{ps_134-5_033-v0007}\var{\inserted[typed]{in}}{in}{typed insertion}{ps_134-5_033-v0008} the case\hand{,} say\hand{,} when $\underline{W}$ is the usual notion of
weak equivalence, the only one for the time being when the notion is
extended from maps in (Cat) to maps between topoi ). In the context of
(Cat), the basic modelizer, we can give a still more general case of
aspheric maps, both instructive and useful :

\hand{\uline{\textcolor{black}{Proposition}}} 4. Let
% Hand-drawn arrows are marked blue; typed objects remain plain.
\diagram{\begin{tikzcd} A \ar[color=blue,r,shift left,"f"] \& B \ar[color=blue,l,shift left,"g"] \end{tikzcd}}
be a pair of adjoint functors between the objects A,B in (Cat), with f
left and g right adjoint. Then f is aspheric.

\var{\deleted{\hand{\uline{Remark 5.}}}}{}{deletion}{ps_134-5_033-v0009} Indeed, by the adjunction formula we immediately get
\[ A_{/b} \hand{\simeq} A_{/g(b)} \]
(as a matter of fact, $f^{*}(b) = g(b)$), hence this category has a final
object and hence is aspheric, qed.

\var{\inserted{\uline{Remark 5.}}}{\hand{\uline{Remark 5.}}}{insertion}{ps_134-5_033-v0010} \unsure{\var{\deleted{The}}{}{deletion}{ps_134-5_033-v0011}|The} conclusion of \var{\deleted{this theorem}}{}{deletion}{ps_134-5_033-v0012}\var{\inserted{prop. 4}}{\hand{prop. 4}}{insertion}{ps_134-5_033-v0013} is mute about g , which will right-%
ly strike as unfair. Dualizing, we could say that $(\var{\over{\unsure{f}}{g}}{\hand{g}}{writing over}{ps_134-5_033-v0014}^{o} ,f^{o})$ is a pair
of adjoint functors between $A^{o}$ and $B^{o}$ , and therefore $g^{o}$ is aspheric.
We will express this fact (by lack of a more suggestive name) by saying
that g is a \hand{\uline{\textcolor{black}{coaspheric}}} map . \var{\inserted{For a functor $i : A \rightarrow B$,}}{\hand{For a functor $i : A \rightarrow B$,}}{insertion}{ps_134-5_033-v0015} In terms of the usual criterion\var{\inserted{for $i^{o}$}}{\hand{for $i^{o}$}}{insertion}{ps_134-5_033-v0016}\hand{,} it just

\margin{33}

%% END SOURCE ps_134-5_033
\ifnum\EditionMode=1
\psapparatusstart
\psentry{ps_134-5_033-v0001}{1}{typewriter deletion}{\xout{\illegible{1}}}{}{The typewriter strikes out this material; it leaves the corrected running text.}
\psentry{ps_134-5_033-v0002}{2}{writing over}{\over{\unsure{an}}{in}}{\hand{in}}{The superposed handwritten reading replaces the underlying reading.}
\psentry{ps_134-5_033-v0003}{3}{typewriter deletion}{\xout{the}}{}{The typewriter strikes out this material; it leaves the corrected running text.}
\psentry{ps_134-5_033-v0004}{4}{typed insertion}{\inserted[typed]{a}}{a}{The typed insertion is incorporated at its indicated position; insertion brackets leave the running text.}
\psentry{ps_134-5_033-v0005}{5}{typewriter deletion}{\xout{category}}{}{The typewriter strikes out this material; it leaves the corrected running text.}
\psentry{ps_134-5_033-v0006}{6}{insertion}{\inserted{extra}}{\hand{extra}}{The handwritten insertion is incorporated at its indicated position; insertion brackets leave the running text.}
\psentry{ps_134-5_033-v0007}{7}{typewriter deletion}{\xout{for}}{}{The typewriter strikes out this material; it leaves the corrected running text.}
\psentry{ps_134-5_033-v0008}{8}{typed insertion}{\inserted[typed]{in}}{in}{The typed insertion is incorporated at its indicated position; insertion brackets leave the running text.}
\psentry{ps_134-5_033-v0009}{9}{deletion}{\deleted{\hand{\uline{Remark 5.}}}}{}{The author cancels this material by hand; it leaves the corrected running text.}
\psentry{ps_134-5_033-v0010}{10}{insertion}{\inserted{\uline{Remark 5.}}}{\hand{\uline{Remark 5.}}}{The handwritten insertion is incorporated at its indicated position; insertion brackets leave the running text.}
\psentry{ps_134-5_033-v0011}{11}{deletion}{\deleted{The}}{}{The author cancels this material by hand; it leaves the corrected running text.}
\psentry{ps_134-5_033-v0012}{12}{deletion}{\deleted{this theorem}}{}{The author cancels this material by hand; it leaves the corrected running text.}
\psentry{ps_134-5_033-v0013}{13}{insertion}{\inserted{prop. 4}}{\hand{prop. 4}}{The handwritten insertion is incorporated at its indicated position; insertion brackets leave the running text.}
\psentry{ps_134-5_033-v0014}{14}{writing over}{$\over{\unsure{f}}{g}$}{$\hand{g}$}{The superposed handwritten reading replaces the underlying reading.}
\psentry{ps_134-5_033-v0015}{15}{insertion}{\inserted{For a functor $i : A \rightarrow B$,}}{\hand{For a functor $i : A \rightarrow B$,}}{The handwritten insertion is incorporated at its indicated position; insertion brackets leave the running text.}
\psentry{ps_134-5_033-v0016}{16}{insertion}{\inserted{for $i^{o}$}}{\hand{for $i^{o}$}}{The handwritten insertion is incorporated at its indicated position; insertion brackets leave the running text.}
\psapparatusend
\fi
\pssource{ps_134-5_034}{134-5/34}{https://grothendieck.umontpellier.fr/web/viewer.html?file=../134-5.pdf\#page=35}
\psnumber{ps_134-5_034-v0001}{1}
\psnumber{ps_134-5_034-v0002}{2}
\psnumber{ps_134-5_034-v0003}{3}
\psnumber{ps_134-5_034-v0004}{4}
\psnumber{ps_134-5_034-v0005}{5}
\psnumber{ps_134-5_034-v0006}{6}
\psnumber{ps_134-5_034-v0007}{7}
\psnumber{ps_134-5_034-v0008}{8}
\psnumber{ps_134-5_034-v0009}{9}
\psnumber{ps_134-5_034-v0010}{10}
\psnumber{ps_134-5_034-v0011}{11}
\psnumber{ps_134-5_034-v0012}{12}
\psnumber{ps_134-5_034-v0013}{13}
\psnumber{ps_134-5_034-v0014}{14}
\psnumber{ps_134-5_034-v0015}{15}
\psnumber{ps_134-5_034-v0016}{16}
\psnumber{ps_134-5_034-v0017}{17}
%% BEGIN SOURCE ps_134-5_034
% Pursuing Stacks (cotes 134-2 to 134-8): archive cote 134-5, site page 34 (PDF page 35), typescript page 221.
% typescript. Draft compared: ps4.tex:1640-1703.
% AI-transcribed by Claude Opus 5.5 (claude-opus-5-5, adaptive thinking, effort high), Message Batches API, batch ps2 (msgbatch_01NVwr1Yuo8qJf5rykeGBXPK), 2026-11-12; reviewed against both supplied scan halves; see CHECK_DECISIONS.tsv.
\scanpage{134-5/34}
\mspage{221}
means that for any b in B, the category
\[ {}_{b\backslash}A \overset{\mathrm{def}}{=} Ax_B({}_{b\backslash}B) \;=\; \text{category of pairs } (x,p)\hand{,} \text{ with } x \text{ in } A \]
\[ \text{and } p : b \hand{\rightarrow} \var{\over{\unsure{g}}{i}}{\hand{i}}{writing over}{ps_134-5_034-v0001}(x) \]
\margin{$i(x)$}
is aspheric. (NB In the case of the functor g from B to A, the corres-%
ponding categories ${}_{a\backslash}B$ \var{\inserted{even}}{\hand{even}}{insertion}{ps_134-5_034-v0002} have \hand{\uline{\textcolor{black}{initial}}} \hand{\uline{\textcolor{black}{objects}}}.) This condition comes in
here rather formally, we'll see later though that it has a \var{\xout{very}}{}{typewriter deletion}{ps_134-5_034-v0003} \var{\inserted[typed]{quite}}{quite}{typed insertion}{ps_134-5_034-v0004} remarkab-%
le interpretation, in terms of a very strong property of cohomological
``cofinality'' of the functor i, implying the usual notion of i being a
``cofinal'' functor\hand{,}or A being ``cofinal''in B (namely i giving rise to
isomorphisms $\underset{\hand{\longrightarrow}}{\lim}_A \hand{\longrightarrow} \underset{\hand{\longrightarrow}}{\lim}_B$ for any \var{\xout{indexed system}}{}{typewriter deletion}{ps_134-5_034-v0005} direct system
B $\hand{\longrightarrow}$M with values in a category M admitting direct limits)\hand{,} as the ``dime\hand{n}-%
sion zero'' shadow of this ``all dimensions'' property. To make an anology
which will acquire more precise meaning later, the asphericity property
for a map i in (Cat) can be viewed as a (slightly weakened) version of
a \hand{\uline{\textcolor{black}{proper map with aspheric fibers}}}, whereas the coasphericity property
appears as the corresponding version of a \hand{\uline{\textcolor{black}{smooth map with aspheric fibers}}}.
The qualification ``slightly weakened'' reflects notably in the fact that
the notions \var{\xout{considered}}{}{typewriter deletion}{ps_134-5_034-v0006} \var{\inserted[typed]{of asphericity and coasphericity}}{of asphericity and coasphericity}{typed insertion}{ps_134-5_034-v0007} are \hand{\uline{\textcolor{black}{not}}} stable under arbitrary base change -
but rather\hand{,} asphericity for a map is stable under base change by \var{\xout{a}}{}{typewriter deletion}{ps_134-5_034-v0008} \hand{\uline{\textcolor{black}{fibra-%
tion functors}}}(more generally, by \var{\xout{a}}{}{typewriter deletion}{ps_134-5_034-v0009} smooth maps in (Cat)), whereas
coasphericity is stable under base change by \hand{\uline{\textcolor{black}{cofibration functors}}}
(more generally, by proper maps in (Cat)). Thus, whereas a functor
i: A $\hand{\longrightarrow}$ B which is an equivalence of categories is clearly both aspheric
and coaspheric, this property \var{\xout{does}}{}{typewriter deletion}{ps_134-5_034-v0010} \var{\inserted[typed]{is}}{is}{typed insertion}{ps_134-5_034-v0011} not preserved by arbitrary base
change e.g. passage to fibers : \var{\inserted{e.g.}}{\hand{e.g.}}{insertion}{ps_134-5_034-v0012} some fibers may be empty, and therefore
are neither adpheric nor coaspheric !

The most comprehensive property \var{\xout{including both}}{}{typewriter deletion}{ps_134-5_034-v0013} for a map in (Cat),
implying both asphericity and coasphericity, is to be in U$\underline{W}$ i.e. a
``\hand{\uline{\textcolor{black}{universal weak equivalence}}}'' - namely it is in $\underline{W}$ and remains so after
any base change. These maps deserve the name of ``\hand{\uline{\textcolor{black}{trivial Serre fibrati-%
ons}}}''. They include\var{\inserted[typed]{all}}{all}{typed insertion}{ps_134-5_034-v0014} maps with aspheric fibers which are either ``proper''
or ``smooth'', for instance those which are either \var{\inserted[typed]{co}}{co}{typed insertion}{ps_134-5_034-v0015}fibering or \var{\xout{co}}{}{typewriter deletion}{ps_134-5_034-v0016}fibering
functors (in the usual sense of category theory). We'll come back upon
these notions in a \var{\xout{more}}{}{typewriter deletion}{ps_134-5_034-v0017} systematic way in the next part of the reflecti\hand{\unsure{ons}}

74. As announced yesterday, I would like still to make an overdue
digression on products of test categories, before embarking on the
notion of asphericity structure. It will be useful to begin with some

\margin{34}

%% END SOURCE ps_134-5_034
\ifnum\EditionMode=1
\psapparatusstart
\psentry{ps_134-5_034-v0001}{1}{writing over}{$\over{\unsure{g}}{i}$}{$\hand{i}$}{The superposed handwritten reading replaces the underlying reading.}
\psentry{ps_134-5_034-v0002}{2}{insertion}{\inserted{even}}{\hand{even}}{The handwritten insertion is incorporated at its indicated position; insertion brackets leave the running text.}
\psentry{ps_134-5_034-v0003}{3}{typewriter deletion}{\xout{very}}{}{The typewriter strikes out this material; it leaves the corrected running text.}
\psentry{ps_134-5_034-v0004}{4}{typed insertion}{\inserted[typed]{quite}}{quite}{The typed insertion is incorporated at its indicated position; insertion brackets leave the running text.}
\psentry{ps_134-5_034-v0005}{5}{typewriter deletion}{\xout{indexed system}}{}{The typewriter strikes out this material; it leaves the corrected running text.}
\psentry{ps_134-5_034-v0006}{6}{typewriter deletion}{\xout{considered}}{}{The typewriter strikes out this material; it leaves the corrected running text.}
\psentry{ps_134-5_034-v0007}{7}{typed insertion}{\inserted[typed]{of asphericity and coasphericity}}{of asphericity and coasphericity}{The typed insertion is incorporated at its indicated position; insertion brackets leave the running text.}
\psentry{ps_134-5_034-v0008}{8}{typewriter deletion}{\xout{a}}{}{The typewriter strikes out this material; it leaves the corrected running text.}
\psentry{ps_134-5_034-v0009}{9}{typewriter deletion}{\xout{a}}{}{The typewriter strikes out this material; it leaves the corrected running text.}
\psentry{ps_134-5_034-v0010}{10}{typewriter deletion}{\xout{does}}{}{The typewriter strikes out this material; it leaves the corrected running text.}
\psentry{ps_134-5_034-v0011}{11}{typed insertion}{\inserted[typed]{is}}{is}{The typed insertion is incorporated at its indicated position; insertion brackets leave the running text.}
\psentry{ps_134-5_034-v0012}{12}{insertion}{\inserted{e.g.}}{\hand{e.g.}}{The handwritten insertion is incorporated at its indicated position; insertion brackets leave the running text.}
\psentry{ps_134-5_034-v0013}{13}{typewriter deletion}{\xout{including both}}{}{The typewriter strikes out this material; it leaves the corrected running text.}
\psentry{ps_134-5_034-v0014}{14}{typed insertion}{\inserted[typed]{all}}{all}{The typed insertion is incorporated at its indicated position; insertion brackets leave the running text.}
\psentry{ps_134-5_034-v0015}{15}{typed insertion}{\inserted[typed]{co}}{co}{The typed insertion is incorporated at its indicated position; insertion brackets leave the running text.}
\psentry{ps_134-5_034-v0016}{16}{typewriter deletion}{\xout{co}}{}{The typewriter strikes out this material; it leaves the corrected running text.}
\psentry{ps_134-5_034-v0017}{17}{typewriter deletion}{\xout{more}}{}{The typewriter strikes out this material; it leaves the corrected running text.}
\psapparatusend
\fi
\pssource{ps_134-5_035}{134-5/35}{https://grothendieck.umontpellier.fr/web/viewer.html?file=../134-5.pdf\#page=36}
\psnumber{ps_134-5_035-v0001}{1}
\psnumber{ps_134-5_035-v0002}{2}
\psnumber{ps_134-5_035-v0003}{3}
\psnumber{ps_134-5_035-v0004}{4}
\psnumber{ps_134-5_035-v0005}{5}
%% BEGIN SOURCE ps_134-5_035
% Pursuing Stacks (cotes 134-2 to 134-8): archive cote 134-5, site page 35 (PDF page 36), typescript page 222.
% typescript. Draft compared: ps4.tex:1687-1754.
% AI-transcribed by Claude Opus 5.5 (claude-opus-5-5, adaptive thinking, effort high), Message Batches API, batch ps2 (msgbatch_01NVwr1Yuo8qJf5rykeGBXPK), 2026-11-12; reviewed against both supplied scan halves; see CHECK_DECISIONS.tsv.
\scanpage{134-5/35}
\mspage{222}
generalities on presheaves on a product category $A\mathrm{x}B$ , where for the
time being $A,B$ are any two objects of (Cat). The following notation
is often useful, for a pair of presheaves
\[ F\hand{\in} \mathrm{Ob}\, A\hat{} \quad , \qquad G\hand{\in} \mathrm{Ob}\, B\hat{} \quad , \]
introducing an``\hand{\uline{\textcolor{black}{external product}}}''
\[ F\hand{\boxtimes} G \hand{\in} \var{\xout{\illegible{1}}}{}{typewriter deletion}{ps_134-5_035-v0001}\mathrm{Ob}(A\mathrm{x}B)\hat{} \]
by the formula
\[ F\hand{\boxtimes} G\ (a,b) = F(a)\mathrm{x}G(b) \ . \]
Introducing the two projections
\[ p_1 : A\mathrm{x}B \hand{\longrightarrow} A \quad , \qquad p_2 : A\mathrm{x}B \hand{\longrightarrow} B \quad , \]
and the corresponding inverse image functors $p_i^{*}$ , we have
\[ F\hand{\boxtimes} G \ = \ p_1^{*}(F)\mathrm{x}p_2^{*}(G) \qquad . \]
Of course $F\hand{\boxtimes} G$ depends bifunctorially on the pair $(F,G)$, and it is ea-%
sily checked\hand{,}by the way\hand{,}that the corresponding functor
\[ A\hat{}\mathrm{x}\ B\hat{} \hand{\longrightarrow} (A\mathrm{x}B)\hat{} \quad , \qquad (F,G)\hand{\longmapsto} F\hand{\boxtimes} G \quad , \]
is fully faithful.

We have a tautological canonical isomorphism
\[ (A\mathrm{x}B)_{/F\hand{\boxtimes} G} \hand{\xrightarrow{\sim}} A_{/F}\ \mathrm{x}\ B_{/G} \qquad , \]
and hence the

\hand{\uline{\textcolor{black}{Proposition}}}\hand{1}\var{\over{.}{:}}{\hand{:}}{writing over}{ps_134-5_035-v0002} a)If $F$ and $G$ are aspheric objects in $A\hat{}$and $B\hat{}$\hand{/}respectively,
then $F\hand{\boxtimes} G$ is an aspheric object of $(A\mathrm{x}B)\hat{}$.
b) If $u : F\hand{\rightarrow} F'$ and $v : G\hand{\rightarrow} G'$ are aspheric maps in $A\hat{}$and $B\hat{}$ res-%
pectively, then
\[ u\hand{\boxtimes} v \ : \ F\hand{\boxtimes} G \hand{\rightarrow} F'\hand{\boxtimes} G' \]
is an aspheric map in $(A\mathrm{x}B)\hat{}$.

Indeed, this follows respectively from the fact that a product of
two aspheric objects (resp. aspheric maps) in (Cat) is again aspheric.

\hand{\uline{\textcolor{black}{Corollary}}} 1\var{\over{.}{:}}{\hand{:}}{writing over}{ps_134-5_035-v0003}Let $F$ in $A\hat{}$ be aspheric over the final object $e_{A\hat{}}$ , then
$F\hand{\boxtimes} e_{B\hat{}} = p_1^{*}(F)$ is aspheric over the final object of $(A\mathrm{x}B)\hat{}$.

\var{\inserted[typed]{\xout{Assumption}}}{}{typed insertion}{ps_134-5_035-v0004}\hand{\uline{\textcolor{black}{Corollary}}} 2. Assume $A$ and $B$ are totally aspheric, then so is $A\mathrm{x}B$.

We have to check\var{\xout{s}}{}{typewriter deletion}{ps_134-5_035-v0005} that the product elements
\[ (a,b)\mathrm{x}(a',b') \ = (a\mathrm{x}a')\hand{\boxtimes}(b\mathrm{x}b') \]
are aspheric, which follows from the assumption (namely $a\mathrm{x}a'$ and $b\mathrm{x}b'$
aspheric) and part a) of the proposition.

\margin{35}

%% END SOURCE ps_134-5_035
\ifnum\EditionMode=1
\psapparatusstart
\psentry{ps_134-5_035-v0001}{1}{typewriter deletion}{$\xout{\illegible{1}}$}{}{The typewriter strikes out this material; it leaves the corrected running text.}
\psentry{ps_134-5_035-v0002}{2}{writing over}{\over{.}{:}}{\hand{:}}{The superposed handwritten reading replaces the underlying reading.}
\psentry{ps_134-5_035-v0003}{3}{writing over}{\over{.}{:}}{\hand{:}}{The superposed handwritten reading replaces the underlying reading.}
\psentry{ps_134-5_035-v0004}{4}{typed insertion}{\inserted[typed]{\xout{Assumption}}}{}{The typed insertion is incorporated at its indicated position; insertion brackets leave the running text.}
\psentry{ps_134-5_035-v0005}{5}{typewriter deletion}{\xout{s}}{}{The typewriter strikes out this material; it leaves the corrected running text.}
\psapparatusend
\fi
\pssource{ps_134-5_036}{134-5/36}{https://grothendieck.umontpellier.fr/web/viewer.html?file=../134-5.pdf\#page=37}
\psnumber{ps_134-5_036-v0001}{1}
\psnumber{ps_134-5_036-v0002}{2}
\psnumber{ps_134-5_036-v0003}{3}
\psnumber{ps_134-5_036-v0004}{4}
\psnumber{ps_134-5_036-v0005}{5}
\psnumber{ps_134-5_036-v0006}{6}
\psnumber{ps_134-5_036-v0007}{7}
\psnumber{ps_134-5_036-v0008}{8}
\psnumber{ps_134-5_036-v0009}{9}
\psnumber{ps_134-5_036-v0010}{10}
\psnumber{ps_134-5_036-v0011}{11}
\psnumber{ps_134-5_036-v0012}{12}
\psnumber{ps_134-5_036-v0013}{13}
\psnumber{ps_134-5_036-v0014}{14}
\psnumber{ps_134-5_036-v0015}{15}
%% BEGIN SOURCE ps_134-5_036
% Pursuing Stacks (cotes 134-2 to 134-8): archive cote 134-5, site page 36 (PDF page 37), typescript page 223.
% typescript. Draft compared: ps4.tex:1735-1810.
% AI-transcribed by Claude Opus 5.5 (claude-opus-5-5, adaptive thinking, effort high), Message Batches API, batch ps2 (msgbatch_01NVwr1Yuo8qJf5rykeGBXPK), 2026-11-12; reviewed against both supplied scan halves; see CHECK_DECISIONS.tsv.
\scanpage{134-5/36}
\mspage{223}

Assume now A is a local test category, i.e.\var{\inserted[typed]{$A^\wedge$}}{$A^\wedge$}{typed insertion}{ps_134-5_036-v0001} admits a separating
interval
\[ \overline{I} = (I,\hand{\delta}_0,\hand{\delta}_1) \]
such that I is aspheric over the final object of $A^\wedge$. Then $p_1^*(\underline{I})$ is of
course a separating interval in $(AxB)^\wedge$, which by cor.1 is aspheric over
the final object. Hence:

\hand{\uline{\textcolor{black}{Proposition}}} 2: If A is a local test category, so is AxB for any B in
(Cat). If A is a test category, then so is AxB for any aspheric B .

The second assertion follows, remembering that a test category is
just an \hand{\uline{\textcolor{black}{aspheric}}} local test category. Using cor. 2, we get :

\hand{\uline{\textcolor{black}{Coro\unsure{\var{\over[typed]{a}{l}}{l}{typewriter overwriting}{ps_134-5_036-v0002}}l\unsure{\var{\over[typed]{l}{a}}{a}{typewriter overwriting}{ps_134-5_036-v0003}}ry}}}: If A is a strict test category, and B totally aspheric,
then AxB is a strict test category. In particular, if A and B are
strict test categories, so is their product.

We need only remember that strict test categories are just test
categories which are totally aspheric.

As an illustration, we get the fact that the categories of multi-%
complexes of various kinds, which we can \var{\inserted[typed]{even}}{even}{typed insertion}{ps_134-5_036-v0004} take mixed (semisimplicial
in some variables, cubical in some others, and hemispherical say in
others still) are strict modelizers(as generally granted), which corres
\hand{-}ponds to the fact that a \var{\xout{xx}}{}{typewriter deletion}{ps_134-5_036-v0005} finite product of \var{\xout{the}}{}{typewriter deletion}{ps_134-5_036-v0006} standard semisimplicial
cubical and hemispherical\var{\inserted[typed]{(strict)}}{(strict)}{typed insertion}{ps_134-5_036-v0007} test categories \hand{$\Delta$} , \hand{$\Box$} and \hand{\unsure{$\bigcirc$}} , is again a
strict test category.

I would like now to dwell a little upon the comparison of ``homotopy
models'', \var{\inserted[typed]{using}}{using}{typed insertion}{ps_134-5_036-v0008}\var{\xout{\illegible{1}}}{}{typewriter deletion}{ps_134-5_036-v0009} respectively \var{\inserted[typed]{\xout{using}}}{}{typed insertion}{ps_134-5_036-v0010}\var{\xout{in}}{}{typewriter deletion}{ps_134-5_036-v0011} two test categories A and B, namely
working in $A^\wedge$ and $B^\wedge$ respectively. More specifically, we have two
descriptions of Hot (short for Hot($\underline{W}$) here ) , namely as \var{\xout{Hot$_A$}}{}{typewriter deletion}{ps_134-5_036-v0012}
\[ \mathrm{Hot}_A = \underline{W}_A{}^{-1}A^\wedge \quad\text{and}\quad \mathrm{Hot}_B = \underline{W}_B{}^{-1}B^\wedge \quad , \]
and we want to describe conveniently the ta\var{\inserted[typed]{u}}{u}{typed insertion}{ps_134-5_036-v0013}tological equivalence between
these two categories ( this equivalence being defined up to unique iso-%
morphism). The most \hand{``}tautological\hand{''} way indeed is to use the basic modeli-%
zer (Cat) and it's localization (Hot) as the int ermediary \var{\xout{step}}{}{typewriter deletion}{ps_134-5_036-v0014} ,
i.e. using the diagram of equivalences of categories
% Hand-drawn arrows are marked blue; typed objects remain plain.
% The hemispherical-category symbol above has an inner arc inside its circle; the bigcirc approximation remains doubtful.
\diagram{
(1)\qquad
\begin{tikzcd}[column sep=small]
\mathrm{Hot}_A \ar[color=blue, dr, "\hand{\approx}", "\overline{\hand{i}}_{\hand{A}}"'] \& \& \mathrm{Hot}_B \ar[color=blue, dl, "\overline{\hand{i}}_{\hand{B}}", "\hand{\approx}"'] \\
\& (\mathrm{Hot}) \&
\end{tikzcd}
\qquad .
}
Remember we have a handy quasi-inverse \hand{$\overline{\textcolor{black}{\var{\over{\unsure{j}}{j}}{\hand{j}}{writing over}{ps_134-5_036-v0015}_B}}$}

\margin{36}

%% END SOURCE ps_134-5_036
\ifnum\EditionMode=1
\psapparatusstart
\psentry{ps_134-5_036-v0001}{1}{typed insertion}{\inserted[typed]{$A^\wedge$}}{$A^\wedge$}{The typed insertion is incorporated at its indicated position; insertion brackets leave the running text.}
\psentry{ps_134-5_036-v0002}{2}{typewriter overwriting}{\over[typed]{a}{l}}{l}{The superposed typed reading replaces the underlying reading.}
\psentry{ps_134-5_036-v0003}{3}{typewriter overwriting}{\over[typed]{l}{a}}{a}{The superposed typed reading replaces the underlying reading.}
\psentry{ps_134-5_036-v0004}{4}{typed insertion}{\inserted[typed]{even}}{even}{The typed insertion is incorporated at its indicated position; insertion brackets leave the running text.}
\psentry{ps_134-5_036-v0005}{5}{typewriter deletion}{\xout{xx}}{}{The typewriter strikes out this material; it leaves the corrected running text.}
\psentry{ps_134-5_036-v0006}{6}{typewriter deletion}{\xout{the}}{}{The typewriter strikes out this material; it leaves the corrected running text.}
\psentry{ps_134-5_036-v0007}{7}{typed insertion}{\inserted[typed]{(strict)}}{(strict)}{The typed insertion is incorporated at its indicated position; insertion brackets leave the running text.}
\psentry{ps_134-5_036-v0008}{8}{typed insertion}{\inserted[typed]{using}}{using}{The typed insertion is incorporated at its indicated position; insertion brackets leave the running text.}
\psentry{ps_134-5_036-v0009}{9}{typewriter deletion}{\xout{\illegible{1}}}{}{The typewriter strikes out this material; it leaves the corrected running text.}
\psentry{ps_134-5_036-v0010}{10}{typed insertion}{\inserted[typed]{\xout{using}}}{}{The typed insertion is incorporated at its indicated position; insertion brackets leave the running text.}
\psentry{ps_134-5_036-v0011}{11}{typewriter deletion}{\xout{in}}{}{The typewriter strikes out this material; it leaves the corrected running text.}
\psentry{ps_134-5_036-v0012}{12}{typewriter deletion}{\xout{Hot$_A$}}{}{The typewriter strikes out this material; it leaves the corrected running text.}
\psentry{ps_134-5_036-v0013}{13}{typed insertion}{\inserted[typed]{u}}{u}{The typed insertion is incorporated at its indicated position; insertion brackets leave the running text.}
\psentry{ps_134-5_036-v0014}{14}{typewriter deletion}{\xout{step}}{}{The typewriter strikes out this material; it leaves the corrected running text.}
\psentry{ps_134-5_036-v0015}{15}{writing over}{$\over{\unsure{j}}{j}$}{$\hand{j}$}{The superposed handwritten reading replaces the underlying reading.}
\psapparatusend
\fi
\pssource{ps_134-5_037}{134-5/37}{https://grothendieck.umontpellier.fr/web/viewer.html?file=../134-5.pdf\#page=38}
\psnumber{ps_134-5_037-v0001}{1}
\psnumber{ps_134-5_037-v0002}{2}
\psnumber{ps_134-5_037-v0003}{3}
\psnumber{ps_134-5_037-v0004}{4}
\psnumber{ps_134-5_037-v0005}{5}
\psnumber{ps_134-5_037-v0006}{6}
\psnumber{ps_134-5_037-v0007}{7}
\psnumber{ps_134-5_037-v0008}{8}
\psnumber{ps_134-5_037-v0009}{9}
\psnumber{ps_134-5_037-v0010}{10}
%% BEGIN SOURCE ps_134-5_037
% Pursuing Stacks (cotes 134-2 to 134-8): archive cote 134-5, site page 37 (PDF page 38), typescript page 224.
% typescript. Draft compared: ps4.tex:1791-1864.
% AI-transcribed by Claude Opus 5.5 (claude-opus-5-5, adaptive thinking, effort high), Message Batches API, batch ps2 (msgbatch_01NVwr1Yuo8qJf5rykeGBXPK), 2026-11-12; reviewed against both supplied scan halves; see CHECK_DECISIONS.tsv.
\scanpage{134-5/37}
\mspage{224}
to $\overline{i_{B}}$ , using the \var{\xout{fxx}}{}{typewriter deletion}{ps_134-5_037-v0001} canonical functor
\[ j_B = i_B^{*} : \ (Cat) \hand{\longrightarrow} B^{\wedge} \ , \qquad X \hand{\longrightarrow} B_{/X} \ \var{\xout{\illegible{2}Hom}}{}{typewriter deletion}{ps_134-5_037-v0002} \ . \]
Thus we get a description of an equivalence
(\var{\over{**}{\unsure{1'}}}{\hand{\unsure{1'}}}{writing over}{ps_134-5_037-v0003}) \[ \overline{j_B}\ \overline{i_A} \ : \ Hot_A \hand{\xrightarrow{\approx}} Hot_B \ , \]
whose quasi-inverse of course is just the similar $\overline{j_A}\ \overline{i_B}$ . We can vary
a little this description, admittedly cumbersome in practice, by repla-%
cing $\overline{\var{\over{i}{j}}{\hand{j}}{writing over}{ps_134-5_037-v0004}_{B}}$ by the isomorphic functor $\overline{i^{*}}$, where $i : B \hand{\to} (Cat)$ is
any test functor from B to (Cat) (while we have to keep however $\overline{i_A}$ as it
is , without the possibility of replacing $i_A$ by a more amenable test
functor).

Another way for comparing $Hot_A$ and $Hot_B$ arises, as we saw yester%
day, whenever we have an \hand{\uline{\textcolor{black}{aspheric}}} functor
\hand{(2)} \[ i : A \hand{\longrightarrow} B \ , \]
by just taking
\hand{(2')} \[ \overline{i^{*}} \ : \ Hot_B \hand{\xrightarrow{\approx}} Hot_A \ . \]
This of course is about the \var{\xout{most}}{}{typewriter deletion}{ps_134-5_037-v0005} simplest way imaginable, all the more
as the functor $i^{*} : B^{\wedge} \hand{\to} A^{\wedge}$ commutes to arbitrary direct and inverse
limits, just perfect for comparing constructions in $B^{\wedge}$and constructions
in $A^{\wedge}$ - whereas the functor $j_B i_A : A^{\wedge} \hand{\longrightarrow} B^{\wedge}$ , giving rise to (*) above,
commutes just to sums and fiber\unsure{\"{e}}d products, not to amalgamated sums nor
to products, sad enough! We could add here that if we got \hand{\uline{\textcolor{black}{two}}}
aspheric functors from A to B, namely i plus
\[ i' : \ A \hand{\longrightarrow} B \ , \]
then (as immediately checked) any map between these functors
\[ u : \ i \hand{\longrightarrow} i' \]
gives rise to an \var{\xout{equivalence}}{}{typewriter deletion}{ps_134-5_037-v0006}\var{\inserted[typed]{isomorphism}}{isomorphism}{typed insertion}{ps_134-5_037-v0007} between the corresponding equivalences
% Hand-drawn arrows are marked blue; typed objects remain plain.
\diagram{\begin{tikzcd}[column sep=large]
Hot_B \ar[color=blue, r, bend left, "\overline{i^{*}}", ""{name=U, below}] \ar[color=blue, r, bend right, "\overline{i'^{*}}"', ""{name=D, above}] \& Hot_A \ar[color=blue, from=D, to=U, "\hand{\overline{u}}"']
\end{tikzcd}} ,

which is nothing but the canonical\var{\xout{x}}{}{typewriter deletion}{ps_134-5_037-v0008}isomorphism referred to yesterday
\margin{(*)} (both functors being canonically isomorphic to (\var{\over{**}{\unsure{1'}}}{\hand{\unsure{1'}}}{writing over}{ps_134-5_037-v0009})), and yields the
most evident way for \var{\xout{x}}{}{typewriter deletion}{ps_134-5_037-v0010}``computing'' the latter. Thus it turns out that
the isomorphism $\overline{u}$ does not depend upon the choice of u. If we want
to ignore this fact and look at the situation sternly, in a wholly compu-%
tational spirit, we could present things by stating that we get a
contravariant functor, from aspheric maps $i: A \hand{\to} B$ to equivalences
$Hot_B \hand{\to} Hot_A$ :
\margin{37}

%% END SOURCE ps_134-5_037
\ifnum\EditionMode=1
\psapparatusstart
\psentry{ps_134-5_037-v0001}{1}{typewriter deletion}{\xout{fxx}}{}{The typewriter strikes out this material; it leaves the corrected running text.}
\psentry{ps_134-5_037-v0002}{2}{typewriter deletion}{$\xout{\illegible{2}Hom}$}{}{The typewriter strikes out this material; it leaves the corrected running text.}
\psentry{ps_134-5_037-v0003}{3}{writing over}{\over{**}{\unsure{1'}}}{\hand{\unsure{1'}}}{The superposed handwritten reading replaces the underlying reading.}
\psentry{ps_134-5_037-v0004}{4}{writing over}{$\over{i}{j}$}{$\hand{j}$}{The superposed handwritten reading replaces the underlying reading.}
\psentry{ps_134-5_037-v0005}{5}{typewriter deletion}{\xout{most}}{}{The typewriter strikes out this material; it leaves the corrected running text.}
\psentry{ps_134-5_037-v0006}{6}{typewriter deletion}{\xout{equivalence}}{}{The typewriter strikes out this material; it leaves the corrected running text.}
\psentry{ps_134-5_037-v0007}{7}{typed insertion}{\inserted[typed]{isomorphism}}{isomorphism}{The typed insertion is incorporated at its indicated position; insertion brackets leave the running text.}
\psentry{ps_134-5_037-v0008}{8}{typewriter deletion}{\xout{x}}{}{The typewriter strikes out this material; it leaves the corrected running text.}
\psentry{ps_134-5_037-v0009}{9}{writing over}{\over{**}{\unsure{1'}}}{\hand{\unsure{1'}}}{The superposed handwritten reading replaces the underlying reading.}
\psentry{ps_134-5_037-v0010}{10}{typewriter deletion}{\xout{x}}{}{The typewriter strikes out this material; it leaves the corrected running text.}
\psapparatusend
\fi
\pssource{ps_134-5_038}{134-5/38}{https://grothendieck.umontpellier.fr/web/viewer.html?file=../134-5.pdf\#page=39}
\psnumber{ps_134-5_038-v0001}{1}
\psnumber{ps_134-5_038-v0002}{2}
\psnumber{ps_134-5_038-v0003}{3}
\psnumber{ps_134-5_038-v0004}{4}
\psnumber{ps_134-5_038-v0005}{5}
\psnumber{ps_134-5_038-v0006}{6}
\psnumber{ps_134-5_038-v0007}{7}
\psnumber{ps_134-5_038-v0008}{8}
\psnumber{ps_134-5_038-v0009}{9}
\psnumber{ps_134-5_038-v0010}{10}
\psnumber{ps_134-5_038-v0011}{11}
\psnumber{ps_134-5_038-v0012}{12}
\psnumber{ps_134-5_038-v0013}{13}
\psnumber{ps_134-5_038-v0014}{14}
\psnumber{ps_134-5_038-v0015}{15}
\psnumber{ps_134-5_038-v0016}{16}
\psnumber{ps_134-5_038-v0017}{17}
\psnumber{ps_134-5_038-v0018}{18}
%% BEGIN SOURCE ps_134-5_038
% Pursuing Stacks (cotes 134-2 to 134-8): archive cote 134-5, site page 38 (PDF page 39), typescript page 225.
% typescript. Draft compared: ps4.tex:1851-1911.
% AI-transcribed by Claude Opus 5.5 (claude-opus-5-5, adaptive thinking, effort high), Message Batches API, batch ps2 (msgbatch_01NVwr1Yuo8qJf5rykeGBXPK), 2026-11-12; reviewed against both supplied scan halves; see CHECK_DECISIONS.tsv.
\scanpage{134-5/38}
\mspage{225}

\[ \underline{Asph}(A,B)^{\circ}\hand{\longrightarrow}\underline{Hom}(Hot_{B},Hot_{A}) \quad , \qquad i \hand{\longmapsto} \overline{i^{*}} \quad , \]
where \uline{Asph} denotes the full subcategory of the functor category \uline{Hom}
made up with aspheric functors. This functor transforms arbitrary
arrows from the left hand side into isomorphisms on the right, and
therefore it factors through the fundamental groupoid (i.e. localization
of the left hand category with respect to the set of \hand{\uline{\textcolor{black}{all}}} it\unsure{(|'}s arrows):
\[ (\hand{\Pi}_{1}(\underline{Asph}(A,B)))^{\circ} \hand{\longrightarrow} \underline{Hom}(Hot_{B},Hot_{A}) \quad . \]
If we now remember that we had assumed A and B to be test categories
(otherwise the functors just written would be still defined, but their
values would not necessarily be \hand{\uline{\textcolor{black}{equivalences}}} between $Hot_{B}$ and $Hot_{A}$ ,
but merely functors between these), we may hope that this might imply
that the category $\underline{Asph}(A,B)$ of aspheric functors from A \var{\inserted{whenever non-empty}}{\hand{whenever non-empty}}{insertion}{ps_134-5_038-v0001} to be \var{\deleted{is}}{}{deletion}{ps_134-5_038-v0002}
1-connected. Whenever this is so, in any case, we get ``a priori''
(namely without any reference to \hand{(}\hand{\uline{\textcolor{black}{Hot}}}\hand{)} itself) a transitive system of
isomorphisms between the functors $\overline{i^{*}}$ , for i in $\underline{Asph}(A,B)$, \var{\xout{As a matter}}{}{typewriter deletion}{ps_134-5_038-v0003}
\var{\xout{of fact}}{}{typewriter deletion}{ps_134-5_038-v0004} hence a \var{\inserted[typed]{``canonical''}}{``canonical''}{typed insertion}{ps_134-5_038-v0005} \hand{\uline{\textcolor{black}{functor}}} $Hot_{B}\hand{\longrightarrow}Hot_{A}$\hand{,} \hand{\uline{\textcolor{black}{defined up to}}}\hand{\unsure{,}} \hand{\uline{\textcolor{black}{unique isomorphism}}}
(and which, in case A and B are test categories, or more generally
pseudo-test categories\var{\xout{*}}{}{typewriter deletion}{ps_134-5_038-v0006}, is an equivalence, and the one precisely
\margin{(1)} stemming from the diagram (\var{\over{*}{1}}{\hand{1}}{writing over}{ps_134-5_038-v0007})).

\hand{\uline{\textcolor{black}{Remark}}} Here the reflection slipped, almost against will, into a rela-%
ted one, about comparison of $Hot_{A}$ and $Hot_{B}$ for arbitrary A and B
(not necessarily test categories nor \hand{\uline{\textcolor{black}{even}}} pseudo test categories), using
aspheric functors $i: A \hand{\longrightarrow} B$ to get $i^{*} : Hot_{B}\hand{\longrightarrow}Hot_{A}$ \var{\over{,}{(}}{\hand{(}}{writing over}{ps_134-5_038-v0008} not necessa-%
rily an equivalence\hand{)}. As seen above, this functor depends rather loose-%
ly on the choice of i, and \var{\xout{our}}{}{typewriter deletion}{ps_134-5_038-v0009} we could develop comprehensive conditions
on A and B, not \var{\xout{at}}{}{typewriter deletion}{ps_134-5_038-v0010} at all of a test-condition nature, implying that
\var{\xout{there}}{}{typewriter deletion}{ps_134-5_038-v0011} just using the isomorphims $\overline{u}$ between these \var{\inserted{functors}}{\hand{functors}}{insertion}{ps_134-5_038-v0012}, we get a canonical
\hand{\uline{\textcolor{black}{transitive}}} system of isomorphisms between the\var{\over{\unsure{m}}{m}}{\hand{m}}{writing over}{ps_134-5_038-v0013}, \var{\deleted{namely}}{}{deletion}{ps_134-5_038-v0014} \var{\inserted{hence}}{\hand{hence}}{insertion}{ps_134-5_038-v0015} a \hand{\uline{\textcolor{black}{canonical}}}
functor $Hot_{B}\hand{\longrightarrow}Hot_{A}$ , not depending on the \var{\inserted{particular}}{\hand{particular}}{insertion}{ps_134-5_038-v0016} choice of any aspheric
functor $i: A\hand{\longrightarrow} B$. As we are mainly interested in the modelizing case
though, I dont think I should dwell on this much longer here. Anyhow,
in case both A and B are pseudo-test categories, and provided only
$\underline{Asph}(A,B)$ is 0-connected (not necessarily 1-connected), it follows \var{\xout{frm}}{}{typewriter deletion}{ps_134-5_038-v0017}
\margin{(1)} from comparison with the diagram (\var{\over{*}{1}}{\hand{1}}{writing over}{ps_134-5_038-v0018}) above that the isomorphisms $\overline{u}$
do give rise \hand{\uline{\textcolor{black}{ind}}}eed to a \hand{\uline{\textcolor{black}{transitive}}} system of isomorphisms between the
equivalences $\hand{\overline{\textcolor{black}{i}^{*}}}$ .

In most cases though, such as $A =\hand{\Delta}$ and $B =\hand{\Box}$ say (the test

\margin{38}

%% END SOURCE ps_134-5_038
\ifnum\EditionMode=1
\psapparatusstart
\psentry{ps_134-5_038-v0001}{1}{insertion}{\inserted{whenever non-empty}}{\hand{whenever non-empty}}{The handwritten insertion is incorporated at its indicated position; insertion brackets leave the running text.}
\psentry{ps_134-5_038-v0002}{2}{deletion}{\deleted{is}}{}{The author cancels this material by hand; it leaves the corrected running text.}
\psentry{ps_134-5_038-v0003}{3}{typewriter deletion}{\xout{As a matter}}{}{The typewriter strikes out this material; it leaves the corrected running text.}
\psentry{ps_134-5_038-v0004}{4}{typewriter deletion}{\xout{of fact}}{}{The typewriter strikes out this material; it leaves the corrected running text.}
\psentry{ps_134-5_038-v0005}{5}{typed insertion}{\inserted[typed]{``canonical''}}{``canonical''}{The typed insertion is incorporated at its indicated position; insertion brackets leave the running text.}
\psentry{ps_134-5_038-v0006}{6}{typewriter deletion}{\xout{*}}{}{The typewriter strikes out this material; it leaves the corrected running text.}
\psentry{ps_134-5_038-v0007}{7}{writing over}{\over{*}{1}}{\hand{1}}{The superposed handwritten reading replaces the underlying reading.}
\psentry{ps_134-5_038-v0008}{8}{writing over}{\over{,}{(}}{\hand{(}}{The superposed handwritten reading replaces the underlying reading.}
\psentry{ps_134-5_038-v0009}{9}{typewriter deletion}{\xout{our}}{}{The typewriter strikes out this material; it leaves the corrected running text.}
\psentry{ps_134-5_038-v0010}{10}{typewriter deletion}{\xout{at}}{}{The typewriter strikes out this material; it leaves the corrected running text.}
\psentry{ps_134-5_038-v0011}{11}{typewriter deletion}{\xout{there}}{}{The typewriter strikes out this material; it leaves the corrected running text.}
\psentry{ps_134-5_038-v0012}{12}{insertion}{\inserted{functors}}{\hand{functors}}{The handwritten insertion is incorporated at its indicated position; insertion brackets leave the running text.}
\psentry{ps_134-5_038-v0013}{13}{writing over}{\over{\unsure{m}}{m}}{\hand{m}}{The superposed handwritten reading replaces the underlying reading.}
\psentry{ps_134-5_038-v0014}{14}{deletion}{\deleted{namely}}{}{The author cancels this material by hand; it leaves the corrected running text.}
\psentry{ps_134-5_038-v0015}{15}{insertion}{\inserted{hence}}{\hand{hence}}{The handwritten insertion is incorporated at its indicated position; insertion brackets leave the running text.}
\psentry{ps_134-5_038-v0016}{16}{insertion}{\inserted{particular}}{\hand{particular}}{The handwritten insertion is incorporated at its indicated position; insertion brackets leave the running text.}
\psentry{ps_134-5_038-v0017}{17}{typewriter deletion}{\xout{frm}}{}{The typewriter strikes out this material; it leaves the corrected running text.}
\psentry{ps_134-5_038-v0018}{18}{writing over}{\over{*}{1}}{\hand{1}}{The superposed handwritten reading replaces the underlying reading.}
\psapparatusend
\fi
\pssource{ps_134-5_039}{134-5/39}{https://grothendieck.umontpellier.fr/web/viewer.html?file=../134-5.pdf\#page=40}
\psnumber{ps_134-5_039-v0001}{1}
\psnumber{ps_134-5_039-v0002}{2}
\psnumber{ps_134-5_039-v0003}{3}
\psnumber{ps_134-5_039-v0004}{4}
\psnumber{ps_134-5_039-v0005}{5}
%% BEGIN SOURCE ps_134-5_039
% Pursuing Stacks (cotes 134-2 to 134-8): archive cote 134-5, site page 39 (PDF page 40), typescript page 226.
% typescript. Draft compared: ps4.tex:1897-1956.
% AI-transcribed by Claude Opus 5.5 (claude-opus-5-5, adaptive thinking, effort high), Message Batches API, batch ps2 (msgbatch_01NVwr1Yuo8qJf5rykeGBXPK), 2026-11-12; reviewed against both supplied scan halves; see CHECK_DECISIONS.tsv.
\scanpage{134-5/39}
\mspage{226}

categories of standard simplices and standard cubes respectively), we
do no have any aspheric functor $A\hand{\to} B$ at hand, and presumably we may
well have that \uline{Asph}$(A,B)$ is empty\hand{,} poor it ! We now assume again that
A and B are test categories, hence $AxB$ is a test category, and the
natural idea for comparison of $\mathrm{Hot}_A$ and $\mathrm{Hot}_B$ is to use the diagram

% Hand-drawn arrows are marked blue; typed objects remain plain.
\diagram{\begin{tikzcd}[column sep=small]
 \& AxB \ar[color=blue,dl,"p_1"'] \ar[color=blue,dr,"p_2"] \& \\
 A \& \& B
\end{tikzcd}}\quad ,

where now $p_1$ and $p_2$ are aspheric (because B and A are aspheric). Thus
we get a corresponding diagram on the corresponding modelizers

\diagram{\begin{tikzcd}[column sep=small]
 \& (AxB)\hat{} \& \\
 A\hat{} \ar[color=blue,ur,"\hand{p_1^*}"] \& \& \var{\deleted{\unsure{Bt}}}{}{deletion}{ps_134-5_039-v0001}\, B\hat{} \ar[color=blue,ul,"\hand{p_2^*}"']
\end{tikzcd}}\quad ,

giving rise in turn to equivalences between the corresponding localiza-%
tions

(2)\qquad
\diagram{\begin{tikzcd}[column sep=small]
 \& \mathrm{Hot}_{AxB} \& \\
 \mathrm{Hot}_A \ar[color=blue,ur,"\hand{\overline{p_1^*}}","\hand{\approx}"'] \& \& \mathrm{Hot}_B \ar[color=blue,ul,"\hand{\overline{p_2^*}}"',"\hand{\approx}"]
\end{tikzcd}}\quad ,

These equivalences are compatible with the canonical \var{\xout{funct\unsure{s}}}{}{typewriter deletion}{ps_134-5_039-v0002} equivalences
with Hot itself, and hence they give rise to an equivalence $\mathrm{Hot}_A\hand{\xrightarrow{\approx}}\mathrm{Hot}_B$
which (up to canonical isomorphism) is indeed the canonical one. This
way for comparing $\mathrm{Hot}_A$ and $\mathrm{Hot}_B$ looks a lot more convenient than the
first one, as the functors $p_1^{*}$ and $p_2^{*}$ which serve as intermediaries
have all desirable exactness properties, and their \var{\inserted{very}}{\hand{very}}{insertion}{ps_134-5_039-v0003} description is the
simplest imaginable.

We are very close here to an Eilenberg-Zilber situation, which
will arise more specifically when A and B are both \hand{\uline{\textcolor{black}{strict}}} test categories\hand{,}
i.e. in the corresponding modelizers $A\hat{}$\hand{,}$B\hat{}$ products of models do cor-%
respond to products of the corresponding homotopy types. As seen above,
this implies the same for $A\var{\over{\unsure{x}}{x}}{\hand{x}}{writing over}{ps_134-5_039-v0004}B$. Thus, if F is in $A\hat{}$, G in $B\hat{}$, the
object $F\hand{\boxtimes} G$ in $(AxB)\hat{}$ is just a \var{\inserted[typed]{(AxB)-}}{(AxB)-}{typed insertion}{ps_134-5_039-v0005}model for the homotopy type described
respectively by F (in terms of A) and G (in terms of B). As a matter
of fact, the relation
\[ (AxB)_{F\hand{\boxtimes} G} \hand{\simeq} A_{/F}\ x\ B_{/G}\ \hand{,} \]
already noticed before\hand{,} implies that this interpretation holds, indepen-%
dently of strictness. In the classical statement of Eilenberg-Zilber
(as I recall it), we got $A=B$ (both categories A and B being just \hand{$\Delta$});

\margin{39}

%% END SOURCE ps_134-5_039
\ifnum\EditionMode=1
\psapparatusstart
\psentry{ps_134-5_039-v0001}{1}{deletion}{$\deleted{\unsure{Bt}}$}{}{The author cancels this material by hand; it leaves the corrected running text.}
\psentry{ps_134-5_039-v0002}{2}{typewriter deletion}{\xout{funct\unsure{s}}}{}{The typewriter strikes out this material; it leaves the corrected running text.}
\psentry{ps_134-5_039-v0003}{3}{insertion}{\inserted{very}}{\hand{very}}{The handwritten insertion is incorporated at its indicated position; insertion brackets leave the running text.}
\psentry{ps_134-5_039-v0004}{4}{writing over}{$\over{\unsure{x}}{x}$}{$\hand{x}$}{The superposed handwritten reading replaces the underlying reading.}
\psentry{ps_134-5_039-v0005}{5}{typed insertion}{\inserted[typed]{(AxB)-}}{(AxB)-}{The typed insertion is incorporated at its indicated position; insertion brackets leave the running text.}
\psapparatusend
\fi
\pssource{ps_134-5_040}{134-5/40}{https://grothendieck.umontpellier.fr/web/viewer.html?file=../134-5.pdf\#page=41}
\psnumber{ps_134-5_040-v0001}{1}
\psnumber{ps_134-5_040-v0002}{2}
\psnumber{ps_134-5_040-v0003}{3}
\psnumber{ps_134-5_040-v0004}{4}
\psnumber{ps_134-5_040-v0005}{5}
\psnumber{ps_134-5_040-v0006}{6}
\psnumber{ps_134-5_040-v0007}{7}
\psnumber{ps_134-5_040-v0008}{8}
%% BEGIN SOURCE ps_134-5_040
% Pursuing Stacks (cotes 134-2 to 134-8): archive cote 134-5, site page 40 (PDF page 41), typescript page 227.
% typescript. Draft compared: ps4.tex:1941-2002.
% AI-transcribed by Claude Opus 5.5 (claude-opus-5-5, adaptive thinking, effort high), Message Batches API, batch ps2 (msgbatch_01NVwr1Yuo8qJf5rykeGBXPK), 2026-11-12; reviewed against both supplied scan halves; see CHECK_DECISIONS.tsv.
\scanpage{134-5/40}
\mspage{227}

on the one hand the AxB - model is used in order to get readily the Kün-%
neth type relations for homology and cohomology, whereas we are interes-%
ted really in the A-model FxG. Using the diagonal map
\[ \hand{\delta} : A \hand{\longrightarrow} A\mathrm{x}A \quad , \]
we have of course
\[ F\mathrm{x}G \hand{\simeq} \hand{\delta}^{*}(F\hand{\boxtimes} G) \ , \]
more generally it is expected that the functor
\[ \hand{\delta}^{*} \ : \ (A\mathrm{x}A)^{\wedge} \hand{\longrightarrow} A^{\wedge} \]
is modelizing , i.e. gives a mean of passing from (AxA)-models to
A-models. This essentially translates into \hand{$\delta$} being an aspheric map -
something which will not be true for arbitrary A in (Cat). We get in
this respect:

\hand{\uline{\textcolor{black}{Proposition}}} 3. Let A be an object in (Cat). Then the diagonal map
\margin{x} \hand{$\delta$} : A $\hand{\to}$ AxA is aspheric iff A is totally aspheric.

This is just a ta\var{\inserted[typed]{u}}{u}{typed insertion}{ps_134-5_040-v0001}tology - one among many which tell us that the
conceptual set-up is OK indeed\var{\over{.}{!}}{\hand{!}}{writing over}{ps_134-5_040-v0002} A related tautology:

\hand{\uline{\textcolor{black}{Corollary}}}. Let i: A $\hand{\to}$ B, i': A $\hand{\to}$ B' be two aspheric functors with
same source A, then the corresponding functor
\[ (i,i') : \ A \hand{\longrightarrow} B\mathrm{x}B' \]
is aspheric, provided A is totally aspheric.

This can be seen either directly, or as a corollary of the
proposition, by viewing (i,i') as a compositi\unsure{on}
\[ A \hand{\xrightarrow{\delta}} A\mathrm{x}A \overset{i\mathrm{x}i'}{\hand{\longrightarrow}} B\mathrm{x}B' \qquad . \]

To sum up: when we got a bunch of strict test categories, and
a (possibly empty) bunch of aspheric functors between some of these,
using \var{\inserted[typed]{finite}}{finite}{typed insertion}{ps_134-5_040-v0003} products we get a (substantially larger !) bunch of strict test
categories, giving rise to correspo\var{\over[typed]{\unsure{g}}{n}}{n}{typewriter overwriting}{ps_134-5_040-v0004}ding strict modelizers for homoto-%
py types; and using projections, diagonal maps\hand{,} and the \var{\inserted[typed]{given}}{given}{typed insertion}{ps_134-5_040-v0005} aspheric func-%
tors, we get an impressive lot of aspheric maps between all these,
namely as many ways to ``commute'' from one type of ``homotopy models'' to
others. The simplest example\hand{:} \var{\deleted{if}}{}{deletion}{ps_134-5_040-v0006} start\var{\deleted{ing}}{}{deletion}{ps_134-5_040-v0007} with just one strict test
category A, taking products
\[ A^{I} \]
where I is any finite set, and the ``simplicial'' maps between these,
expressing contravariance of $A^{I}$ with respect to I. T\var{\over{h}{h}}{\hand{h}}{writing over}{ps_134-5_040-v0008}e case most commo%
% NOTE: the right edge of the scan cuts the end of this line after "commo"
ly used is A = \hand{$\Delta$}, giving rise to the formalism of semisimplicial
multicomplexes.

\margin{40}

%% END SOURCE ps_134-5_040
\ifnum\EditionMode=1
\psapparatusstart
\psentry{ps_134-5_040-v0001}{1}{typed insertion}{\inserted[typed]{u}}{u}{The typed insertion is incorporated at its indicated position; insertion brackets leave the running text.}
\psentry{ps_134-5_040-v0002}{2}{writing over}{\over{.}{!}}{\hand{!}}{The superposed handwritten reading replaces the underlying reading.}
\psentry{ps_134-5_040-v0003}{3}{typed insertion}{\inserted[typed]{finite}}{finite}{The typed insertion is incorporated at its indicated position; insertion brackets leave the running text.}
\psentry{ps_134-5_040-v0004}{4}{typewriter overwriting}{\over[typed]{\unsure{g}}{n}}{n}{The superposed typed reading replaces the underlying reading.}
\psentry{ps_134-5_040-v0005}{5}{typed insertion}{\inserted[typed]{given}}{given}{The typed insertion is incorporated at its indicated position; insertion brackets leave the running text.}
\psentry{ps_134-5_040-v0006}{6}{deletion}{\deleted{if}}{}{The author cancels this material by hand; it leaves the corrected running text.}
\psentry{ps_134-5_040-v0007}{7}{deletion}{\deleted{ing}}{}{The author cancels this material by hand; it leaves the corrected running text.}
\psentry{ps_134-5_040-v0008}{8}{writing over}{\over{h}{h}}{\hand{h}}{The superposed handwritten reading replaces the underlying reading.}
\psapparatusend
\fi
\pssource{ps_134-5_041}{134-5/41}{https://grothendieck.umontpellier.fr/web/viewer.html?file=../134-5.pdf\#page=42}
\psnumber{ps_134-5_041-v0001}{1}
\psnumber{ps_134-5_041-v0002}{2}
\psnumber{ps_134-5_041-v0003}{3}
%% BEGIN SOURCE ps_134-5_041
% Pursuing Stacks (cotes 134-2 to 134-8): archive cote 134-5, site page 41 (PDF page 42), typescript page 228.
% typescript. Draft compared: ps4.tex:1986-2051.
% AI-transcribed by Claude Opus 5.5 (claude-opus-5-5, adaptive thinking, effort high), Message Batches API, batch ps2 (msgbatch_01NVwr1Yuo8qJf5rykeGBXPK), 2026-11-12; reviewed against both supplied scan halves; see CHECK_DECISIONS.tsv.
\scanpage{134-5/41}
\mspage{228}

To finish these generalities on aspheric functors, I would like still
to make some comments on \uline{Asph}(A,B) , in case B admits binary products ,
with A and B otherwise arbitrary in (Cat). We are interested, for two
functors
\[ i,i' : A \hand{\rightrightarrows} B , \]
in the functor
\[ i\mathrm{x} i' : a \hand{\longmapsto} i(a)\mathrm{x} i'(a) \; ; \; A \hand{\longrightarrow} B , \]
which can be viewed indeed as a product object in \uline{Hom}(A,B) of i and i'.
Our interest here is mainly to give conditions ensuring that ixi' is
aspheric. It turns out that there is no point to this end to assume \hand{\uline{\textcolor{black}{both}}}
i and i' to be aspheric, what counts is that one, say i, should be asphe-%
ric. The point is made very clearly in the following

\hand{\uline{\textcolor{black}{Proposition}}} 4. Let $i: A \hand{\longrightarrow} B$ be an aspheric \var{\xout{functor, where B admits}}{}{typewriter deletion}{ps_134-5_041-v0001}
map in (Cat), we assume for simplicity that in B binary products exist.

a) Let $b_o \hand{\in} \mathrm{Ob}\, B$ , $i_{b_o} : A \hand{\to} B$ be the constant functor with value $b_o$,
consider the product functor
\[ i\mathrm{x} i_{b_o} : a \hand{\mapsto} i(a)\mathrm{x} b_o : A \hand{\to} B \quad . \]
In order for this functor to be aspheric, it is necessary and sufficient
that for any object y in B, the object $\underline{Hom}(b_o,y)$ in $B^\wedge$ be aspheric (a con-%
dition which depends only upon $b_o$ , not upon i nor even upon A).

\var{b) Assume this condition satisfied for any $b_o$ in B , \xout{\unsure{and such that for}} \inserted[typed]{Then for}
\hand{\uline{\textcolor{black}{any}}} functor $i': A \hand{\to} B$ \hand{,} \hand{(}\unsure{\over{a}{A}}ssume moreover A totally aspheric\over{.}{,}\hand{)}\hand{$\nwarrow$} \deleted{\xout{Then}}
ixi' is aspheric \xout{too} .}{b) Assume this condition satisfied for any $b_o$ in B , \unsure{\hand{A}}ssume moreover A totally aspheric\hand{,}  Then for
\hand{\uline{\textcolor{black}{any}}} functor $i': A \hand{\to} B$ \hand{,} 
ixi' is aspheric  .}{transposition}{ps_134-5_041-v0002}

\hand{\uline{\textcolor{black}{Proof}}}. a)Asphericity of $i\mathrm{x} i_{b_o}$ means that for any object y in B, the
corresponding \var{\xout{wh}}{}{typewriter deletion}{ps_134-5_041-v0003} presheaf on A
\[ (i\mathrm{x} i_{b_o})^*(y) = ( a \hand{\longmapsto} \mathrm{Hom}(i(a)\mathrm{x} b_o,y)) \]
is aspheric. Now the formula defining the presheaf $\underline{Hom}(b_o,y)$ on $B\unsure{*|\wedge}$
yields
\[ \mathrm{Hom}(i(a)\mathrm{x} b_o,y) \hand{\simeq} \mathrm{Hom}(i(a), \underline{Hom}(b_o,y)) \quad , \]
hence the presheaf we get on A is nothing but
\[ i^*(\underline{Hom}(b_o,y)) \quad . \]
As i is aspheric, the criterion (ii') of prop. 1 (page 214) implies that
this presheaf is aspheric iff $\underline{Hom}(b_o,y)$ is , qed.

b) We may view ixi' as a composition
\[ A \hand{\xrightarrow{\delta}} A\mathrm{x} A \overset{i\hand{\boxed{\textcolor{black}{\mathrm{x}}}} i'}{\hand{\longrightarrow}} B , \]
where the ``external product'' $i\hand{\boxed{\textcolor{black}{\mathrm{x}}}} i' : A\mathrm{x} A \hand{\to} B$ is defined by
\margin{41}

%% END SOURCE ps_134-5_041
\ifnum\EditionMode=1
\psapparatusstart
\psentry{ps_134-5_041-v0001}{1}{typewriter deletion}{\xout{functor, where B admits}}{}{The typewriter strikes out this material; it leaves the corrected running text.}
\psentry{ps_134-5_041-v0002}{2}{transposition}{b) Assume this condition satisfied for any $b_o$ in B , \xout{\unsure{and such that for}} \inserted[typed]{Then for}
\hand{\uline{\textcolor{black}{any}}} functor $i': A \hand{\to} B$ \hand{,} \hand{(}\unsure{\over{a}{A}}ssume moreover A totally aspheric\over{.}{,}\hand{)}\hand{$\nwarrow$} \deleted{\xout{Then}}
ixi' is aspheric \xout{too} .}{b) Assume this condition satisfied for any $b_o$ in B , \unsure{\hand{A}}ssume moreover A totally aspheric\hand{,}  Then for
\hand{\uline{\textcolor{black}{any}}} functor $i': A \hand{\to} B$ \hand{,} 
ixi' is aspheric  .}{The handwritten enclosure and upward locator move the total-asphericity hypothesis after in B in the previous line, before Then for; typed order stays in the diplomatic source.}
\psentry{ps_134-5_041-v0003}{3}{typewriter deletion}{\xout{wh}}{}{The typewriter strikes out this material; it leaves the corrected running text.}
\psapparatusend
\fi
\pssource{ps_134-5_042}{134-5/42}{https://grothendieck.umontpellier.fr/web/viewer.html?file=../134-5.pdf\#page=43}
\psnumber{ps_134-5_042-v0001}{1}
\psnumber{ps_134-5_042-v0002}{2}
\psnumber{ps_134-5_042-v0003}{3}
\psnumber{ps_134-5_042-v0004}{4}
\psnumber{ps_134-5_042-v0005}{5}
\psnumber{ps_134-5_042-v0006}{6}
\psnumber{ps_134-5_042-v0007}{7}
\psnumber{ps_134-5_042-v0008}{8}
\psnumber{ps_134-5_042-v0009}{9}
\psnumber{ps_134-5_042-v0010}{10}
\psnumber{ps_134-5_042-v0011}{11}
\psnumber{ps_134-5_042-v0012}{12}
\psnumber{ps_134-5_042-v0013}{13}
\psnumber{ps_134-5_042-v0014}{14}
\psnumber{ps_134-5_042-v0015}{15}
\psnumber{ps_134-5_042-v0016}{16}
\psnumber{ps_134-5_042-v0017}{17}
\psnumber{ps_134-5_042-v0018}{18}
%% BEGIN SOURCE ps_134-5_042
% Pursuing Stacks (cotes 134-2 to 134-8): archive cote 134-5, site page 42 (PDF page 43), typescript page 229.
% typescript. Draft compared: ps4.tex:2034-2102.
% AI-transcribed by Claude Opus 5.5 (claude-opus-5-5, adaptive thinking, effort high), Message Batches API, batch ps2 (msgbatch_01NVwr1Yuo8qJf5rykeGBXPK), 2026-11-12; reviewed against both supplied scan halves; see CHECK_DECISIONS.tsv.
\scanpage{134-5/42}
\mspage{229}
\[ i\hand{\boxtimes}i'\ (a,a') \overset{\mathrm{def}}{=} i(a)xi'(a') \quad . \]
By prop.3 we know  that the diagonal map  for A is aspheric,  thus we
are left with proving that i\hand{\ensuremath{\boxtimes}}i' is aspheric.  More generally,  we get:

\hand{\uline{\textcolor{black}{Corollary}}}\hand{1}. Let B in (Cat)  satisfy the conditions of a) and b) above,
and let  A,A' be two objects in (Cat),  and
\[ i\colon A \hand{\longrightarrow} B \quad , \quad i'\colon A'\hand{\longrightarrow} B \]
two maps with target B,  hence a map
\[ i\hand{\boxtimes}i' \ : \ (a,a')\hand{\mapsto} i(a)xi'(a') \ : \ AxA' \hand{\longrightarrow} B \quad . \]
If \var{\xout{either}}{}{typewriter deletion}{ps_134-5_042-v0001} i \var{\inserted[typed]{is}}{is}{typed insertion}{ps_134-5_042-v0002}\var{\xout{or $i'$ are}}{}{typewriter deletion}{ps_134-5_042-v0003} aspheric\hand{,} \var{\deleted{\inserted[typed]{and A' is aspheric} \xout{and $\unsure{i'}$ is aspheric}}}{}{deletion}{ps_134-5_042-v0004} then i\hand{\ensuremath{\boxtimes}}i' is
aspheric\hand{,} \var{\xout{iff A' is aspheric}}{}{typewriter deletion}{ps_134-5_042-v0005} \hand{iff A' is aspheric.}

We have to \var{\deleted{check}}{}{deletion}{ps_134-5_042-v0006} \var{\inserted{express}}{\hand{express}}{insertion}{ps_134-5_042-v0007} that for any object b in B,  the presheaf
\var{\xout{iffxA}}{}{typewriter deletion}{ps_134-5_042-v0008} \[ (a,a')\hand{\mapsto} \mathrm{Hom}(i(a)xi'(a'),b) \hand{\simeq} \mathrm{Hom}(i(a), \underline{\mathrm{Hom}}(i'(a'),b)) \]
on AxA' is aspheric.  For a' fixed in A',  the corresponding presheaf
on A is aspheric,  as we saw in a).  The conclusion now follows from the
useful

\hand{\uline{\textcolor{black}{Lemma}}}:  Let F be a presheaf on a product category AxA' ,  with A,A' in
(Cat).  As\unsure{s}ume that for any a' in A',  the corresponding presheaf
a$\hand{\mapsto}$F(a,a') on A be aspheric.  Then the composition
\[ (AxA')_{/F} \hand{\longrightarrow} AxA' \hand{\xrightarrow{\ \textcolor{black}{pr_2}\ }} A' \]
is in $\underline{W}$ ,  and hence F is aspheric iff A' is a\unsure{s}pheric.

Proof of lemma:  The composition is fibering, as both factors are.
This reminds us of the ``fibration condition'' L 5 \var{\inserted[typed]{on $\underline{W}$}}{on $\underline{W}$}{typed insertion}{ps_134-5_042-v0009}(page 164),  as yester-%
day  (page 219)\hand{,} wher\unsure{e} we stated that this condition follows from the
conditions Loc 1)  to Loc 3)  reviewed yesterday.  This condition asserts
that a fibration with aspheric fibers is in $\underline{W}$  -  hence the lemma.

\var{\inserted{\uline{Remarks.}}}{\hand{\uline{Remarks.}}}{insertion}{ps_134-5_042-v0010}
The results stated in prop.4  and it's corollary give a lot of elbow
freedom for getting new aspheric functors in terms of old ones,  with
target category B  -  provided B sa\var{\over{s}{t}}{\hand{t}}{writing over}{ps_134-5_042-v0011}isfies the two assumptions:  \var{\xout{existen\unsure{ce}}}{}{typewriter deletion}{ps_134-5_042-v0012}
stability under binary products (a \var{\inserted[typed]{property}}{property}{typed insertion}{ps_134-5_042-v0013}\var{\xout{condition}}{}{typewriter deletion}{ps_134-5_042-v0014} frequently met with,  al-%
though the standard test categories \hand{$\Delta$},\hand{$\Box$} and \hand{\unsure{$\bigcirc$}} lack it ...),  and the
\var{\xout{fact that}}{}{typewriter deletion}{ps_134-5_042-v0015} asphericity of the \var{\inserted[typed]{presheaves}}{presheaves}{typed insertion}{ps_134-5_042-v0016}\var{\xout{objects}}{}{typewriter deletion}{ps_134-5_042-v0017} $\underline{\mathrm{Hom}}$(b,y),  for any two objects b,y
in B.  A slightly stronger condition (indeed an equivalent one,  for
fixed b,  when b admits already a section over $e_{B\hat{}}$  and if $B\hat{}$ is totally
aspheric,  hence an aspheric object is \var{\inserted{even}}{\hand{even}}{insertion}{ps_134-5_042-v0018} aspheric over $e_{B\hat{}}$ ...) is \hand{\uline{\textcolor{black}{contrac-}}}%
\hand{\uline{\textcolor{black}{tibility}}} of the objects b of B ,  for the homotopy interval structure on
$B\hat{}$ defined by $\underline{W}_B$ \hand{(}i.e. in terms of homotopy intervals aspheric over $e_{B\hat{}}$\hand{)}.

\margin{42}

%% END SOURCE ps_134-5_042
\ifnum\EditionMode=1
\psapparatusstart
\psentry{ps_134-5_042-v0001}{1}{typewriter deletion}{\xout{either}}{}{The typewriter strikes out this material; it leaves the corrected running text.}
\psentry{ps_134-5_042-v0002}{2}{typed insertion}{\inserted[typed]{is}}{is}{The typed insertion is incorporated at its indicated position; insertion brackets leave the running text.}
\psentry{ps_134-5_042-v0003}{3}{typewriter deletion}{\xout{or $i'$ are}}{}{The typewriter strikes out this material; it leaves the corrected running text.}
\psentry{ps_134-5_042-v0004}{4}{deletion}{\deleted{\inserted[typed]{and A' is aspheric} \xout{and $\unsure{i'}$ is aspheric}}}{}{The author cancels this material by hand; it leaves the corrected running text.}
\psentry{ps_134-5_042-v0005}{5}{typewriter deletion}{\xout{iff A' is aspheric}}{}{The typewriter strikes out this material; it leaves the corrected running text.}
\psentry{ps_134-5_042-v0006}{6}{deletion}{\deleted{check}}{}{The author cancels this material by hand; it leaves the corrected running text.}
\psentry{ps_134-5_042-v0007}{7}{insertion}{\inserted{express}}{\hand{express}}{The handwritten insertion is incorporated at its indicated position; insertion brackets leave the running text.}
\psentry{ps_134-5_042-v0008}{8}{typewriter deletion}{\xout{iffxA}}{}{The typewriter strikes out this material; it leaves the corrected running text.}
\psentry{ps_134-5_042-v0009}{9}{typed insertion}{\inserted[typed]{on $\underline{W}$}}{on $\underline{W}$}{The typed insertion is incorporated at its indicated position; insertion brackets leave the running text.}
\psentry{ps_134-5_042-v0010}{10}{insertion}{\inserted{\uline{Remarks.}}}{\hand{\uline{Remarks.}}}{The handwritten insertion is incorporated at its indicated position; insertion brackets leave the running text.}
\psentry{ps_134-5_042-v0011}{11}{writing over}{\over{s}{t}}{\hand{t}}{The superposed handwritten reading replaces the underlying reading.}
\psentry{ps_134-5_042-v0012}{12}{typewriter deletion}{\xout{existen\unsure{ce}}}{}{The typewriter strikes out this material; it leaves the corrected running text.}
\psentry{ps_134-5_042-v0013}{13}{typed insertion}{\inserted[typed]{property}}{property}{The typed insertion is incorporated at its indicated position; insertion brackets leave the running text.}
\psentry{ps_134-5_042-v0014}{14}{typewriter deletion}{\xout{condition}}{}{The typewriter strikes out this material; it leaves the corrected running text.}
\psentry{ps_134-5_042-v0015}{15}{typewriter deletion}{\xout{fact that}}{}{The typewriter strikes out this material; it leaves the corrected running text.}
\psentry{ps_134-5_042-v0016}{16}{typed insertion}{\inserted[typed]{presheaves}}{presheaves}{The typed insertion is incorporated at its indicated position; insertion brackets leave the running text.}
\psentry{ps_134-5_042-v0017}{17}{typewriter deletion}{\xout{objects}}{}{The typewriter strikes out this material; it leaves the corrected running text.}
\psentry{ps_134-5_042-v0018}{18}{insertion}{\inserted{even}}{\hand{even}}{The handwritten insertion is incorporated at its indicated position; insertion brackets leave the running text.}
\psapparatusend
\fi
\pssource{ps_134-5_043}{134-5/43}{https://grothendieck.umontpellier.fr/web/viewer.html?file=../134-5.pdf\#page=44}
\psnumber{ps_134-5_043-v0001}{1}
\psnumber{ps_134-5_043-v0002}{2}
\psnumber{ps_134-5_043-v0003}{3}
\psnumber{ps_134-5_043-v0004}{4}
\psnumber{ps_134-5_043-v0005}{5}
\psnumber{ps_134-5_043-v0006}{6}
\psnumber{ps_134-5_043-v0007}{7}
\psnumber{ps_134-5_043-v0008}{8}
\psnumber{ps_134-5_043-v0009}{9}
\psnumber{ps_134-5_043-v0010}{10}
\psnumber{ps_134-5_043-v0011}{11}
\psnumber{ps_134-5_043-v0012}{12}
\psnumber{ps_134-5_043-v0013}{13}
\psnumber{ps_134-5_043-v0014}{14}
\psnumber{ps_134-5_043-v0015}{15}
\psnumber{ps_134-5_043-v0016}{16}
\psnumber{ps_134-5_043-v0017}{17}
\psnumber{ps_134-5_043-v0018}{18}
\psnumber{ps_134-5_043-v0019}{19}
\psnumber{ps_134-5_043-v0020}{20}
\psnumber{ps_134-5_043-v0021}{21}
\psnumber{ps_134-5_043-v0022}{22}
\psnumber{ps_134-5_043-v0023}{23}
\psnumber{ps_134-5_043-v0024}{24}
\psnumber{ps_134-5_043-v0025}{25}
\psnumber{ps_134-5_043-v0026}{26}
\psnumber{ps_134-5_043-v0027}{27}
\psnumber{ps_134-5_043-v0028}{28}
\psnumber{ps_134-5_043-v0029}{29}
%% BEGIN SOURCE ps_134-5_043
% Pursuing Stacks (cotes 134-2 to 134-8): archive cote 134-5, site page 43 (PDF page 44), typescript page 230.
% typescript. Draft compared: ps4.tex:2089-2156.
% AI-transcribed by Claude Opus 5.5 (claude-opus-5-5, adaptive thinking, effort high), Message Batches API, batch ps2 (msgbatch_01NVwr1Yuo8qJf5rykeGBXPK), 2026-11-12; reviewed against both supplied scan halves; see CHECK_DECISIONS.tsv.
\scanpage{134-5/43}
\mspage{230}

(Compare with propositions on pages 121 and 143\hand{.}) Whereas this \var{\inserted{latter}}{\hand{latter}}{insertion}{ps_134-5_043-v0001} assumption
admittedly is quite a stringent one, it is however of a type which has
become familiar to us in relation with test categories, where it seems a
rather common lot.

If B satisfies these conditions (as in prop.4) and if A is totally
aspheric, we see from prop.4 that the category $\underline{\mathrm{Asph}}(A,B)$ of aspheric
functors from A to B is stable under binary products. Now \var{\xout{itxisxeasily}}{}{typewriter deletion}{ps_134-5_043-v0002}
\var{\inserted[typed]{a non-empty object C of (Cat)}}{a non-empty object C of (Cat)}{typed insertion}{ps_134-5_043-v0003}
\var{\xout{\unsure{that when} a category C}}{}{typewriter deletion}{ps_134-5_043-v0004} stable under binary products gives rise to a
category $C\hat{}$ which is clearly totally aspheric for any basic localizer $\underline{W}$ ,
and in particular for the usual one \unsure{$\underline{W}_e$|$\underline{W}_0$} corresponding to usual weak equi-%
valence. A fortiori such a category C \var{\inserted[typed]{\xout{(ifxnonxempty)}}}{}{typed insertion}{ps_134-5_043-v0005} is 1-connected \hand{$\leftarrow$} \hand{(}which is easily
checked \var{\inserted{too}}{\hand{too}}{insertion}{ps_134-5_043-v0006} by down-to-earth direct arguments\hand{)} \var{\deleted{too}}{}{deletion}{ps_134-5_043-v0007}. Thus the reflections of
page 225 apply, and imply that if $\underline{\mathrm{Asph}}(A,B)$ is non-empty i.e. if there
is at least \hand{\uline{\textcolor{black}{one}}} aspheric functor $i\colon A\hand{\longrightarrow} B$, then there is a canonical
trans\hand{i}tive system of \unsure{\var{\over{si}{is}}{\hand{is}}{writing over}{ps_134-5_043-v0008}}omorphisms between all functors
\[ (\ast) \qquad \mathrm{Hot}_B \hand{\longrightarrow} \mathrm{Hot}_A \]
\margin{$\overline{i^*}$} of the type \var{\over{$i^*$}{$\overline{i^*}$}}{\hand{$\overline{i^*}$}}{writing over}{ps_134-5_043-v0009}, and hence there is a canonical functor $(\ast)$, defined up
to unique isomorphism, as announced in the remark on page 225.
\margin{For a generalization, see section 82; p. 272.}

13.6.

75. The generalities on aspheric maps of the last three sections should
be more than what is needed to develop now the notion of \hand{\uline{\textcolor{black}{asphericity
structure}}} - which, \unsure{\hand{,}}together with the \var{\inserted[typed]{closely related}}{closely related}{typed insertion}{ps_134-5_043-v0010} notion of \hand{\uline{\textcolor{black}{contractibility struc-%
ture}}}\hand{,} tentatively dealt with before, and the various \var{\inserted[typed]{``test-}}{``test-}{typed insertion}{ps_134-5_043-v0011}notions''\hand{(}\var{\over{\unsure{cf}}{e.g}}{\hand{e.g}}{writing over}{ps_134-5_043-v0012}. \hand{\uline{\textcolor{black}{test}}}
\var{\deleted{functors}}{}{deletion}{ps_134-5_043-v0013} \hand{\uline{\textcolor{black}{categories}}} and \hand{\uline{\textcolor{black}{test\var{\xout{x}}{}{typewriter deletion}{ps_134-5_043-v0014}functors}}}) seem\var{\deleted{s}}{}{deletion}{ps_134-5_043-v0015} to me \var{\deleted{at present}}{}{deletion}{ps_134-5_043-v0016} the main
pay off \var{\deleted{till now}}{}{deletion}{ps_134-5_043-v0017} \var{\inserted{so far}}{\hand{so far}}{insertion}{ps_134-5_043-v0018} of our effort to come to a grasp of a general formalism
of ``homotopy models''.

\var{\inserted{In the case of asphericity structures, just}}{\hand{In the case of asphericity structures, just}}{insertion}{ps_134-5_043-v0019} \hand{(}\unsure{\var{\over{A}{a}}{\hand{a}}{writing over}{ps_134-5_043-v0020}}s for the kindred notion of a contractibility structure, in all \var{\deleted{cases}}{}{deletion}{ps_134-5_043-v0021} \var{\inserted{instances}}{\hand{instances}}{insertion}{ps_134-5_043-v0022}
I could think of at present, the underlying category M of an asphericity
structure is \hand{\uline{\textcolor{black}{not}}} an object in (Cat) nor even a ``small category'', but is
``large'' - namely the cardinality of Ob M and F\var{\over{l}{l}}{\hand{l}}{writing over}{ps_134-5_043-v0023}(M) \var{\over{\unsure{is}}{are}}{\hand{are}}{writing over}{ps_134-5_043-v0024} not in the\hand{``}univer-%
se''\var{\inserted{$\underline{U}$}}{\hand{$\underline{U}$}}{insertion}{ps_134-5_043-v0025} we are working in, \var{\xout{whichxhasxxx}}{}{typewriter deletion}{ps_134-5_043-v0026} still less is M an object of $\underline{U}$ -
all we need\var{\inserted{instead}}{\hand{instead}}{insertion}{ps_134-5_043-v0027}, as usual, is that M be a \hand{\uline{\textcolor{black}{$\underline{U}$-category}}}, namely that for any
two objects of $\underline{U}$ , $\underline{\mathrm{Hom}}(x,y)$ be an element of $\underline{U}$ . Till now, the universe
$\underline{U}$ has been present in our reflections in a very much implicit way, in kee-%
ping with the unformal nature of the reflections, which however by and by
have become more formal\var{\deleted{,}}{}{deletion}{ps_134-5_043-v0028} \hand{(}as I finally let myself become invo\var{\over{v}{l}}{\hand{l}}{writing over}{ps_134-5_043-v0029}ved in
\margin{43}

%% END SOURCE ps_134-5_043
\ifnum\EditionMode=1
\psapparatusstart
\psentry{ps_134-5_043-v0001}{1}{insertion}{\inserted{latter}}{\hand{latter}}{The handwritten insertion is incorporated at its indicated position; insertion brackets leave the running text.}
\psentry{ps_134-5_043-v0002}{2}{typewriter deletion}{\xout{itxisxeasily}}{}{The typewriter strikes out this material; it leaves the corrected running text.}
\psentry{ps_134-5_043-v0003}{3}{typed insertion}{\inserted[typed]{a non-empty object C of (Cat)}}{a non-empty object C of (Cat)}{The typed insertion is incorporated at its indicated position; insertion brackets leave the running text.}
\psentry{ps_134-5_043-v0004}{4}{typewriter deletion}{\xout{\unsure{that when} a category C}}{}{The typewriter strikes out this material; it leaves the corrected running text.}
\psentry{ps_134-5_043-v0005}{5}{typed insertion}{\inserted[typed]{\xout{(ifxnonxempty)}}}{}{The typed insertion is incorporated at its indicated position; insertion brackets leave the running text.}
\psentry{ps_134-5_043-v0006}{6}{insertion}{\inserted{too}}{\hand{too}}{The handwritten insertion is incorporated at its indicated position; insertion brackets leave the running text.}
\psentry{ps_134-5_043-v0007}{7}{deletion}{\deleted{too}}{}{The author cancels this material by hand; it leaves the corrected running text.}
\psentry{ps_134-5_043-v0008}{8}{writing over}{\over{si}{is}}{\hand{is}}{The superposed handwritten reading replaces the underlying reading.}
\psentry{ps_134-5_043-v0009}{9}{writing over}{\over{$i^*$}{$\overline{i^*}$}}{\hand{$\overline{i^*}$}}{The superposed handwritten reading replaces the underlying reading.}
\psentry{ps_134-5_043-v0010}{10}{typed insertion}{\inserted[typed]{closely related}}{closely related}{The typed insertion is incorporated at its indicated position; insertion brackets leave the running text.}
\psentry{ps_134-5_043-v0011}{11}{typed insertion}{\inserted[typed]{``test-}}{``test-}{The typed insertion is incorporated at its indicated position; insertion brackets leave the running text.}
\psentry{ps_134-5_043-v0012}{12}{writing over}{\over{\unsure{cf}}{e.g}}{\hand{e.g}}{The superposed handwritten reading replaces the underlying reading.}
\psentry{ps_134-5_043-v0013}{13}{deletion}{\deleted{functors}}{}{The author cancels this material by hand; it leaves the corrected running text.}
\psentry{ps_134-5_043-v0014}{14}{typewriter deletion}{\xout{x}}{}{The typewriter strikes out this material; it leaves the corrected running text.}
\psentry{ps_134-5_043-v0015}{15}{deletion}{\deleted{s}}{}{The author cancels this material by hand; it leaves the corrected running text.}
\psentry{ps_134-5_043-v0016}{16}{deletion}{\deleted{at present}}{}{The author cancels this material by hand; it leaves the corrected running text.}
\psentry{ps_134-5_043-v0017}{17}{deletion}{\deleted{till now}}{}{The author cancels this material by hand; it leaves the corrected running text.}
\psentry{ps_134-5_043-v0018}{18}{insertion}{\inserted{so far}}{\hand{so far}}{The handwritten insertion is incorporated at its indicated position; insertion brackets leave the running text.}
\psentry{ps_134-5_043-v0019}{19}{insertion}{\inserted{In the case of asphericity structures, just}}{\hand{In the case of asphericity structures, just}}{The handwritten insertion is incorporated at its indicated position; insertion brackets leave the running text.}
\psentry{ps_134-5_043-v0020}{20}{writing over}{\over{A}{a}}{\hand{a}}{The superposed handwritten reading replaces the underlying reading.}
\psentry{ps_134-5_043-v0021}{21}{deletion}{\deleted{cases}}{}{The author cancels this material by hand; it leaves the corrected running text.}
\psentry{ps_134-5_043-v0022}{22}{insertion}{\inserted{instances}}{\hand{instances}}{The handwritten insertion is incorporated at its indicated position; insertion brackets leave the running text.}
\psentry{ps_134-5_043-v0023}{23}{writing over}{\over{l}{l}}{\hand{l}}{The superposed handwritten reading replaces the underlying reading.}
\psentry{ps_134-5_043-v0024}{24}{writing over}{\over{\unsure{is}}{are}}{\hand{are}}{The superposed handwritten reading replaces the underlying reading.}
\psentry{ps_134-5_043-v0025}{25}{insertion}{\inserted{$\underline{U}$}}{\hand{$\underline{U}$}}{The handwritten insertion is incorporated at its indicated position; insertion brackets leave the running text.}
\psentry{ps_134-5_043-v0026}{26}{typewriter deletion}{\xout{whichxhasxxx}}{}{The typewriter strikes out this material; it leaves the corrected running text.}
\psentry{ps_134-5_043-v0027}{27}{insertion}{\inserted{instead}}{\hand{instead}}{The handwritten insertion is incorporated at its indicated position; insertion brackets leave the running text.}
\psentry{ps_134-5_043-v0028}{28}{deletion}{\deleted{,}}{}{The author cancels this material by hand; it leaves the corrected running text.}
\psentry{ps_134-5_043-v0029}{29}{writing over}{\over{v}{l}}{\hand{l}}{The superposed handwritten reading replaces the underlying reading.}
\psapparatusend
\fi
\pssource{ps_134-5_044}{134-5/44}{https://grothendieck.umontpellier.fr/web/viewer.html?file=../134-5.pdf\#page=45}
\psnumber{ps_134-5_044-v0001}{1}
\psnumber{ps_134-5_044-v0002}{2}
\psnumber{ps_134-5_044-v0003}{3}
\psnumber{ps_134-5_044-v0004}{4}
\psnumber{ps_134-5_044-v0005}{5}
\psnumber{ps_134-5_044-v0006}{6}
\psnumber{ps_134-5_044-v0007}{7}
\psnumber{ps_134-5_044-v0008}{8}
%% BEGIN SOURCE ps_134-5_044
% Pursuing Stacks (cotes 134-2 to 134-8): archive cote 134-5, site page 44 (PDF page 45), typescript page 231.
% typescript. Draft compared: ps4.tex:2144-2221.
% AI-transcribed by Claude Opus 5.5 (claude-opus-5-5, adaptive thinking, effort high), Message Batches API, batch ps2 (msgbatch_01NVwr1Yuo8qJf5rykeGBXPK), 2026-11-12; reviewed against both supplied scan halves; see CHECK_DECISIONS.tsv.
\scanpage{134-5/44}
\mspage{231}
a mihimum of technical work,  needed for keeping out of the uneasiness
of ``thin air conjecturing''\hand{)}. An attentive reader will have felt occasiona%
-ly this implicit presence of $\underline{U}$,  for instance in the definition of the
basic modelizer (Cat) (which is,  as all modelizers,  a ``large'' category)\hand{,}
and in our occasional reference to various categories as being ``small'' or
``large''.  He will have noticed that whereas modelizers are by necessity
large categories (just as (Hot) itself,  whose set of isomorphism classes
of objects is large),  test categories are supposed to be small (and of-%
ten even to be in (Cat)\hand{)} -  with the effect that $A^{\wedge}$ ,  the category of
presheaves on A, is automatically a $\underline{U}$-category  (which would not be the
case if A was merely assumed to be a $\underline{U}$-category).

An \hand{\uline{\textcolor{black}{asphericity}}} \hand{\uline{\textcolor{black}{structure}}} (with re\var{\over{p}{\unsure{s}}}{\hand{\unsure{s}}}{writing over}{ps_134-5_044-v0001}pect to the badic localizer $\underline{W}$ )
on the $\underline{U}$-category M consists of a subset
\[ \text{(1)} \qquad M_{as} \hand{\subset} \text{Ob } M \quad , \]
whose elements will be called the \hand{\uline{\textcolor{black}{aspheric objects}}} of M  (more specifical%
ly,  the $\underline{W}$-aspheric objects,  if confusion may arise),  this subset being
su\var{\over{\unsure{m}}{b}}{\hand{b}}{writing over}{ps_134-5_044-v0002}mitted to the following condition:

% A lowered h-like mark overlaps the h of There; its relation to the main letter is unclear.
(Asstr)  T\unsure{h}ere exists an object A in (Cat),  and a functor $i: A \hand{\longrightarrow} M$ ,
satisfying the following two conditions:
(i)  For any a in A,  $i(a) \hand{\in} M_{as}$ ,  i.e.  i factors through the
full subcategory (also denoted by $M_{as}$) of M defined by $M_{as}$.
(ii)  Let
\[ i^{*}:\ M \hand{\longrightarrow} A^{\wedge} \quad , \quad x \hand{\mapsto} i^{*}(x) = (a \hand{\mapsto} \text{Hom}(i(a),x)) \ , \]
\var{\inserted[typed]{then we \xout{xatx}}}{then we }{typed insertion}{ps_134-5_044-v0003} \var{\xout{xxxxxxxxxxxx}}{}{typewriter deletion}{ps_134-5_044-v0004}  have
% In (2), unreadable obliterations lie above the hat and below the as subscript.
\[ \text{(2)} \qquad M_{as} = (i^{*})^{-1}(A^{\wedge\var{\deleted{\illegible{1}}}{}{deletion}{ps_134-5_044-v0005}}_{as\var{\deleted{\illegible{1}}}{}{deletion}{ps_134-5_044-v0006}}) \quad , \]
where $A^{\wedge}_{as}$ is the subset of Ob $A^{\wedge}$ of all ($\underline{W}$-)aspheric objects
of $A^{\wedge}$,  i.e. the \var{\xout{objects}}{}{typewriter deletion}{ps_134-5_044-v0007} presheaves F on A such that the
object $A_{/F}$  of (Cat) is $\underline{W}$-aspheric.

In other words, for x in M,  we have the equivalence
\[ \text{(2 bis)} \qquad x \hand{\in} M_{as} \hand{\Longleftrightarrow} i^{*}(x) \hand{\in} A^{\wedge}_{as} \ \text{ i.e. } \ A_{/x}\ (\overset{def}{=} A_{/i^{*}(x)}) \hand{\in} (\text{Cat})_{as} , \]
where $(\text{Cat})_{as}$ is the subset of Ob (Cat) made up with all $\underline{W}$-aspheric ob-%
jects of (Cat)  -  i.e. objects X such that  $X \hand{\to \Delta_{\bullet}}$ \hand{(}= final object\hand{)}
be in $\underline{W}$ .

Thus an asphericity structure on M can always be defined by a func-%
tor  \var{\xout{ixxix}}{}{typewriter deletion}{ps_134-5_044-v0008}
\[ i:\ A \hand{\longrightarrow} M \quad , \]
with A in (Cat),  by the formula  (2);  and conversely,  any such functor

\margin{44}

%% END SOURCE ps_134-5_044
\ifnum\EditionMode=1
\psapparatusstart
\psentry{ps_134-5_044-v0001}{1}{writing over}{\over{p}{\unsure{s}}}{\hand{\unsure{s}}}{The superposed handwritten reading replaces the underlying reading.}
\psentry{ps_134-5_044-v0002}{2}{writing over}{\over{\unsure{m}}{b}}{\hand{b}}{The superposed handwritten reading replaces the underlying reading.}
\psentry{ps_134-5_044-v0003}{3}{typed insertion}{\inserted[typed]{then we \xout{xatx}}}{then we }{The typed insertion is incorporated at its indicated position; insertion brackets leave the running text.}
\psentry{ps_134-5_044-v0004}{4}{typewriter deletion}{\xout{xxxxxxxxxxxx}}{}{The typewriter strikes out this material; it leaves the corrected running text.}
\psentry{ps_134-5_044-v0005}{5}{deletion}{$\deleted{\illegible{1}}$}{}{The author cancels this material by hand; it leaves the corrected running text.}
\psentry{ps_134-5_044-v0006}{6}{deletion}{$\deleted{\illegible{1}}$}{}{The author cancels this material by hand; it leaves the corrected running text.}
\psentry{ps_134-5_044-v0007}{7}{typewriter deletion}{\xout{objects}}{}{The typewriter strikes out this material; it leaves the corrected running text.}
\psentry{ps_134-5_044-v0008}{8}{typewriter deletion}{\xout{ixxix}}{}{The typewriter strikes out this material; it leaves the corrected running text.}
\psapparatusend
\fi
\pssource{ps_134-5_045}{134-5/45}{https://grothendieck.umontpellier.fr/web/viewer.html?file=../134-5.pdf\#page=46}
\psnumber{ps_134-5_045-v0001}{1}
\psnumber{ps_134-5_045-v0002}{2}
\psnumber{ps_134-5_045-v0003}{3}
\psnumber{ps_134-5_045-v0004}{4}
\psnumber{ps_134-5_045-v0005}{5}
\psnumber{ps_134-5_045-v0006}{6}
\psnumber{ps_134-5_045-v0007}{7}
\psnumber{ps_134-5_045-v0008}{8}
\psnumber{ps_134-5_045-v0009}{9}
\psnumber{ps_134-5_045-v0010}{10}
\psnumber{ps_134-5_045-v0011}{11}
\psnumber{ps_134-5_045-v0012}{12}
\psnumber{ps_134-5_045-v0013}{13}
\psnumber{ps_134-5_045-v0014}{14}
\psnumber{ps_134-5_045-v0015}{15}
\psnumber{ps_134-5_045-v0016}{16}
%% BEGIN SOURCE ps_134-5_045
% Pursuing Stacks (cotes 134-2 to 134-8): archive cote 134-5, site page 45 (PDF page 46), typescript page 232.
% typescript. Draft compared: ps4.tex:2202-2277.
% AI-transcribed by Claude Opus 5.5 (claude-opus-5-5, adaptive thinking, effort high), Message Batches API, batch ps2 (msgbatch_01NVwr1Yuo8qJf5rykeGBXPK), 2026-11-12; reviewed against both supplied scan halves; see CHECK_DECISIONS.tsv.
\scanpage{134-5/45}
\mspage{232}

defines an asphericity structure \var{\xout{xx}}{}{typewriter deletion}{ps_134-5_045-v0001} $M_{as}$ on $M$, admitting $i$ as a ``\hand{\uline{\textcolor{black}{testing}}}%
\margin{X}
\hand{\uline{\textcolor{black}{functor}}}'' (i.e. a functor satisfying (i) and (ii) above), provided only
we assume that

(3) For any $a$ in $A$, $i^{*}(i(a))$ is an aspheric object of $A^{\wedge}$ .

This latter condition is automatically satisfied if $i$ is fully faithful,
for instance, if it is the inclusion functor of a full subcategory $A$
of $M$. Thus, \hand{\uline{\textcolor{black}{any small full subcategory}}} $A$ \hand{\uline{\textcolor{black}{of}}} $M$ \hand{\uline{\textcolor{black}{defines an asphericity}}}
\hand{\uline{\textcolor{black}{structure on}}} $M$, and we'll see in a minute that any asphericity structure
on $M$ can be defined this way.

\hand{\uline{\textcolor{black}{Proposition}}} 1 Let $M_{as}$ be an asphericity structure on $M$, and $A,B$ ob-%
jects in (Cat), and
\[ A \hand{\xrightarrow{\textcolor{black}{f}}} B \hand{\xrightarrow{\textcolor{black}{j}}} M_{\var{\xout{xx}}{}{typewriter deletion}{ps_134-5_045-v0002}} \quad \var{\xout{x}}{}{typewriter deletion}{ps_134-5_045-v0003} \]
be functors, with $j$ factoring through $M_{as}$ .

a) If $f$ is aspheric, then j is a testing functor iff $i=jf$ is .

b) If $j$ is fully faithful, and if \var{\xout{bothxjxand}}{}{typewriter deletion}{ps_134-5_045-v0004} $i=jf$ \var{\xout{are}}{}{typewriter deletion}{ps_134-5_045-v0005}\var{\inserted[typed]{is a}}{is a}{typed insertion}{ps_134-5_045-v0006} testing func-%
tor\var{\xout{\unsure{s}}}{}{typewriter deletion}{ps_134-5_045-v0007}, then $f$ is aspheric, and $j$ is a testing functor too.

Proof. a) follows trivially from
\[ B^{\wedge}_{as} = (f^{*})^{-1}(A^{\wedge}_{as}) , \]
which is one of the ways of expressing that $f$ is as\var{\over{\unsure{s}}{p}}{\hand{p}}{writing over}{ps_134-5_045-v0008}heric\var{\xout{\unsure{,}}}{}{typewriter deletion}{ps_134-5_045-v0009} (criterion
(ii') of prop.1, p. 214). And the first assertion in b) is a trivial
consequence of the definition of testing functors, and of aspher\var{\xout{x}}{}{typewriter deletion}{ps_134-5_045-v0010}icity
of $f$ (by criterion (iii) on p. 214) -- and the second assertion of b)
now follows from a).

\hand{\uline{\textcolor{black}{Corollary}}} Let $i$: $A\hand{\longrightarrow} M$ be a testing functor for $(M,M_{as})$, and let
$B$ any \hand{\uline{\textcolor{black}{small}}} \var{\inserted[typed]{full}}{full}{typed insertion}{ps_134-5_045-v0011} subcategory of $M_{as}$ containing $i(A)$. Then the induced functor
$A\hand{\longrightarrow} B$ is aspheric, and the inclusion functor $B\hand{\hookrightarrow} M$ is a testing functor.

\hand{\uline{\textcolor{black}{R\var{\xout{m}}{}{typewriter deletion}{ps_134-5_045-v0012}emarks}}} This shows, as announced above, that any asphericity structu-%
re on $M$ can be defined by a small full subcategory of $M$ (for instance,
the smallest full subcategory of $M$ containing $i(A)$). We have been slight-%
ly floppy though, while we defined testin\var{\over{\unsure{g}}{g}}{\hand{g}}{writing over}{ps_134-5_045-v0013}\var{\xout{xxx}}{}{typewriter deletion}{ps_134-5_045-v0014} functors \var{\xout{xx}}{}{typewriter deletion}{ps_134-5_045-v0015} by insisting
that the source should be in (Cat) (which at times will be convenient),
whereas the $B$ we got here is merely small, namely \hand{\uline{\textcolor{black}{isomorphic}}} to an ob-%
ject of (Cat), but not necessarily in (Cat) itself. This visi\var{\xout{x}}{}{typewriter deletion}{ps_134-5_045-v0016}bly is an
``unessential floppyness'', which could be straightened out trivially, ei-%
ther by enlarging accordingly the notion of testing functor, or by sub-%
mitting $M$ to the (somewhat artificial, admittedly) restriction that all

\margin{45}

%% END SOURCE ps_134-5_045
\ifnum\EditionMode=1
\psapparatusstart
\psentry{ps_134-5_045-v0001}{1}{typewriter deletion}{\xout{xx}}{}{The typewriter strikes out this material; it leaves the corrected running text.}
\psentry{ps_134-5_045-v0002}{2}{typewriter deletion}{$\xout{xx}$}{}{The typewriter strikes out this material; it leaves the corrected running text.}
\psentry{ps_134-5_045-v0003}{3}{typewriter deletion}{$\xout{x}$}{}{The typewriter strikes out this material; it leaves the corrected running text.}
\psentry{ps_134-5_045-v0004}{4}{typewriter deletion}{\xout{bothxjxand}}{}{The typewriter strikes out this material; it leaves the corrected running text.}
\psentry{ps_134-5_045-v0005}{5}{typewriter deletion}{\xout{are}}{}{The typewriter strikes out this material; it leaves the corrected running text.}
\psentry{ps_134-5_045-v0006}{6}{typed insertion}{\inserted[typed]{is a}}{is a}{The typed insertion is incorporated at its indicated position; insertion brackets leave the running text.}
\psentry{ps_134-5_045-v0007}{7}{typewriter deletion}{\xout{\unsure{s}}}{}{The typewriter strikes out this material; it leaves the corrected running text.}
\psentry{ps_134-5_045-v0008}{8}{writing over}{\over{\unsure{s}}{p}}{\hand{p}}{The superposed handwritten reading replaces the underlying reading.}
\psentry{ps_134-5_045-v0009}{9}{typewriter deletion}{\xout{\unsure{,}}}{}{The typewriter strikes out this material; it leaves the corrected running text.}
\psentry{ps_134-5_045-v0010}{10}{typewriter deletion}{\xout{x}}{}{The typewriter strikes out this material; it leaves the corrected running text.}
\psentry{ps_134-5_045-v0011}{11}{typed insertion}{\inserted[typed]{full}}{full}{The typed insertion is incorporated at its indicated position; insertion brackets leave the running text.}
\psentry{ps_134-5_045-v0012}{12}{typewriter deletion}{\xout{m}}{}{The typewriter strikes out this material; it leaves the corrected running text.}
\psentry{ps_134-5_045-v0013}{13}{writing over}{\over{\unsure{g}}{g}}{\hand{g}}{The superposed handwritten reading replaces the underlying reading.}
\psentry{ps_134-5_045-v0014}{14}{typewriter deletion}{\xout{xxx}}{}{The typewriter strikes out this material; it leaves the corrected running text.}
\psentry{ps_134-5_045-v0015}{15}{typewriter deletion}{\xout{xx}}{}{The typewriter strikes out this material; it leaves the corrected running text.}
\psentry{ps_134-5_045-v0016}{16}{typewriter deletion}{\xout{x}}{}{The typewriter strikes out this material; it leaves the corrected running text.}
\psapparatusend
\fi
\pssource{ps_134-5_046}{134-5/46}{https://grothendieck.umontpellier.fr/web/viewer.html?file=../134-5.pdf\#page=47}
\psnumber{ps_134-5_046-v0001}{1}
\psnumber{ps_134-5_046-v0002}{2}
\psnumber{ps_134-5_046-v0003}{3}
\psnumber{ps_134-5_046-v0004}{4}
\psnumber{ps_134-5_046-v0005}{5}
\psnumber{ps_134-5_046-v0006}{6}
\psnumber{ps_134-5_046-v0007}{7}
%% BEGIN SOURCE ps_134-5_046
% Pursuing Stacks (cotes 134-2 to 134-8): archive cote 134-5, site page 46 (PDF page 47), typescript page 233.
% typescript. Draft compared: ps4.tex:2262-2327.
% AI-transcribed by Claude Opus 5.5 (claude-opus-5-5, adaptive thinking, effort high), Message Batches API, batch ps2 (msgbatch_01NVwr1Yuo8qJf5rykeGBXPK), 2026-11-12; reviewed against both supplied scan halves; see CHECK_DECISIONS.tsv.
\scanpage{134-5/46}
\mspage{233}
it's small subcategories should be in (Cat). This presumably will be sa-%
tisfied by most large categories we are going to consider, and it should\hand{\unsure{'|)}}
nt be hard moreover to show that any $\underline{U}$-category is isomorphic to a cate-%
gory $M'$ satisfying the above extra assumption.

A little more serious maybe is the use we are making here of the
name of a ``testing functor'', which seems to be conflicting with an ear-%
lier use (def.5, p. 175\var{\deleted{\xout{)xxx}}}{}{deletion}{ps_134-5_046-v0001} and def.6, p.176), where we insisted \var{\inserted[typed]{(for instance)}}{(for instance)}{typed insertion}{ps_134-5_046-v0002} that
the source $A$ should be a test category. That's why I have been using
here the name ``testing functor'' rather than ``test functor'', to be on
the safe side formally speaking - but this is still playing on words,
namely cheating a little. Maybe the name ``\hand{\uline{\textcolor{black}{aspherical functor}}}'' instead
of``testing functor'' would be less misleading, thinking of the case when
$M$ is small itself, and endowing $M$ with the canonical asphericity struc-%
ture, for which
\[ M_{as} = M \quad , \]
(admitting the identity functor as a testing functor) - in which case
the``testing functors'' $A \hand{\to} M$ are indeed just the aspherical functors.
The drawback is that when $M$ is small, \var{\xout{there}}{}{typewriter deletion}{ps_134-5_046-v0003} it may well be endowed with
an asphericity structure $M_{as}$ different from the previous one, in which
case the proposed extension of the name ``aspheric functor'' again leads to
an ambiguity, unless specified by ``$M_{as}$-\hand{\uline{\textcolor{black}{aspheric}}}''(where, after all,
$M_{as}$ could be \hand{\uline{\textcolor{black}{any}}} full subcategory of $M$) - but then the notion reduces
to the one of an aspheric functor $A \hand{\to} M_{as}$ . But the same after all
holds even for large $M$ - the notion of a ``testing functor'' $A \hand{\to} M$
(with respect to an asphericity structure $M_{as}$ on $M$) does not really de-%
pend on the \hand{\uline{\textcolor{black}{pair}}} $(M,M_{as})$, but rather on the (\var{\xout{generally}}{}{typewriter deletion}{ps_134-5_046-v0004}\var{\inserted[typed]{possibly}}{possibly}{typed insertion}{ps_134-5_046-v0005} large) category
$M_{as}$ itself - namely it is no more no less than a functor
\[ A \hand{\to} M_{as} \]
which satisfies the usual asphericity condition (iii) (of prop.1 p.214)\hand{,}
with the only difference that $M_{as}$ may not be small, and therefore \var{\xout{we}}{}{typewriter deletion}{ps_134-5_046-v0006}
$M_{as}\hat{}$ may not be a $\underline{U}$-category (and we are \var{\inserted{therefore}}{\hand{therefore}}{insertion}{ps_134-5_046-v0007} reluctant to work with this
latter category at all, unless we first pass to the next larger universe
$\underline{U}'$ ...).

This short reflection rather convinces me that the designation of
``testing functors'' as introduced on the page before, by the alternative
name of $M_{as}$-\hand{\uline{\textcolor{black}{aspheric functors}}}, or just \hand{\uline{\textcolor{black}{aspheric functors}}} when no confu-%
sion is likely to arise as to existence and choice of $M_{as}$, is satisfac-%
tory indeed. I'll use it tentatively, as a synonyme to ``testing func-%
tor'', and it will appear soon enough if this terminology is a good one
\margin{46}

%% END SOURCE ps_134-5_046
\ifnum\EditionMode=1
\psapparatusstart
\psentry{ps_134-5_046-v0001}{1}{deletion}{\deleted{\xout{)xxx}}}{}{The author cancels this material by hand; it leaves the corrected running text.}
\psentry{ps_134-5_046-v0002}{2}{typed insertion}{\inserted[typed]{(for instance)}}{(for instance)}{The typed insertion is incorporated at its indicated position; insertion brackets leave the running text.}
\psentry{ps_134-5_046-v0003}{3}{typewriter deletion}{\xout{there}}{}{The typewriter strikes out this material; it leaves the corrected running text.}
\psentry{ps_134-5_046-v0004}{4}{typewriter deletion}{\xout{generally}}{}{The typewriter strikes out this material; it leaves the corrected running text.}
\psentry{ps_134-5_046-v0005}{5}{typed insertion}{\inserted[typed]{possibly}}{possibly}{The typed insertion is incorporated at its indicated position; insertion brackets leave the running text.}
\psentry{ps_134-5_046-v0006}{6}{typewriter deletion}{\xout{we}}{}{The typewriter strikes out this material; it leaves the corrected running text.}
\psentry{ps_134-5_046-v0007}{7}{insertion}{\inserted{therefore}}{\hand{therefore}}{The handwritten insertion is incorporated at its indicated position; insertion brackets leave the running text.}
\psapparatusend
\fi
\pssource{ps_134-5_047}{134-5/47}{https://grothendieck.umontpellier.fr/web/viewer.html?file=../134-5.pdf\#page=48}
\psnumber{ps_134-5_047-v0001}{1}
\psnumber{ps_134-5_047-v0002}{2}
\psnumber{ps_134-5_047-v0003}{3}
\psnumber{ps_134-5_047-v0004}{4}
\psnumber{ps_134-5_047-v0005}{5}
\psnumber{ps_134-5_047-v0006}{6}
\psnumber{ps_134-5_047-v0007}{7}
\psnumber{ps_134-5_047-v0008}{8}
\psnumber{ps_134-5_047-v0009}{9}
\psnumber{ps_134-5_047-v0010}{10}
\psnumber{ps_134-5_047-v0011}{11}
\psnumber{ps_134-5_047-v0012}{12}
%% BEGIN SOURCE ps_134-5_047
% Pursuing Stacks (cotes 134-2 to 134-8): archive cote 134-5, site page 47 (PDF page 48), typescript page 234.
% typescript. Draft compared: ps4.tex:2310-2381.
% AI-transcribed by Claude Opus 5.5 (claude-opus-5-5, adaptive thinking, effort high), Message Batches API, batch ps2 (msgbatch_01NVwr1Yuo8qJf5rykeGBXPK), 2026-11-12; reviewed against both supplied scan halves; see CHECK_DECISIONS.tsv.
\scanpage{134-5/47}
\mspage{234}

or not  -  namely if it is suggestive,  and not too conducive to misunder-%
standings.

\hand{\uline{\textcolor{black}{Proposition}}} 2 .  Let $(M,M_{\mathrm{as}})$ be an asphericity structure,  and let
\[ i\colon A \hand{\to} M \]
be an $M_{\mathrm{as}}$-aspheric functor,  hence $i^{*}\colon M \hand{\to} A\hat{}$.  Let
\[ \text{(4)} \qquad W =(i^{*})^{-1}(\underline{W}_A) \hand{\subset} \mathrm{Fl}(M) \quad . \]
Then M is a \var{\inserted[typed]{mildly}}{mildly}{typed insertion}{ps_134-5_047-v0001} saturated subset of $\mathrm{Fl}(M)$,  \hand{\uline{\textcolor{black}{independant}}} \hand{\uline{\textcolor{black}{of the choice of}}} $(\var{\over{M}{A}}{\hand{A}}{writing over}{ps_134-5_047-v0002},i)$.

Mild saturation of W follows from mild saturation of $\underline{W}_A$ ,  the set
of $\underline{W}$-weak equivalences in $A\hat{}$ (which follows from mild saturation of $\underline{W}$ it-%
self and the defi nition
\[ \underline{W}_A \;=\; i_A^{\;-1}(\underline{W}) \quad ). \]
To prove that $W_i$ for $(A,i)$ is the same as $W^{\var{\xout{x}}{}{typewriter deletion}{ps_134-5_047-v0003}}_{i'}$ for $(A',i')$,  choose a small
\var{\inserted[typed]{full}}{full}{typed insertion}{ps_134-5_047-v0004} subcategory B of $M_{\mathrm{as}}$ such that i and i' factor through  B,  let $j\colon B \hand{\to} M$
be the inclusion,  we only have to check $W_i = W_j$  (and similatrly for $W_{i'}$),
which follows from  $i^{*} = \var{\over{\unsure{f}}{f}}{\hand{f}}{writing over}{ps_134-5_047-v0005}^{*}j^{*}$  (where $f\colon A \hand{\to} B$ is the induced functor)
and the relation
\[ \underline{W}_B = \var{\xout{f^{*}(\underline{W}_A)}}{}{typewriter deletion}{ps_134-5_047-v0006}\ (f^{*})^{-1}(\underline{W}_A) \quad , \]
which follows from f being aspheric by prop.1 b),  and the known property
(iv) (prop.1\hand{,} p. 214) of asphericity,  qed.

\hand{\uline{\textcolor{black}{Remark}\var{\deleted{\unsure{:}}}{}{deletion}{ps_134-5_047-v0007}1}}  Once we prove that $\underline{W}$  is even strongly saturated,  it will follo\var{\over{\unsure{w}}{w}}{\hand{w}}{writing over}{ps_134-5_047-v0008}
of course that the sets $\underline{W}_A$ ,  and hence W above,  are strongly saturated
too.

We'll call W the set of \hand{\uline{\textcolor{black}{weak}}} \hand{\uline{\textcolor{black}{equivalences}}} in M,  for the given asphe-%
ricity structure $M_{\mathrm{as}}$ in M.  We are interested \var{\inserted{now}}{\hand{now}}{insertion}{ps_134-5_047-v0009} in giving a condition on $M_{\mathrm{as}}$
ensuring that \unsure{conversely|Conversely},  $M_{\mathrm{as}}$ is known when the corresponding set W of
weak equivalences is.  We'll \var{\xout{say that (M}}{}{typewriter deletion}{ps_134-5_047-v0010} assume for this that M has a
final object $e_M$ ,  and we\hand{'}ll call the asphericity structure $(M,M_{\mathrm{as}})$
``\hand{\uline{\textcolor{black}{aspheric}}}''  if  $e_M$ is aspheric i.e.  $e_M \hand{\in} M_{\mathrm{as}}$  (which does not depend of
course on the choice of $e_M$ ).  Now we get the tautology:

\hand{\uline{\textcolor{black}{Proposition}}} 3.   Let $(M,M\var{\xout{a\unsure{s}}}{}{typewriter deletion}{ps_134-5_047-v0011}_{\mathrm{as}})$ be an asphericity structure,  and  $u\colon x \hand{\to} y$
a map in M.  If y is aspheric,  then x is aspheric iff \var{\over[typed]{\unsure{x}}{u}}{u}{typewriter overwriting}{ps_134-5_047-v0012} is aspheric.

\hand{\uline{\textcolor{black}{Corollary}}}  Assume M admits a final object $e_M$ ,  and that $e_M$ is aspheric
(i.e. $(M,M_{\mathrm{as}})$ is aspheric).  Then an object x in M is aspheric iff the
map $x \hand{\to} e_M$ is aspheric.

Next question then is to state the conditions on a pair (M,W), with
$W \hand{\subset} \mathrm{Fl}(M)$,  for the existence of an aspheric asphericity structure on M,
\margin{47}

%% END SOURCE ps_134-5_047
\ifnum\EditionMode=1
\psapparatusstart
\psentry{ps_134-5_047-v0001}{1}{typed insertion}{\inserted[typed]{mildly}}{mildly}{The typed insertion is incorporated at its indicated position; insertion brackets leave the running text.}
\psentry{ps_134-5_047-v0002}{2}{writing over}{$\over{M}{A}$}{$\hand{A}$}{The superposed handwritten reading replaces the underlying reading.}
\psentry{ps_134-5_047-v0003}{3}{typewriter deletion}{$\xout{x}$}{}{The typewriter strikes out this material; it leaves the corrected running text.}
\psentry{ps_134-5_047-v0004}{4}{typed insertion}{\inserted[typed]{full}}{full}{The typed insertion is incorporated at its indicated position; insertion brackets leave the running text.}
\psentry{ps_134-5_047-v0005}{5}{writing over}{$\over{\unsure{f}}{f}$}{$\hand{f}$}{The superposed handwritten reading replaces the underlying reading.}
\psentry{ps_134-5_047-v0006}{6}{typewriter deletion}{$\xout{f^{*}(\underline{W}_A)}$}{}{The typewriter strikes out this material; it leaves the corrected running text.}
\psentry{ps_134-5_047-v0007}{7}{deletion}{\deleted{\unsure{:}}}{}{The author cancels this material by hand; it leaves the corrected running text.}
\psentry{ps_134-5_047-v0008}{8}{writing over}{\over{\unsure{w}}{w}}{\hand{w}}{The superposed handwritten reading replaces the underlying reading.}
\psentry{ps_134-5_047-v0009}{9}{insertion}{\inserted{now}}{\hand{now}}{The handwritten insertion is incorporated at its indicated position; insertion brackets leave the running text.}
\psentry{ps_134-5_047-v0010}{10}{typewriter deletion}{\xout{say that (M}}{}{The typewriter strikes out this material; it leaves the corrected running text.}
\psentry{ps_134-5_047-v0011}{11}{typewriter deletion}{$\xout{a\unsure{s}}$}{}{The typewriter strikes out this material; it leaves the corrected running text.}
\psentry{ps_134-5_047-v0012}{12}{typewriter overwriting}{\over[typed]{\unsure{x}}{u}}{u}{The superposed typed reading replaces the underlying reading.}
\psapparatusend
\fi
\pssource{ps_134-5_048}{134-5/48}{https://grothendieck.umontpellier.fr/web/viewer.html?file=../134-5.pdf\#page=49}
\psnumber{ps_134-5_048-v0001}{1}
\psnumber{ps_134-5_048-v0002}{2}
\psnumber{ps_134-5_048-v0003}{3}
\psnumber{ps_134-5_048-v0004}{4}
\psnumber{ps_134-5_048-v0005}{5}
\psnumber{ps_134-5_048-v0006}{6}
\psnumber{ps_134-5_048-v0007}{7}
\psnumber{ps_134-5_048-v0008}{8}
\psnumber{ps_134-5_048-v0009}{9}
\psnumber{ps_134-5_048-v0010}{10}
%% BEGIN SOURCE ps_134-5_048
% Pursuing Stacks (cotes 134-2 to 134-8): archive cote 134-5, site page 48 (PDF page 49), typescript page 235.
% typescript. Draft compared: ps4.tex:2363-2433.
% AI-transcribed by Claude Opus 5.5 (claude-opus-5-5, adaptive thinking, effort high), Message Batches API, batch ps2 (msgbatch_01NVwr1Yuo8qJf5rykeGBXPK), 2026-11-12; reviewed against both supplied scan halves; see CHECK_DECISIONS.tsv.
\scanpage{134-5/48}
\mspage{235}

admitting W as it's set of weak equivalences.  We assume beforehand M ad-%
mits a final object $e_M$ .  A ns condition is the \unsure{existence} of a small cate-%
gory A and a functor
\[ i: A \hand{\longrightarrow} M \]
satisfying the following two conditions:

(i) for x in M of the form i(a) ($a \hand{\in} \mathrm{Ob}\, A$) \hand{\uline{\textcolor{black}{or}}} $e_M$ , $i^{*}(x) \hand{\in} A^{\wedge}_{as}$ ,

(ii) $W = (i^{*})^{-1}(\underline{W}_A)$ .

Another ns condition, which looks more pleasant I guess, is that there
exist a small full subcategory B of M, containing $e_M$ , with inclusion
functor $j: B \hand{\to} M$, such that
\[ (5) \qquad W = (j^{*})^{-1}(\var{\xout{M}}{}{typewriter deletion}{ps_134-5_048-v0001}\ \underline{W}_B) \quad . \]
In any case, if we want a ns condition on W for it to be the set of
weak equivalences for some asphericity structure on M (not necessarily
an aspheric one, and therefore maybe not unique), we get: there should
exist a small full subcategory B of M, with inclusion functor j, such
that (5) above holds. (Here we do not assume $e_M$ to exist, still less B
to contain it) Similarly, the corollary of prop.1 implies that a sub-%
set $M_{as}$ of Ob M is an asphericity structure on M, iff there exists a
small full subcategory B in M, such that
\[ (6) \qquad M_{as} = (j^{*})^{-1}(B^{\wedge}_{as}) \quad ; \]
moreover, if M admits a final object $e_M$ , the asphericity structure $M_{as}$
is aspheric iff B can be chosen to contain $e_M$ .

\hand{\uline{\textcolor{black}{Remarks}}} 2.  The relationship between aspheric asphericity struct\var{\xout{\unsure{m|u}}}{}{typewriter deletion}{ps_134-5_048-v0002}ures
$M_{as}$ on M, and sets $W \hand{\subset} \mathrm{F}\var{\over{\mathrm{l}}{\mathrm{l}}}{\hand{\mathrm{l}}}{writing over}{ps_134-5_048-v0003}(M)$ of ``weak equivalences'' in M \unsure{\var{\over[typed]{(}{s}}{s}{typewriter overwriting}{ps_134-5_048-v0004}}atisfying
the condition above (and which we may call ``weak equivalence structures''
on M), in case M admits a final object $e_M$ , is reminiscent of the rela-%
tionship between \var{\xout{\unsure{thomste}}}{}{typewriter deletion}{ps_134-5_048-v0005} ``contractibility structures'' $M_{c\var{\deleted{\unsure{ont}}}{}{deletion}{ps_134-5_048-v0006}} \hand{\subset} \mathrm{Ob}\, M$ on M,
and those \var{\xout{h}}{}{typewriter deletion}{ps_134-5_048-v0007}``homotopism structures'' \var{\xout{on M}}{}{typewriter deletion}{ps_134-5_048-v0008} $h_M \hand{\subset} \mathrm{F}\var{\over{\mathrm{l}}{\mathrm{l}}}{\hand{\mathrm{l}}}{writing over}{ps_134-5_048-v0009}(M)$ on M which can be
described in terms of such a contractibility structure (cf. sections 51
and 52). Both pairs can be viewed as giving two equivalent ways of
expressing one and the same kind of structure - the structure concerned
by the first pair being centered on asphericity notions, whereas the
second is concerned with typical homotopy notions rather. We'll see
later that any contractibility structure defines in an evident way an
\margin{x} asphericity structure, and in the most interesting cases (e.g. canonical
\var{\xout{contractib\unsure{l}}}{}{typewriter deletion}{ps_134-5_048-v0010} modelizers), it is uniquely determined by the latter.

\margin{48}

%% END SOURCE ps_134-5_048
\ifnum\EditionMode=1
\psapparatusstart
\psentry{ps_134-5_048-v0001}{1}{typewriter deletion}{$\xout{M}$}{}{The typewriter strikes out this material; it leaves the corrected running text.}
\psentry{ps_134-5_048-v0002}{2}{typewriter deletion}{\xout{\unsure{m|u}}}{}{The typewriter strikes out this material; it leaves the corrected running text.}
\psentry{ps_134-5_048-v0003}{3}{writing over}{$\over{\mathrm{l}}{\mathrm{l}}$}{$\hand{\mathrm{l}}$}{The superposed handwritten reading replaces the underlying reading.}
\psentry{ps_134-5_048-v0004}{4}{typewriter overwriting}{\over[typed]{(}{s}}{s}{The superposed typed reading replaces the underlying reading.}
\psentry{ps_134-5_048-v0005}{5}{typewriter deletion}{\xout{\unsure{thomste}}}{}{The typewriter strikes out this material; it leaves the corrected running text.}
\psentry{ps_134-5_048-v0006}{6}{deletion}{$\deleted{\unsure{ont}}$}{}{The author cancels this material by hand; it leaves the corrected running text.}
\psentry{ps_134-5_048-v0007}{7}{typewriter deletion}{\xout{h}}{}{The typewriter strikes out this material; it leaves the corrected running text.}
\psentry{ps_134-5_048-v0008}{8}{typewriter deletion}{\xout{on M}}{}{The typewriter strikes out this material; it leaves the corrected running text.}
\psentry{ps_134-5_048-v0009}{9}{writing over}{$\over{\mathrm{l}}{\mathrm{l}}$}{$\hand{\mathrm{l}}$}{The superposed handwritten reading replaces the underlying reading.}
\psentry{ps_134-5_048-v0010}{10}{typewriter deletion}{\xout{contractib\unsure{l}}}{}{The typewriter strikes out this material; it leaves the corrected running text.}
\psapparatusend
\fi
\pssource{ps_134-5_049}{134-5/49}{https://grothendieck.umontpellier.fr/web/viewer.html?file=../134-5.pdf\#page=50}
\psnumber{ps_134-5_049-v0001}{1}
\psnumber{ps_134-5_049-v0002}{2}
\psnumber{ps_134-5_049-v0003}{3}
\psnumber{ps_134-5_049-v0004}{4}
\psnumber{ps_134-5_049-v0005}{5}
\psnumber{ps_134-5_049-v0006}{6}
\psnumber{ps_134-5_049-v0007}{7}
\psnumber{ps_134-5_049-v0008}{8}
\psnumber{ps_134-5_049-v0009}{9}
\psnumber{ps_134-5_049-v0010}{10}
\psnumber{ps_134-5_049-v0011}{11}
\psnumber{ps_134-5_049-v0012}{12}
\psnumber{ps_134-5_049-v0013}{13}
\psnumber{ps_134-5_049-v0014}{14}
\psnumber{ps_134-5_049-v0015}{15}
\psnumber{ps_134-5_049-v0016}{16}
\psnumber{ps_134-5_049-v0017}{17}
\psnumber{ps_134-5_049-v0018}{18}
\psnumber{ps_134-5_049-v0019}{19}
\psnumber{ps_134-5_049-v0020}{20}
\psnumber{ps_134-5_049-v0021}{21}
\psnumber{ps_134-5_049-v0022}{22}
\psnumber{ps_134-5_049-v0023}{23}
\psnumber{ps_134-5_049-v0024}{24}
\psnumber{ps_134-5_049-v0025}{25}
\psnumber{ps_134-5_049-v0026}{26}
\psnumber{ps_134-5_049-v0027}{27}
%% BEGIN SOURCE ps_134-5_049
% Pursuing Stacks (cotes 134-2 to 134-8): archive cote 134-5, site page 49 (PDF page 50), typescript page 236.
% typescript. Draft compared: ps4.tex:2418-2499.
% AI-transcribed by Claude Opus 5.5 (claude-opus-5-5, adaptive thinking, effort high), Message Batches API, batch ps2 (msgbatch_01NVwr1Yuo8qJf5rykeGBXPK), 2026-11-12; reviewed against both supplied scan halves; see CHECK_DECISIONS.tsv.
\scanpage{134-5/49}
\mspage{236}

To sum up some of the main relationships between the three asphericity
notions just introduced ($M_{as}$ , $W$, $M_{as}$-aspheric functors), \var{\over{L}{l}}{\hand{l}}{writing over}{ps_134-5_049-v0001}et's state
one more tautological proposition, which is very much a paraphrase of
the display given earlier (prop. on p. 214) of the manifold aspects of
the notion of an aspheric map between small categories:

\hand{\uline{\textcolor{black}{Proposition}}} 4. Let $(M,M_{as})$ be an asphericity structure, $A$ a small
category, $i\colon A\hand{\to} M$ a functor, factoring through $M_{as}$ . Consider the
following conditions on $i$ :

(i) $i$ is $M_{as}$-aspheric (or \var{\over[typed]{\illegible{1}}{a}}{a}{typewriter overwriting}{ps_134-5_049-v0002} ``testing functor'' for $M_{as}$), i.e. \var{\xout{forx}}{}{typewriter deletion}{ps_134-5_049-v0003}
\[ M_{as} = (i^{*})^{-1}(A^{\wedge}_{as}) \qquad . \]
(i') For $x$ in $M_{as}$ , $i^{*}(x)$ is aspheric, i.e.
\[ M_{as} \hand{\subset} (i^{*})^{-1}(A^{\wedge}_{as}) \quad . \]
(i'') (Here $B$ is a given small full subcategory \var{\inserted[typed]{of $M_{as}$}}{of $M_{as}$}{typed insertion}{ps_134-5_049-v0004} \var{\inserted{containing $i(A)$ and}}{\hand{containing $i(A)$ and}}{insertion}{ps_134-5_049-v0005} which``\hand{\uline{\textcolor{black}{generates}}}'' the
asphericity structure $M_{as}$ , namely such that the inclusion functor
$j\colon B\hand{\hookrightarrow} M$ is $M_{as}$-aspheric i.e. $M_{as} = (j^{*})^{-1}(B^{\wedge}_{as})$.) For any $x$ in $B$,
$i^{*}(x)$ is aspheric, i.e.
\[ B \hand{\subset} (i^{*})^{-1}(A^{\wedge}_{as}) \quad . \]
(ii) For any map $u$ in $M$, $u$ is \var{\xout{aspheric}}{}{typewriter deletion}{ps_134-5_049-v0006} \var{\inserted[typed]{a weak equivalence}}{a weak equivalence}{typed insertion}{ps_134-5_049-v0007} iff $i^{*}(u)$ is, i.e.
\[ W = (i^{*})^{-1}(\underline{W}_{A}) \quad . \]
(ii') If the map $u$ in $M$ is \var{\xout{aspheric,}}{}{typewriter deletion}{ps_134-5_049-v0008} \var{\inserted[typed]{a weak equivalence}}{a weak equivalence}{typed insertion}{ps_134-5_049-v0009} so is $i^{*}(u)$ , i.e.
\[ W \hand{\subset} (i^{*})^{-1}(\underline{W}_{A}) \quad . \]
(ii'') (Here $B$ is given as in (i'') above) For any map \unsure{u} in $B$, $i^{*}(u)$
is a weak equivalence, \var{\inserted{i.e. $Fl(B) \hand{\subset} (i^{*})^{-1}(\underline{W}_{A})$}}{\hand{i.e. $Fl(B) \hand{\subset} (i^{*})^{-1}(\underline{W}_{A})$}}{insertion}{ps_134-5_049-v0010}. \margin{——}

% A small upright mark above The remains unidentified; the margin strokes are retained.
\margin{——} The conditions (i)(i')(i'') are equivalent and imply all others, and
we have the tautological implications (ii)$\hand{\Rightarrow}$(ii')$\hand{\Rightarrow}$(ii''). If $M$ admits
a final object $e_M$ and \var{\inserted[typed]{if A and}}{if A and}{typed insertion}{ps_134-5_049-v0011} the asphericity structure \var{\inserted{$M_{as}$}}{\hand{$M_{as}$}}{insertion}{ps_134-5_049-v0012} \var{\deleted{is}}{}{deletion}{ps_134-5_049-v0013} \var{\inserted[typed]{are}}{are}{typed insertion}{ps_134-5_049-v0014} aspheric, then \var{\xout{all}}{}{typewriter deletion}{ps_134-5_049-v0015} \var{\inserted{all}}{\hand{all}}{insertion}{ps_134-5_049-v0016}
\var{\inserted[typed]{six}}{six}{typed insertion}{ps_134-5_049-v0017} \var{\xout{\illegible{1}}}{}{typewriter deletion}{ps_134-5_049-v0018} conditions \var{\inserted[typed]{except the last}}{except the last}{typed insertion}{ps_134-5_049-v0019} are equivalent; and all \var{\inserted{six}}{\hand{six}}{insertion}{ps_134-5_049-v0020} are equivalent if \var{\inserted[typed]{moreover}}{moreover}{typed insertion}{ps_134-5_049-v0021} \var{\xout{x}}{}{typewriter deletion}{ps_134-5_049-v0022} $e_M \hand{\in} \mathrm{Ob}\, B$ .

Proof. The imp\var{\xout{p}}{}{typewriter deletion}{ps_134-5_049-v0023}lications (i)$\hand{\Rightarrow}$(i')$\hand{\Rightarrow}$(i'') are tautological, on
the other hand (i'') just means that the induced functor $f\colon A\hand{\to} B$ is as-%
pheric, which by prop.1 a) implies that $i$ is aspheric. On the other
hand (i)$\hand{\Rightarrow}$(ii) by the definition of $W$ (cf. prop. 2). If $e_M$ exists
and \var{\inserted{A and}}{\hand{A and}}{insertion}{ps_134-5_049-v0024} the asphericity structure $M_{as}$ \var{\over{is}{are}}{\hand{are}}{writing over}{ps_134-5_049-v0025} aspheric, and if $B$ contains $e_M$ ,
then (ii'') implies that the maps $x\hand{\to} e_M$ for $x$ in $B$ \var{\xout{xxxxx}}{}{typewriter deletion}{ps_134-5_049-v0026} are trans-%
formed by $i^{*}$ into a weak equivalence, and as $e_{A^{\wedge}}$ is aspheric this
implies $i^{*}(x)$ is aspheric i.e. (i''), which \var{\xout{impl}}{}{typewriter deletion}{ps_134-5_049-v0027} proves the last state-%

\margin{49}

%% END SOURCE ps_134-5_049
\ifnum\EditionMode=1
\psapparatusstart
\psentry{ps_134-5_049-v0001}{1}{writing over}{\over{L}{l}}{\hand{l}}{The superposed handwritten reading replaces the underlying reading.}
\psentry{ps_134-5_049-v0002}{2}{typewriter overwriting}{\over[typed]{\illegible{1}}{a}}{a}{The superposed typed reading replaces the underlying reading.}
\psentry{ps_134-5_049-v0003}{3}{typewriter deletion}{\xout{forx}}{}{The typewriter strikes out this material; it leaves the corrected running text.}
\psentry{ps_134-5_049-v0004}{4}{typed insertion}{\inserted[typed]{of $M_{as}$}}{of $M_{as}$}{The typed insertion is incorporated at its indicated position; insertion brackets leave the running text.}
\psentry{ps_134-5_049-v0005}{5}{insertion}{\inserted{containing $i(A)$ and}}{\hand{containing $i(A)$ and}}{The handwritten insertion is incorporated at its indicated position; insertion brackets leave the running text.}
\psentry{ps_134-5_049-v0006}{6}{typewriter deletion}{\xout{aspheric}}{}{The typewriter strikes out this material; it leaves the corrected running text.}
\psentry{ps_134-5_049-v0007}{7}{typed insertion}{\inserted[typed]{a weak equivalence}}{a weak equivalence}{The typed insertion is incorporated at its indicated position; insertion brackets leave the running text.}
\psentry{ps_134-5_049-v0008}{8}{typewriter deletion}{\xout{aspheric,}}{}{The typewriter strikes out this material; it leaves the corrected running text.}
\psentry{ps_134-5_049-v0009}{9}{typed insertion}{\inserted[typed]{a weak equivalence}}{a weak equivalence}{The typed insertion is incorporated at its indicated position; insertion brackets leave the running text.}
\psentry{ps_134-5_049-v0010}{10}{insertion}{\inserted{i.e. $Fl(B) \hand{\subset} (i^{*})^{-1}(\underline{W}_{A})$}}{\hand{i.e. $Fl(B) \hand{\subset} (i^{*})^{-1}(\underline{W}_{A})$}}{The handwritten insertion is incorporated at its indicated position; insertion brackets leave the running text.}
\psentry{ps_134-5_049-v0011}{11}{typed insertion}{\inserted[typed]{if A and}}{if A and}{The typed insertion is incorporated at its indicated position; insertion brackets leave the running text.}
\psentry{ps_134-5_049-v0012}{12}{insertion}{\inserted{$M_{as}$}}{\hand{$M_{as}$}}{The handwritten insertion is incorporated at its indicated position; insertion brackets leave the running text.}
\psentry{ps_134-5_049-v0013}{13}{deletion}{\deleted{is}}{}{The author cancels this material by hand; it leaves the corrected running text.}
\psentry{ps_134-5_049-v0014}{14}{typed insertion}{\inserted[typed]{are}}{are}{The typed insertion is incorporated at its indicated position; insertion brackets leave the running text.}
\psentry{ps_134-5_049-v0015}{15}{typewriter deletion}{\xout{all}}{}{The typewriter strikes out this material; it leaves the corrected running text.}
\psentry{ps_134-5_049-v0016}{16}{insertion}{\inserted{all}}{\hand{all}}{The handwritten insertion is incorporated at its indicated position; insertion brackets leave the running text.}
\psentry{ps_134-5_049-v0017}{17}{typed insertion}{\inserted[typed]{six}}{six}{The typed insertion is incorporated at its indicated position; insertion brackets leave the running text.}
\psentry{ps_134-5_049-v0018}{18}{typewriter deletion}{\xout{\illegible{1}}}{}{The typewriter strikes out this material; it leaves the corrected running text.}
\psentry{ps_134-5_049-v0019}{19}{typed insertion}{\inserted[typed]{except the last}}{except the last}{The typed insertion is incorporated at its indicated position; insertion brackets leave the running text.}
\psentry{ps_134-5_049-v0020}{20}{insertion}{\inserted{six}}{\hand{six}}{The handwritten insertion is incorporated at its indicated position; insertion brackets leave the running text.}
\psentry{ps_134-5_049-v0021}{21}{typed insertion}{\inserted[typed]{moreover}}{moreover}{The typed insertion is incorporated at its indicated position; insertion brackets leave the running text.}
\psentry{ps_134-5_049-v0022}{22}{typewriter deletion}{\xout{x}}{}{The typewriter strikes out this material; it leaves the corrected running text.}
\psentry{ps_134-5_049-v0023}{23}{typewriter deletion}{\xout{p}}{}{The typewriter strikes out this material; it leaves the corrected running text.}
\psentry{ps_134-5_049-v0024}{24}{insertion}{\inserted{A and}}{\hand{A and}}{The handwritten insertion is incorporated at its indicated position; insertion brackets leave the running text.}
\psentry{ps_134-5_049-v0025}{25}{writing over}{\over{is}{are}}{\hand{are}}{The superposed handwritten reading replaces the underlying reading.}
\psentry{ps_134-5_049-v0026}{26}{typewriter deletion}{\xout{xxxxx}}{}{The typewriter strikes out this material; it leaves the corrected running text.}
\psentry{ps_134-5_049-v0027}{27}{typewriter deletion}{\xout{impl}}{}{The typewriter strikes out this material; it leaves the corrected running text.}
\psapparatusend
\fi
\pssource{ps_134-5_050}{134-5/50}{https://grothendieck.umontpellier.fr/web/viewer.html?file=../134-5.pdf\#page=51}
\psnumber{ps_134-5_050-v0001}{1}
\psnumber{ps_134-5_050-v0002}{2}
\psnumber{ps_134-5_050-v0003}{3}
\psnumber{ps_134-5_050-v0004}{4}
\psnumber{ps_134-5_050-v0005}{5}
\psnumber{ps_134-5_050-v0006}{6}
\psnumber{ps_134-5_050-v0007}{7}
\psnumber{ps_134-5_050-v0008}{8}
\psnumber{ps_134-5_050-v0009}{9}
\psnumber{ps_134-5_050-v0010}{10}
\psnumber{ps_134-5_050-v0011}{11}
\psnumber{ps_134-5_050-v0012}{12}
\psnumber{ps_134-5_050-v0013}{13}
\psnumber{ps_134-5_050-v0014}{14}
\psnumber{ps_134-5_050-v0015}{15}
\psnumber{ps_134-5_050-v0016}{16}
\psnumber{ps_134-5_050-v0017}{17}
\psnumber{ps_134-5_050-v0018}{18}
\psnumber{ps_134-5_050-v0019}{19}
%% BEGIN SOURCE ps_134-5_050
% Pursuing Stacks (cotes 134-2 to 134-8): archive cote 134-5, site page 50 (PDF page 51), typescript page 237.
% typescript. Draft compared: ps4.tex:2483-2552.
% AI-transcribed by Claude Opus 5.5 (claude-opus-5-5, adaptive thinking, effort high), Message Batches API, batch ps2 (msgbatch_01NVwr1Yuo8qJf5rykeGBXPK), 2026-11-12; reviewed against both supplied scan halves; see CHECK_DECISIONS.tsv.
\scanpage{134-5/50}
\mspage{237}

ment of the proposition - all six conditions are equivalent in this case.
If no B is given, but still assuming A and $(M,M_{as})$ aspheric, we can
choose a generating subcategory for $(M,M_{as})$ large enough in order to
contain $i(A)$ and $e_M$ , and we get that \var{\over{c}{c}}{\hand{c}}{writing over}{ps_134-5_050-v0001}onditions (i) to (ii') are equi-%
valent, qed.

76. It's time to give some examples.

1) Take $M = (Cat)$, $M_{as} =$ set of $\underline{W}$-aspheric objects in (Cat). We
then got an asphericity structure, as we see by taking any weak test
category A in (Cat)\var{\inserted[typed]{(def.2 on page 172)}}{(def.2 on page 172)}{typed insertion}{ps_134-5_050-v0002}, and the functor
\[ i_A \ : \ a \hand{\longmapsto} A_{/a} \ : \ A \hand{\longrightarrow} (Cat) \quad , \]
which satisfies indeed (i) and (ii) of p.231 above.\margin{//} We may call
\[ (7) \qquad ((Cat),\ (Cat)_{\underline{W}\text{-as}} \overset{\hand{def}}{=} (Cat)_{as} \ ) \]
the ``basic asphericity structure'' , giving rise (by taking the correspon-%
ding set W of ``weak equivalences'') to the ``basic modelizer''\var{\deleted{\unsure{,}}}{}{deletion}{ps_134-5_050-v0003} $((Cat), \underline{W})$.
Turning attention \var{\xout{from}}{}{typewriter deletion}{ps_134-5_050-v0004} towards the former corresponds to a shift in em-%
phasis; whereas previously our\var{\inserted[typed]{main}}{main}{typed insertion}{ps_134-5_050-v0005} emphasis has \var{\xout{\illegible{1}}}{}{typewriter deletion}{ps_134-5_050-v0006} been dwelling consistent-%
ly with the notion of ``weak equivalence''\hand{,} namely with \var{\xout{given}}{}{typewriter deletion}{ps_134-5_050-v0007} giving on a
category M a bunch of \hand{\uline{\textcolor{black}{arrows}}} W , here we are \var{\xout{in}}{}{typewriter deletion}{ps_134-5_050-v0008} working rather with the
notion of ``aspheric \hand{\uline{\textcolor{black}{objects}}}'' as the basic notion. One hint in this direc-%
tion comes from prop.1 \var{\xout{\unsure{p}}}{}{typewriter deletion}{ps_134-5_050-v0009} on p.214 , when we saw that for a functor
$f\colon A \hand{\longrightarrow} B$ between small categories, giving rise to $f^{*}\colon B\hat{} \hand{\longrightarrow} A\hat{}$ (a map
\margin{?} between asphericity structures, as a matter of fact), asphericity of f
can \var{\xout{be}}{}{typewriter deletion}{ps_134-5_050-v0010} always\var{\inserted{be}}{\hand{be}}{insertion}{ps_134-5_050-v0011} expressed in manifold ways as a property of f relative to
the notion of aspheric \hand{\uline{\textcolor{black}{objects}}} in $A\hat{}$and $B\hat{}$, but not as a property relati-%
ve to the notion of weak equivalences, unless A and B are assumed to be
aspheric.

Coming back to the case of the\var{\over{x}{``}}{\hand{``}}{writing over}{ps_134-5_050-v0012}basic asphericity structure'' (7),
we get more general types of ``aspheric'' functors $A \hand{\longrightarrow} (Cat)$ than functors
$i_A$ , by taking any weak test category A and any weak test functor
$i \ : \ A \hand{\longrightarrow} (Cat)$ (cf def.5 p. 175 ),\var{\inserted[typed]{provided however the objects i(a) are aspheric .}}{provided however the objects i(a) are aspheric .}{typed insertion}{ps_134-5_050-v0013} It would seem though that,\var{\xout{the}}{}{typewriter deletion}{ps_134-5_050-v0014}
for a given A, even assuming A to be a \hand{\uline{\textcolor{black}{strict}}} test category say (\var{\over{or}{and}}{\hand{and}}{writing over}{ps_134-5_050-v0015} even
a ``contractor''\var{\inserted{moreover}}{\hand{moreover}}{insertion}{ps_134-5_050-v0016}), \var{\xout{the notion property}}{}{typewriter deletion}{ps_134-5_050-v0017} and restricting to functors
$i\colon A \hand{\longrightarrow} (Cat)$ which factor through $(Cat)_{as}$ (to make them eligible for
being ``appheric functors''), the \var{\deleted{property}}{}{deletion}{ps_134-5_050-v0018}\var{\inserted{condition}}{\hand{condition}}{insertion}{ps_134-5_050-v0019} for i to be a weak test functor\hand{,}
namely for $i^{*}$ to be ``model preserving''\hand{,} is substantially stronger than
mere ``asphericity'': indeed the latter just means that $i^{*}$ transforms weak

\margin{50}

%% END SOURCE ps_134-5_050
\ifnum\EditionMode=1
\psapparatusstart
\psentry{ps_134-5_050-v0001}{1}{writing over}{\over{c}{c}}{\hand{c}}{The superposed handwritten reading replaces the underlying reading.}
\psentry{ps_134-5_050-v0002}{2}{typed insertion}{\inserted[typed]{(def.2 on page 172)}}{(def.2 on page 172)}{The typed insertion is incorporated at its indicated position; insertion brackets leave the running text.}
\psentry{ps_134-5_050-v0003}{3}{deletion}{\deleted{\unsure{,}}}{}{The author cancels this material by hand; it leaves the corrected running text.}
\psentry{ps_134-5_050-v0004}{4}{typewriter deletion}{\xout{from}}{}{The typewriter strikes out this material; it leaves the corrected running text.}
\psentry{ps_134-5_050-v0005}{5}{typed insertion}{\inserted[typed]{main}}{main}{The typed insertion is incorporated at its indicated position; insertion brackets leave the running text.}
\psentry{ps_134-5_050-v0006}{6}{typewriter deletion}{\xout{\illegible{1}}}{}{The typewriter strikes out this material; it leaves the corrected running text.}
\psentry{ps_134-5_050-v0007}{7}{typewriter deletion}{\xout{given}}{}{The typewriter strikes out this material; it leaves the corrected running text.}
\psentry{ps_134-5_050-v0008}{8}{typewriter deletion}{\xout{in}}{}{The typewriter strikes out this material; it leaves the corrected running text.}
\psentry{ps_134-5_050-v0009}{9}{typewriter deletion}{\xout{\unsure{p}}}{}{The typewriter strikes out this material; it leaves the corrected running text.}
\psentry{ps_134-5_050-v0010}{10}{typewriter deletion}{\xout{be}}{}{The typewriter strikes out this material; it leaves the corrected running text.}
\psentry{ps_134-5_050-v0011}{11}{insertion}{\inserted{be}}{\hand{be}}{The handwritten insertion is incorporated at its indicated position; insertion brackets leave the running text.}
\psentry{ps_134-5_050-v0012}{12}{writing over}{\over{x}{``}}{\hand{``}}{The superposed handwritten reading replaces the underlying reading.}
\psentry{ps_134-5_050-v0013}{13}{typed insertion}{\inserted[typed]{provided however the objects i(a) are aspheric .}}{provided however the objects i(a) are aspheric .}{The typed insertion is incorporated at its indicated position; insertion brackets leave the running text.}
\psentry{ps_134-5_050-v0014}{14}{typewriter deletion}{\xout{the}}{}{The typewriter strikes out this material; it leaves the corrected running text.}
\psentry{ps_134-5_050-v0015}{15}{writing over}{\over{or}{and}}{\hand{and}}{The superposed handwritten reading replaces the underlying reading.}
\psentry{ps_134-5_050-v0016}{16}{insertion}{\inserted{moreover}}{\hand{moreover}}{The handwritten insertion is incorporated at its indicated position; insertion brackets leave the running text.}
\psentry{ps_134-5_050-v0017}{17}{typewriter deletion}{\xout{the notion property}}{}{The typewriter strikes out this material; it leaves the corrected running text.}
\psentry{ps_134-5_050-v0018}{18}{deletion}{\deleted{property}}{}{The author cancels this material by hand; it leaves the corrected running text.}
\psentry{ps_134-5_050-v0019}{19}{insertion}{\inserted{condition}}{\hand{condition}}{The handwritten insertion is incorporated at its indicated position; insertion brackets leave the running text.}
\psapparatusend
\fi
\pssource{ps_134-5_051}{134-5/51}{https://grothendieck.umontpellier.fr/web/viewer.html?file=../134-5.pdf\#page=52}
\psnumber{ps_134-5_051-v0001}{1}
\psnumber{ps_134-5_051-v0002}{2}
\psnumber{ps_134-5_051-v0003}{3}
\psnumber{ps_134-5_051-v0004}{4}
\psnumber{ps_134-5_051-v0005}{5}
\psnumber{ps_134-5_051-v0006}{6}
\psnumber{ps_134-5_051-v0007}{7}
\psnumber{ps_134-5_051-v0008}{8}
\psnumber{ps_134-5_051-v0009}{9}
\psnumber{ps_134-5_051-v0010}{10}
\psnumber{ps_134-5_051-v0011}{11}
\psnumber{ps_134-5_051-v0012}{12}
\psnumber{ps_134-5_051-v0013}{13}
\psnumber{ps_134-5_051-v0014}{14}
\psnumber{ps_134-5_051-v0015}{15}
\psnumber{ps_134-5_051-v0016}{16}
\psnumber{ps_134-5_051-v0017}{17}
\psnumber{ps_134-5_051-v0018}{18}
\psnumber{ps_134-5_051-v0019}{19}
\psnumber{ps_134-5_051-v0020}{20}
%% BEGIN SOURCE ps_134-5_051
% Pursuing Stacks (cotes 134-2 to 134-8): archive cote 134-5, site page 51 (PDF page 52), typescript page 238.
% typescript. Draft compared: ps4.tex:2540-2605.
% AI-transcribed by Claude Opus 5.5 (claude-opus-5-5, adaptive thinking, effort high), Message Batches API, batch ps2 (msgbatch_01NVwr1Yuo8qJf5rykeGBXPK), 2026-11-12; reviewed against both supplied scan halves; see CHECK_DECISIONS.tsv.
\scanpage{134-5/51}
\mspage{238}
equivalences into weak equivalences,  i.e. \var{\xout{\unsure{functo}}}{}{typewriter deletion}{ps_134-5_051-v0001} gives rise to a functor
\[ \mathrm{Hot} = \underline{W}^{-1}(\mathrm{Cat}) \hand{\longrightarrow} \mathrm{Hot}_A = \underline{W}_A^{-1}A^{\wedge} \quad , \]
whereas the latter insists thatnthis functor moreover should be an equi-%
valence of categories.   Theorem 1 (p. 176)  gives a hint though that \var{\xout{th}}{}{typewriter deletion}{ps_134-5_051-v0002}
the two conditions may well be equivalent  -  this being so at any rate
provided the objects i(a) in (Cat) are,  not only aspheric,  but even
\hand{\uline{\textcolor{black}{contractible}}}.  This reminds us at once  of the ``silly question'' of
section 46  (p.95),  which was the starting point \var{\inserted[typed]{for}}{for}{typed insertion}{ps_134-5_051-v0003} \var{\xout{of}}{}{typewriter deletion}{ps_134-5_051-v0004} the subsequent
ref\var{\over{f}{l}}{\hand{l}}{writing over}{ps_134-5_051-v0005}ections leading up to the theorem 1 recalled above;  and\hand{,} \var{\inserted{also,}}{\hand{also,}}{insertion}{ps_134-5_051-v0006} beyond
this still somewhat technical result, \var{\inserted{it was the ultimate motivation}}{\hand{it was the ultimate motivation}}{insertion}{ps_134-5_051-v0007} for the present  reflections on as-%
phericity structures.  T\var{\over{h}{h}}{\hand{h}}{writing over}{ps_134-5_051-v0008}e main purpose for these,  I feel,  is to lead
up to a comprehensive answer of the \hand{``}silly question\hand{''}.  We'll have to come
back to this very soon !

2)   Take  M = (Spaces),  the category of topological spaces,  and
$M_{as}$ the spaces which are weakly equivalent to a point.  We get an asphe-%
ricity structure,  indeed an \hand{\uline{\textcolor{black}{aspheric}}} asphericity structure (I forgot to
make this \var{\inserted[typed]{evident}}{evident}{typed insertion}{ps_134-5_051-v0009} specification  in the case \var{\xout{of the}}{}{typewriter deletion}{ps_134-5_051-v0010} 1) above),  as we see by
taking \var{\xout{\unsure{thexx}}}{}{typewriter deletion}{ps_134-5_051-v0011} $A = \hand{\Delta}$ for instance,  and
\[ i: \hand{\Delta} \hand{\longrightarrow} (\mathrm{Spaces}) \]
the ``geometric rea\var{\over{\unsure{a}}{l}}{\hand{l}}{writing over}{ps_134-5_051-v0012}ization\var{\deleted{''}}{}{deletion}{ps_134-5_051-v0013} functor\hand{''}for simplices,  which satisfies
conditions (i) and (ii) of page 231,  as is well known (cf. the book of
Gabriel-Zisman);  as a matter of fact,  i is even a test functor for
the modelizer $(M,W)$\var{\xout{x}}{}{typewriter deletion}{ps_134-5_051-v0014} (which is \var{\inserted[typed]{one way of stating \xout{the}}}{one way of stating }{typed insertion}{ps_134-5_051-v0015} the main content of GZ's book).

3)   T\var{\over{h}{h}}{\hand{h}}{writing over}{ps_134-5_051-v0016}e two examples  above \var{\xout{where}}{}{typewriter deletion}{ps_134-5_051-v0017} are \hand{\uline{\textcolor{black}{aspheric}}} asphericity structu-%
res,  and such moreover that (M,W) is ($\underline{W}$-)modelizing.  These extra fea-%
tures however are not \var{\inserted{always}}{\hand{always}}{insertion}{ps_134-5_051-v0018} present in the next example
\[ (8) \qquad M = A^{\wedge} \ , \quad M_{as} = A^{\wedge}_{as} \quad , \qquad A \text{ a small category,} \]
which is indeed tautologically an asphericity structure,  by taking
the canonical \var{\inserted[typed]{inclusion}}{inclusion}{typed insertion}{ps_134-5_051-v0019} functor  (which is fully faithful)
\[ A \hand{\longrightarrow} A^{\wedge} \quad , \]
satisfying the conditions (i)(ii) above (p.231).   This shows moreover
that the corresponding notion of weak equivalence is the usual one,
(I forgot to state the similar fact in example 2), sorry):
\[ W = \underline{W}_A \quad . \]
T\var{\over{h}{h}}{\hand{h}}{writing over}{ps_134-5_051-v0020}e asphericity structure is aspheric iff A is aspheric.  A functor
\[ A' \hand{\longrightarrow} A \]
\margin{51}

%% END SOURCE ps_134-5_051
\ifnum\EditionMode=1
\psapparatusstart
\psentry{ps_134-5_051-v0001}{1}{typewriter deletion}{\xout{\unsure{functo}}}{}{The typewriter strikes out this material; it leaves the corrected running text.}
\psentry{ps_134-5_051-v0002}{2}{typewriter deletion}{\xout{th}}{}{The typewriter strikes out this material; it leaves the corrected running text.}
\psentry{ps_134-5_051-v0003}{3}{typed insertion}{\inserted[typed]{for}}{for}{The typed insertion is incorporated at its indicated position; insertion brackets leave the running text.}
\psentry{ps_134-5_051-v0004}{4}{typewriter deletion}{\xout{of}}{}{The typewriter strikes out this material; it leaves the corrected running text.}
\psentry{ps_134-5_051-v0005}{5}{writing over}{\over{f}{l}}{\hand{l}}{The superposed handwritten reading replaces the underlying reading.}
\psentry{ps_134-5_051-v0006}{6}{insertion}{\inserted{also,}}{\hand{also,}}{The handwritten insertion is incorporated at its indicated position; insertion brackets leave the running text.}
\psentry{ps_134-5_051-v0007}{7}{insertion}{\inserted{it was the ultimate motivation}}{\hand{it was the ultimate motivation}}{The handwritten insertion is incorporated at its indicated position; insertion brackets leave the running text.}
\psentry{ps_134-5_051-v0008}{8}{writing over}{\over{h}{h}}{\hand{h}}{The superposed handwritten reading replaces the underlying reading.}
\psentry{ps_134-5_051-v0009}{9}{typed insertion}{\inserted[typed]{evident}}{evident}{The typed insertion is incorporated at its indicated position; insertion brackets leave the running text.}
\psentry{ps_134-5_051-v0010}{10}{typewriter deletion}{\xout{of the}}{}{The typewriter strikes out this material; it leaves the corrected running text.}
\psentry{ps_134-5_051-v0011}{11}{typewriter deletion}{\xout{\unsure{thexx}}}{}{The typewriter strikes out this material; it leaves the corrected running text.}
\psentry{ps_134-5_051-v0012}{12}{writing over}{\over{\unsure{a}}{l}}{\hand{l}}{The superposed handwritten reading replaces the underlying reading.}
\psentry{ps_134-5_051-v0013}{13}{deletion}{\deleted{''}}{}{The author cancels this material by hand; it leaves the corrected running text.}
\psentry{ps_134-5_051-v0014}{14}{typewriter deletion}{\xout{x}}{}{The typewriter strikes out this material; it leaves the corrected running text.}
\psentry{ps_134-5_051-v0015}{15}{typed insertion}{\inserted[typed]{one way of stating \xout{the}}}{one way of stating }{The typed insertion is incorporated at its indicated position; insertion brackets leave the running text.}
\psentry{ps_134-5_051-v0016}{16}{writing over}{\over{h}{h}}{\hand{h}}{The superposed handwritten reading replaces the underlying reading.}
\psentry{ps_134-5_051-v0017}{17}{typewriter deletion}{\xout{where}}{}{The typewriter strikes out this material; it leaves the corrected running text.}
\psentry{ps_134-5_051-v0018}{18}{insertion}{\inserted{always}}{\hand{always}}{The handwritten insertion is incorporated at its indicated position; insertion brackets leave the running text.}
\psentry{ps_134-5_051-v0019}{19}{typed insertion}{\inserted[typed]{inclusion}}{inclusion}{The typed insertion is incorporated at its indicated position; insertion brackets leave the running text.}
\psentry{ps_134-5_051-v0020}{20}{writing over}{\over{h}{h}}{\hand{h}}{The superposed handwritten reading replaces the underlying reading.}
\psapparatusend
\fi
\pssource{ps_134-5_052}{134-5/52}{https://grothendieck.umontpellier.fr/web/viewer.html?file=../134-5.pdf\#page=53}
\psnumber{ps_134-5_052-v0001}{1}
\psnumber{ps_134-5_052-v0002}{2}
\psnumber{ps_134-5_052-v0003}{3}
\psnumber{ps_134-5_052-v0004}{4}
\psnumber{ps_134-5_052-v0005}{5}
\psnumber{ps_134-5_052-v0006}{6}
\psnumber{ps_134-5_052-v0007}{7}
\psnumber{ps_134-5_052-v0008}{8}
%% BEGIN SOURCE ps_134-5_052
% Pursuing Stacks (cotes 134-2 to 134-8): archive cote 134-5, site page 52 (PDF page 53), typescript page 239.
% typescript. Draft compared: ps4.tex:2588-2678.
% AI-transcribed by Claude Opus 5.5 (claude-opus-5-5, adaptive thinking, effort high), Message Batches API, batch ps2 (msgbatch_01NVwr1Yuo8qJf5rykeGBXPK), 2026-11-12; reviewed against both supplied scan halves; see CHECK_DECISIONS.tsv.
\scanpage{134-5/52}
\mspage{239}

where A' is another small category,  is $M_{as}$-aspheric  as a functor from
\var{\xout{A'to}}{}{typewriter deletion}{ps_134-5_052-v0001} A' to $A^{\wedge}$, iff it is aspheric.

This last example suggest to call an asphericity structure $(M,M_{as})$
\var{\xout{``aspheric''}}{}{typewriter deletion}{ps_134-5_052-v0002}  ``\hand{\uline{\textcolor{black}{totally aspheric}}}''  if M is stable under finite products,
and if the final object of M,  as well as the product of any two aspheric
objects of M,  is again aspheric;  in other words,  if any finite product
in M whose factors are aspheric is aspheric.  We have,  in this respect:

\hand{\uline{\textcolor{black}{Proposition}}} 5.  Let $(M,M_{as})$ be an asphericity structure,  \var{\xout{The following}}{}{typewriter deletion}{ps_134-5_052-v0003}
where M is stable under finite products.  The following conditions are
equivalent:

(i)  M is totally aspheric,  i.e. $e_M$  and the product of any two
aspheric objects of M are aspheric.

(ii)  There exists a small \var{\xout{generating}}{}{typewriter deletion}{ps_134-5_052-v0004} subcategory B of M ,  stable
under finite products (i.e. containing a final object $e_M$ of M and,  with
any two objects x and y,  a product \var{\xout{xy}}{}{typewriter deletion}{ps_134-5_052-v0005} $x \mathrm{x} y$  in M),  and which \var{\xout{is}}{}{typewriter deletion}{ps_134-5_052-v0006} gene-%
rates the asphericity structure  (i.e.  $j\colon B\hand{\hookrightarrow} M$  is $M_{as}$-aspheric).

(iii)  There exists a \var{\xout{\unsure{small}}}{}{typewriter deletion}{ps_134-5_052-v0007} small category A such that $A^{\wedge}$\var{\xout{\illegible{1}}}{}{typewriter deletion}{ps_134-5_052-v0008} is totally
aspheric (def.1  p. 170),  and a $M_{as}$-aspheric functor  $A \hand{\longrightarrow} M$ .

% A single pen arrow points upward beside this paragraph, with its lower end curving right.
\margin{$\uparrow$}The proof is immediate.  Of course,  the asphericity structure in
example 3)  is totally aspheric,  iff $A^{\wedge}$ is totally aspheric in the usual
sense referred to above.

\margin{$\ulcorner$}77.    Let
\margin{$\downarrow$}(9) \[ \underline{M} \;=\; (M,M_{as}) \]
be an asphericity structure,  which will be referred to also as merely
M ,  when no ambiguity concerning $M_{as}$  is feared.  We'll write
(10) \[ \mathrm{Hot}_{\underline{M}} \Big(\text{or simply }\ \mathrm{Hot}_{M}\Big) = W^{-1}M \quad , \]
where W of course is the set of weak equivalences in M.  We are now
going to define a canonical functor
(11) \[ \mathrm{Hot}_{\underline{M}} \hand{\longrightarrow} (\mathrm{Hot})_{\underline{W}} \ \ (\text{or simply } (\mathrm{Hot}) \overset{\mathrm{def}}{=} \underline{W}^{-1}(\mathrm{Cat})) \ , \]
defined at any rate up to canonical isomorphism.  For this,  take any
aspheric (namely, $M_{as}$-aspheric) functor
\[ i\colon\ A \hand{\longrightarrow} M \quad , \]
and consider the composition
(12) \[ M \hand{\xrightarrow{\ \textcolor{black}{i^{*}}\ }} A^{\wedge} \hand{\xrightarrow{\ \textcolor{black}{i_{A}}\ }} (\mathrm{Cat}) \quad . \]

\margin{52}

%% END SOURCE ps_134-5_052
\ifnum\EditionMode=1
\psapparatusstart
\psentry{ps_134-5_052-v0001}{1}{typewriter deletion}{\xout{A'to}}{}{The typewriter strikes out this material; it leaves the corrected running text.}
\psentry{ps_134-5_052-v0002}{2}{typewriter deletion}{\xout{``aspheric''}}{}{The typewriter strikes out this material; it leaves the corrected running text.}
\psentry{ps_134-5_052-v0003}{3}{typewriter deletion}{\xout{The following}}{}{The typewriter strikes out this material; it leaves the corrected running text.}
\psentry{ps_134-5_052-v0004}{4}{typewriter deletion}{\xout{generating}}{}{The typewriter strikes out this material; it leaves the corrected running text.}
\psentry{ps_134-5_052-v0005}{5}{typewriter deletion}{\xout{xy}}{}{The typewriter strikes out this material; it leaves the corrected running text.}
\psentry{ps_134-5_052-v0006}{6}{typewriter deletion}{\xout{is}}{}{The typewriter strikes out this material; it leaves the corrected running text.}
\psentry{ps_134-5_052-v0007}{7}{typewriter deletion}{\xout{\unsure{small}}}{}{The typewriter strikes out this material; it leaves the corrected running text.}
\psentry{ps_134-5_052-v0008}{8}{typewriter deletion}{\xout{\illegible{1}}}{}{The typewriter strikes out this material; it leaves the corrected running text.}
\psapparatusend
\fi
\pssource{ps_134-5_053}{134-5/53}{https://grothendieck.umontpellier.fr/web/viewer.html?file=../134-5.pdf\#page=54}
\psnumber{ps_134-5_053-v0001}{1}
\psnumber{ps_134-5_053-v0002}{2}
\psnumber{ps_134-5_053-v0003}{3}
\psnumber{ps_134-5_053-v0004}{4}
\psnumber{ps_134-5_053-v0005}{5}
\psnumber{ps_134-5_053-v0006}{6}
\psnumber{ps_134-5_053-v0007}{7}
\psnumber{ps_134-5_053-v0008}{8}
\psnumber{ps_134-5_053-v0009}{9}
\psnumber{ps_134-5_053-v0010}{10}
\psnumber{ps_134-5_053-v0011}{11}
\psnumber{ps_134-5_053-v0012}{12}
%% BEGIN SOURCE ps_134-5_053
% Pursuing Stacks (cotes 134-2 to 134-8): archive cote 134-5, site page 53 (PDF page 54), typescript page 240.
% typescript. Draft compared: ps4.tex:2652-2724.
% AI-transcribed by Claude Opus 5.5 (claude-opus-5-5, adaptive thinking, effort high), Message Batches API, batch ps2 (msgbatch_01NVwr1Yuo8qJf5rykeGBXPK), 2026-11-12; reviewed against both supplied scan halves; see CHECK_DECISIONS.tsv.
\scanpage{134-5/53}
\mspage{240}

The three categories in (12) are endowed respectively with sets of arrows
$W$ , $\underline{W}_A$ , $\underline{W}$ , and the two functors satisfy the conditions
\[ (13) \qquad W = (i^{*})^{-1}(\underline{W}_A) \quad \text{and} \quad \underline{W}_A = (i_A)^{-1}(\underline{W}) \quad , \quad \text{hence} \quad W = (i_M)^{-1}(\underline{W}) \quad , \]
where
\[ i_{M,i}\colon \ M \hand{\longrightarrow} (\mathrm{Cat}) \]
is the composition $i_A i^{*}$ . Therefore we get a functor
\[ (14) \qquad \hand{\overline{\textcolor{black}{i_{M,i}}}} \ : \ \mathrm{Hot}_M \hand{\longrightarrow} (\mathrm{Hot}) \quad , \]
which a priori depends upon the choice of $(A,i)$. If we admit strong sa-%
turation of $\underline{W}$, it follows that this functor is ``conservative'', namely
an arrow in $\mathrm{Hot}_M$ is an isomorphism, provided it's image in (Hot) is.

To define (11) in terms of (14), we have to describe merely a tran-%
sitive system of isomorphisms between the functors (14), for varying
pair $(A,i)$. Therefore, consider two such $(A,i)$ and $(A',i')$, choose
a small \var{\inserted[typed]{full}}{full}{typed insertion}{ps_134-5_053-v0001} subcategory \var{\inserted[typed]{B}}{B}{typed insertion}{ps_134-5_053-v0002} of $M_{as}$ containing both $i(A)$ and $i'(A')$ (therefore
$B$ is a generating subcategory for the asphericity structure $M_{as}$ ), and
consider the incl\unsure{us}ion functor $j\colon B \hand{\hookrightarrow} M$. From prop.1 b) it follows
that the functors
\[ f\colon A \hand{\longrightarrow} B \quad , \quad f' : A' \hand{\longrightarrow} B \]
induced by $i,i'$ are aspheric. Using this, and the criterion (i) of
asphericity (prop. 1 on p. 214) we immediately get two isomorphisms
% Hand-drawn arrows are marked blue; typed objects remain plain.
\diagram{\begin{tikzcd}[column sep=small]
 \& \hand{\overline{\textcolor{black}{i_{M,\var{\over{i}{j}}{\hand{j}}{writing over}{ps_134-5_053-v0003}}}}} \ar[dl, dash, color=blue, "\hand{\sim}"'] \ar[dr, dash, color=blue, "\hand{\sim}"] \& \\
\hand{\overline{\textcolor{black}{i_{M,\var{\over{i}{i}}{\hand{i}}{writing over}{ps_134-5_053-v0004}}}}} \& \& \hand{\overline{\textcolor{black}{i_{M,i'}}}}
\end{tikzcd}}
,

and hence an isomorphism
\[ \hand{\overline{\textcolor{black}{i_{M,i}}}} \hand{\xrightarrow{\sim}} \hand{\overline{\textcolor{black}{i_{M,i'}}}} \quad , \]
depending only on the choice of $B$. As a matter of fact, it does'nt
depend on this choice\var{\over{,}{.}}{\hand{.}}{writing over}{ps_134-5_053-v0005} \hand{\unsure{$\leftarrow$}} \var{\over{t}{T}}{\hand{T}}{writing over}{ps_134-5_053-v0006}o see this, we are reduced to comparing
the two isomorphi\unsure{sm}s arising from a $B$ and a $B'$ such that $B \hand{\subset} B'$, in \var{\xout{whi}}{}{typewriter deletion}{ps_134-5_053-v0007}
which case the inclusion functor $g\colon B \hand{\to} B'$ is aspheric\var{\inserted{(prop 1)}}{\hand{(prop 1)}}{insertion}{ps_134-5_053-v0008}. We leave the
detail\hand{s} of the verification\var{\deleted{,}}{}{deletion}{ps_134-5_053-v0009} \hand{(}consisting mainly of some diagram chasing
and some compatibilities between the $\hand{\lambda}_i(F)$ -isomorphisms of prop.1 on
page 214\hand{)} to the sceptical reader. As for transitivity for a triple
$(A,i),(A',i'),(A'',i'')$, it \var{\inserted{now}}{\hand{now}}{insertion}{ps_134-5_053-v0010} follows at once, by using a $B$ suitable \var{\deleted{sim}}{}{deletion}{ps_134-5_053-v0011} \var{\inserted{simul-}}{\hand{simul-}}{insertion}{ps_134-5_053-v0012}

\margin{53}

%% END SOURCE ps_134-5_053
\ifnum\EditionMode=1
\psapparatusstart
\psentry{ps_134-5_053-v0001}{1}{typed insertion}{\inserted[typed]{full}}{full}{The typed insertion is incorporated at its indicated position; insertion brackets leave the running text.}
\psentry{ps_134-5_053-v0002}{2}{typed insertion}{\inserted[typed]{B}}{B}{The typed insertion is incorporated at its indicated position; insertion brackets leave the running text.}
\psentry{ps_134-5_053-v0003}{3}{writing over}{$\over{i}{j}$}{$\hand{j}$}{The superposed handwritten reading replaces the underlying reading.}
\psentry{ps_134-5_053-v0004}{4}{writing over}{$\over{i}{i}$}{$\hand{i}$}{The superposed handwritten reading replaces the underlying reading.}
\psentry{ps_134-5_053-v0005}{5}{writing over}{\over{,}{.}}{\hand{.}}{The superposed handwritten reading replaces the underlying reading.}
\psentry{ps_134-5_053-v0006}{6}{writing over}{\over{t}{T}}{\hand{T}}{The superposed handwritten reading replaces the underlying reading.}
\psentry{ps_134-5_053-v0007}{7}{typewriter deletion}{\xout{whi}}{}{The typewriter strikes out this material; it leaves the corrected running text.}
\psentry{ps_134-5_053-v0008}{8}{insertion}{\inserted{(prop 1)}}{\hand{(prop 1)}}{The handwritten insertion is incorporated at its indicated position; insertion brackets leave the running text.}
\psentry{ps_134-5_053-v0009}{9}{deletion}{\deleted{,}}{}{The author cancels this material by hand; it leaves the corrected running text.}
\psentry{ps_134-5_053-v0010}{10}{insertion}{\inserted{now}}{\hand{now}}{The handwritten insertion is incorporated at its indicated position; insertion brackets leave the running text.}
\psentry{ps_134-5_053-v0011}{11}{deletion}{\deleted{sim}}{}{The author cancels this material by hand; it leaves the corrected running text.}
\psentry{ps_134-5_053-v0012}{12}{insertion}{\inserted{simul-}}{\hand{simul-}}{The handwritten insertion is incorporated at its indicated position; insertion brackets leave the running text.}
\psapparatusend
\fi
\pssource{ps_134-5_054}{134-5/54}{https://grothendieck.umontpellier.fr/web/viewer.html?file=../134-5.pdf\#page=55}
\psnumber{ps_134-5_054-v0001}{1}
\psnumber{ps_134-5_054-v0002}{2}
\psnumber{ps_134-5_054-v0003}{3}
\psnumber{ps_134-5_054-v0004}{4}
\psnumber{ps_134-5_054-v0005}{5}
\psnumber{ps_134-5_054-v0006}{6}
\psnumber{ps_134-5_054-v0007}{7}
\psnumber{ps_134-5_054-v0008}{8}
\psnumber{ps_134-5_054-v0009}{9}
\psnumber{ps_134-5_054-v0010}{10}
\psnumber{ps_134-5_054-v0011}{11}
\psnumber{ps_134-5_054-v0012}{12}
\psnumber{ps_134-5_054-v0013}{13}
\psnumber{ps_134-5_054-v0014}{14}
\psnumber{ps_134-5_054-v0015}{15}
\psnumber{ps_134-5_054-v0016}{16}
\psnumber{ps_134-5_054-v0017}{17}
%% BEGIN SOURCE ps_134-5_054
% Pursuing Stacks (cotes 134-2 to 134-8): archive cote 134-5, site page 54 (PDF page 55), typescript page 241.
% typescript. Draft compared: ps4.tex:2711-2789.
% AI-transcribed by Claude Opus 5.5 (claude-opus-5-5, adaptive thinking, effort high), Message Batches API, batch ps2 (msgbatch_01NVwr1Yuo8qJf5rykeGBXPK), 2026-11-12; reviewed against both supplied scan halves; see CHECK_DECISIONS.tsv.
\scanpage{134-5/54}
\mspage{241}
taneously for all three.

Having thus well in hand the \var{\xout{all-important}}{}{typewriter deletion}{ps_134-5_054-v0001}\var{\inserted[typed]{basic}}{basic}{typed insertion}{ps_134-5_054-v0002} functor (11) (which in
case of example \var{\xout{three}}{}{typewriter deletion}{ps_134-5_054-v0003} 3) above with $M=A^{\wedge}$, reduces to the all-important
functor $\mathrm{Hot}_{A}\hand{\to}(\mathrm{Hot})$), \var{\xout{Hot}}{}{typewriter deletion}{ps_134-5_054-v0004} we cannot but define an asphericity structure
to be \hand{\uline{\textcolor{black}{modelizing}}} , as meaning that this canonical functor is an equiva-%
lence of categories. In the case of $M= A^{\wedge}$ above, this means that A
is a pseudo-test category - a relatively weak test notion still. It
appears as just a \var{\xout{bit}}{}{typewriter deletion}{ps_134-5_054-v0005} small bit stronger than merely assuming the pair
$(M,W)$ to be modelizing, \var{\inserted[typed]{i.e. to be a ``\hand{\uline{\textcolor{black}{modelizer}}}'',}}{i.e. to be a ``\hand{\uline{\textcolor{black}{modelizer}}}'',}{typed insertion}{ps_134-5_054-v0006} namely assuming $\mathrm{Hot}_{M}$ to be equivalent (in some
way or other ...) to (Hot). Maybe we should be a little more cautious
with the use of the word ``modelizing'' though, and devise a terminology
which should reflect very closely the \unsure{hierarchy|Hierarchy} of \var{\deleted{progressively stronger}}{}{deletion}{ps_134-5_054-v0007}
test notions
\[
\begin{aligned}&\text{(pseudo-test cat.)}\hand{\supset}\text{(weak test cat.)}\hand{\supset}\text{(test cat.)}\\&\qquad\hand{\supset}\text{(strict test cat.)}\hand{\supset} \var{\inserted{\deleted{\illegible{1}}}}{\hand{}}{insertion}{ps_134-5_054-v0008}\var{\inserted[typed]{(contractors)}}{(contractors)}{typed insertion}{ps_134-5_054-v0009}\end{aligned}
\]
which gradually has peeled out of our earlier reflections, \var{\inserted[typed]{by pinpointing}}{by pinpointing}{typed insertion}{ps_134-5_054-v0010} \var{\xout{with}}{}{typewriter deletion}{ps_134-5_054-v0011} corres-%
ponding qualifications for an asphericity structure,as \var{\inserted[typed]{being}}{being}{typed insertion}{ps_134-5_054-v0012} ``pseudo-modeli\var{\xout{xxx}}{}{typewriter deletion}{ps_134-5_054-v0013}%
zing'' , \var{\xout{``modelizing''xx}}{}{typewriter deletion}{ps_134-5_054-v0014} ``weakly modelizing'', ``modelizing'', ``strictly mode%
lizing''\hand{,} and ??? . It now appears that a little extra reflection is nee-%
ded here - for today it's getting a little late though !

18.6\hand{, and 19.6.}

78. It is about time now to get a comprehensive treatment, in the
context of asphericity structures, of the relationships suggested time
ago in the ``observation'' and the \var{\xout{x}}{}{typewriter deletion}{ps_134-5_054-v0015}``silly question'' of section 46 (pages
94 and 95). The former is concerned with test-functors from test cate-%
gories to modelizers, the second more generally with model-preserving
maps between modelizers, having properties similar to the map $M\hand{\to} A^{\wedge}$
stemming from a test-functor. For the time being, as I found but litt-%
le time for the ``extra reflection'' which seems needed, only the first
situation is by now reasonably clear in my mind.

As no evident \var{\over{p}{s}}{\hand{s}}{writing over}{ps_134-5_054-v0016}eries of ``asphericity structure''-notions has appea-%
red, paralleling the series of test-notions \var{\xout{displayed}}{}{typewriter deletion}{ps_134-5_054-v0017} recalled by the
end of last monday's reflections (see above), I'm going to keep (provi-%
sionally only, maybe) the name of a \hand{\uline{\textcolor{black}{modelizing}}} \hand{\uline{\textcolor{black}{asphericity}}} \hand{\uline{\textcolor{black}{structure}}}
$(M,M_{\mathrm{as}})$ as one for which the canonical functor
\[
(1) \qquad \mathrm{Hot}_{M} = W^{-1}M \hand{\longrightarrow} (\mathrm{Hot})
\]
(cf (14) p. 240) is an equivalence. This notion at any rate is sa-%
\margin{54}
% A small circumflex-like mark is visible above the p of asphericity near the opening paragraph; its function is not securely identifiable.

%% END SOURCE ps_134-5_054
\ifnum\EditionMode=1
\psapparatusstart
\psentry{ps_134-5_054-v0001}{1}{typewriter deletion}{\xout{all-important}}{}{The typewriter strikes out this material; it leaves the corrected running text.}
\psentry{ps_134-5_054-v0002}{2}{typed insertion}{\inserted[typed]{basic}}{basic}{The typed insertion is incorporated at its indicated position; insertion brackets leave the running text.}
\psentry{ps_134-5_054-v0003}{3}{typewriter deletion}{\xout{three}}{}{The typewriter strikes out this material; it leaves the corrected running text.}
\psentry{ps_134-5_054-v0004}{4}{typewriter deletion}{\xout{Hot}}{}{The typewriter strikes out this material; it leaves the corrected running text.}
\psentry{ps_134-5_054-v0005}{5}{typewriter deletion}{\xout{bit}}{}{The typewriter strikes out this material; it leaves the corrected running text.}
\psentry{ps_134-5_054-v0006}{6}{typed insertion}{\inserted[typed]{i.e. to be a ``\hand{\uline{\textcolor{black}{modelizer}}}'',}}{i.e. to be a ``\hand{\uline{\textcolor{black}{modelizer}}}'',}{The typed insertion is incorporated at its indicated position; insertion brackets leave the running text.}
\psentry{ps_134-5_054-v0007}{7}{deletion}{\deleted{progressively stronger}}{}{The author cancels this material by hand; it leaves the corrected running text.}
\psentry{ps_134-5_054-v0008}{8}{insertion}{$\inserted{\deleted{\illegible{1}}}$}{$\hand{}$}{The handwritten insertion is incorporated at its indicated position; insertion brackets leave the running text.}
\psentry{ps_134-5_054-v0009}{9}{typed insertion}{$\inserted[typed]{(contractors)}$}{$(contractors)$}{The typed insertion is incorporated at its indicated position; insertion brackets leave the running text.}
\psentry{ps_134-5_054-v0010}{10}{typed insertion}{\inserted[typed]{by pinpointing}}{by pinpointing}{The typed insertion is incorporated at its indicated position; insertion brackets leave the running text.}
\psentry{ps_134-5_054-v0011}{11}{typewriter deletion}{\xout{with}}{}{The typewriter strikes out this material; it leaves the corrected running text.}
\psentry{ps_134-5_054-v0012}{12}{typed insertion}{\inserted[typed]{being}}{being}{The typed insertion is incorporated at its indicated position; insertion brackets leave the running text.}
\psentry{ps_134-5_054-v0013}{13}{typewriter deletion}{\xout{xxx}}{}{The typewriter strikes out this material; it leaves the corrected running text.}
\psentry{ps_134-5_054-v0014}{14}{typewriter deletion}{\xout{``modelizing''xx}}{}{The typewriter strikes out this material; it leaves the corrected running text.}
\psentry{ps_134-5_054-v0015}{15}{typewriter deletion}{\xout{x}}{}{The typewriter strikes out this material; it leaves the corrected running text.}
\psentry{ps_134-5_054-v0016}{16}{writing over}{\over{p}{s}}{\hand{s}}{The superposed handwritten reading replaces the underlying reading.}
\psentry{ps_134-5_054-v0017}{17}{typewriter deletion}{\xout{displayed}}{}{The typewriter strikes out this material; it leaves the corrected running text.}
\psapparatusend
\fi
\pssource{ps_134-5_055}{134-5/55}{https://grothendieck.umontpellier.fr/web/viewer.html?file=../134-5.pdf\#page=56}
\psnumber{ps_134-5_055-v0001}{1}
\psnumber{ps_134-5_055-v0002}{2}
\psnumber{ps_134-5_055-v0003}{3}
\psnumber{ps_134-5_055-v0004}{4}
\psnumber{ps_134-5_055-v0005}{5}
\psnumber{ps_134-5_055-v0006}{6}
\psnumber{ps_134-5_055-v0007}{7}
\psnumber{ps_134-5_055-v0008}{8}
\psnumber{ps_134-5_055-v0009}{9}
\psnumber{ps_134-5_055-v0010}{10}
\psnumber{ps_134-5_055-v0011}{11}
%% BEGIN SOURCE ps_134-5_055
% Pursuing Stacks (cotes 134-2 to 134-8): archive cote 134-5, site page 55 (PDF page 56), typescript page 242.
% typescript. Draft compared: ps4.tex:2771-2844.
% AI-transcribed by Claude Opus 5.5 (claude-opus-5-5, adaptive thinking, effort high), Message Batches API, batch ps2 (msgbatch_01NVwr1Yuo8qJf5rykeGBXPK), 2026-11-12; reviewed against both supplied scan halves; see CHECK_DECISIONS.tsv.
\scanpage{134-5/55}
\mspage{242}
tisfactory for f\unsure{or}mulating the following statement,  which comes out here
rather tautologically,  and which however appears to me as exactly what
I had been looking for in the ``observation'' recalled above:

\hand{\uline{\textcolor{black}{Theorem}1}} Let $(M,M_{as})$  be a \hand{\uline{\textcolor{black}{modelizing}}} asphericity structure,  \var{\xout{and}}{}{typewriter deletion}{ps_134-5_055-v0001} A a
\var{\xout{smallxtest}}{}{typewriter deletion}{ps_134-5_055-v0002} pseudo-test category,  and
\[ i:\ A \hand{\longrightarrow} M \]
a functor,  factoring through $M_{as}$  (i.e. $i(a)$ is aspheric for any a in A)\hand{.}
\var{\inserted[typed]{We assume moreover M has a final object $e_M$ .}}{We assume moreover M has a final object $e_M$ .}{typed insertion}{ps_134-5_055-v0003}
Then the six conditions (i) to (ii") of prop.4  (p.236)\hand{,} \hand{(}the first
of which expresses that i is $M_{as}$-aspheric\hand{)} are equivalent,  and they are
equival\unsure{e|a}nt to the following condition:
\var{\xout{(iii)}}{}{typewriter deletion}{ps_134-5_055-v0004}  (iii)  The functor  $i^{*}$: $M \hand{\to} A\hat{}$ gives rise to a functor
\[ Hot_M = W^{-1}M \hand{\longrightarrow} Hot_A = \underline{W}_A{}^{-1}A\hat{} \]
(i.e.  (ii') of prop.4 holds,  namely  $i^{*}(W) \hand{\subset} \underline{W}_A$ ),  \hand{\uline{\textcolor{black}{and}}}  this functor
moreover is an equivalence of categories.

Proof.  For the first statement (equivalence of conditiohs (i) to
(ii')), \var{\inserted[typed]{By prop.4}}{By prop.4}{typed insertion}{ps_134-5_055-v0005} we need only show that the asphericity structure  and A are as-%
pheric.  But \var{\xout{itxis}}{}{typewriter deletion}{ps_134-5_055-v0006} this follows from the assumptions and \var{\inserted{from}}{\hand{from}}{insertion}{ps_134-5_055-v0007} the

\hand{\uline{\textcolor{black}{Lemma}}} : Let $(M,M_{as})$  be any modelizing asphericity structure,  then a
final object of M is aspheric\var{\inserted[typed]{(i.e. the structure is ``aspheric'').}}{(i.e. the structure is ``aspheric'').}{typed insertion}{ps_134-5_055-v0008} In particular,  if A is a pseudo-test
category (i.e.  $(A\hat{},A\hat{}_{as})$  is a modelizing asphericity structure),  then A
is aspheric.

We only have to prove the first statement.  \var{\xout{Let}}{}{typewriter deletion}{ps_134-5_055-v0009} It is immediate
that $e_M$  is a final object of $W^{-1}M = Hot_M$  (this is valid whenever we
got a localization $W^{-1}M$  of a category M with final object $e_M$ ),  hence
it's image in \var{\xout{(Cat}}{}{typewriter deletion}{ps_134-5_055-v0010}(Hot) = $\underline{W}^{-1}$(Cat)  is a final object.  Thus we are redu-%
ced to proving the following

\hand{\uline{\textcolor{black}{Corollary}}} :  Let $(M,M_{as})$ be any asphericity structure,  consider the
composition
\[ \hand{\phi}_M\ :\ M \hand{\longrightarrow} Hot_M = W^{-1}M \hand{\xrightarrow{can.}} (Hot) \quad , \]
then we get
\[ M_{as} = \hand{\{} x \hand{\in} Ob\ M \ \hand{|}\ \hand{\phi}_M(x) \text{ is a final object of (Hot)}\hand{\}}. \]

Indeed,  using the construction of  $Hot_M \hand{\longrightarrow} (Hot)$  in terms of a
given $M_{as}$-aspheric functor  $B \hand{\to} M$ ,  we are reduced to the same state-%
ment,  with $(M,M_{as})$  replaced by $(B\hat{},\ B\hat{}_{as})$.  The statement then \var{\over{d}{r}}{\hand{r}}{writing over}{ps_134-5_055-v0011}educes
to :  the object B in (Cat)  is aspheric iff it's image in (Hot) is a
final object  (this can be viewed also as the particular case of the
\margin{55}
% Small displaced marks occur below the n in Indeed and above/below the h area in The statement; their precise correction function is unresolved.

%% END SOURCE ps_134-5_055
\ifnum\EditionMode=1
\psapparatusstart
\psentry{ps_134-5_055-v0001}{1}{typewriter deletion}{\xout{and}}{}{The typewriter strikes out this material; it leaves the corrected running text.}
\psentry{ps_134-5_055-v0002}{2}{typewriter deletion}{\xout{smallxtest}}{}{The typewriter strikes out this material; it leaves the corrected running text.}
\psentry{ps_134-5_055-v0003}{3}{typed insertion}{\inserted[typed]{We assume moreover M has a final object $e_M$ .}}{We assume moreover M has a final object $e_M$ .}{The typed insertion is incorporated at its indicated position; insertion brackets leave the running text.}
\psentry{ps_134-5_055-v0004}{4}{typewriter deletion}{\xout{(iii)}}{}{The typewriter strikes out this material; it leaves the corrected running text.}
\psentry{ps_134-5_055-v0005}{5}{typed insertion}{\inserted[typed]{By prop.4}}{By prop.4}{The typed insertion is incorporated at its indicated position; insertion brackets leave the running text.}
\psentry{ps_134-5_055-v0006}{6}{typewriter deletion}{\xout{itxis}}{}{The typewriter strikes out this material; it leaves the corrected running text.}
\psentry{ps_134-5_055-v0007}{7}{insertion}{\inserted{from}}{\hand{from}}{The handwritten insertion is incorporated at its indicated position; insertion brackets leave the running text.}
\psentry{ps_134-5_055-v0008}{8}{typed insertion}{\inserted[typed]{(i.e. the structure is ``aspheric'').}}{(i.e. the structure is ``aspheric'').}{The typed insertion is incorporated at its indicated position; insertion brackets leave the running text.}
\psentry{ps_134-5_055-v0009}{9}{typewriter deletion}{\xout{Let}}{}{The typewriter strikes out this material; it leaves the corrected running text.}
\psentry{ps_134-5_055-v0010}{10}{typewriter deletion}{\xout{(Cat}}{}{The typewriter strikes out this material; it leaves the corrected running text.}
\psentry{ps_134-5_055-v0011}{11}{writing over}{\over{d}{r}}{\hand{r}}{The superposed handwritten reading replaces the underlying reading.}
\psapparatusend
\fi
\pssource{ps_134-5_056}{134-5/56}{https://grothendieck.umontpellier.fr/web/viewer.html?file=../134-5.pdf\#page=57}
\psnumber{ps_134-5_056-v0001}{1}
\psnumber{ps_134-5_056-v0002}{2}
\psnumber{ps_134-5_056-v0003}{3}
\psnumber{ps_134-5_056-v0004}{4}
\psnumber{ps_134-5_056-v0005}{5}
\psnumber{ps_134-5_056-v0006}{6}
\psnumber{ps_134-5_056-v0007}{7}
\psnumber{ps_134-5_056-v0008}{8}
\psnumber{ps_134-5_056-v0009}{9}
\psnumber{ps_134-5_056-v0010}{10}
\psnumber{ps_134-5_056-v0011}{11}
\psnumber{ps_134-5_056-v0012}{12}
\psnumber{ps_134-5_056-v0013}{13}
\psnumber{ps_134-5_056-v0014}{14}
%% BEGIN SOURCE ps_134-5_056
% Pursuing Stacks (cotes 134-2 to 134-8): archive cote 134-5, site page 56 (PDF page 57), typescript page 243.
% typescript. Draft compared: ps4.tex:2830-2896.
% AI-transcribed by Claude Opus 5.5 (claude-opus-5-5, adaptive thinking, effort high), Message Batches API, batch ps2 (msgbatch_01NVwr1Yuo8qJf5rykeGBXPK), 2026-11-12; reviewed against both supplied scan halves; see CHECK_DECISIONS.tsv.
\scanpage{134-5/56}
\mspage{243}
corollary,  when $M = (\mathrm{Cat})$ ,  $M_{as} = (\mathrm{Cat})_{as}$ ).  Now this follows at once
from \var{\inserted[typed]{strong}}{strong}{typed insertion}{ps_134-5_056-v0001} saturation of $\underline{W}$ ,  which (as we announced earlier) follows from
Loc 1)  to  Loc 3)  (and will be pro\var{\xout{\unsure{s}}}{}{typewriter deletion}{ps_134-5_056-v0002}ved in \var{\xout{the}}{}{typewriter deletion}{ps_134-5_056-v0003} part \hand{V} of the notes).
The reader who fears a vicious circle may till then restrict \var{\xout{eventual}}{}{typewriter deletion}{ps_134-5_056-v0004}
use of the theorem to the case when we assume beforehand that $\var{\over{\illegible{1}}{e}}{\hand{e}}{writing over}{ps_134-5_056-v0005}_M$ and
$e_{A^\wedge}$ are aspheric .

It is now clear that the six conditions of prop. 4  are equiva-%
lent,  and they \var{\xout{imply}}{}{typewriter deletion}{ps_134-5_056-v0006} \var{\inserted[typed]{are}}{are}{typed insertion}{ps_134-5_056-v0007} of course \var{\xout{\unsure{(iii)}}}{}{typewriter deletion}{ps_134-5_056-v0008} implied by (iii).  Conversely,
they imply (iii),  as follows from the fact that in the canonical diagram
(commutative up to canonical isomorphism)
% Hand-drawn arrows are marked blue; typed objects remain plain.
\diagram{\begin{tikzcd}[column sep=tiny]
\mathrm{Hot}_M \ar[rr,color=blue] \ar[dr,color=blue] \& \& \mathrm{Hot}_A \ar[dl,color=blue] \\
\& (\mathrm{Hot}) \&
\end{tikzcd}}
,
the two downwards arrows are equivalences,  qed.

\var{\xout{The theorem}}{}{typewriter deletion}{ps_134-5_056-v0009}
The theorem above seems to me to be exactly the ``something very
simple-minded surely'' which I was feeling to get b\unsure{ur}ningly clos\unsure{e}, by
the end of march,  ne\var{\over[typed]{q}{a}}{a}{typewriter overwriting}{ps_134-5_056-v0010}rly three months ago (p.89);  at least, \var{\inserted[typed]{to be ``just it'' ,}}{to be ``just it'' ,}{typed insertion}{ps_134-5_056-v0011} as far as
the case of \hand{\uline{\textcolor{black}{test-functors}}} is concerned.  We may equally view th\var{\over{\unsure{i}}{i}}{\hand{i}}{writing over}{ps_134-5_056-v0012}s theo-%
rem as giving the precise relationship between the notion of a weak test
functor or a \hand{\uline{\textcolor{black}{test functor}}}, \hand{(}the latest version of which (in the \unsure{c}ontext
of $\underline{W}$-notions)  appears in section 65  (def.5 and 6,  p. 175 and 176)\hand{)},
and the notion of \hand{\uline{\textcolor{black}{aspheric functors}}},  more precisely of $M_{as}$-aspheric
functors,  introduced lately (p. 233).  This now is the moment surely
to check if the terminology of test functors and weak ones  introduced
earlier,  before the relevant notion of asphericity structures was at
hand,  is really satisfactory indeed,  and if needed,  adjust it slightly.

So let again
\[ i\colon A \hand{\longrightarrow} M \]
be a functor,  with A small,  and $(M,M_{as})$ an \var{\inserted{\deleted{aspheric}}}{\hand{}}{insertion}{ps_134-5_056-v0013} asphericity structure.  We
don't  assume beforehand,  neither that A is a test-category or the like,
nor that $(M,M_{as})$ be modelizing.  We now paraphrase def.5  (p.175) of
weak test functors as follows:

\hand{\uline{\textcolor{black}{Definition}}} \var{\deleted{\illegible{1}}}{}{deletion}{ps_134-5_056-v0014}1:  The functor i above is called a \hand{\uline{\textcolor{black}{pse\unsure{u|d}do-test}}} \hand{\uline{\textcolor{black}{functor}}} ,
if it satisfies the following three conditions:

a)  the corresponding functor  $i^{*}\colon M \hand{\to} A^\wedge$ is ``model-preserving'',  by

\margin{56}

%% END SOURCE ps_134-5_056
\ifnum\EditionMode=1
\psapparatusstart
\psentry{ps_134-5_056-v0001}{1}{typed insertion}{\inserted[typed]{strong}}{strong}{The typed insertion is incorporated at its indicated position; insertion brackets leave the running text.}
\psentry{ps_134-5_056-v0002}{2}{typewriter deletion}{\xout{\unsure{s}}}{}{The typewriter strikes out this material; it leaves the corrected running text.}
\psentry{ps_134-5_056-v0003}{3}{typewriter deletion}{\xout{the}}{}{The typewriter strikes out this material; it leaves the corrected running text.}
\psentry{ps_134-5_056-v0004}{4}{typewriter deletion}{\xout{eventual}}{}{The typewriter strikes out this material; it leaves the corrected running text.}
\psentry{ps_134-5_056-v0005}{5}{writing over}{$\over{\illegible{1}}{e}$}{$\hand{e}$}{The superposed handwritten reading replaces the underlying reading.}
\psentry{ps_134-5_056-v0006}{6}{typewriter deletion}{\xout{imply}}{}{The typewriter strikes out this material; it leaves the corrected running text.}
\psentry{ps_134-5_056-v0007}{7}{typed insertion}{\inserted[typed]{are}}{are}{The typed insertion is incorporated at its indicated position; insertion brackets leave the running text.}
\psentry{ps_134-5_056-v0008}{8}{typewriter deletion}{\xout{\unsure{(iii)}}}{}{The typewriter strikes out this material; it leaves the corrected running text.}
\psentry{ps_134-5_056-v0009}{9}{typewriter deletion}{\xout{The theorem}}{}{The typewriter strikes out this material; it leaves the corrected running text.}
\psentry{ps_134-5_056-v0010}{10}{typewriter overwriting}{\over[typed]{q}{a}}{a}{The superposed typed reading replaces the underlying reading.}
\psentry{ps_134-5_056-v0011}{11}{typed insertion}{\inserted[typed]{to be ``just it'' ,}}{to be ``just it'' ,}{The typed insertion is incorporated at its indicated position; insertion brackets leave the running text.}
\psentry{ps_134-5_056-v0012}{12}{writing over}{\over{\unsure{i}}{i}}{\hand{i}}{The superposed handwritten reading replaces the underlying reading.}
\psentry{ps_134-5_056-v0013}{13}{insertion}{\inserted{\deleted{aspheric}}}{\hand{}}{The handwritten insertion is incorporated at its indicated position; insertion brackets leave the running text.}
\psentry{ps_134-5_056-v0014}{14}{deletion}{\deleted{\illegible{1}}}{}{The author cancels this material by hand; it leaves the corrected running text.}
\psapparatusend
\fi
\pssource{ps_134-5_057}{134-5/57}{https://grothendieck.umontpellier.fr/web/viewer.html?file=../134-5.pdf\#page=58}
\psnumber{ps_134-5_057-v0001}{1}
\psnumber{ps_134-5_057-v0002}{2}
\psnumber{ps_134-5_057-v0003}{3}
\psnumber{ps_134-5_057-v0004}{4}
\psnumber{ps_134-5_057-v0005}{5}
\psnumber{ps_134-5_057-v0006}{6}
\psnumber{ps_134-5_057-v0007}{7}
\psnumber{ps_134-5_057-v0008}{8}
\psnumber{ps_134-5_057-v0009}{9}
\psnumber{ps_134-5_057-v0010}{10}
\psnumber{ps_134-5_057-v0011}{11}
%% BEGIN SOURCE ps_134-5_057
% Pursuing Stacks (cotes 134-2 to 134-8): archive cote 134-5, site page 57 (PDF page 58), typescript page 244.
% typescript. Draft compared: ps4.tex:2879-2945.
% AI-transcribed by Claude Opus 5.5 (claude-opus-5-5, adaptive thinking, effort high), Message Batches API, batch ps2 (msgbatch_01NVwr1Yuo8qJf5rykeGBXPK), 2026-11-12; reviewed against both supplied scan halves; see CHECK_DECISIONS.tsv.
\scanpage{134-5/57}
\mspage{244}

which is mean\hand{t} here that
\[ W \ =(i^{*})^{-1}(\underline{W}_A) \qquad , \]
\hand{\uline{\textcolor{black}{and}}} the induced functor
\[ \mathrm{Hot}_M = W^{-1}M \hand{\to} \mathrm{Hot}_A = \underline{W}_A{}^{-1}A^{\wedge} \]
is an equivalence.

b) The functor $i_A : A^{\wedge} \hand{\to} (\mathrm{Cat})$ is model preserving (for $\underline{W}_A,\underline{W}$ ) ,
which reduces to the canonical functor
\[ \mathrm{Hot}_A = \underline{W}_A{}^{-1}A^{\wedge} \hand{\longrightarrow} (\mathrm{Hot}) = \underline{W}^{-1}(\mathrm{Cat}) \]
being an equivalence (as we know already that $\underline{W}_A =(i_A)^{\hand{\unsure{*}}-1}(\underline{W})$ , by defi-%
nition of $\underline{W}_A$ ).

Condition b) just means that A is a pseudo-test category, i.e.
that \var{\xout{\unsure{tAx}}}{}{typewriter deletion}{ps_134-5_057-v0001} the asphericity structure $(A^{\wedge},A^{\wedge}_{as})$ it defines is modelizing.
By the corollary of lemma above, it implies that A is aspheric (which
corresponds to condition c) of def.5 in loc. cit.). Condition a)
comes in two parts, the first just meaning that i is $M_{as}$ -aspheric -
this translation being valid, at any rate, in case \var{\inserted[typed]{we assume already A aspheric, and}}{we assume already A aspheric, and}{typed insertion}{ps_134-5_057-v0002} the given aspherici\var{\xout{t}}{}{typewriter deletion}{ps_134-5_057-v0003}
ty structure is aspheric i.e. M admits a final object $e_M$ , and $e_M$ is
aspheric. \var{\inserted[typed]{We certainly do want a pseudo-test functor to be}}{We certainly do want a pseudo-test functor to be}{typed insertion}{ps_134-5_057-v0004} \var{\xout{At any rate, if we assume at least that $e_M$ exists, then}}{}{typewriter deletion}{ps_134-5_057-v0005}
(at the very least) $M_{as}$-aspheric, so we should either strengthen con-%
dition a) to this effect (which however does'nt look as nice), or assu%
-me beforehand \var{\xout{$e_M$ exists and}}{}{typewriter deletion}{ps_134-5_057-v0006} $(M,M_{as})$ aspheric. At any rate, if we use
the first variant\var{\xout{x}}{}{typewriter deletion}{ps_134-5_057-v0007} of the definition, condition a) in full then \var{\xout{reads}}{}{typewriter deletion}{ps_134-5_057-v0008}
is equ\var{\xout{\unsure{axx}}}{}{typewriter deletion}{ps_134-5_057-v0009}ivalent (granting b) ) to :

a') T\var{\over{\unsure{h}}{h}}{\hand{h}}{writing over}{ps_134-5_057-v0010}e functor i is $M_{as}$-aspheric, \hand{\uline{\textcolor{black}{and}}} $(M,M_{as})$ is modelizing ,
which in turn implies that $(M,M_{as})$ is aspheric. So we may as well assu-%
me asphericity of $(M,M_{as})$ beforehand ! At any rate, we see that the
notion we are after can be decomposed into three conditions, namely:

1) A is a pseudo-test category, i.e. $(A^{\wedge},A^{\wedge}_{as})$ is modelizing.

2) $(M,M_{as})$ is modelizing.

3) The functor i is $M_{as}$-aspheric.

T\var{\over{\unsure{h}}{h}}{\hand{h}}{writing over}{ps_134-5_057-v0011}e two first conditions are just conditions on A and on $(M,M_{as})$
respectively, the third is just the familiar asphericity condition\hand{ on i}.

To get the notion of a \hand{\uline{\textcolor{black}{weak}}} \hand{\uline{\textcolor{black}{test-functor}}} (def.5, p.175) we have to
be just one step more specific in 1), by demanding that A be even a

\margin{57}
% Angular hand marks lie above and below the minus sign of the exponent in condition b); the tentative asterisk reading remains doubtful.
% The cancelled axx fragment is displaced below the beginning of equivalent.

%% END SOURCE ps_134-5_057
\ifnum\EditionMode=1
\psapparatusstart
\psentry{ps_134-5_057-v0001}{1}{typewriter deletion}{\xout{\unsure{tAx}}}{}{The typewriter strikes out this material; it leaves the corrected running text.}
\psentry{ps_134-5_057-v0002}{2}{typed insertion}{\inserted[typed]{we assume already A aspheric, and}}{we assume already A aspheric, and}{The typed insertion is incorporated at its indicated position; insertion brackets leave the running text.}
\psentry{ps_134-5_057-v0003}{3}{typewriter deletion}{\xout{t}}{}{The typewriter strikes out this material; it leaves the corrected running text.}
\psentry{ps_134-5_057-v0004}{4}{typed insertion}{\inserted[typed]{We certainly do want a pseudo-test functor to be}}{We certainly do want a pseudo-test functor to be}{The typed insertion is incorporated at its indicated position; insertion brackets leave the running text.}
\psentry{ps_134-5_057-v0005}{5}{typewriter deletion}{\xout{At any rate, if we assume at least that $e_M$ exists, then}}{}{The typewriter strikes out this material; it leaves the corrected running text.}
\psentry{ps_134-5_057-v0006}{6}{typewriter deletion}{\xout{$e_M$ exists and}}{}{The typewriter strikes out this material; it leaves the corrected running text.}
\psentry{ps_134-5_057-v0007}{7}{typewriter deletion}{\xout{x}}{}{The typewriter strikes out this material; it leaves the corrected running text.}
\psentry{ps_134-5_057-v0008}{8}{typewriter deletion}{\xout{reads}}{}{The typewriter strikes out this material; it leaves the corrected running text.}
\psentry{ps_134-5_057-v0009}{9}{typewriter deletion}{\xout{\unsure{axx}}}{}{The typewriter strikes out this material; it leaves the corrected running text.}
\psentry{ps_134-5_057-v0010}{10}{writing over}{\over{\unsure{h}}{h}}{\hand{h}}{The superposed handwritten reading replaces the underlying reading.}
\psentry{ps_134-5_057-v0011}{11}{writing over}{\over{\unsure{h}}{h}}{\hand{h}}{The superposed handwritten reading replaces the underlying reading.}
\psapparatusend
\fi
\pssource{ps_134-5_058}{134-5/58}{https://grothendieck.umontpellier.fr/web/viewer.html?file=../134-5.pdf\#page=59}
\psnumber{ps_134-5_058-v0001}{1}
\psnumber{ps_134-5_058-v0002}{2}
\psnumber{ps_134-5_058-v0003}{3}
\psnumber{ps_134-5_058-v0004}{4}
\psnumber{ps_134-5_058-v0005}{5}
\psnumber{ps_134-5_058-v0006}{6}
\psnumber{ps_134-5_058-v0007}{7}
\psnumber{ps_134-5_058-v0008}{8}
\psnumber{ps_134-5_058-v0009}{9}
\psnumber{ps_134-5_058-v0010}{10}
\psnumber{ps_134-5_058-v0011}{11}
\psnumber{ps_134-5_058-v0012}{12}
%% BEGIN SOURCE ps_134-5_058
% Pursuing Stacks (cotes 134-2 to 134-8): archive cote 134-5, site page 58 (PDF page 59), typescript page 245.
% typescript. Draft compared: ps4.tex:2928-2996.
% AI-transcribed by Claude Opus 5.5 (claude-opus-5-5, adaptive thinking, effort high), Message Batches API, batch ps2 (msgbatch_01NVwr1Yuo8qJf5rykeGBXPK), 2026-11-12; reviewed against both supplied scan halves; see CHECK_DECISIONS.tsv.
\scanpage{134-5/58}
\mspage{245}

weak test category, namely that the functor
\[ i^{*}_{A} = j_{A} : (\mathrm{Cat}) \hand{\longrightarrow} A^{\wedge} \]
be model-preserving (which implies that $i_{A}$ is too).

Following def.6 (p.176), w'll say that $i$ is a \hand{\uline{\textcolor{black}{test}}} \hand{\uline{\textcolor{black}{functor}}} if $i$
\hand{\uline{\textcolor{black}{and}}} the induced functors $i_{/a} : A_{/a} \hand{\longrightarrow} A \hand{\longrightarrow} M$ are weak test functors.
Using the definition of a test category (def.3, p.173), we see that this
just means that the following conditions hold:

1') $A$ is a test category.

2') $(M,M_{as})$ is modelizing (i.e. same as 2) above).

3') The functor $i$ and the induced functors $i_{/a} : A_{/a} \hand{\rightarrow} M$ are \var{\xout{xxw}}{}{typewriter deletion}{ps_134-5_058-v0001}
$M_{as}$ - aspheric.

This last condition merits a name by itself, independently of the
other assumptions:

\hand{\uline{\textcolor{black}{Definition}}} 2. Let $(M,M_{as})$ be any asphericity structure, $A$ a small cate-%
gory, and $i: A \hand{\rightarrow} M$ a functor. We'll say that $i$ is \hand{\uline{\textcolor{black}{totally}}} $M_{as}$-\hand{\uline{\textcolor{black}{aspheric}}}
(or simply \hand{\uline{\textcolor{black}{totally}}} \hand{\uline{\textcolor{black}{aspheric}}} , if no confusion is feared), if $i$ and \var{\xout{th}}{}{typewriter deletion}{ps_134-5_058-v0002}
the induced functors $i_{a} : A_{/a} \hand{\longrightarrow} M$ are $M_{as}$-aspheric (for any $a$ in $A$),
We'll say that $i$ is \hand{\uline{\textcolor{black}{locally}}} $M_{as}$-\hand{\uline{\textcolor{black}{aspheric}}} (or simply \hand{\uline{\textcolor{black}{locally}}} \hand{\uline{\textcolor{black}{aspheric}}})
if for any $a$ in $A$, the induced functor $i_{/a} : A_{/a} \hand{\longrightarrow} M$ is $M_{as}$-aspheric.

Thus $i$ is totally aspheric iff it is aspheric and locally aspheric.
On the other hand, $i$ is a test functor iff $A$ is a test category, $(M,M_{as})$
is modelizing, and $i$ is totally aspheric.

In order for $i$ to be locally aspheric, it is ns that \var{\inserted[typed]{$i$ factor through $M_{as}$ and}}{$i$ factor through $M_{as}$ and}{typed insertion}{ps_134-5_058-v0003} for any
$x$ in $M_{as}$ , $i^{*}(x)$ be \hand{\uline{\textcolor{black}{aspheric over}}} $e_{A^{\wedge}}$ ; if $B$ is any \var{\xout{generating}}{}{typewriter deletion}{ps_134-5_058-v0004} subca-%
tegory of $M$ containing $i(A)$ and generating the asphericity structure,
it is enough in this latter condition to take $x$ in $B$ .

\var{\inserted{\uline{Remarks.}}}{\hand{\uline{Remarks.}}}{insertion}{ps_134-5_058-v0005}
Thus we see that the three gradations for the test-functor notion,
as suggested by the definitions 5 and 6 of section 65 and now by the
present context of asphericity structures, just amount to gradations \var{\xout{on}}{}{typewriter deletion}{ps_134-5_058-v0006}
for the test conditions on the category $A$ itself (\var{\inserted[typed]{namely}}{namely}{typed insertion}{ps_134-5_058-v0007} to be \var{\inserted[typed]{a}}{a}{typed insertion}{ps_134-5_058-v0008} pseudo-test, a
weak test or just a plain test -category), and a two-step gradation on
the asphericity condition for $i$ (namely, that $i$ be either just aspheric
in the two first cases, or totally aspheric in the third), \var{\inserted{while}}{\hand{while}}{insertion}{ps_134-5_058-v0009} these two
asphericity conditions on $i$ \var{\deleted{being}}{}{deletion}{ps_134-5_058-v0010} \var{\inserted{are}}{\hand{are}}{insertion}{ps_134-5_058-v0011} of significance, \var{\xout{xxxx}}{}{typewriter deletion}{ps_134-5_058-v0012} independently of
any specific assumption which we may make on either $A$ or $(M,M_{as})$ . This

\margin{58}

%% END SOURCE ps_134-5_058
\ifnum\EditionMode=1
\psapparatusstart
\psentry{ps_134-5_058-v0001}{1}{typewriter deletion}{\xout{xxw}}{}{The typewriter strikes out this material; it leaves the corrected running text.}
\psentry{ps_134-5_058-v0002}{2}{typewriter deletion}{\xout{th}}{}{The typewriter strikes out this material; it leaves the corrected running text.}
\psentry{ps_134-5_058-v0003}{3}{typed insertion}{\inserted[typed]{$i$ factor through $M_{as}$ and}}{$i$ factor through $M_{as}$ and}{The typed insertion is incorporated at its indicated position; insertion brackets leave the running text.}
\psentry{ps_134-5_058-v0004}{4}{typewriter deletion}{\xout{generating}}{}{The typewriter strikes out this material; it leaves the corrected running text.}
\psentry{ps_134-5_058-v0005}{5}{insertion}{\inserted{\uline{Remarks.}}}{\hand{\uline{Remarks.}}}{The handwritten insertion is incorporated at its indicated position; insertion brackets leave the running text.}
\psentry{ps_134-5_058-v0006}{6}{typewriter deletion}{\xout{on}}{}{The typewriter strikes out this material; it leaves the corrected running text.}
\psentry{ps_134-5_058-v0007}{7}{typed insertion}{\inserted[typed]{namely}}{namely}{The typed insertion is incorporated at its indicated position; insertion brackets leave the running text.}
\psentry{ps_134-5_058-v0008}{8}{typed insertion}{\inserted[typed]{a}}{a}{The typed insertion is incorporated at its indicated position; insertion brackets leave the running text.}
\psentry{ps_134-5_058-v0009}{9}{insertion}{\inserted{while}}{\hand{while}}{The handwritten insertion is incorporated at its indicated position; insertion brackets leave the running text.}
\psentry{ps_134-5_058-v0010}{10}{deletion}{\deleted{being}}{}{The author cancels this material by hand; it leaves the corrected running text.}
\psentry{ps_134-5_058-v0011}{11}{insertion}{\inserted{are}}{\hand{are}}{The handwritten insertion is incorporated at its indicated position; insertion brackets leave the running text.}
\psentry{ps_134-5_058-v0012}{12}{typewriter deletion}{\xout{xxxx}}{}{The typewriter strikes out this material; it leaves the corrected running text.}
\psapparatusend
\fi
\pssource{ps_134-5_059}{134-5/59}{https://grothendieck.umontpellier.fr/web/viewer.html?file=../134-5.pdf\#page=60}
\psnumber{ps_134-5_059-v0001}{1}
\psnumber{ps_134-5_059-v0002}{2}
\psnumber{ps_134-5_059-v0003}{3}
\psnumber{ps_134-5_059-v0004}{4}
\psnumber{ps_134-5_059-v0005}{5}
\psnumber{ps_134-5_059-v0006}{6}
\psnumber{ps_134-5_059-v0007}{7}
\psnumber{ps_134-5_059-v0008}{8}
\psnumber{ps_134-5_059-v0009}{9}
\psnumber{ps_134-5_059-v0010}{10}
%% BEGIN SOURCE ps_134-5_059
% Pursuing Stacks (cotes 134-2 to 134-8): archive cote 134-5, site page 59 (PDF page 60), typescript page 246.
% typescript. Draft compared: ps4.tex:2984-3036.
% AI-transcribed by Claude Opus 5.5 (claude-opus-5-5, adaptive thinking, effort high), Message Batches API, batch ps2 (msgbatch_01NVwr1Yuo8qJf5rykeGBXPK), 2026-11-12; reviewed against both supplied scan halves; see CHECK_DECISIONS.tsv.
\scanpage{134-5/59}
\mspage{246}

seems to diminish somewhat the \var{\xout{ax}}{}{typewriter deletion}{ps_134-5_059-v0001} emphasis I had put formerly upon the
notion of a test functor and it's weak variant, and enhance accordingly
the notion\var{\deleted{s}}{}{deletion}{ps_134-5_059-v0002} of an aspheric functor (with respect to a given asphericity
structure) and the two related asphericity notions for i which spring
from it (namely the notions of locally and \var{\inserted[typed]{of}}{of}{typed insertion}{ps_134-5_059-v0003} totally aspheric functors),
which now seem to come out as the more relevant and the more general ones.

To be wholly happy, we still need the relevant reformulation, in
terms of asphericity structures, of the main result of section 65, name-%
ly theorem 1 (p.177) characterizing test functors with values in (Cat),
under the assumption that the objects i(a) be contractible. This theorem,
I recall, has been the main outcome of the ``grimding'' reflections taking
their start with the``observation'' on p.94 about ten days earlier. The
approach then followed, as well as the contractibility assumption \var{\over{on}{for}}{\hand{for}}{writing over}{ps_134-5_059-v0004} the
objects i(a) made in the theorem, retrospectively look awkward - it is
clear that the relevant notions of aspher\var{\over[typed]{\unsure{ci}}{ic}}{ic}{typewriter overwriting}{ps_134-5_059-v0005}ity structures, and of asphe-%
ric functors into these, had been lacking. The theorem stated last
(p. 242) looks indeed a lot more satisfactory than the former, except
however in one respect - namely that the asphericity condition on i ,
in theorem 1 of p.177, can be expressed \var{\xout{ax}}{}{typewriter deletion}{ps_134-5_059-v0006} by asphericity over $e_{A}{}^{\wedge}$ \hand{\uline{\textcolor{black}{of}}}
\hand{\uline{\textcolor{black}{just}}} $i^{*}(\hand{\Delta_{1}})$ \hand{\uline{\textcolor{black}{alone}}}, rather than having to take $i^{*}(x)$ for all x in some
fixed generating subcategory of $(M,M_{as})$ containing i(A). To recover such
kind of \var{\inserted[typed]{minimal}}{minimal}{typed insertion}{ps_134-5_059-v0007} criterion in a more general case than (Cat), we'll have to rela-%
te the notion of an asphericity structure with the earlier one of a
\hand{\uline{\textcolor{black}{contractibility}}} \hand{\uline{\textcolor{black}{structure}}} ; \var{\xout{which}}{}{typewriter deletion}{ps_134-5_059-v0008} \var{\inserted[typed]{th\over{\unsure{is}}{e}}}{th\hand{e}}{typed insertion}{ps_134-5_059-v0009} \var{\inserted{latter}}{\hand{latter}}{insertion}{ps_134-5_059-v0010} in our present reflections has faded
somewhat into the background, while at an earlier stage the homotopy and
contractibility notions had invaded the picture to an extent as to over-%
shadow and nearly bring to oblivion the magic of the ``asphericity game''.

% NOTE: the lower part of the page is hidden by a sheet laid over it (torn edge); only the top halves of the letters of the next line are visible
\unsure{We'll come back upon contractibility structures in the next section}

\margin{59}
% Overlapping type impressions remain in the before two related and in the final digit of p.94; the visible final readings are retained.
% Small displaced marks around h occur in The approach and The theorem; their correction function remains unresolved.
% The final line is largely concealed by the laid-over sheet and has not been reconstructed from the witness.

%% END SOURCE ps_134-5_059
\ifnum\EditionMode=1
\psapparatusstart
\psentry{ps_134-5_059-v0001}{1}{typewriter deletion}{\xout{ax}}{}{The typewriter strikes out this material; it leaves the corrected running text.}
\psentry{ps_134-5_059-v0002}{2}{deletion}{\deleted{s}}{}{The author cancels this material by hand; it leaves the corrected running text.}
\psentry{ps_134-5_059-v0003}{3}{typed insertion}{\inserted[typed]{of}}{of}{The typed insertion is incorporated at its indicated position; insertion brackets leave the running text.}
\psentry{ps_134-5_059-v0004}{4}{writing over}{\over{on}{for}}{\hand{for}}{The superposed handwritten reading replaces the underlying reading.}
\psentry{ps_134-5_059-v0005}{5}{typewriter overwriting}{\over[typed]{\unsure{ci}}{ic}}{ic}{The superposed typed reading replaces the underlying reading.}
\psentry{ps_134-5_059-v0006}{6}{typewriter deletion}{\xout{ax}}{}{The typewriter strikes out this material; it leaves the corrected running text.}
\psentry{ps_134-5_059-v0007}{7}{typed insertion}{\inserted[typed]{minimal}}{minimal}{The typed insertion is incorporated at its indicated position; insertion brackets leave the running text.}
\psentry{ps_134-5_059-v0008}{8}{typewriter deletion}{\xout{which}}{}{The typewriter strikes out this material; it leaves the corrected running text.}
\psentry{ps_134-5_059-v0009}{9}{typed insertion}{\inserted[typed]{th\over{\unsure{is}}{e}}}{th\hand{e}}{The typed insertion is incorporated at its indicated position; insertion brackets leave the running text.}
\psentry{ps_134-5_059-v0010}{10}{insertion}{\inserted{latter}}{\hand{latter}}{The handwritten insertion is incorporated at its indicated position; insertion brackets leave the running text.}
\psapparatusend
\fi
\pssource{ps_134-5_060}{134-5/60}{https://grothendieck.umontpellier.fr/web/viewer.html?file=../134-5.pdf\#page=61}
\psnumber{ps_134-5_060-v0001}{1}
\psnumber{ps_134-5_060-v0002}{2}
\psnumber{ps_134-5_060-v0003}{3}
\psnumber{ps_134-5_060-v0004}{4}
\psnumber{ps_134-5_060-v0005}{5}
\psnumber{ps_134-5_060-v0006}{6}
%% BEGIN SOURCE ps_134-5_060
% Pursuing Stacks (cotes 134-2 to 134-8): archive cote 134-5, site page 60 (PDF page 61), typescript page 247.
% typescript. Draft compared: ps4.tex:3018-3090.
% AI-transcribed by Claude Opus 5.5 (claude-opus-5-5, adaptive thinking, effort high), Message Batches API, batch ps2 (msgbatch_01NVwr1Yuo8qJf5rykeGBXPK), 2026-11-12; reviewed against both supplied scan halves; see CHECK_DECISIONS.tsv.
\scanpage{134-5/60}
\mspage{247}
20.6.

79. I \unsure{would} like now to write out the relationship between asphericity
structures, and contractibility structures as defined in section 51 D)
(pages 117-119). First we'll need to rid ourselves of the smallness as\var{\xout{\unsure{s}}}{}{typewriter deletion}{ps_134-5_060-v0001}%
sumption for a generating category of an asphericity structure:

\hand{\uline{\textcolor{black}{Proposition}1}} Let M be a $\underline{U}$-category ($\underline{U}$ being our basic universe), and N
any full subcategory. The following conditions are e\var{\over{\unsure{qm}}{qu}}{\hand{qu}}{writing over}{ps_134-5_060-v0002}ivalent:

(i) There exists an asphericity structure $M_{as}$ in M such that
(a) $M_{as}\hand{\supset} N$ and (b) N admits a \hand{\uline{\textcolor{black}{small}}} full subcategory\var{\inserted{$N_o$}}{\hand{$N_o$}}{insertion}{ps_134-5_060-v0003} generating
the asphericity structure, i.e. such that the inclusion functor $i_o$ from
$N_o$ to M be $M_{as}$-aspheric.

(ii) There exists a small full subcategory $N_o$ of N, such that for
any x in N, $i_o^{*}(x)$ in $N_o^{\wedge}$ be aspheric, where $i_o : N_o \hand{\to} M$ is the inclusion
functor.

(iii) The couple (N,N) is an asphericity structure.

Moreover, when these conditions are satisfied, the asphericity
structure $M_{as}$ in (i) is unique.

Of course, we'll say it is the asphericity structure on M \hand{\uline{\textcolor{black}{genera-%
ted}}} by the full subcategory N of M, and the latter will be called a
\hand{\uline{\textcolor{black}{generating}}} \hand{\uline{\textcolor{black}{subcategory}}} for the given asphericity structure.

The proposition is a tautology, in view of the definitions and of
prop.1 (p. 232) and it's corollary. The fo\unsure{r|o}m (ii) or (iii) of the condi-%
tion shows that it depends only upon the category N, not upon the way
in which N is embedded in a larger category M. We may call a category N
satisfying condition (iii) above an \hand{\uline{\textcolor{black}{aspherator}}} , by which we would like
to express that this category represents a standart way of generating
asphericity structures, through \var{\deleted{the}}{}{deletion}{ps_134-5_060-v0004}\var{\inserted[typed]{any full}}{any full}{typed insertion}{ps_134-5_060-v0005} embedding of N into a category M.
This condition is automatically satisfied if N is small, more generally
it holds if N is equivalent to a small category. It should be kept in
mind that the condition depends both on the choice of \var{\over{\unsure{th}}{th}}{\hand{th}}{writing over}{ps_134-5_060-v0006}e basic universe
$\underline{U}$ and on the choice of the basic localizer $\underline{W}$ in the corresponding large
category (Cat). It looks pretty sure that the condition is not always
satisfied (I doubt it is for N = (Sets) say), but I confess I did'nt
sit down to make an explicit example.

\margin{60}

%% END SOURCE ps_134-5_060
\ifnum\EditionMode=1
\psapparatusstart
\psentry{ps_134-5_060-v0001}{1}{typewriter deletion}{\xout{\unsure{s}}}{}{The typewriter strikes out this material; it leaves the corrected running text.}
\psentry{ps_134-5_060-v0002}{2}{writing over}{\over{\unsure{qm}}{qu}}{\hand{qu}}{The superposed handwritten reading replaces the underlying reading.}
\psentry{ps_134-5_060-v0003}{3}{insertion}{\inserted{$N_o$}}{\hand{$N_o$}}{The handwritten insertion is incorporated at its indicated position; insertion brackets leave the running text.}
\psentry{ps_134-5_060-v0004}{4}{deletion}{\deleted{the}}{}{The author cancels this material by hand; it leaves the corrected running text.}
\psentry{ps_134-5_060-v0005}{5}{typed insertion}{\inserted[typed]{any full}}{any full}{The typed insertion is incorporated at its indicated position; insertion brackets leave the running text.}
\psentry{ps_134-5_060-v0006}{6}{writing over}{\over{\unsure{th}}{th}}{\hand{th}}{The superposed handwritten reading replaces the underlying reading.}
\psapparatusend
\fi
\pssource{ps_134-5_061}{134-5/61}{https://grothendieck.umontpellier.fr/web/viewer.html?file=../134-5.pdf\#page=62}
\psnumber{ps_134-5_061-v0001}{1}
\psnumber{ps_134-5_061-v0002}{2}
\psnumber{ps_134-5_061-v0003}{3}
\psnumber{ps_134-5_061-v0004}{4}
\psnumber{ps_134-5_061-v0005}{5}
\psnumber{ps_134-5_061-v0006}{6}
\psnumber{ps_134-5_061-v0007}{7}
\psnumber{ps_134-5_061-v0008}{8}
\psnumber{ps_134-5_061-v0009}{9}
\psnumber{ps_134-5_061-v0010}{10}
\psnumber{ps_134-5_061-v0011}{11}
\psnumber{ps_134-5_061-v0012}{12}
\psnumber{ps_134-5_061-v0013}{13}
\psnumber{ps_134-5_061-v0014}{14}
\psnumber{ps_134-5_061-v0015}{15}
\psnumber{ps_134-5_061-v0016}{16}
\psnumber{ps_134-5_061-v0017}{17}
\psnumber{ps_134-5_061-v0018}{18}
\psnumber{ps_134-5_061-v0019}{19}
\psnumber{ps_134-5_061-v0020}{20}
\psnumber{ps_134-5_061-v0021}{21}
\psnumber{ps_134-5_061-v0022}{22}
\psnumber{ps_134-5_061-v0023}{23}
\psnumber{ps_134-5_061-v0024}{24}
\psnumber{ps_134-5_061-v0025}{25}
\psnumber{ps_134-5_061-v0026}{26}
\psnumber{ps_134-5_061-v0027}{27}
\psnumber{ps_134-5_061-v0028}{28}
\psnumber{ps_134-5_061-v0029}{29}
\psnumber{ps_134-5_061-v0030}{30}
%% BEGIN SOURCE ps_134-5_061
% Pursuing Stacks (cotes 134-2 to 134-8): archive cote 134-5, site page 61 (PDF page 62), typescript page 248.
% typescript. Draft compared: ps4.tex:3076-3140.
% AI-transcribed by Claude Opus 5.5 (claude-opus-5-5, adaptive thinking, effort high), Message Batches API, batch ps2 (msgbatch_01NVwr1Yuo8qJf5rykeGBXPK), 2026-11-12; reviewed against both supplied scan halves; see CHECK_DECISIONS.tsv.
\scanpage{134-5/61}
\mspage{\var{\deleted{248}}{}{deletion}{ps_134-5_061-v0001} \hand{248}}

\var{\xout{Max}}{}{typewriter deletion}{ps_134-5_061-v0002} Let now
\var{\inserted[typed]{$(M,M_c)$}}{$(M,M_c)$}{typed insertion}{ps_134-5_061-v0003} \var{\inserted[typed]{be a}}{be a}{typed insertion}{ps_134-5_061-v0004}\var{\xout{\hand{\uline{\textcolor{black}{Proposition}}} 2\unsure{.} Let $(M,M_c)$ be a }}{}{typewriter deletion}{ps_134-5_061-v0005}contractibility structure, such that
there exist a \hand{\uline{\textcolor{black}{small}}} \var{\inserted[typed]{full}}{full}{typed insertion}{ps_134-5_061-v0006} subcategory $C$ of $M$ \var{\xout{whosex\allowbreak{}objectsx\allowbreak{}arex\allowbreak{}contractible}}{}{typewriter deletion}{ps_134-5_061-v0007}
\var{\over{W}{w}}{\hand{w}}{writing over}{ps_134-5_061-v0008}hich generates the contractibility structure, i.e. such that \var{\xout{th}}{}{typewriter deletion}{ps_134-5_061-v0009}
(a) the objects of $C$ are ``contractible'' i.e. are in $M_c$ and (b) any
object of $M_c$ \var{\deleted{which}}{}{deletion}{ps_134-5_061-v0010} is $C$-contractible\var{\deleted{\unsure{s}}}{}{deletion}{ps_134-5_061-v0011} (i.e. contractible for the homotopy
interval structure admitting the intervals made up with objects of $C$ as
a generating family)\hand{,} \var{\deleted{is in $M_c$}}{}{deletion}{ps_134-5_061-v0012} . Independently of any smallness assump-%
tion upon $C$, we gave \var{\over{\unsure{in}}{in}}{\hand{in}}{writing over}{ps_134-5_061-v0013} section 51 D) \var{\xout{the}}{}{typewriter deletion}{ps_134-5_061-v0014} (p.118), under the name
of ``basic assumption'' (Bas 4), the ns condition on a full subcategory
$C$ of a given category $M$, in order that $C$ can be viewed as a generating
set of contractible objects, for a suitable contractibility structure
\var{\xout{Mx XX}}{}{typewriter deletion}{ps_134-5_061-v0015}
\[ M_c \hand{\subset} \mathrm{Ob} M \]
in $M$ (which is uniquely defined by $C$). It turns out that in case $C$ \var{\xout{co}}{}{typewriter deletion}{ps_134-5_061-v0016}con-%
\var{\xout{stable}}{}{typewriter deletion}{ps_134-5_061-v0017} tains a final object of $M$, and is stable \var{\inserted[typed]{under}}{under}{typed insertion}{ps_134-5_061-v0018} \var{\deleted{by}}{}{deletion}{ps_134-5_061-v0019} binary products in
$M$, the condition (Bas 4) \var{\xout{doesxnot}}{}{typewriter deletion}{ps_134-5_061-v0020} depends only \var{\inserted[typed]{upon}}{upon}{typed insertion}{ps_134-5_061-v0021} \var{\deleted{on}}{}{deletion}{ps_134-5_061-v0022} the category structure
of $C$ , and not upon the particular way this category is embedded in
an\var{\over{o}{o}}{\hand{o}}{writing over}{ps_134-5_061-v0023}ther one $M$.

In loc. cit. we did not impose, when defining a contractibility
structure, a condition that the\var{\over{n}{r}}{\hand{r}}{writing over}{ps_134-5_061-v0024} should exist a small set of generators
for the structure. From now on, we'll assume that the (possibly large)
categories $M$ we are working with are $\underline{U}$-categories, and that ``contracti-%
bility structure'' means ``$\underline{U}$-contractibility structure'', namely one \var{\inserted[typed]{such that}}{such that}{typed insertion}{ps_134-5_061-v0025} \var{\deleted{admit-}}{}{deletion}{ps_134-5_061-v0026}
\var{\inserted[typed]{M admit}}{M admit}{typed insertion}{ps_134-5_061-v0027} \var{\deleted{ting}}{}{deletion}{ps_134-5_061-v0028} a small full subcategory $C$, generating the structure.

We recall too that in the definition of a contractibility structure
$(M,M_c)$, it has always been understood that $M$ is stable under finite pro-%
ducts.

\hand{\uline{\textcolor{black}{Proposition}}} 2. With the conventions above, let $(M,M_c)$ be any contrac-%
tibility structure. Then:

a) The full subcategory $M_c$ of $M$ generates an asphericity structure
$M_{as}$ in $M$ (cf. prop.1).

b) Any small full subcategory $C$ of $M$ which generates the contractibi
lity structure\var{\inserted{and such that $C^{\wedge}$ be totally aspheric,}}{\hand{and such that $C^{\wedge}$ be totally aspheric,}}{insertion}{ps_134-5_061-v0029} \hand{\uline{\textcolor{black}{generates}}} the asphericity structure $M_{as}$\hand{,} i.e. the inclu-%
sion functor $i: C \hand{\to} M$ is $M_{as}$-aspheric.

c) The asphericity structure $M_{as}$ is totally aspheric (cf. \var{\over{P}{p}}{\hand{p}}{writing over}{ps_134-5_061-v0030}rop.5
p. 239).

\margin{61}

%% END SOURCE ps_134-5_061
\ifnum\EditionMode=1
\psapparatusstart
\psentry{ps_134-5_061-v0001}{1}{deletion}{\deleted{248}}{}{The author cancels this material by hand; it leaves the corrected running text.}
\psentry{ps_134-5_061-v0002}{2}{typewriter deletion}{\xout{Max}}{}{The typewriter strikes out this material; it leaves the corrected running text.}
\psentry{ps_134-5_061-v0003}{3}{typed insertion}{\inserted[typed]{$(M,M_c)$}}{$(M,M_c)$}{The typed insertion is incorporated at its indicated position; insertion brackets leave the running text.}
\psentry{ps_134-5_061-v0004}{4}{typed insertion}{\inserted[typed]{be a}}{be a}{The typed insertion is incorporated at its indicated position; insertion brackets leave the running text.}
\psentry{ps_134-5_061-v0005}{5}{typewriter deletion}{\xout{\hand{\uline{\textcolor{black}{Proposition}}} 2\unsure{.} Let $(M,M_c)$ be a }}{}{The typewriter strikes out this material; it leaves the corrected running text.}
\psentry{ps_134-5_061-v0006}{6}{typed insertion}{\inserted[typed]{full}}{full}{The typed insertion is incorporated at its indicated position; insertion brackets leave the running text.}
\psentry{ps_134-5_061-v0007}{7}{typewriter deletion}{\xout{whosex\allowbreak{}objectsx\allowbreak{}arex\allowbreak{}contractible}}{}{The typewriter strikes out this material; it leaves the corrected running text.}
\psentry{ps_134-5_061-v0008}{8}{writing over}{\over{W}{w}}{\hand{w}}{The superposed handwritten reading replaces the underlying reading.}
\psentry{ps_134-5_061-v0009}{9}{typewriter deletion}{\xout{th}}{}{The typewriter strikes out this material; it leaves the corrected running text.}
\psentry{ps_134-5_061-v0010}{10}{deletion}{\deleted{which}}{}{The author cancels this material by hand; it leaves the corrected running text.}
\psentry{ps_134-5_061-v0011}{11}{deletion}{\deleted{\unsure{s}}}{}{The author cancels this material by hand; it leaves the corrected running text.}
\psentry{ps_134-5_061-v0012}{12}{deletion}{\deleted{is in $M_c$}}{}{The author cancels this material by hand; it leaves the corrected running text.}
\psentry{ps_134-5_061-v0013}{13}{writing over}{\over{\unsure{in}}{in}}{\hand{in}}{The superposed handwritten reading replaces the underlying reading.}
\psentry{ps_134-5_061-v0014}{14}{typewriter deletion}{\xout{the}}{}{The typewriter strikes out this material; it leaves the corrected running text.}
\psentry{ps_134-5_061-v0015}{15}{typewriter deletion}{\xout{Mx XX}}{}{The typewriter strikes out this material; it leaves the corrected running text.}
\psentry{ps_134-5_061-v0016}{16}{typewriter deletion}{\xout{co}}{}{The typewriter strikes out this material; it leaves the corrected running text.}
\psentry{ps_134-5_061-v0017}{17}{typewriter deletion}{\xout{stable}}{}{The typewriter strikes out this material; it leaves the corrected running text.}
\psentry{ps_134-5_061-v0018}{18}{typed insertion}{\inserted[typed]{under}}{under}{The typed insertion is incorporated at its indicated position; insertion brackets leave the running text.}
\psentry{ps_134-5_061-v0019}{19}{deletion}{\deleted{by}}{}{The author cancels this material by hand; it leaves the corrected running text.}
\psentry{ps_134-5_061-v0020}{20}{typewriter deletion}{\xout{doesxnot}}{}{The typewriter strikes out this material; it leaves the corrected running text.}
\psentry{ps_134-5_061-v0021}{21}{typed insertion}{\inserted[typed]{upon}}{upon}{The typed insertion is incorporated at its indicated position; insertion brackets leave the running text.}
\psentry{ps_134-5_061-v0022}{22}{deletion}{\deleted{on}}{}{The author cancels this material by hand; it leaves the corrected running text.}
\psentry{ps_134-5_061-v0023}{23}{writing over}{\over{o}{o}}{\hand{o}}{The superposed handwritten reading replaces the underlying reading.}
\psentry{ps_134-5_061-v0024}{24}{writing over}{\over{n}{r}}{\hand{r}}{The superposed handwritten reading replaces the underlying reading.}
\psentry{ps_134-5_061-v0025}{25}{typed insertion}{\inserted[typed]{such that}}{such that}{The typed insertion is incorporated at its indicated position; insertion brackets leave the running text.}
\psentry{ps_134-5_061-v0026}{26}{deletion}{\deleted{admit-}}{}{The author cancels this material by hand; it leaves the corrected running text.}
\psentry{ps_134-5_061-v0027}{27}{typed insertion}{\inserted[typed]{M admit}}{M admit}{The typed insertion is incorporated at its indicated position; insertion brackets leave the running text.}
\psentry{ps_134-5_061-v0028}{28}{deletion}{\deleted{ting}}{}{The author cancels this material by hand; it leaves the corrected running text.}
\psentry{ps_134-5_061-v0029}{29}{insertion}{\inserted{and such that $C^{\wedge}$ be totally aspheric,}}{\hand{and such that $C^{\wedge}$ be totally aspheric,}}{The handwritten insertion is incorporated at its indicated position; insertion brackets leave the running text.}
\psentry{ps_134-5_061-v0030}{30}{writing over}{\over{P}{p}}{\hand{p}}{The superposed handwritten reading replaces the underlying reading.}
\psapparatusend
\fi
\pssource{ps_134-5_062}{134-5/62}{https://grothendieck.umontpellier.fr/web/viewer.html?file=../134-5.pdf\#page=63}
\psnumber{ps_134-5_062-v0001}{1}
\psnumber{ps_134-5_062-v0002}{2}
\psnumber{ps_134-5_062-v0003}{3}
\psnumber{ps_134-5_062-v0004}{4}
\psnumber{ps_134-5_062-v0005}{5}
\psnumber{ps_134-5_062-v0006}{6}
\psnumber{ps_134-5_062-v0007}{7}
\psnumber{ps_134-5_062-v0008}{8}
\psnumber{ps_134-5_062-v0009}{9}
\psnumber{ps_134-5_062-v0010}{10}
\psnumber{ps_134-5_062-v0011}{11}
\psnumber{ps_134-5_062-v0012}{12}
\psnumber{ps_134-5_062-v0013}{13}
%% BEGIN SOURCE ps_134-5_062
% Pursuing Stacks (cotes 134-2 to 134-8): archive cote 134-5, site page 62 (PDF page 63), typescript page 249.
% typescript. Draft compared: ps4.tex:3125-3188.
% AI-transcribed by Claude Opus 5.5 (claude-opus-5-5, adaptive thinking, effort high), Message Batches API, batch ps2 (msgbatch_01NVwr1Yuo8qJf5rykeGBXPK), 2026-11-12; reviewed against both supplied scan halves; see CHECK_DECISIONS.tsv.
\scanpage{134-5/62}
\mspage{249}

\var{\xout{Condition}}{}{typewriter deletion}{ps_134-5_062-v0001} The first statement a) can be rephrased,  by saying that
\var{\xout{Mx}}{}{typewriter deletion}{ps_134-5_062-v0002} the category $M_c$ of contractible objects of $M$ is an \hand{\uline{\textcolor{black}{aspherator}}} .  To
prove this,  we use the fact that $M_c$ is stable under finite products,
which implies that we can find a small full subcategory \var{\xout{of}}{}{typewriter deletion}{ps_134-5_062-v0003} $C$ of $M_c$ ,
stable under such products,  and which generates the contractibility
structure.  The stability condition upon $C^{\unsure{*}}$ implies that $C^{\wedge}$\hand{/}is totally
aspheric.  Therefore a) and b) will be proved,  if we prove that any
small full subcategory $C$ of $M_c$ ,  such that $C^{\wedge}$ is totally aspheric,
satisfies the conditions of prop.1 (with $N=M_c$) ,  namely that for any
$x$ in $M_c$ ,  $i^{*}(x)$ is an aspheric object of $C^{\wedge}$,  where $i\colon C\hand{\to} M$  (or $C\hand{\to} M_c$,
equivalently) is the inclusion functor.  But from the fact that $i^{*}$ commu-%
tes with finite products,  it follows that $i^{*}(x)$  is contractible in $C^{\wedge}$,
for the \var{\xout{homotopy}}{}{typewriter deletion}{ps_134-5_062-v0004} homotopy interval structure admitting as a generating
family of homotopy intervals the intervals in $C^{\wedge}$ made up with objects of
$C$ .  The assumption of total asphericity upon $C^{\wedge}$ implies that these
homotopy intervals are aspheric over $e_{C^{\wedge}}$ ,  and from this follows \var{\xout{that}}{}{typewriter deletion}{ps_134-5_062-v0005}
(as already used a number of times earlier)  that $i^{*}(x)$ \var{\inserted[typed]{too}}{too}{typed insertion}{ps_134-5_062-v0006} is aspheric
over $e_{C^{\wedge}}$ ,  and hence aspheric as $e_{C^{\wedge}}$ is aspheric (because of the assump-%
tion of total asphericity).  This proves a) and b),  and c) follows via
prop. 5 (p. 239).

We'll call of course the asphericity structure described in prop.2
the \hand{\uline{\textcolor{black}{asphericity}}} \hand{\uline{\textcolor{black}{structure}}} \hand{\uline{\textcolor{black}{generated}}} \hand{\uline{\textcolor{black}{by the given contractibility struc\var{\xout{tur}}{}{typewriter deletion}{ps_134-5_062-v0007}}}}
\hand{\uline{\textcolor{black}{ture}}}.
\hand{\uline{\textcolor{black}{Remark\unsure{s}}}}\var{\inserted{1}}{\hand{1}}{insertion}{ps_134-5_062-v0008} Clearly,  not every asphericity structure can be generated by a
contractibility structure,  as a necessary condition (presumably not a
sufficient one) is total asphericity.  We don't expect either,  in case
\margin{X} it can,  that the \var{\inserted[typed]{generating}}{generating}{typed insertion}{ps_134-5_062-v0009} contractibility structure is uniquely defined;  howe-%
ver,  we do expect in this case that there should exist a canonical
(largest) choice - we\var{\over{\unsure{ll}}{'ll}}{\hand{'ll}}{writing over}{ps_134-5_062-v0010} have to come back upon this in due course. For
the time being,  let's only remark that all modelizing asphericity struc-%
tures met \var{\inserted{with}}{\hand{with}}{insertion}{ps_134-5_062-v0011} so far,  it seems,  do come from contractibility structures.

It occurs to me that \var{\xout{in}}{}{typewriter deletion}{ps_134-5_062-v0012} the last statement  was a little hasty  -
after all we have met with test categories which are not strict ones,
hence ($A^{\wedge}_{\var{\xout{\illegible{1}}}{}{typewriter deletion}{ps_134-5_062-v0013}}$ , $A^{\wedge}_{as}$ )  is a modelizing asphericity structure  (it would be
enough even that $A$ be a pseudo-test category),  which is'nt totally asphe
ric,  and a fortiori does not come from a contractibility structure. \hand{Thus}

\margin{62}
% The long diagonal pen stroke beside C-hat has no visible arrowhead; its function remains unresolved.
% The early cancelled Mx cluster also extends below its baseline with cancellation type; no further old letters are securely recoverable.

%% END SOURCE ps_134-5_062
\ifnum\EditionMode=1
\psapparatusstart
\psentry{ps_134-5_062-v0001}{1}{typewriter deletion}{\xout{Condition}}{}{The typewriter strikes out this material; it leaves the corrected running text.}
\psentry{ps_134-5_062-v0002}{2}{typewriter deletion}{\xout{Mx}}{}{The typewriter strikes out this material; it leaves the corrected running text.}
\psentry{ps_134-5_062-v0003}{3}{typewriter deletion}{\xout{of}}{}{The typewriter strikes out this material; it leaves the corrected running text.}
\psentry{ps_134-5_062-v0004}{4}{typewriter deletion}{\xout{homotopy}}{}{The typewriter strikes out this material; it leaves the corrected running text.}
\psentry{ps_134-5_062-v0005}{5}{typewriter deletion}{\xout{that}}{}{The typewriter strikes out this material; it leaves the corrected running text.}
\psentry{ps_134-5_062-v0006}{6}{typed insertion}{\inserted[typed]{too}}{too}{The typed insertion is incorporated at its indicated position; insertion brackets leave the running text.}
\psentry{ps_134-5_062-v0007}{7}{typewriter deletion}{\xout{tur}}{}{The typewriter strikes out this material; it leaves the corrected running text.}
\psentry{ps_134-5_062-v0008}{8}{insertion}{\inserted{1}}{\hand{1}}{The handwritten insertion is incorporated at its indicated position; insertion brackets leave the running text.}
\psentry{ps_134-5_062-v0009}{9}{typed insertion}{\inserted[typed]{generating}}{generating}{The typed insertion is incorporated at its indicated position; insertion brackets leave the running text.}
\psentry{ps_134-5_062-v0010}{10}{writing over}{\over{\unsure{ll}}{'ll}}{\hand{'ll}}{The superposed handwritten reading replaces the underlying reading.}
\psentry{ps_134-5_062-v0011}{11}{insertion}{\inserted{with}}{\hand{with}}{The handwritten insertion is incorporated at its indicated position; insertion brackets leave the running text.}
\psentry{ps_134-5_062-v0012}{12}{typewriter deletion}{\xout{in}}{}{The typewriter strikes out this material; it leaves the corrected running text.}
\psentry{ps_134-5_062-v0013}{13}{typewriter deletion}{$\xout{\illegible{1}}$}{}{The typewriter strikes out this material; it leaves the corrected running text.}
\psapparatusend
\fi
\pssource{ps_134-5_063}{134-5/63}{https://grothendieck.umontpellier.fr/web/viewer.html?file=../134-5.pdf\#page=64}
\psnumber{ps_134-5_063-v0001}{1}
\psnumber{ps_134-5_063-v0002}{2}
\psnumber{ps_134-5_063-v0003}{3}
\psnumber{ps_134-5_063-v0004}{4}
\psnumber{ps_134-5_063-v0005}{5}
\psnumber{ps_134-5_063-v0006}{6}
\psnumber{ps_134-5_063-v0007}{7}
\psnumber{ps_134-5_063-v0008}{8}
\psnumber{ps_134-5_063-v0009}{9}
\psnumber{ps_134-5_063-v0010}{10}
\psnumber{ps_134-5_063-v0011}{11}
\psnumber{ps_134-5_063-v0012}{12}
\psnumber{ps_134-5_063-v0013}{13}
\psnumber{ps_134-5_063-v0014}{14}
\psnumber{ps_134-5_063-v0015}{15}
\psnumber{ps_134-5_063-v0016}{16}
\psnumber{ps_134-5_063-v0017}{17}
\psnumber{ps_134-5_063-v0018}{18}
\psnumber{ps_134-5_063-v0019}{19}
\psnumber{ps_134-5_063-v0020}{20}
\psnumber{ps_134-5_063-v0021}{21}
%% BEGIN SOURCE ps_134-5_063
% Pursuing Stacks (cotes 134-2 to 134-8): archive cote 134-5, site page 63 (PDF page 64), typescript page 250.
% typescript. Draft compared: ps4.tex:3172-3233.
% AI-transcribed by Claude Opus 5.5 (claude-opus-5-5, adaptive thinking, effort high), Message Batches API, batch ps2 (msgbatch_01NVwr1Yuo8qJf5rykeGBXPK), 2026-11-12; reviewed against both supplied scan halves; see CHECK_DECISIONS.tsv.
\scanpage{134-5/63}
\mspage{250}

I better correct the sta\var{\inserted{t}}{\hand{t}}{insertion}{ps_134-5_063-v0001}ement, to the effect that \var{\inserted{it seems}}{\hand{it seems}}{insertion}{ps_134-5_063-v0002} all modelizing \hand{\uline{\textcolor{black}{totally}}}
\hand{\uline{\textcolor{black}{aspheric}}} asphericity structures met with so far \var{\deleted{seem to be}}{}{deletion}{ps_134-5_063-v0003} \var{\inserted{are}}{\hand{are}}{insertion}{ps_134-5_063-v0004} generated by
suitable contractibility structures.

Let M be a category endowed with a contractibility structure $M_c$ ,
hence also with an asphericity structure $M_{as}$ , and let \var{\xout{\unsure{ix\allowbreak{}x\allowbreak{}Ax\allowbreak{}x\allowbreak{}x\allowbreak{}x\allowbreak{}Mx\allowbreak{}x\allowbreak{}bex\allowbreak{}x\allowbreak{}x\allowbreak{}}}}{}{typewriter deletion}{ps_134-5_063-v0005}
\var{\xout{\unsure{xxxx}}}{}{typewriter deletion}{ps_134-5_063-v0006}
\var{\xout{functor from a small category \unsure{Ix\allowbreak{}x\allowbreak{}Ax\allowbreak{}x\allowbreak{}x\allowbreak{}x\allowbreak{}Mx\allowbreak{}}}}{}{typewriter deletion}{ps_134-5_063-v0007} \var{\xout{\unsure{Ixx}}}{}{typewriter deletion}{ps_134-5_063-v0008}
\[ i : A \hand{\to} M \]
be a functor from a small category A to M. We want to give ns conditions
\var{\over{f}{f}}{\hand{f}}{writing over}{ps_134-5_063-v0009}or i to be a \hand{\uline{\textcolor{black}{test}}} \hand{\uline{\textcolor{black}{functor}}}, in terms of homotopy notions in M and in A\^{}.
To this end, it seems necessary to refresh \var{\hand{(}somewhat memory\hand{)}}{memory somewhat}{transposition}{ps_134-5_063-v0010}
% A continuous handwritten transposition stroke runs beneath somewhat and above memory; the typed word order is preserved.
 and recall
some relevant notions which were developped in part IV of our notes(sec-%
\hand{\uline{\textcolor{black}{tions 54 and 55).}}}

It will be convenient to call an object \var{\inserted[typed]{F}}{F}{typed insertion}{ps_134-5_063-v0011} of A\^{} \hand{\uline{\textcolor{black}{locally}}} \hand{\uline{\textcolor{black}{aspheric}}}
(resp. \hand{\uline{\textcolor{black}{totally}}} \hand{\uline{\textcolor{black}{aspheric}}} ), if it's product \var{\inserted[typed]{in A\^{}}}{in A\^{}}{typed insertion}{ps_134-5_063-v0012} by any object of A, and hence
also it's product by any aspheric object of A\^{}, is aspheric (resp. and
if moreover F itself is aspheric). With this terminology, A\^{} is totally
aspheric iff every aspheric object of A\^{} is totally aspheric, and \var{\inserted[typed]{if}}{if}{typed insertion}{ps_134-5_063-v0013} more-%
over A is aspheric i.e. $e_{A\hat{}}$ is aspheric. \var{\xout{If A is aspheric\unsure{x\allowbreak{}x\allowbreak{}x\allowbreak{}thenx\allowbreak{}Ix\allowbreak{}}}}{}{typewriter deletion}{ps_134-5_063-v0014} Note
that \var{\xout{\unsure{the}}}{}{typewriter deletion}{ps_134-5_063-v0015} \var{\inserted[typed]{F is}}{F is}{typed insertion}{ps_134-5_063-v0016} locally aspheric \var{\xout{\unsure{objectxx}}}{}{typewriter deletion}{ps_134-5_063-v0017} iff the map $F \hand{\to} e_{A\hat{}}$ is aspheric ,
or what amounts to the same, if this map is universally in $\underline{W}_A$ , or
equivalently, if the corresponding functor $A_{/F} \hand{\to} A$ is aspheric.
If A itself is aspheric, and\hand{/}in this case only, this condition implies
already that F is aspheric, i.e. that F is totally aspheric. We'll
denote by \var{\xout{\unsure{it}}}{}{typewriter deletion}{ps_134-5_063-v0018}
\[ \hat{A}_{tot.as} \quad ( \text{ resp. } \quad \hat{A}_{loc.\ as} \ ) \]
the full subcategory of A\^{} made up with the totally aspheric (resp.
locally aspheric) objects of A\^{} . Thus we get
% Hand-drawn arrows are marked blue; typed objects remain plain.
\diagram{
\begin{tikzcd}[column sep=small]
 \& \hat{A}_{tot.as} = \hat{A}_{loc.as} \hand{\cap} \hat{A}_{as} \ar[dl, hook', color=blue] \ar[dr, hook, color=blue] \& \\
\hat{A}_{loc.as} \& \& \hat{A}_{as} \ .
\end{tikzcd}
}

N\var{\over{o}{\unsure{o}}}{\hand{\unsure{o}}}{writing over}{ps_134-5_063-v0019}w recall that (section 54) in terms of the set $\underline{W}_A$ of weak equi-%
valences in A\^{}, we constructed a homotopy structure on A\^{}, more speci-%
fically a homotopy interval structure, admitting as a generating family
of homotopy intervals \var{\xout{\unsure{f}}}{}{typewriter deletion}{ps_134-5_063-v0020}
\[ \underline{I} =(I,\hand{\delta}_0,\hand{\delta}_1) \]
the set of intervals such that I be a locally aspheric object of A\^{}\hand{\unsure{,}}\var{\xout{\unsure{xi}}}{}{typewriter deletion}{ps_134-5_063-v0021}i.e.
% The uncertain comma-like tick is raised above the type line near i.e.; its function is unresolved.
\[ I \hand{\in} \mathrm{Ob}\, \hat{A}_{loc.\ as} \quad . \]

\margin{63}

%% END SOURCE ps_134-5_063
\ifnum\EditionMode=1
\psapparatusstart
\psentry{ps_134-5_063-v0001}{1}{insertion}{\inserted{t}}{\hand{t}}{The handwritten insertion is incorporated at its indicated position; insertion brackets leave the running text.}
\psentry{ps_134-5_063-v0002}{2}{insertion}{\inserted{it seems}}{\hand{it seems}}{The handwritten insertion is incorporated at its indicated position; insertion brackets leave the running text.}
\psentry{ps_134-5_063-v0003}{3}{deletion}{\deleted{seem to be}}{}{The author cancels this material by hand; it leaves the corrected running text.}
\psentry{ps_134-5_063-v0004}{4}{insertion}{\inserted{are}}{\hand{are}}{The handwritten insertion is incorporated at its indicated position; insertion brackets leave the running text.}
\psentry{ps_134-5_063-v0005}{5}{typewriter deletion}{\xout{\unsure{ix\allowbreak{}x\allowbreak{}Ax\allowbreak{}x\allowbreak{}x\allowbreak{}x\allowbreak{}Mx\allowbreak{}x\allowbreak{}bex\allowbreak{}x\allowbreak{}x\allowbreak{}}}}{}{The typewriter strikes out this material; it leaves the corrected running text.}
\psentry{ps_134-5_063-v0006}{6}{typewriter deletion}{\xout{\unsure{xxxx}}}{}{The typewriter strikes out this material; it leaves the corrected running text.}
\psentry{ps_134-5_063-v0007}{7}{typewriter deletion}{\xout{functor from a small category \unsure{Ix\allowbreak{}x\allowbreak{}Ax\allowbreak{}x\allowbreak{}x\allowbreak{}x\allowbreak{}Mx\allowbreak{}}}}{}{The typewriter strikes out this material; it leaves the corrected running text.}
\psentry{ps_134-5_063-v0008}{8}{typewriter deletion}{\xout{\unsure{Ixx}}}{}{The typewriter strikes out this material; it leaves the corrected running text.}
\psentry{ps_134-5_063-v0009}{9}{writing over}{\over{f}{f}}{\hand{f}}{The superposed handwritten reading replaces the underlying reading.}
\psentry{ps_134-5_063-v0010}{10}{transposition}{\hand{(}somewhat memory\hand{)}}{memory somewhat}{A continuous handwritten stroke passes beneath somewhat and above memory, with curved ends; the original typed order is retained in the diplomatic page and the author-indicated exchanged order is given here.}
\psentry{ps_134-5_063-v0011}{11}{typed insertion}{\inserted[typed]{F}}{F}{The typed insertion is incorporated at its indicated position; insertion brackets leave the running text.}
\psentry{ps_134-5_063-v0012}{12}{typed insertion}{\inserted[typed]{in A\^{}}}{in A\^{}}{The typed insertion is incorporated at its indicated position; insertion brackets leave the running text.}
\psentry{ps_134-5_063-v0013}{13}{typed insertion}{\inserted[typed]{if}}{if}{The typed insertion is incorporated at its indicated position; insertion brackets leave the running text.}
\psentry{ps_134-5_063-v0014}{14}{typewriter deletion}{\xout{If A is aspheric\unsure{x\allowbreak{}x\allowbreak{}x\allowbreak{}thenx\allowbreak{}Ix\allowbreak{}}}}{}{The typewriter strikes out this material; it leaves the corrected running text.}
\psentry{ps_134-5_063-v0015}{15}{typewriter deletion}{\xout{\unsure{the}}}{}{The typewriter strikes out this material; it leaves the corrected running text.}
\psentry{ps_134-5_063-v0016}{16}{typed insertion}{\inserted[typed]{F is}}{F is}{The typed insertion is incorporated at its indicated position; insertion brackets leave the running text.}
\psentry{ps_134-5_063-v0017}{17}{typewriter deletion}{\xout{\unsure{objectxx}}}{}{The typewriter strikes out this material; it leaves the corrected running text.}
\psentry{ps_134-5_063-v0018}{18}{typewriter deletion}{\xout{\unsure{it}}}{}{The typewriter strikes out this material; it leaves the corrected running text.}
\psentry{ps_134-5_063-v0019}{19}{writing over}{\over{o}{\unsure{o}}}{\hand{\unsure{o}}}{The superposed handwritten reading replaces the underlying reading.}
\psentry{ps_134-5_063-v0020}{20}{typewriter deletion}{\xout{\unsure{f}}}{}{The typewriter strikes out this material; it leaves the corrected running text.}
\psentry{ps_134-5_063-v0021}{21}{typewriter deletion}{\xout{\unsure{xi}}}{}{The typewriter strikes out this material; it leaves the corrected running text.}
\psapparatusend
\fi
\pssource{ps_134-5_064}{134-5/64}{https://grothendieck.umontpellier.fr/web/viewer.html?file=../134-5.pdf\#page=65}
\psnumber{ps_134-5_064-v0001}{1}
\psnumber{ps_134-5_064-v0002}{2}
\psnumber{ps_134-5_064-v0003}{3}
\psnumber{ps_134-5_064-v0004}{4}
\psnumber{ps_134-5_064-v0005}{5}
\psnumber{ps_134-5_064-v0006}{6}
\psnumber{ps_134-5_064-v0007}{7}
\psnumber{ps_134-5_064-v0008}{8}
\psnumber{ps_134-5_064-v0009}{9}
\psnumber{ps_134-5_064-v0010}{10}
\psnumber{ps_134-5_064-v0011}{11}
\psnumber{ps_134-5_064-v0012}{12}
%% BEGIN SOURCE ps_134-5_064
% Pursuing Stacks (cotes 134-2 to 134-8): archive cote 134-5, site page 64 (PDF page 65), typescript page 251.
% typescript. Draft compared: ps4.tex:3215-3280.
% AI-transcribed by Claude Opus 5.5 (claude-opus-5-5, adaptive thinking, effort high), Message Batches API, batch ps2 (msgbatch_01NVwr1Yuo8qJf5rykeGBXPK), 2026-11-12; reviewed against both supplied scan halves; see CHECK_DECISIONS.tsv.
\scanpage{134-5/64}
\mspage{251}

Let
\[ h \;=\; h_{\underline{W}_A} \]
this homotopy structure, hence a corresponding notion of h-\hand{\uline{\textcolor{black}{equivalence}}} \var{\inserted[typed]{or h-\hand{\uline{\textcolor{black}{homotopy}}}}}{or h-\hand{\uline{\textcolor{black}{homotopy}}}}{typed insertion}{ps_134-5_064-v0001}
$\overset{h}{\hand{\sim}}$ for arrows in $A\hat{}$, a corresponding notion of h-homotopisms i.e.
a set of arrows
\[ W^h \hand{\subset} \mathrm{Fl}\,A\hat{}\ \hand{,} \quad \text{such that} \quad W^h \hand{\subset} \underline{W}_A \quad , \]
a notion  of h-homotopy interval (namely an interval such that $\hand{\delta}_0$ and $\hand{\delta}_1$
be  h-homotopic , for which it is sufficient\hand{,}but not necessary\hand{,}that I
be locally aspheric ...), and last not least, a notion of \hand{\uline{\textcolor{black}{contractible}}}
\hand{\uline{\textcolor{black}{objects}}}, making up a full subcategory
\[ A\hat{}_c \hand{\subset} A\hat{}\ , \quad \text{such that} \quad A\hat{}_c \hand{\subset} A\hat{}_{\mathrm{loc.as}} \quad . \]
The latter inclusion, in case A is aspheric, can be equally written
\[ A\hat{}_c \hand{\subset} A\hat{}_{\mathrm{tot.as}} \quad \text{(if A aspheric).} \]

Coming back now to the contractibility structure $(M,M_c)$\hand{,} and a
functor  i: A $\hand{\longrightarrow}$ M , we are interested in the corresponding functor
\[ u = i^* : M \hand{\longrightarrow} \var{\deleted{\illegible{1}}}{}{deletion}{ps_134-5_064-v0002}\ M' \;=\; A\hat{}\ \hand{,} \]
where both members will be viewed as being endowed with their re\unsure{sp}ective
homotopy structures  -  the one of M being of the most restrictive type
envisioned in section 51, namely it is  defined in terms of a contracti-%
bility structure, whereas the one of \var{\deleted{\illegible{1}}}{}{deletion}{ps_134-5_064-v0003}\var{\inserted[typed]{M'}}{M'}{typed insertion}{ps_134-5_064-v0004} is a priori defined in terms of
a homotopy interval structure, but not necessarily in terms of a contrac
tibility structure. Now this situation has been described in section
5\unsure{3} , as far as compatibility \var{\inserted[typed]{conditions}}{conditions}{typed insertion}{ps_134-5_064-v0005} with homotopy structures are concerned ,
independently by the way of any special assumption on M'  (such as being
a category of presheaves on some small category A), or on the functor
u, except for commuting to finite products. Compatibility of u \var{\xout{xxx}}{}{typewriter deletion}{ps_134-5_064-v0006} with
the homotopy structures on M,M'  can be expressed by either one of the
following six conditions, which are all equivalent:

% A small isolated pen dot above the i of homotopic has no established lexical function.
H 1)  u transforms homotopic arrows of M into homotopic arrows in M'\hand{.}

H 2)  u transforms homotopisms \var{\deleted{into}}{}{deletion}{ps_134-5_064-v0007}\var{\inserted{in}}{\hand{in}}{insertion}{ps_134-5_064-v0008} M into homotopisms in M' .

H 3)  u transforms any homotopy interval $\underline{I} = (I,\hand{\delta}_0,\hand{\delta}_1)$ in M into a
homotopy interval in M' \hand{(} i.e. if two sections of an object
I over $e_M$ are homotopic, so are $u(\hand{\delta}_0)$ and $u(\hand{\delta}_1)$\hand{).}

H 3')  Same as H 3), but $\underline{I}$ being restricted to be in a given \var{\deleted{genera-}}{}{deletion}{ps_134-5_064-v0009}%
\var{\deleted{ting}}{}{deletion}{ps_134-5_064-v0010} family of homotopy intervals\hand{,} \var{\inserted{generating}}{\hand{generating}}{insertion}{ps_134-5_064-v0011} for the homotopy \var{\xout{intervals}}{}{typewriter deletion}{ps_134-5_064-v0012}
structure in M.

\margin{64}

%% END SOURCE ps_134-5_064
\ifnum\EditionMode=1
\psapparatusstart
\psentry{ps_134-5_064-v0001}{1}{typed insertion}{\inserted[typed]{or h-\hand{\uline{\textcolor{black}{homotopy}}}}}{or h-\hand{\uline{\textcolor{black}{homotopy}}}}{The typed insertion is incorporated at its indicated position; insertion brackets leave the running text.}
\psentry{ps_134-5_064-v0002}{2}{deletion}{$\deleted{\illegible{1}}$}{}{The author cancels this material by hand; it leaves the corrected running text.}
\psentry{ps_134-5_064-v0003}{3}{deletion}{\deleted{\illegible{1}}}{}{The author cancels this material by hand; it leaves the corrected running text.}
\psentry{ps_134-5_064-v0004}{4}{typed insertion}{\inserted[typed]{M'}}{M'}{The typed insertion is incorporated at its indicated position; insertion brackets leave the running text.}
\psentry{ps_134-5_064-v0005}{5}{typed insertion}{\inserted[typed]{conditions}}{conditions}{The typed insertion is incorporated at its indicated position; insertion brackets leave the running text.}
\psentry{ps_134-5_064-v0006}{6}{typewriter deletion}{\xout{xxx}}{}{The typewriter strikes out this material; it leaves the corrected running text.}
\psentry{ps_134-5_064-v0007}{7}{deletion}{\deleted{into}}{}{The author cancels this material by hand; it leaves the corrected running text.}
\psentry{ps_134-5_064-v0008}{8}{insertion}{\inserted{in}}{\hand{in}}{The handwritten insertion is incorporated at its indicated position; insertion brackets leave the running text.}
\psentry{ps_134-5_064-v0009}{9}{deletion}{\deleted{genera-}}{}{The author cancels this material by hand; it leaves the corrected running text.}
\psentry{ps_134-5_064-v0010}{10}{deletion}{\deleted{ting}}{}{The author cancels this material by hand; it leaves the corrected running text.}
\psentry{ps_134-5_064-v0011}{11}{insertion}{\inserted{generating}}{\hand{generating}}{The handwritten insertion is incorporated at its indicated position; insertion brackets leave the running text.}
\psentry{ps_134-5_064-v0012}{12}{typewriter deletion}{\xout{intervals}}{}{The typewriter strikes out this material; it leaves the corrected running text.}
\psapparatusend
\fi
\pssource{ps_134-5_065}{134-5/65}{https://grothendieck.umontpellier.fr/web/viewer.html?file=../134-5.pdf\#page=66}
\psnumber{ps_134-5_065-v0001}{1}
\psnumber{ps_134-5_065-v0002}{2}
\psnumber{ps_134-5_065-v0003}{3}
\psnumber{ps_134-5_065-v0004}{4}
\psnumber{ps_134-5_065-v0005}{5}
\psnumber{ps_134-5_065-v0006}{6}
\psnumber{ps_134-5_065-v0007}{7}
\psnumber{ps_134-5_065-v0008}{8}
\psnumber{ps_134-5_065-v0009}{9}
\psnumber{ps_134-5_065-v0010}{10}
\psnumber{ps_134-5_065-v0011}{11}
\psnumber{ps_134-5_065-v0012}{12}
\psnumber{ps_134-5_065-v0013}{13}
\psnumber{ps_134-5_065-v0014}{14}
\psnumber{ps_134-5_065-v0015}{15}
\psnumber{ps_134-5_065-v0016}{16}
\psnumber{ps_134-5_065-v0017}{17}
\psnumber{ps_134-5_065-v0018}{18}
\psnumber{ps_134-5_065-v0019}{19}
%% BEGIN SOURCE ps_134-5_065
% Pursuing Stacks (cotes 134-2 to 134-8): archive cote 134-5, site page 65 (PDF page 66), typescript page 252.
% typescript. Draft compared: ps4.tex:3264-3339.
% AI-transcribed by Claude Opus 5.5 (claude-opus-5-5, adaptive thinking, effort high), Message Batches API, batch ps2 (msgbatch_01NVwr1Yuo8qJf5rykeGBXPK), 2026-11-12; reviewed against both supplied scan halves; see CHECK_DECISIONS.tsv.
\scanpage{134-5/65}
\mspage{252}

H 4)\quad u transforms any contractible object\var{\inserted[typed]{x}}{x}{typed insertion}{ps_134-5_065-v0001} of M into a contractible
object of M' .

H 4')\quad Same as H 4) , with x being restricted to be in a given \var{\xout{\illegible{1}}}{}{typewriter deletion}{ps_134-5_065-v0002}
subcategory C of M, generating for the contractibility struc-%
ture of M.

\hand{\uline{\textcolor{black}{Remark}}}\hand{ 2.} It should be noted that the condition\var{\inserted{upon C}}{\hand{upon C}}{insertion}{ps_134-5_065-v0003} stated\var{\deleted{,}}{}{deletion}{ps_134-5_065-v0004} in H 4') \var{\deleted{upon C,}}{}{deletion}{ps_134-5_065-v0005} name\var{\over{l}{l}}{\hand{l}}{writing over}{ps_134-5_065-v0006}\hand{y}
that the (given) C should be generating for the contractibility structu\var{\xout{r}}{}{typewriter deletion}{ps_134-5_065-v0007}%
re $M_c$ of M, means exactly two things (a) $C \hand{\subset} M_c$ , and (b) the fami-%
ly of all intervals $\underline{I} = (I,\hand{\delta}_0,\hand{\delta}_1)$ made up with objects of C (these in-%
tervals are necessarily homotopy intervals for the\var{\xout{\unsure{given}}}{}{typewriter deletion}{ps_134-5_065-v0008} homotopy struc-%
ture of M) \hand{\uline{\textcolor{black}{generates}}} the homotopy structure of M, i.e. two arrows in M
are homotopic iff they can be joined by a chain a arrows, two consecuti-%
ve among which being related by an elementary homotopy involving an
interval of that family. This \var{\xout{remark}}{}{typewriter deletion}{ps_134-5_065-v0009}\var{\inserted[typed]{reminder}}{reminder}{typed insertion}{ps_134-5_065-v0010} being made, it follows that C
\var{\xout{can}}{}{typewriter deletion}{ps_134-5_065-v0011} is equally eligible for applying criterion H 3'),,which means that
\var{\xout{in H 4')}}{}{typewriter deletion}{ps_134-5_065-v0012} we get still another equivalent formulation of H 4'), by deman-%
ding merely that for any two sections $\hand{\delta}_0,\hand{\delta}_1$ of x over $e_M$ , the corres-%
ponding sections of u(x) should be homotopic.

After these preliminaries, we can state at last \var{\xout{what seems to me}}{}{typewriter deletion}{ps_134-5_065-v0013}
the following generalisation of the main result of section 65 (th.1
p. 176), concerning test functors with values in (Cat):

\margin{$\times$ small unnecessary, cf. also remarks on p. 348, \illegible{1}}
\hand{\uline{\textcolor{black}{Theorem 1}}} . Let $(M,M_c)$ be a\var{\xout{n}}{}{typewriter deletion}{ps_134-5_065-v0014} contractibility structure\var{\deleted{,}}{}{deletion}{ps_134-5_065-v0015}\var{\inserted[typed]{(cf p. 248)}}{(cf p. 248)}{typed insertion}{ps_134-5_065-v0016}\hand{,} C a \hand{\uwave{\textcolor{black}{small}}} full
\margin{\illegible{3}}
subcategory of M generating the contractibility structure (cf. remark 2
above), A a small category, and i : $A \hand{\longrightarrow} M$ a functor, factoring
through $M_c$ . We consider M as endowed equally with the asphericity
structure $M_{as}$ generated by $M_c$ (cf. prop.2). Then the following condi-%
tions on i are equivalent:

(i) The functor $i^{*}$: $M \hand{\to} A^{\wedge}$ deduced from i is compatible with the
homotopy structures on M and on $A^{\wedge}$ (cf. page 250\var{\inserted[typed]{-251}}{-251}{typed insertion}{ps_134-5_065-v0017} for the latter), i.e.
$i^{*}$ satisfies ei\hand{$\downarrow$}ther one of the six equivalent conditions H 1) to H 4')
above.

(ii) The functor $i^{*}$ is locally $M_{as}$-aspheric (def.2 p.245), i.e.
for any x in $M_{as}$ , $i^{*}(x)$ is a locally aspheric object of $A^{\wedge}$ (cf. p.250),
i.e. $i^{*}(x)$ is aspheric over $e_{A^{\wedge}}$ .

(iii) \var{\xout{For any}}{}{typewriter deletion}{ps_134-5_065-v0018}\var{\inserted[typed]{Same as}}{Same as}{typed insertion}{ps_134-5_065-v0019} (ii), but with x restricted to be in C.

\hand{\uline{\textcolor{black}{Corollary 1}}} . In order for i to be totally $M_{as}$-aspheric (def.2 p.245)

\margin{65}

%% END SOURCE ps_134-5_065
\ifnum\EditionMode=1
\psapparatusstart
\psentry{ps_134-5_065-v0001}{1}{typed insertion}{\inserted[typed]{x}}{x}{The typed insertion is incorporated at its indicated position; insertion brackets leave the running text.}
\psentry{ps_134-5_065-v0002}{2}{typewriter deletion}{\xout{\illegible{1}}}{}{The typewriter strikes out this material; it leaves the corrected running text.}
\psentry{ps_134-5_065-v0003}{3}{insertion}{\inserted{upon C}}{\hand{upon C}}{The handwritten insertion is incorporated at its indicated position; insertion brackets leave the running text.}
\psentry{ps_134-5_065-v0004}{4}{deletion}{\deleted{,}}{}{The author cancels this material by hand; it leaves the corrected running text.}
\psentry{ps_134-5_065-v0005}{5}{deletion}{\deleted{upon C,}}{}{The author cancels this material by hand; it leaves the corrected running text.}
\psentry{ps_134-5_065-v0006}{6}{writing over}{\over{l}{l}}{\hand{l}}{The superposed handwritten reading replaces the underlying reading.}
\psentry{ps_134-5_065-v0007}{7}{typewriter deletion}{\xout{r}}{}{The typewriter strikes out this material; it leaves the corrected running text.}
\psentry{ps_134-5_065-v0008}{8}{typewriter deletion}{\xout{\unsure{given}}}{}{The typewriter strikes out this material; it leaves the corrected running text.}
\psentry{ps_134-5_065-v0009}{9}{typewriter deletion}{\xout{remark}}{}{The typewriter strikes out this material; it leaves the corrected running text.}
\psentry{ps_134-5_065-v0010}{10}{typed insertion}{\inserted[typed]{reminder}}{reminder}{The typed insertion is incorporated at its indicated position; insertion brackets leave the running text.}
\psentry{ps_134-5_065-v0011}{11}{typewriter deletion}{\xout{can}}{}{The typewriter strikes out this material; it leaves the corrected running text.}
\psentry{ps_134-5_065-v0012}{12}{typewriter deletion}{\xout{in H 4')}}{}{The typewriter strikes out this material; it leaves the corrected running text.}
\psentry{ps_134-5_065-v0013}{13}{typewriter deletion}{\xout{what seems to me}}{}{The typewriter strikes out this material; it leaves the corrected running text.}
\psentry{ps_134-5_065-v0014}{14}{typewriter deletion}{\xout{n}}{}{The typewriter strikes out this material; it leaves the corrected running text.}
\psentry{ps_134-5_065-v0015}{15}{deletion}{\deleted{,}}{}{The author cancels this material by hand; it leaves the corrected running text.}
\psentry{ps_134-5_065-v0016}{16}{typed insertion}{\inserted[typed]{(cf p. 248)}}{(cf p. 248)}{The typed insertion is incorporated at its indicated position; insertion brackets leave the running text.}
\psentry{ps_134-5_065-v0017}{17}{typed insertion}{\inserted[typed]{-251}}{-251}{The typed insertion is incorporated at its indicated position; insertion brackets leave the running text.}
\psentry{ps_134-5_065-v0018}{18}{typewriter deletion}{\xout{For any}}{}{The typewriter strikes out this material; it leaves the corrected running text.}
\psentry{ps_134-5_065-v0019}{19}{typed insertion}{\inserted[typed]{Same as}}{Same as}{The typed insertion is incorporated at its indicated position; insertion brackets leave the running text.}
\psapparatusend
\fi
\pssource{ps_134-5_066}{134-5/66}{https://grothendieck.umontpellier.fr/web/viewer.html?file=../134-5.pdf\#page=67}
\psnumber{ps_134-5_066-v0001}{1}
\psnumber{ps_134-5_066-v0002}{2}
\psnumber{ps_134-5_066-v0003}{3}
\psnumber{ps_134-5_066-v0004}{4}
\psnumber{ps_134-5_066-v0005}{5}
\psnumber{ps_134-5_066-v0006}{6}
\psnumber{ps_134-5_066-v0007}{7}
\psnumber{ps_134-5_066-v0008}{8}
\psnumber{ps_134-5_066-v0009}{9}
\psnumber{ps_134-5_066-v0010}{10}
%% BEGIN SOURCE ps_134-5_066
% Pursuing Stacks (cotes 134-2 to 134-8): archive cote 134-5, site page 66 (PDF page 67), typescript page 253.
% typescript. Draft compared: ps4.tex:3321-3397.
% AI-transcribed by Claude Opus 5.5 (claude-opus-5-5, adaptive thinking, effort high), Message Batches API, batch ps2 (msgbatch_01NVwr1Yuo8qJf5rykeGBXPK), 2026-11-12; reviewed against both supplied scan halves; see CHECK_DECISIONS.tsv.
\scanpage{134-5/66}
\mspage{253}

it is ns that A be aspheric,  and that the equivalent conditions of th.1
be satisfied.

% A raised capital-like overtype is visible above the n of Indeed; no secure old/new letter pair is established.
Indeed,  if follows at once from the definitions \var{\xout{that}}{}{typewriter deletion}{ps_134-5_066-v0001} and from the
fact that the asphericity structure $(M,M_{as})$ is aspheric i.e. \var{\inserted{that}}{\hand{that}}{insertion}{ps_134-5_066-v0002} $e_M$ is \var{\xout{in M}}{}{typewriter deletion}{ps_134-5_066-v0003}
aspheric,  that  i is totally $M_{as}$-aspheric iff it is locally $M_{as}$-aspheric
(i.e. condition (ii) of theorem 1),  and \var{\inserted{if}}{\hand{if}}{insertion}{ps_134-5_066-v0004} moreover A is aspheric,  hence
the corollary.

\hand{\uline{\textcolor{black}{Corollary}}} 2.  In order for i to be a test functor,  it is ns that A be
a test category,  that $(M,M_{as})$  be modelizing,  and that the equivalent
conditions of theorem 1 be satisfied.

This follows from cor.1  and the reformulation of the notion of a
test-functor,  given \var{\xout{in xth}}{}{typewriter deletion}{ps_134-5_066-v0005} p. 245.

This corollary contains \var{\xout{that}}{}{typewriter deletion}{ps_134-5_066-v0006} the main result of section 65 as a parti-%
cular case,  when taking M = (Cat)  with the usual contractibility struc-%
ture,  \var{\xout{except that}}{}{typewriter deletion}{ps_134-5_066-v0007} and $C = \hand{\{\Delta}_1\hand{\}}$ ,  except that in loc. cit.  we did not
have to assume beforehand that A be a test category,  but only that A is
aspheric:  this condition,  plus condition (iii) above (namely,  $i^{\hand{*}}(\hand{\Delta}_1)$
locally aspheric)  implies already that A is a test category.  In order
to get also this extra result,  we state still another corollary:\hand{(*)}:

\hand{\uline{\textcolor{black}{Corollary}}} 3:  Assume
\[ \underline{I} = (I, \hand{\delta}_0, \hand{\delta}_1\hand{,} \hand{\mu}) \quad , \text{ with } I \hand{\in} \mathrm{Ob}\, C \ , \]
\margin{\unsure{multiplicati-} ty wholly unne- cessary and enough that $I \in M_c$ cf. p. 348 remark 1}
is a \hand{\uline{\textcolor{black}{multiplicative interval}}} in M ,  i.e. an interval endo\var{\xout{\unsure{w}}}{}{typewriter deletion}{ps_134-5_066-v0008}wed with a
multiplication $\hand{\mu}$, admitting $\hand{\delta}_0$ as a left unit and $\hand{\delta}_1$ as a left zero ele-%
ment (cf. section 49, p.108  and section 51, p.120  - where such an inter-%
val was provisionally called a ``contractor'').  Assume moreover that for
any x in $M_c$ ,  the two compositions  $x \hand{\to} e_M \hand{\overset{\delta_0,\delta_1}{\rightrightarrows}} I$  are distinct\hand{(} \var{\xout{for}}{}{typewriter deletion}{ps_134-5_066-v0009} which
is the case for instance if  $\mathrm{Ker}(\hand{\delta}_0,\hand{\delta}_1)$  exists in M and is a strict ini-%
tial object $\hand{\emptyset}_M$  of M  (i.e. an initial object such that any map $x \hand{\to} \hand{\emptyset}_M$ in
M is an isomorphism),  and \var{\inserted{moreover}}{\hand{moreover}}{insertion}{ps_134-5_066-v0010} $\hand{\emptyset}_M \hand{\notin} M_c$ \hand{)}. Then the conditions of th.1  imply
that A is a local test category,  and hence a test category provided A is
aspheric.

Indeed,  $i^*(\underline{I})$  is a multiplicative interval in $A^\wedge$ which is locally
aspheric,  and (as follows immediately from the assumptions on $\underline{I}$) separa-%
ting  -  hence $A^\wedge$ is a local test category.

\margin{66}

%% END SOURCE ps_134-5_066
\ifnum\EditionMode=1
\psapparatusstart
\psentry{ps_134-5_066-v0001}{1}{typewriter deletion}{\xout{that}}{}{The typewriter strikes out this material; it leaves the corrected running text.}
\psentry{ps_134-5_066-v0002}{2}{insertion}{\inserted{that}}{\hand{that}}{The handwritten insertion is incorporated at its indicated position; insertion brackets leave the running text.}
\psentry{ps_134-5_066-v0003}{3}{typewriter deletion}{\xout{in M}}{}{The typewriter strikes out this material; it leaves the corrected running text.}
\psentry{ps_134-5_066-v0004}{4}{insertion}{\inserted{if}}{\hand{if}}{The handwritten insertion is incorporated at its indicated position; insertion brackets leave the running text.}
\psentry{ps_134-5_066-v0005}{5}{typewriter deletion}{\xout{in xth}}{}{The typewriter strikes out this material; it leaves the corrected running text.}
\psentry{ps_134-5_066-v0006}{6}{typewriter deletion}{\xout{that}}{}{The typewriter strikes out this material; it leaves the corrected running text.}
\psentry{ps_134-5_066-v0007}{7}{typewriter deletion}{\xout{except that}}{}{The typewriter strikes out this material; it leaves the corrected running text.}
\psentry{ps_134-5_066-v0008}{8}{typewriter deletion}{\xout{\unsure{w}}}{}{The typewriter strikes out this material; it leaves the corrected running text.}
\psentry{ps_134-5_066-v0009}{9}{typewriter deletion}{\xout{for}}{}{The typewriter strikes out this material; it leaves the corrected running text.}
\psentry{ps_134-5_066-v0010}{10}{insertion}{\inserted{moreover}}{\hand{moreover}}{The handwritten insertion is incorporated at its indicated position; insertion brackets leave the running text.}
\psapparatusend
\fi
\pssource{ps_134-5_067}{134-5/67}{https://grothendieck.umontpellier.fr/web/viewer.html?file=../134-5.pdf\#page=68}
\psnumber{ps_134-5_067-v0001}{1}
\psnumber{ps_134-5_067-v0002}{2}
\psnumber{ps_134-5_067-v0003}{3}
\psnumber{ps_134-5_067-v0004}{4}
\psnumber{ps_134-5_067-v0005}{5}
\psnumber{ps_134-5_067-v0006}{6}
\psnumber{ps_134-5_067-v0007}{7}
\psnumber{ps_134-5_067-v0008}{8}
\psnumber{ps_134-5_067-v0009}{9}
\psnumber{ps_134-5_067-v0010}{10}
\psnumber{ps_134-5_067-v0011}{11}
\psnumber{ps_134-5_067-v0012}{12}
\psnumber{ps_134-5_067-v0013}{13}
\psnumber{ps_134-5_067-v0014}{14}
\psnumber{ps_134-5_067-v0015}{15}
%% BEGIN SOURCE ps_134-5_067
% Pursuing Stacks (cotes 134-2 to 134-8): archive cote 134-5, site page 67 (PDF page 68), typescript page 254.
% typescript. Draft compared: ps4.tex:3377-3441.
% AI-transcribed by Claude Opus 5.5 (claude-opus-5-5, adaptive thinking, effort high), Message Batches API, batch ps2 (msgbatch_01NVwr1Yuo8qJf5rykeGBXPK), 2026-11-12; reviewed against both supplied scan halves; see CHECK_DECISIONS.tsv.
\scanpage{134-5/67}
\mspage{254}

This array of immediate corollaries of th.1 do convince me that this
statement is indeed ``the'' natural generalization of the ``awkward'' main
result of section 51. All we have to do is to prove theorem 1 then.

The three conditions of theorem 1 can be rewritten simply as
\[ \text{(i)} \qquad i^{*}(C) \hand{\subset} A^{\wedge}_{c} \qquad \text{(using H 4')} \ , \]
\[ \text{(ii)} \qquad i^{*}(M_{as}) \hand{\subset} A^{\wedge}_{loc.as} \ , \]
\[ \text{(iii)} \qquad i^{*}(C) \hand{\subset} A^{\wedge}_{loc.as} \ , \]
and because of
\[ C \hand{\subset} M_{as} \ , \quad A^{\wedge}_{c} \hand{\subset} A^{\wedge}_{loc.\,as} \ , \]
it follows tau\var{\over[typed]{\unsure{o}}{t}}{t}{typewriter overwriting}{ps_134-5_067-v0001}ologically that (i) and (ii) both imply (iii). On the
other hand, (iii)$\hand{\Rightarrow}$(i) by criterion H 3'), and the definition of the
homotopy \var{\inserted[typed]{interval}}{interval}{typed insertion}{ps_134-5_067-v0002} structure of $A^{\wedge}$ in terms of \var{\xout{homotopy interval}}{}{typewriter deletion}{ps_134-5_067-v0003} $A^{\wedge}_{loc.as}$ . Thus
we are left with proving (i)$\hand{\Rightarrow}$(ii). But (i) can be rewritten as
\[ i^{*}(M_{c}) \hand{\subset} A^{\wedge}_{c} \ , \]
which implies
\[ (*) \qquad i^{*}(M_{c}) \hand{\subset} A^{\wedge}_{loc.as} \ . \]
\var{\deleted{To prove}}{}{deletion}{ps_134-5_067-v0004} \var{\over{t}{T}}{\hand{T}}{writing over}{ps_134-5_067-v0005}hat the latter condition implies (ii)\var{\deleted{, replacing A by any lo-%
calization, we are reduced to the case when A has a final object,
hence is aspheric, and to prove in this case that (*) implies}}{}{deletion}{ps_134-5_067-v0006}
\var{\deleted{\[ (**) \qquad i^{*}(M_{as}) \quad A^{\wedge}_{as} \ , \]}}{}{deletion}{ps_134-5_067-v0007}
\var{\deleted{i.e. that i is $M_{as}$-aspheric. On the other hand, be\unsure{c}ause of our assump-%
tion on A, $A^{\wedge}_{loc.as}$ \quad $A^{\wedge}_{as}$ , hence (*) implies}}{}{deletion}{ps_134-5_067-v0008}
\var{\deleted{\[ (*') \qquad i^{*}(M_{c}) \quad A^{\wedge}_{as} \ . \]}}{}{deletion}{ps_134-5_067-v0009}
\var{\deleted{The conclusion}}{}{deletion}{ps_134-5_067-v0010} now follows from the \var{\inserted{corollary to the}}{\hand{corollary to the}}{insertion}{ps_134-5_067-v0011} following tautology \var{\inserted{(with $N=M_c$)}}{\hand{(with $N=M_c$)}}{insertion}{ps_134-5_067-v0012}, which should
have been stated as a corollary to prop. 1 above (p. 247):

\hand{\uline{\textcolor{black}{Lemma}}}: Let $(M,M_{as})$ be an asphericity structure, generated by the full
subcategory N of M, let A be a small category, and i: A$\hand{\rightarrow}$M a functor
\hand{\uline{\textcolor{black}{factoring th\unsure{ro}ugh}}} N . Then \var{\xout{for}}{}{typewriter deletion}{ps_134-5_067-v0013} i is $M_{as}$-aspheric iff for any x in N ,
$i^{*}(x)$ is an aspheric object of $A^{\wedge}$ .

\hand{\uline{\textcolor{black}{Corollary}}} The functor i is locally \var{\inserted[typed]{$M_{as}$-}}{$M_{as}$-}{typed insertion}{ps_134-5_067-v0014} aspheric (resp. totally $M_{as}$-aspheric)
iff \var{\xout{for}}{}{typewriter deletion}{ps_134-5_067-v0015} $i^{*}(N) \hand{\subset} A^{\wedge}_{loc.as}$ (resp. $i^{*}(N) \hand{\subset} A^{\wedge}_{tot.as}$ ).

\margin{67}

%% END SOURCE ps_134-5_067
\ifnum\EditionMode=1
\psapparatusstart
\psentry{ps_134-5_067-v0001}{1}{typewriter overwriting}{\over[typed]{\unsure{o}}{t}}{t}{The superposed typed reading replaces the underlying reading.}
\psentry{ps_134-5_067-v0002}{2}{typed insertion}{\inserted[typed]{interval}}{interval}{The typed insertion is incorporated at its indicated position; insertion brackets leave the running text.}
\psentry{ps_134-5_067-v0003}{3}{typewriter deletion}{\xout{homotopy interval}}{}{The typewriter strikes out this material; it leaves the corrected running text.}
\psentry{ps_134-5_067-v0004}{4}{deletion}{\deleted{To prove}}{}{The author cancels this material by hand; it leaves the corrected running text.}
\psentry{ps_134-5_067-v0005}{5}{writing over}{\over{t}{T}}{\hand{T}}{The superposed handwritten reading replaces the underlying reading.}
\psentry{ps_134-5_067-v0006}{6}{deletion}{\deleted{, replacing A by any lo-%
calization, we are reduced to the case when A has a final object,
hence is aspheric, and to prove in this case that (*) implies}}{}{The author cancels this material by hand; it leaves the corrected running text.}
\psentry{ps_134-5_067-v0007}{7}{deletion}{\deleted{\[ (**) \qquad i^{*}(M_{as}) \quad A^{\wedge}_{as} \ , \]}}{}{The author cancels this material by hand; it leaves the corrected running text.}
\psentry{ps_134-5_067-v0008}{8}{deletion}{\deleted{i.e. that i is $M_{as}$-aspheric. On the other hand, be\unsure{c}ause of our assump-%
tion on A, $A^{\wedge}_{loc.as}$ \quad $A^{\wedge}_{as}$ , hence (*) implies}}{}{The author cancels this material by hand; it leaves the corrected running text.}
\psentry{ps_134-5_067-v0009}{9}{deletion}{\deleted{\[ (*') \qquad i^{*}(M_{c}) \quad A^{\wedge}_{as} \ . \]}}{}{The author cancels this material by hand; it leaves the corrected running text.}
\psentry{ps_134-5_067-v0010}{10}{deletion}{\deleted{The conclusion}}{}{The author cancels this material by hand; it leaves the corrected running text.}
\psentry{ps_134-5_067-v0011}{11}{insertion}{\inserted{corollary to the}}{\hand{corollary to the}}{The handwritten insertion is incorporated at its indicated position; insertion brackets leave the running text.}
\psentry{ps_134-5_067-v0012}{12}{insertion}{\inserted{(with $N=M_c$)}}{\hand{(with $N=M_c$)}}{The handwritten insertion is incorporated at its indicated position; insertion brackets leave the running text.}
\psentry{ps_134-5_067-v0013}{13}{typewriter deletion}{\xout{for}}{}{The typewriter strikes out this material; it leaves the corrected running text.}
\psentry{ps_134-5_067-v0014}{14}{typed insertion}{\inserted[typed]{$M_{as}$-}}{$M_{as}$-}{The typed insertion is incorporated at its indicated position; insertion brackets leave the running text.}
\psentry{ps_134-5_067-v0015}{15}{typewriter deletion}{\xout{for}}{}{The typewriter strikes out this material; it leaves the corrected running text.}
\psapparatusend
\fi
\pssource{ps_134-5_068}{134-5/68}{https://grothendieck.umontpellier.fr/web/viewer.html?file=../134-5.pdf\#page=69}
\psnumber{ps_134-5_068-v0001}{1}
\psnumber{ps_134-5_068-v0002}{2}
\psnumber{ps_134-5_068-v0003}{3}
\psnumber{ps_134-5_068-v0004}{4}
\psnumber{ps_134-5_068-v0005}{5}
\psnumber{ps_134-5_068-v0006}{6}
\psnumber{ps_134-5_068-v0007}{7}
\psnumber{ps_134-5_068-v0008}{8}
\psnumber{ps_134-5_068-v0009}{9}
\psnumber{ps_134-5_068-v0010}{10}
\psnumber{ps_134-5_068-v0011}{11}
\psnumber{ps_134-5_068-v0012}{12}
\psnumber{ps_134-5_068-v0013}{13}
\psnumber{ps_134-5_068-v0014}{14}
\psnumber{ps_134-5_068-v0015}{15}
\psnumber{ps_134-5_068-v0016}{16}
\psnumber{ps_134-5_068-v0017}{17}
\psnumber{ps_134-5_068-v0018}{18}
\psnumber{ps_134-5_068-v0019}{19}
\psnumber{ps_134-5_068-v0020}{20}
\psnumber{ps_134-5_068-v0021}{21}
\psnumber{ps_134-5_068-v0022}{22}
\psnumber{ps_134-5_068-v0023}{23}
\psnumber{ps_134-5_068-v0024}{24}
\psnumber{ps_134-5_068-v0025}{25}
\psnumber{ps_134-5_068-v0026}{26}
%% BEGIN SOURCE ps_134-5_068
% Pursuing Stacks (cotes 134-2 to 134-8): archive cote 134-5, site page 68 (PDF page 69), typescript page 255.
% typescript. Draft compared: ps4.tex:3425-3498.
% AI-transcribed by Claude Opus 5.5 (claude-opus-5-5, adaptive thinking, effort high), Message Batches API, batch ps2 (msgbatch_01NVwr1Yuo8qJf5rykeGBXPK), 2026-11-12; reviewed against both supplied scan halves; see CHECK_DECISIONS.tsv.
\scanpage{134-5/68}
\mspage{255}

\hand{\uline{\textcolor{black}{Remark}}} \hand{3}. Assume \var{\xout{that}}{}{typewriter deletion}{ps_134-5_068-v0001} in theorem 1 that A is totally aspheric. Then
\var{\xout{the condition (ii}}{}{typewriter deletion}{ps_134-5_068-v0002} condition (ii) just means that i is $M_{as}$-aspheric, and
condition (iii) that $i^{*}(x)$ is aspheric for any x in C (as \var{\xout{\unsure{for $A^{\wedge}$}}}{}{typewriter deletion}{ps_134-5_068-v0003} we got
$A^{\wedge}_{loc.as} = A^{\wedge}_{as} = A^{\wedge}_{tot.as}$ ). If moreover C satisfies the condition of
cor.\var{\over{2}{3}}{\hand{3}}{writing over}{ps_134-5_068-v0004} above, then these conditions imply that \var{\xout{A}}{}{typewriter deletion}{ps_134-5_068-v0005} A is a strict test cate\var{\xout{g}}{}{typewriter deletion}{ps_134-5_068-v0006}%
gory, \var{\xout{hence}}{}{typewriter deletion}{ps_134-5_068-v0007} and if $(M,M_{as})$ is modelizing, that i is a test functor
as stated in corollary 2. These observations sum up the substance of
the restatement of the main result of section 51, \var{\inserted{given}}{\hand{given}}{insertion}{ps_134-5_068-v0008} in theorem 2 of p.178\var{\xout{\illegible{1}}}{}{typewriter deletion}{ps_134-5_068-v0009}\hand{,}
\margin{compare cor. 2 p. 347 for \unsure{stronger} version}in the present general context.

80. In the preceding section, we associated to any contractibility
structure $M_{c}$ on a category M, an asphericity structure $M_{as}$ ``generated''
by the former in a natural sense. It is this possibility of associating
(in a topologically meaningful way) an asphericity structure to a given
contractibility structure, which singles out the latter structure type,
among the three essential distinct``homotopy flavored'' kind of structures
developped at length in sections 51 \var{\deleted{\illegible{1}}}{}{deletion}{ps_134-5_068-v0010}\var{\inserted[typed]{and}}{and}{typed insertion}{ps_134-5_068-v0011} 52 , in preference to the two
weaker notions of \var{\inserted{a}}{\hand{a}}{insertion}{ps_134-5_068-v0012} homotopy interval structure\hand{,} and of a homotopism struc-%
ture (or, equivalently, of a \var{\xout{notion of}}{}{typewriter deletion}{ps_134-5_068-v0013} \hand{``}homotopy relation\hand{''}). The
association
\[ (*) \qquad M_{c} \hand{\longmapsto} M_{as} \quad \hand{(}\hand{\longleftrightarrow} \var{\inserted{corresponding}}{\hand{corresponding}}{insertion}{ps_134-5_068-v0014} notion of weak equivalence W_{\hand{a}} in M\hand{)} \]
finally carried through \var{\xout{here}}{}{typewriter deletion}{ps_134-5_068-v0015} in the last section, \var{\over{has}{had}}{\hand{had}}{writing over}{ps_134-5_068-v0016} been forshadowed
\margin{110, 135}earlier (cf. \var{\inserted{p. 110, \deleted{135} and}}{\hand{p. 110,  and}}{insertion}{ps_134-5_068-v0017} p.142), but was pushed off for quite a while, in order
to``give precedence'' to the other approach in view \var{\inserted{by then}}{\hand{by then}}{insertion}{ps_134-5_068-v0018} towards more general
test functors\var{\inserted{than before}}{\hand{than before}}{insertion}{ps_134-5_068-v0019}\hand{,} leading up finally to the ``awkward main result'' of section
65. In \var{\deleted{\illegible{1}}}{}{deletion}{ps_134-5_068-v0020} (*) , as the asphericity structure $M_{as}$ is totally aspheric
and a fortiori aspheric, $M_{as}$ and the corresponding notion of weak equi-%
valence $W_{a}= W_{M_{c}}$ (which of course should not be confused with the notion
of homotopism\var{\xout{s}}{}{typewriter deletion}{ps_134-5_068-v0021} associated to $M_{c}$ , giving a considerably smaller set of
arrows $W_{c}$ ) determine each other mutually. Till the writing up of secti
\hand{-}on 51, it was the aspect ``weak equivalence'' $W_{a}$ which was in the fore,
whereas the \var{\xout{technically}}{}{typewriter deletion}{ps_134-5_068-v0022} conceptually more relevant aspect of ``aspheric
objects'' did not \var{\xout{come into the fore \illegible{3}}}{}{typewriter deletion}{ps_134-5_068-v0023}\var{\inserted[typed]{appear in full light}}{appear in full light}{typed insertion}{ps_134-5_068-v0024} before I was through with
grinding out the ``awkward approach'' (cf. p.188).

\var{\deleted{Here}}{}{deletion}{ps_134-5_068-v0025} \var{\over{i}{I}}{\hand{I}}{writing over}{ps_134-5_068-v0026}t is time now to remember the opposite association
\[ (\,**\,) \qquad W_{a} \hand{\longmapsto} \text{homotopy structure } h_{W_{a}} \qquad , \]

\margin{68}

%% END SOURCE ps_134-5_068
\ifnum\EditionMode=1
\psapparatusstart
\psentry{ps_134-5_068-v0001}{1}{typewriter deletion}{\xout{that}}{}{The typewriter strikes out this material; it leaves the corrected running text.}
\psentry{ps_134-5_068-v0002}{2}{typewriter deletion}{\xout{the condition (ii}}{}{The typewriter strikes out this material; it leaves the corrected running text.}
\psentry{ps_134-5_068-v0003}{3}{typewriter deletion}{\xout{\unsure{for $A^{\wedge}$}}}{}{The typewriter strikes out this material; it leaves the corrected running text.}
\psentry{ps_134-5_068-v0004}{4}{writing over}{\over{2}{3}}{\hand{3}}{The superposed handwritten reading replaces the underlying reading.}
\psentry{ps_134-5_068-v0005}{5}{typewriter deletion}{\xout{A}}{}{The typewriter strikes out this material; it leaves the corrected running text.}
\psentry{ps_134-5_068-v0006}{6}{typewriter deletion}{\xout{g}}{}{The typewriter strikes out this material; it leaves the corrected running text.}
\psentry{ps_134-5_068-v0007}{7}{typewriter deletion}{\xout{hence}}{}{The typewriter strikes out this material; it leaves the corrected running text.}
\psentry{ps_134-5_068-v0008}{8}{insertion}{\inserted{given}}{\hand{given}}{The handwritten insertion is incorporated at its indicated position; insertion brackets leave the running text.}
\psentry{ps_134-5_068-v0009}{9}{typewriter deletion}{\xout{\illegible{1}}}{}{The typewriter strikes out this material; it leaves the corrected running text.}
\psentry{ps_134-5_068-v0010}{10}{deletion}{\deleted{\illegible{1}}}{}{The author cancels this material by hand; it leaves the corrected running text.}
\psentry{ps_134-5_068-v0011}{11}{typed insertion}{\inserted[typed]{and}}{and}{The typed insertion is incorporated at its indicated position; insertion brackets leave the running text.}
\psentry{ps_134-5_068-v0012}{12}{insertion}{\inserted{a}}{\hand{a}}{The handwritten insertion is incorporated at its indicated position; insertion brackets leave the running text.}
\psentry{ps_134-5_068-v0013}{13}{typewriter deletion}{\xout{notion of}}{}{The typewriter strikes out this material; it leaves the corrected running text.}
\psentry{ps_134-5_068-v0014}{14}{insertion}{$\inserted{corresponding}$}{$\hand{corresponding}$}{The handwritten insertion is incorporated at its indicated position; insertion brackets leave the running text.}
\psentry{ps_134-5_068-v0015}{15}{typewriter deletion}{\xout{here}}{}{The typewriter strikes out this material; it leaves the corrected running text.}
\psentry{ps_134-5_068-v0016}{16}{writing over}{\over{has}{had}}{\hand{had}}{The superposed handwritten reading replaces the underlying reading.}
\psentry{ps_134-5_068-v0017}{17}{insertion}{\inserted{p. 110, \deleted{135} and}}{\hand{p. 110,  and}}{The handwritten insertion is incorporated at its indicated position; insertion brackets leave the running text.}
\psentry{ps_134-5_068-v0018}{18}{insertion}{\inserted{by then}}{\hand{by then}}{The handwritten insertion is incorporated at its indicated position; insertion brackets leave the running text.}
\psentry{ps_134-5_068-v0019}{19}{insertion}{\inserted{than before}}{\hand{than before}}{The handwritten insertion is incorporated at its indicated position; insertion brackets leave the running text.}
\psentry{ps_134-5_068-v0020}{20}{deletion}{\deleted{\illegible{1}}}{}{The author cancels this material by hand; it leaves the corrected running text.}
\psentry{ps_134-5_068-v0021}{21}{typewriter deletion}{\xout{s}}{}{The typewriter strikes out this material; it leaves the corrected running text.}
\psentry{ps_134-5_068-v0022}{22}{typewriter deletion}{\xout{technically}}{}{The typewriter strikes out this material; it leaves the corrected running text.}
\psentry{ps_134-5_068-v0023}{23}{typewriter deletion}{\xout{come into the fore \illegible{3}}}{}{The typewriter strikes out this material; it leaves the corrected running text.}
\psentry{ps_134-5_068-v0024}{24}{typed insertion}{\inserted[typed]{appear in full light}}{appear in full light}{The typed insertion is incorporated at its indicated position; insertion brackets leave the running text.}
\psentry{ps_134-5_068-v0025}{25}{deletion}{\deleted{Here}}{}{The author cancels this material by hand; it leaves the corrected running text.}
\psentry{ps_134-5_068-v0026}{26}{writing over}{\over{i}{I}}{\hand{I}}{The superposed handwritten reading replaces the underlying reading.}
\psapparatusend
\fi
\pssource{ps_134-5_069}{134-5/69}{https://grothendieck.umontpellier.fr/web/viewer.html?file=../134-5.pdf\#page=70}
\psnumber{ps_134-5_069-v0001}{1}
\psnumber{ps_134-5_069-v0002}{2}
\psnumber{ps_134-5_069-v0003}{3}
\psnumber{ps_134-5_069-v0004}{4}
\psnumber{ps_134-5_069-v0005}{5}
\psnumber{ps_134-5_069-v0006}{6}
\psnumber{ps_134-5_069-v0007}{7}
\psnumber{ps_134-5_069-v0008}{8}
\psnumber{ps_134-5_069-v0009}{9}
\psnumber{ps_134-5_069-v0010}{10}
\psnumber{ps_134-5_069-v0011}{11}
\psnumber{ps_134-5_069-v0012}{12}
\psnumber{ps_134-5_069-v0013}{13}
\psnumber{ps_134-5_069-v0014}{14}
\psnumber{ps_134-5_069-v0015}{15}
\psnumber{ps_134-5_069-v0016}{16}
\psnumber{ps_134-5_069-v0017}{17}
\psnumber{ps_134-5_069-v0018}{18}
\psnumber{ps_134-5_069-v0019}{19}
\psnumber{ps_134-5_069-v0020}{20}
\psnumber{ps_134-5_069-v0021}{21}
\psnumber{ps_134-5_069-v0022}{22}
\psnumber{ps_134-5_069-v0023}{23}
\psnumber{ps_134-5_069-v0024}{24}
\psnumber{ps_134-5_069-v0025}{25}
\psnumber{ps_134-5_069-v0026}{26}
\psnumber{ps_134-5_069-v0027}{27}
\psnumber{ps_134-5_069-v0028}{28}
\psnumber{ps_134-5_069-v0029}{29}
%% BEGIN SOURCE ps_134-5_069
% Pursuing Stacks (cotes 134-2 to 134-8): archive cote 134-5, site page 69 (PDF page 70), typescript page 256.
% typescript. Draft compared: ps4.tex:3480-3536.
% AI-transcribed by Claude Opus 5.5 (claude-opus-5-5, adaptive thinking, effort high), Message Batches API, batch ps2 (msgbatch_01NVwr1Yuo8qJf5rykeGBXPK), 2026-11-12; reviewed against both supplied scan halves; see CHECK_DECISIONS.tsv.
\scanpage{134-5/69}
\mspage{256}

associating to a notion of weak equivalence in M i.e. to any saturated
subset \var{\xout{$W_a$}}{}{typewriter deletion}{ps_134-5_069-v0001}
% A diagonal slash crosses the l in Fl(M); its function does not establish a different notation.
\[ W_a \hand{\subset} \mathrm{Fl}(M) \quad , \]
a corresponding homotopy structure $h_{W_a}$ \hand{,}\var{\deleted{\unsure{,}}}{}{deletion}{ps_134-5_069-v0002} a homotopy interval structure
as a matter of fact (section 54, p. 131).  Here we are primarily interes-%
ted of course in the case when $W_a$ is associated to a\var{\deleted{n}}{}{deletion}{ps_134-5_069-v0003} \var{\inserted[typed]{given}}{given}{typed insertion}{ps_134-5_069-v0004} asphericity structu
-re $M_{as}$ in M , which we may as well assume to be aspheric, so \var{\deleted{\unsure{x}as}}{}{deletion}{ps_134-5_069-v0005} \var{\inserted{as}}{\hand{as}}{insertion}{ps_134-5_069-v0006} to
be sure that $W_a$ and $M_{as}$ determine each other.  I recall that the weak
interval structure $h_{W_a}$ \var{\xout{is}}{}{typewriter deletion}{ps_134-5_069-v0007} can be described by the generating family of
homotopy intervals, consisting of all intervals
\[ \underline{I} = (I,\hand{\delta}_o,\hand{\delta}_1) \]
such that $I \hand{\to} e_M$ be universally in $W_a$ , i.e. such that for any object
x in M, the projection
\[ x \text{x} I \hand{\longrightarrow} x \]
be in $W_a$ .  \var{\deleted{In view of}}{}{deletion}{ps_134-5_069-v0008} \var{\inserted{For}}{\hand{For}}{insertion}{ps_134-5_069-v0009} pinning down further the exact relationship bet-%
ween contractibility structures (which \var{\xout{are}}{}{typewriter deletion}{ps_134-5_069-v0010} \var{\inserted[typed]{may\xout{x} be viewed as}}{may be viewed as}{typed insertion}{ps_134-5_069-v0011} just special types of homoto-%
py interval structures) and asphericity structures, we are thinking of
course \var{\xout{(in}}{}{typewriter deletion}{ps_134-5_069-v0012} more specifically of \hand{\uline{\textcolor{black}{totally}}} \hand{\uline{\textcolor{black}{aspheric}}} asphericity structu-%
res, in view of prop.2 c) of the preceding section (p. 248).  It is \var{\deleted{rea-%
sonable to expect}}{}{deletion}{ps_134-5_069-v0013} \var{\inserted{immediate}}{\hand{immediate}}{insertion}{ps_134-5_069-v0014} in this case that \var{\xout{xxx}}{}{typewriter deletion}{ps_134-5_069-v0015} \var{\inserted[typed]{for}}{for}{typed insertion}{ps_134-5_069-v0016} an object I of M,\var{\xout{is such that}}{}{typewriter deletion}{ps_134-5_069-v0017}
$I \hand{\to} e_M$ is in $UW_a$ (i.e. is universally in $W_a$) \var{\deleted{iff}}{}{deletion}{ps_134-5_069-v0018} if (and only if,
of course) I is aspheric, i.e. $I \hand{\longrightarrow} e_M$ is in $W_a$ .  \var{\deleted{But maybe I'll
leave \xout{this} elucidation of this point for a later moment, in order not
to diverge from \xout{the}}}{}{deletion}{ps_134-5_069-v0019} \var{\over{t}{T}}{\hand{T}}{writing over}{ps_134-5_069-v0020}he mos\var{\over{e}{t}}{\hand{t}}{writing over}{ps_134-5_069-v0021} relevant questions which come up here, \var{\inserted{now}}{\hand{now}}{insertion}{ps_134-5_069-v0022}
\var{\deleted{namely:}}{}{deletion}{ps_134-5_069-v0023} \var{\inserted{seem to me the following:}}{\hand{seem to me the following:}}{insertion}{ps_134-5_069-v0024}

1)  If $M_{as}$ is generated by a contractibility structure $M_c$ ,  is the
homotopy structure $h_{W_a}$ associated to $M_{as}$  (indeed, a \var{\xout{weak}}{}{typewriter deletion}{ps_134-5_069-v0025} homotopy in-%
terval structure as recalled above) also the one defined by $M_c$ ,  using
intervals in $M_c$ as a generating family of homotopy intervals\var{\over{\unsure{.}}{?}}{\hand{?}}{writing over}{ps_134-5_069-v0026}

2)  Conversely, what extra conditions on a given asphericity struc-%
ture $M_{as}$ on M are needed (besides total asphericity) to insure that the
corresponding homotopy structure $h_{W_a}$ on M \var{\xout{is a}}{}{typewriter deletion}{ps_134-5_069-v0027} \var{\inserted[typed]{comes from a}}{comes from a}{typed insertion}{ps_134-5_069-v0028} contractibility structure \var{\inserted{$M_c$}}{\hand{$M_c$}}{insertion}{ps_134-5_069-v0029}
(i.e. admits a generating family of homotopy intervals which are \hand{\uline{\textcolor{black}{contrac-}}}%

\margin{69}

%% END SOURCE ps_134-5_069
\ifnum\EditionMode=1
\psapparatusstart
\psentry{ps_134-5_069-v0001}{1}{typewriter deletion}{\xout{$W_a$}}{}{The typewriter strikes out this material; it leaves the corrected running text.}
\psentry{ps_134-5_069-v0002}{2}{deletion}{\deleted{\unsure{,}}}{}{The author cancels this material by hand; it leaves the corrected running text.}
\psentry{ps_134-5_069-v0003}{3}{deletion}{\deleted{n}}{}{The author cancels this material by hand; it leaves the corrected running text.}
\psentry{ps_134-5_069-v0004}{4}{typed insertion}{\inserted[typed]{given}}{given}{The typed insertion is incorporated at its indicated position; insertion brackets leave the running text.}
\psentry{ps_134-5_069-v0005}{5}{deletion}{\deleted{\unsure{x}as}}{}{The author cancels this material by hand; it leaves the corrected running text.}
\psentry{ps_134-5_069-v0006}{6}{insertion}{\inserted{as}}{\hand{as}}{The handwritten insertion is incorporated at its indicated position; insertion brackets leave the running text.}
\psentry{ps_134-5_069-v0007}{7}{typewriter deletion}{\xout{is}}{}{The typewriter strikes out this material; it leaves the corrected running text.}
\psentry{ps_134-5_069-v0008}{8}{deletion}{\deleted{In view of}}{}{The author cancels this material by hand; it leaves the corrected running text.}
\psentry{ps_134-5_069-v0009}{9}{insertion}{\inserted{For}}{\hand{For}}{The handwritten insertion is incorporated at its indicated position; insertion brackets leave the running text.}
\psentry{ps_134-5_069-v0010}{10}{typewriter deletion}{\xout{are}}{}{The typewriter strikes out this material; it leaves the corrected running text.}
\psentry{ps_134-5_069-v0011}{11}{typed insertion}{\inserted[typed]{may\xout{x} be viewed as}}{may be viewed as}{The typed insertion is incorporated at its indicated position; insertion brackets leave the running text.}
\psentry{ps_134-5_069-v0012}{12}{typewriter deletion}{\xout{(in}}{}{The typewriter strikes out this material; it leaves the corrected running text.}
\psentry{ps_134-5_069-v0013}{13}{deletion}{\deleted{rea-%
sonable to expect}}{}{The author cancels this material by hand; it leaves the corrected running text.}
\psentry{ps_134-5_069-v0014}{14}{insertion}{\inserted{immediate}}{\hand{immediate}}{The handwritten insertion is incorporated at its indicated position; insertion brackets leave the running text.}
\psentry{ps_134-5_069-v0015}{15}{typewriter deletion}{\xout{xxx}}{}{The typewriter strikes out this material; it leaves the corrected running text.}
\psentry{ps_134-5_069-v0016}{16}{typed insertion}{\inserted[typed]{for}}{for}{The typed insertion is incorporated at its indicated position; insertion brackets leave the running text.}
\psentry{ps_134-5_069-v0017}{17}{typewriter deletion}{\xout{is such that}}{}{The typewriter strikes out this material; it leaves the corrected running text.}
\psentry{ps_134-5_069-v0018}{18}{deletion}{\deleted{iff}}{}{The author cancels this material by hand; it leaves the corrected running text.}
\psentry{ps_134-5_069-v0019}{19}{deletion}{\deleted{But maybe I'll
leave \xout{this} elucidation of this point for a later moment, in order not
to diverge from \xout{the}}}{}{The author cancels this material by hand; it leaves the corrected running text.}
\psentry{ps_134-5_069-v0020}{20}{writing over}{\over{t}{T}}{\hand{T}}{The superposed handwritten reading replaces the underlying reading.}
\psentry{ps_134-5_069-v0021}{21}{writing over}{\over{e}{t}}{\hand{t}}{The superposed handwritten reading replaces the underlying reading.}
\psentry{ps_134-5_069-v0022}{22}{insertion}{\inserted{now}}{\hand{now}}{The handwritten insertion is incorporated at its indicated position; insertion brackets leave the running text.}
\psentry{ps_134-5_069-v0023}{23}{deletion}{\deleted{namely:}}{}{The author cancels this material by hand; it leaves the corrected running text.}
\psentry{ps_134-5_069-v0024}{24}{insertion}{\inserted{seem to me the following:}}{\hand{seem to me the following:}}{The handwritten insertion is incorporated at its indicated position; insertion brackets leave the running text.}
\psentry{ps_134-5_069-v0025}{25}{typewriter deletion}{\xout{weak}}{}{The typewriter strikes out this material; it leaves the corrected running text.}
\psentry{ps_134-5_069-v0026}{26}{writing over}{\over{\unsure{.}}{?}}{\hand{?}}{The superposed handwritten reading replaces the underlying reading.}
\psentry{ps_134-5_069-v0027}{27}{typewriter deletion}{\xout{is a}}{}{The typewriter strikes out this material; it leaves the corrected running text.}
\psentry{ps_134-5_069-v0028}{28}{typed insertion}{\inserted[typed]{comes from a}}{comes from a}{The typed insertion is incorporated at its indicated position; insertion brackets leave the running text.}
\psentry{ps_134-5_069-v0029}{29}{insertion}{\inserted{$M_c$}}{\hand{$M_c$}}{The handwritten insertion is incorporated at its indicated position; insertion brackets leave the running text.}
\psapparatusend
\fi
\pssource{ps_134-5_070}{134-5/70}{https://grothendieck.umontpellier.fr/web/viewer.html?file=../134-5.pdf\#page=71}
\psnumber{ps_134-5_070-v0001}{1}
\psnumber{ps_134-5_070-v0002}{2}
\psnumber{ps_134-5_070-v0003}{3}
\psnumber{ps_134-5_070-v0004}{4}
\psnumber{ps_134-5_070-v0005}{5}
\psnumber{ps_134-5_070-v0006}{6}
\psnumber{ps_134-5_070-v0007}{7}
\psnumber{ps_134-5_070-v0008}{8}
\psnumber{ps_134-5_070-v0009}{9}
\psnumber{ps_134-5_070-v0010}{10}
\psnumber{ps_134-5_070-v0011}{11}
\psnumber{ps_134-5_070-v0012}{12}
\psnumber{ps_134-5_070-v0013}{13}
\psnumber{ps_134-5_070-v0014}{14}
\psnumber{ps_134-5_070-v0015}{15}
%% BEGIN SOURCE ps_134-5_070
% Pursuing Stacks (cotes 134-2 to 134-8): archive cote 134-5, site page 70 (PDF page 71), typescript page 257.
% typescript. Draft compared: ps4.tex:3522-3584.
% AI-transcribed by Claude Opus 5.5 (claude-opus-5-5, adaptive thinking, effort high), Message Batches API, batch ps2 (msgbatch_01NVwr1Yuo8qJf5rykeGBXPK), 2026-11-12; reviewed against both supplied scan halves; see CHECK_DECISIONS.tsv.
\scanpage{134-5/70}
\mspage{257}

\hand{\uline{\textcolor{black}{tible}}}), and that moreover $M_c$ \hand{\uline{\textcolor{black}{generates}}} $M_{as}$ ?

\var{\deleted{The question just met with before was:
3) What extra condition (if any) is needed for a\xout{n} totally aspheric
asphericity structure, in order that any aspheric object \inserted[typed]{$I$} of $M$ \xout{is} be
"totally aspheric" i.e. the map $I \hand{\longrightarrow} e_M$ is in $UW_a$ ?}}{}{deletion}{ps_134-5_070-v0001}

Before looking up a little these questions, I would like \var{\xout{first}}{}{typewriter deletion}{ps_134-5_070-v0002} \var{\inserted{however}}{\hand{however}}{insertion}{ps_134-5_070-v0003} to
carry through at once the "idyllic picture" of canonical modelizers,
\var{\deleted{as}}{}{deletion}{ps_134-5_070-v0004} forshadowed in section 50 (p.110), as I feel that this should be pos-%
sible at present at no costs. Taking into, account the reflections of
the later sections 57 and 59, we get the following set-up.

Let $M$ be a $\underline{U}$-category\var{\inserted[typed]{stable under finite products}}{stable under finite products}{typed insertion}{ps_134-5_070-v0005}, endowed with a functor
(1) \quad $\hand{\pi_0^M \text{ or } \pi_0}$ \var{\deleted{\illegible{1}}}{}{deletion}{ps_134-5_070-v0006} $: M \hand{\longrightarrow} (\mathrm{Sets})$ ,

on which we make no assuption for the time being. We are thinking of
the example when $M$ is a totally 0-connected category (cf. prop. on page
142 for this notion) and $\hand{\pi}_o$ is the "connected components" functor, or \var{\inserted{when}}{\hand{when}}{insertion}{ps_134-5_070-v0007}
$M = (\mathrm{Spaces})$ and $\hand{\pi}_o$ corresponds to taking sets of arc-wise connected
components. According to section 54 (p.1\var{\over{2}{3}}{\hand{3}}{writing over}{ps_134-5_070-v0008}1), we introduce a correspon-%
ding homotopy interval structure h on $M$, admitting the generating family
of homotopy intervals made up with those intervals $\underline{I} = (I,\hand{\delta}_o,\hand{\delta}_1)$ \var{\xout{such}}{}{typewriter deletion}{ps_134-5_070-v0009}
for which $I \hand{\longrightarrow} e_M$ is "universally in $W_{\hand{\pi}_o}$" , namely
\[ \var{\xout{a}}{}{typewriter deletion}{ps_134-5_070-v0010}\ \hand{\pi}_o(x \text{ x } I) \hand{\longrightarrow} \hand{\pi}_o(x) \quad \text{bijective for any } x \text{ in } M \ . \]
(Under suitable conditions on $\hand{\pi}_o$ , this \var{\xout{structu}}{}{typewriter deletion}{ps_134-5_070-v0011} homotopy interval
structure h on $M$ is the widest one "compatible with $\hand{\pi}_o$" in the sense
of page 130 , namely such that $\hand{\pi}_o$ transforms homotopisms into homoto-%
pisms - cf . proposition p. 133\var{\unsure{\xout{7}}}{}{typewriter deletion}{ps_134-5_070-v0012}.) We are interested in the case when
this homotopy structure on $M$ can be described by a contractibility struc-%
ture $M_c$ on $M$ , which is then uni\var{\over{\unsure{q}}{q}}{\hand{q}}{writing over}{ps_134-5_070-v0013}qe of course, hence well defined in
terms of $\hand{\pi}_o$. Therefore the asphericity structure \var{\xout{defined by}}{}{typewriter deletion}{ps_134-5_070-v0014} generated
by $M_c$ is equally well defined in terms of $\hand{\pi}_o$ , and likewise the corres-%
ponding notion $W_a$ of "weak equivalence" . We then get a canonical func-%
tor
\[ (2) \qquad \mathrm{Hot}_M = W_a^{-1}M \hand{\longrightarrow} (\mathrm{Hot}) \]
(section 77, p.239). We'll have to find still a suitable \var{\inserted[typed]{extra}}{extra}{typed insertion}{ps_134-5_070-v0015} condition on
the functor $\hand{\pi}_o$ , implying that this functor is canonically isomorphic
to the functor deduced from (2) by composing with the canonical functor

\margin{70}

%% END SOURCE ps_134-5_070
\ifnum\EditionMode=1
\psapparatusstart
\psentry{ps_134-5_070-v0001}{1}{deletion}{\deleted{The question just met with before was:
3) What extra condition (if any) is needed for a\xout{n} totally aspheric
asphericity structure, in order that any aspheric object \inserted[typed]{$I$} of $M$ \xout{is} be
"totally aspheric" i.e. the map $I \hand{\longrightarrow} e_M$ is in $UW_a$ ?}}{}{The author cancels this material by hand; it leaves the corrected running text.}
\psentry{ps_134-5_070-v0002}{2}{typewriter deletion}{\xout{first}}{}{The typewriter strikes out this material; it leaves the corrected running text.}
\psentry{ps_134-5_070-v0003}{3}{insertion}{\inserted{however}}{\hand{however}}{The handwritten insertion is incorporated at its indicated position; insertion brackets leave the running text.}
\psentry{ps_134-5_070-v0004}{4}{deletion}{\deleted{as}}{}{The author cancels this material by hand; it leaves the corrected running text.}
\psentry{ps_134-5_070-v0005}{5}{typed insertion}{\inserted[typed]{stable under finite products}}{stable under finite products}{The typed insertion is incorporated at its indicated position; insertion brackets leave the running text.}
\psentry{ps_134-5_070-v0006}{6}{deletion}{\deleted{\illegible{1}}}{}{The author cancels this material by hand; it leaves the corrected running text.}
\psentry{ps_134-5_070-v0007}{7}{insertion}{\inserted{when}}{\hand{when}}{The handwritten insertion is incorporated at its indicated position; insertion brackets leave the running text.}
\psentry{ps_134-5_070-v0008}{8}{writing over}{\over{2}{3}}{\hand{3}}{The superposed handwritten reading replaces the underlying reading.}
\psentry{ps_134-5_070-v0009}{9}{typewriter deletion}{\xout{such}}{}{The typewriter strikes out this material; it leaves the corrected running text.}
\psentry{ps_134-5_070-v0010}{10}{typewriter deletion}{$\xout{a}$}{}{The typewriter strikes out this material; it leaves the corrected running text.}
\psentry{ps_134-5_070-v0011}{11}{typewriter deletion}{\xout{structu}}{}{The typewriter strikes out this material; it leaves the corrected running text.}
\psentry{ps_134-5_070-v0012}{12}{typewriter deletion}{\unsure{\xout{7}}}{}{All material inside this wrapper is cancelled. Its original doubt and presentation remain in the diplomatic text and apparatus; no empty doubt mark or page label is inserted into the corrected running text.}
\psentry{ps_134-5_070-v0013}{13}{writing over}{\over{\unsure{q}}{q}}{\hand{q}}{The superposed handwritten reading replaces the underlying reading.}
\psentry{ps_134-5_070-v0014}{14}{typewriter deletion}{\xout{defined by}}{}{The typewriter strikes out this material; it leaves the corrected running text.}
\psentry{ps_134-5_070-v0015}{15}{typed insertion}{\inserted[typed]{extra}}{extra}{The typed insertion is incorporated at its indicated position; insertion brackets leave the running text.}
\psapparatusend
\fi
\pssource{ps_134-5_071}{134-5/71}{https://grothendieck.umontpellier.fr/web/viewer.html?file=../134-5.pdf\#page=72}
\psnumber{ps_134-5_071-v0001}{1}
\psnumber{ps_134-5_071-v0002}{2}
\psnumber{ps_134-5_071-v0003}{3}
\psnumber{ps_134-5_071-v0004}{4}
\psnumber{ps_134-5_071-v0005}{5}
\psnumber{ps_134-5_071-v0006}{6}
\psnumber{ps_134-5_071-v0007}{7}
\psnumber{ps_134-5_071-v0008}{8}
\psnumber{ps_134-5_071-v0009}{9}
\psnumber{ps_134-5_071-v0010}{10}
\psnumber{ps_134-5_071-v0011}{11}
\psnumber{ps_134-5_071-v0012}{12}
%% BEGIN SOURCE ps_134-5_071
% Pursuing Stacks (cotes 134-2 to 134-8): archive cote 134-5, site page 71 (PDF page 72), typescript page 258.
% typescript. Draft compared: ps4.tex:3565-3623.
% AI-transcribed by Claude Opus 5.5 (claude-opus-5-5, adaptive thinking, effort high), Message Batches API, batch ps2 (msgbatch_01NVwr1Yuo8qJf5rykeGBXPK), 2026-11-12; reviewed against both supplied scan halves; see CHECK_DECISIONS.tsv.
\scanpage{134-5/71}
\mspage{258}

(3) \qquad (H\unsure{o}t) \hand{$\xrightarrow{\pi_o}$} (Sets) ,

which can be defined using a very mild extra condition on the basic loca-%
lizer $\underline{W}$ (namely f\hand{$\in$}$\underline{W}$ implies \hand{$\pi$}$_{o}$(f) bijective, cf condition L a) on
page 165). Thus there seems to be a little work ahead after all -
in order to deduce something like a one to one correspondance, say, bet-%
ween \var{\xout{functors (1)}}{}{typewriter deletion}{ps_134-5_071-v0001} pairs (M,\hand{$\pi$}$_{o}$) satisfying suitable conditions, and
\var{\inserted[typed]{certain types of}}{certain types of}{typed insertion}{ps_134-5_071-v0002} asphericity structures \var{\xout{on M, \unsure{presumabl} which presumably}}{}{typewriter deletion}{ps_134-5_071-v0003} (M,M$_{as}$)
\var{\xout{which are totally aspheric (which \unsure{presumab}}}{}{typewriter deletion}{ps_134-5_071-v0004} (which will have to be assu-%
med totally aspheric, and presumably \var{\inserted{a little}}{\hand{a little}}{insertion}{ps_134-5_071-v0005} more still).

The case of special interest to us is the one when the asphericity
structure we get on M in terms of the functor \hand{$\pi$}$^{M}_{o}$ is \hand{\uline{\textcolor{black}{modelizing}}}, hence
even strictly modelizing (i.e. (M,W$_{a}$) is a strict modelizer ), as
(M,M$_{as}$) is totally aspheric. If we assume moreover that \var{\xout{M is totally}}{}{typewriter deletion}{ps_134-5_071-v0006}
the category M is totally O-connected and that \hand{$\pi$}$_{o}$ is just the"connected
components" functor, then the modelizing asphericity structure we got
on M is canonically determined by the mere category structure of M,
and deserves therefore to be called \var{\xout{\unsure{the}}}{}{typewriter deletion}{ps_134-5_071-v0007}\var{\inserted{the}}{\hand{the}}{insertion}{ps_134-5_071-v0008} \hand{\uline{\textcolor{black}{canonical}}} \var{\inserted[typed]{(modelizing)}}{(modelizing)}{typed insertion}{ps_134-5_071-v0009} \hand{\uline{\textcolor{black}{asphericity structure}}}
\hand{\uline{\textcolor{black}{on}}} M. A \hand{\uline{\textcolor{black}{canonical}}} \hand{\uline{\textcolor{black}{modelizer}}} (M,W) is a modelizer which can be obtained
from a canonical asphericity structure (M,M$_{as}$) by W = W$_{a}$ = corresponding
set of weak equivalences (for M$_{as}$).

\hand{\uline{\textcolor{black}{Remark}}} The slightly sketchy definition we just gave for \var{\inserted[typed]{the}}{the}{typed insertion}{ps_134-5_071-v0010} canonical \var{\inserted[typed]{modelizing}}{modelizing}{typed insertion}{ps_134-5_071-v0011} as-%
phericity structures, and accordingly for the canonical modelizers, is
essentially complete. The point however which requires clarification is
the relationship between the \var{\xout{\unsure{t}component}}{}{typewriter deletion}{ps_134-5_071-v0012} "connected components" functor
on the corresponding (totally O-connected) category M, and the composi-%
tion of the canonical functors (2) and (3).

\margin{71}

%% END SOURCE ps_134-5_071
\ifnum\EditionMode=1
\psapparatusstart
\psentry{ps_134-5_071-v0001}{1}{typewriter deletion}{\xout{functors (1)}}{}{The typewriter strikes out this material; it leaves the corrected running text.}
\psentry{ps_134-5_071-v0002}{2}{typed insertion}{\inserted[typed]{certain types of}}{certain types of}{The typed insertion is incorporated at its indicated position; insertion brackets leave the running text.}
\psentry{ps_134-5_071-v0003}{3}{typewriter deletion}{\xout{on M, \unsure{presumabl} which presumably}}{}{The typewriter strikes out this material; it leaves the corrected running text.}
\psentry{ps_134-5_071-v0004}{4}{typewriter deletion}{\xout{which are totally aspheric (which \unsure{presumab}}}{}{The typewriter strikes out this material; it leaves the corrected running text.}
\psentry{ps_134-5_071-v0005}{5}{insertion}{\inserted{a little}}{\hand{a little}}{The handwritten insertion is incorporated at its indicated position; insertion brackets leave the running text.}
\psentry{ps_134-5_071-v0006}{6}{typewriter deletion}{\xout{M is totally}}{}{The typewriter strikes out this material; it leaves the corrected running text.}
\psentry{ps_134-5_071-v0007}{7}{typewriter deletion}{\xout{\unsure{the}}}{}{The typewriter strikes out this material; it leaves the corrected running text.}
\psentry{ps_134-5_071-v0008}{8}{insertion}{\inserted{the}}{\hand{the}}{The handwritten insertion is incorporated at its indicated position; insertion brackets leave the running text.}
\psentry{ps_134-5_071-v0009}{9}{typed insertion}{\inserted[typed]{(modelizing)}}{(modelizing)}{The typed insertion is incorporated at its indicated position; insertion brackets leave the running text.}
\psentry{ps_134-5_071-v0010}{10}{typed insertion}{\inserted[typed]{the}}{the}{The typed insertion is incorporated at its indicated position; insertion brackets leave the running text.}
\psentry{ps_134-5_071-v0011}{11}{typed insertion}{\inserted[typed]{modelizing}}{modelizing}{The typed insertion is incorporated at its indicated position; insertion brackets leave the running text.}
\psentry{ps_134-5_071-v0012}{12}{typewriter deletion}{\xout{\unsure{t}component}}{}{The typewriter strikes out this material; it leaves the corrected running text.}
\psapparatusend
\fi
\pssource{ps_134-5_072}{134-5/72}{https://grothendieck.umontpellier.fr/web/viewer.html?file=../134-5.pdf\#page=73}
\psnumber{ps_134-5_072-v0001}{1}
\psnumber{ps_134-5_072-v0002}{2}
\psnumber{ps_134-5_072-v0003}{3}
\psnumber{ps_134-5_072-v0004}{4}
\psnumber{ps_134-5_072-v0005}{5}
\psnumber{ps_134-5_072-v0006}{6}
\psnumber{ps_134-5_072-v0007}{7}
\psnumber{ps_134-5_072-v0008}{8}
\psnumber{ps_134-5_072-v0009}{9}
\psnumber{ps_134-5_072-v0010}{10}
\psnumber{ps_134-5_072-v0011}{11}
\psnumber{ps_134-5_072-v0012}{12}
\psnumber{ps_134-5_072-v0013}{13}
\psnumber{ps_134-5_072-v0014}{14}
\psnumber{ps_134-5_072-v0015}{15}
\psnumber{ps_134-5_072-v0016}{16}
\psnumber{ps_134-5_072-v0017}{17}
\psnumber{ps_134-5_072-v0018}{18}
\psnumber{ps_134-5_072-v0019}{19}
\psnumber{ps_134-5_072-v0020}{20}
\psnumber{ps_134-5_072-v0021}{21}
%% BEGIN SOURCE ps_134-5_072
% Pursuing Stacks (cotes 134-2 to 134-8): archive cote 134-5, site page 72 (PDF page 73), typescript page 259.
% typescript. Draft compared: ps4.tex:3605-3692.
% AI-transcribed by Claude Opus 5.5 (claude-opus-5-5, adaptive thinking, effort high), Message Batches API, batch ps2 (msgbatch_01NVwr1Yuo8qJf5rykeGBXPK), 2026-11-12; reviewed against both supplied scan halves; see CHECK_DECISIONS.tsv.
\scanpage{134-5/72}
\mspage{259}

25.6.

81. I did'nt find much time since monday for mathematical pondering -
the little I got nonetheless has been enough though for convincing myself
that things came out \var{\xout{\unsure{st}}}{}{typewriter deletion}{ps_134-5_072-v0001} more nicely still than I expected by then. \hand{\unsure{--}} \var{\over{O}{\unsure{o}}}{\hand{\unsure{o}}}{writing over}{ps_134-5_072-v0002}ne
main point \var{\(\underset{\hand{\cdot\ \cdot\ \cdot\ \cdot\ \cdot}}{\text{\deleted{being}}}\)}{being}{stet}{ps_134-5_072-v0003} that, provided the basic localizer \var{\inserted[typed]{$\underline{W}$}}{$\underline{W}$}{typed insertion}{ps_134-5_072-v0004} satisfies the mild
extra assumption Loc 4) below, any contractibility structure $M_c$ on a
category $M$ \var{\inserted[typed]{with finite products,}}{with finite products,}{typed insertion}{ps_134-5_072-v0005} \var{\xout{is}}{}{typewriter deletion}{ps_134-5_072-v0006} can be recovered, in the simplest imaginable way, in
terms of the associated asphericity structure $M_a$ , or equivalently, in
terms of the corresponding set $W_a$ of ``weak equivalences'';\var{\xout{x}}{}{typewriter deletion}{ps_134-5_072-v0007} namely, $M_c$
is the set of contractible objects of $M$ , for the \var{\inserted[typed]{homotopy}}{homotopy}{typed insertion}{ps_134-5_072-v0008} \var{\xout{weak}}{}{typewriter deletion}{ps_134-5_072-v0009} interval structu-%
re defined in terms of all intervals made up with objects of $M_a$ \var{\xout{x}}{}{typewriter deletion}{ps_134-5_072-v0010} . This
implies that the canonical map
\[ \hand{(1)} \qquad \mathrm{Homtp}_4(M) = \mathrm{Cont}(M) \hand{\hookrightarrow} \underline{W}\text{-}\mathrm{Asph}(M) \quad , \]
from the set of contractibility structures on $M$ to the set of aspherici-%
ty structures on $M$ relative to the basic localizer $\underline{W}$ (or \hand{``}\hand{\uline{\textcolor{black}{$\underline{W}$-asphericity
structures}}}\hand{''}), is \hand{\uline{\textcolor{black}{injective}}}. In other words, we may view \var{\xout{thexnotionxof}}{}{typewriter deletion}{ps_134-5_072-v0011}
a contractibility structure (on a category $M$ stable under finite products\hand{),}
which is an \hand{\uline{\textcolor{black}{absolute}}} notion (namely independaht of \var{\xout{any}}{}{typewriter deletion}{ps_134-5_072-v0012}\var{\inserted{the}}{\hand{the}}{insertion}{ps_134-5_072-v0013} choice of \var{\xout{the bax}}{}{typewriter deletion}{ps_134-5_072-v0014}
\var{\inserted{a}}{\hand{a}}{insertion}{ps_134-5_072-v0015} basic localizer $\underline{W}$), as a``\hand{\uline{\textcolor{black}{particular}}} \hand{\uline{\textcolor{black}{case}}}'' of a $\underline{W}$-asphericity structure
(depending on the choice of $\underline{W}$), namely as ``equivalent'' to a $\underline{W}$-aspherici-%
ty structure, satisfying some extra conditions which we'll have to writ\hand{e}
down below. \hand{//}
% The paragraph-break mark has two near-vertical strokes with horizontal extensions beneath down below and above the following This; the upper extension lies near satisfying.

This pleasant fact associates immediately with two related
ones. The first is just a reminder of our reflections of sections 51 and
52, namely that the set of contractibility structures on $M$ can be viewed
as one among four similar sets of ``homotopy structures'' on $M$
\[ \hand{(2)} \qquad \mathrm{Homtp}_4(M) \hand{\hookrightarrow} \mathrm{Homtp}_3(M) \hand{\hookrightarrow} \mathrm{Homtp}_2(M) \hand{\xrightarrow{\sim}} \mathrm{Homtp}_1(M) \quad , \]
corresponding to the four \var{\xout{maxim}}{}{typewriter deletion}{ps_134-5_072-v0016} basic ``homotopy notions'' met with so far,
namely (besides contractibility structures $\mathrm{Homtp}_4$) the \var{\xout{weak}}{}{typewriter deletion}{ps_134-5_072-v0017} \hand{\uline{\textcolor{black}{homotopy
interval structures}}} ($\mathrm{Homtp}_3$), the \hand{\uline{\textcolor{black}{homotopism structures}}} \var{\xout{$W_o$}}{}{typewriter deletion}{ps_134-5_072-v0018} ($\mathrm{Homtp}_2$) and
the \hand{\uline{\textcolor{black}{homotopy relations}}} \var{\inserted{between maps}}{\hand{between maps}}{insertion}{ps_134-5_072-v0019} ($\mathrm{Homtp}_1$) . In the sequel, if \var{\xout{$M_c$xxxObxMxxis}}{}{typewriter deletion}{ps_134-5_072-v0020}
\[ \hand{(3)} \qquad M_c \hand{\subset} \mathrm{Ob}\, M \]
is a given homotopy structure on $M$, we'll denote by \var{\xout{\unsure{$J_c$}}}{}{typewriter deletion}{ps_134-5_072-v0021}
\[ \hand{(4)} \qquad J_c \hand{\subset} \mathrm{Int}(M), \quad W_c \hand{\subset} \mathrm{Fl}(M) \ , \quad R_c \hand{\subset} \mathrm{Fl}(M)\text{x}\mathrm{Fl}(M) \]
the corresponding other three homotopy structures on $M$, where $\mathrm{Int}(M)$
denotes the set of all``intervals'' $\underline{I} = (I, \hand{\delta}_o, \hand{\delta}_1)$ in $M$, i.e. objects of
$M$ endowed with two sections $\hand{\delta}_o,\hand{\delta}_1$ over the final object $e_M$ of $M$.

\margin{72}

%% END SOURCE ps_134-5_072
\ifnum\EditionMode=1
\psapparatusstart
\psentry{ps_134-5_072-v0001}{1}{typewriter deletion}{\xout{\unsure{st}}}{}{The typewriter strikes out this material; it leaves the corrected running text.}
\psentry{ps_134-5_072-v0002}{2}{writing over}{\over{O}{\unsure{o}}}{\hand{\unsure{o}}}{The superposed handwritten reading replaces the underlying reading.}
\psentry{ps_134-5_072-v0003}{3}{stet}{\(\underset{\hand{\cdot\ \cdot\ \cdot\ \cdot\ \cdot}}{\text{\deleted{being}}}\)}{being}{Five short pen reinstatement dots appear beneath the crossed typed being. The diplomatic form preserves cancellation and dots; the author-final reading restores the word.}
\psentry{ps_134-5_072-v0004}{4}{typed insertion}{\inserted[typed]{$\underline{W}$}}{$\underline{W}$}{The typed insertion is incorporated at its indicated position; insertion brackets leave the running text.}
\psentry{ps_134-5_072-v0005}{5}{typed insertion}{\inserted[typed]{with finite products,}}{with finite products,}{The typed insertion is incorporated at its indicated position; insertion brackets leave the running text.}
\psentry{ps_134-5_072-v0006}{6}{typewriter deletion}{\xout{is}}{}{The typewriter strikes out this material; it leaves the corrected running text.}
\psentry{ps_134-5_072-v0007}{7}{typewriter deletion}{\xout{x}}{}{The typewriter strikes out this material; it leaves the corrected running text.}
\psentry{ps_134-5_072-v0008}{8}{typed insertion}{\inserted[typed]{homotopy}}{homotopy}{The typed insertion is incorporated at its indicated position; insertion brackets leave the running text.}
\psentry{ps_134-5_072-v0009}{9}{typewriter deletion}{\xout{weak}}{}{The typewriter strikes out this material; it leaves the corrected running text.}
\psentry{ps_134-5_072-v0010}{10}{typewriter deletion}{\xout{x}}{}{The typewriter strikes out this material; it leaves the corrected running text.}
\psentry{ps_134-5_072-v0011}{11}{typewriter deletion}{\xout{thexnotionxof}}{}{The typewriter strikes out this material; it leaves the corrected running text.}
\psentry{ps_134-5_072-v0012}{12}{typewriter deletion}{\xout{any}}{}{The typewriter strikes out this material; it leaves the corrected running text.}
\psentry{ps_134-5_072-v0013}{13}{insertion}{\inserted{the}}{\hand{the}}{The handwritten insertion is incorporated at its indicated position; insertion brackets leave the running text.}
\psentry{ps_134-5_072-v0014}{14}{typewriter deletion}{\xout{the bax}}{}{The typewriter strikes out this material; it leaves the corrected running text.}
\psentry{ps_134-5_072-v0015}{15}{insertion}{\inserted{a}}{\hand{a}}{The handwritten insertion is incorporated at its indicated position; insertion brackets leave the running text.}
\psentry{ps_134-5_072-v0016}{16}{typewriter deletion}{\xout{maxim}}{}{The typewriter strikes out this material; it leaves the corrected running text.}
\psentry{ps_134-5_072-v0017}{17}{typewriter deletion}{\xout{weak}}{}{The typewriter strikes out this material; it leaves the corrected running text.}
\psentry{ps_134-5_072-v0018}{18}{typewriter deletion}{\xout{$W_o$}}{}{The typewriter strikes out this material; it leaves the corrected running text.}
\psentry{ps_134-5_072-v0019}{19}{insertion}{\inserted{between maps}}{\hand{between maps}}{The handwritten insertion is incorporated at its indicated position; insertion brackets leave the running text.}
\psentry{ps_134-5_072-v0020}{20}{typewriter deletion}{\xout{$M_c$xxxObxMxxis}}{}{The typewriter strikes out this material; it leaves the corrected running text.}
\psentry{ps_134-5_072-v0021}{21}{typewriter deletion}{\xout{\unsure{$J_c$}}}{}{The typewriter strikes out this material; it leaves the corrected running text.}
\psapparatusend
\fi
\pssource{ps_134-5_073}{134-5/73}{https://grothendieck.umontpellier.fr/web/viewer.html?file=../134-5.pdf\#page=74}
\psnumber{ps_134-5_073-v0001}{1}
\psnumber{ps_134-5_073-v0002}{2}
\psnumber{ps_134-5_073-v0003}{3}
\psnumber{ps_134-5_073-v0004}{4}
\psnumber{ps_134-5_073-v0005}{5}
\psnumber{ps_134-5_073-v0006}{6}
\psnumber{ps_134-5_073-v0007}{7}
\psnumber{ps_134-5_073-v0008}{8}
%% BEGIN SOURCE ps_134-5_073
% Pursuing Stacks (cotes 134-2 to 134-8): archive cote 134-5, site page 73 (PDF page 74), typescript page 260.
% typescript. Draft compared: ps4.tex:3670-3754.
% AI-transcribed by Claude Opus 5.5 (claude-opus-5-5, adaptive thinking, effort high), Message Batches API, batch ps2 (msgbatch_01NVwr1Yuo8qJf5rykeGBXPK), 2026-11-12; reviewed against both supplied scan halves; see CHECK_DECISIONS.tsv.
\scanpage{134-5/73}
\mspage{260}

% Small strokes above and below the h of this The have no secure lexical interpretation.
The second fact alluded to above is concerned with behaviour of the
$\underline{W}$-asphericity notions,  for varying $\underline{W}$,  more specifically for a pair
\[ (5) \qquad \underline{W} \hand{\subset} \underline{W}' \hand{\subset} Fl((Cat)) \]
of two basic localizers,  $\underline{W}$  and $\underline{W}'$ ,  such that $\underline{W}$  ``refines'' $\underline{W}'$.  It then
follows that for any small category A,  we have
\[ (6) \qquad \underline{W}_A \hand{\subset} \underset{\text{\var{\deleted{\unsure{$W_A$}}}{}{deletion}{ps_134-5_073-v0001}}}{\underline{W}'_A} \quad , \]
and accordingly,  that \hand{\uline{\textcolor{black}{for any \var{\inserted{$\underline{W}$-}}{\hand{$\underline{W}$-}}{insertion}{ps_134-5_073-v0002}asphericity structure}}}
\[ (\var{\over[typed]{\unsure{6}}{7}}{7}{typewriter overwriting}{ps_134-5_073-v0003}) \qquad M_{\underline{W}} \hand{\subset} \text{Ob } M \]
\hand{\uline{\textcolor{black}{on}}} M,  \hand{\uline{\textcolor{black}{there exists a unique}}} \hand{\uline{\textcolor{black}{$\underline{W}'$-asphericity structure}}} \var{\xout{on}}{}{typewriter deletion}{ps_134-5_073-v0004} $M_{\underline{W}'}$ \hand{\uline{\textcolor{black}{on}}} M ,
\[ (8) \qquad M_{\underline{W}} \hand{\subset} \underset{\text{\var{\deleted{\illegible{1}}}{}{deletion}{ps_134-5_073-v0005}}}{M_{\underline{W}'}} \quad , \]
\hand{\uline{\textcolor{black}{such that for any small category}}} A,  \hand{\uline{\textcolor{black}{a functor}}} $A \hand{\to} M$  \unsure{w}hich is $M_{\underline{W}\text{\var{\xout{\illegible{1}}}{}{typewriter deletion}{ps_134-5_073-v0006}}}$-asphe-%
\hand{\uline{\textcolor{black}{ric}}} (\hand{\uline{\textcolor{black}{with re\var{\over{ps}{\unsure{sp}}}{\hand{\unsure{sp}}}{writing over}{ps_134-5_073-v0007}ect to}}} $\underline{W}$)  \hand{\uline{\textcolor{black}{is also}}} $M_{\underline{W}'}$-\hand{\uline{\textcolor{black}{aspheric}}} (\hand{\uline{\textcolor{black}{with respect to}}} $\underline{W}'$).  This
is merely a tautology,  which we did'nt state earlier,  because there was
no compelling reason before to look at what happens when $\underline{W}$ is allowed to
vary.  Thus we get a canonical map
\[ (9) \qquad \underline{W}\text{-Asph}(M) \hand{\longrightarrow} \underline{W}'\text{-Asph}(M) \quad , \]
with the evident trahsitivity property for a triple
\[ \underline{W} \hand{\subset} \underline{W}' \hand{\subset} \underline{W}'' \quad , \]
in other words we get a functor
\[ \underline{W} \hand{\longmapsto} \underline{W}\text{-Asph }(M) \]
from the category of all basic localizers (the arrows between localizers
being inclusions (5))  to the category of sets.  \var{The relation of\hand{$|$}this func%
-torial dependance  of  $\underline{W}$-Asph(M) on $\underline{W}$  \deleted{with} \hand{(}the canonical inclusion (1) \hand{with}\hand{)}}{The relation of the canonical inclusion (1) \hand{with} this func%
-torial dependance  of  $\underline{W}$-Asph(M) on $\underline{W}$}{transposition}{ps_134-5_073-v0008}
is expressed in the commutativity of
% Hand-drawn arrows are marked blue; typed objects remain plain.
% The preceding enclosure is linked upward to the position after of; see the explicit transposition variant.
\diagram{
\begin{tikzcd}[column sep=-6pt]
 \& \mathrm{Homtp}_4(M) = \mathrm{Cont}(M) \ar[dl, hook, color=blue] \ar[dr, hook, color=blue] \& \\
\underline{W}\text{-Asph}(M) \ar[rr, color=blue] \& \& \underline{W}'\text{-Asph}(M)
\end{tikzcd} \quad .
}

It is time to write down the ``mild extra assumption'' on $\underline{W}$ needed to
insure injectivity of (1),  namely the familiar enough condition:

% The h of The has small surrounding strokes; Fl has a slash-like stroke through l, and the preceding condition has a darkened n.
Loc 4)  The set $\underline{W} \hand{\subset} Fl((Cat))$ of weak equivalences in (Cat)  is \hand{``}compatible\hand{''}
with the functor $\hand{\underline{\pi}}_o$: (Cat) $\hand{\to}$ (Sets),  i.e.
\[ (10) \qquad f \hand{\in} \underline{W} \hand{\Rightarrow} \hand{\underline{\pi}}_{\unsure{o}}(f) \text{  bijective .} \]

\margin{73}

%% END SOURCE ps_134-5_073
\ifnum\EditionMode=1
\psapparatusstart
\psentry{ps_134-5_073-v0001}{1}{deletion}{\deleted{\unsure{$W_A$}}}{}{The author cancels this material by hand; it leaves the corrected running text.}
\psentry{ps_134-5_073-v0002}{2}{insertion}{\inserted{$\underline{W}$-}}{\hand{$\underline{W}$-}}{The handwritten insertion is incorporated at its indicated position; insertion brackets leave the running text.}
\psentry{ps_134-5_073-v0003}{3}{typewriter overwriting}{$\over[typed]{\unsure{6}}{7}$}{$7$}{The superposed typed reading replaces the underlying reading.}
\psentry{ps_134-5_073-v0004}{4}{typewriter deletion}{\xout{on}}{}{The typewriter strikes out this material; it leaves the corrected running text.}
\psentry{ps_134-5_073-v0005}{5}{deletion}{\deleted{\illegible{1}}}{}{The author cancels this material by hand; it leaves the corrected running text.}
\psentry{ps_134-5_073-v0006}{6}{typewriter deletion}{\xout{\illegible{1}}}{}{The typewriter strikes out this material; it leaves the corrected running text.}
\psentry{ps_134-5_073-v0007}{7}{writing over}{\over{ps}{\unsure{sp}}}{\hand{\unsure{sp}}}{The superposed handwritten reading replaces the underlying reading.}
\psentry{ps_134-5_073-v0008}{8}{transposition}{The relation of\hand{$|$}this func%
-torial dependance  of  $\underline{W}$-Asph(M) on $\underline{W}$  \deleted{with} \hand{(}the canonical inclusion (1) \hand{with}\hand{)}}{The relation of the canonical inclusion (1) \hand{with} this func%
-torial dependance  of  $\underline{W}$-Asph(M) on $\underline{W}$}{The phrase the canonical inclusion (1), together with the newly written with, is enclosed by a pen line which rises to a locator immediately after The relation of. The original with before the enclosed phrase is crossed out. Preserve typed order diplomatically; move the enclosed unit to the indicated position for the author-final reading.}
\psapparatusend
\fi
\pssource{ps_134-5_074}{134-5/74}{https://grothendieck.umontpellier.fr/web/viewer.html?file=../134-5.pdf\#page=75}
\psnumber{ps_134-5_074-v0001}{1}
\psnumber{ps_134-5_074-v0002}{2}
\psnumber{ps_134-5_074-v0003}{3}
\psnumber{ps_134-5_074-v0004}{4}
\psnumber{ps_134-5_074-v0005}{5}
\psnumber{ps_134-5_074-v0006}{6}
\psnumber{ps_134-5_074-v0007}{7}
%% BEGIN SOURCE ps_134-5_074
% Pursuing Stacks (cotes 134-2 to 134-8): archive cote 134-5, site page 74 (PDF page 75), typescript page 261.
% typescript. Draft compared: ps4.tex:3733-3810.
% AI-transcribed by Claude Opus 5.5 (claude-opus-5-5, adaptive thinking, effort high), Message Batches API, batch ps2 (msgbatch_01NVwr1Yuo8qJf5rykeGBXPK), 2026-11-12; reviewed against both supplied scan halves; see CHECK_DECISIONS.tsv.
\scanpage{134-5/74}
\mspage{261}

(Cf pages 213-214  for the conditions Loc 1)  to  Loc 3)\hand{.})

Among all basic localizers satisfying this extra condition,  there
is one coarsest of all,  which we'll call $\underline{W}_{o}$, defined by the condition
\[ (11) \qquad f \hand{\in} \underline{W}_{o} \hand{\Leftrightarrow} \hand{\underline{\pi}}_{o}(f) \text{ is bijective} \]
for any map f in (Cat).  It is clear too that among all basic localizers
there is a finest,  which we'll call $\underline{W}_{\mathrm{oo}}$ ,  and which can be describes as
\[ (12) \qquad \underline{W}_{\mathrm{oo}} = \hand{\bigcap} \underline{W} \hand{,} \quad \text{intersection of the set of all basic localizers in (Cat) .} \]
We'll see in part \hand{V} of the notes that $\underline{W}_{\mathrm{oo}}$ is none else than just the usu-%
al notion of weak equivalence\var{\deleted{,}}{}{deletion}{ps_134-5_074-v0001}  we started with\hand{,} at the very beginning of
our reflections  (cf. section 17).  Thus functoriality of (1) with respect
to $\underline{W}$ implies that (1) \var{\xout{is}}{}{typewriter deletion}{ps_134-5_074-v0002} for arbitrary $\underline{W}$ can be described in terms of the
particular case $\underline{W}_{\mathrm{oo}}$ ,  as the composition
% An unresolved crossing stroke lies over the hyphen before Asph in the middle term.
\[ (13) \qquad \mathrm{Cont}(M) \hand{\hookrightarrow} \underline{W}_{\mathrm{oo}}\text{\unsure{-}Asph}(M) \hand{\rightarrow} \underline{W}\text{-Asph}(M) \quad . \]
On the other hand,  the strongest version of injectivity of (1),  for
different $\underline{W}$'s ,  is obtained for $\underline{W}_{o}$ ,  i.e.  taking the map
\[ (14) \qquad \mathrm{Cont}(M) \hand{\hookrightarrow} \underline{W}_{o}\text{-Asph}(M) \ . \]

This last map seems to me of special significance,  because the
two sets it relates correspond to ``\hand{\uline{\textcolor{black}{absolute}}}'' notions  (not depending on
the choice of some $\underline{W}$ ),  and which moreover are \var{\over{\unsure{b}}{b}}{\hand{b}}{writing over}{ps_134-5_074-v0003}oth ``\hand{\uline{\textcolor{black}{elementary}}}''\hand{,} in the
sense that they do not depend on anything like consideration of non tri-%
vial homotopy or (co)homology invariants of objects of (Cat)\hand{.}  \var{\over{a}{A}}{\hand{A}}{writing over}{ps_134-5_074-v0004}s a matte
% NOTE: the end of this line ("matte...") is cut off at the right edge of the scan
of fact,  the notion of a contractibility structure corresponds to the
algebraic translation of one of the most elementary and intuitive topolo-%
gical notions,  \var{\xout{whereas the}}{}{typewriter deletion}{ps_134-5_074-v0005} namely contractibility;  whereas the notion
of a $\underline{W}_{o}$-Asphericity structure can be expressed just as ``elementarily'' in
terms of the functor \var{\xout{:}}{}{typewriter deletion}{ps_134-5_074-v0006}
\[ (15) \qquad \hand{\underline{\pi}}_{o} : (\mathrm{Cat}) \hand{\longrightarrow} (\mathrm{Sets}) \quad , \]
which we may call the ``basic'' connected components-functor,  which is
nothing but the algebraic counterpart of the basic intuitive notion of
connected components of a space.  We'll denote by
% The l in Fl(M) carries a diagonal slash-like overstrike; its function is unresolved.
\[ (16) \qquad M_{o} = M_{\underline{W}_{o}} \hand{\subset} \mathrm{Ob}\, M \ , \quad \var{\xout{Wx}}{}{typewriter deletion}{ps_134-5_074-v0007}\ W_{o} \hand{\subset} \mathrm{Fl}(M) \]
the set of $\underline{W}_{o}$-aspheric objects and the set of $\underline{W}_{o}$-weak equivalences,  asso
% NOTE: the end of this line ("asso...") is cut off at the right edge of the scan
ciated to a given contractibility structure $M_{c}$ on M.  The objects of $M_{o}$

\margin{74}

%% END SOURCE ps_134-5_074
\ifnum\EditionMode=1
\psapparatusstart
\psentry{ps_134-5_074-v0001}{1}{deletion}{\deleted{,}}{}{The author cancels this material by hand; it leaves the corrected running text.}
\psentry{ps_134-5_074-v0002}{2}{typewriter deletion}{\xout{is}}{}{The typewriter strikes out this material; it leaves the corrected running text.}
\psentry{ps_134-5_074-v0003}{3}{writing over}{\over{\unsure{b}}{b}}{\hand{b}}{The superposed handwritten reading replaces the underlying reading.}
\psentry{ps_134-5_074-v0004}{4}{writing over}{\over{a}{A}}{\hand{A}}{The superposed handwritten reading replaces the underlying reading.}
\psentry{ps_134-5_074-v0005}{5}{typewriter deletion}{\xout{whereas the}}{}{The typewriter strikes out this material; it leaves the corrected running text.}
\psentry{ps_134-5_074-v0006}{6}{typewriter deletion}{\xout{:}}{}{The typewriter strikes out this material; it leaves the corrected running text.}
\psentry{ps_134-5_074-v0007}{7}{typewriter deletion}{$\xout{Wx}$}{}{The typewriter strikes out this material; it leaves the corrected running text.}
\psapparatusend
\fi
\pssource{ps_134-5_075}{134-5/75}{https://grothendieck.umontpellier.fr/web/viewer.html?file=../134-5.pdf\#page=76}
\psnumber{ps_134-5_075-v0001}{1}
\psnumber{ps_134-5_075-v0002}{2}
\psnumber{ps_134-5_075-v0003}{3}
\psnumber{ps_134-5_075-v0004}{4}
\psnumber{ps_134-5_075-v0005}{5}
\psnumber{ps_134-5_075-v0006}{6}
\psnumber{ps_134-5_075-v0007}{7}
\psnumber{ps_134-5_075-v0008}{8}
\psnumber{ps_134-5_075-v0009}{9}
\psnumber{ps_134-5_075-v0010}{10}
\psnumber{ps_134-5_075-v0011}{11}
\psnumber{ps_134-5_075-v0012}{12}
\psnumber{ps_134-5_075-v0013}{13}
%% BEGIN SOURCE ps_134-5_075
% Pursuing Stacks (cotes 134-2 to 134-8): archive cote 134-5, site page 75 (PDF page 76), typescript page 262.
% typescript. Draft compared: ps4.tex:3795-3872.
% AI-transcribed by Claude Opus 5.5 (claude-opus-5-5, adaptive thinking, effort high), Message Batches API, batch ps2 (msgbatch_01NVwr1Yuo8qJf5rykeGBXPK), 2026-11-12; reviewed against both supplied scan halves; see CHECK_DECISIONS.tsv.
\scanpage{134-5/75}
\margin{(26.6.) (*) This name for \uline{maps} is improper though, as it rather suggests the property of being ``universally in $W_0$''. We could call \unsure{them} ``\uline{0-équivalences}'' with \unsure{respect to the} localizer $W_0$ \illegible{3}}
\margin{\illegible{4}}
\mspage{262}
merit the name of 0-\hand{\uline{\textcolor{black}{connected}}} \hand{\uline{\textcolor{black}{objects}}} of M (with respect to $M_c$ ), and
the arrows of M in $W_o$ merit the name of 0-\hand{\uline{\textcolor{black}{connected}}} \hand{\uline{\textcolor{black}{maps}}}\hand{$^*$}(with respect
to $M_c$ ). \hand{\unsure{\&}} These two notions of 0-connectedness determine each other in an
evident way (valid for any $\underline{W}$\var{\xout{\illegible{1}}}{}{typewriter deletion}{ps_134-5_075-v0001}-aspheric $\underline{W}$\var{\xout{\illegible{1}}}{}{typewriter deletion}{ps_134-5_075-v0002}-asphericity structure, for
any basic localizer $\underline{W}$ ...). Explicitly, this can be expressed by
\[
(17)\quad \hand{\Biggl\{}\begin{array}{l}
x\hand{\in} M_o \hand{\Longleftrightarrow} \hand{\Bigl(}(x\hand{\to} e_M)\hand{\in} W_o\hand{\Bigr)}\\
(f:x\hand{\to} y)\hand{\in} W_o \hand{\Longleftrightarrow} \hand{\Bigl(}\hand{\underline{\pi}}_o(M_{o/x})\hand{\to} \hand{\underline{\pi}}_o(M_{o/y}) \text{ bijective}\hand{\Bigr)}.
\end{array}
\]

The injectivity of (1) can be restated by saying that \hand{\uline{\textcolor{black}{any contracti-}}}%
\hand{\uline{\textcolor{black}{bility structure}}} $M_c$ \hand{\uline{\textcolor{black}{on a category}}} M (with finite products) \hand{\uline{\textcolor{black}{can be recove-}}}%
\hand{\uline{\textcolor{black}{red}}} \hand{\uline{\textcolor{black}{in terms of the corresponding notion of}}} 0-\hand{\uline{\textcolor{black}{connected}}} \hand{\uline{\textcolor{black}{objects of}}} M ,
or, equivalently, in terms of the corresponding notion of 0-connected
maps in M. Still another way of phrasing this result, is in terms of
the canonical functor
\[
M\hand{\to} W_o^{-1}M = \mathrm{Hot}_{(M,M_o),\underline{W}_o} \hand{\longrightarrow} \mathrm{Hot}_{\underline{W}_o} = \underline{W}_o^{-1}(\mathrm{Cat})
\]
\[
\hand{\xrightarrow{\approx}} (\mathrm{Sets}) \quad ,
\]
which is a functor
\[
(18)\qquad \hand{\pi}_o : M \hand{\longrightarrow} (\mathrm{Sets})
\]
canonically associated to the contractibility structure. We may say that
\hand{\uline{\textcolor{black}{the contractibility structure}}} $M_c$ \var{\xout{is determined}}{}{typewriter deletion}{ps_134-5_075-v0003} \hand{\uline{\textcolor{black}{can be recovered in terms}}}
\hand{\uline{\textcolor{black}{of the corresponding functor}}} $\hand{\pi}_o$ , more accurately still in terms of the
isomorphism class of the latter. Indeed , in terms of this functor $\hand{\pi}_o$,
we recover $M_o$ and $W_o$ by the \var{\xout{formulas}}{}{typewriter deletion}{ps_134-5_075-v0004} relations :
\var{\xout{(19)}}{}{typewriter deletion}{ps_134-5_075-v0005}
\[
(19)\qquad
\begin{array}{l}
M_o = \hand{\Bigl\{} x \hand{\hand{\in}} \mathrm{Ob}\, M \hand{\Bigm|} \hand{\pi}_o(x) \text{ a one-point set}\hand{\Bigr\}}\hand{,}\\
W_o = \hand{\Bigl\{} f \hand{\hand{\in}} \mathrm{Fl}(M) \hand{\Bigm|} \hand{\pi}_o(f) \text{ is bijective}\hand{\Bigr\}} \quad .
\end{array}
\]
This shows that the ``nice'' main fact mentioned at the beginning of
today's notes, namely (essentially) injectivity of (1) (and, moreover\unsure{,}
an explicit description of a way how to recover an $M_c$ in terms of the
corresponding \var{\xout{aspher}}{}{typewriter deletion}{ps_134-5_075-v0006} $\underline{W}$-asphericity structure) is \var{\xout{only}}{}{typewriter deletion}{ps_134-5_075-v0007} not really depen-%
dent on \var{\inserted[typed]{relatively}}{relatively}{typed insertion}{ps_134-5_075-v0008} sophisticated notions \var{\deleted{of}}{}{deletion}{ps_134-5_075-v0009} \var{\inserted{such as}}{\hand{such as}}{insertion}{ps_134-5_075-v0010} ``basic localizers'' and corresponding
``asphericity structures'', but can be viewed as an ``e\var{\inserted{l}}{\hand{l}}{insertion}{ps_134-5_075-v0011}ementary'' result
% The initial typed e in elementary has dark overlapping type beneath the handwritten l; no further underlying letter is securely recoverable.
(namely independant of any consideration of ``higher'' homotopy or homology
invariants, apart from $\hand{\pi}_o$) about relationship between contractibility
structures, and corresponding 0-connectedness notions; \var{\deleted{which}}{}{deletion}{ps_134-5_075-v0012} \var{\inserted{the latter}}{\hand{the latter}}{insertion}{ps_134-5_075-v0013} may at
will be expressed in terms of either one of the three structural data
\[
(20)\qquad M_o \ , \ W_o \ \text{ or } \ \hand{\pi}_o \ .
\]
\margin{75}

%% END SOURCE ps_134-5_075
\ifnum\EditionMode=1
\psapparatusstart
\psentry{ps_134-5_075-v0001}{1}{typewriter deletion}{\xout{\illegible{1}}}{}{The typewriter strikes out this material; it leaves the corrected running text.}
\psentry{ps_134-5_075-v0002}{2}{typewriter deletion}{\xout{\illegible{1}}}{}{The typewriter strikes out this material; it leaves the corrected running text.}
\psentry{ps_134-5_075-v0003}{3}{typewriter deletion}{\xout{is determined}}{}{The typewriter strikes out this material; it leaves the corrected running text.}
\psentry{ps_134-5_075-v0004}{4}{typewriter deletion}{\xout{formulas}}{}{The typewriter strikes out this material; it leaves the corrected running text.}
\psentry{ps_134-5_075-v0005}{5}{typewriter deletion}{\xout{(19)}}{}{The typewriter strikes out this material; it leaves the corrected running text.}
\psentry{ps_134-5_075-v0006}{6}{typewriter deletion}{\xout{aspher}}{}{The typewriter strikes out this material; it leaves the corrected running text.}
\psentry{ps_134-5_075-v0007}{7}{typewriter deletion}{\xout{only}}{}{The typewriter strikes out this material; it leaves the corrected running text.}
\psentry{ps_134-5_075-v0008}{8}{typed insertion}{\inserted[typed]{relatively}}{relatively}{The typed insertion is incorporated at its indicated position; insertion brackets leave the running text.}
\psentry{ps_134-5_075-v0009}{9}{deletion}{\deleted{of}}{}{The author cancels this material by hand; it leaves the corrected running text.}
\psentry{ps_134-5_075-v0010}{10}{insertion}{\inserted{such as}}{\hand{such as}}{The handwritten insertion is incorporated at its indicated position; insertion brackets leave the running text.}
\psentry{ps_134-5_075-v0011}{11}{insertion}{\inserted{l}}{\hand{l}}{The handwritten insertion is incorporated at its indicated position; insertion brackets leave the running text.}
\psentry{ps_134-5_075-v0012}{12}{deletion}{\deleted{which}}{}{The author cancels this material by hand; it leaves the corrected running text.}
\psentry{ps_134-5_075-v0013}{13}{insertion}{\inserted{the latter}}{\hand{the latter}}{The handwritten insertion is incorporated at its indicated position; insertion brackets leave the running text.}
\psapparatusend
\fi
\pssource{ps_134-5_076}{134-5/76}{https://grothendieck.umontpellier.fr/web/viewer.html?file=../134-5.pdf\#page=77}
\psnumber{ps_134-5_076-v0001}{1}
\psnumber{ps_134-5_076-v0002}{2}
\psnumber{ps_134-5_076-v0003}{3}
\psnumber{ps_134-5_076-v0004}{4}
\psnumber{ps_134-5_076-v0005}{5}
\psnumber{ps_134-5_076-v0006}{6}
\psnumber{ps_134-5_076-v0007}{7}
\psnumber{ps_134-5_076-v0008}{8}
\psnumber{ps_134-5_076-v0009}{9}
\psnumber{ps_134-5_076-v0010}{10}
\psnumber{ps_134-5_076-v0011}{11}
\psnumber{ps_134-5_076-v0012}{12}
\psnumber{ps_134-5_076-v0013}{13}
\psnumber{ps_134-5_076-v0014}{14}
\psnumber{ps_134-5_076-v0015}{15}
\psnumber{ps_134-5_076-v0016}{16}
\psnumber{ps_134-5_076-v0017}{17}
\psnumber{ps_134-5_076-v0018}{18}
\psnumber{ps_134-5_076-v0019}{19}
\psnumber{ps_134-5_076-v0020}{20}
\psnumber{ps_134-5_076-v0021}{21}
\psnumber{ps_134-5_076-v0022}{22}
\psnumber{ps_134-5_076-v0023}{23}
\psnumber{ps_134-5_076-v0024}{24}
\psnumber{ps_134-5_076-v0025}{25}
\psnumber{ps_134-5_076-v0026}{26}
\psnumber{ps_134-5_076-v0027}{27}
\psnumber{ps_134-5_076-v0028}{28}
\psnumber{ps_134-5_076-v0029}{29}
\psnumber{ps_134-5_076-v0030}{30}
\psnumber{ps_134-5_076-v0031}{31}
\psnumber{ps_134-5_076-v0032}{32}
\psnumber{ps_134-5_076-v0033}{33}
\psnumber{ps_134-5_076-v0034}{34}
\psnumber{ps_134-5_076-v0035}{35}
%% BEGIN SOURCE ps_134-5_076
% Pursuing Stacks (cotes 134-2 to 134-8): archive cote 134-5, site page 76 (PDF page 77), typescript page 263.
% typescript. Draft compared: ps4.tex:3857-3918.
% AI-transcribed by Claude Opus 5.5 (claude-opus-5-5, adaptive thinking, effort high), Message Batches API, batch ps2 (msgbatch_01NVwr1Yuo8qJf5rykeGBXPK), 2026-11-12; reviewed against both supplied scan halves; see CHECK_DECISIONS.tsv.
\scanpage{134-5/76}
\margin{à la ligne}
\mspage{2\var{\over{0}{6}}{\hand{6}}{writing over}{ps_134-5_076-v0001}3}

This relationship has been "in the air" \var{\xout{xxxx}}{}{typewriter deletion}{ps_134-5_076-v0002} since section 50 (p. 109-%
110),  and I kind of turned around it consistently up to section 60,
without really getting to the core.  One reason for this "turning around"
has been,  I guess,  that I let myself be distracted,  not to say hypno-%
tized,  by the"canonical" $\hand{\pi}_0$ functor on a category (which makes really
good sense only when M is"\unsure{totally} 0-connected"\var{\xout{though}}{}{typewriter deletion}{ps_134-5_076-v0003} as a category, \var{\inserted{a condition which}}{\hand{a condition which}}{insertion}{ps_134-5_076-v0004} \var{\deleted{which}}{}{deletion}{ps_134-5_076-v0005}
\var{\over{w}{sh}}{\hand{sh}}{writing over}{ps_134-5_076-v0006}ould mean \var{\inserted[typed]{more or less}}{more or less}{typed insertion}{ps_134-5_076-v0007} I suppose that the 0-connected objects of M,  defined in terms
of the mere category structure of M,  define a $\underline{W}_{\unsure{\bullet|0}}$-asphericity structure
on M ),  Even after realizing (in section 59)  that one should generalize
the \var{\deleted{definition}}{}{deletion}{ps_134-5_076-v0008} \var{\inserted{description}}{\hand{description}}{insertion}{ps_134-5_076-v0009} of a contractibility structure in terms of \var{\xout{axxx}}{}{typewriter deletion}{ps_134-5_076-v0010} a "con-%
nected components functor" $\hand{\pi}_0$ ,  to the case of a functor $\hand{\pi}_0$ : M $\hand{\rightarrow}$(Sets)
given beforehand,  and satisfying suitable restrictions (which I did not
try to elucidate),  I still was under the impression\var{\xout{s}}{}{typewriter deletion}{ps_134-5_076-v0011} that the contracti-%
bility structures one could get this \var{\xout{xx}}{}{typewriter deletion}{ps_134-5_076-v0012} way must be of an extremely speci-%
al nature.  In order to become aware of the fact that this is by no means
so,  namely that any contractibility structure could be obtained from a
suitable functor $\hand{\pi}_0$ ,  it \var{\xout{sh}}{}{typewriter deletion}{ps_134-5_076-v0013}would have been necessary to notice that such
a structure $M_c$  defines in a natural way a functor $\hand{\pi}_0$ .  There was \var{\xout{indeed}}{}{typewriter deletion}{ps_134-5_076-v0014}
indeed the realization that $M_c$ should \var{\inserted[typed]{allow to}}{allow to}{typed insertion}{ps_134-5_076-v0015} define a notion of weak equivalence
(cf. page 136),  but it was'nt clearly realized by then that at the same
time as a notion of weak equivalence W in M,  we should also get a canoni%
cal functor
\[ M \hand{\longrightarrow} W^{-1}M \hand{\longrightarrow} \text{(Hot)} \quad , \]
namely something a lot more precise still than a functor with values in
(Sets)  merely !  But rather than push ahead in this direction, I then
% A separate faint typed xx cancellation lies below "this direction", before the next line.
\var{\xout{xx}}{}{typewriter deletion}{ps_134-5_076-v0016}
decided (p.138)  that it would be "unreasonable" \var{\xout{xxx}}{}{typewriter deletion}{ps_134-5_076-v0017} to go on \var{\inserted[typed]{still longer}}{still longer}{typed insertion}{ps_134-5_076-v0018} pushing off
investigation of test functors with \unsure{values|valuas} in (Cat),  following \var{\over{an}{the}}{\hand{the}}{writing over}{ps_134-5_076-v0019} app-%
roach which had \var{\inserted{\deleted{\illegible{1}}}}{\hand{}}{insertion}{ps_134-5_076-v0020} been on my mind for quite a while by then,  and \var{\deleted{written}}{}{deletion}{ps_134-5_076-v0021} \var{\inserted{finally}}{\hand{finally}}{insertion}{ps_134-5_076-v0022}
\var{\deleted{down}}{}{deletion}{ps_134-5_076-v0023} \var{\inserted{sketched}}{\hand{sketched}}{insertion}{ps_134-5_076-v0024} (with the promise of a corresponding generalization of the former
"key result" on testnfunctors)  in section 47.  Retrospectively,  the
whole "grinding" part III of these notes  now looks as a rather heavy
and \var{\inserted[typed]{long-winded}}{long-winded}{typed insertion}{ps_134-5_076-v0025} \var{\xout{awkwardx}}{}{typewriter deletion}{ps_134-5_076-v0026}digression,prompted by this approach\var{\xout{,}}{}{typewriter deletion}{ps_134-5_076-v0027}  \var{\xout{which}}{}{typewriter deletion}{ps_134-5_076-v0028} to still a parti-%
cular case of test functors (namely with values in (Cat),  and more strin-%
gently still \var{\inserted{factoring through}}{\hand{factoring through}}{insertion}{ps_134-5_076-v0029} \var{\deleted{in}}{}{deletion}{ps_134-5_076-v0030} (Cat)$_{as}$ ).  \var{\xout{which}}{}{typewriter deletion}{ps_134-5_076-v0031}
% NOTE: a handwritten line under the struck "which" and a double vertical stroke before "This" mark a new paragraph here
This particular case has been \var{\xout{by}}{}{typewriter deletion}{ps_134-5_076-v0032} \var{\inserted[typed]{of}}{of}{typed insertion}{ps_134-5_076-v0033} no
\var{\xout{means}}{}{typewriter deletion}{ps_134-5_076-v0034} use in the present part IV of the notes,  developping the really
relevant notions in terms of asphericity structures.  \hand{\uline{\textcolor{black}{Technically spea\var{\xout{k}}{}{typewriter deletion}{ps_134-5_076-v0035}-}}}%
\hand{\uline{\textcolor{black}{king}}},  \hand{\uline{\textcolor{black}{it now appears that most of the reflections of part III}}}  \hand{\uline{\textcolor{black}{are super}}}%
\hand{\uline{\textcolor{black}{seded by part IV}}}  -  the main exception being the development of the

\margin{76}

%% END SOURCE ps_134-5_076
\ifnum\EditionMode=1
\psapparatusstart
\psentry{ps_134-5_076-v0001}{1}{writing over}{\over{0}{6}}{\hand{6}}{The superposed handwritten reading replaces the underlying reading.}
\psentry{ps_134-5_076-v0002}{2}{typewriter deletion}{\xout{xxxx}}{}{The typewriter strikes out this material; it leaves the corrected running text.}
\psentry{ps_134-5_076-v0003}{3}{typewriter deletion}{\xout{though}}{}{The typewriter strikes out this material; it leaves the corrected running text.}
\psentry{ps_134-5_076-v0004}{4}{insertion}{\inserted{a condition which}}{\hand{a condition which}}{The handwritten insertion is incorporated at its indicated position; insertion brackets leave the running text.}
\psentry{ps_134-5_076-v0005}{5}{deletion}{\deleted{which}}{}{The author cancels this material by hand; it leaves the corrected running text.}
\psentry{ps_134-5_076-v0006}{6}{writing over}{\over{w}{sh}}{\hand{sh}}{The superposed handwritten reading replaces the underlying reading.}
\psentry{ps_134-5_076-v0007}{7}{typed insertion}{\inserted[typed]{more or less}}{more or less}{The typed insertion is incorporated at its indicated position; insertion brackets leave the running text.}
\psentry{ps_134-5_076-v0008}{8}{deletion}{\deleted{definition}}{}{The author cancels this material by hand; it leaves the corrected running text.}
\psentry{ps_134-5_076-v0009}{9}{insertion}{\inserted{description}}{\hand{description}}{The handwritten insertion is incorporated at its indicated position; insertion brackets leave the running text.}
\psentry{ps_134-5_076-v0010}{10}{typewriter deletion}{\xout{axxx}}{}{The typewriter strikes out this material; it leaves the corrected running text.}
\psentry{ps_134-5_076-v0011}{11}{typewriter deletion}{\xout{s}}{}{The typewriter strikes out this material; it leaves the corrected running text.}
\psentry{ps_134-5_076-v0012}{12}{typewriter deletion}{\xout{xx}}{}{The typewriter strikes out this material; it leaves the corrected running text.}
\psentry{ps_134-5_076-v0013}{13}{typewriter deletion}{\xout{sh}}{}{The typewriter strikes out this material; it leaves the corrected running text.}
\psentry{ps_134-5_076-v0014}{14}{typewriter deletion}{\xout{indeed}}{}{The typewriter strikes out this material; it leaves the corrected running text.}
\psentry{ps_134-5_076-v0015}{15}{typed insertion}{\inserted[typed]{allow to}}{allow to}{The typed insertion is incorporated at its indicated position; insertion brackets leave the running text.}
\psentry{ps_134-5_076-v0016}{16}{typewriter deletion}{\xout{xx}}{}{The typewriter strikes out this material; it leaves the corrected running text.}
\psentry{ps_134-5_076-v0017}{17}{typewriter deletion}{\xout{xxx}}{}{The typewriter strikes out this material; it leaves the corrected running text.}
\psentry{ps_134-5_076-v0018}{18}{typed insertion}{\inserted[typed]{still longer}}{still longer}{The typed insertion is incorporated at its indicated position; insertion brackets leave the running text.}
\psentry{ps_134-5_076-v0019}{19}{writing over}{\over{an}{the}}{\hand{the}}{The superposed handwritten reading replaces the underlying reading.}
\psentry{ps_134-5_076-v0020}{20}{insertion}{\inserted{\deleted{\illegible{1}}}}{\hand{}}{The handwritten insertion is incorporated at its indicated position; insertion brackets leave the running text.}
\psentry{ps_134-5_076-v0021}{21}{deletion}{\deleted{written}}{}{The author cancels this material by hand; it leaves the corrected running text.}
\psentry{ps_134-5_076-v0022}{22}{insertion}{\inserted{finally}}{\hand{finally}}{The handwritten insertion is incorporated at its indicated position; insertion brackets leave the running text.}
\psentry{ps_134-5_076-v0023}{23}{deletion}{\deleted{down}}{}{The author cancels this material by hand; it leaves the corrected running text.}
\psentry{ps_134-5_076-v0024}{24}{insertion}{\inserted{sketched}}{\hand{sketched}}{The handwritten insertion is incorporated at its indicated position; insertion brackets leave the running text.}
\psentry{ps_134-5_076-v0025}{25}{typed insertion}{\inserted[typed]{long-winded}}{long-winded}{The typed insertion is incorporated at its indicated position; insertion brackets leave the running text.}
\psentry{ps_134-5_076-v0026}{26}{typewriter deletion}{\xout{awkwardx}}{}{The typewriter strikes out this material; it leaves the corrected running text.}
\psentry{ps_134-5_076-v0027}{27}{typewriter deletion}{\xout{,}}{}{The typewriter strikes out this material; it leaves the corrected running text.}
\psentry{ps_134-5_076-v0028}{28}{typewriter deletion}{\xout{which}}{}{The typewriter strikes out this material; it leaves the corrected running text.}
\psentry{ps_134-5_076-v0029}{29}{insertion}{\inserted{factoring through}}{\hand{factoring through}}{The handwritten insertion is incorporated at its indicated position; insertion brackets leave the running text.}
\psentry{ps_134-5_076-v0030}{30}{deletion}{\deleted{in}}{}{The author cancels this material by hand; it leaves the corrected running text.}
\psentry{ps_134-5_076-v0031}{31}{typewriter deletion}{\xout{which}}{}{The typewriter strikes out this material; it leaves the corrected running text.}
\psentry{ps_134-5_076-v0032}{32}{typewriter deletion}{\xout{by}}{}{The typewriter strikes out this material; it leaves the corrected running text.}
\psentry{ps_134-5_076-v0033}{33}{typed insertion}{\inserted[typed]{of}}{of}{The typed insertion is incorporated at its indicated position; insertion brackets leave the running text.}
\psentry{ps_134-5_076-v0034}{34}{typewriter deletion}{\xout{means}}{}{The typewriter strikes out this material; it leaves the corrected running text.}
\psentry{ps_134-5_076-v0035}{35}{typewriter deletion}{\xout{k}}{}{The typewriter strikes out this material; it leaves the corrected running text.}
\psapparatusend
\fi
\pssource{ps_134-5_077}{134-5/77}{https://grothendieck.umontpellier.fr/web/viewer.html?file=../134-5.pdf\#page=78}
\psnumber{ps_134-5_077-v0001}{1}
\psnumber{ps_134-5_077-v0002}{2}
\psnumber{ps_134-5_077-v0003}{3}
\psnumber{ps_134-5_077-v0004}{4}
\psnumber{ps_134-5_077-v0005}{5}
\psnumber{ps_134-5_077-v0006}{6}
\psnumber{ps_134-5_077-v0007}{7}
\psnumber{ps_134-5_077-v0008}{8}
\psnumber{ps_134-5_077-v0009}{9}
\psnumber{ps_134-5_077-v0010}{10}
\psnumber{ps_134-5_077-v0011}{11}
\psnumber{ps_134-5_077-v0012}{12}
\psnumber{ps_134-5_077-v0013}{13}
\psnumber{ps_134-5_077-v0014}{14}
\psnumber{ps_134-5_077-v0015}{15}
\psnumber{ps_134-5_077-v0016}{16}
\psnumber{ps_134-5_077-v0017}{17}
\psnumber{ps_134-5_077-v0018}{18}
\psnumber{ps_134-5_077-v0019}{19}
\psnumber{ps_134-5_077-v0020}{20}
\psnumber{ps_134-5_077-v0021}{21}
\psnumber{ps_134-5_077-v0022}{22}
\psnumber{ps_134-5_077-v0023}{23}
%% BEGIN SOURCE ps_134-5_077
% Pursuing Stacks (cotes 134-2 to 134-8): archive cote 134-5, site page 77 (PDF page 78), typescript page 264.
% typescript. Draft compared: ps4.tex:3904-3965.
% AI-transcribed by Claude Opus 5.5 (claude-opus-5-5, adaptive thinking, effort high), Message Batches API, batch ps2 (msgbatch_01NVwr1Yuo8qJf5rykeGBXPK), 2026-11-12; reviewed against both supplied scan halves; see CHECK_DECISIONS.tsv.
\scanpage{134-5/77}
\mspage{264}

various homotopy notions in section 51 \var{\inserted[typed]{and 52}}{and 52}{typed insertion}{ps_134-5_077-v0001}.  On the other hand,  it is clear
that the main ideas which are coming to fruition in part IV  all origina-%
ted \var{\xout{in}}{}{typewriter deletion}{ps_134-5_077-v0002} during the awkward grinding process in part III  !

The  ``modelizing story''  so far has \var{\xout{centered around}}{}{typewriter deletion}{ps_134-5_077-v0003}\var{\inserted[typed]{turned out as}}{turned out as}{typed insertion}{ps_134-5_077-v0004} the interplay
of three main \var{\xout{types}}{}{typewriter deletion}{ps_134-5_077-v0005}\var{\inserted[typed]{sets}}{sets}{typed insertion}{ps_134-5_077-v0006} of notions.  One \var{\xout{set}}{}{typewriter deletion}{ps_134-5_077-v0007} is made up with the \var{\xout{``test-notions}}{}{typewriter deletion}{ps_134-5_077-v0008}
``\hand{\uline{\textcolor{black}{test}}}-\hand{\uline{\textcolor{black}{notions}}}'',  centering around the notion of a \hand{\uline{\textcolor{black}{test}}}-\hand{\uline{\textcolor{black}{category}}} ,  as
one giving rise to the most elementary type of ``\hand{\uline{\textcolor{black}{modelizers}}}'' ,  namely the
so called ``\hand{\uline{\textcolor{black}{elementary modelizers}}}'' $(A\hat{}\,,\underline{W}_A)$.  \var{\xout{\unsure{This} seemed}}{}{typewriter deletion}{ps_134-5_077-v0009} This was develop-%
ped in part II  (while part I was concerned with the initial motivation
of the reflections,  namely stacks,  forgotten for the time being !).
The second set of notions concerns the so-called ``\hand{\uline{\textcolor{black}{homotopy}}} \hand{\uline{\textcolor{black}{notions}}}'' ,
developped at some length  in sections 51 to 55,  summarized in the dia-%
gram (2) above.  They constitute the main technical content of part III
of the notes,  with however one major shortcoming:  the relationship
between these notions,  and 0-connectedness notions  (20),  was only par-%
tially understood in part III, \var{\inserted{namely}}{\hand{namely}}{insertion}{ps_134-5_077-v0010} as a one-way relationship merely,  asso-%
ciating to \var{\inserted[typed]{suitable}}{suitable}{typed insertion}{ps_134-5_077-v0011} 0-connectedness notions in a category M,  a corresponding ho-%
motopy structure in M.  The third set of notions \var{\xout{centers around the}}{}{typewriter deletion}{ps_134-5_077-v0012} may
% The handwritten opening quotation strokes overlap a dark earlier mark; no underlying character is securely recoverable.
be called \hand{``}\hand{\uline{\textcolor{black}{asphericity}}} \hand{\uline{\textcolor{black}{notions}}}'',  they center around the notions of
\hand{\uline{\textcolor{black}{aspheric}}} \hand{\uline{\textcolor{black}{objects}}} and \hand{\uline{\textcolor{black}{aspheric}}} \hand{\uline{\textcolor{black}{maps}}} (in a category endowed with a\var{\xout{n}}{}{typewriter deletion}{ps_134-5_077-v0013} so-%
called \var{\xout{a}}{}{typewriter deletion}{ps_134-5_077-v0014} \hand{\uline{\textcolor{black}{asphericity}}} \hand{\uline{\textcolor{black}{structure}}} ),  and more specifically around the
notion of \var{\inserted[typed]{an}}{an}{typed insertion}{ps_134-5_077-v0015} aspheric map in (Cat),  whose formal properties turn out to be
the \var{\xout{main source}}{}{typewriter deletion}{ps_134-5_077-v0016} key for the development of a theory of asphericity struc-%
tures.  The first and the third set of notions  (namely test notions  and
asphericity notions)  depend on the choice of a ``\hand{\uline{\textcolor{black}{basic localizer}}}'' \var{\inserted[typed]{$\underline{W}$}}{$\underline{W}$}{typed insertion}{ps_134-5_077-v0017} in (Cat)\hand{,}
whereas the second set,  namely homotopy notions,  is ``absolute'' \var{\deleted{namely}}{}{deletion}{ps_134-5_077-v0018}\hand{i.e.}
does not depend on any such choice,  nor an any knowledge of homotopy or
(co)homology invariants.

Whereas the test notions are essentially concerned with modelizers,
namely getting descriptions of the category of homotopy types (Hot)  in
terms of elementary modelizers $A\hat{}$ (as \var{\inserted{being}}{\hand{being}}{insertion}{ps_134-5_077-v0019} $\underline{W}_A^{\ -1}A\hat{}$),  it appears that the
homotopy notions\hand{,} as well as the asphericity notions\hand{,} are independant of
any modelizing \var{\inserted{notions and}}{\hand{notions and}}{insertion}{ps_134-5_077-v0020} assumptions.  In a deductive presentation of the theory,
the % The stroke beside fact below the interlinear general is an insertion locator, not a deletion.
test notions would come last,  whereas they came first in these notes
-  as an illustration of the \var{\inserted[typed]{general}}{general}{typed insertion}{ps_134-5_077-v0021} fact that the deductive approach will present
things roughly in opposite order in which they have been discovered!  The
test notions,  as an outcome \var{\inserted{of the attempt}}{\hand{of the attempt}}{insertion}{ps_134-5_077-v0022} to get a picture of modelizers,  have kept
acting as a constant \var{\xout{source of}}{}{typewriter deletion}{ps_134-5_077-v0023} guideline in the whole reflection,  even
though technically speaking they are ``irrelevant'' for the development of
\margin{77}

%% END SOURCE ps_134-5_077
\ifnum\EditionMode=1
\psapparatusstart
\psentry{ps_134-5_077-v0001}{1}{typed insertion}{\inserted[typed]{and 52}}{and 52}{The typed insertion is incorporated at its indicated position; insertion brackets leave the running text.}
\psentry{ps_134-5_077-v0002}{2}{typewriter deletion}{\xout{in}}{}{The typewriter strikes out this material; it leaves the corrected running text.}
\psentry{ps_134-5_077-v0003}{3}{typewriter deletion}{\xout{centered around}}{}{The typewriter strikes out this material; it leaves the corrected running text.}
\psentry{ps_134-5_077-v0004}{4}{typed insertion}{\inserted[typed]{turned out as}}{turned out as}{The typed insertion is incorporated at its indicated position; insertion brackets leave the running text.}
\psentry{ps_134-5_077-v0005}{5}{typewriter deletion}{\xout{types}}{}{The typewriter strikes out this material; it leaves the corrected running text.}
\psentry{ps_134-5_077-v0006}{6}{typed insertion}{\inserted[typed]{sets}}{sets}{The typed insertion is incorporated at its indicated position; insertion brackets leave the running text.}
\psentry{ps_134-5_077-v0007}{7}{typewriter deletion}{\xout{set}}{}{The typewriter strikes out this material; it leaves the corrected running text.}
\psentry{ps_134-5_077-v0008}{8}{typewriter deletion}{\xout{``test-notions}}{}{The typewriter strikes out this material; it leaves the corrected running text.}
\psentry{ps_134-5_077-v0009}{9}{typewriter deletion}{\xout{\unsure{This} seemed}}{}{The typewriter strikes out this material; it leaves the corrected running text.}
\psentry{ps_134-5_077-v0010}{10}{insertion}{\inserted{namely}}{\hand{namely}}{The handwritten insertion is incorporated at its indicated position; insertion brackets leave the running text.}
\psentry{ps_134-5_077-v0011}{11}{typed insertion}{\inserted[typed]{suitable}}{suitable}{The typed insertion is incorporated at its indicated position; insertion brackets leave the running text.}
\psentry{ps_134-5_077-v0012}{12}{typewriter deletion}{\xout{centers around the}}{}{The typewriter strikes out this material; it leaves the corrected running text.}
\psentry{ps_134-5_077-v0013}{13}{typewriter deletion}{\xout{n}}{}{The typewriter strikes out this material; it leaves the corrected running text.}
\psentry{ps_134-5_077-v0014}{14}{typewriter deletion}{\xout{a}}{}{The typewriter strikes out this material; it leaves the corrected running text.}
\psentry{ps_134-5_077-v0015}{15}{typed insertion}{\inserted[typed]{an}}{an}{The typed insertion is incorporated at its indicated position; insertion brackets leave the running text.}
\psentry{ps_134-5_077-v0016}{16}{typewriter deletion}{\xout{main source}}{}{The typewriter strikes out this material; it leaves the corrected running text.}
\psentry{ps_134-5_077-v0017}{17}{typed insertion}{\inserted[typed]{$\underline{W}$}}{$\underline{W}$}{The typed insertion is incorporated at its indicated position; insertion brackets leave the running text.}
\psentry{ps_134-5_077-v0018}{18}{deletion}{\deleted{namely}}{}{The author cancels this material by hand; it leaves the corrected running text.}
\psentry{ps_134-5_077-v0019}{19}{insertion}{\inserted{being}}{\hand{being}}{The handwritten insertion is incorporated at its indicated position; insertion brackets leave the running text.}
\psentry{ps_134-5_077-v0020}{20}{insertion}{\inserted{notions and}}{\hand{notions and}}{The handwritten insertion is incorporated at its indicated position; insertion brackets leave the running text.}
\psentry{ps_134-5_077-v0021}{21}{typed insertion}{\inserted[typed]{general}}{general}{The typed insertion is incorporated at its indicated position; insertion brackets leave the running text.}
\psentry{ps_134-5_077-v0022}{22}{insertion}{\inserted{of the attempt}}{\hand{of the attempt}}{The handwritten insertion is incorporated at its indicated position; insertion brackets leave the running text.}
\psentry{ps_134-5_077-v0023}{23}{typewriter deletion}{\xout{source of}}{}{The typewriter strikes out this material; it leaves the corrected running text.}
\psapparatusend
\fi
\pssource{ps_134-5_078}{134-5/78}{https://grothendieck.umontpellier.fr/web/viewer.html?file=../134-5.pdf\#page=79}
\psnumber{ps_134-5_078-v0001}{1}
\psnumber{ps_134-5_078-v0002}{2}
\psnumber{ps_134-5_078-v0003}{3}
\psnumber{ps_134-5_078-v0004}{4}
\psnumber{ps_134-5_078-v0005}{5}
\psnumber{ps_134-5_078-v0006}{6}
\psnumber{ps_134-5_078-v0007}{7}
\psnumber{ps_134-5_078-v0008}{8}
\psnumber{ps_134-5_078-v0009}{9}
\psnumber{ps_134-5_078-v0010}{10}
\psnumber{ps_134-5_078-v0011}{11}
\psnumber{ps_134-5_078-v0012}{12}
\psnumber{ps_134-5_078-v0013}{13}
\psnumber{ps_134-5_078-v0014}{14}
\psnumber{ps_134-5_078-v0015}{15}
\psnumber{ps_134-5_078-v0016}{16}
\psnumber{ps_134-5_078-v0017}{17}
\psnumber{ps_134-5_078-v0018}{18}
\psnumber{ps_134-5_078-v0019}{19}
\psnumber{ps_134-5_078-v0020}{20}
\psnumber{ps_134-5_078-v0021}{21}
\psnumber{ps_134-5_078-v0022}{22}
\psnumber{ps_134-5_078-v0023}{23}
\psnumber{ps_134-5_078-v0024}{24}
\psnumber{ps_134-5_078-v0025}{25}
\psnumber{ps_134-5_078-v0026}{26}
\psnumber{ps_134-5_078-v0027}{27}
\psnumber{ps_134-5_078-v0028}{28}
%% BEGIN SOURCE ps_134-5_078
% Pursuing Stacks (cotes 134-2 to 134-8): archive cote 134-5, site page 78 (PDF page 79), typescript page 265.
% typescript. Draft compared: ps4.tex:3952-4018.
% AI-transcribed by Claude Opus 5.5 (claude-opus-5-5, adaptive thinking, effort high), Message Batches API, batch ps2 (msgbatch_01NVwr1Yuo8qJf5rykeGBXPK), 2026-11-12; reviewed against both supplied scan halves; see CHECK_DECISIONS.tsv.
\scanpage{134-5/78}
\mspage{265}

the main properties of homotopy \var{\xout{notions}}{}{typewriter deletion}{ps_134-5_078-v0001} and asphericity notions and their
interplay.

By the end of part III and the beginning of part IV, there has
been some floating in my mind as to w\var{\over{\unsure{h}}{he}}{\hand{he}}{writing over}{ps_134-5_078-v0002}ther which among the two structures\hand{,}
namely contractibility structures or asphericity structures, should be
considered as ``the'' key structure for an understanding of the modelizing
story. There was a (justified) feeling, expressed first at the begin-%
ning of section 67, that in some sense, asphericity structures were
``more general'' than contractibility structures, which caused me for a
while to view them as the more ``basic'' ones. I would be more tempted at
present to hold the opposite view. \hand{\uline{\textcolor{black}{The notion of a contractibility}}}
\hand{\uline{\textcolor{black}{structure now appears as a kind of hinge between the two main sets of}}}
\hand{\uline{\textcolor{black}{notions besides the test notions,}}} \hand{\uline{\textcolor{black}{namely between}}} \var{\xout{the}}{}{typewriter deletion}{ps_134-5_078-v0003} \hand{\uline{\textcolor{black}{homotopy notions}}}
\hand{\uline{\textcolor{black}{and asphericity notions}}}. On the one hand, as displayed in diagram (2),
the notion of a contractibility structure \var{\xout{is the}}{}{typewriter deletion}{ps_134-5_078-v0004} appears as the most
stringent one among the four main \var{\deleted{sets}}{}{deletion}{ps_134-5_078-v0005}\var{\inserted{types}}{\hand{types}}{insertion}{ps_134-5_078-v0006} of homotopy \var{\xout{notions}}{}{typewriter deletion}{ps_134-5_078-v0007} structures.
On the other hand, by (1) it can be equally viewed as being a special
case of \var{\inserted{an}}{\hand{an}}{insertion}{ps_134-5_078-v0008} asphericity structure, and as such it \var{\xout{can be expressed b}}{}{typewriter deletion}{ps_134-5_078-v0009} gives
rise to (and can be expressed by) either one of the following \var{\deleted{four}}{}{deletion}{ps_134-5_078-v0010} \hand{three}
asphericity-flavored data on a category M (for any given basic localizer
$\underline{W}$ satisfying Loc 4)):

\margin{(21} \var{\deleted{(21)}}{}{deletion}{ps_134-5_078-v0011}
\[
\hand{\Biggl\{}
\begin{array}{l}
M_{\underline{W}} \hand{\subset} \mathrm{Ob}\,M \quad , \quad \underline{W}_{M_c} \text{ or simply } \underline{W}_M \hand{\subset} F\ell(M)\\
\hand{\phi}^{\underline{W}}_{M} \text{ or } \hand{\phi}_M : \var{\deleted{M \hand{\to} \illegible{1}\ \mathrm{Hot}_{\underline{W}}}}{}{deletion}{ps_134-5_078-v0012} \hand{M \longrightarrow \mathrm{Hot}_{\underline{W}}}\\
\var{\deleted{\hand{\pi}_0 = \hand{\phi}^{\hand{\underline{W}}\,\underline{W}_0}_{M} : M \hand{\longrightarrow} (\mathrm{Sets})}}{}{deletion}{ps_134-5_078-v0013} \var{\deleted{\illegible{1}\ \hand{M \longrightarrow (\unsure{\mathrm{Hot}})}}}{}{deletion}{ps_134-5_078-v0014} \hand{).}
\end{array}
% The lower extension of the handwritten grouping brace is hatched beside the cancelled third line.
\]
\var{\margin{\deleted{$\varphi^{\underline{W}_o}_{M}$}}}{}{marginal note revised}{ps_134-5_078-v0015}

Restricting to the case when $\underline{W}$ is either $\underline{W}_{0}$ or $\underline{W}_{\mathrm{oo}}$, we get the corres-%
ponding structures on M , namely the three structures $M_0, W_0, \hand{\pi}_0$ of (20)
plus the \var{\over{two}{three}}{\hand{three}}{writing over}{ps_134-5_078-v0016} extra structures :
\[
(22) \qquad M_\mathrm{oo} = M_{\underline{W}_\mathrm{oo}} \hand{\subset} \mathrm{Ob}\,M \ , \quad W_\mathrm{oo} = (\underline{W}_\mathrm{oo})_M \hand{\subset} F\ell\,(M) \ , \ \hand{\varphi_M : M \longrightarrow (\mathrm{Hot}\var{\deleted{_{\underline{W}}}}{}{deletion}{ps_134-5_078-v0017})}
\]
which can be referred to as $\mathrm{oo}$-\hand{\uline{\textcolor{black}{connected}}} \hand{\uline{\textcolor{black}{objects}}} \hand{\uline{\textcolor{black}{of}}} M, \var{\xout{and}}{}{typewriter deletion}{ps_134-5_078-v0018} $\mathrm{oo}$-\hand{\uline{\textcolor{black}{connected}}}
\margin{$\times$} \hand{\uline{\textcolor{black}{arrows}}} of M\var{\over{.}{,}}{\hand{,}}{writing over}{ps_134-5_078-v0019} \var{\inserted{(and ``\uline{canonical functor}'' from M to Hot).}}{\hand{(and ``\uline{canonical functor}'' from M to Hot).}}{insertion}{ps_134-5_078-v0020} Putting together (4), (20)\hand{,}(22), we see that a contractibi-%
lity structure \var{\inserted{$M_c$}}{\hand{$M_c$}}{insertion}{ps_134-5_078-v0021} gives rise to \var{\deleted{eight other types of}}{}{deletion}{ps_134-5_078-v0022} \var{\inserted{\deleted{nine} ten different}}{\hand{ ten different}}{insertion}{ps_134-5_078-v0023} structures \var{\inserted{(including $M_c$ itself)}}{\hand{(including $M_c$ itself)}}{insertion}{ps_134-5_078-v0024} and is deter-%
mined by each one o\var{\over{n}{f}}{\hand{f}}{writing over}{ps_134-5_078-v0025} these, \var{\deleted{namely:}}{}{deletion}{ps_134-5_078-v0026}

\var{\deleted{\xout{(23)}x \quad $J_c$ Int(M) , $W_c$ $F\ell(M)$ , $R_c$ $F\ell(M)$x$F\ell(M)$}}{}{deletion}{ps_134-5_078-v0027}
\var{\deleted{$M_0$ \quad Ob M , \quad $W_0$}}{}{deletion}{ps_134-5_078-v0028}

\margin{78}

%% END SOURCE ps_134-5_078
\ifnum\EditionMode=1
\psapparatusstart
\psentry{ps_134-5_078-v0001}{1}{typewriter deletion}{\xout{notions}}{}{The typewriter strikes out this material; it leaves the corrected running text.}
\psentry{ps_134-5_078-v0002}{2}{writing over}{\over{\unsure{h}}{he}}{\hand{he}}{The superposed handwritten reading replaces the underlying reading.}
\psentry{ps_134-5_078-v0003}{3}{typewriter deletion}{\xout{the}}{}{The typewriter strikes out this material; it leaves the corrected running text.}
\psentry{ps_134-5_078-v0004}{4}{typewriter deletion}{\xout{is the}}{}{The typewriter strikes out this material; it leaves the corrected running text.}
\psentry{ps_134-5_078-v0005}{5}{deletion}{\deleted{sets}}{}{The author cancels this material by hand; it leaves the corrected running text.}
\psentry{ps_134-5_078-v0006}{6}{insertion}{\inserted{types}}{\hand{types}}{The handwritten insertion is incorporated at its indicated position; insertion brackets leave the running text.}
\psentry{ps_134-5_078-v0007}{7}{typewriter deletion}{\xout{notions}}{}{The typewriter strikes out this material; it leaves the corrected running text.}
\psentry{ps_134-5_078-v0008}{8}{insertion}{\inserted{an}}{\hand{an}}{The handwritten insertion is incorporated at its indicated position; insertion brackets leave the running text.}
\psentry{ps_134-5_078-v0009}{9}{typewriter deletion}{\xout{can be expressed b}}{}{The typewriter strikes out this material; it leaves the corrected running text.}
\psentry{ps_134-5_078-v0010}{10}{deletion}{\deleted{four}}{}{The author cancels this material by hand; it leaves the corrected running text.}
\psentry{ps_134-5_078-v0011}{11}{deletion}{\deleted{(21)}}{}{The author cancels this material by hand; it leaves the corrected running text.}
\psentry{ps_134-5_078-v0012}{12}{deletion}{$\deleted{M \hand{\to} \illegible{1}\ \mathrm{Hot}_{\underline{W}}}$}{}{The author cancels this material by hand; it leaves the corrected running text.}
\psentry{ps_134-5_078-v0013}{13}{deletion}{$\deleted{\hand{\pi}_0 = \hand{\phi}^{\hand{\underline{W}}\,\underline{W}_0}_{M} : M \hand{\longrightarrow} (\mathrm{Sets})}$}{}{The author cancels this material by hand; it leaves the corrected running text.}
\psentry{ps_134-5_078-v0014}{14}{deletion}{$\deleted{\illegible{1}\ \hand{M \longrightarrow (\unsure{\mathrm{Hot}})}}$}{}{The author cancels this material by hand; it leaves the corrected running text.}
\psentry{ps_134-5_078-v0015}{15}{marginal note revised}{\margin{\deleted{$\varphi^{\underline{W}_o}_{M}$}}}{}{The author's internal cancellations and replacements in this marginal note are applied at its source-line call.}
\psentry{ps_134-5_078-v0016}{16}{writing over}{\over{two}{three}}{\hand{three}}{The superposed handwritten reading replaces the underlying reading.}
\psentry{ps_134-5_078-v0017}{17}{deletion}{$\deleted{_{\underline{W}}}$}{}{The author cancels this material by hand; it leaves the corrected running text.}
\psentry{ps_134-5_078-v0018}{18}{typewriter deletion}{\xout{and}}{}{The typewriter strikes out this material; it leaves the corrected running text.}
\psentry{ps_134-5_078-v0019}{19}{writing over}{\over{.}{,}}{\hand{,}}{The superposed handwritten reading replaces the underlying reading.}
\psentry{ps_134-5_078-v0020}{20}{insertion}{\inserted{(and ``\uline{canonical functor}'' from M to Hot).}}{\hand{(and ``\uline{canonical functor}'' from M to Hot).}}{The handwritten insertion is incorporated at its indicated position; insertion brackets leave the running text.}
\psentry{ps_134-5_078-v0021}{21}{insertion}{\inserted{$M_c$}}{\hand{$M_c$}}{The handwritten insertion is incorporated at its indicated position; insertion brackets leave the running text.}
\psentry{ps_134-5_078-v0022}{22}{deletion}{\deleted{eight other types of}}{}{The author cancels this material by hand; it leaves the corrected running text.}
\psentry{ps_134-5_078-v0023}{23}{insertion}{\inserted{\deleted{nine} ten different}}{\hand{ ten different}}{The handwritten insertion is incorporated at its indicated position; insertion brackets leave the running text.}
\psentry{ps_134-5_078-v0024}{24}{insertion}{\inserted{(including $M_c$ itself)}}{\hand{(including $M_c$ itself)}}{The handwritten insertion is incorporated at its indicated position; insertion brackets leave the running text.}
\psentry{ps_134-5_078-v0025}{25}{writing over}{\over{n}{f}}{\hand{f}}{The superposed handwritten reading replaces the underlying reading.}
\psentry{ps_134-5_078-v0026}{26}{deletion}{\deleted{namely:}}{}{The author cancels this material by hand; it leaves the corrected running text.}
\psentry{ps_134-5_078-v0027}{27}{deletion}{\deleted{\xout{(23)}x \quad $J_c$ Int(M) , $W_c$ $F\ell(M)$ , $R_c$ $F\ell(M)$x$F\ell(M)$}}{}{The author cancels this material by hand; it leaves the corrected running text.}
\psentry{ps_134-5_078-v0028}{28}{deletion}{\deleted{$M_0$ \quad Ob M , \quad $W_0$}}{}{The author cancels this material by hand; it leaves the corrected running text.}
\psapparatusend
\fi
\pssource{ps_134-5_079}{134-5/79}{https://grothendieck.umontpellier.fr/web/viewer.html?file=../134-5.pdf\#page=80}
\psnumber{ps_134-5_079-v0001}{1}
\psnumber{ps_134-5_079-v0002}{2}
\psnumber{ps_134-5_079-v0003}{3}
\psnumber{ps_134-5_079-v0004}{4}
\psnumber{ps_134-5_079-v0005}{5}
\psnumber{ps_134-5_079-v0006}{6}
\psnumber{ps_134-5_079-v0007}{7}
\psnumber{ps_134-5_079-v0008}{8}
%% BEGIN SOURCE ps_134-5_079
% Pursuing Stacks (cotes 134-2 to 134-8): archive cote 134-5, site page 79 (PDF page 80), typescript page 266.
% typescript. Draft compared: ps4.tex:4000-4101.
% AI-transcribed by Claude Opus 5.5 (claude-opus-5-5, adaptive thinking, effort high), Message Batches API, batch ps2 (msgbatch_01NVwr1Yuo8qJf5rykeGBXPK), 2026-11-12; reviewed against both supplied scan halves; see CHECK_DECISIONS.tsv.
\scanpage{134-5/79}
\mspage{266}

\[
(23)\hand{\Biggl\{}\begin{array}{llll}
M_c\hand{\subset}\mathrm{Ob}\ M\ , & J_c\hand{\subset}\mathrm{Int}(M)\ , & W_c\hand{\subset}\mathrm{F}\ell(M)\ , & R_c\hand{\subset}\mathrm{F}\ell(M)\mathrm{x}\mathrm{F}\ell(M)\\
M_o\hand{\subset}\mathrm{Ob}\ M\ \ , & & W_o\hand{\subset}\mathrm{F}\ell(M)\ , & \hand{\pi}_o\ :\ M\hand{\longrightarrow}(\mathrm{Sets})\\
M_\mathrm{oo}\hand{\subset}\mathrm{Ob}\ M\ \ , & & W_\mathrm{oo}\hand{\subset}\mathrm{F}\ell(M)\ , & \hand{\phi}_M:\ M\hand{\longrightarrow}(\mathrm{Hot})=\underline{W}_\mathrm{oo}^{-1}(\mathrm{Cat}),
\end{array}
\]

namely: \hand{\uline{\textcolor{black}{contractible objects}}}, \hand{\uline{\textcolor{black}{homotopy intervals}}}, \hand{\uline{\textcolor{black}{homotopisms}}}, \hand{\uline{\textcolor{black}{homotopy}}}
\hand{\uline{\textcolor{black}{relation for maps}}}, 0-\hand{\uline{\textcolor{black}{connected objects}}}, 0-\hand{\uline{\textcolor{black}{connected maps}}},\hand{(*)} the \hand{\uline{\textcolor{black}{connected}}}
\hand{\uline{\textcolor{black}{components functor}}}, $\mathrm{oo}$-\hand{\uline{\textcolor{black}{connected}}} (or ``aspheric'' , more accurately $\underline{W}_\mathrm{oo}$-as-%
pheric) \hand{\uline{\textcolor{black}{objects}}}, \var{\xout{oo-connected \unsure{(}}}{}{typewriter deletion}{ps_134-5_079-v0001} $\underline{W}_\mathrm{oo}$-\hand{\uline{\textcolor{black}{equivalences}}} (or simply ``weak
equivalences''), and the (would-be ``modelizing'') \hand{\uline{\textcolor{black}{canonical functor}}} \hand{\uline{\textcolor{black}{from}}} M
\hand{\uline{\textcolor{black}{to homotopy types}}}. Moreover, it should be remembered that, just as
the structures $J_c$,$W_c$,$R_c$ are by no means unrestricted (the fact that they
stem from a contractibility structure being a substantial restriction),
the two asphericity structures (namely, \var{\inserted{the}}{\hand{the}}{insertion}{ps_134-5_079-v0002} 0-asphericity structure $M_o$ and
the $\mathrm{oo}$-asphericity structure $M_\mathrm{oo}$ ) are \var{\xout{xx}}{}{typewriter deletion}{ps_134-5_079-v0003} subject to extra-conditions which
will be written down below, implying \var{\xout{thatxxx}}{}{typewriter deletion}{ps_134-5_079-v0004} among others that they are
totally aspheric \var{\inserted{(hence $\pi_0$ and $\varphi_M$ are compatible with finite products)}}{\hand{(hence $\pi_0$ and $\varphi_M$ are compatible with finite products)}}{insertion}{ps_134-5_079-v0005}.

We may complement the ten structural data (23) above, by the following \var{\inserted[typed]{two}}{two}{typed insertion}{ps_134-5_079-v0006}
\[ (24) \qquad R_o\ ,\ R_\mathrm{oo}\ \hand{\subset}\ \mathrm{F}\ell(M)\ \mathrm{x}\ \mathrm{F}\ell(M) \qquad , \]
defined by
\[ (f,g)\hand{\in}R_o\hand{\Longleftrightarrow}\hand{\Big(}\hand{\pi}_o(f)=\hand{\pi}_o(g)\hand{\Big)}\ \ ,\quad (f,g)\hand{\in}R_\mathrm{oo}\ \var{\over{\illegible{1}}{\Longleftrightarrow}}{\hand{\Longleftrightarrow}}{writing over}{ps_134-5_079-v0007}\hand{\Big(}\hand{\phi}_M(f)=\hand{\phi}_M(g)\hand{\Big)}\ . \]
More generally, for any basic localizer $\underline{W}$ , we may define
\[ (25) \qquad R_{\underline{W}}\hand{\subset}\ \mathrm{F}\ell(M)\mathrm{x}\mathrm{F}\ell(M)\ ,\qquad (f,g)\hand{\in}R_{\underline{W}}\ \hand{\Longleftrightarrow}\ \hand{\Big(}\hand{\phi}_M^{\hand{\underline{W}}}(f)=\hand{\phi}_M^{\hand{\underline{W}}}(g)\hand{\Big)}\ . \]
It is not clear however that $M_c$ , or equivalently $M_{\underline{W}}$ or $\underline{W}_M$ , can be
recovered\var{\xout{d}}{}{typewriter deletion}{ps_134-5_079-v0008} in terms of the equivalence relation $R_{\underline{W}}$ among maps of M.

Among the three series of structures appearing in (23), we have
the tautological relations:
% Hand-drawn arrows are marked blue; typed objects remain plain.
\[
(26)\ \hand{\Biggl\{}\begin{array}{l}
M_c\hand{\subset}M_\mathrm{oo}\hand{\subset}M_o\\
W_c\hand{\subset}W_\mathrm{oo}\hand{\subset}W_o\\
R_c\hand{\subset}R_\mathrm{oo}\hand{\subset}R_o\\
\diagram{\begin{tikzcd}
M \ar[r,"\hand{\varphi_M}", color=blue] \ar[rr, bend right, "\hand{\pi_0}"', color=blue] \& (\mathrm{Hot}) \ar[r,"\hand{\underline{\pi}_0}", color=blue] \& (\mathrm{Sets})
\end{tikzcd}}
\qquad \hand{(commutative).}
\end{array}
\]
Finally, we may also display some of the main functors defined in terms
of a given contractibility structure $M_c$ :
\[
(27)\qquad
\diagram{\begin{tikzcd}[column sep=small]
M \ar[r, color=blue] \& \hand{\overline{\textcolor{black}{M}}} \ar[r, color=blue] \ar[d, phantom, "\hand{\parallel}"] \& \mathrm{Hot}_M \ar[r, color=blue] \ar[d, phantom, "\hand{\parallel}"] \& (\mathrm{Hot}) \ar[r, color=blue] \ar[d, phantom, "\hand{\parallel}"] \& \mathrm{Hot}_{\underline{W}} \ar[r, color=blue] \& (\mathrm{Sets}) \ar[d, phantom, "\hand{\wr}"] \\
 \& \hand{W_c^{-1}M} \& \hand{W_\infty^{-1}M} \& \hand{\mathrm{Hot}_{\underline{W}_\infty}} \& \& \hand{\mathrm{Hot}_{\underline{W}_o}}
\end{tikzcd}}
\ .
\]

\margin{(*) This name is inadequate, cf. note p. 262.}

\margin{79}

%% END SOURCE ps_134-5_079
\ifnum\EditionMode=1
\psapparatusstart
\psentry{ps_134-5_079-v0001}{1}{typewriter deletion}{\xout{oo-connected \unsure{(}}}{}{The typewriter strikes out this material; it leaves the corrected running text.}
\psentry{ps_134-5_079-v0002}{2}{insertion}{\inserted{the}}{\hand{the}}{The handwritten insertion is incorporated at its indicated position; insertion brackets leave the running text.}
\psentry{ps_134-5_079-v0003}{3}{typewriter deletion}{\xout{xx}}{}{The typewriter strikes out this material; it leaves the corrected running text.}
\psentry{ps_134-5_079-v0004}{4}{typewriter deletion}{\xout{thatxxx}}{}{The typewriter strikes out this material; it leaves the corrected running text.}
\psentry{ps_134-5_079-v0005}{5}{insertion}{\inserted{(hence $\pi_0$ and $\varphi_M$ are compatible with finite products)}}{\hand{(hence $\pi_0$ and $\varphi_M$ are compatible with finite products)}}{The handwritten insertion is incorporated at its indicated position; insertion brackets leave the running text.}
\psentry{ps_134-5_079-v0006}{6}{typed insertion}{\inserted[typed]{two}}{two}{The typed insertion is incorporated at its indicated position; insertion brackets leave the running text.}
\psentry{ps_134-5_079-v0007}{7}{writing over}{$\over{\illegible{1}}{\Longleftrightarrow}$}{$\hand{\Longleftrightarrow}$}{The superposed handwritten reading replaces the underlying reading.}
\psentry{ps_134-5_079-v0008}{8}{typewriter deletion}{\xout{d}}{}{The typewriter strikes out this material; it leaves the corrected running text.}
\psapparatusend
\fi
\pssource{ps_134-5_080}{134-5/80}{https://grothendieck.umontpellier.fr/web/viewer.html?file=../134-5.pdf\#page=81}
\psnumber{ps_134-5_080-v0001}{1}
\psnumber{ps_134-5_080-v0002}{2}
\psnumber{ps_134-5_080-v0003}{3}
\psnumber{ps_134-5_080-v0004}{4}
\psnumber{ps_134-5_080-v0005}{5}
\psnumber{ps_134-5_080-v0006}{6}
\psnumber{ps_134-5_080-v0007}{7}
\psnumber{ps_134-5_080-v0008}{8}
\psnumber{ps_134-5_080-v0009}{9}
\psnumber{ps_134-5_080-v0010}{10}
%% BEGIN SOURCE ps_134-5_080
% Pursuing Stacks (cotes 134-2 to 134-8): archive cote 134-5, site page 80 (PDF page 81), typescript page 267.
% typescript. Draft compared: ps4.tex:4077-4157.
% AI-transcribed by Claude Opus 5.5 (claude-opus-5-5, adaptive thinking, effort high), Message Batches API, batch ps2 (msgbatch_01NVwr1Yuo8qJf5rykeGBXPK), 2026-11-12; reviewed against both supplied scan halves; see CHECK_DECISIONS.tsv.
\scanpage{134-5/80}
\mspage{267}

26.6.

82. Yesterday I have been busy mainly with the readjustment of \var{\xout{general}}{}{typewriter deletion}{ps_134-5_080-v0001} \var{\inserted[typed]{the overall}}{the overall}{typed insertion}{ps_134-5_080-v0002}
perspective on the main notions developped so far, which has sprung from
the new fact stated at the beginning of yesterday's notes: namely that an
arbitrary contractibility structure $M_c$ (on a category M stable under finite
products) can be recovered in terms of the associated $\underline{W}$-asphericity struc-%
ture, where $\underline{W}$ is any basic localizer satisfying (besides the conditions
Loc 1) to Loc 3) of p. 213-214) the extra assumption Loc 4) of p.260. It
seems about time now to enter into some \var{\xout{technical}}{}{typewriter deletion}{ps_134-5_080-v0003} a little more technical
specifications along the same lines - and to start with, give a proof
of this ``new fact'' ! Let's state it again in full:

\hand{\uline{\textcolor{black}{Theorem}}} 1 . Let M be a $\underline{U}$-category stable under finite products, $M_c$ \var{\inserted[typed]{$\hand{\subset}$ Ob M}}{$\hand{\subset}$ Ob M}{typed insertion}{ps_134-5_080-v0004} a con-%
tractibility structure on M, admitting a small full subcategory C which
generates the structure. Let moreover $\underline{W}$ be any basic localizer satisfying
condition Loc 4) (compatibility with the $\hand{\underline{\pi}}_0$-functor (Cat)$\hand{\longrightarrow}$(Sets)), and
let $M_{\underline{W}} \hand{\subset}$ Ob M the $\underline{W}$-asphericity structure generated by $M_c$ , $W = \underline{W}_{M_c}$ the
corresponding set of ``$\underline{W}$-equivalences'' or ``weak equivalences'' in $F\ell$(M). Con%
% The l in Fl(M) carries the visible diagonal slash-like overstrike.
-sider the homotopy structure $h_W$ associated to W, i.e. (cf. section 54)
the homotopy \var{\xout{interval}}{}{typewriter deletion}{ps_134-5_080-v0005} structure associated to the homotopy interval struc-%
ture \var{\inserted{J}}{\hand{J}}{insertion}{ps_134-5_080-v0006} generated by the set \var{\inserted{$J_0$}}{\hand{$J_0$}}{insertion}{ps_134-5_080-v0007} of all intervals
\[ \underline{I} = (I, \hand{\delta}_0, \hand{\delta}_1) \]
in M such that
\[ I \hand{\in} M_{\underline{W}} \qquad . \]
Then $h_W$ \var{\xout{\unsure{s}ame}}{}{typewriter deletion}{ps_134-5_080-v0008} is the homotopy structure on M associated to the contractibi-%
lity structure $M_c$ , and hence $M_c$ can be described in terms of $M_{\underline{W}}$ (or of
$\underline{W}_{M_c} = W$ ) as the set of objects which are contractible for $h_W$ , i.e. such
that the map $x \hand{\longrightarrow} e_M$ is an $h_W$-homotopism.

\hand{\uline{\textcolor{black}{Proof}}}. Let $M'_c$ be the set of $h_W$-contractible objects of M, clearly we have
\[ M_c \hand{\subset} M'_c \hand{\subset} M_{\underline{W}} \overset{\mathrm{def}}{=} M_a . \]
The theorem amounts to saying that $M_c$ generates the homotopy interval
structure J (by which we mean that the set of intervals of M made up with
objects of $M_c$ generates the structure J). Indeed, because of $M_c \hand{\subset} M'_c$
this will imply that J is associated to a contractibility structure, \var{\xout{namely}}{}{typewriter deletion}{ps_134-5_080-v0009}
\var{\xout{ly}}{}{typewriter deletion}{ps_134-5_080-v0010} namely to $M'_c$ . But for an object x of M to be in $M'_c$ i.e. to be contrac%
tible for the structure J, amounts to be contractible for $M_c$ , and hence
by saturation of $M_c$ , to be in $M_c$ , hence $M'_c = M_c$ - which yields what we

\margin{80}

%% END SOURCE ps_134-5_080
\ifnum\EditionMode=1
\psapparatusstart
\psentry{ps_134-5_080-v0001}{1}{typewriter deletion}{\xout{general}}{}{The typewriter strikes out this material; it leaves the corrected running text.}
\psentry{ps_134-5_080-v0002}{2}{typed insertion}{\inserted[typed]{the overall}}{the overall}{The typed insertion is incorporated at its indicated position; insertion brackets leave the running text.}
\psentry{ps_134-5_080-v0003}{3}{typewriter deletion}{\xout{technical}}{}{The typewriter strikes out this material; it leaves the corrected running text.}
\psentry{ps_134-5_080-v0004}{4}{typed insertion}{\inserted[typed]{$\hand{\subset}$ Ob M}}{$\hand{\subset}$ Ob M}{The typed insertion is incorporated at its indicated position; insertion brackets leave the running text.}
\psentry{ps_134-5_080-v0005}{5}{typewriter deletion}{\xout{interval}}{}{The typewriter strikes out this material; it leaves the corrected running text.}
\psentry{ps_134-5_080-v0006}{6}{insertion}{\inserted{J}}{\hand{J}}{The handwritten insertion is incorporated at its indicated position; insertion brackets leave the running text.}
\psentry{ps_134-5_080-v0007}{7}{insertion}{\inserted{$J_0$}}{\hand{$J_0$}}{The handwritten insertion is incorporated at its indicated position; insertion brackets leave the running text.}
\psentry{ps_134-5_080-v0008}{8}{typewriter deletion}{\xout{\unsure{s}ame}}{}{The typewriter strikes out this material; it leaves the corrected running text.}
\psentry{ps_134-5_080-v0009}{9}{typewriter deletion}{\xout{namely}}{}{The typewriter strikes out this material; it leaves the corrected running text.}
\psentry{ps_134-5_080-v0010}{10}{typewriter deletion}{\xout{ly}}{}{The typewriter strikes out this material; it leaves the corrected running text.}
\psapparatusend
\fi
\pssource{ps_134-5_081}{134-5/81}{https://grothendieck.umontpellier.fr/web/viewer.html?file=../134-5.pdf\#page=82}
\psnumber{ps_134-5_081-v0001}{1}
\psnumber{ps_134-5_081-v0002}{2}
\psnumber{ps_134-5_081-v0003}{3}
\psnumber{ps_134-5_081-v0004}{4}
\psnumber{ps_134-5_081-v0005}{5}
\psnumber{ps_134-5_081-v0006}{6}
\psnumber{ps_134-5_081-v0007}{7}
\psnumber{ps_134-5_081-v0008}{8}
\psnumber{ps_134-5_081-v0009}{9}
\psnumber{ps_134-5_081-v0010}{10}
%% BEGIN SOURCE ps_134-5_081
% Pursuing Stacks (cotes 134-2 to 134-8): archive cote 134-5, site page 81 (PDF page 82), typescript page 268.
% typescript. Draft compared: ps4.tex:4141-4208.
% AI-transcribed by Claude Opus 5.5 (claude-opus-5-5, adaptive thinking, effort high), Message Batches API, batch ps2 (msgbatch_01NVwr1Yuo8qJf5rykeGBXPK), 2026-11-12; reviewed against both supplied scan halves; see CHECK_DECISIONS.tsv.
\scanpage{134-5/81}
\mspage{268}

want. By the description of J in terms of $J_o$ , we are \var{\inserted{now}}{\hand{now}}{insertion}{ps_134-5_081-v0001} reduced to proving
the following
\hand{\uline{\textcolor{black}{Lemma}}}\hand{1}: Let I be an object of $M_a = M_{\underline{W}}$ . Then any two sections of I (over
$e_M$ ) are $M_c$ - homotopic.

Let, as in the previous section, $\underline{W}_o \hand{\subset} \underline{W}$ be the largest of all basic
% The l in Fl((Cat)) carries a diagonal slash-like stroke.
localizers satisfying Loc 4), i.e.
\[ \underline{W}_o = \hand{\{} f \hand{\in} Fl((Cat)) \hand{|} \hand{\underline{\Pi}}_o(f) \text{ a bijection } \hand{\}} , \]
therefore we have
\[ M_{\underline{W}} \hand{\subset} M_{\underline{W}_o} \overset{def}{=} M_o , \]
and we are reduced to proving the lemma for $\underline{W}_o$ instead of $\underline{W}$ , i.e. for $M_o$
instead of $M_{\underline{W}}$ . We'll use the \var{\inserted[typed]{small}}{small}{typed insertion}{ps_134-5_081-v0002} full subcategory C of $M_c$ generating the
contractibility structure $M_c$ , we may assume that C is stable under finite
products. Hence $C\hat{}$ is totally $\underline{W}_{oC}$ -connected i.e. totally 0-connected.
Moreover, as $C \hand{\subset} M_c$ , \var{\xout{xxxf}}{}{typewriter deletion}{ps_134-5_081-v0003} and $e_M$ is in C, it follows that every object
of C has a section (over $e_M = e_C$ ) - which implies that every non empty
object of $C\hat{}$ has a section, i.e. $C\hat{}$ is ``strictly totally 0-connected''
(cf. p. 144 and 149). Note that (by prop.2 b) of p. \var{\over[typed]{\unsure{3}}{2}}{2}{typewriter overwriting}{ps_134-5_081-v0004}48) we have
\[ M_o = \hand{\{} x \hand{\in} Ob\, M \hand{|} i^{\hand{*}}(x) \text{ is 0-connected in } C\hat{} \hand{\}} , \]
where i is the \var{\xout{x}}{}{typewriter deletion}{ps_134-5_081-v0005}inclusion functor:
\[ i: C \hand{\longrightarrow} M . \]
We have to prove that any two section $\hand{\delta}_o,\hand{\delta}_1$ of an object I of $M_o$ are $M_c$-
homotopic, or what amounts to the same, C-homotopic. This translates
readily into the statement that $i^*(\hand{\delta}_o)$ and $i^*(\hand{\delta}_1)$ are C-homotopic in $C\hat{}$.
Thus we are reduced to the following lemma (in the case of the topos $C\hat{}$):
\hand{\uline{\textcolor{black}{Lemma}}} 2. Let $\hand{\mathcal{C}}$ \var{\inserted{be}}{\hand{be}}{insertion}{ps_134-5_081-v0006} a totally 0-connected topos such that any non-empty object
of $\hand{\mathcal{C}}$ has a section, and let C be a small full generating subcategory,
whose elements are 0-connected. Then for any 0-connected object I of $\hand{\mathcal{C}}$ ,
and any two sections $\hand{\delta}_o,\hand{\delta}_1$ of I , these are \var{\inserted[typed]{C-}}{C-}{typed insertion}{ps_134-5_081-v0007}homotopic, i.e. they can be
joined by a \var{\inserted[typed]{finite}}{finite}{typed insertion}{ps_134-5_081-v0008} chain of sections, any two consecutive among which can be ob-%
tained as the images of two sections \var{\inserted[typed]{$s_i$, $t_i$}}{$s_i$, $t_i$}{typed insertion}{ps_134-5_081-v0009} of an object $x_i$ of C , by means of a
% The h of h_i is densely overprinted; no distinct prior letter can be recovered.
map $h_i : x_i \hand{\longrightarrow} I$ .

This lemma is essentially a restatement (cleaned from extraneous \var{\xout{assum}}{}{typewriter deletion}{ps_134-5_081-v0010}
assumptions due to an awkward conceptual background) of the proposition of
page 149 (section 60), and the proof will be left to the reader.

\margin{81}

%% END SOURCE ps_134-5_081
\ifnum\EditionMode=1
\psapparatusstart
\psentry{ps_134-5_081-v0001}{1}{insertion}{\inserted{now}}{\hand{now}}{The handwritten insertion is incorporated at its indicated position; insertion brackets leave the running text.}
\psentry{ps_134-5_081-v0002}{2}{typed insertion}{\inserted[typed]{small}}{small}{The typed insertion is incorporated at its indicated position; insertion brackets leave the running text.}
\psentry{ps_134-5_081-v0003}{3}{typewriter deletion}{\xout{xxxf}}{}{The typewriter strikes out this material; it leaves the corrected running text.}
\psentry{ps_134-5_081-v0004}{4}{typewriter overwriting}{\over[typed]{\unsure{3}}{2}}{2}{The superposed typed reading replaces the underlying reading.}
\psentry{ps_134-5_081-v0005}{5}{typewriter deletion}{\xout{x}}{}{The typewriter strikes out this material; it leaves the corrected running text.}
\psentry{ps_134-5_081-v0006}{6}{insertion}{\inserted{be}}{\hand{be}}{The handwritten insertion is incorporated at its indicated position; insertion brackets leave the running text.}
\psentry{ps_134-5_081-v0007}{7}{typed insertion}{\inserted[typed]{C-}}{C-}{The typed insertion is incorporated at its indicated position; insertion brackets leave the running text.}
\psentry{ps_134-5_081-v0008}{8}{typed insertion}{\inserted[typed]{finite}}{finite}{The typed insertion is incorporated at its indicated position; insertion brackets leave the running text.}
\psentry{ps_134-5_081-v0009}{9}{typed insertion}{\inserted[typed]{$s_i$, $t_i$}}{$s_i$, $t_i$}{The typed insertion is incorporated at its indicated position; insertion brackets leave the running text.}
\psentry{ps_134-5_081-v0010}{10}{typewriter deletion}{\xout{assum}}{}{The typewriter strikes out this material; it leaves the corrected running text.}
\psapparatusend
\fi
\pssource{ps_134-5_082}{134-5/82}{https://grothendieck.umontpellier.fr/web/viewer.html?file=../134-5.pdf\#page=83}
\psnumber{ps_134-5_082-v0001}{1}
\psnumber{ps_134-5_082-v0002}{2}
\psnumber{ps_134-5_082-v0003}{3}
\psnumber{ps_134-5_082-v0004}{4}
\psnumber{ps_134-5_082-v0005}{5}
\psnumber{ps_134-5_082-v0006}{6}
\psnumber{ps_134-5_082-v0007}{7}
\psnumber{ps_134-5_082-v0008}{8}
\psnumber{ps_134-5_082-v0009}{9}
\psnumber{ps_134-5_082-v0010}{10}
\psnumber{ps_134-5_082-v0011}{11}
%% BEGIN SOURCE ps_134-5_082
% Pursuing Stacks (cotes 134-2 to 134-8): archive cote 134-5, site page 82 (PDF page 83), typescript page 269.
% typescript. Draft compared: ps4.tex:4190-4271.
% AI-transcribed by Claude Opus 5.5 (claude-opus-5-5, adaptive thinking, effort high), Message Batches API, batch ps2 (msgbatch_01NVwr1Yuo8qJf5rykeGBXPK), 2026-11-12; reviewed against both supplied scan halves; see CHECK_DECISIONS.tsv.
\scanpage{134-5/82}
\mspage{269}

\hand{\uline{\textcolor{black}{Application to relation between contractibility and objects}}} \hand{\uline{\textcolor{black}{Hom}}}(X,Y).

I would like to review here a few things along these lines, which were
somewhat scattered in the notes before (section 51 E) p.121 and section 57
p. 143 notably), and at times came out awkwardly because of unadequate
conceptual background. We assume M endowed with a contractibility structure
$\mathrm{M}_c$ , and use the notations of the previous section\unsure{;|,} especially concerning
% The l in Fl(M) has a diagonal slash-like overstrike.
the subsets of Ob M and of $F\ell$(M)
\[ (1) \qquad \mathrm{M}_c \hand{\subset} \mathrm{M}_{\mathrm{oo}} \hand{\subset} \mathrm{M}_{\underline{W}} \hand{\subset} \mathrm{M}_o \ , \quad \mathrm{W}_c \hand{\subset} \mathrm{W}_{\mathrm{oo}} \hand{\subset} \underline{\mathrm{W}}_{\mathrm{M}} \hand{\subset} \var{\xout{x}}{}{typewriter deletion}{ps_134-5_082-v0001}\, \mathrm{W}_o \quad . \]
Let X be an object of M, such that the object \var{\xout{Hom}}{}{typewriter deletion}{ps_134-5_082-v0002}
\[ (2) \qquad \mathrm{I} = \underline{\mathrm{Hom}}(\mathrm{X},\mathrm{X}) \]
exists in M (NB a priori it is an object in M\^{}, we assume it to be repre-%
sentable). We assume that X is endowed with a section c , hence a section
\var{\xout{of}}{}{typewriter deletion}{ps_134-5_082-v0003} $\hand{\delta}_1$ of I , corresponding to the constant endomorphism of X with value c,
We'll denote by $\hand{\delta}_o$ the section of I corresponding to the identity of X.
Thus
\[ (3) \qquad \underline{\mathrm{I}} = (\mathrm{I},\hand{\delta}_o,\hand{\delta}_1) \]
becomes an interval of M, and the composition law of I = \hand{\uline{\textcolor{black}{Hom}}}(X,X) turns
it into a \hand{\uline{\textcolor{black}{multiplicative}}} \hand{\uline{\textcolor{black}{interval}}} , admitting respectively $\hand{\pi}_o$ and $\hand{\pi}_1$ as
left unit and left zero element. Then the following conditions are equi-%
valent:

(\var{\xout{xx}}{}{typewriter deletion}{ps_134-5_082-v0004}i) X is contractible.

(ii) \hand{\uline{\textcolor{black}{Hom}}}(X,X) is contractible, \hand{i.e. $\mathrm{I}\in \mathrm{M}_c$ .}

(ii') \hand{\uline{\textcolor{black}{Hom}}}(X,X) is 0-connected, \hand{i.e. $\mathrm{I}\in \mathrm{M}_o$ .}

(ii") \var{\xout{Hom}}{}{typewriter deletion}{ps_134-5_082-v0005} The two sections $\hand{\delta}_o,\hand{\delta}_1$ of $\underline{\mathrm{I}}$=\hand{\uline{\textcolor{black}{Hom}}}(X,X) are homotopic.

(iii) For any object Y in M such that \hand{\uline{\textcolor{black}{Hom}}}(Y,X) exists in M, \hand{\uline{\textcolor{black}{Hom}}}(Y,X)
is contractible.

(iii') For any Y as in (iii), \hand{\uline{\textcolor{black}{Hom}}}(Y,X) is \var{\inserted{0-}}{\hand{0-}}{insertion}{ps_134-5_082-v0006}connected.

(iv) For any Y in M such that \hand{\uline{\textcolor{black}{Hom}}}(X,Y) exists in M, the canonical
map
\[ (4) \qquad \mathrm{Y} \hand{\longrightarrow} \underline{\mathrm{Hom}}(\mathrm{X},\mathrm{Y}) \]
is a homotopism (more accurately still, Y as a subobject of
\hand{\uline{\textcolor{black}{Hom}}}(X,Y) is a deformation ret\var{\xout{ra}}{}{typewriter deletion}{ps_134-5_082-v0007}act). Moreover, X is 0-connected.

(iv') For any Y in M as in (iv), the map (4) induces a bijection \var{\xout{xx}}{}{typewriter deletion}{ps_134-5_082-v0008}
\[ \var{\xout{x}}{}{typewriter deletion}{ps_134-5_082-v0009} \qquad \hand{\pi}_o(\mathrm{Y}) \hand{\to} \hand{\pi}_o(\underline{\mathrm{Hom}}(\mathrm{X},\mathrm{Y})) \quad , \]
moreover X is 0-connected.

NB Of co\var{\over[typed]{y}{u}}{u}{typewriter overwriting}{ps_134-5_082-v0010}rse, the 0-connectedness notion and the functor $\hand{\Pi}_o$ used \var{\xout{her}}{}{typewriter deletion}{ps_134-5_082-v0011}

\margin{82}

%% END SOURCE ps_134-5_082
\ifnum\EditionMode=1
\psapparatusstart
\psentry{ps_134-5_082-v0001}{1}{typewriter deletion}{$\xout{x}$}{}{The typewriter strikes out this material; it leaves the corrected running text.}
\psentry{ps_134-5_082-v0002}{2}{typewriter deletion}{\xout{Hom}}{}{The typewriter strikes out this material; it leaves the corrected running text.}
\psentry{ps_134-5_082-v0003}{3}{typewriter deletion}{\xout{of}}{}{The typewriter strikes out this material; it leaves the corrected running text.}
\psentry{ps_134-5_082-v0004}{4}{typewriter deletion}{\xout{xx}}{}{The typewriter strikes out this material; it leaves the corrected running text.}
\psentry{ps_134-5_082-v0005}{5}{typewriter deletion}{\xout{Hom}}{}{The typewriter strikes out this material; it leaves the corrected running text.}
\psentry{ps_134-5_082-v0006}{6}{insertion}{\inserted{0-}}{\hand{0-}}{The handwritten insertion is incorporated at its indicated position; insertion brackets leave the running text.}
\psentry{ps_134-5_082-v0007}{7}{typewriter deletion}{\xout{ra}}{}{The typewriter strikes out this material; it leaves the corrected running text.}
\psentry{ps_134-5_082-v0008}{8}{typewriter deletion}{\xout{xx}}{}{The typewriter strikes out this material; it leaves the corrected running text.}
\psentry{ps_134-5_082-v0009}{9}{typewriter deletion}{$\xout{x}$}{}{The typewriter strikes out this material; it leaves the corrected running text.}
\psentry{ps_134-5_082-v0010}{10}{typewriter overwriting}{\over[typed]{y}{u}}{u}{The superposed typed reading replaces the underlying reading.}
\psentry{ps_134-5_082-v0011}{11}{typewriter deletion}{\xout{her}}{}{The typewriter strikes out this material; it leaves the corrected running text.}
\psapparatusend
\fi
\pssource{ps_134-5_083}{134-5/83}{https://grothendieck.umontpellier.fr/web/viewer.html?file=../134-5.pdf\#page=84}
\psnumber{ps_134-5_083-v0001}{1}
\psnumber{ps_134-5_083-v0002}{2}
\psnumber{ps_134-5_083-v0003}{3}
\psnumber{ps_134-5_083-v0004}{4}
\psnumber{ps_134-5_083-v0005}{5}
\psnumber{ps_134-5_083-v0006}{6}
\psnumber{ps_134-5_083-v0007}{7}
\psnumber{ps_134-5_083-v0008}{8}
\psnumber{ps_134-5_083-v0009}{9}
\psnumber{ps_134-5_083-v0010}{10}
\psnumber{ps_134-5_083-v0011}{11}
%% BEGIN SOURCE ps_134-5_083
% Pursuing Stacks (cotes 134-2 to 134-8): archive cote 134-5, site page 83 (PDF page 84), typescript page 270.
% typescript. Draft compared: ps4.tex:4257-4328.
% AI-transcribed by Claude Opus 5.5 (claude-opus-5-5, adaptive thinking, effort high), Message Batches API, batch ps2 (msgbatch_01NVwr1Yuo8qJf5rykeGBXPK), 2026-11-12; reviewed against both supplied scan halves; see CHECK_DECISIONS.tsv.
\scanpage{134-5/83}
\mspage{270}

here are those associated to the given structure $M_c$ . The case dealt with
in section 57 (p. 143) is \var{\xout{the one}}{}{typewriter deletion}{ps_134-5_083-v0001}\var{\inserted{essentially}}{\hand{essentially}}{insertion}{ps_134-5_083-v0002} (it seems) \var{\deleted{essentially}}{}{deletion}{ps_134-5_083-v0003} the one when
these notions are the ones canonically defined in terms of the category
structure of M alone. In view of the inclusions (1), we could throw in
a handfull more \var{\inserted[typed]{obviously}}{obviously}{typed insertion}{ps_134-5_083-v0004} equivalent conditions, using $M_{\mathrm{oo}}$ , $M_{\underline{W}}$ or $W_{\mathrm{oo}}$ , $\underline{W}_M$ \var{\inserted{instead of $M_c$ etc}}{\hand{instead of $M_c$ etc}}{insertion}{ps_134-5_083-v0005}, but
it seems this would confuse the picture rather than complete it .

Next thing is to look at the canonical map
\[ (5) \qquad \hand{\overline{\Gamma}}(X) \hand{\to} \hand{\pi}_o(X) \]
defined for any object X of M , where $\hand{\Gamma}(X)$ denotes the set of sections
of X, and $\hand{\overline{\Gamma}}(X)$ denotes the set of co\var{\over[typed]{t}{r}}{r}{typewriter overwriting}{ps_134-5_083-v0006}responding homotopy classes of sec-%
tions. (This map was considered in a \var{\xout{special case}}{}{typewriter deletion}{ps_134-5_083-v0007} slightly different case
on page 144.) The map (5) can be viewed as the particular case of
\[ \hand{\overline{\textcolor{black}{\mathrm{Hom}}}}(Y,X) \hand{\to} \mathrm{Hom}_{\mathrm{Sets}}(\hand{\pi}_o(Y),\hand{\pi}_o(X)) , \]
obtained when $Y = e_M$ hence $\hand{\pi}_o(Y) =$ one-point set. By the standard des-%
cription
\[ (6) \qquad \hand{\pi}_o(X)\hand{\simeq} \hand{\underline{\pi}}_o(C_{/X}) \quad ( = \hand{\pi}_o(i^*(X))) \]
we immediately get that (5) is always surjective, as any element in $\hand{\pi}_o(X)$
is induced by an element $x \hand{\to} X$ of $C_{/X}$ , hence by a section
\[ e_M \hand{\to} x \hand{\to} X \quad , \]
where $e_M \hand{\to} x$ is a section of the (cont\var{\over[typed]{t}{r}}{r}{typewriter overwriting}{ps_134-5_083-v0008}actible) object x in C. But (5)
is in fact \hand{\uline{\textcolor{black}{bijective}}} . To see this, let's note that the map (5) is
isomorphic to the corresponding map \var{\inserted{in (Sets)}}{\hand{in (Sets)}}{insertion}{ps_134-5_083-v0009}, with M replaced by C\^{} and X by $i^*(X)$,
C remaining the same, as follows from \var{\xout{the}}{}{typewriter deletion}{ps_134-5_083-v0010} (6) and the formula
\[ \hand{\Gamma}(X)\hand{\simeq} \hand{\overline{\Gamma}}(i^*(X)) \quad . \]
The latter is a particular case of the

\hand{\uline{\textcolor{black}{Lemma}}} 3. Let Y,X be two objects of M, with Y in C. Then the natural map
\[ (7) \qquad \hand{\overline{\textcolor{black}{\mathrm{Hom}}}}(\var{\over{X}{Y}}{\hand{Y}}{writing over}{ps_134-5_083-v0011},X) \hand{\to} \hand{\overline{\textcolor{black}{\mathrm{Hom}}}}(i^*(Y),i^*(X)) \]
between sets of homotopy classes of maps (in M and in C\^{}respectively, the
latter being endowed with the contractibility structure generated by C)
is bijective.

The verification is tautological, due to the fact that for any
object I of C (which is going to play the role of a homotopy interval for
deformations), we got
\[ \mathrm{Hom}_M(\mathrm{IxY},X) \hand{\xrightarrow{\sim}} \mathrm{Hom}_{C\hat{}}(\mathrm{IxY},i^*(X)). \]

\margin{83}

%% END SOURCE ps_134-5_083
\ifnum\EditionMode=1
\psapparatusstart
\psentry{ps_134-5_083-v0001}{1}{typewriter deletion}{\xout{the one}}{}{The typewriter strikes out this material; it leaves the corrected running text.}
\psentry{ps_134-5_083-v0002}{2}{insertion}{\inserted{essentially}}{\hand{essentially}}{The handwritten insertion is incorporated at its indicated position; insertion brackets leave the running text.}
\psentry{ps_134-5_083-v0003}{3}{deletion}{\deleted{essentially}}{}{The author cancels this material by hand; it leaves the corrected running text.}
\psentry{ps_134-5_083-v0004}{4}{typed insertion}{\inserted[typed]{obviously}}{obviously}{The typed insertion is incorporated at its indicated position; insertion brackets leave the running text.}
\psentry{ps_134-5_083-v0005}{5}{insertion}{\inserted{instead of $M_c$ etc}}{\hand{instead of $M_c$ etc}}{The handwritten insertion is incorporated at its indicated position; insertion brackets leave the running text.}
\psentry{ps_134-5_083-v0006}{6}{typewriter overwriting}{\over[typed]{t}{r}}{r}{The superposed typed reading replaces the underlying reading.}
\psentry{ps_134-5_083-v0007}{7}{typewriter deletion}{\xout{special case}}{}{The typewriter strikes out this material; it leaves the corrected running text.}
\psentry{ps_134-5_083-v0008}{8}{typewriter overwriting}{\over[typed]{t}{r}}{r}{The superposed typed reading replaces the underlying reading.}
\psentry{ps_134-5_083-v0009}{9}{insertion}{\inserted{in (Sets)}}{\hand{in (Sets)}}{The handwritten insertion is incorporated at its indicated position; insertion brackets leave the running text.}
\psentry{ps_134-5_083-v0010}{10}{typewriter deletion}{\xout{the}}{}{The typewriter strikes out this material; it leaves the corrected running text.}
\psentry{ps_134-5_083-v0011}{11}{writing over}{$\over{X}{Y}$}{$\hand{Y}$}{The superposed handwritten reading replaces the underlying reading.}
\psapparatusend
\fi
\pssource{ps_134-5_084}{134-5/84}{https://grothendieck.umontpellier.fr/web/viewer.html?file=../134-5.pdf\#page=85}
\psnumber{ps_134-5_084-v0001}{1}
\psnumber{ps_134-5_084-v0002}{2}
\psnumber{ps_134-5_084-v0003}{3}
\psnumber{ps_134-5_084-v0004}{4}
\psnumber{ps_134-5_084-v0005}{5}
\psnumber{ps_134-5_084-v0006}{6}
\psnumber{ps_134-5_084-v0007}{7}
\psnumber{ps_134-5_084-v0008}{8}
\psnumber{ps_134-5_084-v0009}{9}
\psnumber{ps_134-5_084-v0010}{10}
\psnumber{ps_134-5_084-v0011}{11}
\psnumber{ps_134-5_084-v0012}{12}
\psnumber{ps_134-5_084-v0013}{13}
\psnumber{ps_134-5_084-v0014}{14}
\psnumber{ps_134-5_084-v0015}{15}
\psnumber{ps_134-5_084-v0016}{16}
\psnumber{ps_134-5_084-v0017}{17}
\psnumber{ps_134-5_084-v0018}{18}
\psnumber{ps_134-5_084-v0019}{19}
%% BEGIN SOURCE ps_134-5_084
% Pursuing Stacks (cotes 134-2 to 134-8): archive cote 134-5, site page 84 (PDF page 85), typescript page 271.
% typescript. Draft compared: ps4.tex:4310-4388.
% AI-transcribed by Claude Opus 5.5 (claude-opus-5-5, adaptive thinking, effort high), Message Batches API, batch ps2 (msgbatch_01NVwr1Yuo8qJf5rykeGBXPK), 2026-11-12; reviewed against both supplied scan halves; see CHECK_DECISIONS.tsv.
\scanpage{134-5/84}
\mspage{271}

\margin{(here} N\var{\over[typed]{O}{o}}{o}{typewriter overwriting}{ps_134-5_084-v0001}w, \var{\xout{the case of}}{}{typewriter deletion}{ps_134-5_084-v0002} the map (5) in case of a strictly \var{\inserted[typed]{totally}}{totally}{typed insertion}{ps_134-5_084-v0003} 0-connected category
\hand{(}$C^{\wedge}$\hand{)} has been dealt with in section 58 (p.114), it follows easily that the
map is bijective - hence the

\hand{\uline{\textcolor{black}{Proposition}}} 1 . Let $(M,M_c)$ be any contractibility structure, \var{\xout{Xxxx$_o$}}{}{typewriter deletion}{ps_134-5_084-v0004}
\[ \hand{\pi}_o : M \hand{\longrightarrow} \text{(Sets)} \qquad , \hand{\pi}_o(X) \hand{\simeq} \underline{\hand{\pi}}_o(\var{\over{\illegible{1}}{M_c}}{\hand{M_c}}{writing over}{ps_134-5_084-v0005}/X) \] \margin{\tikz[baseline=(m.base)]{\node[draw,circle,inner sep=2pt] (m) {$M_c/X$};}}
the corresponding ``connected components functor'', $X$ any object of $M$. Then
the canonical map (5) above is bijective.

\hand{\uline{\textcolor{black}{Corollary}}} Let $X,Y$ be two objects of $M$, such that $\underline{\mathrm{Hom}}(X,Y)$ exists in $M$.
Then the natural map
\[ (8) \qquad \hand{\overline{\textcolor{black}{\mathrm{Hom}}}}(X,Y) \hand{\longrightarrow} \hand{\pi}_o(\underline{\mathrm{Hom}}(X,Y)) \]
is bijective. \margin{$\lfloor$\rule{2.5cm}{0.4pt}}

\margin{$\ulcorner$\rule{2.5cm}{0.4pt}} To finish these generalities on $\underline{\mathrm{Hom}}$'s , I would like to generalize
to the context of contractibility structures the results of prop.4 p.228
(section 74) on products of aspheric functors. We start within a context
of asphericity structures:
\hand{\uline{\textcolor{black}{Proposition}}} 2. Let $(M,M_a)$ be a $\underline{W}$-asphericity structure (for a basic locali-%
zer $\underline{W}$), \unsure{$C$|$B$} a full subcategory of $M$ generating this structure i.e. such that
\unsure{$C$|$B$} \hand{$\to$} $M$ is $M_a$-$\underline{W}$-aspheric, $A$ a small category,
\[ i: A \hand{\longrightarrow} M \]
a \var{\xout{functor}}{}{typewriter deletion}{ps_134-5_084-v0006} $M_a$-$\underline{W}$-aspheric functor. We assume $M$ stable under binary products.

a) Let $b_o$ in $M$ be such that for any $y$ in $C$, $\underline{\mathrm{Hom}}_M(b_o,y)$ exists in $M$,
and for any $x$ in $M_a$ , \var{\over{\unsure{x}}{x}}{\hand{x}}{writing over}{ps_134-5_084-v0007} x $b_o$ is in $M_a$ . Let $i_{b_o} : A \hand{\longrightarrow} M$ be the constant
functor with value $b_o$ , and consider the product functor
\[ ixi_{b_o} : a \hand{\mapsto} i(a)xb_o : A \hand{\longrightarrow} M \ . \]
This functor is $M_a$-$\underline{W}$-aspheric iff for any $y$ in $C$, $\underline{\mathrm{Hom}}(b_o,y)$ is in $M_a$
(a condition which does not depend on $i$ nor even on $A$).

b) Let $\var{\over{\unsure{C}}{B}}{\hand{B}}{writing over}{ps_134-5_084-v0008}$ be a full subcategory of $M$, such that for any $b_o$ in $\var{\over{\unsure{C}}{B}}{\hand{B}}{writing over}{ps_134-5_084-v0009}$ \var{\inserted{x in $M_a$}}{\hand{x in $M_a$}}{insertion}{ps_134-5_084-v0010} and
and $y$ in $C$, \var{\inserted{we have $x \times b_o \in M_a$ and}}{\hand{we have $x \times b_o \in M_a$ and}}{insertion}{ps_134-5_084-v0011} $\underline{\mathrm{Hom}}(b_o,y)$ \var{\xout{is}}{}{typewriter deletion}{ps_134-5_084-v0012} exists in $M$, and is in $M_a$. \var{\inserted{Assume $A$ totally $\underline{W}$-aspheric}}{\hand{Assume $A$ totally $\underline{W}$-aspheric}}{insertion}{ps_134-5_084-v0013} Let $i'$ : $A \hand{\to} M$
be \var{\over{a}{any}}{\hand{any}}{writing over}{ps_134-5_084-v0014} functor factoring through $\var{\over{\unsure{C}}{B}}{\hand{B}}{writing over}{ps_134-5_084-v0015}$ , \var{\over{T}{t}}{\hand{t}}{writing over}{ps_134-5_084-v0016}hen the product functor \var{\xout{ixix}}{}{typewriter deletion}{ps_134-5_084-v0017}
\[ ixi' : \hand{\unsure{a}} \hand{\longmapsto} i(a)xi'(a) \ : \ A \hand{\longrightarrow} M \]
is $M_a$-$\underline{W}$-aspheric\hand{.} \var{\deleted{iff $A$ is \xout{asphe} $\underline{W}$-aspheric (a condition which does not
depend upon $i'$\over{)}{,} \hand{nor even upon $A$ if $M$ has a final object $e_M$)}}}{}{deletion}{ps_134-5_084-v0018}

The proof is word by word the same as for the analoguous statements
p. 228-229, and therefore left to the reader. (The analoguous statement
to cor.1 \var{\xout{propx}}{}{typewriter deletion}{ps_134-5_084-v0019} p. 229 is equally valid.)

\margin{84}

%% END SOURCE ps_134-5_084
\ifnum\EditionMode=1
\psapparatusstart
\psentry{ps_134-5_084-v0001}{1}{typewriter overwriting}{\over[typed]{O}{o}}{o}{The superposed typed reading replaces the underlying reading.}
\psentry{ps_134-5_084-v0002}{2}{typewriter deletion}{\xout{the case of}}{}{The typewriter strikes out this material; it leaves the corrected running text.}
\psentry{ps_134-5_084-v0003}{3}{typed insertion}{\inserted[typed]{totally}}{totally}{The typed insertion is incorporated at its indicated position; insertion brackets leave the running text.}
\psentry{ps_134-5_084-v0004}{4}{typewriter deletion}{\xout{Xxxx$_o$}}{}{The typewriter strikes out this material; it leaves the corrected running text.}
\psentry{ps_134-5_084-v0005}{5}{writing over}{$\over{\illegible{1}}{M_c}$}{$\hand{M_c}$}{The superposed handwritten reading replaces the underlying reading.}
\psentry{ps_134-5_084-v0006}{6}{typewriter deletion}{\xout{functor}}{}{The typewriter strikes out this material; it leaves the corrected running text.}
\psentry{ps_134-5_084-v0007}{7}{writing over}{\over{\unsure{x}}{x}}{\hand{x}}{The superposed handwritten reading replaces the underlying reading.}
\psentry{ps_134-5_084-v0008}{8}{writing over}{$\over{\unsure{C}}{B}$}{$\hand{B}$}{The superposed handwritten reading replaces the underlying reading.}
\psentry{ps_134-5_084-v0009}{9}{writing over}{$\over{\unsure{C}}{B}$}{$\hand{B}$}{The superposed handwritten reading replaces the underlying reading.}
\psentry{ps_134-5_084-v0010}{10}{insertion}{\inserted{x in $M_a$}}{\hand{x in $M_a$}}{The handwritten insertion is incorporated at its indicated position; insertion brackets leave the running text.}
\psentry{ps_134-5_084-v0011}{11}{insertion}{\inserted{we have $x \times b_o \in M_a$ and}}{\hand{we have $x \times b_o \in M_a$ and}}{The handwritten insertion is incorporated at its indicated position; insertion brackets leave the running text.}
\psentry{ps_134-5_084-v0012}{12}{typewriter deletion}{\xout{is}}{}{The typewriter strikes out this material; it leaves the corrected running text.}
\psentry{ps_134-5_084-v0013}{13}{insertion}{\inserted{Assume $A$ totally $\underline{W}$-aspheric}}{\hand{Assume $A$ totally $\underline{W}$-aspheric}}{The handwritten insertion is incorporated at its indicated position; insertion brackets leave the running text.}
\psentry{ps_134-5_084-v0014}{14}{writing over}{\over{a}{any}}{\hand{any}}{The superposed handwritten reading replaces the underlying reading.}
\psentry{ps_134-5_084-v0015}{15}{writing over}{$\over{\unsure{C}}{B}$}{$\hand{B}$}{The superposed handwritten reading replaces the underlying reading.}
\psentry{ps_134-5_084-v0016}{16}{writing over}{\over{T}{t}}{\hand{t}}{The superposed handwritten reading replaces the underlying reading.}
\psentry{ps_134-5_084-v0017}{17}{typewriter deletion}{\xout{ixix}}{}{The typewriter strikes out this material; it leaves the corrected running text.}
\psentry{ps_134-5_084-v0018}{18}{deletion}{\deleted{iff $A$ is \xout{asphe} $\underline{W}$-aspheric (a condition which does not
depend upon $i'$\over{)}{,} \hand{nor even upon $A$ if $M$ has a final object $e_M$)}}}{}{The author cancels this material by hand; it leaves the corrected running text.}
\psentry{ps_134-5_084-v0019}{19}{typewriter deletion}{\xout{propx}}{}{The typewriter strikes out this material; it leaves the corrected running text.}
\psapparatusend
\fi
\pssource{ps_134-5_085}{134-5/85}{https://grothendieck.umontpellier.fr/web/viewer.html?file=../134-5.pdf\#page=86}
\psnumber{ps_134-5_085-v0001}{1}
\psnumber{ps_134-5_085-v0002}{2}
\psnumber{ps_134-5_085-v0003}{3}
\psnumber{ps_134-5_085-v0004}{4}
\psnumber{ps_134-5_085-v0005}{5}
\psnumber{ps_134-5_085-v0006}{6}
\psnumber{ps_134-5_085-v0007}{7}
\psnumber{ps_134-5_085-v0008}{8}
\psnumber{ps_134-5_085-v0009}{9}
\psnumber{ps_134-5_085-v0010}{10}
\psnumber{ps_134-5_085-v0011}{11}
%% BEGIN SOURCE ps_134-5_085
% Pursuing Stacks (cotes 134-2 to 134-8): archive cote 134-5, site page 85 (PDF page 86), typescript page 272.
% typescript. Draft compared: ps4.tex:4371-4437.
% AI-transcribed by Claude Opus 5.5 (claude-opus-5-5, adaptive thinking, effort high), Message Batches API, batch ps2 (msgbatch_01NVwr1Yuo8qJf5rykeGBXPK), 2026-11-12; reviewed against both supplied scan halves; see CHECK_DECISIONS.tsv.
\scanpage{134-5/85}
\mspage{272}

The case I've in mind is when $M_a$ is generated by a contractibility
structure $M_c$ , and $C=B = M_c$ . The condition in b) boils down to exis-%
tence of $\underline{Hom}(b_o,y)$ in $M$, when $b_o$ and $y$ are both contractible objects of $M$
- as a matter of fact, in all cases I'm having in mind, \var{\xout{Hom(}}{}{typewriter deletion}{ps_134-5_085-v0001} the $\underline{Hom}$
exists even wi thout any contractibility assumption. Thus we get:

\hand{\uline{\textcolor{black}{Corollary}}} Let $(M,M_c)$ a contractibility structure, such that for any
two contractible objects $x,y$ , $\underline{Hom}_M(x,y)$ exists in $M$\var{\xout{$_{\illegible{1}}$}}{}{typewriter deletion}{ps_134-5_085-v0002} (a condition which
presumably is even superfluous ...). Consider the $\underline{W}$-asphericity structure
$M_{\underline{W}}$ on $M$ generated by $M_c$ . Then the product of a $M_{\underline{W}}$-$\underline{W}$-aspheric functor
$i: A \hand{\longrightarrow} \var{\over[typed]{\unsure{B}}{M}}{M}{typewriter overwriting}{ps_134-5_085-v0003}$ with a functor $i' : A \hand{\longrightarrow} M$ factoring through $M_c$ is again
$M_{\underline{W}}$-$\underline{W}$-aspheric. In particular, the category \var{\xout{$\underline{W}$-Asph(A,(M,M$_{\underline{W}}$)}}{}{typewriter deletion}{ps_134-5_085-v0004} of all asphe-%
ric functors from $A$ to $M$ which factor through $M_c$ is stable under binary pro%
-ducts.

From this one can deduce as on p.230 that if the latter category is
non empty, i.e. if there does exist an aspheric functor $A\hand{\to} M$ factoring
through $M_c$ , then one can define a canonical functor
\[ (9) \qquad Hot_{M,\underline{W}} \hand{\longrightarrow} Hot_{A,\underline{W}} \quad , \]
defined up to unique isomorphism, via a transitive system of isomorphisms
between the functors deduced from aspheric functors $A \hand{\to} M$ factoring
through $M_c$ .

83. After \var{\xout{s}}{}{typewriter deletion}{ps_134-5_085-v0005}stating and proving theorem 1 of previous section, I forgot
to give an answer to the most natural question arising from it - namely
\var{\inserted[typed]{how to}}{how to}{typed insertion}{ps_134-5_085-v0006} characterize those $\underline{W}$-asphericity structures on a category $M$ stable under
finite products, which can be generated by a contractibility structure.
The reason for this is surely that I have nothing better to offer than a
tautology: let
\[ M_a \hand{\subset} \mathrm{Ob}\, M \]
be the given $\underline{W}$-asphericity structure, \var{\xout{consider the corresponding}}{}{typewriter deletion}{ps_134-5_085-v0007} \var{\inserted[typed]{we assume beforehand that}}{we assume beforehand that}{typed insertion}{ps_134-5_085-v0008} this
structure is totally aspheric (which is a necessary condition for $M_a$ to
come from a structure $M_c$). \unsure{Tf|If} $W_a$ is the corresponding set of $\underline{W}$-equivalen-%
ces, it follows that the homotopy structure $h_{W_a}$ is also the one defined
by the homotopy interval structure \var{\inserted[typed]{$J$}}{$J$}{typed insertion}{ps_134-5_085-v0009} generated by the set \var{\inserted[typed]{$J_o$}}{$J_o$}{typed insertion}{ps_134-5_085-v0010} of intervals of $M$
made up with objects of $M_a$ . Let\var{\xout{$M_c$ be the corresponding}}{}{typewriter deletion}{ps_134-5_085-v0011}

\margin{85}

%% END SOURCE ps_134-5_085
\ifnum\EditionMode=1
\psapparatusstart
\psentry{ps_134-5_085-v0001}{1}{typewriter deletion}{\xout{Hom(}}{}{The typewriter strikes out this material; it leaves the corrected running text.}
\psentry{ps_134-5_085-v0002}{2}{typewriter deletion}{\xout{$_{\illegible{1}}$}}{}{The typewriter strikes out this material; it leaves the corrected running text.}
\psentry{ps_134-5_085-v0003}{3}{typewriter overwriting}{$\over[typed]{\unsure{B}}{M}$}{$M$}{The superposed typed reading replaces the underlying reading.}
\psentry{ps_134-5_085-v0004}{4}{typewriter deletion}{\xout{$\underline{W}$-Asph(A,(M,M$_{\underline{W}}$)}}{}{The typewriter strikes out this material; it leaves the corrected running text.}
\psentry{ps_134-5_085-v0005}{5}{typewriter deletion}{\xout{s}}{}{The typewriter strikes out this material; it leaves the corrected running text.}
\psentry{ps_134-5_085-v0006}{6}{typed insertion}{\inserted[typed]{how to}}{how to}{The typed insertion is incorporated at its indicated position; insertion brackets leave the running text.}
\psentry{ps_134-5_085-v0007}{7}{typewriter deletion}{\xout{consider the corresponding}}{}{The typewriter strikes out this material; it leaves the corrected running text.}
\psentry{ps_134-5_085-v0008}{8}{typed insertion}{\inserted[typed]{we assume beforehand that}}{we assume beforehand that}{The typed insertion is incorporated at its indicated position; insertion brackets leave the running text.}
\psentry{ps_134-5_085-v0009}{9}{typed insertion}{\inserted[typed]{$J$}}{$J$}{The typed insertion is incorporated at its indicated position; insertion brackets leave the running text.}
\psentry{ps_134-5_085-v0010}{10}{typed insertion}{\inserted[typed]{$J_o$}}{$J_o$}{The typed insertion is incorporated at its indicated position; insertion brackets leave the running text.}
\psentry{ps_134-5_085-v0011}{11}{typewriter deletion}{\xout{$M_c$ be the corresponding}}{}{The typewriter strikes out this material; it leaves the corrected running text.}
\psapparatusend
\fi
\pssource{ps_134-5_086}{134-5/86}{https://grothendieck.umontpellier.fr/web/viewer.html?file=../134-5.pdf\#page=87}
\psnumber{ps_134-5_086-v0001}{1}
\psnumber{ps_134-5_086-v0002}{2}
\psnumber{ps_134-5_086-v0003}{3}
\psnumber{ps_134-5_086-v0004}{4}
\psnumber{ps_134-5_086-v0005}{5}
\psnumber{ps_134-5_086-v0006}{6}
%% BEGIN SOURCE ps_134-5_086
% Pursuing Stacks (cotes 134-2 to 134-8): archive cote 134-5, site page 86 (PDF page 87), typescript page 273.
% typescript. Draft compared: ps4.tex:4421-4476.
% AI-transcribed by Claude Opus 5.5 (claude-opus-5-5, adaptive thinking, effort high), Message Batches API, batch ps2 (msgbatch_01NVwr1Yuo8qJf5rykeGBXPK), 2026-11-12; reviewed against both supplied scan halves; see CHECK_DECISIONS.tsv.
\scanpage{134-5/86}
\mspage{273}

\[ M_c \hand{\subset} \text{Ob } M \]
be the corresponding set of contractible objects (which may not be a con-%
tractibility structure on M). The conditions now are the following :

a) $M_c$ generate\var{\over[typed]{\unsure{s}}{d}}{d}{typewriter overwriting}{ps_134-5_086-v0001} the homotopy interval structure J, or equivalently,
any two sections of an object of $M_a$ are $M_c$-homotopic.

This condition is clearly equivalent to saying that J is indeed gene-%
rated by a contractibility structure, and the latter is necessarily $M_c$ .

b) $M_c$ generates the $\underline{W}$-asphericity structure $M_a$ , i.e. there exists
a small full subcategory B of $M_c$ such that the inclusion functor \var{\inserted[typed]{i:}}{i:}{typed insertion}{ps_134-5_086-v0002}$B \hand{\longrightarrow} M$\var{\xout{$_{\unsure{a}}$}}{}{typewriter deletion}{ps_134-5_086-v0003}
be \var{\xout{$\underline{W}$}}{}{typewriter deletion}{ps_134-5_086-v0004} $M_a$-$\underline{W}$-aspheric, i.e. for any x in $M_a$ , $i^*(x)$ in $B\hat{}$ is $\underline{W}$-aspheric \hand{(}i.e
$B_{/x}$ is $\underline{W}$-aspheric in (Cat)\hand{)}.

One would like too a ns condition on a functor
\[ \hand{\phi} : M \hand{\longrightarrow} \text{Hot}_{\underline{W}} \quad , \]
for this functor to be \var{\xout{associated}}{}{typewriter deletion}{ps_134-5_086-v0005} isomorphic to the one associated to a
contractibility structure on M - with special interest in the case when
$\underline{W} = \underline{W}_o$ , i.e. when the functor reduces to \hand{a functor}
\[ \hand{\pi}_o : M \hand{\to} (\text{Sets}) \quad . \]
Here again, I have nothing to offer except a tautological statement, whic\unsure{h}
is'nt worth the trouble writing it down. Nor do I have at present a com-%
pelling feeling that there should exist such a c\var{\over{\illegible{1}}{h}}{\hand{h}}{writing over}{ps_134-5_086-v0006}aracterization, under sui-%
table exactness assumptions on M say, and possibly assuming too that M is
stable under the operation $\underline{\text{Hom}}$ . Here also arises the question \hand{\uline{\textcolor{black}{whether a}}}
\hand{\uline{\textcolor{black}{functor}}} $\hand{\phi} : M \hand{\longrightarrow} \text{Hot}_{\underline{W}}$ \hand{\uline{\textcolor{black}{stemming from a contractibility structure}}} $M_c$ \hand{\uline{\textcolor{black}{on}}} M% The hook at the right of the pen underline is continuous with its terminal stroke, not a separately placed comma.
\hand{\uline{\textcolor{black}{can have non trivial automorphisms}}} - a question closely connected to the
``inspiring assumption'' of section 28.

\margin{86}

%% END SOURCE ps_134-5_086
\ifnum\EditionMode=1
\psapparatusstart
\psentry{ps_134-5_086-v0001}{1}{typewriter overwriting}{\over[typed]{\unsure{s}}{d}}{d}{The superposed typed reading replaces the underlying reading.}
\psentry{ps_134-5_086-v0002}{2}{typed insertion}{\inserted[typed]{i:}}{i:}{The typed insertion is incorporated at its indicated position; insertion brackets leave the running text.}
\psentry{ps_134-5_086-v0003}{3}{typewriter deletion}{\xout{$_{\unsure{a}}$}}{}{The typewriter strikes out this material; it leaves the corrected running text.}
\psentry{ps_134-5_086-v0004}{4}{typewriter deletion}{\xout{$\underline{W}$}}{}{The typewriter strikes out this material; it leaves the corrected running text.}
\psentry{ps_134-5_086-v0005}{5}{typewriter deletion}{\xout{associated}}{}{The typewriter strikes out this material; it leaves the corrected running text.}
\psentry{ps_134-5_086-v0006}{6}{writing over}{\over{\illegible{1}}{h}}{\hand{h}}{The superposed handwritten reading replaces the underlying reading.}
\psapparatusend
\fi
\pssource{ps_134-5_087}{134-5/87}{https://grothendieck.umontpellier.fr/web/viewer.html?file=../134-5.pdf\#page=88}
\psnumber{ps_134-5_087-v0001}{1}
\psnumber{ps_134-5_087-v0002}{2}
\psnumber{ps_134-5_087-v0003}{3}
\psnumber{ps_134-5_087-v0004}{4}
\psnumber{ps_134-5_087-v0005}{5}
\psnumber{ps_134-5_087-v0006}{6}
\psnumber{ps_134-5_087-v0007}{7}
\psnumber{ps_134-5_087-v0008}{8}
\psnumber{ps_134-5_087-v0009}{9}
\psnumber{ps_134-5_087-v0010}{10}
%% BEGIN SOURCE ps_134-5_087
% Pursuing Stacks (cotes 134-2 to 134-8): archive cote 134-5, site page 87 (PDF page 88), typescript page 274.
% typescript. Draft compared: ps4.tex:4458-4524.
% AI-transcribed by Claude Opus 5.5 (claude-opus-5-5, adaptive thinking, effort high), Message Batches API, batch ps2 (msgbatch_01NVwr1Yuo8qJf5rykeGBXPK), 2026-11-12; reviewed against both supplied scan halves; see CHECK_DECISIONS.tsv.
\scanpage{134-5/87}
\mspage{274}

29.6.

84. Two days have passed without writing an\hand{y} \var{\deleted{the}}{}{deletion}{ps_134-5_087-v0001} notes. Much of the
time I spent on writing mathematical letters - one pretty long one to
Gerd Falting\hand{s}, who (on my request) had sent me preprints of his \var{\xout{late}}{}{typewriter deletion}{ps_134-5_087-v0002}
recent work, \var{\inserted{notably}}{\hand{notably}}{insertion}{ps_134-5_087-v0003} on the Tate conjectures for abelian varieties and on the
Mordell conjecture, and had expressed interest hearing about some ideas
and conjectures on ``anabelian algebraic geometry''. I had been impressed,
\var{\over[typed]{\illegible{2}}{fr}}{fr}{typewriter overwriting}{ps_134-5_087-v0004}om a glance upon the last \var{\inserted{of his}}{\hand{of his}}{insertion}{ps_134-5_087-v0005} manuscript\hand{s}, to see three \var{\inserted[typed]{key}}{key}{typed insertion}{ps_134-5_087-v0006} conjectures \var{\xout{which}}{}{typewriter deletion}{ps_134-5_087-v0007}
proved in about forty pages, while they were being considered as quite
out of reach by the people supposed to know. Some ``anabelian'' conjectures
of mine are closely related to the Tate and Mordell conjectures just pro-%
ved by Falting - Deligne had pointed out to me about two years ago that
a certain fixed-point conjecture (which I like to view as the basic conjec-%
ture at present in the anabelian program) implied Mordell's, so why loose
one's time on it ! I have the feeling Falting is the kind of chap who
may become interested in things which are supposed to be too far off to
be worth looking at, that's why I had written him a few words, under
the moment's inspiration. - Another letter, not quite as long, was an
answer to a very long and patient letter of Tim Porter, te\var{\over[typed]{a}{l}}{l}{typewriter overwriting}{ps_134-5_087-v0008}ling me about
a number of things which have been done by homotopy people, \var{\inserted{and}}{\hand{and}}{insertion}{ps_134-5_087-v0009} which I was
\var{\inserted{of course}}{\hand{of course}}{insertion}{ps_134-5_087-v0010} wholly ignorant of. His letter has been the first echo I got from some-%
one who read part of the notes on ``Pursuing stacks'' , and I was glad he
could make some sense of what he read so far, and conversely - that not
all he was telling me was going wholly ``above my head'' !

Apart from this and the (not unpleasant !) daily routine, I spent
a fair bunch of hours on scratchwork, centering around trying to figure
out the right notion of morphism for asphericity structures. It was a sur-%
prise that the notion should be so reticent for revealing itself - as a
matter of fact, I am not quite sure yet if I got the right notion, in
sufficient generality I mean. Time will tell - for the time being,
the notion of morphism I have to offer, while maybe too restrictive,
looks really seducing, because it parallels so perfectly the formalism
of aspheric functors $i: A \hand{\longrightarrow} M$ and the corresponding (would-be modeli-%
zing) functor $i^{\hand{*}}: M \hand{\to} A^{\wedge}$. Again, it has been a surprise, right now,
that after uncounted hours of unconvincing efforts today and yesterday,
and almost wholly unrelated to these, this pretty set up would come out
within ten minutes reflection !

\margin{87}

%% END SOURCE ps_134-5_087
\ifnum\EditionMode=1
\psapparatusstart
\psentry{ps_134-5_087-v0001}{1}{deletion}{\deleted{the}}{}{The author cancels this material by hand; it leaves the corrected running text.}
\psentry{ps_134-5_087-v0002}{2}{typewriter deletion}{\xout{late}}{}{The typewriter strikes out this material; it leaves the corrected running text.}
\psentry{ps_134-5_087-v0003}{3}{insertion}{\inserted{notably}}{\hand{notably}}{The handwritten insertion is incorporated at its indicated position; insertion brackets leave the running text.}
\psentry{ps_134-5_087-v0004}{4}{typewriter overwriting}{\over[typed]{\illegible{2}}{fr}}{fr}{The superposed typed reading replaces the underlying reading.}
\psentry{ps_134-5_087-v0005}{5}{insertion}{\inserted{of his}}{\hand{of his}}{The handwritten insertion is incorporated at its indicated position; insertion brackets leave the running text.}
\psentry{ps_134-5_087-v0006}{6}{typed insertion}{\inserted[typed]{key}}{key}{The typed insertion is incorporated at its indicated position; insertion brackets leave the running text.}
\psentry{ps_134-5_087-v0007}{7}{typewriter deletion}{\xout{which}}{}{The typewriter strikes out this material; it leaves the corrected running text.}
\psentry{ps_134-5_087-v0008}{8}{typewriter overwriting}{\over[typed]{a}{l}}{l}{The superposed typed reading replaces the underlying reading.}
\psentry{ps_134-5_087-v0009}{9}{insertion}{\inserted{and}}{\hand{and}}{The handwritten insertion is incorporated at its indicated position; insertion brackets leave the running text.}
\psentry{ps_134-5_087-v0010}{10}{insertion}{\inserted{of course}}{\hand{of course}}{The handwritten insertion is incorporated at its indicated position; insertion brackets leave the running text.}
\psapparatusend
\fi
\pssource{ps_134-5_088}{134-5/88}{https://grothendieck.umontpellier.fr/web/viewer.html?file=../134-5.pdf\#page=89}
\psnumber{ps_134-5_088-v0001}{1}
\psnumber{ps_134-5_088-v0002}{2}
\psnumber{ps_134-5_088-v0003}{3}
\psnumber{ps_134-5_088-v0004}{4}
\psnumber{ps_134-5_088-v0005}{5}
\psnumber{ps_134-5_088-v0006}{6}
\psnumber{ps_134-5_088-v0007}{7}
\psnumber{ps_134-5_088-v0008}{8}
\psnumber{ps_134-5_088-v0009}{9}
\psnumber{ps_134-5_088-v0010}{10}
\psnumber{ps_134-5_088-v0011}{11}
%% BEGIN SOURCE ps_134-5_088
% Pursuing Stacks (cotes 134-2 to 134-8): archive cote 134-5, site page 88 (PDF page 89), typescript page 275.
% typescript. Draft compared: ps4.tex:4509-4587.
% AI-transcribed by Claude Opus 5.5 (claude-opus-5-5, adaptive thinking, effort high), Message Batches API, batch ps2 (msgbatch_01NVwr1Yuo8qJf5rykeGBXPK), 2026-11-12; reviewed against both supplied scan halves; see CHECK_DECISIONS.tsv.
\scanpage{134-5/88}
\mspage{\hand{275}}

\var{\inserted[typed]{\hand{\uline{\textcolor{black}{Proposition.}}}}}{\hand{\uline{\textcolor{black}{Proposition.}}}}{typed insertion}{ps_134-5_088-v0001} \var{\xout{\hand{\uline{\textcolor{black}{Theorem}}}}}{}{typewriter deletion}{ps_134-5_088-v0002} Let $(M,M_a)$, $(M',M'_a)$ be two $\underline{W}$-asphericity structures ($\underline{W}$ a basic
localizer), $B \hand{\subset} M_a$ a full subcategory of $M$ generating the asphericity
structure $M_a$ , and
\[ f^* : M \hand{\longrightarrow} M' \]
a functor, admitting a left adjoint \hand{$f_!$} (hence $f^*$ commutes to inverse limits)
We consider the following six conditions, paralleling those of prop.4 of
section 75 (p.236):
\[ \text{(i)} \quad M_a = (f^*)^{-1}(M'_a ) , \]
\[ \text{(i')} \quad M_a \hand{\subset} (f^*)^{-1}(M'_a) , \quad \text{i.e.} \quad f^*(M_a) \hand{\subset} M'_a , \]
\[ \text{(i")} \quad B \hand{\subset} (f^*)^{-1}(M'_a) , \quad \text{i.e.} \quad f^*(B) \hand{\subset} M'_a , \]
\[ \text{(ii)} \quad W_a =(f^*)^{-1}(W'_a) , \]
\[ \text{(ii')} \quad W_a \hand{\subset} (f^*)^{-1}(W'_a) , \quad \text{i.e.} \quad f^*(W_a) \hand{\subset} W'_a , \]
\[ \text{(ii")} \quad \mathrm{F}\var{\over{l}{\ell}}{\hand{\ell}}{writing over}{ps_134-5_088-v0003}(B) \hand{\subset} (f^*)^{-1}(W'_a) \quad \text{i.e.} \quad f^*(\mathrm{F}\var{\over{l}{\ell}}{\hand{\ell}}{writing over}{ps_134-5_088-v0004}(B)) \hand{\subset} W'_a , \]
where $W_a$, $W'_a$ are the sets of $\underline{W}$-equivalences in $M$ and $M'$ respectively.
We'll make the extra assumption (I nearly forgot, sorry !) :
(Awk) There exists a \var{\xout{test functor}}{}{typewriter deletion}{ps_134-5_088-v0005} $M'_a$-$\underline{W}$-aspheric functor
\[ i' : A \hand{\longrightarrow} M' \]
(A a small category), such that \var{\inserted[typed]{$i =$}}{$i =$}{typed insertion}{ps_134-5_088-v0006} $f\hand{_!}i' : A \hand{\rightarrow} M$ factors through C.
Under these conditions and with these notations, the following holds:
The conditions (i),(i'),(i") are equivalent and imply the three others,
which satisfy the tautological implications (ii)$\hand{\Rightarrow}$(ii')$\hand{\Rightarrow}$(ii") . If
M and M' admit final objects and their asphericity structures are aspheric\hand{,}
then \var{\xout{six}}{}{typewriter deletion}{ps_134-5_088-v0007} \var{\inserted[typed]{the five}}{the five}{typed insertion}{ps_134-5_088-v0008} \var{\inserted{first}}{\hand{first}}{insertion}{ps_134-5_088-v0009} conditions (i) to (ii') are equivalent, and if moreover $e_M \hand{\in} \mathrm{Ob}\, B$\hand{,}
all six are equivalent.

Proof: reduction to loc. cit. , using the commutative diagram
\margin{(1)}
% Hand-drawn arrows are marked blue; typed objects remain plain.
\diagram{\begin{tikzcd}[column sep=tiny]
M \ar[rr,"f^*",color=blue] \ar[dr,"i^*"',color=blue] \& \& M' \ar[dl,"i'^*",color=blue] \\
\& \hat{A} \&
\end{tikzcd}}
(up to isomorphism), and the relations
\[ M'_a =(i'^*)^{-1}(\hat{A}_a) \quad , \quad W'_a = (i'^*)^{-1}(\underline{W}_A) \quad . \]
This proof shows moreover that the conditions (i),(i'),(i") are equivalent
each to $i$ being $M_a$-$\underline{W}$-aspheric . Thus, if these conditions are satis-%
\margin{cf p.280\,$\times$}
fied (let's say then that $f^*$ is a \hand{\uline{\textcolor{black}{morphism}}}\var{\inserted{(*)}}{\hand{(*)}}{insertion}{ps_134-5_088-v0010} for the asphericity structu\var{\xout{xx}}{}{typewriter deletion}{ps_134-5_088-v0011}%
res ), we may use $i^{\hand{*}}$ and $i'^{\hand{*}}$ for describing the canonical functors from

\margin{88}

%% END SOURCE ps_134-5_088
\ifnum\EditionMode=1
\psapparatusstart
\psentry{ps_134-5_088-v0001}{1}{typed insertion}{\inserted[typed]{\hand{\uline{\textcolor{black}{Proposition.}}}}}{\hand{\uline{\textcolor{black}{Proposition.}}}}{The typed insertion is incorporated at its indicated position; insertion brackets leave the running text.}
\psentry{ps_134-5_088-v0002}{2}{typewriter deletion}{\xout{\hand{\uline{\textcolor{black}{Theorem}}}}}{}{The typewriter strikes out this material; it leaves the corrected running text.}
\psentry{ps_134-5_088-v0003}{3}{writing over}{$\over{l}{\ell}$}{$\hand{\ell}$}{The superposed handwritten reading replaces the underlying reading.}
\psentry{ps_134-5_088-v0004}{4}{writing over}{$\over{l}{\ell}$}{$\hand{\ell}$}{The superposed handwritten reading replaces the underlying reading.}
\psentry{ps_134-5_088-v0005}{5}{typewriter deletion}{\xout{test functor}}{}{The typewriter strikes out this material; it leaves the corrected running text.}
\psentry{ps_134-5_088-v0006}{6}{typed insertion}{\inserted[typed]{$i =$}}{$i =$}{The typed insertion is incorporated at its indicated position; insertion brackets leave the running text.}
\psentry{ps_134-5_088-v0007}{7}{typewriter deletion}{\xout{six}}{}{The typewriter strikes out this material; it leaves the corrected running text.}
\psentry{ps_134-5_088-v0008}{8}{typed insertion}{\inserted[typed]{the five}}{the five}{The typed insertion is incorporated at its indicated position; insertion brackets leave the running text.}
\psentry{ps_134-5_088-v0009}{9}{insertion}{\inserted{first}}{\hand{first}}{The handwritten insertion is incorporated at its indicated position; insertion brackets leave the running text.}
\psentry{ps_134-5_088-v0010}{10}{insertion}{\inserted{(*)}}{\hand{(*)}}{The handwritten insertion is incorporated at its indicated position; insertion brackets leave the running text.}
\psentry{ps_134-5_088-v0011}{11}{typewriter deletion}{\xout{xx}}{}{The typewriter strikes out this material; it leaves the corrected running text.}
\psapparatusend
\fi
\pssource{ps_134-5_089}{134-5/89}{https://grothendieck.umontpellier.fr/web/viewer.html?file=../134-5.pdf\#page=90}
\psnumber{ps_134-5_089-v0001}{1}
\psnumber{ps_134-5_089-v0002}{2}
\psnumber{ps_134-5_089-v0003}{3}
\psnumber{ps_134-5_089-v0004}{4}
\psnumber{ps_134-5_089-v0005}{5}
\psnumber{ps_134-5_089-v0006}{6}
\psnumber{ps_134-5_089-v0007}{7}
\psnumber{ps_134-5_089-v0008}{8}
\psnumber{ps_134-5_089-v0009}{9}
\psnumber{ps_134-5_089-v0010}{10}
\psnumber{ps_134-5_089-v0011}{11}
%% BEGIN SOURCE ps_134-5_089
% Pursuing Stacks (cotes 134-2 to 134-8): archive cote 134-5, site page 89 (PDF page 90), typescript page 276.
% typescript. Draft compared: ps4.tex:4570-4639.
% AI-transcribed by Claude Opus 5.5 (claude-opus-5-5, adaptive thinking, effort high), Message Batches API, batch ps2 (msgbatch_01NVwr1Yuo8qJf5rykeGBXPK), 2026-11-12; reviewed against both supplied scan halves; see CHECK_DECISIONS.tsv.
\scanpage{134-5/89}
\mspage{276}

the localizations of M,M' to $\mathrm{Hot}_{\underline{W}}$ , and therefore from (1) get a
\hand{\uline{\textcolor{black}{commutative}}} diagram (\var{\over[typed]{\illegible{1}}{u}}{u}{typewriter overwriting}{ps_134-5_089-v0001}p to canonical isomorphism)
% Hand-drawn arrows are marked blue; typed objects remain plain.
\diagram{
\begin{tikzcd}[column sep=small]
\mathrm{Hot}_{M,\underline{W}} \ar[rr, "\hand{\overline{\textcolor{black}{f^{*}}}}",color=blue] \ar[dr,color=blue] \& \& \mathrm{Hot}_{M',\underline{W}} \ar[dl,color=blue] \\
\& \mathrm{Hot}_{\underline{W}} \& \quad . % An incomplete small oval above the o is visible in the scan; no separate textual value is secure.
\end{tikzcd}
}
(2)

From this follows:
\hand{\uline{\textcolor{black}{Corollary}}}: Let $f^{*}$ be a morphism for the given asphericity structures,
and assume these to be \hand{\uline{\textcolor{black}{modelizing}}} . Then the induced functor for the
localizations
\[ \hand{\overline{\textcolor{black}{f^{*}}}} : \; W_{a}^{-1}M = \mathrm{Hot}_{M,\underline{W}} \hand{\longrightarrow} {W'_{a}}^{-1}M' = \mathrm{Hot}_{M',\underline{W}} \]
is an equivalence of categories.

This is the answer, at last, of the ``silly question\var{\xout{s}}{}{typewriter deletion}{ps_134-5_089-v0002}'' of section
45 (p. 95)\hand{!}\var{\deleted{\illegible{1}}}{}{deletion}{ps_134-5_089-v0003}

We have still to comment though on the restrictive conditions we had
to make, for getting the equivalences stated in the proposition, and the
result stated in the corollary. These conditions are twofold: a) exis-%
tence of a left adjoint $f_{!}$ , which \var{\inserted{presumably,}}{\hand{presumably,}}{insertion}{ps_134-5_089-v0004} in most circumstances we are going
to meet, will \var{\deleted{presumably}}{}{deletion}{ps_134-5_089-v0005} be equivalent with $f^{*}$ commuting to inverse li-%
mits. It's a pretty restrictive condition, but of a rather natural kind,
often met with in the modelizing situations; b) this is the ``awkward''
condition (Awk), which can be equally stated as follows: there exists
a \var{\xout{full}}{}{typewriter deletion}{ps_134-5_089-v0006} small full subcategory B' of M' (NB we may take the full subcate-%
gory defined by i'(Ob A)), \hand{\uline{\textcolor{black}{generating}}} the asphericity structure $M'_{a}$ ,
and such moreover that
\[ f_{!}(B') \hand{\subset} M_a^{\var{\deleted{\illegible{1}}}{}{deletion}{ps_134-5_089-v0007}} \quad , \]
(so that we can choose C , a \var{\xout{generating}}{}{typewriter deletion}{ps_134-5_089-v0008} full subcategory of M generating
the asphericity structure, such that C contains $f_{!}(B')$). Another way of
phrasing this condition on $f^{*}$ (preliminary to the choice of C)\var{\xout{namely}}{}{typewriter deletion}{ps_134-5_089-v0009} is
that the full subcategory of M'
\[ (f_{!})^{-1}(M_{a}) \hand{\cap} M'_{a} \tag{3} \]
\var{\xout{of M'}}{}{typewriter deletion}{ps_134-5_089-v0010} % The typed line is enclosed by a pen loop with an upward locator to display (3).
\var{\hand{\tikz[baseline=(m.base)]{\node[text=black,inner xsep=4pt,inner ysep=3pt] (m) {generates the asphericity structure of M'.}; \draw (m.west) to[out=90,in=180] (m.north) to[out=0,in=90] (m.east) to[out=-90,in=0] (m.south) to[out=180,in=-90] (m.west); \draw (m.north) -- ++(0.12,0.28) -- ++(0.12,-0.28);}}}{generates the asphericity structure of M'.}{placement}{ps_134-5_089-v0011}

If this condition, plus the condition (i) say, which we express
jointly by saying that $f^{*}$ is a ``morphism'' of asphericity structures'', did
imply the condition (stronger than (3))

\margin{89}

%% END SOURCE ps_134-5_089
\ifnum\EditionMode=1
\psapparatusstart
\psentry{ps_134-5_089-v0001}{1}{typewriter overwriting}{\over[typed]{\illegible{1}}{u}}{u}{The superposed typed reading replaces the underlying reading.}
\psentry{ps_134-5_089-v0002}{2}{typewriter deletion}{\xout{s}}{}{The typewriter strikes out this material; it leaves the corrected running text.}
\psentry{ps_134-5_089-v0003}{3}{deletion}{\deleted{\illegible{1}}}{}{The author cancels this material by hand; it leaves the corrected running text.}
\psentry{ps_134-5_089-v0004}{4}{insertion}{\inserted{presumably,}}{\hand{presumably,}}{The handwritten insertion is incorporated at its indicated position; insertion brackets leave the running text.}
\psentry{ps_134-5_089-v0005}{5}{deletion}{\deleted{presumably}}{}{The author cancels this material by hand; it leaves the corrected running text.}
\psentry{ps_134-5_089-v0006}{6}{typewriter deletion}{\xout{full}}{}{The typewriter strikes out this material; it leaves the corrected running text.}
\psentry{ps_134-5_089-v0007}{7}{deletion}{$\deleted{\illegible{1}}$}{}{The author cancels this material by hand; it leaves the corrected running text.}
\psentry{ps_134-5_089-v0008}{8}{typewriter deletion}{\xout{generating}}{}{The typewriter strikes out this material; it leaves the corrected running text.}
\psentry{ps_134-5_089-v0009}{9}{typewriter deletion}{\xout{namely}}{}{The typewriter strikes out this material; it leaves the corrected running text.}
\psentry{ps_134-5_089-v0010}{10}{typewriter deletion}{\xout{of M'}}{}{The typewriter strikes out this material; it leaves the corrected running text.}
\psentry{ps_134-5_089-v0011}{11}{placement}{\hand{\tikz[baseline=(m.base)]{\node[text=black,inner xsep=4pt,inner ysep=3pt] (m) {generates the asphericity structure of M'.}; \draw (m.west) to[out=90,in=180] (m.north) to[out=0,in=90] (m.east) to[out=-90,in=0] (m.south) to[out=180,in=-90] (m.west); \draw (m.north) -- ++(0.12,0.28) -- ++(0.12,-0.28);}}}{generates the asphericity structure of M'.}{The typed sentence is enclosed by a pen loop whose upward locator attaches it to display (3). The reading text retains the typed wording and order without the placement graphic.}
\psapparatusend
\fi
\pssource{ps_134-5_090}{134-5/90}{https://grothendieck.umontpellier.fr/web/viewer.html?file=../134-5.pdf\#page=91}
\psnumber{ps_134-5_090-v0001}{1}
\psnumber{ps_134-5_090-v0002}{2}
\psnumber{ps_134-5_090-v0003}{3}
\psnumber{ps_134-5_090-v0004}{4}
\psnumber{ps_134-5_090-v0005}{5}
\psnumber{ps_134-5_090-v0006}{6}
\psnumber{ps_134-5_090-v0007}{7}
\psnumber{ps_134-5_090-v0008}{8}
\psnumber{ps_134-5_090-v0009}{9}
\psnumber{ps_134-5_090-v0010}{10}
%% BEGIN SOURCE ps_134-5_090
% Pursuing Stacks (cotes 134-2 to 134-8): archive cote 134-5, site page 90 (PDF page 91), typescript page 277.
% typescript. Draft compared: ps4.tex:4619-4686.
% AI-transcribed by Claude Opus 5.5 (claude-opus-5-5, adaptive thinking, effort high), Message Batches API, batch ps2 (msgbatch_01NVwr1Yuo8qJf5rykeGBXPK), 2026-11-12; reviewed against both supplied scan halves; see CHECK_DECISIONS.tsv.
\scanpage{134-5/90}
\mspage{277}

(4) \[ f_{!}(M'_a) \hand{\subset} M_a \quad , \]
we would replace (Awk) by this condition (4), which does'nt look awkward
any longer, and does not refer to any A or i'whatsoever (and \var{\xout{if in the}}{}{typewriter deletion}{ps_134-5_090-v0001}
the \var{\xout{conditions}}{}{typewriter deletion}{ps_134-5_090-v0002} six conditions (i) to (ii"), with the exception of (i")
and (ii"), do not make any reference to any given aspheric functor ).
In any case, one might think of defining a \hand{\uline{\textcolor{black}{morphism}}} $f^*$ of asphericity
structures as a functor admitting a left adjoint $f_{!}$ satisfying (4), and
such moreover that (i) above, or equivalently (i'), namely the nicely
symmetric relation to (4)

(5) \[ f^*(M_a) \hand{\subset} M'_a \quad , \]
is satisfied. This notion for a morphism, stricter than the one we adop-%
ted provisionally, looks a lot nicer indeed - the trouble is that I
could not make up my mind if the restriction (4) is not \var{\xout{\unsure{an}}}{}{typewriter deletion}{ps_134-5_090-v0003} an unreasonable
one. It would be reasonable indeed, if for those cases which are the most
interesting for us, and primarily \var{\xout{of all}}{}{typewriter deletion}{ps_134-5_090-v0004}\var{\inserted[typed]{for any}}{for any}{typed insertion}{ps_134-5_090-v0005} functor\var{\xout{s}}{}{typewriter deletion}{ps_134-5_090-v0006}
\[ f^* = i^* : M \hand{\longrightarrow} A^{\wedge} \]
associated to a $M_a$% A curved oblique stroke crosses the left side of the typed M; it does not form an underline, and its separate textual value is unclear.
-$\underline{W}$-aspheric functor
\[ i: A \hand{\to} M \]
for \var{\inserted{any}}{\hand{any}}{insertion}{ps_134-5_090-v0007} \var{\deleted{a}}{}{deletion}{ps_134-5_090-v0008} given "nice" asphericity structure $(M,M_a)$ , this condition is satis-%
fied. If course, in this case, $i^*$ does admit a left adjoint
\[ i_{!} : A^{\wedge} \hand{\longrightarrow} M \quad , \]
provided only M is stable under (s\var{\over{\unsure{am}}{ma}}{\hand{ma}}{writing over}{ps_134-5_090-v0009}ll) direct limits, which we'll assume
without any reluctance ! Thus we are led to the following

\hand{\uline{\textcolor{black}{Question}}}\hand{1}. Under which conditions is it true, for a $M_a$-$\underline{W}$-aspheric
functor $i: A \hand{\to} M$ ( where M is endowed with an asphericity structure $M_a$),
that

(6) \[ i_{!}(A^{\wedge}_a \ ) \hand{\subset} M_a \quad , \]
i.e. $i_{!}$ (the canonical extension of i to $A^{\wedge}$, commuting with direct
limits) takes aspheric objects into aspheric objects ?

\var{\xout{Of special}}{}{typewriter deletion}{ps_134-5_090-v0010} Yesterday and today I pondered mainly about the typical
case when $M = B^{\wedge}$, endowed with it's canonical asphericity structure,

\margin{90}

%% END SOURCE ps_134-5_090
\ifnum\EditionMode=1
\psapparatusstart
\psentry{ps_134-5_090-v0001}{1}{typewriter deletion}{\xout{if in the}}{}{The typewriter strikes out this material; it leaves the corrected running text.}
\psentry{ps_134-5_090-v0002}{2}{typewriter deletion}{\xout{conditions}}{}{The typewriter strikes out this material; it leaves the corrected running text.}
\psentry{ps_134-5_090-v0003}{3}{typewriter deletion}{\xout{\unsure{an}}}{}{The typewriter strikes out this material; it leaves the corrected running text.}
\psentry{ps_134-5_090-v0004}{4}{typewriter deletion}{\xout{of all}}{}{The typewriter strikes out this material; it leaves the corrected running text.}
\psentry{ps_134-5_090-v0005}{5}{typed insertion}{\inserted[typed]{for any}}{for any}{The typed insertion is incorporated at its indicated position; insertion brackets leave the running text.}
\psentry{ps_134-5_090-v0006}{6}{typewriter deletion}{\xout{s}}{}{The typewriter strikes out this material; it leaves the corrected running text.}
\psentry{ps_134-5_090-v0007}{7}{insertion}{\inserted{any}}{\hand{any}}{The handwritten insertion is incorporated at its indicated position; insertion brackets leave the running text.}
\psentry{ps_134-5_090-v0008}{8}{deletion}{\deleted{a}}{}{The author cancels this material by hand; it leaves the corrected running text.}
\psentry{ps_134-5_090-v0009}{9}{writing over}{\over{\unsure{am}}{ma}}{\hand{ma}}{The superposed handwritten reading replaces the underlying reading.}
\psentry{ps_134-5_090-v0010}{10}{typewriter deletion}{\xout{Of special}}{}{The typewriter strikes out this material; it leaves the corrected running text.}
\psapparatusend
\fi
\pssource{ps_134-5_091}{134-5/91}{https://grothendieck.umontpellier.fr/web/viewer.html?file=../134-5.pdf\#page=92}
\psnumber{ps_134-5_091-v0001}{1}
\psnumber{ps_134-5_091-v0002}{2}
\psnumber{ps_134-5_091-v0003}{3}
\psnumber{ps_134-5_091-v0004}{4}
\psnumber{ps_134-5_091-v0005}{5}
\psnumber{ps_134-5_091-v0006}{6}
\psnumber{ps_134-5_091-v0007}{7}
\psnumber{ps_134-5_091-v0008}{8}
\psnumber{ps_134-5_091-v0009}{9}
\psnumber{ps_134-5_091-v0010}{10}
\psnumber{ps_134-5_091-v0011}{11}
\psnumber{ps_134-5_091-v0012}{12}
\psnumber{ps_134-5_091-v0013}{13}
%% BEGIN SOURCE ps_134-5_091
% Pursuing Stacks (cotes 134-2 to 134-8): archive cote 134-5, site page 91 (PDF page 92), typescript page 278.
% typescript. Draft compared: ps4.tex:4666-4799.
% AI-transcribed by Claude Opus 5.5 (claude-opus-5-5, adaptive thinking, effort high), Message Batches API, batch ps2 (msgbatch_01NVwr1Yuo8qJf5rykeGBXPK), 2026-11-12; reviewed against both supplied scan halves; see CHECK_DECISIONS.tsv.
\scanpage{134-5/91}
\mspage{\hand{278}}

where B is a small category, \var{\inserted{and moreover $i$ is a functor $i: A\to B\subset B^\wedge$.}}{\hand{and moreover $i$ is a functor $i: A\to B\subset B^\wedge$.}}{insertion}{ps_134-5_091-v0001} This then brings us to the related

\hand{\uline{\textcolor{black}{Question}}} 1': Let i: $A\hand{\to} B$ be an aspherical map in (Cat), under which
restrictive conditions on A,B,i (if any) does \var{\xout{$i_!$\illegible{1}}}{}{typewriter deletion}{ps_134-5_091-v0002}
\[ i_! : A^\wedge\hand{\to} B^\wedge \]
take aspheric objects into aspheric ones, i.e. do we have
\[ (7) \qquad i_!(A^\wedge_{as}) \hand{\subset} B^\wedge_{as} \qquad ? \]

The question makes a sense even without assuming i to be aspheric.
If we drop this asphericity assumption on i, we are reduced, for a given
B , (\var{\xout{and}}{}{typewriter deletion}{ps_134-5_091-v0003}\var{\inserted[typed]{by}}{by}{typed insertion}{ps_134-5_091-v0004} replacing A by $A_{/F}$ ) where F is a given \var{\xout{object}}{}{typewriter deletion}{ps_134-5_091-v0005} aspheric object of
$A^\wedge$) to the case when A is aspheric and $F=e_{A^\wedge}$ , i.e. to l\var{\over[typed]{l}{o}}{o}{typewriter overwriting}{ps_134-5_091-v0006}oking at whe-%
ther the element
\[ (8) \qquad i_!(e_{A^\wedge}) \;=\; \mathop{\underset{\hand{\longrightarrow}}{\mathrm{lim}}}\limits_{A}{}_{(B^\wedge)}\, i(a) \]
in $B^\wedge$\hand{/}is aspheric. This element appears as the direct limit in $B^\wedge$ of
aspheric elements i(a), the indexing category A being itself aspheric.
Thus it is rather tempting to hope that the limit might well be aspheric
too. Let's assume $\underline{W}=\underline{W}_{\mathrm{oo}}$ = usual weak equivalence, it turns out that when
A is 1-connected so is (8), which would seem to give some support to the
hope that (8) is aspheric if A is. However this is defintely not so, un-%
less at the very least we assume B to be equally aspheric (which is howe-%
ver a rather natural assumption, as it follows from A aspheric if we assu%
-me moreover i \var{\inserted{to be}}{\hand{to be}}{insertion}{ps_134-5_091-v0007} aspheric). To see this, we take i to be cofibering with
0-connected fibers - the cofibration assumption implies that the direct
limit\var{\deleted{s}}{}{deletion}{ps_134-5_091-v0008} \var{\inserted{over A}}{\hand{over A}}{insertion}{ps_134-5_091-v0009} can be computed first fiberwise and then take a limit over B
(the most general version of associativity for direct limits ! ), whereas
the 0-connectedness assumption on the fibers implies that the limit taken
over the fiber $A_b$ is just b itself, hence (8) is isomorphic to
\[ \mathop{\underset{\hand{\longrightarrow}}{\mathrm{lim}}}\limits_{b}(B^{\var{\over{\illegible{1}}{\wedge}}{\hand{\wedge}}{writing over}{ps_134-5_091-v0010}})\,\var{\deleted{\illegible{1}}}{}{deletion}{ps_134-5_091-v0011}\var{\inserted[typed]{b}}{b}{typed insertion}{ps_134-5_091-v0012} \;=\; e_{B^\wedge} \quad , \]
which is'nt aspheric except precisely when B is; now for topological
reasons it is easy to find an aspheric A cofibered with 0-connected fibers
over a non-aspheric B.

Thus in the question of asphericity of (8\var{\over{p}{)}}{\hand{)}}{writing over}{ps_134-5_091-v0013} we'll better assume both
A \hand{\uline{\textcolor{black}{and}}} B aspheric. If either A has a final object or B has an initial ob-%
ject, asphericity of (8) is more or less trivial in any case. For cer-%
tain categories B even without initial object, (8) is always aspheric
when A is, this I checked at any rate for the ordered category
% Hand-drawn arrows are marked blue; typed objects remain plain.
\[ B \;=\; \diagram{\begin{tikzcd}[row sep=small, column sep=small] \hand{\alpha} \arrow[dr,color=blue] \& \\ \& \hand{\gamma} \\ \hand{\beta} \arrow[ur,color=blue] \& \end{tikzcd}} \quad , \]

\margin{91}

%% END SOURCE ps_134-5_091
\ifnum\EditionMode=1
\psapparatusstart
\psentry{ps_134-5_091-v0001}{1}{insertion}{\inserted{and moreover $i$ is a functor $i: A\to B\subset B^\wedge$.}}{\hand{and moreover $i$ is a functor $i: A\to B\subset B^\wedge$.}}{The handwritten insertion is incorporated at its indicated position; insertion brackets leave the running text.}
\psentry{ps_134-5_091-v0002}{2}{typewriter deletion}{\xout{$i_!$\illegible{1}}}{}{The typewriter strikes out this material; it leaves the corrected running text.}
\psentry{ps_134-5_091-v0003}{3}{typewriter deletion}{\xout{and}}{}{The typewriter strikes out this material; it leaves the corrected running text.}
\psentry{ps_134-5_091-v0004}{4}{typed insertion}{\inserted[typed]{by}}{by}{The typed insertion is incorporated at its indicated position; insertion brackets leave the running text.}
\psentry{ps_134-5_091-v0005}{5}{typewriter deletion}{\xout{object}}{}{The typewriter strikes out this material; it leaves the corrected running text.}
\psentry{ps_134-5_091-v0006}{6}{typewriter overwriting}{\over[typed]{l}{o}}{o}{The superposed typed reading replaces the underlying reading.}
\psentry{ps_134-5_091-v0007}{7}{insertion}{\inserted{to be}}{\hand{to be}}{The handwritten insertion is incorporated at its indicated position; insertion brackets leave the running text.}
\psentry{ps_134-5_091-v0008}{8}{deletion}{\deleted{s}}{}{The author cancels this material by hand; it leaves the corrected running text.}
\psentry{ps_134-5_091-v0009}{9}{insertion}{\inserted{over A}}{\hand{over A}}{The handwritten insertion is incorporated at its indicated position; insertion brackets leave the running text.}
\psentry{ps_134-5_091-v0010}{10}{writing over}{$\over{\illegible{1}}{\wedge}$}{$\hand{\wedge}$}{The superposed handwritten reading replaces the underlying reading.}
\psentry{ps_134-5_091-v0011}{11}{deletion}{$\deleted{\illegible{1}}$}{}{The author cancels this material by hand; it leaves the corrected running text.}
\psentry{ps_134-5_091-v0012}{12}{typed insertion}{$\inserted[typed]{b}$}{$b$}{The typed insertion is incorporated at its indicated position; insertion brackets leave the running text.}
\psentry{ps_134-5_091-v0013}{13}{writing over}{\over{p}{)}}{\hand{)}}{The superposed handwritten reading replaces the underlying reading.}
\psapparatusend
\fi
\pssource{ps_134-5_092}{134-5/92}{https://grothendieck.umontpellier.fr/web/viewer.html?file=../134-5.pdf\#page=93}
\psnumber{ps_134-5_092-v0001}{1}
\psnumber{ps_134-5_092-v0002}{2}
\psnumber{ps_134-5_092-v0003}{3}
\psnumber{ps_134-5_092-v0004}{4}
\psnumber{ps_134-5_092-v0005}{5}
\psnumber{ps_134-5_092-v0006}{6}
\psnumber{ps_134-5_092-v0007}{7}
\psnumber{ps_134-5_092-v0008}{8}
\psnumber{ps_134-5_092-v0009}{9}
\psnumber{ps_134-5_092-v0010}{10}
\psnumber{ps_134-5_092-v0011}{11}
\psnumber{ps_134-5_092-v0012}{12}
\psnumber{ps_134-5_092-v0013}{13}
\psnumber{ps_134-5_092-v0014}{14}
\psnumber{ps_134-5_092-v0015}{15}
%% BEGIN SOURCE ps_134-5_092
% Pursuing Stacks (cotes 134-2 to 134-8): archive cote 134-5, site page 92 (PDF page 93), typescript page 279.
% typescript. Draft compared: ps4.tex:4666-4799.
% AI-transcribed by Claude Opus 5.5 (claude-opus-5-5, adaptive thinking, effort high), Message Batches API, batch ps2 (msgbatch_01NVwr1Yuo8qJf5rykeGBXPK), 2026-11-12; reviewed against both supplied scan halves; see CHECK_DECISIONS.tsv.
\scanpage{134-5/92}
\mspage{279}

giving rise to a rather interesting computation.   This suggested,  for
arbitrary B again ,  to look at the dual of the category above \var{\xout{for}}{}{typewriter deletion}{ps_134-5_092-v0001} as A :
% Hand-drawn arrows are marked blue; typed objects remain plain.
\[ (9) \qquad A = \diagram{\begin{tikzcd}[row sep=-3pt,column sep=small] \hand{\alpha}\ar[dr,color=blue] \& \\ \& \hand{\gamma} \\ \hand{\beta}\ar[ur,color=blue] \& \end{tikzcd}} . \]
Here the question then amounts to \hand{\uline{\textcolor{black}{whether for a diagram}}}
\[ \diagram{\begin{tikzcd}[row sep=-3pt,column sep=small] \& a \\ c\ar[ur,color=blue]\ar[dr,color=blue] \& \\ \& b \end{tikzcd}} \]

\hand{\uline{\textcolor{black}{in}}} B ,  \hand{\uline{\textcolor{black}{the amalgamated sum}}}
\[ (\var{\over{9}{10}}{\hand{10}}{writing over}{ps_134-5_092-v0002}) \qquad a \hand{\amalg}_c b \]
\hand{\uline{\textcolor{black}{in}}} B$\hat{\ }$ \hand{\uline{\textcolor{black}{is aspheric}}}.  This I know to be true if either  $c \hand{\to} a$  or $c \hand{\to} B$  are
monomorphisms (as stated in section 70),  Now I don't really expect this
\margin{$\times$} to be true in general,  even if B is such an excellent category as \hand{$\Delta$} say,
however I did'nt push \hand{t}\var{\over{p}{\unsure{h}}}{\hand{\unsure{h}}}{writing over}{ps_134-5_092-v0003}rough and make a counter\var{\over[typed]{\unsure{t}}{e}}{e}{typewriter overwriting}{ps_134-5_092-v0004}xample.  If we now want
to  remember the asphericity condition on i which we dropped,  this condi-%
tion,  when A has an initial object \var{\over{a}{$\gamma$}}{\hand{$\gamma$}}{writing over}{ps_134-5_092-v0005},  just means that \var{\inserted{c=}}{\hand{c=}}{insertion}{ps_134-5_092-v0006}i(\var{\over{a}{$\gamma$}}{\hand{$\gamma$}}{writing over}{ps_134-5_092-v0007}) is an initial
object of B .   Even with this extra condition,  I do not really expect
\margin{! $\times$} (\var{\over{9}{10}}{\hand{10}}{writing over}{ps_134-5_092-v0008})\hand{\illegible{1}} to be necessarily aspheric.  Therefore,  I do not expect the inclu-%
sion (7) to hold for any aspheric functor i,  even when A (and hence B)
are assumed to be aspheric,  without some extra condition,  on A or B say.

% The h in The has two small horizontal impressions above and an h-like impression below; their textual role is unclear.
The condition that A$\hat{\ }$ be \hand{\uline{\textcolor{black}{totally aspheric}}}  seems in this context a
rather natural one  -  for instance,  when reducing the question whether
(7)  holds to the question of asphericity of an object (8) (where A stands
for  $A_{/F}$  ),  and if we do not want to loose the \var{\xout{\illegible{5}}}{}{typewriter deletion}{ps_134-5_092-v0009} asphericity assump-%
tion for i when taking the composition  $A_{/F} \hand{\to} A \hand{\xrightarrow{\textcolor{black}{i}}} B$ , we would like
that asphericity of F imply asphericity of $A_{/F} \hand{\to} A$  -  which means preci-%
sely that A is totally aspheric.  The trouble is that in the would-be
counterexample above with A given by (9),  A is stable under Inf  i.e.
under binary product,  hence A$\hat{\ }$ \hand{\uline{\textcolor{black}{is}}}  totally aspheric indeed  -  thus it
is doubtful \var{\xout{wha}}{}{typewriter deletion}{ps_134-5_092-v0010}that this extra condition in the \hand{``}question 1'\hand{''} above is
quite enough.  If it does fail indeed,  the next best wo\var{\over[typed]{y}{u}}{u}{typewriter overwriting}{ps_134-5_092-v0011}ld be to \var{\xout{assume}}{}{typewriter deletion}{ps_134-5_092-v0012}
try the stronger condition `` A a strict test category''  (NB the category
(9) above is'nt a test category !),  which leads to the question whether
(8) is aspheric when A is a test category (no\var{\xout{t}}{}{typewriter deletion}{ps_134-5_092-v0013} longer a strict one
though !)  and  $i : A \hand{\to} B$ an aspheric functor.   But I confess I have
no idea at present how to handle this question,  and I am dubious there
will come out any positive result along these lines,  even when assuming
A and B to be both stable under \var{\xout{\unsure{finite}}}{}{typewriter deletion}{ps_134-5_092-v0014}\var{\inserted[typed]{binary}}{binary}{typed insertion}{ps_134-5_092-v0015} products  say\hand{,} and to be contrac-%
tors and what not !

\margin{92}

%% END SOURCE ps_134-5_092
\ifnum\EditionMode=1
\psapparatusstart
\psentry{ps_134-5_092-v0001}{1}{typewriter deletion}{\xout{for}}{}{The typewriter strikes out this material; it leaves the corrected running text.}
\psentry{ps_134-5_092-v0002}{2}{writing over}{$\over{9}{10}$}{$\hand{10}$}{The superposed handwritten reading replaces the underlying reading.}
\psentry{ps_134-5_092-v0003}{3}{writing over}{\over{p}{\unsure{h}}}{\hand{\unsure{h}}}{The superposed handwritten reading replaces the underlying reading.}
\psentry{ps_134-5_092-v0004}{4}{typewriter overwriting}{\over[typed]{\unsure{t}}{e}}{e}{The superposed typed reading replaces the underlying reading.}
\psentry{ps_134-5_092-v0005}{5}{writing over}{\over{a}{$\gamma$}}{\hand{$\gamma$}}{The superposed handwritten reading replaces the underlying reading.}
\psentry{ps_134-5_092-v0006}{6}{insertion}{\inserted{c=}}{\hand{c=}}{The handwritten insertion is incorporated at its indicated position; insertion brackets leave the running text.}
\psentry{ps_134-5_092-v0007}{7}{writing over}{\over{a}{$\gamma$}}{\hand{$\gamma$}}{The superposed handwritten reading replaces the underlying reading.}
\psentry{ps_134-5_092-v0008}{8}{writing over}{\over{9}{10}}{\hand{10}}{The superposed handwritten reading replaces the underlying reading.}
\psentry{ps_134-5_092-v0009}{9}{typewriter deletion}{\xout{\illegible{5}}}{}{The typewriter strikes out this material; it leaves the corrected running text.}
\psentry{ps_134-5_092-v0010}{10}{typewriter deletion}{\xout{wha}}{}{The typewriter strikes out this material; it leaves the corrected running text.}
\psentry{ps_134-5_092-v0011}{11}{typewriter overwriting}{\over[typed]{y}{u}}{u}{The superposed typed reading replaces the underlying reading.}
\psentry{ps_134-5_092-v0012}{12}{typewriter deletion}{\xout{assume}}{}{The typewriter strikes out this material; it leaves the corrected running text.}
\psentry{ps_134-5_092-v0013}{13}{typewriter deletion}{\xout{t}}{}{The typewriter strikes out this material; it leaves the corrected running text.}
\psentry{ps_134-5_092-v0014}{14}{typewriter deletion}{\xout{\unsure{finite}}}{}{The typewriter strikes out this material; it leaves the corrected running text.}
\psentry{ps_134-5_092-v0015}{15}{typed insertion}{\inserted[typed]{binary}}{binary}{The typed insertion is incorporated at its indicated position; insertion brackets leave the running text.}
\psapparatusend
\fi
\pssource{ps_134-5_093}{134-5/93}{https://grothendieck.umontpellier.fr/web/viewer.html?file=../134-5.pdf\#page=94}
\psnumber{ps_134-5_093-v0001}{1}
\psnumber{ps_134-5_093-v0002}{2}
\psnumber{ps_134-5_093-v0003}{3}
\psnumber{ps_134-5_093-v0004}{4}
\psnumber{ps_134-5_093-v0005}{5}
\psnumber{ps_134-5_093-v0006}{6}
\psnumber{ps_134-5_093-v0007}{7}
\psnumber{ps_134-5_093-v0008}{8}
%% BEGIN SOURCE ps_134-5_093
% Pursuing Stacks (cotes 134-2 to 134-8): archive cote 134-5, site page 93 (PDF page 94), typescript page 280.
% typescript. Draft compared: ps4.tex:4783-4850.
% AI-transcribed by Claude Opus 5.5 (claude-opus-5-5, adaptive thinking, effort high), Message Batches API, batch ps2 (msgbatch_01NVwr1Yuo8qJf5rykeGBXPK), 2026-11-12; reviewed against both supplied scan halves; see CHECK_DECISIONS.tsv.
\scanpage{134-5/93}
\mspage{280}

Another typical case for ``Question 1'' (p.277) is the case when
\[ M = (\mathrm{Cat}) \quad , \]
endowed with the usual asphericity structure, giving rise to \var{\deleted{a}}{}{deletion}{ps_134-5_093-v0001}\var{\inserted[typed]{the}}{the}{typed insertion}{ps_134-5_093-v0002} notion of
``test functor''
\[ i\colon A \hand{\longrightarrow} (\mathrm{Cat}) \]
we have been working with almost from the very start. The most important
case of all is of course the canonical functor
\[ i_A : a \hand{\longmapsto} A_{/a} : A \hand{\longrightarrow} (\mathrm{Cat}) \quad , \]
and we well know that this functor is aspheric (for the natural aspherici-%
ty structure of (Cat)) if\var{\deleted{f}}{}{deletion}{ps_134-5_093-v0003} A is a \var{\inserted[typed]{weak}}{weak}{typed insertion}{ps_134-5_093-v0004} test category , and in this case it
is true indeed that not only $i_A^{*}$ , but equally $i_{A\,!}$ is modelizing, and
transforms aspheric objects into aspheric objects. But we know too that
for more general test functors, for instance the standard inclusion
\[ (11) \qquad i\colon \hand{\Delta} \hand{\hookrightarrow} (\mathrm{Cat}) \quad , \]
it is no longer true in general that $i_!$ be modelizing - so is it at all
reasonable to expect
\[ (1\var{\over[typed]{1}{2}}{2}{typewriter overwriting}{ps_134-5_093-v0005}) \qquad i_!(A^{\hand{\wedge}}_a) \hand{\subset} (\mathrm{Cat})_a \qquad ? \]
Definitely I'll have to find out the answer in the typical case (11),
whether I like it or not !


30.6.

85. Just got an impressive heap of reprints and preprints by Tim
Porter (announced in his letter two days ago). Many titles have to do
with ``shape theory'', ``coherence'' and``homotopy limits'' (\var{\over[typed]{i}{c}}{c}{typewriter overwriting}{ps_134-5_093-v0006}orresponding to
proobjects in various modelizers, as I understand it from his letter).
The one which most attracts my attention though is ``\hand{\uline{\textcolor{black}{Cat}}} as a closed model
category'' by R.W. Thomason - the title comes \var{\over{\unsure{p}}{q}}{\hand{q}}{writing over}{ps_134-5_093-v0007}uite as a surprise, as
my ponderings on the homotopy structure of (Cat) had left me with the de-%
finite feeling that there was'nt a closed model structure on (Cat), with
the usual notion of weak equivalence. I'll have to \var{\xout{look}}{}{typewriter deletion}{ps_134-5_093-v0008} have a closer
\margin{$\times$} look at this paper definitely, before starting on part V of these notes !

It was getting prohibitively late yesterday, so I had to stop. It
then occured to me that the (admittedly provisional) notion of \hand{\uline{\textcolor{black}{morphism}}}
of asphericity structures I had proposed yesterday (p.275) is definitely
stupid, because there is no reason whatever that it should be stable un-%
der composition ! This surely was the reason for the feeling of uneasi-%

\margin{93}

%% END SOURCE ps_134-5_093
\ifnum\EditionMode=1
\psapparatusstart
\psentry{ps_134-5_093-v0001}{1}{deletion}{\deleted{a}}{}{The author cancels this material by hand; it leaves the corrected running text.}
\psentry{ps_134-5_093-v0002}{2}{typed insertion}{\inserted[typed]{the}}{the}{The typed insertion is incorporated at its indicated position; insertion brackets leave the running text.}
\psentry{ps_134-5_093-v0003}{3}{deletion}{\deleted{f}}{}{The author cancels this material by hand; it leaves the corrected running text.}
\psentry{ps_134-5_093-v0004}{4}{typed insertion}{\inserted[typed]{weak}}{weak}{The typed insertion is incorporated at its indicated position; insertion brackets leave the running text.}
\psentry{ps_134-5_093-v0005}{5}{typewriter overwriting}{$\over[typed]{1}{2}$}{$2$}{The superposed typed reading replaces the underlying reading.}
\psentry{ps_134-5_093-v0006}{6}{typewriter overwriting}{\over[typed]{i}{c}}{c}{The superposed typed reading replaces the underlying reading.}
\psentry{ps_134-5_093-v0007}{7}{writing over}{\over{\unsure{p}}{q}}{\hand{q}}{The superposed handwritten reading replaces the underlying reading.}
\psentry{ps_134-5_093-v0008}{8}{typewriter deletion}{\xout{look}}{}{The typewriter strikes out this material; it leaves the corrected running text.}
\psapparatusend
\fi
\pssource{ps_134-5_094}{134-5/94}{https://grothendieck.umontpellier.fr/web/viewer.html?file=../134-5.pdf\#page=95}
\psnumber{ps_134-5_094-v0001}{1}
\psnumber{ps_134-5_094-v0002}{2}
\psnumber{ps_134-5_094-v0003}{3}
\psnumber{ps_134-5_094-v0004}{4}
\psnumber{ps_134-5_094-v0005}{5}
\psnumber{ps_134-5_094-v0006}{6}
\psnumber{ps_134-5_094-v0007}{7}
\psnumber{ps_134-5_094-v0008}{8}
\psnumber{ps_134-5_094-v0009}{9}
\psnumber{ps_134-5_094-v0010}{10}
%% BEGIN SOURCE ps_134-5_094
% Pursuing Stacks (cotes 134-2 to 134-8): archive cote 134-5, site page 94 (PDF page 95), typescript page 281.
% typescript. Draft compared: ps4.tex:4837-4898.
% AI-transcribed by Claude Opus 5.5 (claude-opus-5-5, adaptive thinking, effort high), Message Batches API, batch ps2 (msgbatch_01NVwr1Yuo8qJf5rykeGBXPK), 2026-11-12; reviewed against both supplied scan halves; see CHECK_DECISIONS.tsv.
\scanpage{134-5/94}
\mspage{281}

ness caused by the condition (Awk),  which surely deserves it's name !
This flaw of course disappears,  if we strengthen the condition,  as sug-%
gested on p.277,  by condition (4)  -  that $f_!$ should take aspheric ob-%
jects into aspheric ones.  This makes all the more imperative the task of
finding out whether this condition is reasonable,  and otherwise, what to
put in as a substitute.  % The h of The has small extra impressions above and below; no separate textual reading is secure.
The test case now seems to be really the one when
$f^*$ is the usual nerve functor
\[ f^* = i^* :\ (\mathrm{Cat}) \hand{\longrightarrow} \hand{\Delta}^{\wedge} \ , \]
associated to the standard inclusion
\[ i :\ \hand{\Delta} \hand{\longrightarrow} (\mathrm{Cat}) \quad . \]
I never before had a closer look upon the corresponding functor
\[ i_{\var{\over{\illegible{1}}{!}}{\hand{!}}{writing over}{ps_134-5_094-v0001}} :\ \hand{\Delta}^{\wedge} \hand{\longrightarrow} (\mathrm{Cat}) \]
left adjoint to the nerve functor,  except just \var{\xout{noticing}}{}{typewriter deletion}{ps_134-5_094-v0002} rectifying a big
blunder (p.22),  and convincing myself that $i_!$ did \hand{\uline{\textcolor{black}{not}}} take weak equiva-%
lences into weak equivalences.  But why not aspheric objects into asphe-%
ric ones ?

Pondering a bit over this matter today,  and trying to see whether
$i_!$  \var{\over{\unsure{took}}{takes}}{\hand{takes}}{writing over}{ps_134-5_094-v0003} \hand{\uline{\textcolor{black}{contractible}}} objects into contractible one's,  this brought up
the question whether $i_!$ commutes with finite products -  which will allow
then to make use of the generalities on morphisms of contractibility
structures,  and will imply that $i_!$  is indeed such a morphism (as it
transforms $\hand{\Delta}$ ,  which generates the contractibility structure of $\hand{\Delta}^{\wedge}$,
into contractible elements of (Cat)).  As a matter of fact,  a number of
times in my scratchwork during the last four months I've met with the
question of when a functor of the type $f_!$ commutes to various types of
finite inverse limits,  and this now is the occasion for writing down
some useful general facts in this respect,  which had remained \var{\inserted{somewhat}}{\hand{somewhat}}{insertion}{ps_134-5_094-v0004} in the air
so far.  It seems that the relevant facts can be summed up in two steps.

\hand{\uline{\textcolor{black}{Proposition}}} 1 .  Let M be a category stable under (small) direct limits,
A a small category,  f: A $\hand{\longrightarrow}$ M a functor,  and
\[ f_! :\ A^{\wedge} \hand{\longrightarrow} M \quad , \quad f^* :\ M \hand{\to} A^{\wedge} \]
the corresponding pair of adjoint functors.

a)  Assume M stable under binary products,  and that \var{\deleted{products}}{}{deletion}{ps_134-5_094-v0005}\var{\inserted{for any $x$}}{\hand{for any $x$}}{insertion}{ps_134-5_094-v0006} in M\hand{,}
\var{\deleted{are distributive with respect to arbitrary sums}}{}{deletion}{ps_134-5_094-v0007}\var{\inserted{the product functor $y \mapsto y\times x$ commutes to arbitrary direct limits}}{\hand{the product functor $y \mapsto y\times x$ commutes to arbitrary direct limits}}{insertion}{ps_134-5_094-v0008} (with small \var{\deleted{sets}}{}{deletion}{ps_134-5_094-v0009}\var{\inserted{categories}}{\hand{categories}}{insertion}{ps_134-5_094-v0010} of indi-%
ces of course).  Then,  in order for $f_{\hand{!}}$ to commute with binary products,

\margin{94}

%% END SOURCE ps_134-5_094
\ifnum\EditionMode=1
\psapparatusstart
\psentry{ps_134-5_094-v0001}{1}{writing over}{$\over{\illegible{1}}{!}$}{$\hand{!}$}{The superposed handwritten reading replaces the underlying reading.}
\psentry{ps_134-5_094-v0002}{2}{typewriter deletion}{\xout{noticing}}{}{The typewriter strikes out this material; it leaves the corrected running text.}
\psentry{ps_134-5_094-v0003}{3}{writing over}{\over{\unsure{took}}{takes}}{\hand{takes}}{The superposed handwritten reading replaces the underlying reading.}
\psentry{ps_134-5_094-v0004}{4}{insertion}{\inserted{somewhat}}{\hand{somewhat}}{The handwritten insertion is incorporated at its indicated position; insertion brackets leave the running text.}
\psentry{ps_134-5_094-v0005}{5}{deletion}{\deleted{products}}{}{The author cancels this material by hand; it leaves the corrected running text.}
\psentry{ps_134-5_094-v0006}{6}{insertion}{\inserted{for any $x$}}{\hand{for any $x$}}{The handwritten insertion is incorporated at its indicated position; insertion brackets leave the running text.}
\psentry{ps_134-5_094-v0007}{7}{deletion}{\deleted{are distributive with respect to arbitrary sums}}{}{The author cancels this material by hand; it leaves the corrected running text.}
\psentry{ps_134-5_094-v0008}{8}{insertion}{\inserted{the product functor $y \mapsto y\times x$ commutes to arbitrary direct limits}}{\hand{the product functor $y \mapsto y\times x$ commutes to arbitrary direct limits}}{The handwritten insertion is incorporated at its indicated position; insertion brackets leave the running text.}
\psentry{ps_134-5_094-v0009}{9}{deletion}{\deleted{sets}}{}{The author cancels this material by hand; it leaves the corrected running text.}
\psentry{ps_134-5_094-v0010}{10}{insertion}{\inserted{categories}}{\hand{categories}}{The handwritten insertion is incorporated at its indicated position; insertion brackets leave the running text.}
\psapparatusend
\fi
\pssource{ps_134-5_095}{134-5/95}{https://grothendieck.umontpellier.fr/web/viewer.html?file=../134-5.pdf\#page=96}
\psnumber{ps_134-5_095-v0001}{1}
\psnumber{ps_134-5_095-v0002}{2}
\psnumber{ps_134-5_095-v0003}{3}
\psnumber{ps_134-5_095-v0004}{4}
\psnumber{ps_134-5_095-v0005}{5}
\psnumber{ps_134-5_095-v0006}{6}
\psnumber{ps_134-5_095-v0007}{7}
\psnumber{ps_134-5_095-v0008}{8}
%% BEGIN SOURCE ps_134-5_095
% Pursuing Stacks (cotes 134-2 to 134-8): archive cote 134-5, site page 95 (PDF page 96), typescript page 282.
% typescript. Draft compared: ps4.tex:4881-4936.
% AI-transcribed by Claude Opus 5.5 (claude-opus-5-5, adaptive thinking, effort high), Message Batches API, batch ps2 (msgbatch_01NVwr1Yuo8qJf5rykeGBXPK), 2026-11-12; reviewed against both supplied scan halves; see CHECK_DECISIONS.tsv.
\scanpage{134-5/95}
\mspage{282}

it is ns that for any two objects a,b of A, the map
\margin{(*)} \[ f_{\hand{!}}(axb)\hand{\longrightarrow} f(a)xf(b) \]
in M be an isomorphism.

b) Assume M stable under fibered products, and that base change in
M commutes to arbitrary direct limits (with small indexing categories of
course). Then, in order for f to commute with fibered products, it is
ns that it does so for any diagram
% Hand-drawn arrows are marked blue; typed objects remain plain.
\diagram{\begin{tikzcd}[column sep=tiny,row sep=small] b \ar[dr,color=blue] \& \& c \ar[dl,color=blue] \\ \& \var{\over{\illegible{1}}{F}}{\hand{F}}{writing over}{ps_134-5_095-v0001} \& \end{tikzcd}} \margin{)}
with \var{\inserted{b,c}}{\hand{b,c}}{insertion}{ps_134-5_095-v0002} objects in A\hand{, $F$ in $A^\wedge$}.

\hand{\uline{\textcolor{black}{Corollary}}} Assume M stable under finite inverse limits, and that base
change in M commutes with arbitrary direct limits. Then $f_{!}$ is left exact
iff \var{\xout{\unsure{it}}}{}{typewriter deletion}{ps_134-5_095-v0003} it satisfies the conditions in a) and b) above, and moreover
transforms $e_{A^\wedge}$ into a final object of M (which also means, if A has a
final object $e_A$ , that $f(e_A)$ is a final object of M).

Proof of p\var{\over[typed]{\illegible{2}}{\unsure{ro}}}{\unsure{ro}}{typewriter overwriting}{ps_134-5_095-v0004}position. \var{\deleted{a) To prove that for F and G in $A^\wedge$,
\[ f_{!}(FxG) \quad f_{!}(F)xf_{!}(G) \]
is bijective, we are reduced, using \inserted[typed]{the} distributivity assumption in M, to
the case when F and G are 0-connected. \xout{xAs}
\[ \xout{A_{/FxG} \quad A_{/F} \quad A_{/G}} \]
\xout{we get}
\[ f_{!}(FxG) = \underset{A_{/F}xA_{/G}}{\mathrm{lim}} \quad . \]}}{}{deletion}{ps_134-5_095-v0005}

a) As we have, by definition of $f_{!}$ , for any F and G in $A^\wedge$
\[ f_{!}(F) = \underset{A_{/F}}{\underset{\hand{\longrightarrow}}{\mathrm{lim}}} f(a) \quad , \quad f_{!}(G) = \underset{A_{/G}}{\underset{\hand{\longrightarrow}}{\mathrm{lim}}} f(b) \quad , \]
deduced from the corresponding relations in $A^\wedge$
\[ F \var{\over{=}{\simeq}}{\hand{\simeq}}{writing over}{ps_134-5_095-v0006} \underset{A_{/F}}{\underset{\hand{\longrightarrow}}{\mathrm{lim}}}\ a \quad , \quad G \var{\over{=}{\simeq}}{\hand{\simeq}}{writing over}{ps_134-5_095-v0007} \underset{A_{/G}}{\underset{\hand{\longrightarrow}}{\mathrm{Lim}}}\ b \quad , \]
the looked for bijectivity of
\[ f_{!}(FxG) \hand{\to} f_{!}(F)xf_{!}(G) \]
will follow from the assumption\var{\inserted{(*)}}{\hand{(*)}}{insertion}{ps_134-5_095-v0008}, and from the following lemma\hand{,} \hand{(}applied
in both categories $A^\wedge$ and M \hand{)} \hand{the proof of which is immediate:}

\margin{95}

%% END SOURCE ps_134-5_095
\ifnum\EditionMode=1
\psapparatusstart
\psentry{ps_134-5_095-v0001}{1}{writing over}{$\over{\illegible{1}}{F}$}{$\hand{F}$}{The superposed handwritten reading replaces the underlying reading.}
\psentry{ps_134-5_095-v0002}{2}{insertion}{\inserted{b,c}}{\hand{b,c}}{The handwritten insertion is incorporated at its indicated position; insertion brackets leave the running text.}
\psentry{ps_134-5_095-v0003}{3}{typewriter deletion}{\xout{\unsure{it}}}{}{The typewriter strikes out this material; it leaves the corrected running text.}
\psentry{ps_134-5_095-v0004}{4}{typewriter overwriting}{\over[typed]{\illegible{2}}{\unsure{ro}}}{\unsure{ro}}{The superposed typed reading replaces the underlying reading.}
\psentry{ps_134-5_095-v0005}{5}{deletion}{\deleted{a) To prove that for F and G in $A^\wedge$,
\[ f_{!}(FxG) \quad f_{!}(F)xf_{!}(G) \]
is bijective, we are reduced, using \inserted[typed]{the} distributivity assumption in M, to
the case when F and G are 0-connected. \xout{xAs}
\[ \xout{A_{/FxG} \quad A_{/F} \quad A_{/G}} \]
\xout{we get}
\[ f_{!}(FxG) = \underset{A_{/F}xA_{/G}}{\mathrm{lim}} \quad . \]}}{}{The author cancels this material by hand; it leaves the corrected running text.}
\psentry{ps_134-5_095-v0006}{6}{writing over}{$\over{=}{\simeq}$}{$\hand{\simeq}$}{The superposed handwritten reading replaces the underlying reading.}
\psentry{ps_134-5_095-v0007}{7}{writing over}{$\over{=}{\simeq}$}{$\hand{\simeq}$}{The superposed handwritten reading replaces the underlying reading.}
\psentry{ps_134-5_095-v0008}{8}{insertion}{\inserted{(*)}}{\hand{(*)}}{The handwritten insertion is incorporated at its indicated position; insertion brackets leave the running text.}
\psapparatusend
\fi
\pssource{ps_134-5_096}{134-5/96}{https://grothendieck.umontpellier.fr/web/viewer.html?file=../134-5.pdf\#page=97}
\psnumber{ps_134-5_096-v0001}{1}
\psnumber{ps_134-5_096-v0002}{2}
\psnumber{ps_134-5_096-v0003}{3}
\psnumber{ps_134-5_096-v0004}{4}
\psnumber{ps_134-5_096-v0005}{5}
\psnumber{ps_134-5_096-v0006}{6}
\psnumber{ps_134-5_096-v0007}{7}
\psnumber{ps_134-5_096-v0008}{8}
\psnumber{ps_134-5_096-v0009}{9}
\psnumber{ps_134-5_096-v0010}{10}
\psnumber{ps_134-5_096-v0011}{11}
\psnumber{ps_134-5_096-v0012}{12}
\psnumber{ps_134-5_096-v0013}{13}
\psnumber{ps_134-5_096-v0014}{14}
%% BEGIN SOURCE ps_134-5_096
% Pursuing Stacks (cotes 134-2 to 134-8): archive cote 134-5, site page 96 (PDF page 97), typescript page 283.
% typescript. Draft compared: ps4.tex:4920-4986.
% AI-transcribed by Claude Opus 5.5 (claude-opus-5-5, adaptive thinking, effort high), Message Batches API, batch ps2 (msgbatch_01NVwr1Yuo8qJf5rykeGBXPK), 2026-11-12; reviewed against both supplied scan halves; see CHECK_DECISIONS.tsv.
\scanpage{134-5/96}
\mspage{283}

\hand{\uline{\textcolor{black}{Lemma.}}} Let M be a category stable under direct limits and under binary
products, and such that the product functors $y\hand{\mapsto} y\mathrm{x}x$ commute to di-%
rect limits. Then binary products are distributive with respedt to
direct limits, i.e. if
\[ u:\ I\hand{\to} M \quad , \quad v:\ J\hand{\to} M \]
are two functors of small categories with valuas in M, we have
\[ \underset{\hand{I\times J}}{\underset{\hand{\longrightarrow}}{\mathrm{lim}}}\ u(i)\mathrm{x}v(j)\ \hand{\xrightarrow{\sim}}\ (\underset{\hand{I}}{\underset{\hand{\longrightarrow}}{\mathrm{lim}}}\, u(i))\mathrm{x}(\underset{\hand{J}}{\underset{\hand{\longrightarrow}}{\mathrm{lim}}}\, v(j)) \quad . \]

Proof of b). Follows from a), applying it to the induced functor
\[ f_{/F}\ :\ A_{/F}\hand{\longrightarrow} M_{/f_{\hand{!}}(F)} \quad , \]
in order to prove commutation of $f_{\hand{!}}$ with a fibered product corresponding
to a diagram
% Hand-drawn arrows are marked blue; typed objects remain plain.
\diagram{
(* *) \qquad
\begin{tikzcd}[column sep=small]
G \ar[dr,color=blue] \& \& H \ar[dl,color=blue] \\
\& F \&
\end{tikzcd}
}
in $A^\wedge$.

\hand{\uline{\textcolor{black}{Remark}}}\hand{1}. Part b) does not look as nice as part a), as we would like to
be able to take F equally in A - what I first expected to get out. If
we make \var{\inserted{only}}{\hand{only}}{insertion}{ps_134-5_096-v0001} this weaker assumption, it will follow at once that $f_{\hand{!}}$ commutes
\margin{cf section 89}
to fibered products corresponding to diagrams (* *) \hand{\uline{\textcolor{black}{with}}} F \hand{\uline{\textcolor{black}{in}}} A. It does'%
nt seem finally that from this follows commutation of $f_{\hand{!}}$ to all fibered
products, unless making some extra assumptions, such as $f_{\hand{!}}$ transforming
monomorphisms into monomorphisms (which is a necessary condition anyhow).
We'll have to come back upon this - for the time being we need only to \var{\xout{de}}{}{typewriter deletion}{ps_134-5_096-v0002}
deal with products, for which a) is adequate.

The ``second step'' now is a tautology\hand{;} \var{\deleted{\hand{\unsure{using} prop. 1}}}{}{deletion}{ps_134-5_096-v0003}

\hand{\uline{\textcolor{black}{Proposition}}} 2. Under the preliminary assumptions of prop.1 above, assume
\var{\inserted[typed]{\deleted{and \xout{ff(A) generates} and f(A) generates M by binary}}}{}{typed insertion}{ps_134-5_096-v0004}
moreover that f is fully faithful\hand{,} \var{\deleted{Then $f_{\hand{!}}$ commutes to products and
fibered products (in $A^\wedge$ ) of elements \hand{\uline{\textcolor{black}{in}}} A .

We may assume f to be the inclusion functor of a full subcategory
A in M. The assumption on}}{}{deletion}{ps_134-5_096-v0005} and that f(A) \hand{``}generates M by \var{\xout{s}}{}{typewriter deletion}{ps_134-5_096-v0006}strict epimor-%
phisms\hand{,''} i.e. that for any object P in M, we have
\[ \hand{(***)}\qquad P\ \hand{\xleftarrow{\sim}}\ \underset{A_{/P}}{\underset{\hand{\longrightarrow}}{\mathrm{lim}}}\ f(x) \qquad . \]
Then $f_{\hand{!}}$ commutes to \var{\inserted{finite}}{\hand{finite}}{insertion}{ps_134-5_096-v0007} \var{\deleted{binary}}{}{deletion}{ps_134-5_096-v0008} products\hand{,} \var{\inserted{in $A^\wedge$ of objects of $A$,}}{\hand{in $A^\wedge$ of objects of $A$,}}{insertion}{ps_134-5_096-v0009}% This pen insertion is enclosed by a loop below the line and linked by a locator to products.
 and \var{\inserted{to}}{\hand{to}}{insertion}{ps_134-5_096-v0010} fibered products \var{\inserted{in $A^\wedge$}}{\hand{in $A^\wedge$}}{insertion}{ps_134-5_096-v0011} of objects \var{\deleted{\unsure{under}}}{}{deletion}{ps_134-5_096-v0012} \var{\inserted{of A}}{\hand{of A}}{insertion}{ps_134-5_096-v0013} .
\var{\deleted{\hand{an object a in $A^\wedge$}}}{}{deletion}{ps_134-5_096-v0014}.

\margin{96}

%% END SOURCE ps_134-5_096
\ifnum\EditionMode=1
\psapparatusstart
\psentry{ps_134-5_096-v0001}{1}{insertion}{\inserted{only}}{\hand{only}}{The handwritten insertion is incorporated at its indicated position; insertion brackets leave the running text.}
\psentry{ps_134-5_096-v0002}{2}{typewriter deletion}{\xout{de}}{}{The typewriter strikes out this material; it leaves the corrected running text.}
\psentry{ps_134-5_096-v0003}{3}{deletion}{\deleted{\hand{\unsure{using} prop. 1}}}{}{The author cancels this material by hand; it leaves the corrected running text.}
\psentry{ps_134-5_096-v0004}{4}{typed insertion}{\inserted[typed]{\deleted{and \xout{ff(A) generates} and f(A) generates M by binary}}}{}{The typed insertion is incorporated at its indicated position; insertion brackets leave the running text.}
\psentry{ps_134-5_096-v0005}{5}{deletion}{\deleted{Then $f_{\hand{!}}$ commutes to products and
fibered products (in $A^\wedge$ ) of elements \hand{\uline{\textcolor{black}{in}}} A .

We may assume f to be the inclusion functor of a full subcategory
A in M. The assumption on}}{}{The author cancels this material by hand; it leaves the corrected running text.}
\psentry{ps_134-5_096-v0006}{6}{typewriter deletion}{\xout{s}}{}{The typewriter strikes out this material; it leaves the corrected running text.}
\psentry{ps_134-5_096-v0007}{7}{insertion}{\inserted{finite}}{\hand{finite}}{The handwritten insertion is incorporated at its indicated position; insertion brackets leave the running text.}
\psentry{ps_134-5_096-v0008}{8}{deletion}{\deleted{binary}}{}{The author cancels this material by hand; it leaves the corrected running text.}
\psentry{ps_134-5_096-v0009}{9}{insertion}{\inserted{in $A^\wedge$ of objects of $A$,}}{\hand{in $A^\wedge$ of objects of $A$,}}{The handwritten insertion is incorporated at its indicated position; insertion brackets leave the running text.}
\psentry{ps_134-5_096-v0010}{10}{insertion}{\inserted{to}}{\hand{to}}{The handwritten insertion is incorporated at its indicated position; insertion brackets leave the running text.}
\psentry{ps_134-5_096-v0011}{11}{insertion}{\inserted{in $A^\wedge$}}{\hand{in $A^\wedge$}}{The handwritten insertion is incorporated at its indicated position; insertion brackets leave the running text.}
\psentry{ps_134-5_096-v0012}{12}{deletion}{\deleted{\unsure{under}}}{}{The author cancels this material by hand; it leaves the corrected running text.}
\psentry{ps_134-5_096-v0013}{13}{insertion}{\inserted{of A}}{\hand{of A}}{The handwritten insertion is incorporated at its indicated position; insertion brackets leave the running text.}
\psentry{ps_134-5_096-v0014}{14}{deletion}{\deleted{\hand{an object a in $A^\wedge$}}}{}{The author cancels this material by hand; it leaves the corrected running text.}
\psapparatusend
\fi
\pssource{ps_134-5_097}{134-5/97}{https://grothendieck.umontpellier.fr/web/viewer.html?file=../134-5.pdf\#page=98}
\psnumber{ps_134-5_097-v0001}{1}
\psnumber{ps_134-5_097-v0002}{2}
\psnumber{ps_134-5_097-v0003}{3}
\psnumber{ps_134-5_097-v0004}{4}
\psnumber{ps_134-5_097-v0005}{5}
%% BEGIN SOURCE ps_134-5_097
% Pursuing Stacks (cotes 134-2 to 134-8): archive cote 134-5, site page 97 (PDF page 98), typescript page 284.
% typescript. Draft compared: ps4.tex:4971-5035.
% AI-transcribed by Claude Opus 5.5 (claude-opus-5-5, adaptive thinking, effort high), Message Batches API, batch ps2 (msgbatch_01NVwr1Yuo8qJf5rykeGBXPK), 2026-11-12; reviewed against both supplied scan halves; see CHECK_DECISIONS.tsv.
\scanpage{134-5/97}
\mspage{284}

\var{\xout{Pw}}{}{typewriter deletion}{ps_134-5_097-v0001} \hand{\uline{\textcolor{black}{Proof}}}. That $f_{!}$ transforms final object into final object follows
from the formula $(* * *)$ above, by taking $P$ to be the final object of $M$.
To get commutation to binary products, we apply the formula for $P=ax_{M}b$,
and notice that
\[ A_{/P} \hand{\simeq} A_{/P'} \quad , \quad \text{with} \quad P' = ax_{A^{\wedge}}b \quad , \]
because the inclusion $f$: $A \hand{\hookrightarrow} M$ is full (I forgot to say we may assume $f$
to be the inclusion fu\unsure{nc}tor of a full subcategory). But we have too
\[ f_{!}(P') = \underset{A_{/P'}}{\underset{\hand{\longrightarrow}}{\mathrm{lim}}}\ f(x) \qquad , \]
\margin{\tikz[x=1em,y=1em,baseline=-.4ex]{\draw (1,1.4) .. controls (1.15,.75) and (.7,1.2) .. (.35,.9) .. controls (-.25,.15) and (.2,-.1) .. (.8,-.1) .. controls (1,-.1) and (.85,-.35) .. (.85,-.55);}}hence $f_{!}(P') \hand{\simeq} P$ \hand{,} \var{\xout{and.}}{}{typewriter deletion}{ps_134-5_097-v0002} The proof for fibered products is similar, or
can be reduced to the case of products by considering the induced functor
\[ f_{/a} : A_{\hand{/}a} \hand{\to} M_{/a} \quad . \]

\hand{\uline{\textcolor{black}{Remark}}}\hand{2}.The various variants of the notion of a generating subcategory $A$
of a category $M$ have been dealt with in some detail in SGA 4 vol 1 exp I
par. 7 (p.45-60) (Springer LN 269) , cf more specifically prop. 7.2 on
page 47, summarizing the main relationships. The str\var{\over[typed]{\illegible{1}}{o}}{o}{typewriter overwriting}{ps_134-5_097-v0003}ngest of all noti-%
% The i/d impressions in consid ered overlap; the visible word gap is retained.
ons consid ered there is generation \var{\over[typed]{p}{b}}{b}{typewriter overwriting}{ps_134-5_097-v0004}y strict epimorphisms, which can be
expressed by the formula $(* * *)$ above, and is also equivalent to the
\hand{\uline{\textcolor{black}{functor}}} $f^{*}$ \hand{\uline{\textcolor{black}{being fully faithful}}} .

% A dark angular impression lies over the first e of immediate; no distinct earlier letter is recoverable.
We come back now to the case $M$=(Cat) - it is immediate that the
preliminary exactness conditions on $M$ of prop.1 are satisfied. Thus,
putting together prop. 1 and 2, we get:

\hand{\uline{\textcolor{black}{Proposition}}} 3. Let $A$ be a full subcategory of (Cat), generat ing (Cat)
by strict epimorphisms, $i$: $A \hand{\hookrightarrow}$(Cat) the inclusion functor. Then the
corresponding functor
\[ i_{!}: A^{\wedge} \hand{\longrightarrow} \text{(Cat)} \]
commutes to finite products, and also to fibered products in $A^{\wedge}$ over an
objedt of $A^{\wedge}$ coming from $A$.

% The h in The has a displaced typewriter impression above it.
The most familiar case when this applies is indeed the case of the
canonical inclusion of \hand{$\Delta$} into (Cat), as it is well-known that the cor-%
responding $i^{*}$ , namely the nerve functor, is fully faithful. From
this follows that \hand{\uline{\textcolor{black}{a fortiori}}}, for any full subcategory $A$ of (Cat) con-%
\margin{hasty!}taining \hand{$\Delta$} , $A$ generates \var{\xout{by\illegible{5}}}{}{typewriter deletion}{ps_134-5_097-v0005} (Cat) by strict epimorphisms, and hence
proposition 3 applies.

\hand{\uline{\textcolor{black}{Remark}}} 3. I doubt however, even for $A$ = \hand{$\Delta$}, the most typical case,

\margin{97}

%% END SOURCE ps_134-5_097
\ifnum\EditionMode=1
\psapparatusstart
\psentry{ps_134-5_097-v0001}{1}{typewriter deletion}{\xout{Pw}}{}{The typewriter strikes out this material; it leaves the corrected running text.}
\psentry{ps_134-5_097-v0002}{2}{typewriter deletion}{\xout{and.}}{}{The typewriter strikes out this material; it leaves the corrected running text.}
\psentry{ps_134-5_097-v0003}{3}{typewriter overwriting}{\over[typed]{\illegible{1}}{o}}{o}{The superposed typed reading replaces the underlying reading.}
\psentry{ps_134-5_097-v0004}{4}{typewriter overwriting}{\over[typed]{p}{b}}{b}{The superposed typed reading replaces the underlying reading.}
\psentry{ps_134-5_097-v0005}{5}{typewriter deletion}{\xout{by\illegible{5}}}{}{The typewriter strikes out this material; it leaves the corrected running text.}
\psapparatusend
\fi
\pssource{ps_134-5_098}{134-5/98}{https://grothendieck.umontpellier.fr/web/viewer.html?file=../134-5.pdf\#page=99}
\psnumber{ps_134-5_098-v0001}{1}
\psnumber{ps_134-5_098-v0002}{2}
\psnumber{ps_134-5_098-v0003}{3}
\psnumber{ps_134-5_098-v0004}{4}
\psnumber{ps_134-5_098-v0005}{5}
\psnumber{ps_134-5_098-v0006}{6}
\psnumber{ps_134-5_098-v0007}{7}
\psnumber{ps_134-5_098-v0008}{8}
\psnumber{ps_134-5_098-v0009}{9}
\psnumber{ps_134-5_098-v0010}{10}
\psnumber{ps_134-5_098-v0011}{11}
\psnumber{ps_134-5_098-v0012}{12}
\psnumber{ps_134-5_098-v0013}{13}
\psnumber{ps_134-5_098-v0014}{14}
\psnumber{ps_134-5_098-v0015}{15}
\psnumber{ps_134-5_098-v0016}{16}
\psnumber{ps_134-5_098-v0017}{17}
%% BEGIN SOURCE ps_134-5_098
% Pursuing Stacks (cotes 134-2 to 134-8): archive cote 134-5, site page 98 (PDF page 99), typescript page 285.
% typescript. Draft compared: ps4.tex:5020-5090.
% AI-transcribed by Claude Opus 5.5 (claude-opus-5-5, adaptive thinking, effort high), Message Batches API, batch ps2 (msgbatch_01NVwr1Yuo8qJf5rykeGBXPK), 2026-11-12; reviewed against both supplied scan halves; see CHECK_DECISIONS.tsv.
\scanpage{134-5/98}
\mspage{285}
\margin{see below -- $i_!$ indeed is \uline{not} left exact!}
that $i_!$ is even left exact,  because it looks unlikely to me, in view
of the explicit description of $i_!(K.)$ for a given ss complex K. , in
terms of ``generators and relations'' (cf. Gabriel-Zisman's book),  that
$i_!$ should transform monomorphisms into monomorphisms.  For what follows
this is \unsure{irrelevant} anyhow.

8\var{\over{5}{6}}{\hand{6}}{writing over}{ps_134-5_098-v0001}.  After these somewhat painful preliminaries,  it seems to me that
firm ground is in sight at last !  The main feeling that finally comes
out of these reflections,  is that we have a very good hold on checking
whether or not a functor $f_! : A^\wedge \hand{\to} M$ commutes to finite products,  and
moreover,  that this condition is satisfied in many ``good'' cases,  \var{\deleted{and}}{}{deletion}{ps_134-5_098-v0002} \var{\inserted{in}}{\hand{in}}{insertion}{ps_134-5_098-v0003}
more cases at any rate than I suspected.  In case of the inclusion
functor of \hand{$\Delta$} into (Cat),  more generally of any \var{\inserted[typed]{full}}{full}{typed insertion}{ps_134-5_098-v0004} subcategory \var{\xout{of}}{}{typewriter deletion}{ps_134-5_098-v0005} A
of (Cat) consisting of contractible objects,  containing $\hand{\Delta}_1$ ,  and
large enough in order to generate (Cat) by strict epimorphisms (or, as
we'll say,  to \hand{\uline{\textcolor{black}{generate strictly}}} the category (Cat)),  it follows that
we have the inclusion
\[ (1) \qquad i_!(A^\wedge_c) \hand{\subset} (\mathrm{Cat})_c \quad , \]
i.e.  $i_!$ transforms \hand{\uline{\textcolor{black}{contractible}}} objects into contractible objects,
or,  equivalently,  is a \hand{\uline{\textcolor{black}{morphism for the}}} \var{\xout{corresponding}}{}{typewriter deletion}{ps_134-5_098-v0006} \hand{\uline{\textcolor{black}{homotopy struc-}}}%
\hand{\uline{\textcolor{black}{tures}}} on $A^\wedge$ and (Cat).  (Note that the two first assumptions on the
subcategory A of (Cat)  imply that the \var{\xout{canonical}}{}{typewriter deletion}{ps_134-5_098-v0007} homotopy structure
on $A^\wedge$ generated by the ``intervals'' made up with elements of A,  is a
contractibility structure,  admitting A as a generating set of contrac-%
tible objects.  The situation is really nice only when assuming that $A^\wedge$
is totally aspheric though,  which will imply that A\var{\xout{$^\wedge$}}{}{typewriter deletion}{ps_134-5_098-v0008} is a \var{\xout{str}}{}{typewriter deletion}{ps_134-5_098-v0009} test
category (even a strict one) and that the $\underline{W}$-asphericity structure \var{\inserted[typed]{on $A^\wedge$}}{on $A^\wedge$}{typed insertion}{ps_134-5_098-v0010} asso-%
ciated to the contractibility structure just described is the usual one
(cf prop. 2 b)  p. 248,  where we take C to be A).)

The \var{\inserted[typed]{inclusion}}{inclusion}{typed insertion}{ps_134-5_098-v0011} \var{\xout{relation}}{}{typewriter deletion}{ps_134-5_098-v0012} (1) of course strongly resembles the problematic
inclusion (4) (p.277) of yesterdays notes,  namely\hand{,} \var{\xout{to}}{}{typewriter deletion}{ps_134-5_098-v0013} here to $i_!$ taking
aspheric objects into aspheric ones.  I doubt though this \var{\inserted{latter}}{\hand{latter}}{insertion}{ps_134-5_098-v0014} is true even
for the standard case A = \hand{$\Delta$}\hand{?}  But before trying at all costs to see
whether it holds or not,  the point I wish to make is that for the pur-%
pose we have in view,  namely defining \var{\inserted[typed]{suitable}}{suitable}{typed insertion}{ps_134-5_098-v0015} conditions,  stable under compo-%
sition,  \var{\xout{between adjoint}}{}{typewriter deletion}{ps_134-5_098-v0016} on a pair $(f_!,f^*)$ of adjoint functors between
two asphericity structures M and M' ,  \var{\inserted{-- for this purpose}}{\hand{-- for this purpose}}{insertion}{ps_134-5_098-v0017} a relation of the type  (1),
namely

\margin{98}

%% END SOURCE ps_134-5_098
\ifnum\EditionMode=1
\psapparatusstart
\psentry{ps_134-5_098-v0001}{1}{writing over}{\over{5}{6}}{\hand{6}}{The superposed handwritten reading replaces the underlying reading.}
\psentry{ps_134-5_098-v0002}{2}{deletion}{\deleted{and}}{}{The author cancels this material by hand; it leaves the corrected running text.}
\psentry{ps_134-5_098-v0003}{3}{insertion}{\inserted{in}}{\hand{in}}{The handwritten insertion is incorporated at its indicated position; insertion brackets leave the running text.}
\psentry{ps_134-5_098-v0004}{4}{typed insertion}{\inserted[typed]{full}}{full}{The typed insertion is incorporated at its indicated position; insertion brackets leave the running text.}
\psentry{ps_134-5_098-v0005}{5}{typewriter deletion}{\xout{of}}{}{The typewriter strikes out this material; it leaves the corrected running text.}
\psentry{ps_134-5_098-v0006}{6}{typewriter deletion}{\xout{corresponding}}{}{The typewriter strikes out this material; it leaves the corrected running text.}
\psentry{ps_134-5_098-v0007}{7}{typewriter deletion}{\xout{canonical}}{}{The typewriter strikes out this material; it leaves the corrected running text.}
\psentry{ps_134-5_098-v0008}{8}{typewriter deletion}{\xout{$^\wedge$}}{}{The typewriter strikes out this material; it leaves the corrected running text.}
\psentry{ps_134-5_098-v0009}{9}{typewriter deletion}{\xout{str}}{}{The typewriter strikes out this material; it leaves the corrected running text.}
\psentry{ps_134-5_098-v0010}{10}{typed insertion}{\inserted[typed]{on $A^\wedge$}}{on $A^\wedge$}{The typed insertion is incorporated at its indicated position; insertion brackets leave the running text.}
\psentry{ps_134-5_098-v0011}{11}{typed insertion}{\inserted[typed]{inclusion}}{inclusion}{The typed insertion is incorporated at its indicated position; insertion brackets leave the running text.}
\psentry{ps_134-5_098-v0012}{12}{typewriter deletion}{\xout{relation}}{}{The typewriter strikes out this material; it leaves the corrected running text.}
\psentry{ps_134-5_098-v0013}{13}{typewriter deletion}{\xout{to}}{}{The typewriter strikes out this material; it leaves the corrected running text.}
\psentry{ps_134-5_098-v0014}{14}{insertion}{\inserted{latter}}{\hand{latter}}{The handwritten insertion is incorporated at its indicated position; insertion brackets leave the running text.}
\psentry{ps_134-5_098-v0015}{15}{typed insertion}{\inserted[typed]{suitable}}{suitable}{The typed insertion is incorporated at its indicated position; insertion brackets leave the running text.}
\psentry{ps_134-5_098-v0016}{16}{typewriter deletion}{\xout{between adjoint}}{}{The typewriter strikes out this material; it leaves the corrected running text.}
\psentry{ps_134-5_098-v0017}{17}{insertion}{\inserted{-- for this purpose}}{\hand{-- for this purpose}}{The handwritten insertion is incorporated at its indicated position; insertion brackets leave the running text.}
\psapparatusend
\fi
\pssource{ps_134-5_099}{134-5/99}{https://grothendieck.umontpellier.fr/web/viewer.html?file=../134-5.pdf\#page=100}
\psnumber{ps_134-5_099-v0001}{1}
\psnumber{ps_134-5_099-v0002}{2}
\psnumber{ps_134-5_099-v0003}{3}
\psnumber{ps_134-5_099-v0004}{4}
\psnumber{ps_134-5_099-v0005}{5}
\psnumber{ps_134-5_099-v0006}{6}
\psnumber{ps_134-5_099-v0007}{7}
\psnumber{ps_134-5_099-v0008}{8}
\psnumber{ps_134-5_099-v0009}{9}
\psnumber{ps_134-5_099-v0010}{10}
\psnumber{ps_134-5_099-v0011}{11}
%% BEGIN SOURCE ps_134-5_099
% Pursuing Stacks (cotes 134-2 to 134-8): archive cote 134-5, site page 99 (PDF page 100), typescript page 286.
% typescript. Draft compared: ps4.tex:5075-5141.
% AI-transcribed by Claude Opus 5.5 (claude-opus-5-5, adaptive thinking, effort high), Message Batches API, batch ps2 (msgbatch_01NVwr1Yuo8qJf5rykeGBXPK), 2026-11-12; reviewed against both supplied scan halves; see CHECK_DECISIONS.tsv.
\scanpage{134-5/99}
\mspage{286}

(2) \[ f_!(M'_c) \hand{\subset} M_c \quad , \]
is just as good as the relation (4) on p.277\hand{,} provided we make the evi-%
dent assumption needed for (2) to make a sense,  namely \hand{\uline{\textcolor{black}{the given asphe-}}}%
\hand{\uline{\textcolor{black}{ricity structures on}}} M,M' \hand{\uline{\textcolor{black}{to be associated to contractibility structures}}}.
This condition (2) is \hand{(}considerably\hand{!)} weaker than
% The a near the prime in M'_a is repeated at a higher typing position; the clean subscript a remains below.
(2') \[ f_!(M'_a) \hand{\subset} M_a \quad , \]
due to the inclusion $M'_c \hand{\subset} M'_a$ % A small u-like impression sits below the first letter of and; it is not read as an additional word.
and to the fact that (with the notations
and using the results of section 81)  the inclusion (2) is implied by
the apparently weaker inclusion (under the assumption  Loc 4)) :
\[ f_!(M'_c) \hand{\subset} M_o \quad , \]
where  $M_c \hand{\subset} M_a \hand{\subset} M_o$ .  On the other hand,  (2)  \hand{\uline{\textcolor{black}{is evidently stable}}}
\hand{\uline{\textcolor{black}{under composition}}}.

To be specific\var{\inserted{about ``as good as'',}}{\hand{about ``as good as'',}}{insertion}{ps_134-5_099-v0001} let's come back to the statement of the ``pretty''
proposition on p.275,  somewhat mared by the ``awkward'' condition ``Awk''
we had to throw in, \var{\inserted{which looks so}}{\hand{which looks so}}{insertion}{ps_134-5_099-v0002} ugly because of it's lack of stability under com-%
positions.  Let's replace this condition by the slightly stronger one
(2),  which is stable under composition.  Condition (2) does imply
(Awk) indeed,  as we see by taking for A a small full subcategory of
$M_c$ ,  generating the \var{\xout{\unsure{contractibility}}}{}{typewriter deletion}{ps_134-5_099-v0003}\var{\inserted[typed]{asphericity}}{asphericity}{typed insertion}{ps_134-5_099-v0004} structure of M,  and for \var{\over[typed]{C}{B}}{B}{typewriter overwriting}{ps_134-5_099-v0005} any
full subcategory of M contained in $M_a$ and big enough in order  a) to
contain $f_!(A)$  and  b)  to generate the asphericity structure of M.
We may take for instance for A any subcategory of $M'_c$  closed under
finite products in M' and generating the contractibility structure,
and accordingly for B,  so that the pair (A,B) will match for any
choice of the basic localizer $\underline{W}$ .  Thus the \var{\inserted[typed]{six}}{six}{typed insertion}{ps_134-5_099-v0006} conditions \var{\xout{of}}{}{typewriter deletion}{ps_134-5_099-v0007} (i) to (i'')
of the proposition are equivalent,  provided in the last (ii'') we assu-%
me moreover that B contain a final object $e_M$ of M.  Moreover,  it is
clear now that these conditions can be viewed also as a property of the
functor
(3) \[ f_{!\,c} \ : \ M'_c \hand{\longrightarrow} M_c \]
induce\var{\over[typed]{\unsure{e}}{d}}{d}{typewriter overwriting}{ps_134-5_099-v0008} by $f_!$ , \var{\xout{and}}{}{typewriter deletion}{ps_134-5_099-v0009} which we may \var{\over[typed]{p}{a}}{a}{typewriter overwriting}{ps_134-5_099-v0010}dequately express by saying that this
functor (3)  is $\underline{W}$-\hand{\uline{\textcolor{black}{aspheric}}} (where the two sides of (3) are categories
which \var{\over{may}{need}}{\hand{need}}{writing over}{ps_134-5_099-v0011} not be small,  but which are at any rate $\underline{W}$-\hand{\uline{\textcolor{black}{aspherators}}}  in
the sense of section 79, p.247).

\margin{99}

%% END SOURCE ps_134-5_099
\ifnum\EditionMode=1
\psapparatusstart
\psentry{ps_134-5_099-v0001}{1}{insertion}{\inserted{about ``as good as'',}}{\hand{about ``as good as'',}}{The handwritten insertion is incorporated at its indicated position; insertion brackets leave the running text.}
\psentry{ps_134-5_099-v0002}{2}{insertion}{\inserted{which looks so}}{\hand{which looks so}}{The handwritten insertion is incorporated at its indicated position; insertion brackets leave the running text.}
\psentry{ps_134-5_099-v0003}{3}{typewriter deletion}{\xout{\unsure{contractibility}}}{}{The typewriter strikes out this material; it leaves the corrected running text.}
\psentry{ps_134-5_099-v0004}{4}{typed insertion}{\inserted[typed]{asphericity}}{asphericity}{The typed insertion is incorporated at its indicated position; insertion brackets leave the running text.}
\psentry{ps_134-5_099-v0005}{5}{typewriter overwriting}{\over[typed]{C}{B}}{B}{The superposed typed reading replaces the underlying reading.}
\psentry{ps_134-5_099-v0006}{6}{typed insertion}{\inserted[typed]{six}}{six}{The typed insertion is incorporated at its indicated position; insertion brackets leave the running text.}
\psentry{ps_134-5_099-v0007}{7}{typewriter deletion}{\xout{of}}{}{The typewriter strikes out this material; it leaves the corrected running text.}
\psentry{ps_134-5_099-v0008}{8}{typewriter overwriting}{\over[typed]{\unsure{e}}{d}}{d}{The superposed typed reading replaces the underlying reading.}
\psentry{ps_134-5_099-v0009}{9}{typewriter deletion}{\xout{and}}{}{The typewriter strikes out this material; it leaves the corrected running text.}
\psentry{ps_134-5_099-v0010}{10}{typewriter overwriting}{\over[typed]{p}{a}}{a}{The superposed typed reading replaces the underlying reading.}
\psentry{ps_134-5_099-v0011}{11}{writing over}{\over{may}{need}}{\hand{need}}{The superposed handwritten reading replaces the underlying reading.}
\psapparatusend
\fi
\pssource{ps_134-5_100}{134-5/100}{https://grothendieck.umontpellier.fr/web/viewer.html?file=../134-5.pdf\#page=101}
\psnumber{ps_134-5_100-v0001}{1}
\psnumber{ps_134-5_100-v0002}{2}
\psnumber{ps_134-5_100-v0003}{3}
\psnumber{ps_134-5_100-v0004}{4}
\psnumber{ps_134-5_100-v0005}{5}
\psnumber{ps_134-5_100-v0006}{6}
\psnumber{ps_134-5_100-v0007}{7}
\psnumber{ps_134-5_100-v0008}{8}
\psnumber{ps_134-5_100-v0009}{9}
\psnumber{ps_134-5_100-v0010}{10}
\psnumber{ps_134-5_100-v0011}{11}
\psnumber{ps_134-5_100-v0012}{12}
%% BEGIN SOURCE ps_134-5_100
% Pursuing Stacks (cotes 134-2 to 134-8): archive cote 134-5, site page 100 (PDF page 101), typescript page 287.
% typescript. Draft compared: ps4.tex:5123-5202.
% AI-transcribed by Claude Opus 5.5 (claude-opus-5-5, adaptive thinking, effort high), Message Batches API, batch ps2 (msgbatch_01NVwr1Yuo8qJf5rykeGBXPK), 2026-11-12; reviewed against both supplied scan halves; see CHECK_DECISIONS.tsv.
\scanpage{134-5/100}
\mspage{287}

Stated this way, the condition just obtained on the pair of adjoint
functors $(f_!,f^*)$\hand{,} between the two \var{\xout{contractibility structures (}}{}{typewriter deletion}{ps_134-5_100-v0001} categories
M, M' endowed with contractibility structures, namely (2) and
\[ (4) \qquad f^*(M_a) \hand{\subset} M'_a \quad , \]
are not quite symmetric, as the condition on $f_!$ is expressed in terms
of contractible objects, whereas the condition on $f^*$ is expressed in
terms of aspheric \var{\deleted{objects}}{}{deletion}{ps_134-5_100-v0002} \var{\inserted{ones}}{\hand{ones}}{insertion}{ps_134-5_100-v0003}. % One continuous oblique pen stroke runs from assuming/that toward the gap after (4) on the next line.
However, assuming\hand{/}that $\underline{W}$ satisfies Loc 4)
and using theorem 1 p. 252, we see that (4) is equivalent to the
condition
\[ (5) \qquad f^*(M_c) \hand{\subset} M'_c \quad , \]
which does not depend any longer on the choice of $\underline{W}$ !

Finally, we are led to a notion of pure homotopy theory, in
terms of contractibility structures alone, without the intr\var{\xout{o}}{}{typewriter deletion}{ps_134-5_100-v0004}usion of
the choice of a basic localizer $\underline{W}$ and corresponding asphericity notions\hand{.}
We may call \var{\xout{it}}{}{typewriter deletion}{ps_134-5_100-v0005} the pair a ``\hand{\uline{\textcolor{black}{bimorphism}}}'' of contractibility structures,
and introduce it via the following \hand{summing-up statement:}

\var{\inserted{``Scholie''}}{\hand{``Scholie''}}{insertion}{ps_134-5_100-v0006} \var{\deleted{\hand{\uline{\textcolor{black}{Proposition}}}}}{}{deletion}{ps_134-5_100-v0007} : Let $(M,M_c)$ and $(M',M'_c)$ be two contractibility structure,
each admitting a small generating subcategory for the contractibility
structures, say C and C' respectively. Let $(f_!,f^*)$ be a pair of ad-%
joint functors
\[ f_! : M' \hand{\longrightarrow} M \quad , \qquad f^* : M \hand{\longrightarrow} M' \quad . \]
Let's consider the following inclusion conditions
\[ (!) \ \var{\xout{(!\illegible{1}}}{}{typewriter deletion}{ps_134-5_100-v0008} \qquad f_!(M'_c) \hand{\subset} M_c \quad , \]
\margin{(!')}
\[ (!\hand{'}) \qquad f_!(C') \hand{\subset} M_o \quad , \]
\[ (*) \qquad f^*(M_c) \hand{\subset} M'_c \quad , \]
\[ (*') \qquad f^*(C) \hand{\subset} M'_o \quad , \]
where $M_o$ and $M'_o$ are (as in section 81) the sets of ``0-connected''
objects in M and M' respectively, for the given contractibility struc-%
tures.
% A pen step/caret marks the new paragraph boundary before a).
\par\hand{a)} \var{\deleted{(}}{}{deletion}{ps_134-5_100-v0009}Clearly, in view of
\[ C \hand{\subset} M_c \hand{\subset} M_o \quad , \quad C' \hand{\subset} M'_c \hand{\subset} M'_o \quad , \]
condition (!) implies (!') and condition (*) implies (*')\var{\deleted{)}}{}{deletion}{ps_134-5_100-v0010}. Moreover,
it is even true that conditions (*) and (*') are equivalent, and the
same holds for (!) and (!') provided $f_!$ commutes with finite pro\var{\xout{\unsure{ducts}}}{}{typewriter deletion}{ps_134-5_100-v0011}\hand{-}%
ducts. We'll say that $(f_!\ ,\ f^*\ )$ is a \hand{\uline{\textcolor{black}{bimorphism}}} for the given contrac-%
tibility structures, if \var{\inserted{the}}{\hand{the}}{insertion}{ps_134-5_100-v0012} inclusions (!) and (*) do hold.

\margin{100}

%% END SOURCE ps_134-5_100
\ifnum\EditionMode=1
\psapparatusstart
\psentry{ps_134-5_100-v0001}{1}{typewriter deletion}{\xout{contractibility structures (}}{}{The typewriter strikes out this material; it leaves the corrected running text.}
\psentry{ps_134-5_100-v0002}{2}{deletion}{\deleted{objects}}{}{The author cancels this material by hand; it leaves the corrected running text.}
\psentry{ps_134-5_100-v0003}{3}{insertion}{\inserted{ones}}{\hand{ones}}{The handwritten insertion is incorporated at its indicated position; insertion brackets leave the running text.}
\psentry{ps_134-5_100-v0004}{4}{typewriter deletion}{\xout{o}}{}{The typewriter strikes out this material; it leaves the corrected running text.}
\psentry{ps_134-5_100-v0005}{5}{typewriter deletion}{\xout{it}}{}{The typewriter strikes out this material; it leaves the corrected running text.}
\psentry{ps_134-5_100-v0006}{6}{insertion}{\inserted{``Scholie''}}{\hand{``Scholie''}}{The handwritten insertion is incorporated at its indicated position; insertion brackets leave the running text.}
\psentry{ps_134-5_100-v0007}{7}{deletion}{\deleted{\hand{\uline{\textcolor{black}{Proposition}}}}}{}{The author cancels this material by hand; it leaves the corrected running text.}
\psentry{ps_134-5_100-v0008}{8}{typewriter deletion}{$\xout{(!\illegible{1}}$}{}{The typewriter strikes out this material; it leaves the corrected running text.}
\psentry{ps_134-5_100-v0009}{9}{deletion}{\deleted{(}}{}{The author cancels this material by hand; it leaves the corrected running text.}
\psentry{ps_134-5_100-v0010}{10}{deletion}{\deleted{)}}{}{The author cancels this material by hand; it leaves the corrected running text.}
\psentry{ps_134-5_100-v0011}{11}{typewriter deletion}{\xout{\unsure{ducts}}}{}{The typewriter strikes out this material; it leaves the corrected running text.}
\psentry{ps_134-5_100-v0012}{12}{insertion}{\inserted{the}}{\hand{the}}{The handwritten insertion is incorporated at its indicated position; insertion brackets leave the running text.}
\psapparatusend
\fi
\pssource{ps_134-5_101}{134-5/101}{https://grothendieck.umontpellier.fr/web/viewer.html?file=../134-5.pdf\#page=102}
\psnumber{ps_134-5_101-v0001}{1}
\psnumber{ps_134-5_101-v0002}{2}
\psnumber{ps_134-5_101-v0003}{3}
\psnumber{ps_134-5_101-v0004}{4}
\psnumber{ps_134-5_101-v0005}{5}
\psnumber{ps_134-5_101-v0006}{6}
\psnumber{ps_134-5_101-v0007}{7}
\psnumber{ps_134-5_101-v0008}{8}
\psnumber{ps_134-5_101-v0009}{9}
\psnumber{ps_134-5_101-v0010}{10}
\psnumber{ps_134-5_101-v0011}{11}
\psnumber{ps_134-5_101-v0012}{12}
\psnumber{ps_134-5_101-v0013}{13}
\psnumber{ps_134-5_101-v0014}{14}
\psnumber{ps_134-5_101-v0015}{15}
%% BEGIN SOURCE ps_134-5_101
% Pursuing Stacks (cotes 134-2 to 134-8): archive cote 134-5, site page 101 (PDF page 102), typescript page 288.
% typescript. Draft compared: ps4.tex:5189-5310.
% AI-transcribed by Claude Opus 5.5 (claude-opus-5-5, adaptive thinking, effort high), Message Batches API, batch ps2 (msgbatch_01NVwr1Yuo8qJf5rykeGBXPK), 2026-11-12; reviewed against both supplied scan halves; see CHECK_DECISIONS.tsv.
\scanpage{134-5/101}
\mspage{288}

b) Assume we got a bimorphism $(f_!,f^*)$ , i.e. (!) and ($*$) above
hold. Let $\underline{W}$ be a basic localizer, hence \var{\xout{\illegible{1}}}{}{typewriter deletion}{ps_134-5_101-v0001} \var{\inserted[typed]{the}}{the}{typed insertion}{ps_134-5_101-v0002} sets $M_{\underline{W}}$ and $M_{\underline{W}'}$ of \var{\xout{asphe}}{}{typewriter deletion}{ps_134-5_101-v0003}
$\underline{W}$ -aspheric objects in M and M' respectively, and the sets $\underline{W}_M$ and $\underline{W}_{M'}$
of $\underline{W}$-equivalences in M and M' . Then the following relations hold
\[ (\underline{W}) \qquad M_{\underline{W}} = (f^*)^{-1}(M'_{\underline{W}}) \ , \]
\[ (\underline{W}') \qquad \underline{W}_M = (f^*)^{-1}(\underline{W}_{M'}) \ , \]

c) Under the assumption of b) , hence ($\underline{W}'$) holds and $f^*$ gives
rise to a functor \var{\inserted[typed]{$\hand{\overline{\textcolor{black}{f^*}}}$}}{$\hand{\overline{\textcolor{black}{f^*}}}$}{typed insertion}{ps_134-5_101-v0004} between the localizations \var{\xout{\unsure{$\underline{W}_M$}}}{}{typewriter deletion}{ps_134-5_101-v0005}
\[ \mathrm{Hot}_{M,\underline{W}} = \underline{W}_M^{\ -1}M \quad\text{and}\quad \mathrm{Hot}_{M',\underline{W}} = \underline{W}_{M'}^{\ -1}M' \qquad , \]
the following diagram is commutative up to \var{\xout{unique}}{}{typewriter deletion}{ps_134-5_101-v0006} \var{\inserted[typed]{canonical}}{canonical}{typed insertion}{ps_134-5_101-v0007} isomorphism :
% Hand-drawn arrows are marked blue; typed objects remain plain.
\diagram{
\begin{tikzcd}[column sep=small]
\mathrm{Hot}_{M,\underline{W}} \ar[rr,color=blue,"\hand{f^*}"] \ar[dr,color=blue] \& \& \mathrm{Hot}_{M',\underline{W}} \ar[dl,color=blue] \\
\& \mathrm{Hot}_{\underline{W}} \qquad , \&
\end{tikzcd}
}
where the vertical functors are the canonical functors of section 77.
In particular, if the lat\hand{\uline{\textcolor{black}{te}}}r are \var{\xout{\unsure{e}}}{}{typewriter deletion}{ps_134-5_101-v0008}equivalences\hand{(}i.e. M and M' are
$\underline{W}$-modelizing\hand{)} , then so is $f^*$ .

d) Assume merely that the inclusion (!) holds, and let $\underline{W}$ as in b\hand{)}
c) be a basic localizer, satisfying moreover Loc 4). Then the \var{\xout{pair}}{}{typewriter deletion}{ps_134-5_101-v0009}
\var{\xout{$(f_!,f^*)$}}{}{typewriter deletion}{ps_134-5_101-v0010} following conditions on the pair $(f_!,f^*)$ are equivalent :
(i) The pair is a bimorphism, i.e. ($*$) (or equivalently, ($*'$))
holds.
(ii) The inclusion \var{\inserted{$\subset$ in}}{\hand{$\subset$ in}}{insertion}{ps_134-5_101-v0011} ($\underline{W}$) above holds\hand{,} \var{\deleted{\illegible{1}}}{}{deletion}{ps_134-5_101-v0012}
(iii) The inclusion \var{\inserted{$\subset$ in}}{\hand{$\subset$ in}}{insertion}{ps_134-5_101-v0013} ($\underline{W}'$) above holds.
(iv) The functor induced by $f_!$ ,
\[ f_{!\,c} : M'_c \hand{\to} M_c \]
is a\var{\xout{n}}{}{typewriter deletion}{ps_134-5_101-v0014} $\underline{W}$-aspheric functor between the aspherators $M'_c$ , $M_c$ .
(v) (If A is a given small category, and \var{\xout{i}}{}{typewriter deletion}{ps_134-5_101-v0015}
\[ i' : A \hand{\to} M' \]
a given functor, \hand{\uline{\textcolor{black}{factoring}}} \hand{\uline{\textcolor{black}{through}}} $M'_c$ , and $M'_{\underline{W}}$-$\underline{W}$-aspheric.) The
composition
\[ i = f_!\, i' : A \hand{\to} M \]
is $\underline{W}_M$-$\underline{W}$-aspheric.

\margin{101}

%% END SOURCE ps_134-5_101
\ifnum\EditionMode=1
\psapparatusstart
\psentry{ps_134-5_101-v0001}{1}{typewriter deletion}{\xout{\illegible{1}}}{}{The typewriter strikes out this material; it leaves the corrected running text.}
\psentry{ps_134-5_101-v0002}{2}{typed insertion}{\inserted[typed]{the}}{the}{The typed insertion is incorporated at its indicated position; insertion brackets leave the running text.}
\psentry{ps_134-5_101-v0003}{3}{typewriter deletion}{\xout{asphe}}{}{The typewriter strikes out this material; it leaves the corrected running text.}
\psentry{ps_134-5_101-v0004}{4}{typed insertion}{\inserted[typed]{$\hand{\overline{\textcolor{black}{f^*}}}$}}{$\hand{\overline{\textcolor{black}{f^*}}}$}{The typed insertion is incorporated at its indicated position; insertion brackets leave the running text.}
\psentry{ps_134-5_101-v0005}{5}{typewriter deletion}{\xout{\unsure{$\underline{W}_M$}}}{}{The typewriter strikes out this material; it leaves the corrected running text.}
\psentry{ps_134-5_101-v0006}{6}{typewriter deletion}{\xout{unique}}{}{The typewriter strikes out this material; it leaves the corrected running text.}
\psentry{ps_134-5_101-v0007}{7}{typed insertion}{\inserted[typed]{canonical}}{canonical}{The typed insertion is incorporated at its indicated position; insertion brackets leave the running text.}
\psentry{ps_134-5_101-v0008}{8}{typewriter deletion}{\xout{\unsure{e}}}{}{The typewriter strikes out this material; it leaves the corrected running text.}
\psentry{ps_134-5_101-v0009}{9}{typewriter deletion}{\xout{pair}}{}{The typewriter strikes out this material; it leaves the corrected running text.}
\psentry{ps_134-5_101-v0010}{10}{typewriter deletion}{\xout{$(f_!,f^*)$}}{}{The typewriter strikes out this material; it leaves the corrected running text.}
\psentry{ps_134-5_101-v0011}{11}{insertion}{\inserted{$\subset$ in}}{\hand{$\subset$ in}}{The handwritten insertion is incorporated at its indicated position; insertion brackets leave the running text.}
\psentry{ps_134-5_101-v0012}{12}{deletion}{\deleted{\illegible{1}}}{}{The author cancels this material by hand; it leaves the corrected running text.}
\psentry{ps_134-5_101-v0013}{13}{insertion}{\inserted{$\subset$ in}}{\hand{$\subset$ in}}{The handwritten insertion is incorporated at its indicated position; insertion brackets leave the running text.}
\psentry{ps_134-5_101-v0014}{14}{typewriter deletion}{\xout{n}}{}{The typewriter strikes out this material; it leaves the corrected running text.}
\psentry{ps_134-5_101-v0015}{15}{typewriter deletion}{\xout{i}}{}{The typewriter strikes out this material; it leaves the corrected running text.}
\psapparatusend
\fi
\pssource{ps_134-5_102}{134-5/102}{https://grothendieck.umontpellier.fr/web/viewer.html?file=../134-5.pdf\#page=103}
\psnumber{ps_134-5_102-v0001}{1}
\psnumber{ps_134-5_102-v0002}{2}
\psnumber{ps_134-5_102-v0003}{3}
\psnumber{ps_134-5_102-v0004}{4}
\psnumber{ps_134-5_102-v0005}{5}
\psnumber{ps_134-5_102-v0006}{6}
\psnumber{ps_134-5_102-v0007}{7}
\psnumber{ps_134-5_102-v0008}{8}
\psnumber{ps_134-5_102-v0009}{9}
\psnumber{ps_134-5_102-v0010}{10}
\psnumber{ps_134-5_102-v0011}{11}
\psnumber{ps_134-5_102-v0012}{12}
\psnumber{ps_134-5_102-v0013}{13}
\psnumber{ps_134-5_102-v0014}{14}
\psnumber{ps_134-5_102-v0015}{15}
%% BEGIN SOURCE ps_134-5_102
% Pursuing Stacks (cotes 134-2 to 134-8): archive cote 134-5, site page 102 (PDF page 103), typescript page 289.
% typescript. Draft compared: ps4.tex:5189-5310.
% AI-transcribed by Claude Opus 5.5 (claude-opus-5-5, adaptive thinking, effort high), Message Batches API, batch ps2 (msgbatch_01NVwr1Yuo8qJf5rykeGBXPK), 2026-11-12; reviewed against both supplied scan halves; see CHECK_DECISIONS.tsv.
\scanpage{134-5/102}
\mspage{289}

\hand{\uline{\textcolor{black}{Comments}}}. In terms of what is known to us since sections 81 and 82,
and notably that any two sections of a 0-connected object of $M$ o\var{\over[typed]{\unsure{f}}{r}}{r}{typewriter overwriting}{ps_134-5_102-v0001} of $M'$
are homotopic, the Scholie is just a long tautology. I have taken
great care \var{\inserted[typed]{however}}{however}{typed insertion}{ps_134-5_102-v0002} in stating it , so as to get as clear a view as possible of
exactly what the relevant relationships are. In a) the conditions
$(*)$ and $(*')$ are just different ways of stating that $f^*$ is a morphism
of contractibility structures, which can still be expressed in various
other ways, compare p.251-252. Similarly, if we assume that $f_!$ com-%
mutes with finite products, then (!) or (!') can be viewed as expres-%
sing that $f_!$ is a morphism of homotopy structures (in opposite directi-%
on), which again could be expressed in various other ways, for instance
in terms of the homotopy relation for maps, or in terms of homotopisms.
This implies that in any case $f^*$ induces a \var{\xout{map for the}}{}{typewriter deletion}{ps_134-5_102-v0003} functor between
the strict localizations
\[ \hand{\overline{\textcolor{black}{f^*}}}^{\hand{c}}_{\var{\deleted{\illegible{1}}}{}{deletion}{ps_134-5_102-v0004}} : W_c^{-1}M = \overline{M} \hand{\longrightarrow} W'^{-1}_cM' = \overline{M}' \quad , \]
and similarly for $f_!$ when\var{\inserted{\deleted{\illegible{1}}}}{\hand{}}{insertion}{ps_134-5_102-v0005} $f_!$ commutes to finite products. I doubt
however that the latter \var{\inserted{localized}}{\hand{localized}}{insertion}{ps_134-5_102-v0006} map is of geometric relevance, except \var{\inserted{maybe}}{\hand{maybe}}{insertion}{ps_134-5_102-v0007} in the
cases when \var{\xout{it}}{}{typewriter deletion}{ps_134-5_102-v0008} $f_!$ gives rise to relations similar to $(\underline{W})$ and $(\underline{W}')$ in
b) above, and both $f_!$ and $f^*$ are model-preserving with respect to
$\underline{W}$-equivalences, and define quasi-inverse equivalences between the
localizations $\mathrm{Hot}_{M,\underline{W}}$ and $\mathrm{Hot}_{M',\underline{W}}$ - a rather exceptional case indeed.

After stating the Scholie, there is scarcely a doubt left in my
mind about the notion of bimorphism which finally peeled out of our
reflections being a relevant one. This \var{\over{s}{S}}{\hand{S}}{writing over}{ps_134-5_102-v0009}cholie\hand{,} \var{\deleted{now,}}{}{deletion}{ps_134-5_102-v0010} rather than
the ``seducing'' proposition of section 84 (p.275)\hand{,} now seems to me the
adequate \hand{``}answer\hand{''} of the not-so-silly-after-all question of section 46\var{\deleted{.}}{}{deletion}{ps_134-5_102-v0011} \var{\inserted{(cf page 95 -- nearly 200 pages ago!).}}{\hand{(cf page 95 -- nearly 200 pages ago!).}}{insertion}{ps_134-5_102-v0012}
% Displaced small impressions lie below the h in The and above the h in which terminology; neither is read as additional text.
The only minor uncertainty remaining in my mind is whether or not in
the notion of a bimorphism for contractibility structures, we should
insist that $f_!$ should commute to finite products. It seems that it
will be hard to check condition (!), except via an apparently weaker
form such as (!'), and using commutation of $f_!$ to products. But on the
other hand the Scholie makes a good sense, without assuming such commu-%
tation property. If the notion of \var{\xout{a}}{}{typewriter deletion}{ps_134-5_102-v0013} bimorphism is going to be useful,
time \var{\inserted{only}}{\hand{only}}{insertion}{ps_134-5_102-v0014} will tell which terminology use is the best.

More puzzling \var{\xout{is}}{}{typewriter deletion}{ps_134-5_102-v0015} however is this facit, that we still have not
been able to give a satisfactory definition of a \hand{\uline{\textcolor{black}{morphism}}} of \hand{\uline{\textcolor{black}{asphericity}}}
structures, independently of any restrictive assumption on these, such

\margin{102}

%% END SOURCE ps_134-5_102
\ifnum\EditionMode=1
\psapparatusstart
\psentry{ps_134-5_102-v0001}{1}{typewriter overwriting}{\over[typed]{\unsure{f}}{r}}{r}{The superposed typed reading replaces the underlying reading.}
\psentry{ps_134-5_102-v0002}{2}{typed insertion}{\inserted[typed]{however}}{however}{The typed insertion is incorporated at its indicated position; insertion brackets leave the running text.}
\psentry{ps_134-5_102-v0003}{3}{typewriter deletion}{\xout{map for the}}{}{The typewriter strikes out this material; it leaves the corrected running text.}
\psentry{ps_134-5_102-v0004}{4}{deletion}{$\deleted{\illegible{1}}$}{}{The author cancels this material by hand; it leaves the corrected running text.}
\psentry{ps_134-5_102-v0005}{5}{insertion}{\inserted{\deleted{\illegible{1}}}}{\hand{}}{The handwritten insertion is incorporated at its indicated position; insertion brackets leave the running text.}
\psentry{ps_134-5_102-v0006}{6}{insertion}{\inserted{localized}}{\hand{localized}}{The handwritten insertion is incorporated at its indicated position; insertion brackets leave the running text.}
\psentry{ps_134-5_102-v0007}{7}{insertion}{\inserted{maybe}}{\hand{maybe}}{The handwritten insertion is incorporated at its indicated position; insertion brackets leave the running text.}
\psentry{ps_134-5_102-v0008}{8}{typewriter deletion}{\xout{it}}{}{The typewriter strikes out this material; it leaves the corrected running text.}
\psentry{ps_134-5_102-v0009}{9}{writing over}{\over{s}{S}}{\hand{S}}{The superposed handwritten reading replaces the underlying reading.}
\psentry{ps_134-5_102-v0010}{10}{deletion}{\deleted{now,}}{}{The author cancels this material by hand; it leaves the corrected running text.}
\psentry{ps_134-5_102-v0011}{11}{deletion}{\deleted{.}}{}{The author cancels this material by hand; it leaves the corrected running text.}
\psentry{ps_134-5_102-v0012}{12}{insertion}{\inserted{(cf page 95 -- nearly 200 pages ago!).}}{\hand{(cf page 95 -- nearly 200 pages ago!).}}{The handwritten insertion is incorporated at its indicated position; insertion brackets leave the running text.}
\psentry{ps_134-5_102-v0013}{13}{typewriter deletion}{\xout{a}}{}{The typewriter strikes out this material; it leaves the corrected running text.}
\psentry{ps_134-5_102-v0014}{14}{insertion}{\inserted{only}}{\hand{only}}{The handwritten insertion is incorporated at its indicated position; insertion brackets leave the running text.}
\psentry{ps_134-5_102-v0015}{15}{typewriter deletion}{\xout{is}}{}{The typewriter strikes out this material; it leaves the corrected running text.}
\psapparatusend
\fi
\pssource{ps_134-5_103}{134-5/103}{https://grothendieck.umontpellier.fr/web/viewer.html?file=../134-5.pdf\#page=104}
\psnumber{ps_134-5_103-v0001}{1}
\psnumber{ps_134-5_103-v0002}{2}
\psnumber{ps_134-5_103-v0003}{3}
\psnumber{ps_134-5_103-v0004}{4}
\psnumber{ps_134-5_103-v0005}{5}
\psnumber{ps_134-5_103-v0006}{6}
\psnumber{ps_134-5_103-v0007}{7}
\psnumber{ps_134-5_103-v0008}{8}
\psnumber{ps_134-5_103-v0009}{9}
\psnumber{ps_134-5_103-v0010}{10}
\psnumber{ps_134-5_103-v0011}{11}
%% BEGIN SOURCE ps_134-5_103
% Pursuing Stacks (cotes 134-2 to 134-8): archive cote 134-5, site page 103 (PDF page 104), typescript page 290.
% typescript. Draft compared: ps4.tex:5297-5325.
% AI-transcribed by Claude Opus 5.5 (claude-opus-5-5, adaptive thinking, effort high), Message Batches API, batch ps2 (msgbatch_01NVwr1Yuo8qJf5rykeGBXPK), 2026-11-12; reviewed against both supplied scan halves; see CHECK_DECISIONS.tsv.
\scanpage{134-5/103}
\mspage{290}
as \var{\xout{com}}{}{typewriter deletion}{ps_134-5_103-v0001} being generated by a contractibility structure.  Even when making
such an assumption,  it is not wholly clear yet that there is\hand{'}nt a good
notion of a morphism $f^{*}$ of contractibility structures,  without having
to assume there exists a left adjoint $f_{!}$ .  Maybe a little more ponde-%
ring on the situation would be useful.  If it should become clear that
(except for the obvious notion of \hand{\uline{\textcolor{black}{equivalence}}} of asphericity structures)
there was \hand{\uline{\textcolor{black}{not}}} any reasonable notion of a morphism of asphericity struc-%
tures,  this would probably mean that the notion of an asphericity
structure should not be viewed as a main structure type in homotopy
theory in it's own right,  but rather,  mainly,  as \var{\deleted{an}}{}{deletion}{ps_134-5_103-v0002} \var{\inserted{important}}{\hand{important}}{insertion}{ps_134-5_103-v0003} by\var{\deleted{\unsure{e}}}{}{deletion}{ps_134-5_103-v0004}-product \var{\deleted{\unsure{or|on}}}{}{deletion}{ps_134-5_103-v0005}
\var{\deleted{an interesting aspect}}{}{deletion}{ps_134-5_103-v0006} of a cont\var{\over[typed]{\illegible{1}}{r}}{r}{typewriter overwriting}{ps_134-5_103-v0007}actibility s\var{\over[typed]{g}{t}}{t}{typewriter overwriting}{ps_134-5_103-v0008}ructure.  At any rate,
since section 81,  I feel that the main emphasis has \var{\inserted{definitely}}{\hand{definitely}}{insertion}{ps_134-5_103-v0009} been shifting
towards contractibility structures,  which seems now \hand{\uline{\textcolor{black}{the}}} main type of
structure dominating \var{\xout{the}}{}{typewriter deletion}{ps_134-5_103-v0010} in the modelizing story,  as this \var{\inserted{``story''}}{\hand{``story''}}{insertion}{ps_134-5_103-v0011} is gradually
% The curved locator below the inserted story does not cancel the typed is.
emerging into light.

\margin{103}

%% END SOURCE ps_134-5_103
\ifnum\EditionMode=1
\psapparatusstart
\psentry{ps_134-5_103-v0001}{1}{typewriter deletion}{\xout{com}}{}{The typewriter strikes out this material; it leaves the corrected running text.}
\psentry{ps_134-5_103-v0002}{2}{deletion}{\deleted{an}}{}{The author cancels this material by hand; it leaves the corrected running text.}
\psentry{ps_134-5_103-v0003}{3}{insertion}{\inserted{important}}{\hand{important}}{The handwritten insertion is incorporated at its indicated position; insertion brackets leave the running text.}
\psentry{ps_134-5_103-v0004}{4}{deletion}{\deleted{\unsure{e}}}{}{The author cancels this material by hand; it leaves the corrected running text.}
\psentry{ps_134-5_103-v0005}{5}{deletion}{\deleted{\unsure{or|on}}}{}{The author cancels this material by hand; it leaves the corrected running text.}
\psentry{ps_134-5_103-v0006}{6}{deletion}{\deleted{an interesting aspect}}{}{The author cancels this material by hand; it leaves the corrected running text.}
\psentry{ps_134-5_103-v0007}{7}{typewriter overwriting}{\over[typed]{\illegible{1}}{r}}{r}{The superposed typed reading replaces the underlying reading.}
\psentry{ps_134-5_103-v0008}{8}{typewriter overwriting}{\over[typed]{g}{t}}{t}{The superposed typed reading replaces the underlying reading.}
\psentry{ps_134-5_103-v0009}{9}{insertion}{\inserted{definitely}}{\hand{definitely}}{The handwritten insertion is incorporated at its indicated position; insertion brackets leave the running text.}
\psentry{ps_134-5_103-v0010}{10}{typewriter deletion}{\xout{the}}{}{The typewriter strikes out this material; it leaves the corrected running text.}
\psentry{ps_134-5_103-v0011}{11}{insertion}{\inserted{``story''}}{\hand{``story''}}{The handwritten insertion is incorporated at its indicated position; insertion brackets leave the running text.}
\psapparatusend
\fi
\pssource{ps_134-5_104}{134-5/104}{https://grothendieck.umontpellier.fr/web/viewer.html?file=../134-5.pdf\#page=105}
\psnumber{ps_134-5_104-v0001}{1}
\psnumber{ps_134-5_104-v0002}{2}
\psnumber{ps_134-5_104-v0003}{3}
\psnumber{ps_134-5_104-v0004}{4}
\psnumber{ps_134-5_104-v0005}{5}
\psnumber{ps_134-5_104-v0006}{6}
\psnumber{ps_134-5_104-v0007}{7}
\psnumber{ps_134-5_104-v0008}{8}
\psnumber{ps_134-5_104-v0009}{9}
\psnumber{ps_134-5_104-v0010}{10}
%% BEGIN SOURCE ps_134-5_104
% Pursuing Stacks (cotes 134-2 to 134-8): archive cote 134-5, site page 104 (PDF page 105), typescript page 291.
% typescript. Draft compared: ps5.tex:1-64.
% AI-transcribed by Claude Opus 5.5 (claude-opus-5-5, adaptive thinking, effort high), Message Batches API, batch ps2 (msgbatch_01NVwr1Yuo8qJf5rykeGBXPK), 2026-11-12; reviewed against both supplied scan halves; see CHECK_DECISIONS.tsv.
\scanpage{134-5/104}
\mspage{291}

1.7.

87.    I could'nt resist last night and had to look through Thomason's
% The h in The has displaced impressions; namely, the second The, and well also have crowded typewriter impressions without secure earlier readings.
preprint on the closed model structure of (Cat).  The paper is  really
pleasant reading  -  and it gives exactly what had been la\var{\over[typed]{\unsure{c}}{k}}{k}{typewriter overwriting}{ps_134-5_104-v0001}king me in my
reflections lately on the homotopy theory of (Cat)  -  namely a class of
neat monomorphisms  $Y\hand{\longrightarrow}X$  which \var{\inserted{all}}{\hand{all}}{insertion}{ps_134-5_104-v0002} have the property that cobase change
by\var{\xout{\illegible{1}}}{}{typewriter deletion}{ps_134-5_104-v0003} these preserves weak equivalence,  the so-called \hand{\uline{\textcolor{black}{Dwyer}}} \hand{\uline{\textcolor{black}{maps}}}.  I had
hoped for a while that``open immersions''  and their duals,  the ``closed
immersions''  (namely sieve and cosieve maps,  in Thomason's wording), have
this property,  and when it turned out they had'nt,  I had been at a loss
of what stronger property put in instead,  wide enough however to allow
for the standard factorisation statements for a map to go through.  The
definition of a Dwyer map is an extremely pretty one  -  it is an open
immersion $Y\hand{\to}X$  such that the induced map from Y into it's closure \hand{$\overline{Y}$}
should have a right adjoint.  Now this implies that $Y\hand{\to}\overline{Y}$  is aspheric,
and I suspect that this \var{\inserted[typed]{extra}}{extra}{typed insertion}{ps_134-5_104-v0004} condition on an open immersion $Y\hand{\longrightarrow}X$ should be
sufficient,  to imply that it has the ``cofibration property'' above with
respect to weak equivalence.  It would mean in a sense that the given
open immersion is very close to be a closed immersion too,  without how-%
ever being a direct summand necessarily.  The dual notion  is that of a
\margin{\tikz[baseline=(m.base)]{\node[draw,circle,inner sep=2pt] (m) {$\mathring{Y}$};}}closed immersion such that the corresponding interior  $\hand{\mathring{\phantom{Y}}}\!\!Y$  of Y in X gives
a map  $\hand{\mathring{\textcolor{black}{Y}}}\hand{\longrightarrow}Y$  which has a left adjoint,  or \var{\inserted{only}}{\hand{only}}{insertion}{ps_134-5_104-v0005} which is ``coaspheric'' .
By \var{\xout{the}}{}{typewriter deletion}{ps_134-5_104-v0006} Qui\var{\over{ll}{ll}}{\hand{ll}}{writing over}{ps_134-5_104-v0007}en's duality principle,  if one notion works well for push-%
outs,  so does the other.  -  With this notion in hands,  it should'nt
be difficu\var{\over[typed]{b}{l}}{l}{typewriter overwriting}{ps_134-5_104-v0008}t now to get a closed model structure on (Cat) a lot simpler
than Thomason's.  Visibly he was hampered by the standing reflex :
homotopy = semi-simplicial algebra,  which caused him to pass by the de-%
tour of the category $\hand{\Delta}\hat{}$ of semi-simplicial complexes,  rather than just
working in (Cat) itself.  I'll \var{\inserted{have to}}{\hand{have to}}{insertion}{ps_134-5_104-v0009} come back upon this in part V  of the notes\hand{,}
where I intend to investigate the homotopy properties of (Cat) and elemen-%
tary modelizers $A\hat{}$,  including the existence of closed model structures.

    Some comments of Thomason's at the end of his preprint,  about
applications to algebraic K-theory,  seem to indicate that the notion
of ``integration'' and ``cointegration'' of homotopy types I have been inte-%
rested in,  has been studied (under the name holim and \var{\over{\unsure{c}}{h}}{\hand{h}}{writing over}{ps_134-5_104-v0010}ocolim) in the
context of closed model categories by Anderson  (his paper appeared in
1978).   As I am going to develop some ideas along these lines in part VI
on derivators,  I should have a look at what Anderson does,  notably

\margin{104}

%% END SOURCE ps_134-5_104
\ifnum\EditionMode=1
\psapparatusstart
\psentry{ps_134-5_104-v0001}{1}{typewriter overwriting}{\over[typed]{\unsure{c}}{k}}{k}{The superposed typed reading replaces the underlying reading.}
\psentry{ps_134-5_104-v0002}{2}{insertion}{\inserted{all}}{\hand{all}}{The handwritten insertion is incorporated at its indicated position; insertion brackets leave the running text.}
\psentry{ps_134-5_104-v0003}{3}{typewriter deletion}{\xout{\illegible{1}}}{}{The typewriter strikes out this material; it leaves the corrected running text.}
\psentry{ps_134-5_104-v0004}{4}{typed insertion}{\inserted[typed]{extra}}{extra}{The typed insertion is incorporated at its indicated position; insertion brackets leave the running text.}
\psentry{ps_134-5_104-v0005}{5}{insertion}{\inserted{only}}{\hand{only}}{The handwritten insertion is incorporated at its indicated position; insertion brackets leave the running text.}
\psentry{ps_134-5_104-v0006}{6}{typewriter deletion}{\xout{the}}{}{The typewriter strikes out this material; it leaves the corrected running text.}
\psentry{ps_134-5_104-v0007}{7}{writing over}{\over{ll}{ll}}{\hand{ll}}{The superposed handwritten reading replaces the underlying reading.}
\psentry{ps_134-5_104-v0008}{8}{typewriter overwriting}{\over[typed]{b}{l}}{l}{The superposed typed reading replaces the underlying reading.}
\psentry{ps_134-5_104-v0009}{9}{insertion}{\inserted{have to}}{\hand{have to}}{The handwritten insertion is incorporated at its indicated position; insertion brackets leave the running text.}
\psentry{ps_134-5_104-v0010}{10}{writing over}{\over{\unsure{c}}{h}}{\hand{h}}{The superposed handwritten reading replaces the underlying reading.}
\psapparatusend
\fi
\pssource{ps_134-5_105}{134-5/105}{https://grothendieck.umontpellier.fr/web/viewer.html?file=../134-5.pdf\#page=106}
\psnumber{ps_134-5_105-v0001}{1}
\psnumber{ps_134-5_105-v0002}{2}
\psnumber{ps_134-5_105-v0003}{3}
\psnumber{ps_134-5_105-v0004}{4}
\psnumber{ps_134-5_105-v0005}{5}
\psnumber{ps_134-5_105-v0006}{6}
\psnumber{ps_134-5_105-v0007}{7}
\psnumber{ps_134-5_105-v0008}{8}
\psnumber{ps_134-5_105-v0009}{9}
\psnumber{ps_134-5_105-v0010}{10}
\psnumber{ps_134-5_105-v0011}{11}
%% BEGIN SOURCE ps_134-5_105
% Pursuing Stacks (cotes 134-2 to 134-8): archive cote 134-5, site page 105 (PDF page 106), typescript page 292.
% typescript. Draft compared: ps5.tex:52-116.
% AI-transcribed by Claude Opus 5.5 (claude-opus-5-5, adaptive thinking, effort high), Message Batches API, batch ps2 (msgbatch_01NVwr1Yuo8qJf5rykeGBXPK), 2026-11-12; reviewed against both supplied scan halves; see CHECK_DECISIONS.tsv.
\scanpage{134-5/105}
\mspage{292}

% Displaced or crowded typewriter impressions occur at n in Indeed, o/w in Now, era in operation, and o in cointegration; their earlier readings are not securely separable.
what his assumptions on the indexing categories are.  Thomason seems to
believe that the closed model structure of (Cat) is essential for being
able to take homotopy limits  -  whereas it is clear a priori to me that
the notion \var{\xout{should}}{}{typewriter deletion}{ps_134-5_105-v0001} depends only on the notion of weak equivalence.  Indeed,
he seems to consider the possibility of taking homotopy limits in (Cat)
as the main application of his theorem,  and in order to apply Anderson's
results,  believes it is necessary to be able to give concrete caracteri-%
zations for ``fibrations'' and ``cofibrations'' of his closed model structure.
Now it turns out that the case he is interested in (for proving \hand{``}Lichten-%
baum's conjecture\hand{''}) is \var{\xout{the case of}}{}{typewriter deletion}{ps_134-5_105-v0002} a typical case of \hand{\uline{\textcolor{black}{direct}}} homotopy
limits,  namely ``\hand{\uwave{\textcolor{black}{integration}}}''  -  which can be described directly in (Cat)
in such an amazingly simple way (as sketched in section 69, p.198-199).
Thus I feel for this application,  the closed model structure is wholly
irrelevant.  As for cointegration,  I do not expect that there is a com-%
parably simple construction of this operation  within the modelizer (Cat),
but presumably there is when taking $\mathrm{oo}$-Gr-stacks as models  (as suggested
by the ``geometric'' approach to cohomology invariants,  via stacks, \var{\xout{precisely}}{}{typewriter deletion}{ps_134-5_105-v0003}
where the operation of ``direct image''  namely cohomology precisely,  is
the obvious one,  whereas inverse images are more delicate to define, by
an adjunction property with respect to direct images ...).  When working
in (Cat),  cointegration of homotopy types should be not more no\var{\over{t}{r}}{\hand{r}}{writing over}{ps_134-5_105-v0004} less
invol\var{\xout{o}}{}{typewriter deletion}{ps_134-5_105-v0005}ved than in  any closed model category say,  and involve intensive
recourse to ``fibrations'' (in the sense of the closed model structure,  or
more intrinsically,  in the sense that base change by these should preser-%
ve weak equivalence).  Now  the latter have become quite familiar \var{\inserted[typed]{to me}}{to me}{typed insertion}{ps_134-5_105-v0006} \var{\xout{to me}}{}{typewriter deletion}{ps_134-5_105-v0007}
during my long scratchwork on cohomology properties of maps in (Cat),
and I'll \var{\xout{repeat}}{}{typewriter deletion}{ps_134-5_105-v0008} have to try put it down nicely in part V of the notes.

88.     The present ``part IV''  on asphericity structures (and their rela-%
tions to contractibility structures)  turns out a lot longer than I had
anticipated,  and the end is not yet quite in sight !  Therefore,  before
pursuing,  I wou\var{\over[typed]{b}{l}}{l}{typewriter overwriting}{ps_134-5_105-v0009}d like to make a review of the questions \var{\xout{which}}{}{typewriter deletion}{ps_134-5_105-v0010} along
these lines which seem to require elucidation,  and then decide which I'm
going to deal with,  before going over to part V .

1)   Whereas the relevant notions of ``morphisms''  and ``bimorphisms''
for \var{\over[typed]{b}{c}}{c}{typewriter overwriting}{ps_134-5_105-v0011}ontractibility structures seem to me well understood,  there remains

\margin{105}

%% END SOURCE ps_134-5_105
\ifnum\EditionMode=1
\psapparatusstart
\psentry{ps_134-5_105-v0001}{1}{typewriter deletion}{\xout{should}}{}{The typewriter strikes out this material; it leaves the corrected running text.}
\psentry{ps_134-5_105-v0002}{2}{typewriter deletion}{\xout{the case of}}{}{The typewriter strikes out this material; it leaves the corrected running text.}
\psentry{ps_134-5_105-v0003}{3}{typewriter deletion}{\xout{precisely}}{}{The typewriter strikes out this material; it leaves the corrected running text.}
\psentry{ps_134-5_105-v0004}{4}{writing over}{\over{t}{r}}{\hand{r}}{The superposed handwritten reading replaces the underlying reading.}
\psentry{ps_134-5_105-v0005}{5}{typewriter deletion}{\xout{o}}{}{The typewriter strikes out this material; it leaves the corrected running text.}
\psentry{ps_134-5_105-v0006}{6}{typed insertion}{\inserted[typed]{to me}}{to me}{The typed insertion is incorporated at its indicated position; insertion brackets leave the running text.}
\psentry{ps_134-5_105-v0007}{7}{typewriter deletion}{\xout{to me}}{}{The typewriter strikes out this material; it leaves the corrected running text.}
\psentry{ps_134-5_105-v0008}{8}{typewriter deletion}{\xout{repeat}}{}{The typewriter strikes out this material; it leaves the corrected running text.}
\psentry{ps_134-5_105-v0009}{9}{typewriter overwriting}{\over[typed]{b}{l}}{l}{The superposed typed reading replaces the underlying reading.}
\psentry{ps_134-5_105-v0010}{10}{typewriter deletion}{\xout{which}}{}{The typewriter strikes out this material; it leaves the corrected running text.}
\psentry{ps_134-5_105-v0011}{11}{typewriter overwriting}{\over[typed]{b}{c}}{c}{The superposed typed reading replaces the underlying reading.}
\psapparatusend
\fi
\pssource{ps_134-5_106}{134-5/106}{https://grothendieck.umontpellier.fr/web/viewer.html?file=../134-5.pdf\#page=107}
\psnumber{ps_134-5_106-v0001}{1}
\psnumber{ps_134-5_106-v0002}{2}
\psnumber{ps_134-5_106-v0003}{3}
\psnumber{ps_134-5_106-v0004}{4}
\psnumber{ps_134-5_106-v0005}{5}
\psnumber{ps_134-5_106-v0006}{6}
\psnumber{ps_134-5_106-v0007}{7}
\psnumber{ps_134-5_106-v0008}{8}
%% BEGIN SOURCE ps_134-5_106
% Pursuing Stacks (cotes 134-2 to 134-8): archive cote 134-5, site page 106 (PDF page 107), typescript page 293.
% typescript. Draft compared: ps5.tex:102-159.
% AI-transcribed by Claude Opus 5.5 (claude-opus-5-5, adaptive thinking, effort high), Message Batches API, batch ps2 (msgbatch_01NVwr1Yuo8qJf5rykeGBXPK), 2026-11-12; reviewed against both supplied scan halves; see CHECK_DECISIONS.tsv.
\scanpage{134-5/106}
\mspage{293}

a certain feeling of uneasiness with respect to asphericity structures,
which hav'nt got yet a reasonable notion of morphism. Thus I have left
unanswered the questions raised in section 84 around the inclusion \var{\inserted{condition}}{\hand{condition}}{insertion}{ps_134-5_106-v0001}
\[ f_{!}(M_{a}) \hand{\subset} M'_{a} \qquad , \]
so it is not impossible that, while relying on the mere feeling that
the inclusion is just not reasonable and that the answer to the specific
questions are presumably negative, I am about to miss some unexpected \hand{im-}
\var{\over{\illegible{1}}{portant}}{\hand{portant}}{writing over}{ps_134-5_106-v0002} fact ! It seems I developped kind of a block against checking -
maybe the answers are well known and Tim Porter will tell me ... Maybe
I better leave the question for a later moment, as it will ripen by it-%
self if I leave it alone ...

2) I should at last introduce contractors, and morphisms of such.
When I set out on part IV of the notes, I expected that the notion of a
contractor would be one main notion, alongside with the notion of a
canonical modelizer - it turns out that so far I did'nt have much use
yet for one or the other. Contractors \var{c\hand{\tikz[x=.55em,y=.35em,baseline=-.5ex]{\draw (0,.5) .. controls (.35,0) and (.7,.1) .. (1,.4);}}an}{can}{placement}{ps_134-5_106-v0003} be viewed as \var{\xout{small}}{}{typewriter deletion}{ps_134-5_106-v0004} categories
generating contractibility structures, just as aspherators are there
for generating asphericity structures.
% A detached U-like impression is visible above the o of However. How\over{\illegible{4}}{ever}, whereas any small cate-%
gory is an aspherator, the same is definitely not so for contractors,
as we demand that every object in C should be contractible, for the
homotopy structure in C\^{} generated by C itself. If we except the case
of a contractor equivalent to the final category (\var{\xout{the}}{}{typewriter deletion}{ps_134-5_106-v0005} a so-called \hand{\uline{\textcolor{black}{trivial}}}
contractor), a contractor is a strict test category - thus the notion
appears somewhat as a hinge notion between the test notions, and the
\var{\xout{\illegible{1}}}{}{typewriter deletion}{ps_134-5_106-v0006}``pure'' homotopy notions and more specifically, contractibility structu-%
res. Writing up some scratchnotes I got should be a pure routine matter.

3) A lot more interesting seems to me to try and resolve a persis-%
tent feeling of uneasiness which has been floating, throughout the
long-winded reflections on homotopy and asphericity structures in parts
III and IV of the notes. This is tied up with this fact, that \var{\over{\illegible{2}}{my}}{\hand{my}}{writing over}{ps_134-5_106-v0007} treat-%
ment of the main notions, namely contractibility and asphericity, has
been consistently \hand{\uline{\textcolor{black}{non}}} \hand{\uline{\textcolor{black}{auto}}}-\hand{\uline{\textcolor{black}{dual}}}. More specifically, when a category M
is endowed with either a contractibility or an asphericity structure,
it does \hand{\uline{\textcolor{black}{not}}} follow that the opposit category M$^{\circ}$ is too in a natural way.
When defining homotopy relations and homotopism structures (section 51,
A) and B)), these were autodual notions, but the notion of a\var{\xout{n}}{}{typewriter deletion}{ps_134-5_106-v0008} homotopy
interval structure, which we used in order to pass from a contractibility

\margin{106}

%% END SOURCE ps_134-5_106
\ifnum\EditionMode=1
\psapparatusstart
\psentry{ps_134-5_106-v0001}{1}{insertion}{\inserted{condition}}{\hand{condition}}{The handwritten insertion is incorporated at its indicated position; insertion brackets leave the running text.}
\psentry{ps_134-5_106-v0002}{2}{writing over}{\over{\illegible{1}}{portant}}{\hand{portant}}{The superposed handwritten reading replaces the underlying reading.}
\psentry{ps_134-5_106-v0003}{3}{placement}{c\hand{\tikz[x=.55em,y=.35em,baseline=-.5ex]{\draw (0,.5) .. controls (.35,0) and (.7,.1) .. (1,.4);}}an}{can}{The pen curve joins the accidentally separated typed parts c and an; it is a nonlexical joining locator.}
\psentry{ps_134-5_106-v0004}{4}{typewriter deletion}{\xout{small}}{}{The typewriter strikes out this material; it leaves the corrected running text.}
\psentry{ps_134-5_106-v0005}{5}{typewriter deletion}{\xout{the}}{}{The typewriter strikes out this material; it leaves the corrected running text.}
\psentry{ps_134-5_106-v0006}{6}{typewriter deletion}{\xout{\illegible{1}}}{}{The typewriter strikes out this material; it leaves the corrected running text.}
\psentry{ps_134-5_106-v0007}{7}{writing over}{\over{\illegible{2}}{my}}{\hand{my}}{The superposed handwritten reading replaces the underlying reading.}
\psentry{ps_134-5_106-v0008}{8}{typewriter deletion}{\xout{n}}{}{The typewriter strikes out this material; it leaves the corrected running text.}
\psapparatusend
\fi
\pssource{ps_134-5_107}{134-5/107}{https://grothendieck.umontpellier.fr/web/viewer.html?file=../134-5.pdf\#page=108}
\psnumber{ps_134-5_107-v0001}{1}
\psnumber{ps_134-5_107-v0002}{2}
\psnumber{ps_134-5_107-v0003}{3}
\psnumber{ps_134-5_107-v0004}{4}
\psnumber{ps_134-5_107-v0005}{5}
\psnumber{ps_134-5_107-v0006}{6}
\psnumber{ps_134-5_107-v0007}{7}
\psnumber{ps_134-5_107-v0008}{8}
\psnumber{ps_134-5_107-v0009}{9}
\psnumber{ps_134-5_107-v0010}{10}
\psnumber{ps_134-5_107-v0011}{11}
\psnumber{ps_134-5_107-v0012}{12}
\psnumber{ps_134-5_107-v0013}{13}
\psnumber{ps_134-5_107-v0014}{14}
\psnumber{ps_134-5_107-v0015}{15}
\psnumber{ps_134-5_107-v0016}{16}
\psnumber{ps_134-5_107-v0017}{17}
\psnumber{ps_134-5_107-v0018}{18}
\psnumber{ps_134-5_107-v0019}{19}
%% BEGIN SOURCE ps_134-5_107
% Pursuing Stacks (cotes 134-2 to 134-8): archive cote 134-5, site page 107 (PDF page 108), typescript page 294.
% typescript. Draft compared: ps5.tex:148-203.
% AI-transcribed by Claude Opus 5.5 (claude-opus-5-5, adaptive thinking, effort high), Message Batches API, batch ps2 (msgbatch_01NVwr1Yuo8qJf5rykeGBXPK), 2026-11-12; reviewed against both supplied scan halves; see CHECK_DECISIONS.tsv.
\scanpage{134-5/107}
\mspage{294}
structure to the corresponding notions of homotopy equivalence between
maps and of homotopisms,  is highly non autodual too.  It breaks down alto-%
gether when $M$ is a ``pointed'' category,  namely \var{\xout{has}}{}{typewriter deletion}{ps_134-5_107-v0001} contains an object which
is both initial and final \var{\deleted{\illegible{1}}}{}{deletion}{ps_134-5_107-v0002}\hand{-}  in this case,  \var{\xout{the}}{}{typewriter deletion}{ps_134-5_107-v0003} for any homotopy interval
structure on $M$,  any two maps in $M$ are homotopic,  and any map is a homoto-%
pism,  hence any object is contractible !

% Small displaced impressions occur at e in categories, e in features, and the T/A/h cluster in The most; their earlier readings are not securely recoverable.
Our initial motivation,  namely understanding ``modelizers'' for ordi-%
nary homotopy types,  made it very natural to get involved in non autodual
situation,  as the homotopy category (Hot) itself,  and the usual model
categories for it,  displays strongly non autodual features. \hand{(}Thus,  where-%
as the usual test categories all have a final object,  it is easy to see
that a test-category cannot possibly have an initial object.\hand{)} However,
the homotopy and asphericity notions we then came to develop make a sense
and are familiar indeed not only in the ``modelizing story'',  but \var{\xout{practical-}}{}{typewriter deletion}{ps_134-5_107-v0004}
\var{\xout{ly}}{}{typewriter deletion}{ps_134-5_107-v0005} in any situation whatever which turns up,  giving rise to anything like
a ``homology'' or a ``homotopy'' theory.  To give one specific example,  star-%
ting with an abelian category,  the corresponding homology theory is con-%
cerned with the category of A-valued \var{\inserted[typed]{\xout{chain}}}{}{typed insertion}{ps_134-5_107-v0006} complexes,  say $K$\hand{$^{\bullet}$(A)}.  \unsure{The} most
basic notions here are the three homotopy notions and two asphericity
notions,  namely:  homotopy equivalence between maps,  homotopisms,  contra\unsure{c}%
% NOTE: the end of this line is cut off at the right edge of the scan
tible objects, \var{\inserted{and}}{\hand{and}}{insertion}{ps_134-5_107-v0007} quasi-isomorphism\hand{s} (= ``weak equivalence\hand{s}''),  acyclic (=
``aspheric'') objects.  The two first homotopy notions determine each other
mutually in the ususal way,  and define the third,  namely an object is
contractible \var{\xout{iff}}{}{typewriter deletion}{ps_134-5_107-v0008} \hand{(}or null-homotopic\hand{)} iff the map $0\hand{\to} X$,  or equivalently $X\hand{\to} 0$,
is a homotopism.  On the\hand{/}other hand,  if we use the mapping-cylinder constru%
% NOTE: the end of this line is cut off at the right edge of the scan
-tion for a map,  the \var{\xout{notion}}{}{typewriter deletion}{ps_134-5_107-v0009} set of contractible objects determines the set
of homotopisms,  as the maps whose mapping cylinder is contractible.  Like-%
wise,  weak equivalences determine\var{\xout{s}}{}{typewriter deletion}{ps_134-5_107-v0010} \var{\xout{the notion of}}{}{typewriter deletion}{ps_134-5_107-v0011} aspheric objects,  and con-%
versely if the mapping cylinder construction is given.  The question that
now comes to mind immediately \var{\inserted[typed]{is}}{is}{typed insertion}{ps_134-5_107-v0012} \var{\deleted{wather}}{}{deletion}{ps_134-5_107-v0013} \var{\inserted{whether}}{\hand{whether}}{insertion}{ps_134-5_107-v0014} the two sets of notions,  the \var{\inserted[typed]{three}}{three}{typed insertion}{ps_134-5_107-v0015} pure
homotopy notions (determining each other)\hand{,} and the \var{\inserted[typed]{two}}{two}{typed insertion}{ps_134-5_107-v0016} ``asphericity notions''
(determining each other too),  mutually determine each other,  as in the
non commutative set-up we have \var{\over{w}{w}}{\hand{w}}{writing over}{ps_134-5_107-v0017}orked in so far.  We can also remark that
the functor
\[ H^{0} : K^{\hand{\bullet}}(A) \hand{\longrightarrow} A \]
visibly plays the part of the functor \hand{$\pi$}$_{0}$ in the non commutative set up,
it gives rise moreover to the $H^{i}$ functors (any inte\var{\xout{\unsure{g}}}{}{typewriter deletion}{ps_134-5_107-v0018}ger $i$) by composing
with the iterated shift functor  (where the shift of $X$ is just the ma\var{\over{pp}{pp}}{\hand{pp}}{writing over}{ps_134-5_107-v0019}ing

\margin{107}

%% END SOURCE ps_134-5_107
\ifnum\EditionMode=1
\psapparatusstart
\psentry{ps_134-5_107-v0001}{1}{typewriter deletion}{\xout{has}}{}{The typewriter strikes out this material; it leaves the corrected running text.}
\psentry{ps_134-5_107-v0002}{2}{deletion}{\deleted{\illegible{1}}}{}{The author cancels this material by hand; it leaves the corrected running text.}
\psentry{ps_134-5_107-v0003}{3}{typewriter deletion}{\xout{the}}{}{The typewriter strikes out this material; it leaves the corrected running text.}
\psentry{ps_134-5_107-v0004}{4}{typewriter deletion}{\xout{practical-}}{}{The typewriter strikes out this material; it leaves the corrected running text.}
\psentry{ps_134-5_107-v0005}{5}{typewriter deletion}{\xout{ly}}{}{The typewriter strikes out this material; it leaves the corrected running text.}
\psentry{ps_134-5_107-v0006}{6}{typed insertion}{\inserted[typed]{\xout{chain}}}{}{The typed insertion is incorporated at its indicated position; insertion brackets leave the running text.}
\psentry{ps_134-5_107-v0007}{7}{insertion}{\inserted{and}}{\hand{and}}{The handwritten insertion is incorporated at its indicated position; insertion brackets leave the running text.}
\psentry{ps_134-5_107-v0008}{8}{typewriter deletion}{\xout{iff}}{}{The typewriter strikes out this material; it leaves the corrected running text.}
\psentry{ps_134-5_107-v0009}{9}{typewriter deletion}{\xout{notion}}{}{The typewriter strikes out this material; it leaves the corrected running text.}
\psentry{ps_134-5_107-v0010}{10}{typewriter deletion}{\xout{s}}{}{The typewriter strikes out this material; it leaves the corrected running text.}
\psentry{ps_134-5_107-v0011}{11}{typewriter deletion}{\xout{the notion of}}{}{The typewriter strikes out this material; it leaves the corrected running text.}
\psentry{ps_134-5_107-v0012}{12}{typed insertion}{\inserted[typed]{is}}{is}{The typed insertion is incorporated at its indicated position; insertion brackets leave the running text.}
\psentry{ps_134-5_107-v0013}{13}{deletion}{\deleted{wather}}{}{The author cancels this material by hand; it leaves the corrected running text.}
\psentry{ps_134-5_107-v0014}{14}{insertion}{\inserted{whether}}{\hand{whether}}{The handwritten insertion is incorporated at its indicated position; insertion brackets leave the running text.}
\psentry{ps_134-5_107-v0015}{15}{typed insertion}{\inserted[typed]{three}}{three}{The typed insertion is incorporated at its indicated position; insertion brackets leave the running text.}
\psentry{ps_134-5_107-v0016}{16}{typed insertion}{\inserted[typed]{two}}{two}{The typed insertion is incorporated at its indicated position; insertion brackets leave the running text.}
\psentry{ps_134-5_107-v0017}{17}{writing over}{\over{w}{w}}{\hand{w}}{The superposed handwritten reading replaces the underlying reading.}
\psentry{ps_134-5_107-v0018}{18}{typewriter deletion}{\xout{\unsure{g}}}{}{The typewriter strikes out this material; it leaves the corrected running text.}
\psentry{ps_134-5_107-v0019}{19}{writing over}{\over{pp}{pp}}{\hand{pp}}{The superposed handwritten reading replaces the underlying reading.}
\psapparatusend
\fi
\pssource{ps_134-5_108}{134-5/108}{https://grothendieck.umontpellier.fr/web/viewer.html?file=../134-5.pdf\#page=109}
\psnumber{ps_134-5_108-v0001}{1}
\psnumber{ps_134-5_108-v0002}{2}
\psnumber{ps_134-5_108-v0003}{3}
\psnumber{ps_134-5_108-v0004}{4}
\psnumber{ps_134-5_108-v0005}{5}
\psnumber{ps_134-5_108-v0006}{6}
\psnumber{ps_134-5_108-v0007}{7}
\psnumber{ps_134-5_108-v0008}{8}
\psnumber{ps_134-5_108-v0009}{9}
\psnumber{ps_134-5_108-v0010}{10}
\psnumber{ps_134-5_108-v0011}{11}
\psnumber{ps_134-5_108-v0012}{12}
\psnumber{ps_134-5_108-v0013}{13}
\psnumber{ps_134-5_108-v0014}{14}
\psnumber{ps_134-5_108-v0015}{15}
%% BEGIN SOURCE ps_134-5_108
% Pursuing Stacks (cotes 134-2 to 134-8): archive cote 134-5, site page 108 (PDF page 109), typescript page 295.
% typescript. Draft compared: ps5.tex:189-246.
% AI-transcribed by Claude Opus 5.5 (claude-opus-5-5, adaptive thinking, effort high), Message Batches API, batch ps2 (msgbatch_01NVwr1Yuo8qJf5rykeGBXPK), 2026-11-12; reviewed against both supplied scan halves; see CHECK_DECISIONS.tsv.
\scanpage{134-5/108}
\mspage{295}

cylinder of $0\hand{\to} X$ ),  thus the \var{\xout{notion}}{}{typewriter deletion}{ps_134-5_108-v0001} set of aspheric objects follows
formally from the functor $H^{o}$  (as the objects such that $H^{i}(X) = 0$ for all
$i$), much as in the non commutative set-up the functor $\hand{\pi}_{o}$ for a contractibi-%
lity structure determines the latter,  and hence also the corresponding
asphericity structure.

Thus,  the question arises of formulating the basic structures,
namely \var{\xout{``pure}}{}{typewriter deletion}{ps_134-5_108-v0002} contractibility structures and asphericity structures,  in an
autodual way,applying both to the autodual situation just described,  and
to the non autodual one we have been working out in the notes  -  and if
possible even,  \var{\over{\unsure{in}}{to}}{\hand{to}}{writing over}{ps_134-5_108-v0003} all \var{\xout{similar}}{}{typewriter deletion}{ps_134-5_108-v0004} situations met with so far where a homology
or homotopy \var{\inserted{theory}}{\hand{theory}}{insertion}{ps_134-5_108-v0005} of some kind or other has turned up.  Of course,  it is by no
means sure a priori that we can do so,  by keeping first nicely apart the
two sets of notions (contractibility and asphericity) namely defining
them separately,  and then showing that a contractibility structure de-%
termines an asphericity structure\hand{,} and is determined by the latter.  Maybe
we'll have to define from the outset a richer kind of structure,  were
both ``pure'' homotopy notions and asphericity notions are involved.  Also,
the familiar generalization of mapping cylinders,  namely homotopy
fibers and cofibers,  and the corresponding long exact sequences, will
evidently play an important role in the structure to be described.  Now
this again ties in with the corresponding structure of a derivator,  as
contemplated in section 69,  namely ``integration'' and ``cointegration'' of
diagrams in the given category.

Definitely,  this ref\var{\over{f}{l}}{\hand{l}}{writing over}{ps_134-5_108-v0006}ection is \var{\xout{\unsure{leading to}}}{}{typewriter deletion}{ps_134-5_108-v0007} \var{\inserted[typed]{going to}}{going to}{typed insertion}{ps_134-5_108-v0008} lead well beyond the
scope of the present part IV -  it relates rather to part C) of the
working program envisioned by the end of Mai (section 71,  p.207-210),
\margin{$\times$} and rather belongs to \hand{\uline{\textcolor{black}{part VI}}} of these notes,  which presumably will cen-%
ter around the notion of a derivator.  % Small displaced impressions occur around h in both occurrences of There; they are not read as added text.
There is  however a more technical
question, \var{\inserted{and}}{\hand{and}}{insertion}{ps_134-5_108-v0009} of more limited scope,  which deserves some thought and goes
somewhat in the same direction,  namely:  how to define (for a given
contractibility or asphericity structure on M)\var{\deleted{\illegible{1}}}{}{deletion}{ps_134-5_108-v0010} an\var{\deleted{t}}{}{deletion}{ps_134-5_108-v0011} \hand{\uline{\textcolor{black}{induced}}} structure on
a category $M_{/a}$ ,  where a is in M\var{\deleted{\illegible{1}}}{}{deletion}{ps_134-5_108-v0012} ?   There is a little perplexity in
my mind,  even when M is of the type $C^{\wedge}$ say\hand{,} with C a contractor\hand{,} and a \var{\deleted{is}}{}{deletion}{ps_134-5_108-v0013}
in C ,  taking the canonical structures on $C^{\wedge}$ -  because with the most
evident choice of an``induced''  asphericity structure on $(C^{\wedge})_{/a} \hand{\simeq} (C_{/a})^{\wedge}$,
\var{\xout{this structure}}{}{typewriter deletion}{ps_134-5_108-v0014} namely the usual notion of aspheric objects,  this structu\var{\inserted{re}}{\hand{re}}{insertion}{ps_134-5_108-v0015}

\margin{108}

%% END SOURCE ps_134-5_108
\ifnum\EditionMode=1
\psapparatusstart
\psentry{ps_134-5_108-v0001}{1}{typewriter deletion}{\xout{notion}}{}{The typewriter strikes out this material; it leaves the corrected running text.}
\psentry{ps_134-5_108-v0002}{2}{typewriter deletion}{\xout{``pure}}{}{The typewriter strikes out this material; it leaves the corrected running text.}
\psentry{ps_134-5_108-v0003}{3}{writing over}{\over{\unsure{in}}{to}}{\hand{to}}{The superposed handwritten reading replaces the underlying reading.}
\psentry{ps_134-5_108-v0004}{4}{typewriter deletion}{\xout{similar}}{}{The typewriter strikes out this material; it leaves the corrected running text.}
\psentry{ps_134-5_108-v0005}{5}{insertion}{\inserted{theory}}{\hand{theory}}{The handwritten insertion is incorporated at its indicated position; insertion brackets leave the running text.}
\psentry{ps_134-5_108-v0006}{6}{writing over}{\over{f}{l}}{\hand{l}}{The superposed handwritten reading replaces the underlying reading.}
\psentry{ps_134-5_108-v0007}{7}{typewriter deletion}{\xout{\unsure{leading to}}}{}{The typewriter strikes out this material; it leaves the corrected running text.}
\psentry{ps_134-5_108-v0008}{8}{typed insertion}{\inserted[typed]{going to}}{going to}{The typed insertion is incorporated at its indicated position; insertion brackets leave the running text.}
\psentry{ps_134-5_108-v0009}{9}{insertion}{\inserted{and}}{\hand{and}}{The handwritten insertion is incorporated at its indicated position; insertion brackets leave the running text.}
\psentry{ps_134-5_108-v0010}{10}{deletion}{\deleted{\illegible{1}}}{}{The author cancels this material by hand; it leaves the corrected running text.}
\psentry{ps_134-5_108-v0011}{11}{deletion}{\deleted{t}}{}{The author cancels this material by hand; it leaves the corrected running text.}
\psentry{ps_134-5_108-v0012}{12}{deletion}{\deleted{\illegible{1}}}{}{The author cancels this material by hand; it leaves the corrected running text.}
\psentry{ps_134-5_108-v0013}{13}{deletion}{\deleted{is}}{}{The author cancels this material by hand; it leaves the corrected running text.}
\psentry{ps_134-5_108-v0014}{14}{typewriter deletion}{\xout{this structure}}{}{The typewriter strikes out this material; it leaves the corrected running text.}
\psentry{ps_134-5_108-v0015}{15}{insertion}{\inserted{re}}{\hand{re}}{The handwritten insertion is incorporated at its indicated position; insertion brackets leave the running text.}
\psapparatusend
\fi
\pssource{ps_134-5_109}{134-5/109}{https://grothendieck.umontpellier.fr/web/viewer.html?file=../134-5.pdf\#page=110}
\psnumber{ps_134-5_109-v0001}{1}
\psnumber{ps_134-5_109-v0002}{2}
\psnumber{ps_134-5_109-v0003}{3}
\psnumber{ps_134-5_109-v0004}{4}
\psnumber{ps_134-5_109-v0005}{5}
\psnumber{ps_134-5_109-v0006}{6}
\psnumber{ps_134-5_109-v0007}{7}
\psnumber{ps_134-5_109-v0008}{8}
\psnumber{ps_134-5_109-v0009}{9}
\psnumber{ps_134-5_109-v0010}{10}
\psnumber{ps_134-5_109-v0011}{11}
\psnumber{ps_134-5_109-v0012}{12}
\psnumber{ps_134-5_109-v0013}{13}
\psnumber{ps_134-5_109-v0014}{14}
\psnumber{ps_134-5_109-v0015}{15}
\psnumber{ps_134-5_109-v0016}{16}
\psnumber{ps_134-5_109-v0017}{17}
\psnumber{ps_134-5_109-v0018}{18}
\psnumber{ps_134-5_109-v0019}{19}
\psnumber{ps_134-5_109-v0020}{20}
\psnumber{ps_134-5_109-v0021}{21}
\psnumber{ps_134-5_109-v0022}{22}
\psnumber{ps_134-5_109-v0023}{23}
\psnumber{ps_134-5_109-v0024}{24}
\psnumber{ps_134-5_109-v0025}{25}
\psnumber{ps_134-5_109-v0026}{26}
%% BEGIN SOURCE ps_134-5_109
% Pursuing Stacks (cotes 134-2 to 134-8): archive cote 134-5, site page 109 (PDF page 110), typescript page 296.
% typescript. Draft compared: ps5.tex:234-301.
% AI-transcribed by Claude Opus 5.5 (claude-opus-5-5, adaptive thinking, effort high), Message Batches API, batch ps2 (msgbatch_01NVwr1Yuo8qJf5rykeGBXPK), 2026-11-12; reviewed against both supplied scan halves; see CHECK_DECISIONS.tsv.
\scanpage{134-5/109}
\mspage{296}
% Crowded or displaced impressions occur in contractibility, stronger, object, given, generating and without; earlier individual letters are not securely separable.
will practically never be totally aspheric (unless we take a to be a final
object of C), hence will not be associated to a contractibility structure
-- whereas we expect that the contractibility structure of $C^\wedge$ should indu-%
ce one on $C^\wedge{}_{/a}$ . Presumably, the ``correct'' notions of induced struc-%
ture, in the case of asphericity structure and the corresponding notion
of weak equivalence, should be \var{\xout{substant}}{}{typewriter deletion}{ps_134-5_109-v0001} considerably stronger than the
one I just envisioned, and correspond to the intuition of \hand{\uline{\textcolor{black}{``fiberwise}}}
homotopy types'' over the object a (visualized as a space-like object).
A careful description of such induced structures seems to me to be\var{\xout{ing}}{}{typewriter deletion}{ps_134-5_109-v0002}
needed, and \var{\over{\unsure{a}}{t}}{\hand{t}}{writing over}{ps_134-5_109-v0003}he natural place to be the present part IV of the notes.

4) A little reflection on semi-simplicial homotopy notions (and their
analogons when \hand{$\Delta$} is replaced by a general test category\var{\inserted{$\Delta$}}{\hand{$\Delta$}}{insertion}{ps_134-5_109-v0004}) seems needed,
in order to situate the following fact: ss homotopy notions, namely for
ss objects in any category A, behave well with respect to \hand{\uline{\textcolor{black}{any}}} functor
\[ (1) \qquad A \hand{\longrightarrow} \var{\over{\unsure{A}}{B}}{\hand{B}}{writing over}{ps_134-5_109-v0005} , \]
without having to assume that this functor commutes with finite products,
whereas \var{\xout{the homotopy}}{}{typewriter deletion}{ps_134-5_109-v0006} in the context of homotopy structures, when we
have a functor between categories endowed say with homotopy interval
structures (for instance, with contractibility structures), \var{\xout{this}}{}{typewriter deletion}{ps_134-5_109-v0007} \var{\inserted[typed]{such}}{such}{typed insertion}{ps_134-5_109-v0008} func\hand{tor}
\var{\deleted{\unsure{tor|tion}}}{}{deletion}{ps_134-5_109-v0009}
\[ (2) \qquad M \hand{\longrightarrow} \var{\over{\illegible{1}}{N}}{\hand{N}}{writing over}{ps_134-5_109-v0010} \]
behaves well with respect to homotopy notions only, it would seem, if
we assume beforehand it commutes with finite products (plus, of course,
that it transforms a given generating family of homotopy intervals \var{\xout{into}}{}{typewriter deletion}{ps_134-5_109-v0011} of
M into homotopy intervals\var{\xout{\illegible{2}}}{}{typewriter deletion}{ps_134-5_109-v0012} of $\var{\over{\illegible{1}}{N}}{\hand{N}}{writing over}{ps_134-5_109-v0013}$). In case
\[ M = \underline{Hom}(\hand{\Delta},\var{\over{\illegible{1}}{A}}{\hand{A}}{writing over}{ps_134-5_109-v0014}) \quad , \quad \var{\over{M}{N}}{\hand{N}}{writing over}{ps_134-5_109-v0015} = \underline{Hom}(\hand{\Delta},\var{\over{M}{B}}{\hand{B}}{writing over}{ps_134-5_109-v0016}) \quad , \]
and \var{\inserted{when}}{\hand{when}}{insertion}{ps_134-5_109-v0017} (2) comes from a functor (1) which does not commute to finite pro-%
ducts, \var{\inserted{then}}{\hand{then}}{insertion}{ps_134-5_109-v0018} neither does (2) - and still (2) is well-behaved with respect
to \var{\inserted{ss}}{\hand{ss}}{insertion}{ps_134-5_109-v0019} homotopy notions! It should be noted of course that the semi-simpli-%
cial homotopy notions \var{\inserted[typed]{in M}}{in M}{typed insertion}{ps_134-5_109-v0020} can be defined, even without assuming that in A
finite products exist, namely in situations when M does not admit finite
products - and hence, strictly speaking, the set-up of section 51 does
not apply . \var{\inserted[typed]{All}}{All}{typed insertion}{ps_134-5_109-v0021} \var{\over{T}{t}}{\hand{t}}{writing over}{ps_134-5_109-v0022}his causes a slightly awkward feeling, which I would like
to clarify and see what's going on. I suspect it should be simple enough
to do it here and now.

First, assume that A is stable under finite products,and under
\var{\xout{finite}}{}{typewriter deletion}{ps_134-5_109-v0023} \var{\inserted[typed]{direct}}{direct}{typed insertion}{ps_134-5_109-v0024} sums (with small indexing set say - for what we want to do \var{\inserted{with $\Delta$,}}{\hand{with $\Delta$,}}{insertion}{ps_134-5_109-v0025}
finite direct sums even would be enough). We'll use sums only with sum\hand{-}\var{\inserted{mands}}{\hand{mands}}{insertion}{ps_134-5_109-v0026}

\margin{109}

%% END SOURCE ps_134-5_109
\ifnum\EditionMode=1
\psapparatusstart
\psentry{ps_134-5_109-v0001}{1}{typewriter deletion}{\xout{substant}}{}{The typewriter strikes out this material; it leaves the corrected running text.}
\psentry{ps_134-5_109-v0002}{2}{typewriter deletion}{\xout{ing}}{}{The typewriter strikes out this material; it leaves the corrected running text.}
\psentry{ps_134-5_109-v0003}{3}{writing over}{\over{\unsure{a}}{t}}{\hand{t}}{The superposed handwritten reading replaces the underlying reading.}
\psentry{ps_134-5_109-v0004}{4}{insertion}{\inserted{$\Delta$}}{\hand{$\Delta$}}{The handwritten insertion is incorporated at its indicated position; insertion brackets leave the running text.}
\psentry{ps_134-5_109-v0005}{5}{writing over}{$\over{\unsure{A}}{B}$}{$\hand{B}$}{The superposed handwritten reading replaces the underlying reading.}
\psentry{ps_134-5_109-v0006}{6}{typewriter deletion}{\xout{the homotopy}}{}{The typewriter strikes out this material; it leaves the corrected running text.}
\psentry{ps_134-5_109-v0007}{7}{typewriter deletion}{\xout{this}}{}{The typewriter strikes out this material; it leaves the corrected running text.}
\psentry{ps_134-5_109-v0008}{8}{typed insertion}{\inserted[typed]{such}}{such}{The typed insertion is incorporated at its indicated position; insertion brackets leave the running text.}
\psentry{ps_134-5_109-v0009}{9}{deletion}{\deleted{\unsure{tor|tion}}}{}{The author cancels this material by hand; it leaves the corrected running text.}
\psentry{ps_134-5_109-v0010}{10}{writing over}{$\over{\illegible{1}}{N}$}{$\hand{N}$}{The superposed handwritten reading replaces the underlying reading.}
\psentry{ps_134-5_109-v0011}{11}{typewriter deletion}{\xout{into}}{}{The typewriter strikes out this material; it leaves the corrected running text.}
\psentry{ps_134-5_109-v0012}{12}{typewriter deletion}{\xout{\illegible{2}}}{}{The typewriter strikes out this material; it leaves the corrected running text.}
\psentry{ps_134-5_109-v0013}{13}{writing over}{$\over{\illegible{1}}{N}$}{$\hand{N}$}{The superposed handwritten reading replaces the underlying reading.}
\psentry{ps_134-5_109-v0014}{14}{writing over}{$\over{\illegible{1}}{A}$}{$\hand{A}$}{The superposed handwritten reading replaces the underlying reading.}
\psentry{ps_134-5_109-v0015}{15}{writing over}{$\over{M}{N}$}{$\hand{N}$}{The superposed handwritten reading replaces the underlying reading.}
\psentry{ps_134-5_109-v0016}{16}{writing over}{$\over{M}{B}$}{$\hand{B}$}{The superposed handwritten reading replaces the underlying reading.}
\psentry{ps_134-5_109-v0017}{17}{insertion}{\inserted{when}}{\hand{when}}{The handwritten insertion is incorporated at its indicated position; insertion brackets leave the running text.}
\psentry{ps_134-5_109-v0018}{18}{insertion}{\inserted{then}}{\hand{then}}{The handwritten insertion is incorporated at its indicated position; insertion brackets leave the running text.}
\psentry{ps_134-5_109-v0019}{19}{insertion}{\inserted{ss}}{\hand{ss}}{The handwritten insertion is incorporated at its indicated position; insertion brackets leave the running text.}
\psentry{ps_134-5_109-v0020}{20}{typed insertion}{\inserted[typed]{in M}}{in M}{The typed insertion is incorporated at its indicated position; insertion brackets leave the running text.}
\psentry{ps_134-5_109-v0021}{21}{typed insertion}{\inserted[typed]{All}}{All}{The typed insertion is incorporated at its indicated position; insertion brackets leave the running text.}
\psentry{ps_134-5_109-v0022}{22}{writing over}{\over{T}{t}}{\hand{t}}{The superposed handwritten reading replaces the underlying reading.}
\psentry{ps_134-5_109-v0023}{23}{typewriter deletion}{\xout{finite}}{}{The typewriter strikes out this material; it leaves the corrected running text.}
\psentry{ps_134-5_109-v0024}{24}{typed insertion}{\inserted[typed]{direct}}{direct}{The typed insertion is incorporated at its indicated position; insertion brackets leave the running text.}
\psentry{ps_134-5_109-v0025}{25}{insertion}{\inserted{with $\Delta$,}}{\hand{with $\Delta$,}}{The handwritten insertion is incorporated at its indicated position; insertion brackets leave the running text.}
\psentry{ps_134-5_109-v0026}{26}{insertion}{\inserted{mands}}{\hand{mands}}{The handwritten insertion is incorporated at its indicated position; insertion brackets leave the running text.}
\psapparatusend
\fi
\pssource{ps_134-5_110}{134-5/110}{https://grothendieck.umontpellier.fr/web/viewer.html?file=../134-5.pdf\#page=111}
\psnumber{ps_134-5_110-v0001}{1}
\psnumber{ps_134-5_110-v0002}{2}
\psnumber{ps_134-5_110-v0003}{3}
\psnumber{ps_134-5_110-v0004}{4}
\psnumber{ps_134-5_110-v0005}{5}
\psnumber{ps_134-5_110-v0006}{6}
\psnumber{ps_134-5_110-v0007}{7}
\psnumber{ps_134-5_110-v0008}{8}
\psnumber{ps_134-5_110-v0009}{9}
\psnumber{ps_134-5_110-v0010}{10}
\psnumber{ps_134-5_110-v0011}{11}
\psnumber{ps_134-5_110-v0012}{12}
%% BEGIN SOURCE ps_134-5_110
% Pursuing Stacks (cotes 134-2 to 134-8): archive cote 134-5, site page 110 (PDF page 111), typescript page 297.
% typescript. Draft compared: ps5.tex:284-345.
% AI-transcribed by Claude Opus 5.5 (claude-opus-5-5, adaptive thinking, effort high), Message Batches API, batch ps2 (msgbatch_01NVwr1Yuo8qJf5rykeGBXPK), 2026-11-12; reviewed against both supplied scan halves; see CHECK_DECISIONS.tsv.
\scanpage{134-5/110}
\mspage{297}
% The first displayed I has a doubled impression; no secure earlier letter is recoverable.
equal to a chosen final object e of A,  in order to get a functor
\[ I \hand{\longmapsto} I_A \ : \ \text{(Sets)} \hand{\longrightarrow} A \ , \]
where $I_A$  is the ``constant'' object of A with value I,  namely a sum of I
copies of e (sometimes also written  Ixe ).  Using this functor,  we get
a functor
\[ \hand{\Delta}^{\wedge} = \underline{\mathrm{Hom}}(\hand{\Delta}^{\circ}, \text{(Sets)}) \hand{\longrightarrow} M = \underline{\mathrm{Hom}}(\var{\over{\illegible{1}}{\Delta}}{\hand{\Delta}}{writing over}{ps_134-5_110-v0001}^{\circ} , A) \ , \]
which I d\var{\over[typed]{\illegible{2}}{en}}{en}{typewriter overwriting}{ps_134-5_110-v0002}ote by
\[ K \hand{\longmapsto} K_A \qquad , \]
associating to any ss set the corresponding ``constant'' \var{\inserted[typed]{(relative to A)}}{(relative to A)}{typed insertion}{ps_134-5_110-v0003} ss object of A.
On the other hand,  because A admits finite products,  so does M,  which
enables us to make use of the homotopy notions developped in sections 51
etc.  Thus if
\[ \underline{I} = (I, \hand{\delta}_0, \hand{\delta}_1) \]
is any interval in $\hand{\Delta}^{\wedge}$,  considering the corresponding ``\var{\inserted[typed]{A-}}{A-}{typed insertion}{ps_134-5_110-v0004}constant''  inter-%
val  $\underline{I}_A$  in M,  we get homotopy notions in M,  \unsure{which} we may refer to as
$\underline{I}$ - \hand{\uline{\textcolor{black}{homotopy}}} (dropping the subscript A).  They can all be deduced from the
elementary \var{\inserted[typed]{$\underline{I}$-}}{$\underline{I}$-}{typed insertion}{ps_134-5_110-v0005} homotopy \var{\xout{relation}}{}{typewriter deletion}{ps_134-5_110-v0006} between maps \var{\inserted[typed]{in M}}{in M}{typed insertion}{ps_134-5_110-v0007},  which  is expressed in the
known way,  in terms of a map in  M
\[ h\colon \underline{I}_A \ x \ X \hand{\longrightarrow} Y \quad , \]
where X and Y are the source and target in M of the two considered maps,
between which we want to find an \var{\xout{homotopy}}{}{typewriter deletion}{ps_134-5_110-v0008} elementary $\underline{I}$-homotopy h.  This
map h  decomposes componentwise into
\[ (*\ ) \qquad h_n : \hand{(}I_n\hand{)}_A \ x \ X_n \hand{\longrightarrow} Y_n \quad , \]
and each $h_n$ can be interpreted , in view of the definition of $I_n$ ,  as a
\var{\inserted[typed]{map}}{map}{typed insertion}{ps_134-5_110-v0009} \var{\xout{map}}{}{typewriter deletion}{ps_134-5_110-v0010}
\[ (*\ ') \qquad h'_n \ : \ I_n \hand{\longrightarrow} \mathrm{Hom}_A(X_n,Y_n) \quad , \]
provided we assume that taking products in A is distributive with respect
to the sums we are taking,  whence
\[ I_{n\ A} \ x \ X_n \hand{\simeq} (I_n)_{X_n} \ = \ \text{direct sum of } I_n \text{ copies of } \var{\inserted[typed]{X_n.}}{X_n.}{typed insertion}{ps_134-5_110-v0011} \]
Now to give  h ,  or equivalently a sequence of maps $h_n$ in (*)  ``functo-%
rial in n for variable n'' ,  amounts to giving a sequence of maps (* '),
satisfying a corresponding compatibility relation  for variable n.  The
point of course is that  ( for any I in $\hand{\Delta}^{\wedge}$ and  X,Y  in M)  \hand{\uline{\textcolor{black}{the set of}}}
\hand{\uline{\textcolor{black}{data}}} (*' )  \hand{\uline{\textcolor{black}{plus the compatibility condition make}}}\var{\xout{s}}{}{typewriter deletion}{ps_134-5_110-v0012} \hand{\uline{\textcolor{black}{a sense}}},  \hand{\uline{\textcolor{black}{formally}}},
\hand{\uline{\textcolor{black}{independently of}}}  \hand{\uline{\textcolor{black}{any exactness assumptions on A}}}.  Thus it can be taken

\margin{110}

%% END SOURCE ps_134-5_110
\ifnum\EditionMode=1
\psapparatusstart
\psentry{ps_134-5_110-v0001}{1}{writing over}{$\over{\illegible{1}}{\Delta}$}{$\hand{\Delta}$}{The superposed handwritten reading replaces the underlying reading.}
\psentry{ps_134-5_110-v0002}{2}{typewriter overwriting}{\over[typed]{\illegible{2}}{en}}{en}{The superposed typed reading replaces the underlying reading.}
\psentry{ps_134-5_110-v0003}{3}{typed insertion}{\inserted[typed]{(relative to A)}}{(relative to A)}{The typed insertion is incorporated at its indicated position; insertion brackets leave the running text.}
\psentry{ps_134-5_110-v0004}{4}{typed insertion}{\inserted[typed]{A-}}{A-}{The typed insertion is incorporated at its indicated position; insertion brackets leave the running text.}
\psentry{ps_134-5_110-v0005}{5}{typed insertion}{\inserted[typed]{$\underline{I}$-}}{$\underline{I}$-}{The typed insertion is incorporated at its indicated position; insertion brackets leave the running text.}
\psentry{ps_134-5_110-v0006}{6}{typewriter deletion}{\xout{relation}}{}{The typewriter strikes out this material; it leaves the corrected running text.}
\psentry{ps_134-5_110-v0007}{7}{typed insertion}{\inserted[typed]{in M}}{in M}{The typed insertion is incorporated at its indicated position; insertion brackets leave the running text.}
\psentry{ps_134-5_110-v0008}{8}{typewriter deletion}{\xout{homotopy}}{}{The typewriter strikes out this material; it leaves the corrected running text.}
\psentry{ps_134-5_110-v0009}{9}{typed insertion}{\inserted[typed]{map}}{map}{The typed insertion is incorporated at its indicated position; insertion brackets leave the running text.}
\psentry{ps_134-5_110-v0010}{10}{typewriter deletion}{\xout{map}}{}{The typewriter strikes out this material; it leaves the corrected running text.}
\psentry{ps_134-5_110-v0011}{11}{typed insertion}{$\inserted[typed]{X_n.}$}{$X_n.$}{The typed insertion is incorporated at its indicated position; insertion brackets leave the running text.}
\psentry{ps_134-5_110-v0012}{12}{typewriter deletion}{\xout{s}}{}{The typewriter strikes out this material; it leaves the corrected running text.}
\psapparatusend
\fi
\pssource{ps_134-5_111}{134-5/111}{https://grothendieck.umontpellier.fr/web/viewer.html?file=../134-5.pdf\#page=112}
\psnumber{ps_134-5_111-v0001}{1}
\psnumber{ps_134-5_111-v0002}{2}
\psnumber{ps_134-5_111-v0003}{3}
\psnumber{ps_134-5_111-v0004}{4}
\psnumber{ps_134-5_111-v0005}{5}
\psnumber{ps_134-5_111-v0006}{6}
\psnumber{ps_134-5_111-v0007}{7}
\psnumber{ps_134-5_111-v0008}{8}
\psnumber{ps_134-5_111-v0009}{9}
\psnumber{ps_134-5_111-v0010}{10}
\psnumber{ps_134-5_111-v0011}{11}
\psnumber{ps_134-5_111-v0012}{12}
\psnumber{ps_134-5_111-v0013}{13}
%% BEGIN SOURCE ps_134-5_111
% Pursuing Stacks (cotes 134-2 to 134-8): archive cote 134-5, site page 111 (PDF page 112), typescript page 298.
% typescript. Draft compared: ps5.tex:333-387.
% AI-transcribed by Claude Opus 5.5 (claude-opus-5-5, adaptive thinking, effort high), Message Batches API, batch ps2 (msgbatch_01NVwr1Yuo8qJf5rykeGBXPK), 2026-11-12; reviewed against both supplied scan halves; see CHECK_DECISIONS.tsv.
\scanpage{134-5/111}
\mspage{298}

as the formal ingredient of a definition of ``elementary I-homotopy'' bet-%
ween two maps in M, without any assumptions \var{\(\underset{\hand{\tikz[x=.55em,y=.35em]{\draw (0,0)--(1,1)--(1.6,0)--cycle; \draw (-.1,.15)--(1.8,.8);}}}{\text{whatever}}\)}{whatever}{placement}{ps_134-5_111-v0001} on \var{\xout{M}}{}{typewriter deletion}{ps_134-5_111-v0002} the category
A we start with. The standart case is the one when I = \hand{$\Delta$}$_1$ , the ``unit
interval'', but never mind. The definition works just as well, when \hand{$\Delta$}
is replaced by any (let's say small) category \hand{$\Delta$} whatever. On the other
hand, it is immediate that for a functor M $\hand{\to}$ \var{\over{\unsure{M}}{N}}{\hand{N}}{writing over}{ps_134-5_111-v0003} as above, induced by
a functor A $\hand{\longrightarrow}$ B, \var{\xout{two maps in M which are}}{}{typewriter deletion}{ps_134-5_111-v0004} for two map\hand{s} in M, any elemen-%
tary homotopy between them gives rise to an elementary homotopy \var{\over{\unsure{on}}{of}}{\hand{of}}{writing over}{ps_134-5_111-v0005} their
images in \hand{N} - which \var{\xout{impli}}{}{typewriter deletion}{ps_134-5_111-v0006} is just the well known fact (in case \hand{$\Delta$}=\hand{$\Delta$} ,
I = \hand{$\Delta$}$_1$ ) that M $\hand{\to}$ \hand{N} is compatible with simplicial homotopy notions.

In order to fit this into the general framework of section 51, \var{\xout{the}}{}{typewriter deletion}{ps_134-5_111-v0007}
let's remark that if A is a full subcategory of a category A', then for
a pair of maps in M, the elementary $\underline{I}$-homotopies between these are the
same as when considering the given maps as maps in M' = \uline{Hom}(\hand{$\Delta$} ,A'), in
which M is embedded as a full subcategory. Now any (small, say) catego-%
ry A can be embedded canonically into A' = A\^{} as a full subcategory,
and any functor \var{\xout{\inserted[typed]{f\illegible{1}} (1) A \illegible{1} B embeds}}{}{typewriter deletion}{ps_134-5_111-v0008}
\[ f:\ A \hand{\to} B \]
embeds in the corresponding functor
\[ f_{!} :\ A\hat{} \hand{\to} B\hat{} \quad , \]
hence the functor
\[ \hand{\varphi}:\quad M \hand{\longrightarrow} N \]
embeds in the corresponding functor
\[ \hand{\varphi'}:\quad M' \hand{\longrightarrow} N' \qquad , \quad M' = \underline{Hom}(\hand{\Delta}^{\circ},A\hat{}\,), \ N' = \underline{Hom}(\hand{\Delta}^{\circ},B\hat{}\,) \ . \]
\var{\xout{The functor}}{}{typewriter deletion}{ps_134-5_111-v0009} As A\^{}, B\^{} satisfy the required exactness properties, it fol-%
lows that the $\underline{I}$-homotopy notions in M',N' can be interpreted in terms of
the notions of section 51, with respect to $\underline{I}_A$ and $\underline{I}_B$ , \var{\xout{interpreted now}}{}{typewriter deletion}{ps_134-5_111-v0010}
defined now as \var{\inserted[typed]{componentwise}}{componentwise}{typed insertion}{ps_134-5_111-v0011} constant objects of A\^{} and B\^{} respectively. Still, f$_{!}$
commutes to finite products only if f does, so we are still left with ex-%
plaining why M' $\hand{\to}$ N' is well behaved with respect to $\underline{I}$ - homotopy
notions. Equivalently, we \var{\inserted[typed]{need only}}{need only}{typed insertion}{ps_134-5_111-v0012} \var{\xout{have to}}{}{typewriter deletion}{ps_134-5_111-v0013} see this in the case of (2) M$\hand{\to}$ N,
when A and B are supposed to have the required exactness properties to
allow for the interpretation given above of the $\underline{I}$-homotopy notions in
terms of the formalism of section 51, and when moreover f (as f$_{!}$ above)

\margin{111}

%% END SOURCE ps_134-5_111
\ifnum\EditionMode=1
\psapparatusstart
\psentry{ps_134-5_111-v0001}{1}{placement}{\(\underset{\hand{\tikz[x=.55em,y=.35em]{\draw (0,0)--(1,1)--(1.6,0)--cycle; \draw (-.1,.15)--(1.8,.8);}}}{\text{whatever}}\)}{whatever}{The crossed triangular mark lies below the typed word and has no established lexical function; preserve it diplomatically and leave the word plain in reading.}
\psentry{ps_134-5_111-v0002}{2}{typewriter deletion}{\xout{M}}{}{The typewriter strikes out this material; it leaves the corrected running text.}
\psentry{ps_134-5_111-v0003}{3}{writing over}{\over{\unsure{M}}{N}}{\hand{N}}{The superposed handwritten reading replaces the underlying reading.}
\psentry{ps_134-5_111-v0004}{4}{typewriter deletion}{\xout{two maps in M which are}}{}{The typewriter strikes out this material; it leaves the corrected running text.}
\psentry{ps_134-5_111-v0005}{5}{writing over}{\over{\unsure{on}}{of}}{\hand{of}}{The superposed handwritten reading replaces the underlying reading.}
\psentry{ps_134-5_111-v0006}{6}{typewriter deletion}{\xout{impli}}{}{The typewriter strikes out this material; it leaves the corrected running text.}
\psentry{ps_134-5_111-v0007}{7}{typewriter deletion}{\xout{the}}{}{The typewriter strikes out this material; it leaves the corrected running text.}
\psentry{ps_134-5_111-v0008}{8}{typewriter deletion}{\xout{\inserted[typed]{f\illegible{1}} (1) A \illegible{1} B embeds}}{}{The typewriter strikes out this material; it leaves the corrected running text.}
\psentry{ps_134-5_111-v0009}{9}{typewriter deletion}{\xout{The functor}}{}{The typewriter strikes out this material; it leaves the corrected running text.}
\psentry{ps_134-5_111-v0010}{10}{typewriter deletion}{\xout{interpreted now}}{}{The typewriter strikes out this material; it leaves the corrected running text.}
\psentry{ps_134-5_111-v0011}{11}{typed insertion}{\inserted[typed]{componentwise}}{componentwise}{The typed insertion is incorporated at its indicated position; insertion brackets leave the running text.}
\psentry{ps_134-5_111-v0012}{12}{typed insertion}{\inserted[typed]{need only}}{need only}{The typed insertion is incorporated at its indicated position; insertion brackets leave the running text.}
\psentry{ps_134-5_111-v0013}{13}{typewriter deletion}{\xout{have to}}{}{The typewriter strikes out this material; it leaves the corrected running text.}
\psapparatusend
\fi
\pssource{ps_134-5_112}{134-5/112}{https://grothendieck.umontpellier.fr/web/viewer.html?file=../134-5.pdf\#page=113}
\psnumber{ps_134-5_112-v0001}{1}
\psnumber{ps_134-5_112-v0002}{2}
\psnumber{ps_134-5_112-v0003}{3}
\psnumber{ps_134-5_112-v0004}{4}
\psnumber{ps_134-5_112-v0005}{5}
\psnumber{ps_134-5_112-v0006}{6}
\psnumber{ps_134-5_112-v0007}{7}
\psnumber{ps_134-5_112-v0008}{8}
\psnumber{ps_134-5_112-v0009}{9}
\psnumber{ps_134-5_112-v0010}{10}
\psnumber{ps_134-5_112-v0011}{11}
\psnumber{ps_134-5_112-v0012}{12}
\psnumber{ps_134-5_112-v0013}{13}
\psnumber{ps_134-5_112-v0014}{14}
\psnumber{ps_134-5_112-v0015}{15}
\psnumber{ps_134-5_112-v0016}{16}
\psnumber{ps_134-5_112-v0017}{17}
\psnumber{ps_134-5_112-v0018}{18}
%% BEGIN SOURCE ps_134-5_112
% Pursuing Stacks (cotes 134-2 to 134-8): archive cote 134-5, site page 112 (PDF page 113), typescript page 299.
% typescript. Draft compared: ps5.tex:373-434.
% AI-transcribed by Claude Opus 5.5 (claude-opus-5-5, adaptive thinking, effort high), Message Batches API, batch ps2 (msgbatch_01NVwr1Yuo8qJf5rykeGBXPK), 2026-11-12; reviewed against both supplied scan halves; see CHECK_DECISIONS.tsv.
\scanpage{134-5/112}
\mspage{299}
commutes with sums.  This now is readily expressed by the relations
\[ \varphi'(\underline{I}_A) \hand{\simeq} \underline{I}_B \quad , \]
\[ \varphi'(\underline{I}_A \mathbin{\mathrm{x}} X) \var{\over{\hand{\simeq}}{\xrightarrow{\sim}}}{\hand{\xrightarrow{\sim}}}{writing over}{ps_134-5_112-v0001} \var{\xout{\unsure{\varphi}}}{}{typewriter deletion}{ps_134-5_112-v0002} \varphi'(\underline{I}_A)\mathbin{\mathrm{x}}\varphi'(X) \qquad , \]
i.e.  while $\varphi'$  does \hand{\uline{\textcolor{black}{not}}} commute to finite products in general,  however
it \hand{\uline{\textcolor{black}{does}}} commute to the products which enter in the description of elemen-%
tary homotopies (as these products can be expressed in terms of direct
sums in \var{\xout{M, M'\illegible{5} and \unsure{$\varphi$'}}}{}{typewriter deletion}{ps_134-5_112-v0003} A, B ,  and f commutes to these).

Th\var{\over{i}{e}}{\hand{e}}{writing over}{ps_134-5_112-v0004}s\hand{e} \var{\inserted{reflexions}}{\hand{reflexions}}{insertion}{ps_134-5_112-v0005} suggest\var{\deleted{s}}{}{deletion}{ps_134-5_112-v0006} that the notions of homotopy interval structures and
contractibility structures may be generalized,  in a way that the underly-%
ing category need no longer  be stable under finite products nor even ad-%
mit a final object;  and likewise,  the notion of a mo\hand{r}phism of such struc-%
tures may be generalized,  without assuming that the underlying functor
should commute with finite products. The thought that this kind of genera-%
lization may be needed had already occured before \var{\inserted{in these notes}}{\hand{in these notes}}{insertion}{ps_134-5_112-v0007}, in connection with
the corresponding situation about twenty five years ago,  when the notion
of a site was developped.  But at present the extension does'nt seem ur-%
gent yet,  and I better stop \var{\xout{this}}{}{typewriter deletion}{ps_134-5_112-v0008} here this long digression !

% Displaced impressions occur around h in The and n in One, including below The remaining; no added text is inferred.
The remaining questions  possibly to deal with in part IV  are all
concerned with modelizers.   I'll try to be brief\hand{!}

5)   Consider an ``algebraic structure type'',  and the category \var{\inserted[typed]{M}}{M}{typed insertion}{ps_134-5_112-v0009} of  it's
set-theoretic realizations.  I am looking \var{\over{\unsure{on}}{for}}{\hand{for}}{writing over}{ps_134-5_112-v0010} a comprehensive set of suffi-%
cient conditions on M to insure that M is a``canonical modelizer''.   It seems
natural to assume beforehand that in M  (where at any rate small direct and
inverse limits must exist) internal \hand{\uline{\textcolor{black}{Hom}}}'s exist,  and more generally,
for  X,Y two objects over an object S of M,  that $\underline{\mathrm{Hom}}_S(X,Y)$  exists  -
this implies that base change $S'\hand{\to} S$ in M commutes with small direct limits
and a fortiori,  that direct sums are universal  -  we may as well suppose
them disjoint \var{\inserted{too}}{\hand{too}}{insertion}{ps_134-5_112-v0011}.  One feels quite willing \var{\inserted{too}}{\hand{too}}{insertion}{ps_134-5_112-v0012} to throw in the total 0-connected-%
ness assumption  (cf section 58),  and that every non-empty object has a
section over the final object.    This preliminary set of conditions on an
algebraic structure is of course highly unusual,  however it is satisfied
for \var{\over{M}{m}}{\hand{m}}{writing over}{ps_134-5_112-v0013}ost ``elementary'' algebraic structures\var{\deleted{,}}{}{deletion}{ps_134-5_112-v0014} \hand{(}by which I mean $M\hand{\simeq}A\hat{}$ for some
small category A\hand{)},  as well as for  n-stacks or \hand{$\infty$}-Gr-stacks,  for any n
between \var{\over{\illegible{1}}{0}}{\hand{0}}{writing over}{ps_134-5_112-v0015} and $\mathrm{oo}$ .  The hope \var{\over{\illegible{1}}{now}}{\hand{now}}{writing over}{ps_134-5_112-v0016} is,  in terms of these assumptions,  to
give a necessary and sufficient condition in order that  a)  the ``canonical''
homotopy structure  on M  \var{\over{\unsure{is}}{be}}{\hand{be}}{writing over}{ps_134-5_112-v0017} a contractibility structure,  and moreover
b)  the latter structure \var{\over{\unsure{is}}{be}}{\hand{be}}{writing over}{ps_134-5_112-v0018} ``modelizing'',  by which we mean that the

\margin{112}

%% END SOURCE ps_134-5_112
\ifnum\EditionMode=1
\psapparatusstart
\psentry{ps_134-5_112-v0001}{1}{writing over}{$\over{\hand{\simeq}}{\xrightarrow{\sim}}$}{$\hand{\xrightarrow{\sim}}$}{The superposed handwritten reading replaces the underlying reading.}
\psentry{ps_134-5_112-v0002}{2}{typewriter deletion}{$\xout{\unsure{\varphi}}$}{}{The typewriter strikes out this material; it leaves the corrected running text.}
\psentry{ps_134-5_112-v0003}{3}{typewriter deletion}{\xout{M, M'\illegible{5} and \unsure{$\varphi$'}}}{}{The typewriter strikes out this material; it leaves the corrected running text.}
\psentry{ps_134-5_112-v0004}{4}{writing over}{\over{i}{e}}{\hand{e}}{The superposed handwritten reading replaces the underlying reading.}
\psentry{ps_134-5_112-v0005}{5}{insertion}{\inserted{reflexions}}{\hand{reflexions}}{The handwritten insertion is incorporated at its indicated position; insertion brackets leave the running text.}
\psentry{ps_134-5_112-v0006}{6}{deletion}{\deleted{s}}{}{The author cancels this material by hand; it leaves the corrected running text.}
\psentry{ps_134-5_112-v0007}{7}{insertion}{\inserted{in these notes}}{\hand{in these notes}}{The handwritten insertion is incorporated at its indicated position; insertion brackets leave the running text.}
\psentry{ps_134-5_112-v0008}{8}{typewriter deletion}{\xout{this}}{}{The typewriter strikes out this material; it leaves the corrected running text.}
\psentry{ps_134-5_112-v0009}{9}{typed insertion}{\inserted[typed]{M}}{M}{The typed insertion is incorporated at its indicated position; insertion brackets leave the running text.}
\psentry{ps_134-5_112-v0010}{10}{writing over}{\over{\unsure{on}}{for}}{\hand{for}}{The superposed handwritten reading replaces the underlying reading.}
\psentry{ps_134-5_112-v0011}{11}{insertion}{\inserted{too}}{\hand{too}}{The handwritten insertion is incorporated at its indicated position; insertion brackets leave the running text.}
\psentry{ps_134-5_112-v0012}{12}{insertion}{\inserted{too}}{\hand{too}}{The handwritten insertion is incorporated at its indicated position; insertion brackets leave the running text.}
\psentry{ps_134-5_112-v0013}{13}{writing over}{\over{M}{m}}{\hand{m}}{The superposed handwritten reading replaces the underlying reading.}
\psentry{ps_134-5_112-v0014}{14}{deletion}{\deleted{,}}{}{The author cancels this material by hand; it leaves the corrected running text.}
\psentry{ps_134-5_112-v0015}{15}{writing over}{\over{\illegible{1}}{0}}{\hand{0}}{The superposed handwritten reading replaces the underlying reading.}
\psentry{ps_134-5_112-v0016}{16}{writing over}{\over{\illegible{1}}{now}}{\hand{now}}{The superposed handwritten reading replaces the underlying reading.}
\psentry{ps_134-5_112-v0017}{17}{writing over}{\over{\unsure{is}}{be}}{\hand{be}}{The superposed handwritten reading replaces the underlying reading.}
\psentry{ps_134-5_112-v0018}{18}{writing over}{\over{\unsure{is}}{be}}{\hand{be}}{The superposed handwritten reading replaces the underlying reading.}
\psapparatusend
\fi
\pssource{ps_134-5_113}{134-5/113}{https://grothendieck.umontpellier.fr/web/viewer.html?file=../134-5.pdf\#page=114}
\psnumber{ps_134-5_113-v0001}{1}
\psnumber{ps_134-5_113-v0002}{2}
\psnumber{ps_134-5_113-v0003}{3}
\psnumber{ps_134-5_113-v0004}{4}
\psnumber{ps_134-5_113-v0005}{5}
\psnumber{ps_134-5_113-v0006}{6}
\psnumber{ps_134-5_113-v0007}{7}
\psnumber{ps_134-5_113-v0008}{8}
\psnumber{ps_134-5_113-v0009}{9}
\psnumber{ps_134-5_113-v0010}{10}
\psnumber{ps_134-5_113-v0011}{11}
%% BEGIN SOURCE ps_134-5_113
% Pursuing Stacks (cotes 134-2 to 134-8): archive cote 134-5, site page 113 (PDF page 114), typescript page 300.
% typescript. Draft compared: ps5.tex:421-483.
% AI-transcribed by Claude Opus 5.5 (claude-opus-5-5, adaptive thinking, effort high), Message Batches API, batch ps2 (msgbatch_01NVwr1Yuo8qJf5rykeGBXPK), 2026-11-12; reviewed against both supplied scan halves; see CHECK_DECISIONS.tsv.
\scanpage{134-5/113}
\mspage{300}

associated $\underline{W}_{\mathrm{oo}}$-asphericity structure ($\underline{W}_{\mathrm{oo}}$ = usual weak equivalence)  be
modelizing, which will imply that for \hand{\uline{\textcolor{black}{any}}} basic localizer $\underline{W}$ , the corres%
% NOTE: right edge of the scan cuts the end of this line; a hyphen after "corres" may be hidden
ponding $\underline{W}$-asphericity structure is modelizing.

Even if I do'nt look into this question now, it'll turn up soon
enough in a similar shape, when it comes to prove modelizing properties
for categories of stacks of various kinds. The best we could hope for
would be a statement in terms of the category structure of M alone, \var{\deleted{and}}{}{deletion}{ps_134-5_113-v0001} \var{\inserted{with}}{\hand{with}}{insertion}{ps_134-5_113-v0002}
no \var{\inserted[typed]{assumption}}{assumption}{typed insertion}{ps_134-5_113-v0003} \var{\xout{explicit use of the fact}}{}{typewriter deletion}{ps_134-5_113-v0004} that M be defined in terms of an algebraic
structure type. If I try to formulate anything by way of wishful
testing conjecture, what comes to my mind is: is it enough that there
should exist a separating contractible interval? So the first I would try
to get an idea, is to see how to make a counterexample to this ...

3.7.

89. In connection with the left exactness properties of a $f_{\hand{!}}$ functor\hand{,}
considered three days ago (section 85), I have been befallen by some
doubts whether any subcategory of (Cat) containing the subcategory
of standard simplices is \var{\inserted[typed]{strictly}}{strictly}{typed insertion}{ps_134-5_113-v0005} generating. I wrot e there (p.284) that as
this is true for \hand{$\Delta$} itself, it ``follows a fortiori'' for any subcategory
A of (Cat) containing \hand{$\Delta$}. Assuming A to be full and denoting by i the
inclusion functor, this is known to be equi\var{\over{\illegible{1}}{valent}}{\hand{valent}}{writing over}{ps_134-5_113-v0006} (cf. remark 1 same page)
to $i^{*}$ : (Cat)$\hand{\longrightarrow}$ $A^{\wedge}$ being fully faithful, and in this form, it does'nt
look so obvious that when this is true for one full subcategory, \hand{$A_0$} say,
it should be true for any larger one \var{\inserted{$A$}}{\hand{$A$}}{insertion}{ps_134-5_113-v0007}. This thought had \var{\xout{\illegible{1}k\illegible{3}}}{}{typewriter deletion}{ps_134-5_113-v0008} been linge-%
ring for a second while writing the ``a fortiori'' and I then brushed it
aside, because of the formulation of being generating in terms of strict
epimorphisms. Only the next day did it occur to me that it is by no
means clear that if a family of maps $X_i \hand{\longrightarrow} X$ in a category M is
strictly epimorphic, any larger family with same tar get X should be ``a
fortiori'' strictly epimorphic too - the ``a fortiori'' is known to apply
\var{\inserted{only}}{\hand{only}}{insertion}{ps_134-5_113-v0009} in the case of \var{\xout{amilies which}}{}{typewriter deletion}{ps_134-5_113-v0010} the similar notions of epimorphic, or
universally strictly epimorphic, families of maps. After a little per-%
p\var{\over{\unsure{a}}{l}}{\hand{l}}{writing over}{ps_134-5_113-v0011}exity, I found the situation was saved, in the case I was interested
in, through the fact that it was known from Giraud's article on descent
(Bull. Soc. Math France, Memoire 2, 1964, prop.2.5 p.28) that \hand{$\Delta$},
and even the smaller subcategory of simplices of dimension \hand{$\leqslant$} 2, is even

% Small production overlaps occur in The, should, first, true, for, formulation and known, and near Cat/case; their underlying readings are not independently recoverable.
\margin{113}

%% END SOURCE ps_134-5_113
\ifnum\EditionMode=1
\psapparatusstart
\psentry{ps_134-5_113-v0001}{1}{deletion}{\deleted{and}}{}{The author cancels this material by hand; it leaves the corrected running text.}
\psentry{ps_134-5_113-v0002}{2}{insertion}{\inserted{with}}{\hand{with}}{The handwritten insertion is incorporated at its indicated position; insertion brackets leave the running text.}
\psentry{ps_134-5_113-v0003}{3}{typed insertion}{\inserted[typed]{assumption}}{assumption}{The typed insertion is incorporated at its indicated position; insertion brackets leave the running text.}
\psentry{ps_134-5_113-v0004}{4}{typewriter deletion}{\xout{explicit use of the fact}}{}{The typewriter strikes out this material; it leaves the corrected running text.}
\psentry{ps_134-5_113-v0005}{5}{typed insertion}{\inserted[typed]{strictly}}{strictly}{The typed insertion is incorporated at its indicated position; insertion brackets leave the running text.}
\psentry{ps_134-5_113-v0006}{6}{writing over}{\over{\illegible{1}}{valent}}{\hand{valent}}{The superposed handwritten reading replaces the underlying reading.}
\psentry{ps_134-5_113-v0007}{7}{insertion}{\inserted{$A$}}{\hand{$A$}}{The handwritten insertion is incorporated at its indicated position; insertion brackets leave the running text.}
\psentry{ps_134-5_113-v0008}{8}{typewriter deletion}{\xout{\illegible{1}k\illegible{3}}}{}{The typewriter strikes out this material; it leaves the corrected running text.}
\psentry{ps_134-5_113-v0009}{9}{insertion}{\inserted{only}}{\hand{only}}{The handwritten insertion is incorporated at its indicated position; insertion brackets leave the running text.}
\psentry{ps_134-5_113-v0010}{10}{typewriter deletion}{\xout{amilies which}}{}{The typewriter strikes out this material; it leaves the corrected running text.}
\psentry{ps_134-5_113-v0011}{11}{writing over}{\over{\unsure{a}}{l}}{\hand{l}}{The superposed handwritten reading replaces the underlying reading.}
\psapparatusend
\fi
\pssource{ps_134-5_114}{134-5/114}{https://grothendieck.umontpellier.fr/web/viewer.html?file=../134-5.pdf\#page=115}
\psnumber{ps_134-5_114-v0001}{1}
\psnumber{ps_134-5_114-v0002}{2}
\psnumber{ps_134-5_114-v0003}{3}
\psnumber{ps_134-5_114-v0004}{4}
\psnumber{ps_134-5_114-v0005}{5}
\psnumber{ps_134-5_114-v0006}{6}
\psnumber{ps_134-5_114-v0007}{7}
\psnumber{ps_134-5_114-v0008}{8}
\psnumber{ps_134-5_114-v0009}{9}
\psnumber{ps_134-5_114-v0010}{10}
\psnumber{ps_134-5_114-v0011}{11}
\psnumber{ps_134-5_114-v0012}{12}
\psnumber{ps_134-5_114-v0013}{13}
%% BEGIN SOURCE ps_134-5_114
% Pursuing Stacks (cotes 134-2 to 134-8): archive cote 134-5, site page 114 (PDF page 115), typescript page 301.
% typescript. Draft compared: ps5.tex:470-535.
% AI-transcribed by Claude Opus 5.5 (claude-opus-5-5, adaptive thinking, effort high), Message Batches API, batch ps2 (msgbatch_01NVwr1Yuo8qJf5rykeGBXPK), 2026-11-12; reviewed against both supplied scan halves; see CHECK_DECISIONS.tsv.
\scanpage{134-5/114}
\mspage{301}
\var{\over[typed]{\illegible{3}}{gen}}{gen}{typewriter overwriting}{ps_134-5_114-v0001}erating by \var{\xout{strict}}{}{typewriter deletion}{ps_134-5_114-v0002} ``\hand{\uline{\textcolor{black}{universally}}} strict epimorphisms'',  a notion which
is stable under enlargement of the family of maps,  as recalled above.
Thus the statement made on p.284 does hold true.  And I just checked today
that,  while this stability property by enlargement is surely not \var{\inserted[typed]{always}}{always}{typed insertion}{ps_134-5_114-v0003} true
for a family of maps which is strictly epimorphic,  however,  it \hand{\uline{\textcolor{black}{is}}}
true that if a full subcategory $A_o$ of a c\var{\over[typed]{\illegible{2}}{at}}{at}{typewriter overwriting}{ps_134-5_114-v0004}egory M is generat ing by
strict epimorphisms  (or,  as we'll say,  is ``strictly generating''),
then so is any larger full subcategory A.  This is seen by an easy direct
argument,  in terms of the initial definition,  as meaning that for any
object X in M,  the family of maps  $a_i \hand{\longrightarrow} X$  with target X and source in
the given subcategory (A say) should be strictly epimorphic. (For the
definition of the \var{\xout{various}}{}{typewriter deletion}{ps_134-5_114-v0005} common variants  of the notion of epimorphism,
\margin{$\times$} see the ``Glossaire'' at the end of chapter 1,  SGA 4  vol 1 .)

It occured also to me that  (as suspected in remark 3,  loc. cit)
the functor
\[ i_{!} : \hand{\Delta}^{\wedge} \hand{\longrightarrow} (Cat) \]
coming from the inclusion functor i: \hand{$\Delta$} $\hand{\to}$ (Cat)  is \hand{\uline{\textcolor{black}{not}}}  left exact\var{\unsure{\deleted{,}}}{}{deletion}{ps_134-5_114-v0006}
\hand{(}for another reason though than first contemplated\hand{)}, namely because (Cat)
is known \hand{\uline{\textcolor{black}{not to be a topos}}} \var{\unsure{\deleted{,}}}{}{deletion}{ps_134-5_114-v0007}\hand{(}for instance,  an epimorphism need not
be \hand{\uline{\textcolor{black}{strict}}} (or,  what amounts here to the same,  \hand{\uline{\textcolor{black}{effective}}}) -  as stated
in the cited result of Giraud\hand{)}. Indeed,  we have the following

\hand{\uline{\textcolor{black}{Proposition}}} (which should belong to section 85 !).  Let M be a \underline{U}-category
stable under small direct limits,  A a small full subcategory,  i: A$\hand{\to}$M
the inclusion functor,  hence a functor
\[ i_{!} :\ A^{\wedge} \hand{\longrightarrow} M\ . \]
If A is strictly generating  (i.e.  $i^{*}$: M $\hand{\longrightarrow} A^{\wedge}$ fully faithful),  then \var{\xout{\illegible{1}}}{}{typewriter deletion}{ps_134-5_114-v0008}
$i_{!}$  is left exact \var{\xout{\illegible{1}}}{}{typewriter deletion}{ps_134-5_114-v0009}\var{\inserted[typed]{iff}}{iff}{typed insertion}{ps_134-5_114-v0010} M is a topos.

Indeed,  the inclusion functor \var{\inserted{into $A^{\wedge}$}}{\hand{into $A^{\wedge}$}}{insertion}{ps_134-5_114-v0011} of M' ,  the essential image of $i^{*}$
in $A^{\wedge}$,  \var{\deleted{into $A^{\wedge}$}}{}{deletion}{ps_134-5_114-v0012} admits a left adjoint ($i_{!}$  essentially).  By the criteri-%
on of Giraud,  left exactness of this adjoint,  or equivalently of $i_{!}$ ,
means that M' is the category of sheaves on A for a suitable site struc-%
ture on A,  qed.

\hand{\uline{\textcolor{black}{Corollary}}}:  If M is \hand{\uline{\textcolor{black}{not}}} a topos,  then $i_{!}$  does \hand{\uline{\textcolor{black}{not}}}  commute to fibered
products \var{\xout{of}}{}{typewriter deletion}{ps_134-5_114-v0013} in $A^{\wedge}$ of diagrams of the type
% Hand-drawn arrows are marked blue; typed objects remain plain.
(*) \diagram{\begin{tikzcd}[column sep=tiny,row sep=small] b \ar[dr,color=blue] \& \& c \ar[dl,color=blue] \\ \& F \& \end{tikzcd}} ,
\hand{with b,c in A and F in $A^{\wedge}$,}

\margin{114}

%% END SOURCE ps_134-5_114
\ifnum\EditionMode=1
\psapparatusstart
\psentry{ps_134-5_114-v0001}{1}{typewriter overwriting}{\over[typed]{\illegible{3}}{gen}}{gen}{The superposed typed reading replaces the underlying reading.}
\psentry{ps_134-5_114-v0002}{2}{typewriter deletion}{\xout{strict}}{}{The typewriter strikes out this material; it leaves the corrected running text.}
\psentry{ps_134-5_114-v0003}{3}{typed insertion}{\inserted[typed]{always}}{always}{The typed insertion is incorporated at its indicated position; insertion brackets leave the running text.}
\psentry{ps_134-5_114-v0004}{4}{typewriter overwriting}{\over[typed]{\illegible{2}}{at}}{at}{The superposed typed reading replaces the underlying reading.}
\psentry{ps_134-5_114-v0005}{5}{typewriter deletion}{\xout{various}}{}{The typewriter strikes out this material; it leaves the corrected running text.}
\psentry{ps_134-5_114-v0006}{6}{deletion}{\unsure{\deleted{,}}}{}{All material inside this wrapper is cancelled. Its original doubt and presentation remain in the diplomatic text and apparatus; no empty doubt mark or page label is inserted into the corrected running text.}
\psentry{ps_134-5_114-v0007}{7}{deletion}{\unsure{\deleted{,}}}{}{All material inside this wrapper is cancelled. Its original doubt and presentation remain in the diplomatic text and apparatus; no empty doubt mark or page label is inserted into the corrected running text.}
\psentry{ps_134-5_114-v0008}{8}{typewriter deletion}{\xout{\illegible{1}}}{}{The typewriter strikes out this material; it leaves the corrected running text.}
\psentry{ps_134-5_114-v0009}{9}{typewriter deletion}{\xout{\illegible{1}}}{}{The typewriter strikes out this material; it leaves the corrected running text.}
\psentry{ps_134-5_114-v0010}{10}{typed insertion}{\inserted[typed]{iff}}{iff}{The typed insertion is incorporated at its indicated position; insertion brackets leave the running text.}
\psentry{ps_134-5_114-v0011}{11}{insertion}{\inserted{into $A^{\wedge}$}}{\hand{into $A^{\wedge}$}}{The handwritten insertion is incorporated at its indicated position; insertion brackets leave the running text.}
\psentry{ps_134-5_114-v0012}{12}{deletion}{\deleted{into $A^{\wedge}$}}{}{The author cancels this material by hand; it leaves the corrected running text.}
\psentry{ps_134-5_114-v0013}{13}{typewriter deletion}{\xout{of}}{}{The typewriter strikes out this material; it leaves the corrected running text.}
\psapparatusend
\fi
\pssource{ps_134-5_115}{134-5/115}{https://grothendieck.umontpellier.fr/web/viewer.html?file=../134-5.pdf\#page=116}
\psnumber{ps_134-5_115-v0001}{1}
\psnumber{ps_134-5_115-v0002}{2}
\psnumber{ps_134-5_115-v0003}{3}
\psnumber{ps_134-5_115-v0004}{4}
\psnumber{ps_134-5_115-v0005}{5}
\psnumber{ps_134-5_115-v0006}{6}
\psnumber{ps_134-5_115-v0007}{7}
\psnumber{ps_134-5_115-v0008}{8}
\psnumber{ps_134-5_115-v0009}{9}
\psnumber{ps_134-5_115-v0010}{10}
\psnumber{ps_134-5_115-v0011}{11}
\psnumber{ps_134-5_115-v0012}{12}
\psnumber{ps_134-5_115-v0013}{13}
%% BEGIN SOURCE ps_134-5_115
% Pursuing Stacks (cotes 134-2 to 134-8): archive cote 134-5, site page 115 (PDF page 116), typescript page 302.
% typescript. Draft compared: ps5.tex:517-586.
% AI-transcribed by Claude Opus 5.5 (claude-opus-5-5, adaptive thinking, effort high), Message Batches API, batch ps3 (msgbatch_01XdFmuVVBAcgc7EFxdtCqsj), 2026-11-12; reviewed against both supplied scan halves; see CHECK_DECISIONS.tsv.
\scanpage{134-5/115}
\mspage{302}

while it does commute to finite products, and to fibered \unsure{proddcts|products} of
any two objects of $A^{\wedge}$ over an object of $A$.

The ``while'' comes from prop.1 and prop.2 of section 85, which
imply too that, if $M$ is strictly generating and whether or not $M$ is a
topos, left exactness of $i_{!}$ is equivalent with commutation to fibered
products of the diagrams (*). Hence the corollary.

This corollary answers also the perplexity raised in remark 1
(p.283), as to a hypothetical sharper version of part b) of prop.1 ,
concerning fibered products. As anticipated there, it turns out that
this sharper version is not valid, \var{\inserted[typed]{-not}}{-not}{typed insertion}{ps_134-5_115-v0001} without additional assumptions at
any rate.

\var{\xout{90.}}{}{typewriter deletion}{ps_134-5_115-v0002}
90. After this digression on exactness properties of $f_{!}$ functors,
let\hand{'}s come back to the review of \var{\inserted{those}}{\hand{those}}{insertion}{ps_134-5_115-v0003} questions n\var{\over{\illegible{1}}{ot}}{\hand{ot}}{writing over}{ps_134-5_115-v0004} yet dealt with, which
seem more or less to belong to the present part IV of the notes. We
had stopped two days ago with the question 5) of finding some simple
characterization of canonical modelizers, comparable maybe in simplicity
to the characterization we found \var{\deleted{in part II}}{}{deletion}{ps_134-5_115-v0005} \var{\over{of}{for}}{\hand{for}}{writing over}{ps_134-5_115-v0006} test categories \var{\inserted{(in part II)}}{\hand{(in part II)}}{insertion}{ps_134-5_115-v0007}. This
question may well turn out to be related to the following one.

6) This is the question of finding handy existence theorems for
test functors, whereas so far our attention \var{\deleted{towards}}{}{deletion}{ps_134-5_115-v0008} \var{\inserted{to}}{\hand{to}}{insertion}{ps_134-5_115-v0009} test functors had
been turned towards a thorough understanding of the very notion of a test
functor and it's variants. I have the feeling that, after the reflec-%
tions of sections 78 and 86\var{\unsure{\xout{,}}}{}{typewriter deletion}{ps_134-5_115-v0010} notably, the notion in itself is about un-%
derstood now, so that time is getting ripe for asking for existence
theorems. As all modelizers we have been meeting so far were associated
to asphericity structures, it seems reasonable to restrict to these,
namely to the case of a given modelizing asphericity structure
\[ (M,M_{a}) \qquad , \]
and, if need be, even restrict to the case when this structure is asso
-ciated to a contractibility structure $M_{c}$ . We suppose given moreover
a test category $A$ , which we may (if needed) assume to be strict even,
or even a contractor (i.e. the objects of $A$ in $A^{\wedge}$ are moreover contracti-%
\margin{\{}ble, for the homotopy interval structure \var{\inserted[typed]{in $A^{\wedge}$}}{in $A^{\wedge}$}{typed insertion}{ps_134-5_115-v0011} defined by all intervals coming
from $A$). The question then is whether there exists a test functor
\[ A \hand{\longrightarrow} M \quad . \]
\var{\deleted{which}}{}{deletion}{ps_134-5_115-v0012} \var{\inserted{This}}{\hand{This}}{insertion}{ps_134-5_115-v0013} (under the assumptions made) just reduces to the existence of

\margin{115}

%% END SOURCE ps_134-5_115
\ifnum\EditionMode=1
\psapparatusstart
\psentry{ps_134-5_115-v0001}{1}{typed insertion}{\inserted[typed]{-not}}{-not}{The typed insertion is incorporated at its indicated position; insertion brackets leave the running text.}
\psentry{ps_134-5_115-v0002}{2}{typewriter deletion}{\xout{90.}}{}{The typewriter strikes out this material; it leaves the corrected running text.}
\psentry{ps_134-5_115-v0003}{3}{insertion}{\inserted{those}}{\hand{those}}{The handwritten insertion is incorporated at its indicated position; insertion brackets leave the running text.}
\psentry{ps_134-5_115-v0004}{4}{writing over}{\over{\illegible{1}}{ot}}{\hand{ot}}{The superposed handwritten reading replaces the underlying reading.}
\psentry{ps_134-5_115-v0005}{5}{deletion}{\deleted{in part II}}{}{The author cancels this material by hand; it leaves the corrected running text.}
\psentry{ps_134-5_115-v0006}{6}{writing over}{\over{of}{for}}{\hand{for}}{The superposed handwritten reading replaces the underlying reading.}
\psentry{ps_134-5_115-v0007}{7}{insertion}{\inserted{(in part II)}}{\hand{(in part II)}}{The handwritten insertion is incorporated at its indicated position; insertion brackets leave the running text.}
\psentry{ps_134-5_115-v0008}{8}{deletion}{\deleted{towards}}{}{The author cancels this material by hand; it leaves the corrected running text.}
\psentry{ps_134-5_115-v0009}{9}{insertion}{\inserted{to}}{\hand{to}}{The handwritten insertion is incorporated at its indicated position; insertion brackets leave the running text.}
\psentry{ps_134-5_115-v0010}{10}{typewriter deletion}{\unsure{\xout{,}}}{}{All material inside this wrapper is cancelled. Its original doubt and presentation remain in the diplomatic text and apparatus; no empty doubt mark or page label is inserted into the corrected running text.}
\psentry{ps_134-5_115-v0011}{11}{typed insertion}{\inserted[typed]{in $A^{\wedge}$}}{in $A^{\wedge}$}{The typed insertion is incorporated at its indicated position; insertion brackets leave the running text.}
\psentry{ps_134-5_115-v0012}{12}{deletion}{\deleted{which}}{}{The author cancels this material by hand; it leaves the corrected running text.}
\psentry{ps_134-5_115-v0013}{13}{insertion}{\inserted{This}}{\hand{This}}{The handwritten insertion is incorporated at its indicated position; insertion brackets leave the running text.}
\psapparatusend
\fi
\pssource{ps_134-5_116}{134-5/116}{https://grothendieck.umontpellier.fr/web/viewer.html?file=../134-5.pdf\#page=117}
\psnumber{ps_134-5_116-v0001}{1}
\psnumber{ps_134-5_116-v0002}{2}
\psnumber{ps_134-5_116-v0003}{3}
\psnumber{ps_134-5_116-v0004}{4}
\psnumber{ps_134-5_116-v0005}{5}
\psnumber{ps_134-5_116-v0006}{6}
\psnumber{ps_134-5_116-v0007}{7}
\psnumber{ps_134-5_116-v0008}{8}
\psnumber{ps_134-5_116-v0009}{9}
\psnumber{ps_134-5_116-v0010}{10}
\psnumber{ps_134-5_116-v0011}{11}
\psnumber{ps_134-5_116-v0012}{12}
\psnumber{ps_134-5_116-v0013}{13}
\psnumber{ps_134-5_116-v0014}{14}
\psnumber{ps_134-5_116-v0015}{15}
%% BEGIN SOURCE ps_134-5_116
% Pursuing Stacks (cotes 134-2 to 134-8): archive cote 134-5, site page 116 (PDF page 117), typescript page 303.
% typescript. Draft compared: ps5.tex:572-630.
% AI-transcribed by Claude Opus 5.5 (claude-opus-5-5, adaptive thinking, effort high), Message Batches API, batch ps3 (msgbatch_01XdFmuVVBAcgc7EFxdtCqsj), 2026-11-12; reviewed against both supplied scan halves; see CHECK_DECISIONS.tsv.
\scanpage{134-5/116}
\mspage{303}

a functor which be $M_a$-$\underline{W}$-aspheric.  Here $\underline{W}$  is a given basic localizer,
with respect to which we got an asphericity structure.  The most impor-%
tant case for us surely is the one when $\underline{W}=\underline{W}_{\mathrm{oo}}$ namely usual weak equiva-%
lence.  It is immediate indeed that an $M_a$-$\underline{W}$-\var{\xout{test}}{}{typewriter deletion}{ps_134-5_116-v0001}\var{\inserted[typed]{aspheric}}{aspheric}{typed insertion}{ps_134-5_116-v0002} functor is equally \var{\xout{a}}{}{typewriter deletion}{ps_134-5_116-v0003}
aspheric for the corresponding $\underline{W}'$-asphericity structure of $M$,  for any
basic localizer $\underline{W}'\hand{\supset}\underline{W}$ .  Thus if we got an aspheric functor for $\underline{W}_{\mathrm{oo}}$ ,
the finest basic \unsure{localizer} of all,  we get ipso facto an aspheric func-%
tor for any basic localizer $\underline{W}$ .  (Note also that if an asphericity struc
% NOTE: the end of this line runs off the right edge of the scan after "struc"
ture is modelizing for a given $\underline{W}$,  \var{\xout{it is \illegible{2}}}{}{typewriter deletion}{ps_134-5_116-v0004} the corresponding $\underline{W}'$-struc-%
ture is modelizing too,  for any $\underline{W}'\hand{\supset}\underline{W}$ ;  and the analoguous fact holds
for the notion of a test category  -  namely a $\underline{W}$-test category is also
a $\underline{W}'$ test category, and similarly for \var{\xout{stri}}{}{typewriter deletion}{ps_134-5_116-v0005} total asphericity of $A^\wedge$ and
hence for the condition of being a \hand{\uline{\textcolor{black}{strict}}} test category.)

In case  $M$ is even endowed with a contract\var{\over[typed]{\illegible{1}}{i}}{i}{typewriter overwriting}{ps_134-5_116-v0006}bility structure,  we wil
% NOTE: the end of this line runs off the right edge of the scan after "wil"
be \var{\inserted[typed]{interested}}{interested}{typed insertion}{ps_134-5_116-v0007}, more specifically still,  in aspheric functors factoring not only
through $M_a$ ,  but even through $M_c$ :
\[ \unsure{i:}\ A \hand{\longrightarrow} M_c \quad , \]
\var{\deleted{and}}{}{deletion}{ps_134-5_116-v0008}\var{\inserted{while}}{\hand{while}}{insertion}{ps_134-5_116-v0009} replacing the asphericity requirement on this functor,  by the stron-%
ger one that for  any $x$ in $M_c$ ,  the object $i^*(x)$ in $A^\wedge$ be \hand{\uline{\textcolor{black}{contractible}}}
\hand{(}for the homotopy structure in $A^\wedge$ \var{\xout{associated to}}{}{typewriter deletion}{ps_134-5_116-v0010} defined by homotopy in-%
tervals coming from objects in $A$,  say).  In other words,  we are inte-%
rested in the question of existence of \hand{\uline{\textcolor{black}{bimorphisms}}}  of contractibility
structures  (in the sense of section 86)  from $(M,M_c)$  to  $(A^\wedge,A^\wedge_c\ )$.
It may be noticed that in both cases (working with asphericity structures
or with the contractibility structures instead),  in this existence ques-%
tion,  we may altogether forget $M$ \var{\inserted[typed]{itself}}{itself}{typed insertion}{ps_134-5_116-v0011},  and consider it as an existence ques-%
tion for functors from $A$ into either $M_a$  ,  or $M_c$ ,  with the property
that for any object $x$ in the target category $M_a$ or $M_c$ ,  the object
$i^*(x)$  in $A^\wedge$ \var{\over{is}{be}}{\hand{be}}{writing over}{ps_134-5_116-v0012} either aspheric,  or contractible.  In the \var{\xout{first}}{}{typewriter deletion}{ps_134-5_116-v0013} second
case,  we may even restrict $x$ to be in any \var{\over{\illegible{2}}{gi}}{\hand{gi}}{writing over}{ps_134-5_116-v0014}ven subcategory $C$ of $M_c$
generating the contractibility structure  -  and in the cases met with so
far,  we can find such a $C$ reduced to just one object  $I$ .  In the case
of asphericity structures,  \var{\xout{we have}}{}{typewriter deletion}{ps_134-5_116-v0015} the same holds when taking for $C$ a
subcategory generating the asphericity structure,  provided however $C$
contains the image of $A$ by $A$  (which gives little hope to have $C$ restric-%
ted to just one element ! ).

The interest of finding criteria for existence of $\underline{W}$-aspheric or

\margin{116}

%% END SOURCE ps_134-5_116
\ifnum\EditionMode=1
\psapparatusstart
\psentry{ps_134-5_116-v0001}{1}{typewriter deletion}{\xout{test}}{}{The typewriter strikes out this material; it leaves the corrected running text.}
\psentry{ps_134-5_116-v0002}{2}{typed insertion}{\inserted[typed]{aspheric}}{aspheric}{The typed insertion is incorporated at its indicated position; insertion brackets leave the running text.}
\psentry{ps_134-5_116-v0003}{3}{typewriter deletion}{\xout{a}}{}{The typewriter strikes out this material; it leaves the corrected running text.}
\psentry{ps_134-5_116-v0004}{4}{typewriter deletion}{\xout{it is \illegible{2}}}{}{The typewriter strikes out this material; it leaves the corrected running text.}
\psentry{ps_134-5_116-v0005}{5}{typewriter deletion}{\xout{stri}}{}{The typewriter strikes out this material; it leaves the corrected running text.}
\psentry{ps_134-5_116-v0006}{6}{typewriter overwriting}{\over[typed]{\illegible{1}}{i}}{i}{The superposed typed reading replaces the underlying reading.}
\psentry{ps_134-5_116-v0007}{7}{typed insertion}{\inserted[typed]{interested}}{interested}{The typed insertion is incorporated at its indicated position; insertion brackets leave the running text.}
\psentry{ps_134-5_116-v0008}{8}{deletion}{\deleted{and}}{}{The author cancels this material by hand; it leaves the corrected running text.}
\psentry{ps_134-5_116-v0009}{9}{insertion}{\inserted{while}}{\hand{while}}{The handwritten insertion is incorporated at its indicated position; insertion brackets leave the running text.}
\psentry{ps_134-5_116-v0010}{10}{typewriter deletion}{\xout{associated to}}{}{The typewriter strikes out this material; it leaves the corrected running text.}
\psentry{ps_134-5_116-v0011}{11}{typed insertion}{\inserted[typed]{itself}}{itself}{The typed insertion is incorporated at its indicated position; insertion brackets leave the running text.}
\psentry{ps_134-5_116-v0012}{12}{writing over}{\over{is}{be}}{\hand{be}}{The superposed handwritten reading replaces the underlying reading.}
\psentry{ps_134-5_116-v0013}{13}{typewriter deletion}{\xout{first}}{}{The typewriter strikes out this material; it leaves the corrected running text.}
\psentry{ps_134-5_116-v0014}{14}{writing over}{\over{\illegible{2}}{gi}}{\hand{gi}}{The superposed handwritten reading replaces the underlying reading.}
\psentry{ps_134-5_116-v0015}{15}{typewriter deletion}{\xout{we have}}{}{The typewriter strikes out this material; it leaves the corrected running text.}
\psapparatusend
\fi
\pssource{ps_134-5_117}{134-5/117}{https://grothendieck.umontpellier.fr/web/viewer.html?file=../134-5.pdf\#page=118}
\psnumber{ps_134-5_117-v0001}{1}
\psnumber{ps_134-5_117-v0002}{2}
\psnumber{ps_134-5_117-v0003}{3}
\psnumber{ps_134-5_117-v0004}{4}
\psnumber{ps_134-5_117-v0005}{5}
\psnumber{ps_134-5_117-v0006}{6}
\psnumber{ps_134-5_117-v0007}{7}
\psnumber{ps_134-5_117-v0008}{8}
\psnumber{ps_134-5_117-v0009}{9}
\psnumber{ps_134-5_117-v0010}{10}
\psnumber{ps_134-5_117-v0011}{11}
\psnumber{ps_134-5_117-v0012}{12}
\psnumber{ps_134-5_117-v0013}{13}
\psnumber{ps_134-5_117-v0014}{14}
\psnumber{ps_134-5_117-v0015}{15}
\psnumber{ps_134-5_117-v0016}{16}
\psnumber{ps_134-5_117-v0017}{17}
\psnumber{ps_134-5_117-v0018}{18}
\psnumber{ps_134-5_117-v0019}{19}
\psnumber{ps_134-5_117-v0020}{20}
\psnumber{ps_134-5_117-v0021}{21}
\psnumber{ps_134-5_117-v0022}{22}
%% BEGIN SOURCE ps_134-5_117
% Pursuing Stacks (cotes 134-2 to 134-8): archive cote 134-5, site page 117 (PDF page 118), typescript page 304.
% typescript. Draft compared: ps5.tex:615-683.
% AI-transcribed by Claude Opus 5.5 (claude-opus-5-5, adaptive thinking, effort high), Message Batches API, batch ps3 (msgbatch_01XdFmuVVBAcgc7EFxdtCqsj), 2026-11-12; reviewed against both supplied scan halves; see CHECK_DECISIONS.tsv.
\scanpage{134-5/117}
\mspage{304}

more stringently still\var{\xout{\unsure{s}}}{}{typewriter deletion}{ps_134-5_117-v0001} of \hand{``}c-\hand{\uline{\textcolor{black}{aspheric}}} \hand{\uline{\textcolor{black}{functors}}}\hand{''}(as we may call them\var{\over{\unsure{l}}{)}}{\hand{)}}{writing over}{ps_134-5_117-v0002} is
rather evident,  as it gives a way,  via $i^{*}$ ,  for any homotopy type
described in terms of a ``model'' x in M,  to find a corresponding model
$i^{*}(x)$ in $A^{\wedge}$ .  The situation would be more satisfactory still if we \var{\xout{coul}}{}{typewriter deletion}{ps_134-5_117-v0003}
could find the test functor i \unsure{such|sach} that the corresponding functor
\[ \hand{(1)} \qquad i_{!} : A^{\wedge} \hand{\longrightarrow} M \]
be modelizing too  (assuming M to be stable under small direct limits,
so that $i_{!}$ is defined as the left adjoint of \var{\xout{\illegible{1}}}{}{typewriter deletion}{ps_134-5_117-v0004}
\[ \hand{(2)} \qquad i^{*} : M \hand{\longrightarrow} A^{\wedge} \quad .) \]
In this case,  for a homotopy type described by a model \var{\deleted{\unsure{X}}}{}{deletion}{ps_134-5_117-v0005}\var{\inserted{K}}{\hand{K}}{insertion}{ps_134-5_117-v0006} in $A^{\wedge}$,
$i_{!}(K)$  gives a description of the same by a model in M.

Maybe we should remember though that even if we do not know about
any test functor from A to M,  still we always can find in three steps
a modelizing functor
\[ M \hand{\longrightarrow} A^{\wedge} \ , \]
namely a composition
\[ (*) \qquad M \hand{\xrightarrow{\textcolor{black}{\hand{j}^{*}}}} B^{\wedge} \hand{\xrightarrow{\textcolor{black}{\hand{i}_{B}}}} (\mathrm{Cat}) \hand{\xrightarrow{\textcolor{black}{j_{A} = i_{A}^{*}}}} A^{\wedge} \quad , \]
where  $\hand{j}: B \hand{\to} M$  is an $M_{a}$-aspheric functor from an auxiliary small
category,  which  we may assume to be a \var{\xout{weak}}{}{typewriter deletion}{ps_134-5_117-v0007} test category\hand{,} \var{\xout{i\illegible{1}e\illegible{1} such}}{}{typewriter deletion}{ps_134-5_117-v0008}
\var{\deleted{as}}{}{deletion}{ps_134-5_117-v0009}\var{\inserted{by}}{\hand{by}}{insertion}{ps_134-5_117-v0010} a mild extra assumption on \var{\over{\unsure{M}}{M}}{\hand{M}}{writing over}{ps_134-5_117-v0011} (cf cor. 3  p. 253).  The modelizing
functor we thus get has  the disadvantage of not being left exact,
whereas the looked-for functor $i^{*}$ commutes to small \unsure{dnverse|inverse} limits.
Still,  the composition (*)  is pretty near to being left exact,  it
commutes to fibered products (because $i_{B}$  does)  which is the next best
-  we can view it as a\var{\deleted{n}}{}{deletion}{ps_134-5_117-v0012} \var{\inserted[typed]{left}}{left}{typed insertion}{ps_134-5_117-v0013} exact functor from  M  to  $A^{\wedge}_{/E}$ ,  where \var{\over[typed]{B}{E}}{E}{typewriter overwriting}{ps_134-5_117-v0014} is
the image in $A^{\wedge}$ of the final object of M (assuming \var{\xout{there is}}{}{typewriter deletion}{ps_134-5_117-v0015} $e_{M}$ e\var{\over{\illegible{4}}{xist}}{\hand{xist}}{writing over}{ps_134-5_117-v0016}s).

There  is another advantage still of having a test functor
$i : A \hand{\longrightarrow} M$ , \var{\deleted{with respect to}}{}{deletion}{ps_134-5_117-v0017} \var{\inserted{rather than}}{\hand{rather than}}{insertion}{ps_134-5_117-v0018} merely using (*),  namely it al\var{\xout{b}}{}{typewriter deletion}{ps_134-5_117-v0019}lows us to
``enrich''  the category structure of M,  in such a\var{\xout{\illegible{1}}}{}{typewriter deletion}{ps_134-5_117-v0020} way as to get
``external Hom's''  of objects of M,  with \hand{``}values in $A^{\wedge}$\hand{''},  by defining,
for  x,y  in M ,  the object  \hand{\uline{\textcolor{black}{Hom}}}(A)(x,y)  of $A^{\wedge}$ as
\[ \hand{(3)} \qquad \hand{\underline{\textcolor{black}{\mathrm{Hom}}}}(A)(x,y) \ = \ \hand{\{} a \hand{\longmapsto} \mathrm{Hom}_{M}(i(a) \mathrm{x}\hand{x} \, , \ y) \hand{\}} \ . \]
\var{\xout{This}}{}{typewriter deletion}{ps_134-5_117-v0021} Such enr\var{\xout{a}}{}{typewriter deletion}{ps_134-5_117-v0022}iched structure,  when  $A = \hand{\Delta}$\hand{\unsure{,}}  plays an important part
in the second part of Quillen's treatment of homotopical algebra,  un-%
der the name of \hand{(}semi-\hand{)}simplicial categories,  especially wi th the notion
of \hand{(}semi-\hand{)}simplicial model categories,  which looks quite handy indeed.

\margin{117}

%% END SOURCE ps_134-5_117
\ifnum\EditionMode=1
\psapparatusstart
\psentry{ps_134-5_117-v0001}{1}{typewriter deletion}{\xout{\unsure{s}}}{}{The typewriter strikes out this material; it leaves the corrected running text.}
\psentry{ps_134-5_117-v0002}{2}{writing over}{\over{\unsure{l}}{)}}{\hand{)}}{The superposed handwritten reading replaces the underlying reading.}
\psentry{ps_134-5_117-v0003}{3}{typewriter deletion}{\xout{coul}}{}{The typewriter strikes out this material; it leaves the corrected running text.}
\psentry{ps_134-5_117-v0004}{4}{typewriter deletion}{\xout{\illegible{1}}}{}{The typewriter strikes out this material; it leaves the corrected running text.}
\psentry{ps_134-5_117-v0005}{5}{deletion}{\deleted{\unsure{X}}}{}{The author cancels this material by hand; it leaves the corrected running text.}
\psentry{ps_134-5_117-v0006}{6}{insertion}{\inserted{K}}{\hand{K}}{The handwritten insertion is incorporated at its indicated position; insertion brackets leave the running text.}
\psentry{ps_134-5_117-v0007}{7}{typewriter deletion}{\xout{weak}}{}{The typewriter strikes out this material; it leaves the corrected running text.}
\psentry{ps_134-5_117-v0008}{8}{typewriter deletion}{\xout{i\illegible{1}e\illegible{1} such}}{}{The typewriter strikes out this material; it leaves the corrected running text.}
\psentry{ps_134-5_117-v0009}{9}{deletion}{\deleted{as}}{}{The author cancels this material by hand; it leaves the corrected running text.}
\psentry{ps_134-5_117-v0010}{10}{insertion}{\inserted{by}}{\hand{by}}{The handwritten insertion is incorporated at its indicated position; insertion brackets leave the running text.}
\psentry{ps_134-5_117-v0011}{11}{writing over}{\over{\unsure{M}}{M}}{\hand{M}}{The superposed handwritten reading replaces the underlying reading.}
\psentry{ps_134-5_117-v0012}{12}{deletion}{\deleted{n}}{}{The author cancels this material by hand; it leaves the corrected running text.}
\psentry{ps_134-5_117-v0013}{13}{typed insertion}{\inserted[typed]{left}}{left}{The typed insertion is incorporated at its indicated position; insertion brackets leave the running text.}
\psentry{ps_134-5_117-v0014}{14}{typewriter overwriting}{\over[typed]{B}{E}}{E}{The superposed typed reading replaces the underlying reading.}
\psentry{ps_134-5_117-v0015}{15}{typewriter deletion}{\xout{there is}}{}{The typewriter strikes out this material; it leaves the corrected running text.}
\psentry{ps_134-5_117-v0016}{16}{writing over}{\over{\illegible{4}}{xist}}{\hand{xist}}{The superposed handwritten reading replaces the underlying reading.}
\psentry{ps_134-5_117-v0017}{17}{deletion}{\deleted{with respect to}}{}{The author cancels this material by hand; it leaves the corrected running text.}
\psentry{ps_134-5_117-v0018}{18}{insertion}{\inserted{rather than}}{\hand{rather than}}{The handwritten insertion is incorporated at its indicated position; insertion brackets leave the running text.}
\psentry{ps_134-5_117-v0019}{19}{typewriter deletion}{\xout{b}}{}{The typewriter strikes out this material; it leaves the corrected running text.}
\psentry{ps_134-5_117-v0020}{20}{typewriter deletion}{\xout{\illegible{1}}}{}{The typewriter strikes out this material; it leaves the corrected running text.}
\psentry{ps_134-5_117-v0021}{21}{typewriter deletion}{\xout{This}}{}{The typewriter strikes out this material; it leaves the corrected running text.}
\psentry{ps_134-5_117-v0022}{22}{typewriter deletion}{\xout{a}}{}{The typewriter strikes out this material; it leaves the corrected running text.}
\psapparatusend
\fi
\pssource{ps_134-5_118}{134-5/118}{https://grothendieck.umontpellier.fr/web/viewer.html?file=../134-5.pdf\#page=119}
\psnumber{ps_134-5_118-v0001}{1}
\psnumber{ps_134-5_118-v0002}{2}
\psnumber{ps_134-5_118-v0003}{3}
\psnumber{ps_134-5_118-v0004}{4}
\psnumber{ps_134-5_118-v0005}{5}
\psnumber{ps_134-5_118-v0006}{6}
\psnumber{ps_134-5_118-v0007}{7}
%% BEGIN SOURCE ps_134-5_118
% Pursuing Stacks (cotes 134-2 to 134-8): archive cote 134-5, site page 118 (PDF page 119), typescript page 305.
% typescript. Draft compared: ps5.tex:668-738.
% AI-transcribed by Claude Opus 5.5 (claude-opus-5-5, adaptive thinking, effort high), Message Batches API, batch ps3 (msgbatch_01XdFmuVVBAcgc7EFxdtCqsj), 2026-11-12; reviewed against both supplied scan halves; see CHECK_DECISIONS.tsv.
\scanpage{134-5/118}
\mspage{305}
We should of course define composition of the \hand{\uline{\textcolor{black}{Hom}}}(A)'s ,  as required
too in Quillen's  set up.  This is done by relating the \hand{\uline{\textcolor{black}{Hom}}}(A)'s  to
the well-known interhal  \hand{\uline{\textcolor{black}{Hom}}}'s  in $M^{\wedge}$  -  which will show at the same
time that \var{\deleted{in}}{}{deletion}{ps_134-5_118-v0001}\var{\inserted[typed]{for}}{for}{typed insertion}{ps_134-5_118-v0002} formula (3) to make sense,  we do not really have to assume
that M be stable under binary products,  as we can interprete the pro-%
ducts  $i(a)\mathrm{x}\hand{x}$  as being taken in $M^{\wedge}$,  as we\var{\over{\unsure{ll}}{ll}}{\hand{ll}}{writing over}{ps_134-5_118-v0003} as the Hom ,  so as to
get
\[ \mathrm{Hom}_{M^{\wedge}}(i(a)\mathrm{x}\hand{x} \, , \, y) \hand{\simeq} \mathrm{Hom}_{M^{\wedge}}(i(a), \, \hand{\underline{\textcolor{black}{\mathrm{Hom}}}}_{M^{\wedge}}(x,y)) \quad , \]
hence
\[ (4) \qquad \hand{\underline{\textcolor{black}{\mathrm{Hom}}}}(A)(x,y) \hand{\simeq} i^{*}(\hand{\underline{\textcolor{black}{\mathrm{Hom}}}}_{M^{\wedge}}(x,y)) \quad , \]
where $i^{*}$  in the right hand side is interpreted as a functor
\[ i^{*} \ : \ M^{\wedge} \hand{\longrightarrow} A^{\wedge} \hand{,} \]
rather than $M \hand{\longrightarrow} A^{\wedge}$ .  (I leave to the reader the task of enlarging the
basic Universe,  as need may be ...)  As $i^{*}$ commutes to products,  the
evident composition of t\var{\over{\unsure{h}}{h}}{\hand{h}}{writing over}{ps_134-5_118-v0004}e internal \hand{\uline{\textcolor{black}{Hom}}}'s  in $M^{\wedge}$ gives rise to the
looked-for composition of the \hand{\uline{\textcolor{black}{Hom}}}(A)'s,  with the required associativity
properties.  Of course,  in case M is stable under binary products and
\hand{\uline{\textcolor{black}{Hom}}}'s,  \var{\xout{in the}}{}{typewriter deletion}{ps_134-5_118-v0005} which apparently is going to be the case in all modeli-%
zing situations,  there is no need in the interpretation (4) to intro-%
duce the prohibitively large $M^{\wedge}$,  and we can work in M throughout.

There is an important relation though on the external \hand{\uline{\textcolor{black}{Hom}}}(A)'s
which we would like to be true for a satisfactory formalism,  namely
\[ (5) \qquad \hand{\Gamma}_{A^{\wedge}}(\hand{\underline{\textcolor{black}{\mathrm{Hom}}}}(A)(x,y)) \hand{\xleftarrow{\ \sim\ }} \mathrm{Hom}_{M}(x,y) \quad , \]
where $\hand{\Gamma}_{A^{\wedge}}$ just means  $\mathrm{Hom}_{A^{\wedge}}(e_{A^{\wedge}} \, , \, \ldots)$ .   This is equivalent  to
the requirement
\[ (6) \qquad i_{!}(e_{A^{\wedge}} \, ) \hand{\simeq} e_{M} \]
(assuming a final object $e_{M}$ \var{\xout{\unsure{)}}}{}{typewriter deletion}{ps_134-5_118-v0006}in M to exist),  or equivalently
\[ (7) \qquad i(e_{A}) = e_{M} \]
if we assume moreover $e_{A}$ to exist.   Thus,  it will be natural to \var{\xout{try}}{}{typewriter deletion}{ps_134-5_118-v0007}
ask for test functors satisfying the extra condition (5) or (6)  -
and when trying to construct test functors in various situations
(even without being aware of constructing test functors,  as Mr Jour-%
dain was ``doing prose without knowing it'' ...),  the very first thing
everybody has been doing instinctively was to write down formula (7),

\margin{118}

%% END SOURCE ps_134-5_118
\ifnum\EditionMode=1
\psapparatusstart
\psentry{ps_134-5_118-v0001}{1}{deletion}{\deleted{in}}{}{The author cancels this material by hand; it leaves the corrected running text.}
\psentry{ps_134-5_118-v0002}{2}{typed insertion}{\inserted[typed]{for}}{for}{The typed insertion is incorporated at its indicated position; insertion brackets leave the running text.}
\psentry{ps_134-5_118-v0003}{3}{writing over}{\over{\unsure{ll}}{ll}}{\hand{ll}}{The superposed handwritten reading replaces the underlying reading.}
\psentry{ps_134-5_118-v0004}{4}{writing over}{\over{\unsure{h}}{h}}{\hand{h}}{The superposed handwritten reading replaces the underlying reading.}
\psentry{ps_134-5_118-v0005}{5}{typewriter deletion}{\xout{in the}}{}{The typewriter strikes out this material; it leaves the corrected running text.}
\psentry{ps_134-5_118-v0006}{6}{typewriter deletion}{\xout{\unsure{)}}}{}{The typewriter strikes out this material; it leaves the corrected running text.}
\psentry{ps_134-5_118-v0007}{7}{typewriter deletion}{\xout{try}}{}{The typewriter strikes out this material; it leaves the corrected running text.}
\psapparatusend
\fi
\pssource{ps_134-5_119}{134-5/119}{https://grothendieck.umontpellier.fr/web/viewer.html?file=../134-5.pdf\#page=120}
\psnumber{ps_134-5_119-v0001}{1}
\psnumber{ps_134-5_119-v0002}{2}
\psnumber{ps_134-5_119-v0003}{3}
\psnumber{ps_134-5_119-v0004}{4}
\psnumber{ps_134-5_119-v0005}{5}
\psnumber{ps_134-5_119-v0006}{6}
\psnumber{ps_134-5_119-v0007}{7}
\psnumber{ps_134-5_119-v0008}{8}
\psnumber{ps_134-5_119-v0009}{9}
%% BEGIN SOURCE ps_134-5_119
% Pursuing Stacks (cotes 134-2 to 134-8): archive cote 134-5, site page 119 (PDF page 120), typescript page 306.
% typescript. Draft compared: ps5.tex:724-788.
% AI-transcribed by Claude Opus 5.5 (claude-opus-5-5, adaptive thinking, effort high), Message Batches API, batch ps3 (msgbatch_01XdFmuVVBAcgc7EFxdtCqsj), 2026-11-12; reviewed against both supplied scan halves; see CHECK_DECISIONS.tsv.
\scanpage{134-5/119}
\mspage{306}

I would bet !

The motivation for wanting to find test functors being reasonably
clear by now,  what kind of existence theorems may we hope for ? When
A is such a beautiful test category as \hand{$\Delta$,$\Box$} or \hand{\unsure{$\bigcirc$}} ,  I would expect
that for practically any M endowed with a modelizing contractibility
structure say,  under mild \var{\inserted[typed]{\xout{exactness}}}{}{typed insertion}{ps_134-5_119-v0001} restrictions  (such as the exactness assump-%
tions which are nat\var{\over[typed]{t}{u}}{u}{typewriter overwriting}{ps_134-5_119-v0002}ral in the modelizing story),  there should exist a
test functor indeed.  What I feel less definite about is whether it is
reasonable to expect we can find i even such that $i_{!}$ be modelizing too,
in which case we would expect of course that the pair of \var{\xout{quasi}}{}{typewriter deletion}{ps_134-5_119-v0003} equivalen%
ces of categories
(8)
% Hand-drawn arrows are marked blue; typed objects remain plain.
% The third test-category sign has an internal arc inside its outer loop; the bigcirc approximation is doubtful.
\diagram{\begin{tikzcd} \mathrm{Hot}_A \ar[r, shift left,color=blue] \& \mathrm{Hot}_M \ar[l, shift left,color=blue] \end{tikzcd}}
defined in terms of $i_{!}$ and $i^{*}$ should be quasi-inverse to each other,
and the adjunction maps deduced from those between the functors $i_{!}$ and
$i^{*}$ themselves.  This  in turn is equivalent with the adjunction mor-%
phism
(9)
\[ F \hand{\longrightarrow} i^{*}i_{!}(F) \]
being a weak equivalence,  for any object F in $A^{\wedge}$.  A test functor satis%
fying this exacting extra property merits a name of it's own,  we may
call it a \hand{\uline{\textcolor{black}{perfect}}} \hand{\uline{\textcolor{black}{test}}} \hand{\uline{\textcolor{black}{functor}}}  (or a \hand{\uline{\textcolor{black}{perfect}}} \hand{\uline{\textcolor{black}{aspheric}}} \hand{\uline{\textcolor{black}{functor}}} ,  when
not making any modelizing assumptions on A or M ).  Thus the existence
problem of finding test functors can be sharpened to the one of finding
perfect ones.  Remember though that the \var{\xout{best known}}{}{typewriter deletion}{ps_134-5_119-v0004} most familiar test
functor\var{\xout{s}}{}{typewriter deletion}{ps_134-5_119-v0005} of all (besides the geometric realization functor \hand{$\Delta\to$}(Spaces)\hand{),}
namely the inclusion
\hand{(10)}
\[ \hand{i} : \hand{\Delta} \hand{\longrightarrow} \text{(Cat)} \]
giving rise to the nerve functor (introduced for the first time,  I be-%
lieve,  in a Bourbaki talk of mine\hand{,} on passage to quotient by a preequi-%
valence relation in the category of schemes ...),  is  \hand{\uline{\textcolor{black}{not}}} perfect.
The most natural perfect test functor from \hand{$\Delta$} into (Cat),  more general-%
ly from any \var{\inserted[typed]{weak}}{weak}{typed insertion}{ps_134-5_119-v0006} test category A  into (Cat),  is of course  $i_A$  -  the
functor indeed which has been dominating the whole modelizing picture
in our reflections from the start.  In the case of A = \hand{$\Delta$},  Thomason
discovered another \var{\xout{such}}{}{typewriter deletion}{ps_134-5_119-v0007} perfect test functor,  conceptually less simple
surely,  namely  $i_{!}\mathrm{Sd}^{2}$ ,  where Sd is the ``barycentric subdivision
functor''.  I suspect there \var{\xout{is}}{}{typewriter deletion}{ps_134-5_119-v0008} \var{\inserted[typed]{must be}}{must be}{typed insertion}{ps_134-5_119-v0009} an impressive bunch of perfect test func-%
tors from \hand{$\Delta$} with values in more or less any given modelizer,  not only
the basic one  -  and the question here is to get a clear picture of

\margin{119}

%% END SOURCE ps_134-5_119
\ifnum\EditionMode=1
\psapparatusstart
\psentry{ps_134-5_119-v0001}{1}{typed insertion}{\inserted[typed]{\xout{exactness}}}{}{The typed insertion is incorporated at its indicated position; insertion brackets leave the running text.}
\psentry{ps_134-5_119-v0002}{2}{typewriter overwriting}{\over[typed]{t}{u}}{u}{The superposed typed reading replaces the underlying reading.}
\psentry{ps_134-5_119-v0003}{3}{typewriter deletion}{\xout{quasi}}{}{The typewriter strikes out this material; it leaves the corrected running text.}
\psentry{ps_134-5_119-v0004}{4}{typewriter deletion}{\xout{best known}}{}{The typewriter strikes out this material; it leaves the corrected running text.}
\psentry{ps_134-5_119-v0005}{5}{typewriter deletion}{\xout{s}}{}{The typewriter strikes out this material; it leaves the corrected running text.}
\psentry{ps_134-5_119-v0006}{6}{typed insertion}{\inserted[typed]{weak}}{weak}{The typed insertion is incorporated at its indicated position; insertion brackets leave the running text.}
\psentry{ps_134-5_119-v0007}{7}{typewriter deletion}{\xout{such}}{}{The typewriter strikes out this material; it leaves the corrected running text.}
\psentry{ps_134-5_119-v0008}{8}{typewriter deletion}{\xout{is}}{}{The typewriter strikes out this material; it leaves the corrected running text.}
\psentry{ps_134-5_119-v0009}{9}{typed insertion}{\inserted[typed]{must be}}{must be}{The typed insertion is incorporated at its indicated position; insertion brackets leave the running text.}
\psapparatusend
\fi
\pssource{ps_134-5_120}{134-5/120}{https://grothendieck.umontpellier.fr/web/viewer.html?file=../134-5.pdf\#page=121}
\psnumber{ps_134-5_120-v0001}{1}
\psnumber{ps_134-5_120-v0002}{2}
\psnumber{ps_134-5_120-v0003}{3}
\psnumber{ps_134-5_120-v0004}{4}
\psnumber{ps_134-5_120-v0005}{5}
\psnumber{ps_134-5_120-v0006}{6}
\psnumber{ps_134-5_120-v0007}{7}
\psnumber{ps_134-5_120-v0008}{8}
\psnumber{ps_134-5_120-v0009}{9}
\psnumber{ps_134-5_120-v0010}{10}
\psnumber{ps_134-5_120-v0011}{11}
\psnumber{ps_134-5_120-v0012}{12}
%% BEGIN SOURCE ps_134-5_120
% Pursuing Stacks (cotes 134-2 to 134-8): archive cote 134-5, site page 120 (PDF page 121), typescript page 307.
% typescript. Draft compared: ps5.tex:776-828.
% AI-transcribed by Claude Opus 5.5 (claude-opus-5-5, adaptive thinking, effort high), Message Batches API, batch ps3 (msgbatch_01XdFmuVVBAcgc7EFxdtCqsj), 2026-11-12; reviewed against both supplied scan halves; see CHECK_DECISIONS.tsv.
\scanpage{134-5/120}
\mspage{307}

how to get them, and the same of course for test functors which need not
be perfect, including (10).

Next question then would be to see whether the existence theorems
we may get for \hand{$\Delta$}\hand{,} \var{\xout{and}}{}{typewriter deletion}{ps_134-5_120-v0001}\var{\inserted[typed]{or}}{or}{typed insertion}{ps_134-5_120-v0002} his brothers \hand{$\square$} and \hand{\unsure{$\bigcirc$}} and the like, still hold
true for more or less arbitrary test category, or contractor. If so,
this would be a very strong confirmation of the feeling which has been
prompting the reflections in part II, namely that for the purpose of
having "all purpose"-models for homotopy types (inasfar as this is fea-%
sible), any strict test category, or any contractor at any rate, is
just as good as simplices or cubes, which people have kept working with
for the last twenty-five years. If not, it will be quite interesting
indeed to come to a grasp of what the relevant extra features of \hand{$\Delta$} and
the like are, and how restrictive they are.

I doubt I will dive into these questions, still less come to a
clear picture, in the present part of the notes. Still, before lea-%
ving the topic now, I would like to write down some hints\var{\deleted{\unsure{,}}}{}{deletion}{ps_134-5_120-v0003} I came
upon while doing my scratchwork on homotopy properties of $A\hat{}$ categories.
When looking for functors
\[ A \hand{\longrightarrow} M \]
having some specified properties (such as being a test functor, or a
perfect one, etc), we may view this question as meaning that we are
looking for an object \var{\inserted{with specified properties}}{\hand{with specified properties}}{insertion}{ps_134-5_120-v0004} in the category
\[ M^A \;=\; \hand{\underline{\textcolor{black}{Hom}}}(A,M) \quad . \]
Presumably, this category is endowed with an asphericity o\var{\over{f}{r}}{\hand{r}}{writing over}{ps_134-5_120-v0005} cont\var{\over[typed]{\unsure{t}}{r}}{r}{typewriter overwriting}{ps_134-5_120-v0006}acti-%
bility structure if M is\var{\deleted{,}}{}{deletion}{ps_134-5_120-v0007} \hand{(}as we assume\hand{)}\hand{,} presumably even a modelizing
one. This reminds me that as far as the notion of weak equivalence
goes, there \var{\xout{are}}{}{typewriter deletion}{ps_134-5_120-v0008}\var{\inserted[typed]{may be}}{may be}{typed insertion}{ps_134-5_120-v0009} even several non equivalent ways of finding such struc-%
ture on $M^A$ , one being modelizing, whereas another, more useful one
in some respects, is not. Thus, if M is of the type $B\hat{}$ , we may
rewrite $M^A$ as
\[ M^A \hand{\approxeq} (A^{o}\mathrm{x}B)\hat{} \hand{\approxeq} P\hat{} \quad ; \]
hence \var{\inserted{we get}}{\hand{we get}}{insertion}{ps_134-5_120-v0010} the weak equivalence notion coming from $P\hat{}$ , disregarding it's
product structure, which is modelizing indeed if B is a test functor
and A aspheric, hence $A^{o}\mathrm{x}B$ a test \var{\deleted{functor}}{}{deletion}{ps_134-5_120-v0011}\var{\inserted{category}}{\hand{category}}{insertion}{ps_134-5_120-v0012}. The structure which should
be of more relevance though for our purpose should be a considerably
finer one (namely with smaller set of weak equivalences), which we may

% The third category sign has an inner arc within its outer loop; the doubtful bigcirc is an approximation.
% Looking, having and the M in the first inline M^A have dense production overlaps whose underlying readings are not recoverable.
\margin{120}

%% END SOURCE ps_134-5_120
\ifnum\EditionMode=1
\psapparatusstart
\psentry{ps_134-5_120-v0001}{1}{typewriter deletion}{\xout{and}}{}{The typewriter strikes out this material; it leaves the corrected running text.}
\psentry{ps_134-5_120-v0002}{2}{typed insertion}{\inserted[typed]{or}}{or}{The typed insertion is incorporated at its indicated position; insertion brackets leave the running text.}
\psentry{ps_134-5_120-v0003}{3}{deletion}{\deleted{\unsure{,}}}{}{The author cancels this material by hand; it leaves the corrected running text.}
\psentry{ps_134-5_120-v0004}{4}{insertion}{\inserted{with specified properties}}{\hand{with specified properties}}{The handwritten insertion is incorporated at its indicated position; insertion brackets leave the running text.}
\psentry{ps_134-5_120-v0005}{5}{writing over}{\over{f}{r}}{\hand{r}}{The superposed handwritten reading replaces the underlying reading.}
\psentry{ps_134-5_120-v0006}{6}{typewriter overwriting}{\over[typed]{\unsure{t}}{r}}{r}{The superposed typed reading replaces the underlying reading.}
\psentry{ps_134-5_120-v0007}{7}{deletion}{\deleted{,}}{}{The author cancels this material by hand; it leaves the corrected running text.}
\psentry{ps_134-5_120-v0008}{8}{typewriter deletion}{\xout{are}}{}{The typewriter strikes out this material; it leaves the corrected running text.}
\psentry{ps_134-5_120-v0009}{9}{typed insertion}{\inserted[typed]{may be}}{may be}{The typed insertion is incorporated at its indicated position; insertion brackets leave the running text.}
\psentry{ps_134-5_120-v0010}{10}{insertion}{\inserted{we get}}{\hand{we get}}{The handwritten insertion is incorporated at its indicated position; insertion brackets leave the running text.}
\psentry{ps_134-5_120-v0011}{11}{deletion}{\deleted{functor}}{}{The author cancels this material by hand; it leaves the corrected running text.}
\psentry{ps_134-5_120-v0012}{12}{insertion}{\inserted{category}}{\hand{category}}{The handwritten insertion is incorporated at its indicated position; insertion brackets leave the running text.}
\psapparatusend
\fi
\pssource{ps_134-5_121}{134-5/121}{https://grothendieck.umontpellier.fr/web/viewer.html?file=../134-5.pdf\#page=122}
\psnumber{ps_134-5_121-v0001}{1}
\psnumber{ps_134-5_121-v0002}{2}
\psnumber{ps_134-5_121-v0003}{3}
\psnumber{ps_134-5_121-v0004}{4}
\psnumber{ps_134-5_121-v0005}{5}
\psnumber{ps_134-5_121-v0006}{6}
\psnumber{ps_134-5_121-v0007}{7}
\psnumber{ps_134-5_121-v0008}{8}
\psnumber{ps_134-5_121-v0009}{9}
\psnumber{ps_134-5_121-v0010}{10}
\psnumber{ps_134-5_121-v0011}{11}
%% BEGIN SOURCE ps_134-5_121
% Pursuing Stacks (cotes 134-2 to 134-8): archive cote 134-5, site page 121 (PDF page 122), typescript page 308.
% typescript. Draft compared: ps5.tex:815-869.
% AI-transcribed by Claude Opus 5.5 (claude-opus-5-5, adaptive thinking, effort high), Message Batches API, batch ps3 (msgbatch_01XdFmuVVBAcgc7EFxdtCqsj), 2026-11-12; reviewed against both supplied scan halves; see CHECK_DECISIONS.tsv.
\scanpage{134-5/121}
\mspage{308}

visualize best maybe by writing
\[ M^{A} = (\hand{\underline{\textcolor{black}{Hom}}}(A^{o},M^{o}))\var{\over[typed]{\unsure{x}}{{}^{o}}}{{}^{o}}{typewriter overwriting}{ps_134-5_121-v0001} \quad , \]
i.e. interpreting the dual of $M^{A}$ as the category of A-objects of $M^{o}$ ,
for instance (if $A = \hand{\Delta}$) \var{\inserted{as}}{\hand{as}}{insertion}{ps_134-5_121-v0002} the \var{\xout{cat}}{}{typewriter deletion}{ps_134-5_121-v0003} dual of the category of ss objects of
$M^{o}$ . Now Quillen has given handy conditions, in terms of projectives
of $M^{o} = N$ namely in terms of injectives of M, for the category
$\hand{\underline{\textcolor{black}{Hom}}}(\hand{\Delta}^{o},N)$ of ss objects of N to be a closed model category - hence
the dual category $M^{A}$ will turn out as \var{\over[typed]{\illegible{1}}{a}}{a}{typewriter overwriting}{ps_134-5_121-v0004} closed model category too, un-%
der suitable conditions involving existence of injective objects in M.
These conditions are satisfied for instance when M is a topos, and no-%
tably when M is of the type $B^{\wedge}$ - quite an interesting particular case
indeed ! Assuming that M is stable under both types of limits, so is
$M^{A}$ , hence there is an initial and final object, and according to
Quillen's factorization axiom, the \var{\xout{canonical}}{}{typewriter deletion}{ps_134-5_121-v0005} map from the former to
the latter can be factored through an object
\[ F \text{ in } M^{A} \ , \quad \text{i.e.} \quad F \colon A \hand{\longrightarrow} M \ , \]
which is \hand{\uline{\textcolor{black}{cofibering}}}, and such that $F \hand{\longrightarrow} e$ is a \hand{\uline{\textcolor{black}{trivial fibration}}}.
The idea is that these conditions mean more or less, at any rate imply,
that F is a test functor.

I hit upon this ``way out'', while trying to construct test func-%
tors from $\hand{\Delta}$ to any elementary modelizer $B^{\wedge}$ , in order to try and check
that $B^{\wedge}$ is a (semi)simplicial model category in the sense of Quillen.
The intuitive idea of constructing inductively the components $F_{n}$ of F
was simple enough, still I got stuck \var{\over[typed]{o}{i}}{i}{typewriter overwriting}{ps_134-5_121-v0006}n some messyness and did not try
to push through this way, all the \var{\xout{more}}{}{typewriter deletion}{ps_134-5_121-v0007}\var{\inserted[typed]{less}}{less}{typed insertion}{ps_134-5_121-v0008} as this naive approach had no
chance of generalizing to the case of a more or less general test cate-%
gory A. Of course, for the time being Quillen's theorems\hand{,} about
c\var{\over[typed]{d}{e}}{e}{typewriter overwriting}{ps_134-5_121-v0009}rtain categories $\hand{\underline{\textcolor{black}{Hom}}}(A^{o},N)$ being closed model categories, is equal-%
ly restricted to the case when $A = \hand{\Delta}$ , which looks as usual like a
rather arbitrary assumption. Thus, to ``test'' whether the feeling \var{\xout{abou}}{}{typewriter deletion}{ps_134-5_121-v0010}
about \hand{\uline{\textcolor{black}{any}}} test category more or less being ``just as good'' as $\hand{\Delta}$ , a
second point would be to see whether Quillen's theorems extend, which
presumably is going to be very close to the first point I raised.

I take this occasion to raise a third point - where there is no
reason to restrict to an M which is modelizing (neither was there such
reason before, when phrasing everything in terms of aspheric or c-as-%
pheric functors, rather than test functors ...). Namel\var{\over{\unsure{x}}{y}}{\hand{y}}{writing over}{ps_134-5_121-v0011}, assuming as

% Small production marks occur above and below h of The idea; no legible transposition instruction is recoverable.
\margin{121}

%% END SOURCE ps_134-5_121
\ifnum\EditionMode=1
\psapparatusstart
\psentry{ps_134-5_121-v0001}{1}{typewriter overwriting}{$\over[typed]{\unsure{x}}{{}^{o}}$}{${}^{o}$}{The superposed typed reading replaces the underlying reading.}
\psentry{ps_134-5_121-v0002}{2}{insertion}{\inserted{as}}{\hand{as}}{The handwritten insertion is incorporated at its indicated position; insertion brackets leave the running text.}
\psentry{ps_134-5_121-v0003}{3}{typewriter deletion}{\xout{cat}}{}{The typewriter strikes out this material; it leaves the corrected running text.}
\psentry{ps_134-5_121-v0004}{4}{typewriter overwriting}{\over[typed]{\illegible{1}}{a}}{a}{The superposed typed reading replaces the underlying reading.}
\psentry{ps_134-5_121-v0005}{5}{typewriter deletion}{\xout{canonical}}{}{The typewriter strikes out this material; it leaves the corrected running text.}
\psentry{ps_134-5_121-v0006}{6}{typewriter overwriting}{\over[typed]{o}{i}}{i}{The superposed typed reading replaces the underlying reading.}
\psentry{ps_134-5_121-v0007}{7}{typewriter deletion}{\xout{more}}{}{The typewriter strikes out this material; it leaves the corrected running text.}
\psentry{ps_134-5_121-v0008}{8}{typed insertion}{\inserted[typed]{less}}{less}{The typed insertion is incorporated at its indicated position; insertion brackets leave the running text.}
\psentry{ps_134-5_121-v0009}{9}{typewriter overwriting}{\over[typed]{d}{e}}{e}{The superposed typed reading replaces the underlying reading.}
\psentry{ps_134-5_121-v0010}{10}{typewriter deletion}{\xout{abou}}{}{The typewriter strikes out this material; it leaves the corrected running text.}
\psentry{ps_134-5_121-v0011}{11}{writing over}{\over{\unsure{x}}{y}}{\hand{y}}{The superposed handwritten reading replaces the underlying reading.}
\psapparatusend
\fi
\pssource{ps_134-5_122}{134-5/122}{https://grothendieck.umontpellier.fr/web/viewer.html?file=../134-5.pdf\#page=123}
\psnumber{ps_134-5_122-v0001}{1}
\psnumber{ps_134-5_122-v0002}{2}
\psnumber{ps_134-5_122-v0003}{3}
\psnumber{ps_134-5_122-v0004}{4}
\psnumber{ps_134-5_122-v0005}{5}
\psnumber{ps_134-5_122-v0006}{6}
\psnumber{ps_134-5_122-v0007}{7}
\psnumber{ps_134-5_122-v0008}{8}
\psnumber{ps_134-5_122-v0009}{9}
%% BEGIN SOURCE ps_134-5_122
% Pursuing Stacks (cotes 134-2 to 134-8): archive cote 134-5, site page 122 (PDF page 123), typescript page 309.
% typescript. Draft compared: ps5.tex:856-913.
% AI-transcribed by Claude Opus 5.5 (claude-opus-5-5, adaptive thinking, effort high), Message Batches API, batch ps3 (msgbatch_01XdFmuVVBAcgc7EFxdtCqsj), 2026-11-12; reviewed against both supplied scan halves; see CHECK_DECISIONS.tsv.
\scanpage{134-5/122}
\mspage{309}
above that for any test category or contractor A\var{\xout{x}}{}{typewriter deletion}{ps_134-5_122-v0001} , $M^A$ or $M^{A^o}$ can
be endowed with an asphericity structure or a closed model structure,
or at any rate with a set of weak equivalences, hence a localization
or corresponding ``homotopy category''. What one would expect now is
that up to (canonical) equivalence, the latter does not depend upon
the choice of A, and hence is the same for arbitrary A as when using
\hand{$\Delta$} i.e. simplices. This should be true at any rate \var{\xout{f}}{}{typewriter deletion}{ps_134-5_122-v0002}for $M^{A^o}$ \var{\xout{namely}}{}{typewriter deletion}{ps_134-5_122-v0003}
and when M is a topos - which means that \hand{\uline{\textcolor{black}{Illusie's derived category}}}
D.(X) \hand{\uline{\textcolor{black}{of the category of semisimplicial sheaves on a topos}}} X, \hand{\uline{\textcolor{black}{c\var{\over{an}{ould}}{\hand{ould}}{writing over}{ps_134-5_122-v0004}}}}
\hand{\uline{\textcolor{black}{be constructed by using}}}, \hand{\uline{\textcolor{black}{instead of ss objects}}}, A-\hand{\uline{\textcolor{black}{objects}}} \hand{\uline{\textcolor{black}{of the ca-}}}%
\hand{\uline{\textcolor{black}{tegory of sheaves on}}} X, \hand{\uline{\textcolor{black}{where}}} A \hand{\uline{\textcolor{black}{is any test category}}}. This should be
\margin{\tikz[x=.45em,y=.55em]{\draw (0,2)--(0,1.2)--(.8,1.2); \draw (0,0)--(0,.8)--(2,1.2);}}\hand{\uline{\textcolor{black}{one of}}} the main points to settle in part VII of the notes.

There is a slight discrepancy though between the first point,
about existence of test functors with values in M, depending on a given
asphericity o\var{\over{f}{r}}{\hand{r}}{writing over}{ps_134-5_122-v0005} contractibility structure of M, and the second and the
third, which seem to depend only on the category structure of M. This
is further evidence that the set of questions raised here is still far
from being clear in my mind yet. Stating them \var{\xout{\illegible{1}}}{}{typewriter deletion}{ps_134-5_122-v0006}\var{\inserted[typed]{now}}{now}{typed insertion}{ps_134-5_122-v0007}, however confused-%
ly, is a first step towards clarification !

4.7.

Just still two comments about the existence questions for test
functors, before going over to the next questions in our present re-%
view. One is that for given A,\var{\xout{and M,}}{}{typewriter deletion}{ps_134-5_122-v0008} to prove that for rather general
modelizing M there exists a test functor from A to M, we are reduced to
the case when M is of the type $B^{\wedge}$, where B is a test category - namely
it is enough to take a B such that there exists a test functor $B \hand{\longrightarrow} M$.
If we can even find a perfect test functor from \var{\inserted{some}}{\hand{some}}{insertion}{ps_134-5_122-v0009} B to M, than likewise
the existence question for perfect test functor from A to M is reduced
to the case when $M = B^{\wedge}$. These comments may be useful for applying to
the situation Quillen's closed model theory, as envisioned on previous
page - as his criteria for $\hand{\underline{\textcolor{black}{\mathrm{Hom}}}}(\hand{\Delta}^{o}, N)$ to be a closed model category
apply when N is the dual of a topos, for instance the dual of $B^{\wedge}$.

The second comment is about Thomason's result concerning the standar\unsure{d}
inclusion
\[ i:\ \hand{\Delta} \hand{\hookrightarrow} (\mathrm{Cat}) \quad , \]
which can be expressed by saying that, although i itself is not a \hand{\uline{\textcolor{black}{per-}}}%
\hand{\uline{\textcolor{black}{fect}}} test functor, however for any integer $n \hand{\geq} 2$, the composition
\margin{122}

%% END SOURCE ps_134-5_122
\ifnum\EditionMode=1
\psapparatusstart
\psentry{ps_134-5_122-v0001}{1}{typewriter deletion}{\xout{x}}{}{The typewriter strikes out this material; it leaves the corrected running text.}
\psentry{ps_134-5_122-v0002}{2}{typewriter deletion}{\xout{f}}{}{The typewriter strikes out this material; it leaves the corrected running text.}
\psentry{ps_134-5_122-v0003}{3}{typewriter deletion}{\xout{namely}}{}{The typewriter strikes out this material; it leaves the corrected running text.}
\psentry{ps_134-5_122-v0004}{4}{writing over}{\over{an}{ould}}{\hand{ould}}{The superposed handwritten reading replaces the underlying reading.}
\psentry{ps_134-5_122-v0005}{5}{writing over}{\over{f}{r}}{\hand{r}}{The superposed handwritten reading replaces the underlying reading.}
\psentry{ps_134-5_122-v0006}{6}{typewriter deletion}{\xout{\illegible{1}}}{}{The typewriter strikes out this material; it leaves the corrected running text.}
\psentry{ps_134-5_122-v0007}{7}{typed insertion}{\inserted[typed]{now}}{now}{The typed insertion is incorporated at its indicated position; insertion brackets leave the running text.}
\psentry{ps_134-5_122-v0008}{8}{typewriter deletion}{\xout{and M,}}{}{The typewriter strikes out this material; it leaves the corrected running text.}
\psentry{ps_134-5_122-v0009}{9}{insertion}{\inserted{some}}{\hand{some}}{The handwritten insertion is incorporated at its indicated position; insertion brackets leave the running text.}
\psapparatusend
\fi
\pssource{ps_134-5_123}{134-5/123}{https://grothendieck.umontpellier.fr/web/viewer.html?file=../134-5.pdf\#page=124}
\psnumber{ps_134-5_123-v0001}{1}
\psnumber{ps_134-5_123-v0002}{2}
\psnumber{ps_134-5_123-v0003}{3}
\psnumber{ps_134-5_123-v0004}{4}
\psnumber{ps_134-5_123-v0005}{5}
\psnumber{ps_134-5_123-v0006}{6}
\psnumber{ps_134-5_123-v0007}{7}
\psnumber{ps_134-5_123-v0008}{8}
%% BEGIN SOURCE ps_134-5_123
% Pursuing Stacks (cotes 134-2 to 134-8): archive cote 134-5, site page 123 (PDF page 124), typescript page 310.
% typescript. Draft compared: ps5.tex:899-964.
% AI-transcribed by Claude Opus 5.5 (claude-opus-5-5, adaptive thinking, effort high), Message Batches API, batch ps3 (msgbatch_01XdFmuVVBAcgc7EFxdtCqsj), 2026-11-12; reviewed against both supplied scan halves; see CHECK_DECISIONS.tsv.
\scanpage{134-5/123}
\mspage{310}
\[i_n = i_{\hand{!}}\mathrm{Sd}^n\hand{\alpha} : \hand{\Delta\to\Delta}^{\wedge}\hand{\to\Delta}^{\wedge}\hand{\longrightarrow}(\mathrm{Cat})\]
is a perfect test functor, where \hand{$\alpha\textcolor{black}{:}\Delta\to\Delta^{\textcolor{black}{\wedge}}$} is the canonical inclusion
It is tempting to surmise that this result is not special to i alone,
but that it holds for a large class, if not all, test functors from
\hand{$\Delta$} to asphericity or contractibility modelizers. Here Sd$^n$ denotes the
n.th iterate of the barycentric subdivision functor Sd (following
now the notation in Thomason's paper, which presumably is standard,
while I have been using ``Bar'' in part II of the notes). Presumably \var{\inserted{functors}}{\hand{functors}}{insertion}{ps_134-5_123-v0001} ana-%
loguous \var{\deleted{functors}}{}{deletion}{ps_134-5_123-v0002} to Sd can be defined in any elementary modelizer A$^\wedge$,
as suggested by the natural constructions arising\var{\unsure{\xout{,}}}{}{typewriter deletion}{ps_134-5_123-v0003} in connection with
the factorization property for a closed model structure on A$^\wedge$. Thus
possibly there is a general method in view for deducing perfect test
functors from ordinary one's. However, it definitely seems to me that
the natural place for these existence questions is in part V of the
notes, as they seem intimately related to a\hand{n} \var{\deleted{closed}}{}{deletion}{ps_134-5_123-v0004} understanding of
the homotopy structures of \var{\xout{categories}}{}{typewriter deletion}{ps_134-5_123-v0005} elementary modelizers, and more
\margin{\tikz[x=.45em,y=.55em]{\draw[->,color=blue] (1,0)--(0,0)--(0,2.2);}}specifically to the closed model structures to which such modelizers
give rise in various ways.

\hand{\(\begin{array}{ccc}*& &*\\ &*&\end{array}\)}

\margin{\tikz[x=.45em,y=.55em]{\draw[->,color=blue] (1.5,2)--(0,2)--(0,0);}}Before proceeding, I would like to state still another after\var{\xout{th}}{}{typewriter deletion}{ps_134-5_123-v0006}-%
thought to the reflections of section 89, about strictly generating
subcategories of a category M. I recall that a family of arrows in M
with same target X
\[u_i : X_i \hand{\longrightarrow} X\]
is called \hand{\uline{\textcolor{black}{strictly}}} \hand{\uline{\textcolor{black}{epimorphic}}} , if for every object Y of M, the cor-%
responding map
\[(*)\qquad \mathrm{Hom}(X,Y)\hand{\longrightarrow} \hand{\prod_{i}}\mathrm{Hom}(X_i,Y)\]
is injective (which is expressed by saying that the family $(u_i)$ is
\hand{\uline{\textcolor{black}{epimorphic}}} \var{\over{\unsure{a}}{)}}{\hand{)}}{writing over}{ps_134-5_123-v0007}, and if \hand{\uline{\textcolor{black}{moreover}}} the \var{\xout{image of (*)}}{}{typewriter deletion}{ps_134-5_123-v0008} following obviously
necessary condition for an element $(f_i)$ of the product set of (*) to
be in the image of the map (*), is also sufficient:

(Comp) For any two indices i,j (possibly equal) and any commutat\unsure{i-}%
ve square
% Hand-drawn arrows are marked blue; typed objects remain plain.
\diagram{\begin{tikzcd}[sep=small]
 \& T \ar[dl, color=blue, "\hand{v_i}"'] \ar[dr, color=blue, "\hand{v_j}"] \& \\
X_i \ar[dr, color=blue, "\hand{u_i}"'] \& \& X_j \ar[dl, color=blue, "\hand{u_j}"] \\
 \& X \&
\end{tikzcd}}
in M , the relation $f_iv_i = f_jv_j$ holds.

It is immediate that the condition for $(u_i)$ to be strictly epimorphic

\margin{123}

%% END SOURCE ps_134-5_123
\ifnum\EditionMode=1
\psapparatusstart
\psentry{ps_134-5_123-v0001}{1}{insertion}{\inserted{functors}}{\hand{functors}}{The handwritten insertion is incorporated at its indicated position; insertion brackets leave the running text.}
\psentry{ps_134-5_123-v0002}{2}{deletion}{\deleted{functors}}{}{The author cancels this material by hand; it leaves the corrected running text.}
\psentry{ps_134-5_123-v0003}{3}{typewriter deletion}{\unsure{\xout{,}}}{}{All material inside this wrapper is cancelled. Its original doubt and presentation remain in the diplomatic text and apparatus; no empty doubt mark or page label is inserted into the corrected running text.}
\psentry{ps_134-5_123-v0004}{4}{deletion}{\deleted{closed}}{}{The author cancels this material by hand; it leaves the corrected running text.}
\psentry{ps_134-5_123-v0005}{5}{typewriter deletion}{\xout{categories}}{}{The typewriter strikes out this material; it leaves the corrected running text.}
\psentry{ps_134-5_123-v0006}{6}{typewriter deletion}{\xout{th}}{}{The typewriter strikes out this material; it leaves the corrected running text.}
\psentry{ps_134-5_123-v0007}{7}{writing over}{\over{\unsure{a}}{)}}{\hand{)}}{The superposed handwritten reading replaces the underlying reading.}
\psentry{ps_134-5_123-v0008}{8}{typewriter deletion}{\xout{image of (*)}}{}{The typewriter strikes out this material; it leaves the corrected running text.}
\psapparatusend
\fi
\pssource{ps_134-5_124}{134-5/124}{https://grothendieck.umontpellier.fr/web/viewer.html?file=../134-5.pdf\#page=125}
\psnumber{ps_134-5_124-v0001}{1}
\psnumber{ps_134-5_124-v0002}{2}
\psnumber{ps_134-5_124-v0003}{3}
\psnumber{ps_134-5_124-v0004}{4}
\psnumber{ps_134-5_124-v0005}{5}
\psnumber{ps_134-5_124-v0006}{6}
\psnumber{ps_134-5_124-v0007}{7}
\psnumber{ps_134-5_124-v0008}{8}
\psnumber{ps_134-5_124-v0009}{9}
\psnumber{ps_134-5_124-v0010}{10}
\psnumber{ps_134-5_124-v0011}{11}
\psnumber{ps_134-5_124-v0012}{12}
\psnumber{ps_134-5_124-v0013}{13}
\psnumber{ps_134-5_124-v0014}{14}
%% BEGIN SOURCE ps_134-5_124
% Pursuing Stacks (cotes 134-2 to 134-8): archive cote 134-5, site page 124 (PDF page 125), typescript page 311.
% typescript. Draft compared: ps5.tex:946-1008.
% AI-transcribed by Claude Opus 5.5 (claude-opus-5-5, adaptive thinking, effort high), Message Batches API, batch ps3 (msgbatch_01XdFmuVVBAcgc7EFxdtCqsj), 2026-11-12; reviewed against both supplied scan halves; see CHECK_DECISIONS.tsv.
% The w of when is darkly retyped without a recoverable prior letter.
% The e of The overlaps the adjacent cancelled position; no correction order is assigned.
\scanpage{134-5/124}
\mspage{311}

depends only on the \hand{\uline{\textcolor{black}{sieve}}} $X_o$ of $X$ in $M$ (namely, the subobject of $X$\hand{,} vie-%
wed as an object of $M^\wedge$) generated by the $u_i$'s -- we'll say also that
this \hand{\uline{\textcolor{black}{sieve}}} \hand{\uline{\textcolor{black}{is strictly epimorphic}}}. One should beware that this does not
mean of course that $X_o \hand{\rightarrow} X$ is epimorphic in $M^\wedge$ (which would imply $X_o=X$
i.e. that one of the $u_i$'s admits a section i.e. a right inverse); nor
is it true that if a sieve is \var{\inserted[typed]{strictly}}{strictly}{typed insertion}{ps_134-5_124-v0001} epimorphic, a larger one should be \var{\xout{ep}}{}{typewriter deletion}{ps_134-5_124-v0002}
so too -- which means that when adding more arrows to a strictly epi-%
morphic family, the family need \var{\xout{t\unsure{o}}}{}{typewriter deletion}{ps_134-5_124-v0003} not stay str. ep.

We'll say that the family $(u_i)$ is \hand{\uline{\textcolor{black}{universally}}} \hand{\uline{\textcolor{black}{strictly}}} \hand{\uline{\textcolor{black}{epimorphic}}}
if it is strictly epimorphic, i.e. the corresponding crible $X_o$ is, and
if the latter remains so by arbitrary base change $X' \hand{\rightarrow} X$ in $M$, i.e.
if the corresponding \var{\xout{crible}}{}{typewriter deletion}{ps_134-5_124-v0004} \var{\inserted[typed]{sieve}}{sieve}{typed insertion}{ps_134-5_124-v0005} $X'_o$ \var{\over[typed]{in}{of}}{of}{typewriter overwriting}{ps_134-5_124-v0006} $X'$ is strictly epimorphic too. If the
fibered products \var{\xout{X}}{}{typewriter deletion}{ps_134-5_124-v0007}
\[ X'_i = X_i \mathrm{x}_X X' \]
exist in $M$, this condition also means that the corresponding family
of maps
\[ u'_i \; : \; X'_i \hand{\longrightarrow} X' \]
is strictly epimorphic. The\var{\xout{\illegible{1}}}{}{typewriter deletion}{ps_134-5_124-v0008} condition \var{\xout{aga}}{}{typewriter deletion}{ps_134-5_124-v0009} that $(u_i)$ be \var{\xout{\unsure{s}t}}{}{typewriter deletion}{ps_134-5_124-v0010} univ. str.
epimorphic again depends only o\var{\over[typed]{j}{n}}{n}{typewriter overwriting}{ps_134-5_124-v0011} the generated crible, it is moreover
\hand{\uline{\textcolor{black}{stable under base change}}}, and equally \hand{\uline{\textcolor{black}{stable under adding new arrows}}}
i.e. replacing a sieve in $X$ by a larger one.

It should be noted that if the fibered products \var{\xout{$X_i \mathrm{x}_X \overset{X}{X}_j$}}{}{typewriter deletion}{ps_134-5_124-v0012} $X_i \mathrm{x}_X X_j$
exist in $M$, then the compatibility condition (Comp) above is equiva-%
lent to the one obtained by restricting to \var{\xout{\unsure{T}}}{}{typewriter deletion}{ps_134-5_124-v0013}
\[ T = X_i \mathrm{x}_X X_j \quad , \]
with $v_i$ and $v_j$ the two projections.

\var{\xout{Let}}{}{typewriter deletion}{ps_134-5_124-v0014} Assume the indices $i$ are the objects of a category $I$, and the
$X_i$ are the values of a functor
\[ I \hand{\longrightarrow} M \quad , \]
and that the family of arrows $(u_i)$ turns $X$ into the direct limit in $M$
of the $X_i$ :
\[ X = \underset{i}{\underset{\hand{\longrightarrow}}{\mathrm{lim}}} X_i \quad , \]
then it follows immediately that the family $(u_i)$ is strictly epimorphic\hand{.}

After these terminological preliminaries, we're ready to give the

\margin{124}

%% END SOURCE ps_134-5_124
\ifnum\EditionMode=1
\psapparatusstart
\psentry{ps_134-5_124-v0001}{1}{typed insertion}{\inserted[typed]{strictly}}{strictly}{The typed insertion is incorporated at its indicated position; insertion brackets leave the running text.}
\psentry{ps_134-5_124-v0002}{2}{typewriter deletion}{\xout{ep}}{}{The typewriter strikes out this material; it leaves the corrected running text.}
\psentry{ps_134-5_124-v0003}{3}{typewriter deletion}{\xout{t\unsure{o}}}{}{The typewriter strikes out this material; it leaves the corrected running text.}
\psentry{ps_134-5_124-v0004}{4}{typewriter deletion}{\xout{crible}}{}{The typewriter strikes out this material; it leaves the corrected running text.}
\psentry{ps_134-5_124-v0005}{5}{typed insertion}{\inserted[typed]{sieve}}{sieve}{The typed insertion is incorporated at its indicated position; insertion brackets leave the running text.}
\psentry{ps_134-5_124-v0006}{6}{typewriter overwriting}{\over[typed]{in}{of}}{of}{The superposed typed reading replaces the underlying reading.}
\psentry{ps_134-5_124-v0007}{7}{typewriter deletion}{\xout{X}}{}{The typewriter strikes out this material; it leaves the corrected running text.}
\psentry{ps_134-5_124-v0008}{8}{typewriter deletion}{\xout{\illegible{1}}}{}{The typewriter strikes out this material; it leaves the corrected running text.}
\psentry{ps_134-5_124-v0009}{9}{typewriter deletion}{\xout{aga}}{}{The typewriter strikes out this material; it leaves the corrected running text.}
\psentry{ps_134-5_124-v0010}{10}{typewriter deletion}{\xout{\unsure{s}t}}{}{The typewriter strikes out this material; it leaves the corrected running text.}
\psentry{ps_134-5_124-v0011}{11}{typewriter overwriting}{\over[typed]{j}{n}}{n}{The superposed typed reading replaces the underlying reading.}
\psentry{ps_134-5_124-v0012}{12}{typewriter deletion}{\xout{$X_i \mathrm{x}_X \overset{X}{X}_j$}}{}{The typewriter strikes out this material; it leaves the corrected running text.}
\psentry{ps_134-5_124-v0013}{13}{typewriter deletion}{\xout{\unsure{T}}}{}{The typewriter strikes out this material; it leaves the corrected running text.}
\psentry{ps_134-5_124-v0014}{14}{typewriter deletion}{\xout{Let}}{}{The typewriter strikes out this material; it leaves the corrected running text.}
\psapparatusend
\fi
\pssource{ps_134-5_125}{134-5/125}{https://grothendieck.umontpellier.fr/web/viewer.html?file=../134-5.pdf\#page=126}
\psnumber{ps_134-5_125-v0001}{1}
\psnumber{ps_134-5_125-v0002}{2}
\psnumber{ps_134-5_125-v0003}{3}
\psnumber{ps_134-5_125-v0004}{4}
\psnumber{ps_134-5_125-v0005}{5}
\psnumber{ps_134-5_125-v0006}{6}
\psnumber{ps_134-5_125-v0007}{7}
\psnumber{ps_134-5_125-v0008}{8}
\psnumber{ps_134-5_125-v0009}{9}
\psnumber{ps_134-5_125-v0010}{10}
%% BEGIN SOURCE ps_134-5_125
% Pursuing Stacks (cotes 134-2 to 134-8): archive cote 134-5, site page 125 (PDF page 126), typescript page 312.
% typescript. Draft compared: ps5.tex:987-1065.
% AI-transcribed by Claude Opus 5.5 (claude-opus-5-5, adaptive thinking, effort high), Message Batches API, batch ps3 (msgbatch_01XdFmuVVBAcgc7EFxdtCqsj), 2026-11-12; reviewed against both supplied scan halves; see CHECK_DECISIONS.tsv.
% Hand-drawn arrows are marked blue; typed objects remain plain.
\scanpage{134-5/125}
\mspage{312}

following useful statement,  which is lacking in SGA IV Chap.I
(compare loc. cit.  prop.7.2 , page 47,  giving part of the story):
\hand{\uline{\textcolor{black}{Proposition}}}  Let M be a U-category,  A a small \var{\inserted[typed]{full}}{full}{typed insertion}{ps_134-5_125-v0001} subcategory,  i: A $\hand{\to}$ M
the inclusion functor,  hence a functor
\[ i^{*} : M \hand{\longrightarrow} A^{\wedge} \quad . \]
For any object X of M,  we consider the family \var{\xout{\unsure{$\Phi$}}}{}{typewriter deletion}{ps_134-5_125-v0002} $F_X$  of all arrows
in M with target X,  source in A.  The following conditions are equiva-%
lent:

(i)  For any X in M,  the family $F_X$ is strictly epimorphic.

(ii) For any X in M,  the family $F_X$ is universally strictly epimor-%
phic.

(iii)  For any X in M ,  $F_X$ \var{\xout{induces an isomorphism}}{}{typewriter deletion}{ps_134-5_125-v0003}\var{\inserted[typed]{turns X into a direct}}{turns X into a direct}{typed insertion}{ps_134-5_125-v0004}  limit of the
composition functor  $A_{/X} \hand{\to} A \hand{\to} M$ ,  i.e.
\[ X \hand{\xleftarrow{\sim}} \underset{A/X}{\underset{\hand{\longrightarrow}}{\mathrm{lim}}} a \qquad . \]
(iv)  The functor $i^{*}$ is fully faithful.

Proof left to the reader (who may consult loc. cit. for \var{\xout{(i)}}{}{typewriter deletion}{ps_134-5_125-v0005}
(i)$\hand{\Longrightarrow}$(iii)$\hand{\Longleftrightarrow}$(iv) ,  so that only (i)$\hand{\Rightarrow}$(ii) is left to prove).

\hand{\uline{\textcolor{black}{Definition}}}  When the equivalent conditions above are satisfied, we'll
say that A is a \hand{\uline{\textcolor{black}{strictly generating}}} subcategory of M.

NB. The notion makes a sense too without assuming  A to be small,
nor M to be a U-category,  by passing to a larger universe  (it is imme-%
on (i) or (iii)  that the condition does not depend on the choice of the
universe\hand{)}.

\hand{\uline{\textcolor{black}{\var{\over[typed]{O}{C}}{C}{typewriter overwriting}{ps_134-5_125-v0006}orollary}}}  If A is strictly generating in M,  then so is any larger
full subcategory B.

This is clear by criterion (ii)  (whereas it is'nt by any one of
\margin{\tikz[x=.45em,y=.55em]{\draw[->,color=blue] (1.6,0)--(0,0)--(0,2.2);}}the other three criteria !).

Digressions are ended,  let's resume our review!

\var{\xout{\unsure{91}}}{}{typewriter deletion}{ps_134-5_125-v0007}
\margin{\tikz[x=.45em,y=.55em]{\draw[->,color=blue] (1.8,2)--(0,2)--(0,0);}}91.   I see three more questions to review -  presumably \var{\deleted{it}}{}{deletion}{ps_134-5_125-v0008}\var{\inserted{they}}{\hand{they}}{insertion}{ps_134-5_125-v0009} will be
a lot shorter than the last !

7)  \hand{\uline{\textcolor{black}{Description of homotopy types}}} ``\hand{\uline{\textcolor{black}{of finite type}}}''\hand{,} in terms of an
elementary modelizer $A^{\wedge}$.  \var{In\hand{\tikz[x=.22em,y=.55em,baseline=-.2ex]{\draw (0,1.4)--(1,-.5);}}terms}{In terms}{placement}{ps_134-5_125-v0010} of the modelizer (Spaces), a natural
finiteness condition on a homotopy type is that it may be described (up

\margin{125}

%% END SOURCE ps_134-5_125
\ifnum\EditionMode=1
\psapparatusstart
\psentry{ps_134-5_125-v0001}{1}{typed insertion}{\inserted[typed]{full}}{full}{The typed insertion is incorporated at its indicated position; insertion brackets leave the running text.}
\psentry{ps_134-5_125-v0002}{2}{typewriter deletion}{\xout{\unsure{$\Phi$}}}{}{The typewriter strikes out this material; it leaves the corrected running text.}
\psentry{ps_134-5_125-v0003}{3}{typewriter deletion}{\xout{induces an isomorphism}}{}{The typewriter strikes out this material; it leaves the corrected running text.}
\psentry{ps_134-5_125-v0004}{4}{typed insertion}{\inserted[typed]{turns X into a direct}}{turns X into a direct}{The typed insertion is incorporated at its indicated position; insertion brackets leave the running text.}
\psentry{ps_134-5_125-v0005}{5}{typewriter deletion}{\xout{(i)}}{}{The typewriter strikes out this material; it leaves the corrected running text.}
\psentry{ps_134-5_125-v0006}{6}{typewriter overwriting}{\over[typed]{O}{C}}{C}{The superposed typed reading replaces the underlying reading.}
\psentry{ps_134-5_125-v0007}{7}{typewriter deletion}{\xout{\unsure{91}}}{}{The typewriter strikes out this material; it leaves the corrected running text.}
\psentry{ps_134-5_125-v0008}{8}{deletion}{\deleted{it}}{}{The author cancels this material by hand; it leaves the corrected running text.}
\psentry{ps_134-5_125-v0009}{9}{insertion}{\inserted{they}}{\hand{they}}{The handwritten insertion is incorporated at its indicated position; insertion brackets leave the running text.}
\psentry{ps_134-5_125-v0010}{10}{placement}{In\hand{\tikz[x=.22em,y=.55em,baseline=-.2ex]{\draw (0,1.4)--(1,-.5);}}terms}{In terms}{A long slanted pen separator lies at the boundary of the typed words In and terms; the reading preserves the two words with a normal space.}
\psapparatusend
\fi
\pssource{ps_134-5_126}{134-5/126}{https://grothendieck.umontpellier.fr/web/viewer.html?file=../134-5.pdf\#page=127}
\psnumber{ps_134-5_126-v0001}{1}
\psnumber{ps_134-5_126-v0002}{2}
\psnumber{ps_134-5_126-v0003}{3}
\psnumber{ps_134-5_126-v0004}{4}
\psnumber{ps_134-5_126-v0005}{5}
\psnumber{ps_134-5_126-v0006}{6}
\psnumber{ps_134-5_126-v0007}{7}
\psnumber{ps_134-5_126-v0008}{8}
\psnumber{ps_134-5_126-v0009}{9}
\psnumber{ps_134-5_126-v0010}{10}
\psnumber{ps_134-5_126-v0011}{11}
\psnumber{ps_134-5_126-v0012}{12}
\psnumber{ps_134-5_126-v0013}{13}
\psnumber{ps_134-5_126-v0014}{14}
\psnumber{ps_134-5_126-v0015}{15}
\psnumber{ps_134-5_126-v0016}{16}
%% BEGIN SOURCE ps_134-5_126
% Pursuing Stacks (cotes 134-2 to 134-8): archive cote 134-5, site page 126 (PDF page 127), typescript page 313.
% typescript. Draft compared: ps5.tex:1051-1108.
% AI-transcribed by Claude Opus 5.5 (claude-opus-5-5, adaptive thinking, effort high), Message Batches API, batch ps3 (msgbatch_01XdFmuVVBAcgc7EFxdtCqsj), 2026-11-12; reviewed against both supplied scan halves; see CHECK_DECISIONS.tsv.
% Two neighboring m/o impressions in the opening word remain unresolved.
% A separate pen h lies below the typed h of The; two small upper traces have no assigned function.
\scanpage{134-5/126}
\mspage{313}
to \unsure{ismmorphism|isomorphism})  as the homotopy type of a \hand{\uline{\textcolor{black}{space admitting a finite}}}
\hand{\uline{\textcolor{black}{triangulation}}}.  In terms of \var{\over[typed]{S}{s}}{s}{typewriter overwriting}{ps_134-5_126-v0001}s sets i.e. \var{\inserted{of}}{\hand{of}}{insertion}{ps_134-5_126-v0002} the modelizer \hand{$\Delta$}$^\wedge$,  the natu-%
ral finiteness condition,  suggested by the algebraic formalism,  is
that the homotopy type be isomorphic to one defined by an object ``\hand{\uline{\textcolor{black}{of}}}
\hand{\uline{\textcolor{black}{finite presentation}}}'' \hand{\uline{\textcolor{black}{in}}} \hand{$\Delta$}$^\wedge$ ,  namely one which is a direct limit of
a \hand{\uline{\textcolor{black}{finite}}} diagram in \hand{$\Delta$}$^\wedge$,  made up with simplices  i.e. coming from a
diagram in \hand{$\Delta$}.  It is clear that the first finiteness condition implies
the second,  by using a total order on the set of vertices of the trian-%
gulation.  \var{T\(\underset{\hand{h}}{\text{h}}\)e}{The}{placement}{ps_134-5_126-v0003} converse  should'nt be hard either,  using an induction
argument on the number of simplices occuring in the diagram,  and using
the fact that any \var{\xout{subobject}}{}{typewriter deletion}{ps_134-5_126-v0004} quotient object in \hand{$\Delta$}$^\wedge$ of a simplex is
again a simplex,  hence also a subobject of a simplex is a union of
subsimplices;  from this should follow by induction that the geometric
realization of a  ss set of finite presentation is endowed with a
natural \var{\inserted{compact}}{\hand{compact}}{insertion}{ps_134-5_126-v0005} \hand{\uline{\textcolor{black}{piecewise}}} \hand{\uline{\textcolor{black}{linear structure}}},  and hence can be \var{\inserted{finitely}}{\hand{finitely}}{insertion}{ps_134-5_126-v0006} triangulated. Pre-%
sumably,  one can even find a canonical triangulation,  using twofold
barycentric subdivision Sd$^2$ (again !)  on any simplex.  All this is su-%
rely standard knowledge,  and I do'nt feel like diving into technicali-%
ties on this matter,  unless I am forced to.

If we start with an arbitrary test category A,  the notion of an
object of finite presentation in A$^\wedge$ still makes sense.  Indeed,  in any
category M,  stable \var{\deleted{under}}{}{deletion}{ps_134-5_126-v0007} \var{\inserted[typed]{filtering}}{filtering}{typed insertion}{ps_134-5_126-v0008} small direct limits,  we may define object\hand{s} of
finite presentation as those for which the corresponding \hand{\uline{\textcolor{black}{covariant}}}
functor
\[ Y \hand{\longmapsto} \mathrm{Hom}(X,Y) \]
commutes to filt\var{\xout{r}}{}{typewriter deletion}{ps_134-5_126-v0009}ering  direct limits.  If M is stable under \var{\xout{any typ}}{}{typewriter deletion}{ps_134-5_126-v0010}
some type of finite direct limits,  say \var{\inserted{under}}{\hand{under}}{insertion}{ps_134-5_126-v0011} any finite direct limits,  then
so is the \var{\inserted[typed]{full}}{full}{typed insertion}{ps_134-5_126-v0012} subcategory $M_{fp}$  of objects of finite presentation of M .
In the case when M=A$^\wedge$,  A any small category,  it is obvious that ob-%
jects of A,  and hence finite direct limits of such,  are of finite pre-%
sentation,  and it is not difficult to show that the converse equally
holds.  As a matter of fact,  A$^\wedge_{fp}$  \hand{\uline{\textcolor{black}{can be viewed as the solution of the}}}
2-\hand{\uline{\textcolor{black}{universal}}} \hand{\uline{\textcolor{black}{problem of}}} ``\hand{\uline{\textcolor{black}{adding finite direct limits to A}}}'' .

Coming back to the case when A is a test category\hand{,} and hence A$^\wedge$ is
modelizing,  one may ask for conditions \var{\inserted{upon A}}{\hand{upon A}}{insertion}{ps_134-5_126-v0013} which ensure that the homotopy
types of finite presentation are exactly those isomorphic to the homoto-%
py types defined by objects of A$^\wedge_{fp}$ .  \var{\over{\unsure{It}}{One}}{\hand{One}}{writing over}{ps_134-5_126-v0014} \var{\deleted{is clear}}{}{deletion}{ps_134-5_126-v0015} \var{\inserted{expects}}{\hand{expects}}{insertion}{ps_134-5_126-v0016} that some stringent

\margin{126}

%% END SOURCE ps_134-5_126
\ifnum\EditionMode=1
\psapparatusstart
\psentry{ps_134-5_126-v0001}{1}{typewriter overwriting}{\over[typed]{S}{s}}{s}{The superposed typed reading replaces the underlying reading.}
\psentry{ps_134-5_126-v0002}{2}{insertion}{\inserted{of}}{\hand{of}}{The handwritten insertion is incorporated at its indicated position; insertion brackets leave the running text.}
\psentry{ps_134-5_126-v0003}{3}{placement}{T\(\underset{\hand{h}}{\text{h}}\)e}{The}{The typed h has an additional pen h below it; the diplomatic form preserves both positions while the reading gives the word once.}
\psentry{ps_134-5_126-v0004}{4}{typewriter deletion}{\xout{subobject}}{}{The typewriter strikes out this material; it leaves the corrected running text.}
\psentry{ps_134-5_126-v0005}{5}{insertion}{\inserted{compact}}{\hand{compact}}{The handwritten insertion is incorporated at its indicated position; insertion brackets leave the running text.}
\psentry{ps_134-5_126-v0006}{6}{insertion}{\inserted{finitely}}{\hand{finitely}}{The handwritten insertion is incorporated at its indicated position; insertion brackets leave the running text.}
\psentry{ps_134-5_126-v0007}{7}{deletion}{\deleted{under}}{}{The author cancels this material by hand; it leaves the corrected running text.}
\psentry{ps_134-5_126-v0008}{8}{typed insertion}{\inserted[typed]{filtering}}{filtering}{The typed insertion is incorporated at its indicated position; insertion brackets leave the running text.}
\psentry{ps_134-5_126-v0009}{9}{typewriter deletion}{\xout{r}}{}{The typewriter strikes out this material; it leaves the corrected running text.}
\psentry{ps_134-5_126-v0010}{10}{typewriter deletion}{\xout{any typ}}{}{The typewriter strikes out this material; it leaves the corrected running text.}
\psentry{ps_134-5_126-v0011}{11}{insertion}{\inserted{under}}{\hand{under}}{The handwritten insertion is incorporated at its indicated position; insertion brackets leave the running text.}
\psentry{ps_134-5_126-v0012}{12}{typed insertion}{\inserted[typed]{full}}{full}{The typed insertion is incorporated at its indicated position; insertion brackets leave the running text.}
\psentry{ps_134-5_126-v0013}{13}{insertion}{\inserted{upon A}}{\hand{upon A}}{The handwritten insertion is incorporated at its indicated position; insertion brackets leave the running text.}
\psentry{ps_134-5_126-v0014}{14}{writing over}{\over{\unsure{It}}{One}}{\hand{One}}{The superposed handwritten reading replaces the underlying reading.}
\psentry{ps_134-5_126-v0015}{15}{deletion}{\deleted{is clear}}{}{The author cancels this material by hand; it leaves the corrected running text.}
\psentry{ps_134-5_126-v0016}{16}{insertion}{\inserted{expects}}{\hand{expects}}{The handwritten insertion is incorporated at its indicated position; insertion brackets leave the running text.}
\psapparatusend
\fi
\pssource{ps_134-5_127}{134-5/127}{https://grothendieck.umontpellier.fr/web/viewer.html?file=../134-5.pdf\#page=128}
\psnumber{ps_134-5_127-v0001}{1}
\psnumber{ps_134-5_127-v0002}{2}
\psnumber{ps_134-5_127-v0003}{3}
\psnumber{ps_134-5_127-v0004}{4}
\psnumber{ps_134-5_127-v0005}{5}
\psnumber{ps_134-5_127-v0006}{6}
\psnumber{ps_134-5_127-v0007}{7}
\psnumber{ps_134-5_127-v0008}{8}
\psnumber{ps_134-5_127-v0009}{9}
\psnumber{ps_134-5_127-v0010}{10}
\psnumber{ps_134-5_127-v0011}{11}
\psnumber{ps_134-5_127-v0012}{12}
\psnumber{ps_134-5_127-v0013}{13}
\psnumber{ps_134-5_127-v0014}{14}
\psnumber{ps_134-5_127-v0015}{15}
\psnumber{ps_134-5_127-v0016}{16}
\psnumber{ps_134-5_127-v0017}{17}
\psnumber{ps_134-5_127-v0018}{18}
\psnumber{ps_134-5_127-v0019}{19}
\psnumber{ps_134-5_127-v0020}{20}
\psnumber{ps_134-5_127-v0021}{21}
\psnumber{ps_134-5_127-v0022}{22}
%% BEGIN SOURCE ps_134-5_127
% Pursuing Stacks (cotes 134-2 to 134-8): archive cote 134-5, site page 127 (PDF page 128), typescript page 314.
% typescript. Draft compared: ps5.tex:1094-1144.
% AI-transcribed by Claude Opus 5.5 (claude-opus-5-5, adaptive thinking, effort high), Message Batches API, batch ps3 (msgbatch_01XdFmuVVBAcgc7EFxdtCqsj), 2026-11-12; reviewed against both supplied scan halves; see CHECK_DECISIONS.tsv.
\scanpage{134-5/127}
\mspage{314}

extra condition is needed on A to ensure this, \var{\deleted{as we may take for A
any full subcategory of (Cat)$_{\mathrm{as}}$ stable under products and containing
but not $\mathord{\textcolor{black}{0}\llap{\hand{\ensuremath{/}}}}$ , which will ensure that A is \inserted[typed]{strict} a test category, \xout{Sorry \unsure{,} I}
\xout{was going to add that Then and the inclusion functor}}}{}{deletion}{ps_134-5_127-v0001} To see this, let's
take a finite group G, and an aspherical topological space $E_G$ upon
which G operates freely, with quotient $B_G$ , a classifying space of G.
I\var{\over{s}{f}}{\hand{f}}{writing over}{ps_134-5_127-v0002} G $\hand{\neq}$ 1\hand{,} the homotopy type of $B_G$ is'nt of finite type, because $B_G$ has
non vanishing cohomology groups in arbitrarily high dimensions, as
well known. We could transpose the \var{\inserted{following}}{\hand{following}}{insertion}{ps_134-5_127-v0003} construction in either (Cat) or
\hand{$\Delta^{\textcolor{black}{\wedge}}$} say, \var{\xout{h}}{}{typewriter deletion}{ps_134-5_127-v0004}but we may as \var{\xout{q}}{}{typewriter deletion}{ps_134-5_127-v0005}well work in the modelizer \var{\inserted[typed]{$M=$}}{$M=$}{typed insertion}{ps_134-5_127-v0006}(Spaces), and take
any \var{\inserted[typed]{small}}{small}{typed insertion}{ps_134-5_127-v0007} full subcategory A of M , containing the unit interval \var{\(\underset{\hand{-\quad-\quad-}}{\text{\deleted{but not}}}\)}{}{placement}{ps_134-5_127-v0008}\var{\inserted{but not}}{\hand{but not}}{insertion}{ps_134-5_127-v0009}
\var{\(\underset{\hand{-\quad-\quad-\quad-\quad-}}{\text{\deleted{the empty space}}}\)}{}{placement}{ps_134-5_127-v0010}\var{\inserted{the empty space}}{\hand{the empty space}}{insertion}{ps_134-5_127-v0011}, and stable under finite products - which implies
that A is a strict test category. We'll take A large enough to contain
$E_G$ \var{\xout{as an object}}{}{typewriter deletion}{ps_134-5_127-v0012} , and small enough to be made up with aspheric spaces,
hence the inclusion functor
\[ i:\var{\deleted{\unsure{f}}}{}{deletion}{ps_134-5_127-v0013}A \hand{\longrightarrow} \text{(Spaces)} \]
is a test functor. Now consider the quotient object \hand{\uline{\textcolor{black}{in}}} A$^{\wedge}$
\[ F = E_G \,\hand{/}\, G \quad \text{in } A^{\wedge} \quad , \]
i.e. the \var{\xout{functor\illegible{5}A\illegible{2}F}}{}{typewriter deletion}{ps_134-5_127-v0014}\var{\inserted[typed]{presheaf}}{presheaf}{typed insertion}{ps_134-5_127-v0015} on A
\[ F: \quad T \hand{\longmapsto} \mathrm{Hom}(T,E_G)\hand{/}G \ \hand{\simeq}\ \mathrm{Hom}(T,B_G) \qquad , \]
\var{\xout{and}}{}{typewriter deletion}{ps_134-5_127-v0016}\var{\inserted[typed]{where}}{where}{typed insertion}{ps_134-5_127-v0017} the last isomorphism comes from the fact that the object T of A
is \var{\over[typed]{\illegible{1}}{a}}{a}{typewriter overwriting}{ps_134-5_127-v0018}spheric and hence 1-connected. Thus we get
\[ F \hand{\unsure{\simeq|\approx}} i^{*}(B_G) \quad , \]
hence the homotopy type \var{\xout{of}}{}{typewriter deletion}{ps_134-5_127-v0019} defined by F is the homotopy type of $B_G$ ,
which is not of finite type, despite the fact that F is of finite pre-%
sentation.

In order to insure that the homotopy type defined by any object
in $A^{\wedge}{}_{\mathrm{fp}}$ be of finite type, it may be useful perhaps to make \var{\over[typed]{\unsure{i}}{o}}{o}{typewriter overwriting}{ps_134-5_127-v0020}n A the
assumption that any quotient in A$^{\wedge}$ of an obje\var{\over[typed]{c}{c}}{c}{typewriter overwriting}{ps_134-5_127-v0021}t in A is isomorphic to an
object in A, and that the set of all subobjects in A$^{\wedge}$ of an object \var{\inserted[typed]{a}}{a}{typed insertion}{ps_134-5_127-v0022} in
A (i.e. the set of all sieves on a ) is finite - possibly too that
for any two objects a and b of A, Hom(a,b) is finite. As for the
opposite inclusion, namely that any homotopy type of finite type can
be described by an object of $A^{\wedge}{}_{\mathrm{fp}}$ , this would follow from the existen-%
ce of a \hand{\uline{\textcolor{black}{perfect}}} test functor from \hand{$\Delta$} into A$^{\wedge}$ , factoring through

\margin{127}

%% END SOURCE ps_134-5_127
\ifnum\EditionMode=1
\psapparatusstart
\psentry{ps_134-5_127-v0001}{1}{deletion}{\deleted{as we may take for A
any full subcategory of (Cat)$_{\mathrm{as}}$ stable under products and containing
but not $\mathord{\textcolor{black}{0}\llap{\hand{\ensuremath{/}}}}$ , which will ensure that A is \inserted[typed]{strict} a test category, \xout{Sorry \unsure{,} I}
\xout{was going to add that Then and the inclusion functor}}}{}{The author cancels this material by hand; it leaves the corrected running text.}
\psentry{ps_134-5_127-v0002}{2}{writing over}{\over{s}{f}}{\hand{f}}{The superposed handwritten reading replaces the underlying reading.}
\psentry{ps_134-5_127-v0003}{3}{insertion}{\inserted{following}}{\hand{following}}{The handwritten insertion is incorporated at its indicated position; insertion brackets leave the running text.}
\psentry{ps_134-5_127-v0004}{4}{typewriter deletion}{\xout{h}}{}{The typewriter strikes out this material; it leaves the corrected running text.}
\psentry{ps_134-5_127-v0005}{5}{typewriter deletion}{\xout{q}}{}{The typewriter strikes out this material; it leaves the corrected running text.}
\psentry{ps_134-5_127-v0006}{6}{typed insertion}{\inserted[typed]{$M=$}}{$M=$}{The typed insertion is incorporated at its indicated position; insertion brackets leave the running text.}
\psentry{ps_134-5_127-v0007}{7}{typed insertion}{\inserted[typed]{small}}{small}{The typed insertion is incorporated at its indicated position; insertion brackets leave the running text.}
\psentry{ps_134-5_127-v0008}{8}{placement}{\(\underset{\hand{-\quad-\quad-}}{\text{\deleted{but not}}}\)}{}{Three short pen marks lie below the crossed typed phrase, which is separately recopied above. Their chronology and function remain unassigned; the reading follows the live copied phrase without duplicating the cancelled one.}
\psentry{ps_134-5_127-v0009}{9}{insertion}{\inserted{but not}}{\hand{but not}}{The handwritten insertion is incorporated at its indicated position; insertion brackets leave the running text.}
\psentry{ps_134-5_127-v0010}{10}{placement}{\(\underset{\hand{-\quad-\quad-\quad-\quad-}}{\text{\deleted{the empty space}}}\)}{}{Five short pen marks lie below the crossed typed phrase, separately recopied above. Preserve the graphic diplomatically without inventing a reinstatement sequence.}
\psentry{ps_134-5_127-v0011}{11}{insertion}{\inserted{the empty space}}{\hand{the empty space}}{The handwritten insertion is incorporated at its indicated position; insertion brackets leave the running text.}
\psentry{ps_134-5_127-v0012}{12}{typewriter deletion}{\xout{as an object}}{}{The typewriter strikes out this material; it leaves the corrected running text.}
\psentry{ps_134-5_127-v0013}{13}{deletion}{$\deleted{\unsure{f}}$}{}{The author cancels this material by hand; it leaves the corrected running text.}
\psentry{ps_134-5_127-v0014}{14}{typewriter deletion}{\xout{functor\illegible{5}A\illegible{2}F}}{}{The typewriter strikes out this material; it leaves the corrected running text.}
\psentry{ps_134-5_127-v0015}{15}{typed insertion}{\inserted[typed]{presheaf}}{presheaf}{The typed insertion is incorporated at its indicated position; insertion brackets leave the running text.}
\psentry{ps_134-5_127-v0016}{16}{typewriter deletion}{\xout{and}}{}{The typewriter strikes out this material; it leaves the corrected running text.}
\psentry{ps_134-5_127-v0017}{17}{typed insertion}{\inserted[typed]{where}}{where}{The typed insertion is incorporated at its indicated position; insertion brackets leave the running text.}
\psentry{ps_134-5_127-v0018}{18}{typewriter overwriting}{\over[typed]{\illegible{1}}{a}}{a}{The superposed typed reading replaces the underlying reading.}
\psentry{ps_134-5_127-v0019}{19}{typewriter deletion}{\xout{of}}{}{The typewriter strikes out this material; it leaves the corrected running text.}
\psentry{ps_134-5_127-v0020}{20}{typewriter overwriting}{\over[typed]{\unsure{i}}{o}}{o}{The superposed typed reading replaces the underlying reading.}
\psentry{ps_134-5_127-v0021}{21}{typewriter overwriting}{\over[typed]{c}{c}}{c}{The superposed typed reading replaces the underlying reading.}
\psentry{ps_134-5_127-v0022}{22}{typed insertion}{\inserted[typed]{a}}{a}{The typed insertion is incorporated at its indicated position; insertion brackets leave the running text.}
\psapparatusend
\fi
\pssource{ps_134-5_128}{134-5/128}{https://grothendieck.umontpellier.fr/web/viewer.html?file=../134-5.pdf\#page=129}
\psnumber{ps_134-5_128-v0001}{1}
\psnumber{ps_134-5_128-v0002}{2}
\psnumber{ps_134-5_128-v0003}{3}
\psnumber{ps_134-5_128-v0004}{4}
\psnumber{ps_134-5_128-v0005}{5}
\psnumber{ps_134-5_128-v0006}{6}
\psnumber{ps_134-5_128-v0007}{7}
\psnumber{ps_134-5_128-v0008}{8}
\psnumber{ps_134-5_128-v0009}{9}
\psnumber{ps_134-5_128-v0010}{10}
\psnumber{ps_134-5_128-v0011}{11}
\psnumber{ps_134-5_128-v0012}{12}
\psnumber{ps_134-5_128-v0013}{13}
\psnumber{ps_134-5_128-v0014}{14}
\psnumber{ps_134-5_128-v0015}{15}
\psnumber{ps_134-5_128-v0016}{16}
\psnumber{ps_134-5_128-v0017}{17}
\psnumber{ps_134-5_128-v0018}{18}
\psnumber{ps_134-5_128-v0019}{19}
\psnumber{ps_134-5_128-v0020}{20}
\psnumber{ps_134-5_128-v0021}{21}
\psnumber{ps_134-5_128-v0022}{22}
\psnumber{ps_134-5_128-v0023}{23}
\psnumber{ps_134-5_128-v0024}{24}
\psnumber{ps_134-5_128-v0025}{25}
\psnumber{ps_134-5_128-v0026}{26}
\psnumber{ps_134-5_128-v0027}{27}
\psnumber{ps_134-5_128-v0028}{28}
\psnumber{ps_134-5_128-v0029}{29}
%% BEGIN SOURCE ps_134-5_128
% Pursuing Stacks (cotes 134-2 to 134-8): archive cote 134-5, site page 128 (PDF page 129), typescript page 315.
% typescript. Draft compared: ps5.tex:1130-1184.
% AI-transcribed by Claude Opus 5.5 (claude-opus-5-5, adaptive thinking, effort high), Message Batches API, batch ps3 (msgbatch_01XdFmuVVBAcgc7EFxdtCqsj), 2026-11-12; reviewed against both supplied scan halves; see CHECK_DECISIONS.tsv.
% Both hemispherical signs have an inner curved stroke; the bigcirc identification remains tentative.
\scanpage{134-5/128}
\mspage{315}

$A^{\wedge}_{\mathrm{fp}}$ . Thus the \var{\over[typed]{\illegible{1}}{p}}{p}{typewriter overwriting}{ps_134-5_128-v0001}\var{\over[typed]{\illegible{1}}{r}}{r}{typewriter overwriting}{ps_134-5_128-v0002}esent question about finiteness conditions, seems
to be related (possibly) to the previous one about existence of various
types of test functors.

\var{T\(\underset{\hand{h}}{\text{h}}\)e}{The}{placement}{ps_134-5_128-v0003} condition for a strict test category A we are lookihg at is
surely satisfied, besides \hand{$\Delta$}, by the cubical and hemispherical test
categories \hand{$\Box$} and \hand{$\unsure{\bigcirc}$}, and surely also by any finite products of these.
I add this comment, of course, in order to ``push through'' the point
that \var{\inserted{not any more}}{\hand{not any more}}{insertion}{ps_134-5_128-v0004} with re\var{\over[typed]{\illegible{1}}{s}}{s}{typewriter overwriting}{ps_134-5_128-v0005}\var{\over[typed]{\illegible{1}}{p}}{p}{typewriter overwriting}{ps_134-5_128-v0006}ect to finiteness conditions on homotopy types, \var{\over{as}{than}}{\hand{than}}{writing over}{ps_134-5_128-v0007}
(presumably at least, for the time being) in \var{\deleted{no}}{}{deletion}{ps_134-5_128-v0008}\var{\inserted{any}}{\hand{any}}{insertion}{ps_134-5_128-v0009} other essential respect
concerning the ability \var{\xout{of}}{}{typewriter deletion}{ps_134-5_128-v0010} for expressing basic situations and facts in
homotopy or cohomology theory, the category of simplices s\var{\over[typed]{a}{t}}{t}{typewriter overwriting}{ps_134-5_128-v0011}\var{\over[typed]{t}{a}}{a}{typewriter overwriting}{ps_134-5_128-v0012}nds singled
out by itself from all other test categories. \var{N\(\underset{\hand{o}}{\text{o}}\)r}{Nor}{placement}{ps_134-5_128-v0013} does it seem that
the ``trinity''
\[ \hand{\Delta}\ ,\ \hand{\Box}\ ,\ \hand{\unsure{\bigcirc}} \]
has this property, with the only exception so far, possibly, of the
Dold-Puppe theorem\var{\over{.}{,}}{\hand{,}}{writing over}{ps_134-5_128-v0014}\hand{(} as no other test category except these is known to
me for which a Dold-Puppe theorem in it's strict form holds true
(compare reflections section 71)\hand{)}.

8) I could make the same point in favor of more general test cate-%
gories than the trinity above, when it comes to the existence of an
algorithm\var{\xout{\unsure{s}}}{}{typewriter deletion}{ps_134-5_128-v0015} for computing homology and cohomology groups, using suitab-%
le \hand{\uline{\textcolor{black}{boundary}}} \hand{\uline{\textcolor{black}{operators}}}. What is meant by these is clear of course for
the three types of complexes, but then it extends in an obvious way to
multicomplexes too - which means that \var{\xout{\illegible{2}}}{}{typewriter deletion}{ps_134-5_128-v0016}\var{\inserted[typed]{for}}{for}{typed insertion}{ps_134-5_128-v0017} the test categories deduced
from the trinity by taking finite products, we still get an algorithm\var{\xout{\unsure{s}}}{}{typewriter deletion}{ps_134-5_128-v0018}
for cohomology via boundary operators. Of course, for any test catego-%
ry A , using a test functor (if we can find one) from one of the
three above (say) into $A^{\wedge}$ will allow us to reduce ``computation'' of
homology and cohomology invariants in terms of a model in $A^{\wedge}$, to the
case of the corresponding type of complexes - hence \unsure{again|agian} an algo-%
rithm\var{\deleted{\unsure{s}}}{}{deletion}{ps_134-5_128-v0019} (similar to the \var{\over[typed]{o}{f}}{f}{typewriter overwriting}{ps_134-5_128-v0020}\var{\over[typed]{\illegible{1}}{a}}{a}{typewriter overwriting}{ps_134-5_128-v0021}miliar one of computing the cohomomology of an
object of (Cat) \var{\xout{\illegible{1}i}}{}{typewriter deletion}{ps_134-5_128-v0022} semisi\var{\over{\illegible{1}}{m}}{\hand{m}}{writing over}{ps_134-5_128-v0023}plicially, via the nerve\hand{)}. But this is
cheating of course ! The question I want to raise here is about
\var{\inserted[typed]{existence of}}{existence of}{typed insertion}{ps_134-5_128-v0024} \hand{``}boundary operations\hand{''} \var{\xout{from A to}}{}{typewriter deletion}{ps_134-5_128-v0025}\var{\deleted{\inserted[typed]{in}}}{}{deletion}{ps_134-5_128-v0026}\var{\inserted{in}}{\hand{in}}{insertion}{ps_134-5_128-v0027} A, similar to the familiar ones used for
the three \var{\inserted[typed]{basic}}{basic}{typed insertion}{ps_134-5_128-v0028} types of complexes, and allowing to compute the homology and
cohomology groups of an object of $A^{\wedge}$ in the usual way, involving sui-%
table \hand{\uline{\textcolor{black}{signs}}} + or - associated to the various boundary operatio\var{\over[typed]{j}{n}}{n}{typewriter overwriting}{ps_134-5_128-v0029}s.

\margin{128}

%% END SOURCE ps_134-5_128
\ifnum\EditionMode=1
\psapparatusstart
\psentry{ps_134-5_128-v0001}{1}{typewriter overwriting}{\over[typed]{\illegible{1}}{p}}{p}{The superposed typed reading replaces the underlying reading.}
\psentry{ps_134-5_128-v0002}{2}{typewriter overwriting}{\over[typed]{\illegible{1}}{r}}{r}{The superposed typed reading replaces the underlying reading.}
\psentry{ps_134-5_128-v0003}{3}{placement}{T\(\underset{\hand{h}}{\text{h}}\)e}{The}{A pen h lies beneath the typed h; preserve both positions diplomatically and read the word once.}
\psentry{ps_134-5_128-v0004}{4}{insertion}{\inserted{not any more}}{\hand{not any more}}{The handwritten insertion is incorporated at its indicated position; insertion brackets leave the running text.}
\psentry{ps_134-5_128-v0005}{5}{typewriter overwriting}{\over[typed]{\illegible{1}}{s}}{s}{The superposed typed reading replaces the underlying reading.}
\psentry{ps_134-5_128-v0006}{6}{typewriter overwriting}{\over[typed]{\illegible{1}}{p}}{p}{The superposed typed reading replaces the underlying reading.}
\psentry{ps_134-5_128-v0007}{7}{writing over}{\over{as}{than}}{\hand{than}}{The superposed handwritten reading replaces the underlying reading.}
\psentry{ps_134-5_128-v0008}{8}{deletion}{\deleted{no}}{}{The author cancels this material by hand; it leaves the corrected running text.}
\psentry{ps_134-5_128-v0009}{9}{insertion}{\inserted{any}}{\hand{any}}{The handwritten insertion is incorporated at its indicated position; insertion brackets leave the running text.}
\psentry{ps_134-5_128-v0010}{10}{typewriter deletion}{\xout{of}}{}{The typewriter strikes out this material; it leaves the corrected running text.}
\psentry{ps_134-5_128-v0011}{11}{typewriter overwriting}{\over[typed]{a}{t}}{t}{The superposed typed reading replaces the underlying reading.}
\psentry{ps_134-5_128-v0012}{12}{typewriter overwriting}{\over[typed]{t}{a}}{a}{The superposed typed reading replaces the underlying reading.}
\psentry{ps_134-5_128-v0013}{13}{placement}{N\(\underset{\hand{o}}{\text{o}}\)r}{Nor}{A pen o lies beneath the typed o; preserve both positions diplomatically and read the word once.}
\psentry{ps_134-5_128-v0014}{14}{writing over}{\over{.}{,}}{\hand{,}}{The superposed handwritten reading replaces the underlying reading.}
\psentry{ps_134-5_128-v0015}{15}{typewriter deletion}{\xout{\unsure{s}}}{}{The typewriter strikes out this material; it leaves the corrected running text.}
\psentry{ps_134-5_128-v0016}{16}{typewriter deletion}{\xout{\illegible{2}}}{}{The typewriter strikes out this material; it leaves the corrected running text.}
\psentry{ps_134-5_128-v0017}{17}{typed insertion}{\inserted[typed]{for}}{for}{The typed insertion is incorporated at its indicated position; insertion brackets leave the running text.}
\psentry{ps_134-5_128-v0018}{18}{typewriter deletion}{\xout{\unsure{s}}}{}{The typewriter strikes out this material; it leaves the corrected running text.}
\psentry{ps_134-5_128-v0019}{19}{deletion}{\deleted{\unsure{s}}}{}{The author cancels this material by hand; it leaves the corrected running text.}
\psentry{ps_134-5_128-v0020}{20}{typewriter overwriting}{\over[typed]{o}{f}}{f}{The superposed typed reading replaces the underlying reading.}
\psentry{ps_134-5_128-v0021}{21}{typewriter overwriting}{\over[typed]{\illegible{1}}{a}}{a}{The superposed typed reading replaces the underlying reading.}
\psentry{ps_134-5_128-v0022}{22}{typewriter deletion}{\xout{\illegible{1}i}}{}{The typewriter strikes out this material; it leaves the corrected running text.}
\psentry{ps_134-5_128-v0023}{23}{writing over}{\over{\illegible{1}}{m}}{\hand{m}}{The superposed handwritten reading replaces the underlying reading.}
\psentry{ps_134-5_128-v0024}{24}{typed insertion}{\inserted[typed]{existence of}}{existence of}{The typed insertion is incorporated at its indicated position; insertion brackets leave the running text.}
\psentry{ps_134-5_128-v0025}{25}{typewriter deletion}{\xout{from A to}}{}{The typewriter strikes out this material; it leaves the corrected running text.}
\psentry{ps_134-5_128-v0026}{26}{deletion}{\deleted{\inserted[typed]{in}}}{}{The author cancels this material by hand; it leaves the corrected running text.}
\psentry{ps_134-5_128-v0027}{27}{insertion}{\inserted{in}}{\hand{in}}{The handwritten insertion is incorporated at its indicated position; insertion brackets leave the running text.}
\psentry{ps_134-5_128-v0028}{28}{typed insertion}{\inserted[typed]{basic}}{basic}{The typed insertion is incorporated at its indicated position; insertion brackets leave the running text.}
\psentry{ps_134-5_128-v0029}{29}{typewriter overwriting}{\over[typed]{j}{n}}{n}{The superposed typed reading replaces the underlying reading.}
\psapparatusend
\fi
\pssource{ps_134-5_129}{134-5/129}{https://grothendieck.umontpellier.fr/web/viewer.html?file=../134-5.pdf\#page=130}
\psnumber{ps_134-5_129-v0001}{1}
\psnumber{ps_134-5_129-v0002}{2}
\psnumber{ps_134-5_129-v0003}{3}
\psnumber{ps_134-5_129-v0004}{4}
\psnumber{ps_134-5_129-v0005}{5}
\psnumber{ps_134-5_129-v0006}{6}
\psnumber{ps_134-5_129-v0007}{7}
\psnumber{ps_134-5_129-v0008}{8}
\psnumber{ps_134-5_129-v0009}{9}
\psnumber{ps_134-5_129-v0010}{10}
\psnumber{ps_134-5_129-v0011}{11}
%% BEGIN SOURCE ps_134-5_129
% Pursuing Stacks (cotes 134-2 to 134-8): archive cote 134-5, site page 129 (PDF page 130), typescript page 316.
% typescript. Draft compared: ps5.tex:1171-1256.
% AI-transcribed by Claude Opus 5.5 (claude-opus-5-5, adaptive thinking, effort high), Message Batches API, batch ps3 (msgbatch_01XdFmuVVBAcgc7EFxdtCqsj), 2026-11-12; reviewed against both supplied scan halves; see CHECK_DECISIONS.tsv.
\scanpage{134-5/129}
\mspage{316}

It should'nt be hard,  I feel,  to pin down exactly what is needed for
getting such a formalism.  The intuitive idea behind it (suggested by
the example of \var{\xout{the}}{}{typewriter deletion}{ps_134-5_129-v0001} standard complexes and multicomplexes)  is that
\hand{\uline{\textcolor{black}{such a formalism should be associated to cellular decompositions of}}}
$n$-\hand{\uline{\textcolor{black}{cells}}} \hand{\uline{\textcolor{black}{for variable}}} $n$,  such that the interior of each $n$-cell should
be an open cell of the subdivision.  \var{T\(\underset{\hand{h}}{\text{h}}\)ere}{There}{placement}{ps_134-5_129-v0002} may of course be several
$n$-cells\var{\xout{\unsure{s}}}{}{typewriter deletion}{ps_134-5_129-v0003} for the same $n$,  which are not combinatorially isomorphic.
When trying to express this idea in a precise way,  we are led to assu-%
me,  as an extra structure on the would-be test cate\var{\over[typed]{t}{g}}{g}{typewriter overwriting}{ps_134-5_129-v0004}ory $A$,  a functor
\[ (1) \qquad i:\ A \hand{\longrightarrow} (\mathrm{Ord}) \]
of $A$ into the category of ordered sets,  such that for any $a$ in $A$,
$i(a)$ be a \hand{\uline{\textcolor{black}{finite}}} ordered set,  whose geometric \unsure{realization|relaization}  (cf. section
22)  is an $n$-cell for suitable  $n \overset{\hand{def}}{=} \mathrm{dim}(a)$.  We assume moreover that
$i(a)$ has a largest element  $e(a)$,  and that the geometric realization
of  $i(a) \hand{\setminus\{} e(a) \hand{\}}$ is the bounding  $(n-1)$-sphere of the $n$-cell  $\hand{|} i(a) \hand{|}$ :
\[ (2) \qquad \hand{|} i(a)^{*} \hand{|}\ \var{\deleted{home\xout{\unsure{m}}omorphic to}}{}{deletion}{ps_134-5_129-v0005}\ \hand{\simeq}\ S^{n-1} \hand{,}\ \hand{\text{where}}\ i(a)^{*} \overset{\hand{def}}{=} i(a) \hand{\setminus\{}e(a)\hand{\}} \hand{,} \]
which will imply the precedent condition namely
\[ (3) \qquad \hand{|} i(a) \hand{|} \hand{\simeq} B^{n} \quad , \]
as $B^{n}$ can be identified\var{\deleted{\unsure{e}}}{}{deletion}{ps_134-5_129-v0006} to the cone over $S^{n-1}$.  As another condition,
we need that \var{\xout{for any}}{}{typewriter deletion}{ps_134-5_129-v0007}

(4) \quad For any $a$ in $A$ and  $x \hand{\in} i(a)$,  there exists $b$ in $A$ and an iso-%
morphism
\[ i(b) \hand{\xrightarrow{\sim}} i(a)_{/x} \overset{\mathrm{def}}{=} \hand{\{}\, y \hand{\in} i(a) \hand{|}\ y \hand{\leq} x \,\hand{\}} \ \hand{\hookrightarrow i(a)}, \]
induced by a map
\[ b \hand{\longrightarrow} a \]
in $A$.

It is enough to make this assumption for $x$ of codimension 1 in $i(a)$,
which will imply thatnit is true for any $x$.  It seems reasonable on the
other hand to make the assumption that for a given $x$ ,  the object  $b$
in (4),  viewed as an object of  $A_{/a}$ \hand{,}  is determined up to a unique
isomorphism ,  we may call it $\var{\over[typed]{b}{a}}{a}{typewriter overwriting}{ps_134-5_129-v0008}_{x}$ ,  and \var{\inserted{$\partial_x$}}{\hand{$\partial_x$}}{insertion}{ps_134-5_129-v0009} the canonical map of $b$ into $a$
\margin{\textcircled{$a_x$}}
\[ (5) \qquad \hand{\partial}_{x} :\ \var{\over[typed]{b}{a}}{a}{typewriter overwriting}{ps_134-5_129-v0010}_{x} \hand{\longrightarrow} a \qquad ( x \hand{\in} i(a)\hand{,} \text{ of codim.1 in } i(a)\var{\xout{\illegible{1}}}{}{typewriter deletion}{ps_134-5_129-v0011}\hand{,} \]
\[ \text{i.e. } \mathrm{dim}(x)=\mathrm{dim}(a)-1 \ ) \quad . \]
As an extra structure,  we need for any $a$ in $A$
\[ (6) \qquad \hand{\omega}_{a} \ ,\ \text{an \hand{\uline{\textcolor{black}{orientation}}} of the $n$-cell } i(a) \ \ (n=\mathrm{dim}(a)). \]
This allows us,  for any $x$ as in (5),  to define a signature
\[ (7) \qquad \hand{\varepsilon}(x) \ \text{ or } \ \hand{\varepsilon}_{a}(x) \hand{\in} \hand{\{} +1,-1 \hand{\}} \ , \]

\margin{129}

%% END SOURCE ps_134-5_129
\ifnum\EditionMode=1
\psapparatusstart
\psentry{ps_134-5_129-v0001}{1}{typewriter deletion}{\xout{the}}{}{The typewriter strikes out this material; it leaves the corrected running text.}
\psentry{ps_134-5_129-v0002}{2}{placement}{T\(\underset{\hand{h}}{\text{h}}\)ere}{There}{A pen h lies beneath the typed h; the reading gives There once.}
\psentry{ps_134-5_129-v0003}{3}{typewriter deletion}{\xout{\unsure{s}}}{}{The typewriter strikes out this material; it leaves the corrected running text.}
\psentry{ps_134-5_129-v0004}{4}{typewriter overwriting}{\over[typed]{t}{g}}{g}{The superposed typed reading replaces the underlying reading.}
\psentry{ps_134-5_129-v0005}{5}{deletion}{$\deleted{home\xout{\unsure{m}}omorphic to}$}{}{The author cancels this material by hand; it leaves the corrected running text.}
\psentry{ps_134-5_129-v0006}{6}{deletion}{\deleted{\unsure{e}}}{}{The author cancels this material by hand; it leaves the corrected running text.}
\psentry{ps_134-5_129-v0007}{7}{typewriter deletion}{\xout{for any}}{}{The typewriter strikes out this material; it leaves the corrected running text.}
\psentry{ps_134-5_129-v0008}{8}{typewriter overwriting}{$\over[typed]{b}{a}$}{$a$}{The superposed typed reading replaces the underlying reading.}
\psentry{ps_134-5_129-v0009}{9}{insertion}{\inserted{$\partial_x$}}{\hand{$\partial_x$}}{The handwritten insertion is incorporated at its indicated position; insertion brackets leave the running text.}
\psentry{ps_134-5_129-v0010}{10}{typewriter overwriting}{$\over[typed]{b}{a}$}{$a$}{The superposed typed reading replaces the underlying reading.}
\psentry{ps_134-5_129-v0011}{11}{typewriter deletion}{$\xout{\illegible{1}}$}{}{The typewriter strikes out this material; it leaves the corrected running text.}
\psapparatusend
\fi
\pssource{ps_134-5_130}{134-5/130}{https://grothendieck.umontpellier.fr/web/viewer.html?file=../134-5.pdf\#page=131}
\psnumber{ps_134-5_130-v0001}{1}
\psnumber{ps_134-5_130-v0002}{2}
\psnumber{ps_134-5_130-v0003}{3}
\psnumber{ps_134-5_130-v0004}{4}
\psnumber{ps_134-5_130-v0005}{5}
\psnumber{ps_134-5_130-v0006}{6}
\psnumber{ps_134-5_130-v0007}{7}
\psnumber{ps_134-5_130-v0008}{8}
\psnumber{ps_134-5_130-v0009}{9}
\psnumber{ps_134-5_130-v0010}{10}
\psnumber{ps_134-5_130-v0011}{11}
\psnumber{ps_134-5_130-v0012}{12}
\psnumber{ps_134-5_130-v0013}{13}
\psnumber{ps_134-5_130-v0014}{14}
\psnumber{ps_134-5_130-v0015}{15}
\psnumber{ps_134-5_130-v0016}{16}
\psnumber{ps_134-5_130-v0017}{17}
%% BEGIN SOURCE ps_134-5_130
% Pursuing Stacks (cotes 134-2 to 134-8): archive cote 134-5, site page 130 (PDF page 131), typescript page 317.
% typescript. Draft compared: ps5.tex:1234-1294.
% AI-transcribed by Claude Opus 5.5 (claude-opus-5-5, adaptive thinking, effort high), Message Batches API, batch ps3 (msgbatch_01XdFmuVVBAcgc7EFxdtCqsj), 2026-11-12; reviewed against both supplied scan halves; see CHECK_DECISIONS.tsv.
\scanpage{134-5/130}
\mspage{317}
which will be +1 or -1, depending on whether \var{\inserted[typed]{in}}{in}{typed insertion}{ps_134-5_130-v0001} the inclusion
\[ \hand{|\partial}_x\hand{|} \ : \ \hand{|}i(a_x)\hand{|} \hand{\longrightarrow} \hand{|}i(a)\hand{|}\hand{,} \]
the orientation $\hand{\omega}(a_x)$ is induced "à la Stokes" by the orientation $\hand{\omega}(a)$
of the ambient n-cell, or not. Having the boundary \var{\over[typed]{i}{o}}{o}{typewriter overwriting}{ps_134-5_130-v0002}perations (5)
with their signatures (7), \var{\xout{we}}{}{typewriter deletion}{ps_134-5_130-v0003} and the decomposition
(8) \[ A = \hand{\coprod}_{n \hand{\geq} 0} A_n \quad , \quad \text{where} \quad A_n = \hand{\{} a \hand{\in} \mathrm{Ob}\, A \hand{|} \dim a = n \hand{\}}, \]
we get in the usual way, for any contravariant functor \var{\inserted[typed]{K.}}{K.}{typed insertion}{ps_134-5_130-v0004} from A with
values in an additive category,\var{\xout{stable under sums}}{}{typewriter deletion}{ps_134-5_130-v0005} a corresponding
chain complex in this category, with components
(9) \[ K_n = \hand{\coprod}_{a \hand{\in} A_n} K.(a) \qquad , \]
and boundary operators defined in the usual way via (5) and (7). Apply-%
ing this to the case of the category (Ab) of abelian groups, or to it's
dual, \var{\xout{we}}{}{typewriter deletion}{ps_134-5_130-v0006} and to the abelianization of an object X of $A^{\wedge}$ , we obtain
a tentative way for computing the homology and cohomology groups of
the homotopy type of $\underline{A}_{/X}$ , and similarly for any system of twisted
coefficients on X. \var{T\(\underset{\hand{h}}{\text{h}}\)e}{The}{placement}{ps_134-5_130-v0007} question which arises here is to write down
\var{\xout{some}}{}{typewriter deletion}{ps_134-5_130-v0008} a set of natural extra conditions on the data\hand{,} \var{\xout{(if needed\illegible{1}}}{}{typewriter deletion}{ps_134-5_130-v0009} which \var{\xout{\illegible{1}}}{}{typewriter deletion}{ps_134-5_130-v0010}
will ensure that we do get a canonical isomorphism between the "homolo-%
gy" and "cohomology" groups thus constructed, and the usual homology
and cohomology invariants of the object $A_{/X}$ of (Cat). Moreover, \var{\over[typed]{we}{one}}{one}{typewriter overwriting}{ps_134-5_130-v0011}
would like too to have conditions to ensure that A is a test category,
or even a strict one.\var{\hand{\tikz[x=.5em,y=.6em,baseline=.2ex]{\draw (-1,0)--(0,0)--(0,1)--(1,1); \draw (.12,0)--(.12,1);}}}{}{placement}{ps_134-5_130-v0012}

One difficulty here, if one really wants a test
category and not just a weak one (which may not be without any problem
either), is that presumably for this we'll need suitable \hand{\uline{\textcolor{black}{degeneracy
operations}}}, which may \var{\xout{be a}}{}{typewriter deletion}{ps_134-5_130-v0013} well turn out a very exacting condition
indeed ! The sceptical reader may wonder, as I \var{\inserted{am}}{\hand{am}}{insertion}{ps_134-5_130-v0014} just \var{\deleted{did}}{}{deletion}{ps_134-5_130-v0015}\var{\inserted{doing}}{\hand{doing}}{insertion}{ps_134-5_130-v0016} myself,
whether there will be any example within the set-up I propose, which
does not reduce to a finite product of test categories in our trinity.
\var{\deleted{But there \deleted{are}\inserted{should be} obvious ways to satisfy the\xout{se conditions} preliminary con-%
ditions, it occurs to me, in such a way that A be a \xout{test} str\over[typed]{a}{i}ct test
category and the composition
\[ i: A \qquad (\mathrm{Ord}) \qquad (\mathrm{Cat}) \]
a test functor (maybe not a perfect one). To see this, let $A_o$ be
any \inserted{small full} subcategory of the category of finite ordered sets with a largest
element and satisfying the assumption (2) for i(a), such that $A_o$}}{}{deletion}{ps_134-5_130-v0017}

\margin{130}

%% END SOURCE ps_134-5_130
\ifnum\EditionMode=1
\psapparatusstart
\psentry{ps_134-5_130-v0001}{1}{typed insertion}{\inserted[typed]{in}}{in}{The typed insertion is incorporated at its indicated position; insertion brackets leave the running text.}
\psentry{ps_134-5_130-v0002}{2}{typewriter overwriting}{\over[typed]{i}{o}}{o}{The superposed typed reading replaces the underlying reading.}
\psentry{ps_134-5_130-v0003}{3}{typewriter deletion}{\xout{we}}{}{The typewriter strikes out this material; it leaves the corrected running text.}
\psentry{ps_134-5_130-v0004}{4}{typed insertion}{\inserted[typed]{K.}}{K.}{The typed insertion is incorporated at its indicated position; insertion brackets leave the running text.}
\psentry{ps_134-5_130-v0005}{5}{typewriter deletion}{\xout{stable under sums}}{}{The typewriter strikes out this material; it leaves the corrected running text.}
\psentry{ps_134-5_130-v0006}{6}{typewriter deletion}{\xout{we}}{}{The typewriter strikes out this material; it leaves the corrected running text.}
\psentry{ps_134-5_130-v0007}{7}{placement}{T\(\underset{\hand{h}}{\text{h}}\)e}{The}{A pen h lies beneath the typed h; the reading gives The once.}
\psentry{ps_134-5_130-v0008}{8}{typewriter deletion}{\xout{some}}{}{The typewriter strikes out this material; it leaves the corrected running text.}
\psentry{ps_134-5_130-v0009}{9}{typewriter deletion}{\xout{(if needed\illegible{1}}}{}{The typewriter strikes out this material; it leaves the corrected running text.}
\psentry{ps_134-5_130-v0010}{10}{typewriter deletion}{\xout{\illegible{1}}}{}{The typewriter strikes out this material; it leaves the corrected running text.}
\psentry{ps_134-5_130-v0011}{11}{typewriter overwriting}{\over[typed]{we}{one}}{one}{The superposed typed reading replaces the underlying reading.}
\psentry{ps_134-5_130-v0012}{12}{placement}{\hand{\tikz[x=.5em,y=.6em,baseline=.2ex]{\draw (-1,0)--(0,0)--(0,1)--(1,1); \draw (.12,0)--(.12,1);}}}{}{The angular pen paragraph divider is preserved diplomatically; the following paragraph break already represents its separation in the reading.}
\psentry{ps_134-5_130-v0013}{13}{typewriter deletion}{\xout{be a}}{}{The typewriter strikes out this material; it leaves the corrected running text.}
\psentry{ps_134-5_130-v0014}{14}{insertion}{\inserted{am}}{\hand{am}}{The handwritten insertion is incorporated at its indicated position; insertion brackets leave the running text.}
\psentry{ps_134-5_130-v0015}{15}{deletion}{\deleted{did}}{}{The author cancels this material by hand; it leaves the corrected running text.}
\psentry{ps_134-5_130-v0016}{16}{insertion}{\inserted{doing}}{\hand{doing}}{The handwritten insertion is incorporated at its indicated position; insertion brackets leave the running text.}
\psentry{ps_134-5_130-v0017}{17}{deletion}{\deleted{But there \deleted{are}\inserted{should be} obvious ways to satisfy the\xout{se conditions} preliminary con-%
ditions, it occurs to me, in such a way that A be a \xout{test} str\over[typed]{a}{i}ct test
category and the composition
\[ i: A \qquad (\mathrm{Ord}) \qquad (\mathrm{Cat}) \]
a test functor (maybe not a perfect one). To see this, let $A_o$ be
any \inserted{small full} subcategory of the category of finite ordered sets with a largest
element and satisfying the assumption (2) for i(a), such that $A_o$}}{}{The author cancels this material by hand; it leaves the corrected running text.}
\psapparatusend
\fi
\pssource{ps_134-5_131}{134-5/131}{https://grothendieck.umontpellier.fr/web/viewer.html?file=../134-5.pdf\#page=132}
\psnumber{ps_134-5_131-v0001}{1}
\psnumber{ps_134-5_131-v0002}{2}
\psnumber{ps_134-5_131-v0003}{3}
\psnumber{ps_134-5_131-v0004}{4}
\psnumber{ps_134-5_131-v0005}{5}
\psnumber{ps_134-5_131-v0006}{6}
\psnumber{ps_134-5_131-v0007}{7}
\psnumber{ps_134-5_131-v0008}{8}
%% BEGIN SOURCE ps_134-5_131
% Pursuing Stacks (cotes 134-2 to 134-8): archive cote 134-5, site page 131 (PDF page 132), typescript page 318.
% typescript. Draft compared: ps5.tex:1279-1339.
% AI-transcribed by Claude Opus 5.5 (claude-opus-5-5, adaptive thinking, effort high), Message Batches API, batch ps3 (msgbatch_01XdFmuVVBAcgc7EFxdtCqsj), 2026-11-12; reviewed against both supplied scan halves; see CHECK_DECISIONS.tsv.
% Hand-drawn arrows are marked blue; typed objects remain plain.
% The hemispherical signs have double oval outlines; bigcirc is the supplied glyph approximation.
\scanpage{134-5/131}
\mspage{318}

I just spent a while trying to find some convincing example,  by
using a suitable full subcategory $A_o$ of (Ord),  made up with finite sets
satisfying the assumption (2) above for i(a),  under some additional
assumption on $A_o$  such as stability under finite products and under \var{\xout{the}}{}{typewriter deletion}{ps_134-5_131-v0001}
\var{\xout{operation}}{}{typewriter deletion}{ps_134-5_131-v0002} passage from a  to an object $a_{/x}$ ,  and that $A_o$  contain\var{\deleted{s}}{}{deletion}{ps_134-5_131-v0003}
the ordered set
\[ I = \diagram{\begin{tikzcd}[row sep=small,column sep=small] \bullet \ar[dr,color=blue] \& \\ \& \bullet \\ \bullet \ar[ur,color=blue] \& \end{tikzcd}} \quad , \]

whose geometric realization is the segment $B^1$  with it's usual cellu\var{\over{f}{l}}{\hand{l}}{writing over}{ps_134-5_131-v0004}ar
decomposition.  In terms of $A_o$ and introducing \hand{``}orientations\hand{''} of objects
of $A_o$,  the idea was to define an\var{\over[typed]{\illegible{1}}{o}}{o}{typewriter overwriting}{ps_134-5_131-v0005}ther category A (of pairs $(a,\hand{\omega})$,
with  a in $A_o$ and $\hand{\omega}$ an orientation of $\hand{|}a\hand{|}$),  \hand{\uline{\textcolor{black}{stable under finite pro-}}}%
\hand{\uline{\textcolor{black}{ducts}}} so that $A^{\wedge}$ is totally aspheric,  together with a functor
$i\colon A \hand{\longrightarrow} (Ord) \hand{\hookrightarrow} (Cat)$  such that $i^{*}(I)$ should be representable,  and
hence furnish the homotopy interval needed to insure that A is a test
category.  T\var{\over{\unsure{n}}{h}}{\hand{h}}{writing over}{ps_134-5_131-v0006}e first idea that comes to mind,  namely define a map
from $(a,\hand{\omega})$  to $(a',\hand{\omega}')$  as merely a map from a to a' in $A_o$ ,  is
nonsense unfortunately,  as in the data (6),  the orientations will not
be stable under isomorphisms,  a condition which I forgot to state
before,  and which is visibly needed in order to be able to define the
differential between the $K_n$'s .  If we try define A taking into account
this compatibility condition,  we loose existence of products,  anyhow
$i^{*}(I)$ is'nt representable any more,  so why should it be aspheric over
the final object, \var{\deleted{\unsure{6|\&}}}{}{deletion}{ps_134-5_131-v0007}  so why should the functor i be a test functor ?

\var{T\(\underset{\hand{h}}{\text{h}}\)e difficulty}{The difficulty}{placement}{ps_134-5_131-v0008} I find in carrying through any explicit example
for a ``test category with boundary operations'',  except those which stem
from our trinity,  is rather intriguing I feel.  The question is whether
maybe in this direction,  one might get at an intrinsic description of
the trinity,  in terms of the rather natural structure species of a
``test category with boundary operations''.  This is the second instance
where the thought arises that the three standard test categories
$\hand{\Delta}$, $\hand{\square}$ and $\hand{\bigcirc}$ may be distinguished in some respects  -  the first instan\unsure{e}
ce was in relation to the Dold-Puppe theorem.

9)  \hand{\uline{\textcolor{black}{Miscellaneous residual questions from part}}} II.  One of these
was about the category $\hand{\Delta}^f$ of  simplices without degeneracies being a
weak test category (cf. section 43)  -  while it is definitely \hand{\uline{\textcolor{black}{not}}} a test
category.  It seems worth while to write down a proof for this,  maybe
too for the analoguous statements for $\hand{\square}^f$  and $\hand{\bigcirc}^f$ .  This reminds me

\margin{131}

%% END SOURCE ps_134-5_131
\ifnum\EditionMode=1
\psapparatusstart
\psentry{ps_134-5_131-v0001}{1}{typewriter deletion}{\xout{the}}{}{The typewriter strikes out this material; it leaves the corrected running text.}
\psentry{ps_134-5_131-v0002}{2}{typewriter deletion}{\xout{operation}}{}{The typewriter strikes out this material; it leaves the corrected running text.}
\psentry{ps_134-5_131-v0003}{3}{deletion}{\deleted{s}}{}{The author cancels this material by hand; it leaves the corrected running text.}
\psentry{ps_134-5_131-v0004}{4}{writing over}{\over{f}{l}}{\hand{l}}{The superposed handwritten reading replaces the underlying reading.}
\psentry{ps_134-5_131-v0005}{5}{typewriter overwriting}{\over[typed]{\illegible{1}}{o}}{o}{The superposed typed reading replaces the underlying reading.}
\psentry{ps_134-5_131-v0006}{6}{writing over}{\over{\unsure{n}}{h}}{\hand{h}}{The superposed handwritten reading replaces the underlying reading.}
\psentry{ps_134-5_131-v0007}{7}{deletion}{\deleted{\unsure{6|\&}}}{}{The author cancels this material by hand; it leaves the corrected running text.}
\psentry{ps_134-5_131-v0008}{8}{placement}{T\(\underset{\hand{h}}{\text{h}}\)e difficulty}{The difficulty}{The second The has a pen h below its typed h; preserve both positions diplomatically and read the word once.}
\psapparatusend
\fi
\pssource{ps_134-5_132}{134-5/132}{https://grothendieck.umontpellier.fr/web/viewer.html?file=../134-5.pdf\#page=133}
\psnumber{ps_134-5_132-v0001}{1}
\psnumber{ps_134-5_132-v0002}{2}
\psnumber{ps_134-5_132-v0003}{3}
\psnumber{ps_134-5_132-v0004}{4}
\psnumber{ps_134-5_132-v0005}{5}
\psnumber{ps_134-5_132-v0006}{6}
\psnumber{ps_134-5_132-v0007}{7}
\psnumber{ps_134-5_132-v0008}{8}
\psnumber{ps_134-5_132-v0009}{9}
\psnumber{ps_134-5_132-v0010}{10}
\psnumber{ps_134-5_132-v0011}{11}
\psnumber{ps_134-5_132-v0012}{12}
\psnumber{ps_134-5_132-v0013}{13}
\psnumber{ps_134-5_132-v0014}{14}
\psnumber{ps_134-5_132-v0015}{15}
%% BEGIN SOURCE ps_134-5_132
% Pursuing Stacks (cotes 134-2 to 134-8): archive cote 134-5, site page 132 (PDF page 133), typescript page 319.
% typescript. Draft compared: ps5.tex:1327-1377.
% AI-transcribed by Claude Opus 5.5 (claude-opus-5-5, adaptive thinking, effort high), Message Batches API, batch ps3 (msgbatch_01XdFmuVVBAcgc7EFxdtCqsj), 2026-11-12; reviewed against both supplied scan halves; see CHECK_DECISIONS.tsv.
% The hemispherical sign has a double oval outline; bigcirc is the supplied glyph approximation.
\scanpage{134-5/132}
\mspage{319}

too that I never got around introducing formall\var{\over{y}{y}}{\hand{y}}{writing over}{ps_134-5_132-v0001} the hemispherical test
category, which presumably will be \var{\over[typed]{b}{v}}{v}{typewriter overwriting}{ps_134-5_132-v0002}\unsure{a|e}ry useful when it comes to study-%
ing stacks - this too could be done in part IV, as well as proving
of course that \hand{$\bigcirc$} is a strict test category indeed, or better still,
a contractor. It may be fun too constructing test functors from any
one of the three basic test categories in the trinity, to the category
of complexes defined by the two others - six cases altogether to con-%
sider! But as I am not in the process of writing the ``Elements d'Algè-%
bre Homotopique'' , maybe I will skip this !

In the same section 43, I raised the question as to whether the
ordered set of all non-empty finite subsets of a given infinite set,
viewed as \var{\over[typed]{g}{a}}{a}{typewriter overwriting}{ps_134-5_132-v0003} category in the usual way \var{\over{f}{f}}{\hand{f}}{writing over}{ps_134-5_132-v0004}or an ordered set, was a weak
test category (on page 78 it was seen not to be a test category, and
it is immediate then that it is not totally aspheric either). \var{O\(\underset{\text{n}}{\text{n}}\)e}{One}{placement}{ps_134-5_132-v0005} in-%
teresting application, as noticed there, would be to the effect that
(Ord), the subcategory of (Cat) defined by ordered sets, is a modeli\var{\deleted{z\unsure{a}}}{}{deletion}{ps_134-5_132-v0006}%
zer (for the induced notion of weak equivalences). \var{\xout{Here}}{}{typewriter deletion}{ps_134-5_132-v0007}\var{\inserted[typed]{Now}}{Now}{typed insertion}{ps_134-5_132-v0008} the question
arises moreover whether this modelizing structure comes from an aspheri-%
city structure, or even from a contractibility structure - and the
same questio\var{\over[typed]{j}{n}}{n}{typewriter overwriting}{ps_134-5_132-v0009} arises in the more general situation described in the
proposition of page 74.

A last question \var{\over[typed]{I}{a}}{a}{typewriter overwriting}{ps_134-5_132-v0010}long these lines I would like to clear up,
is the relation of total \var{\inserted[typed]{$\underline{W}$-}}{$\underline{W}$-}{typed insertion}{ps_134-5_132-v0011}asphericity for an asphericity structure,
for variable $\underline{W}$ \var{\(\underset{\hand{\smile}}{\text{,}}\)}{,}{placement}{ps_134-5_132-v0012} when $\underline{W} \hand{\subset} \underline{W}'$ . Assuming the localizers satisfy
the condition \var{L\(\underset{\hand{\illegible{1}}}{\text{o}}\)c}{Loc}{placement}{ps_134-5_132-v0013} 4 ), is it true that \var{\inserted[typed]{total}}{total}{typed insertion}{ps_134-5_132-v0014} $\underline{W}$-asphericity is equivalent
to total $\underline{W}'$-asphericity - \var{\xout{and}}{}{typewriter deletion}{ps_134-5_132-v0015} or equivalently, is it equivalent
to total 0-connectedness ?

\margin{132}

%% END SOURCE ps_134-5_132
\ifnum\EditionMode=1
\psapparatusstart
\psentry{ps_134-5_132-v0001}{1}{writing over}{\over{y}{y}}{\hand{y}}{The superposed handwritten reading replaces the underlying reading.}
\psentry{ps_134-5_132-v0002}{2}{typewriter overwriting}{\over[typed]{b}{v}}{v}{The superposed typed reading replaces the underlying reading.}
\psentry{ps_134-5_132-v0003}{3}{typewriter overwriting}{\over[typed]{g}{a}}{a}{The superposed typed reading replaces the underlying reading.}
\psentry{ps_134-5_132-v0004}{4}{writing over}{\over{f}{f}}{\hand{f}}{The superposed handwritten reading replaces the underlying reading.}
\psentry{ps_134-5_132-v0005}{5}{placement}{O\(\underset{\text{n}}{\text{n}}\)e}{One}{A separate typed n impression lies beneath the baseline n of One; read the word once.}
\psentry{ps_134-5_132-v0006}{6}{deletion}{\deleted{z\unsure{a}}}{}{The author cancels this material by hand; it leaves the corrected running text.}
\psentry{ps_134-5_132-v0007}{7}{typewriter deletion}{\xout{Here}}{}{The typewriter strikes out this material; it leaves the corrected running text.}
\psentry{ps_134-5_132-v0008}{8}{typed insertion}{\inserted[typed]{Now}}{Now}{The typed insertion is incorporated at its indicated position; insertion brackets leave the running text.}
\psentry{ps_134-5_132-v0009}{9}{typewriter overwriting}{\over[typed]{j}{n}}{n}{The superposed typed reading replaces the underlying reading.}
\psentry{ps_134-5_132-v0010}{10}{typewriter overwriting}{\over[typed]{I}{a}}{a}{The superposed typed reading replaces the underlying reading.}
\psentry{ps_134-5_132-v0011}{11}{typed insertion}{\inserted[typed]{$\underline{W}$-}}{$\underline{W}$-}{The typed insertion is incorporated at its indicated position; insertion brackets leave the running text.}
\psentry{ps_134-5_132-v0012}{12}{placement}{\(\underset{\hand{\smile}}{\text{,}}\)}{,}{A shallow pen arc lies below the typed comma; the reading preserves the comma without the nonlexical mark.}
\psentry{ps_134-5_132-v0013}{13}{placement}{L\(\underset{\hand{\illegible{1}}}{\text{o}}\)c}{Loc}{An unreadable pen mark lies below the typed o; the reading preserves Loc, while the diplomatic source records the unresolved mark.}
\psentry{ps_134-5_132-v0014}{14}{typed insertion}{\inserted[typed]{total}}{total}{The typed insertion is incorporated at its indicated position; insertion brackets leave the running text.}
\psentry{ps_134-5_132-v0015}{15}{typewriter deletion}{\xout{and}}{}{The typewriter strikes out this material; it leaves the corrected running text.}
\psapparatusend
\fi

\end{document}
